
\ifdefined\rhabNClassesBeatAddRm\else% AUTO-GENERATED by writeup/make_rehaul_ablations.py -- do not edit; re-run the script.
% results/ablations_2014start_w3000 (RULE.txt): column-subset arms of the final 13-bar design,
% frozen hyperparameters, every bar scored as Table 1; DM on daily (average) QLIKE.
\newcommand{\rhabEnetAddEwAllQlike}{0.1871}  % QLIKE enet add_ew all
\newcommand{\rhabEnetAddEwAllMse}{1.47}  % MSE x 1e10 enet add_ew all
\newcommand{\rhabEnetAddEwAllPct}{$-$1.1}  % QLIKE change % vs har
\newcommand{\rhabEnetAddEwAllDm}{$-$2.81}  % DM vs the reference arm
\newcommand{\rhabEnetAddEwAllP}{0.005}  % p-value
\newcommand{\rhabEnetAddEwCloseQlike}{0.1249}  % QLIKE enet add_ew c16
\newcommand{\rhabEnetAddEwCloseMse}{2.10}  % MSE x 1e10 enet add_ew c16
\newcommand{\rhabEnetAddEwClosePct}{$-$1.5}  % QLIKE change % vs har
\newcommand{\rhabEnetAddEwCloseDm}{$-$1.37}  % DM vs the reference arm
\newcommand{\rhabEnetAddEwCloseP}{0.170}  % p-value
\newcommand{\rhabEnetAddEwCloseLateQlike}{0.1120}  % QLIKE enet add_ew c16_2020
\newcommand{\rhabEnetAddEwCloseLateMse}{4.56}  % MSE x 1e10 enet add_ew c16_2020
\newcommand{\rhabEnetAddEwCloseLatePct}{$-$1.5}  % QLIKE change % vs har
\newcommand{\rhabEnetAddEwCloseLateDm}{$-$0.90}  % DM vs the reference arm
\newcommand{\rhabEnetAddEwCloseLateP}{0.369}  % p-value
\newcommand{\rhabEnetAddFomcAllQlike}{0.1847}  % QLIKE enet add_fomc all
\newcommand{\rhabEnetAddFomcAllMse}{1.13}  % MSE x 1e10 enet add_fomc all
\newcommand{\rhabEnetAddFomcAllPct}{$-$2.4}  % QLIKE change % vs har
\newcommand{\rhabEnetAddFomcAllDm}{$-$5.85}  % DM vs the reference arm
\newcommand{\rhabEnetAddFomcAllP}{\ensuremath{<}0.001}  % p-value
\newcommand{\rhabEnetAddFomcCloseQlike}{0.1274}  % QLIKE enet add_fomc c16
\newcommand{\rhabEnetAddFomcCloseMse}{1.85}  % MSE x 1e10 enet add_fomc c16
\newcommand{\rhabEnetAddFomcClosePct}{0.4}  % QLIKE change % vs har
\newcommand{\rhabEnetAddFomcCloseDm}{1.66}  % DM vs the reference arm
\newcommand{\rhabEnetAddFomcCloseP}{0.097}  % p-value
\newcommand{\rhabEnetAddFomcCloseLateQlike}{0.1136}  % QLIKE enet add_fomc c16_2020
\newcommand{\rhabEnetAddFomcCloseLateMse}{3.96}  % MSE x 1e10 enet add_fomc c16_2020
\newcommand{\rhabEnetAddFomcCloseLatePct}{0.0}  % QLIKE change % vs har
\newcommand{\rhabEnetAddFomcCloseLateDm}{$-$0.09}  % DM vs the reference arm
\newcommand{\rhabEnetAddFomcCloseLateP}{0.927}  % p-value
\newcommand{\rhabEnetAddIvAllQlike}{0.1894}  % QLIKE enet add_iv all
\newcommand{\rhabEnetAddIvAllMse}{1.32}  % MSE x 1e10 enet add_iv all
\newcommand{\rhabEnetAddIvAllPct}{0.1}  % QLIKE change % vs har
\newcommand{\rhabEnetAddIvAllDm}{0.23}  % DM vs the reference arm
\newcommand{\rhabEnetAddIvAllP}{0.817}  % p-value
\newcommand{\rhabEnetAddIvCloseQlike}{0.1272}  % QLIKE enet add_iv c16
\newcommand{\rhabEnetAddIvCloseMse}{1.76}  % MSE x 1e10 enet add_iv c16
\newcommand{\rhabEnetAddIvClosePct}{0.3}  % QLIKE change % vs har
\newcommand{\rhabEnetAddIvCloseDm}{0.28}  % DM vs the reference arm
\newcommand{\rhabEnetAddIvCloseP}{0.780}  % p-value
\newcommand{\rhabEnetAddIvCloseLateQlike}{0.1152}  % QLIKE enet add_iv c16_2020
\newcommand{\rhabEnetAddIvCloseLateMse}{3.75}  % MSE x 1e10 enet add_iv c16_2020
\newcommand{\rhabEnetAddIvCloseLatePct}{1.4}  % QLIKE change % vs har
\newcommand{\rhabEnetAddIvCloseLateDm}{0.70}  % DM vs the reference arm
\newcommand{\rhabEnetAddIvCloseLateP}{0.485}  % p-value
\newcommand{\rhabEnetAddLiqAllQlike}{0.1881}  % QLIKE enet add_liq all
\newcommand{\rhabEnetAddLiqAllMse}{1.16}  % MSE x 1e10 enet add_liq all
\newcommand{\rhabEnetAddLiqAllPct}{$-$0.6}  % QLIKE change % vs har
\newcommand{\rhabEnetAddLiqAllDm}{$-$1.54}  % DM vs the reference arm
\newcommand{\rhabEnetAddLiqAllP}{0.124}  % p-value
\newcommand{\rhabEnetAddLiqCloseQlike}{0.1299}  % QLIKE enet add_liq c16
\newcommand{\rhabEnetAddLiqCloseMse}{1.99}  % MSE x 1e10 enet add_liq c16
\newcommand{\rhabEnetAddLiqClosePct}{2.4}  % QLIKE change % vs har
\newcommand{\rhabEnetAddLiqCloseDm}{1.77}  % DM vs the reference arm
\newcommand{\rhabEnetAddLiqCloseP}{0.077}  % p-value
\newcommand{\rhabEnetAddLiqCloseLateQlike}{0.1192}  % QLIKE enet add_liq c16_2020
\newcommand{\rhabEnetAddLiqCloseLateMse}{4.29}  % MSE x 1e10 enet add_liq c16_2020
\newcommand{\rhabEnetAddLiqCloseLatePct}{4.9}  % QLIKE change % vs har
\newcommand{\rhabEnetAddLiqCloseLateDm}{1.90}  % DM vs the reference arm
\newcommand{\rhabEnetAddLiqCloseLateP}{0.058}  % p-value
\newcommand{\rhabEnetAddRmAllQlike}{0.1842}  % QLIKE enet add_rm all
\newcommand{\rhabEnetAddRmAllMse}{1.72}  % MSE x 1e10 enet add_rm all
\newcommand{\rhabEnetAddRmAllPct}{$-$2.7}  % QLIKE change % vs har
\newcommand{\rhabEnetAddRmAllDm}{$-$4.25}  % DM vs the reference arm
\newcommand{\rhabEnetAddRmAllP}{\ensuremath{<}0.001}  % p-value
\newcommand{\rhabEnetAddRmCloseQlike}{0.1223}  % QLIKE enet add_rm c16
\newcommand{\rhabEnetAddRmCloseMse}{2.92}  % MSE x 1e10 enet add_rm c16
\newcommand{\rhabEnetAddRmClosePct}{$-$3.6}  % QLIKE change % vs har
\newcommand{\rhabEnetAddRmCloseDm}{$-$2.25}  % DM vs the reference arm
\newcommand{\rhabEnetAddRmCloseP}{0.024}  % p-value
\newcommand{\rhabEnetAddRmCloseLateQlike}{0.1147}  % QLIKE enet add_rm c16_2020
\newcommand{\rhabEnetAddRmCloseLateMse}{6.56}  % MSE x 1e10 enet add_rm c16_2020
\newcommand{\rhabEnetAddRmCloseLatePct}{1.0}  % QLIKE change % vs har
\newcommand{\rhabEnetAddRmCloseLateDm}{0.33}  % DM vs the reference arm
\newcommand{\rhabEnetAddRmCloseLateP}{0.741}  % p-value
\newcommand{\rhabEnetAddSentAllQlike}{0.1893}  % QLIKE enet add_sent all
\newcommand{\rhabEnetAddSentAllMse}{1.11}  % MSE x 1e10 enet add_sent all
\newcommand{\rhabEnetAddSentAllPct}{0.0}  % QLIKE change % vs har
\newcommand{\rhabEnetAddSentAllDm}{0.04}  % DM vs the reference arm
\newcommand{\rhabEnetAddSentAllP}{0.971}  % p-value
\newcommand{\rhabEnetAddSentCloseQlike}{0.1300}  % QLIKE enet add_sent c16
\newcommand{\rhabEnetAddSentCloseMse}{1.98}  % MSE x 1e10 enet add_sent c16
\newcommand{\rhabEnetAddSentClosePct}{2.5}  % QLIKE change % vs har
\newcommand{\rhabEnetAddSentCloseDm}{3.62}  % DM vs the reference arm
\newcommand{\rhabEnetAddSentCloseP}{\ensuremath{<}0.001}  % p-value
\newcommand{\rhabEnetAddSentCloseLateQlike}{0.1194}  % QLIKE enet add_sent c16_2020
\newcommand{\rhabEnetAddSentCloseLateMse}{4.27}  % MSE x 1e10 enet add_sent c16_2020
\newcommand{\rhabEnetAddSentCloseLatePct}{5.0}  % QLIKE change % vs har
\newcommand{\rhabEnetAddSentCloseLateDm}{2.88}  % DM vs the reference arm
\newcommand{\rhabEnetAddSentCloseLateP}{0.004}  % p-value
\newcommand{\rhabEnetAddVdAllQlike}{0.1900}  % QLIKE enet add_vd all
\newcommand{\rhabEnetAddVdAllMse}{1.11}  % MSE x 1e10 enet add_vd all
\newcommand{\rhabEnetAddVdAllPct}{0.4}  % QLIKE change % vs har
\newcommand{\rhabEnetAddVdAllDm}{2.92}  % DM vs the reference arm
\newcommand{\rhabEnetAddVdAllP}{0.003}  % p-value
\newcommand{\rhabEnetAddVdCloseQlike}{0.1277}  % QLIKE enet add_vd c16
\newcommand{\rhabEnetAddVdCloseMse}{1.96}  % MSE x 1e10 enet add_vd c16
\newcommand{\rhabEnetAddVdClosePct}{0.6}  % QLIKE change % vs har
\newcommand{\rhabEnetAddVdCloseDm}{1.85}  % DM vs the reference arm
\newcommand{\rhabEnetAddVdCloseP}{0.065}  % p-value
\newcommand{\rhabEnetAddVdCloseLateQlike}{0.1160}  % QLIKE enet add_vd c16_2020
\newcommand{\rhabEnetAddVdCloseLateMse}{4.21}  % MSE x 1e10 enet add_vd c16_2020
\newcommand{\rhabEnetAddVdCloseLatePct}{2.1}  % QLIKE change % vs har
\newcommand{\rhabEnetAddVdCloseLateDm}{2.57}  % DM vs the reference arm
\newcommand{\rhabEnetAddVdCloseLateP}{0.010}  % p-value
\newcommand{\rhabEnetAddVwAllQlike}{0.1872}  % QLIKE enet add_vw all
\newcommand{\rhabEnetAddVwAllMse}{1.99}  % MSE x 1e10 enet add_vw all
\newcommand{\rhabEnetAddVwAllPct}{$-$1.1}  % QLIKE change % vs har
\newcommand{\rhabEnetAddVwAllDm}{$-$2.03}  % DM vs the reference arm
\newcommand{\rhabEnetAddVwAllP}{0.042}  % p-value
\newcommand{\rhabEnetAddVwCloseQlike}{0.1256}  % QLIKE enet add_vw c16
\newcommand{\rhabEnetAddVwCloseMse}{2.37}  % MSE x 1e10 enet add_vw c16
\newcommand{\rhabEnetAddVwClosePct}{$-$1.0}  % QLIKE change % vs har
\newcommand{\rhabEnetAddVwCloseDm}{$-$0.75}  % DM vs the reference arm
\newcommand{\rhabEnetAddVwCloseP}{0.454}  % p-value
\newcommand{\rhabEnetAddVwCloseLateQlike}{0.1134}  % QLIKE enet add_vw c16_2020
\newcommand{\rhabEnetAddVwCloseLateMse}{5.22}  % MSE x 1e10 enet add_vw c16_2020
\newcommand{\rhabEnetAddVwCloseLatePct}{$-$0.2}  % QLIKE change % vs har
\newcommand{\rhabEnetAddVwCloseLateDm}{$-$0.13}  % DM vs the reference arm
\newcommand{\rhabEnetAddVwCloseLateP}{0.897}  % p-value
\newcommand{\rhabEnetFullAllQlike}{0.1778}  % QLIKE enet full all
\newcommand{\rhabEnetFullAllMse}{2.43}  % MSE x 1e10 enet full all
\newcommand{\rhabEnetFullAllPct}{$-$6.0}  % QLIKE change % vs har
\newcommand{\rhabEnetFullAllDm}{$-$6.98}  % DM vs the reference arm
\newcommand{\rhabEnetFullAllP}{\ensuremath{<}0.001}  % p-value
\newcommand{\rhabEnetFullCloseQlike}{0.1262}  % QLIKE enet full c16
\newcommand{\rhabEnetFullCloseMse}{3.94}  % MSE x 1e10 enet full c16
\newcommand{\rhabEnetFullClosePct}{$-$0.5}  % QLIKE change % vs har
\newcommand{\rhabEnetFullCloseDm}{$-$0.21}  % DM vs the reference arm
\newcommand{\rhabEnetFullCloseP}{0.832}  % p-value
\newcommand{\rhabEnetFullCloseLateQlike}{0.1234}  % QLIKE enet full c16_2020
\newcommand{\rhabEnetFullCloseLateMse}{9.01}  % MSE x 1e10 enet full c16_2020
\newcommand{\rhabEnetFullCloseLatePct}{8.5}  % QLIKE change % vs har
\newcommand{\rhabEnetFullCloseLateDm}{1.91}  % DM vs the reference arm
\newcommand{\rhabEnetFullCloseLateP}{0.056}  % p-value
\newcommand{\rhabEnetHarAllQlike}{0.1893}  % QLIKE enet har all
\newcommand{\rhabEnetHarAllMse}{1.11}  % MSE x 1e10 enet har all
\newcommand{\rhabEnetHarCloseQlike}{0.1269}  % QLIKE enet har c16
\newcommand{\rhabEnetHarCloseMse}{1.97}  % MSE x 1e10 enet har c16
\newcommand{\rhabEnetHarCloseLateQlike}{0.1137}  % QLIKE enet har c16_2020
\newcommand{\rhabEnetHarCloseLateMse}{4.24}  % MSE x 1e10 enet har c16_2020
\newcommand{\rhabEnetHarqAllQlike}{0.1893}  % QLIKE enet harq all
\newcommand{\rhabEnetHarqAllMse}{1.16}  % MSE x 1e10 enet harq all
\newcommand{\rhabEnetHarqAllPct}{0.0}  % QLIKE change % vs har
\newcommand{\rhabEnetHarqAllDm}{0.63}  % DM vs the reference arm
\newcommand{\rhabEnetHarqAllP}{0.531}  % p-value
\newcommand{\rhabEnetHarqCloseQlike}{0.1267}  % QLIKE enet harq c16
\newcommand{\rhabEnetHarqCloseMse}{1.80}  % MSE x 1e10 enet harq c16
\newcommand{\rhabEnetHarqClosePct}{$-$0.1}  % QLIKE change % vs har
\newcommand{\rhabEnetHarqCloseDm}{$-$1.31}  % DM vs the reference arm
\newcommand{\rhabEnetHarqCloseP}{0.191}  % p-value
\newcommand{\rhabEnetHarqCloseLateQlike}{0.1134}  % QLIKE enet harq c16_2020
\newcommand{\rhabEnetHarqCloseLateMse}{3.84}  % MSE x 1e10 enet harq c16_2020
\newcommand{\rhabEnetHarqCloseLatePct}{$-$0.3}  % QLIKE change % vs har
\newcommand{\rhabEnetHarqCloseLateDm}{$-$1.15}  % DM vs the reference arm
\newcommand{\rhabEnetHarqCloseLateP}{0.251}  % p-value
\newcommand{\rhabEnetLharAllQlike}{0.1841}  % QLIKE enet lhar all
\newcommand{\rhabEnetLharAllMse}{1.71}  % MSE x 1e10 enet lhar all
\newcommand{\rhabEnetLharAllPct}{$-$2.7}  % QLIKE change % vs har
\newcommand{\rhabEnetLharAllDm}{$-$6.66}  % DM vs the reference arm
\newcommand{\rhabEnetLharAllP}{\ensuremath{<}0.001}  % p-value
\newcommand{\rhabEnetLharCloseQlike}{0.1217}  % QLIKE enet lhar c16
\newcommand{\rhabEnetLharCloseMse}{2.76}  % MSE x 1e10 enet lhar c16
\newcommand{\rhabEnetLharClosePct}{$-$4.1}  % QLIKE change % vs har
\newcommand{\rhabEnetLharCloseDm}{$-$3.16}  % DM vs the reference arm
\newcommand{\rhabEnetLharCloseP}{0.002}  % p-value
\newcommand{\rhabEnetLharCloseLateQlike}{0.1124}  % QLIKE enet lhar c16_2020
\newcommand{\rhabEnetLharCloseLateMse}{6.19}  % MSE x 1e10 enet lhar c16_2020
\newcommand{\rhabEnetLharCloseLatePct}{$-$1.1}  % QLIKE change % vs har
\newcommand{\rhabEnetLharCloseLateDm}{$-$0.45}  % DM vs the reference arm
\newcommand{\rhabEnetLharCloseLateP}{0.656}  % p-value
\newcommand{\rhabEnetLoEwAllQlike}{0.1776}  % QLIKE enet lo_ew all
\newcommand{\rhabEnetLoEwAllMse}{2.54}  % MSE x 1e10 enet lo_ew all
\newcommand{\rhabEnetLoEwAllPct}{$-$0.1}  % QLIKE change % vs full
\newcommand{\rhabEnetLoEwAllDm}{$-$1.41}  % DM vs the reference arm
\newcommand{\rhabEnetLoEwAllP}{0.157}  % p-value
\newcommand{\rhabEnetLoEwCloseQlike}{0.1267}  % QLIKE enet lo_ew c16
\newcommand{\rhabEnetLoEwCloseMse}{4.02}  % MSE x 1e10 enet lo_ew c16
\newcommand{\rhabEnetLoEwClosePct}{0.3}  % QLIKE change % vs full
\newcommand{\rhabEnetLoEwCloseDm}{1.18}  % DM vs the reference arm
\newcommand{\rhabEnetLoEwCloseP}{0.238}  % p-value
\newcommand{\rhabEnetLoEwCloseLateQlike}{0.1235}  % QLIKE enet lo_ew c16_2020
\newcommand{\rhabEnetLoEwCloseLateMse}{9.21}  % MSE x 1e10 enet lo_ew c16_2020
\newcommand{\rhabEnetLoEwCloseLatePct}{0.1}  % QLIKE change % vs full
\newcommand{\rhabEnetLoEwCloseLateDm}{0.23}  % DM vs the reference arm
\newcommand{\rhabEnetLoEwCloseLateP}{0.822}  % p-value
\newcommand{\rhabEnetLoFomcAllQlike}{0.1820}  % QLIKE enet lo_fomc all
\newcommand{\rhabEnetLoFomcAllMse}{2.42}  % MSE x 1e10 enet lo_fomc all
\newcommand{\rhabEnetLoFomcAllPct}{2.3}  % QLIKE change % vs full
\newcommand{\rhabEnetLoFomcAllDm}{5.74}  % DM vs the reference arm
\newcommand{\rhabEnetLoFomcAllP}{\ensuremath{<}0.001}  % p-value
\newcommand{\rhabEnetLoFomcCloseQlike}{0.1261}  % QLIKE enet lo_fomc c16
\newcommand{\rhabEnetLoFomcCloseMse}{3.99}  % MSE x 1e10 enet lo_fomc c16
\newcommand{\rhabEnetLoFomcClosePct}{$-$0.1}  % QLIKE change % vs full
\newcommand{\rhabEnetLoFomcCloseDm}{$-$0.51}  % DM vs the reference arm
\newcommand{\rhabEnetLoFomcCloseP}{0.607}  % p-value
\newcommand{\rhabEnetLoFomcCloseLateQlike}{0.1234}  % QLIKE enet lo_fomc c16_2020
\newcommand{\rhabEnetLoFomcCloseLateMse}{9.13}  % MSE x 1e10 enet lo_fomc c16_2020
\newcommand{\rhabEnetLoFomcCloseLatePct}{0.1}  % QLIKE change % vs full
\newcommand{\rhabEnetLoFomcCloseLateDm}{0.13}  % DM vs the reference arm
\newcommand{\rhabEnetLoFomcCloseLateP}{0.894}  % p-value
\newcommand{\rhabEnetLoIvAllQlike}{0.1778}  % QLIKE enet lo_iv all
\newcommand{\rhabEnetLoIvAllMse}{2.49}  % MSE x 1e10 enet lo_iv all
\newcommand{\rhabEnetLoIvAllPct}{0.0}  % QLIKE change % vs full
\newcommand{\rhabEnetLoIvAllDm}{$-$0.13}  % DM vs the reference arm
\newcommand{\rhabEnetLoIvAllP}{0.898}  % p-value
\newcommand{\rhabEnetLoIvCloseQlike}{0.1243}  % QLIKE enet lo_iv c16
\newcommand{\rhabEnetLoIvCloseMse}{4.03}  % MSE x 1e10 enet lo_iv c16
\newcommand{\rhabEnetLoIvClosePct}{$-$1.5}  % QLIKE change % vs full
\newcommand{\rhabEnetLoIvCloseDm}{$-$1.51}  % DM vs the reference arm
\newcommand{\rhabEnetLoIvCloseP}{0.132}  % p-value
\newcommand{\rhabEnetLoIvCloseLateQlike}{0.1193}  % QLIKE enet lo_iv c16_2020
\newcommand{\rhabEnetLoIvCloseLateMse}{9.21}  % MSE x 1e10 enet lo_iv c16_2020
\newcommand{\rhabEnetLoIvCloseLatePct}{$-$3.3}  % QLIKE change % vs full
\newcommand{\rhabEnetLoIvCloseLateDm}{$-$1.72}  % DM vs the reference arm
\newcommand{\rhabEnetLoIvCloseLateP}{0.086}  % p-value
\newcommand{\rhabEnetLoLiqAllQlike}{0.1779}  % QLIKE enet lo_liq all
\newcommand{\rhabEnetLoLiqAllMse}{2.43}  % MSE x 1e10 enet lo_liq all
\newcommand{\rhabEnetLoLiqAllPct}{0.1}  % QLIKE change % vs full
\newcommand{\rhabEnetLoLiqAllDm}{0.35}  % DM vs the reference arm
\newcommand{\rhabEnetLoLiqAllP}{0.724}  % p-value
\newcommand{\rhabEnetLoLiqCloseQlike}{0.1251}  % QLIKE enet lo_liq c16
\newcommand{\rhabEnetLoLiqCloseMse}{4.29}  % MSE x 1e10 enet lo_liq c16
\newcommand{\rhabEnetLoLiqClosePct}{$-$0.9}  % QLIKE change % vs full
\newcommand{\rhabEnetLoLiqCloseDm}{$-$1.54}  % DM vs the reference arm
\newcommand{\rhabEnetLoLiqCloseP}{0.123}  % p-value
\newcommand{\rhabEnetLoLiqCloseLateQlike}{0.1223}  % QLIKE enet lo_liq c16_2020
\newcommand{\rhabEnetLoLiqCloseLateMse}{9.84}  % MSE x 1e10 enet lo_liq c16_2020
\newcommand{\rhabEnetLoLiqCloseLatePct}{$-$0.9}  % QLIKE change % vs full
\newcommand{\rhabEnetLoLiqCloseLateDm}{$-$0.92}  % DM vs the reference arm
\newcommand{\rhabEnetLoLiqCloseLateP}{0.358}  % p-value
\newcommand{\rhabEnetLoRmAllQlike}{0.1813}  % QLIKE enet lo_rm all
\newcommand{\rhabEnetLoRmAllMse}{2.00}  % MSE x 1e10 enet lo_rm all
\newcommand{\rhabEnetLoRmAllPct}{1.9}  % QLIKE change % vs full
\newcommand{\rhabEnetLoRmAllDm}{4.38}  % DM vs the reference arm
\newcommand{\rhabEnetLoRmAllP}{\ensuremath{<}0.001}  % p-value
\newcommand{\rhabEnetLoRmCloseQlike}{0.1285}  % QLIKE enet lo_rm c16
\newcommand{\rhabEnetLoRmCloseMse}{2.35}  % MSE x 1e10 enet lo_rm c16
\newcommand{\rhabEnetLoRmClosePct}{1.8}  % QLIKE change % vs full
\newcommand{\rhabEnetLoRmCloseDm}{1.31}  % DM vs the reference arm
\newcommand{\rhabEnetLoRmCloseP}{0.190}  % p-value
\newcommand{\rhabEnetLoRmCloseLateQlike}{0.1217}  % QLIKE enet lo_rm c16_2020
\newcommand{\rhabEnetLoRmCloseLateMse}{5.18}  % MSE x 1e10 enet lo_rm c16_2020
\newcommand{\rhabEnetLoRmCloseLatePct}{$-$1.3}  % QLIKE change % vs full
\newcommand{\rhabEnetLoRmCloseLateDm}{$-$0.62}  % DM vs the reference arm
\newcommand{\rhabEnetLoRmCloseLateP}{0.538}  % p-value
\newcommand{\rhabEnetLoSentAllQlike}{0.1782}  % QLIKE enet lo_sent all
\newcommand{\rhabEnetLoSentAllMse}{2.43}  % MSE x 1e10 enet lo_sent all
\newcommand{\rhabEnetLoSentAllPct}{0.2}  % QLIKE change % vs full
\newcommand{\rhabEnetLoSentAllDm}{1.49}  % DM vs the reference arm
\newcommand{\rhabEnetLoSentAllP}{0.136}  % p-value
\newcommand{\rhabEnetLoSentCloseQlike}{0.1253}  % QLIKE enet lo_sent c16
\newcommand{\rhabEnetLoSentCloseMse}{3.94}  % MSE x 1e10 enet lo_sent c16
\newcommand{\rhabEnetLoSentClosePct}{$-$0.7}  % QLIKE change % vs full
\newcommand{\rhabEnetLoSentCloseDm}{$-$2.12}  % DM vs the reference arm
\newcommand{\rhabEnetLoSentCloseP}{0.034}  % p-value
\newcommand{\rhabEnetLoSentCloseLateQlike}{0.1212}  % QLIKE enet lo_sent c16_2020
\newcommand{\rhabEnetLoSentCloseLateMse}{9.01}  % MSE x 1e10 enet lo_sent c16_2020
\newcommand{\rhabEnetLoSentCloseLatePct}{$-$1.7}  % QLIKE change % vs full
\newcommand{\rhabEnetLoSentCloseLateDm}{$-$2.21}  % DM vs the reference arm
\newcommand{\rhabEnetLoSentCloseLateP}{0.027}  % p-value
\newcommand{\rhabEnetLoVdAllQlike}{0.1777}  % QLIKE enet lo_vd all
\newcommand{\rhabEnetLoVdAllMse}{2.43}  % MSE x 1e10 enet lo_vd all
\newcommand{\rhabEnetLoVdAllPct}{$-$0.1}  % QLIKE change % vs full
\newcommand{\rhabEnetLoVdAllDm}{$-$1.13}  % DM vs the reference arm
\newcommand{\rhabEnetLoVdAllP}{0.260}  % p-value
\newcommand{\rhabEnetLoVdCloseQlike}{0.1252}  % QLIKE enet lo_vd c16
\newcommand{\rhabEnetLoVdCloseMse}{3.94}  % MSE x 1e10 enet lo_vd c16
\newcommand{\rhabEnetLoVdClosePct}{$-$0.8}  % QLIKE change % vs full
\newcommand{\rhabEnetLoVdCloseDm}{$-$1.94}  % DM vs the reference arm
\newcommand{\rhabEnetLoVdCloseP}{0.053}  % p-value
\newcommand{\rhabEnetLoVdCloseLateQlike}{0.1207}  % QLIKE enet lo_vd c16_2020
\newcommand{\rhabEnetLoVdCloseLateMse}{9.01}  % MSE x 1e10 enet lo_vd c16_2020
\newcommand{\rhabEnetLoVdCloseLatePct}{$-$2.1}  % QLIKE change % vs full
\newcommand{\rhabEnetLoVdCloseLateDm}{$-$2.18}  % DM vs the reference arm
\newcommand{\rhabEnetLoVdCloseLateP}{0.029}  % p-value
\newcommand{\rhabEnetLoVwAllQlike}{0.1777}  % QLIKE enet lo_vw all
\newcommand{\rhabEnetLoVwAllMse}{2.08}  % MSE x 1e10 enet lo_vw all
\newcommand{\rhabEnetLoVwAllPct}{$-$0.1}  % QLIKE change % vs full
\newcommand{\rhabEnetLoVwAllDm}{$-$0.65}  % DM vs the reference arm
\newcommand{\rhabEnetLoVwAllP}{0.516}  % p-value
\newcommand{\rhabEnetLoVwCloseQlike}{0.1254}  % QLIKE enet lo_vw c16
\newcommand{\rhabEnetLoVwCloseMse}{3.77}  % MSE x 1e10 enet lo_vw c16
\newcommand{\rhabEnetLoVwClosePct}{$-$0.7}  % QLIKE change % vs full
\newcommand{\rhabEnetLoVwCloseDm}{$-$1.66}  % DM vs the reference arm
\newcommand{\rhabEnetLoVwCloseP}{0.097}  % p-value
\newcommand{\rhabEnetLoVwCloseLateQlike}{0.1218}  % QLIKE enet lo_vw c16_2020
\newcommand{\rhabEnetLoVwCloseLateMse}{8.59}  % MSE x 1e10 enet lo_vw c16_2020
\newcommand{\rhabEnetLoVwCloseLatePct}{$-$1.3}  % QLIKE change % vs full
\newcommand{\rhabEnetLoVwCloseLateDm}{$-$2.65}  % DM vs the reference arm
\newcommand{\rhabEnetLoVwCloseLateP}{0.008}  % p-value
\newcommand{\rhabEnetSharAllQlike}{0.1877}  % QLIKE enet shar all
\newcommand{\rhabEnetSharAllMse}{1.38}  % MSE x 1e10 enet shar all
\newcommand{\rhabEnetSharAllPct}{$-$0.9}  % QLIKE change % vs har
\newcommand{\rhabEnetSharAllDm}{$-$3.68}  % DM vs the reference arm
\newcommand{\rhabEnetSharAllP}{\ensuremath{<}0.001}  % p-value
\newcommand{\rhabEnetSharCloseQlike}{0.1265}  % QLIKE enet shar c16
\newcommand{\rhabEnetSharCloseMse}{1.83}  % MSE x 1e10 enet shar c16
\newcommand{\rhabEnetSharClosePct}{$-$0.3}  % QLIKE change % vs har
\newcommand{\rhabEnetSharCloseDm}{$-$0.54}  % DM vs the reference arm
\newcommand{\rhabEnetSharCloseP}{0.586}  % p-value
\newcommand{\rhabEnetSharCloseLateQlike}{0.1141}  % QLIKE enet shar c16_2020
\newcommand{\rhabEnetSharCloseLateMse}{3.92}  % MSE x 1e10 enet shar c16_2020
\newcommand{\rhabEnetSharCloseLatePct}{0.3}  % QLIKE change % vs har
\newcommand{\rhabEnetSharCloseLateDm}{0.34}  % DM vs the reference arm
\newcommand{\rhabEnetSharCloseLateP}{0.731}  % p-value
\newcommand{\rhabEnetSharqlAllQlike}{0.1840}  % QLIKE enet sharql all
\newcommand{\rhabEnetSharqlAllMse}{1.69}  % MSE x 1e10 enet sharql all
\newcommand{\rhabEnetSharqlAllPct}{$-$2.8}  % QLIKE change % vs har
\newcommand{\rhabEnetSharqlAllDm}{$-$6.95}  % DM vs the reference arm
\newcommand{\rhabEnetSharqlAllP}{\ensuremath{<}0.001}  % p-value
\newcommand{\rhabEnetSharqlCloseQlike}{0.1216}  % QLIKE enet sharql c16
\newcommand{\rhabEnetSharqlCloseMse}{2.64}  % MSE x 1e10 enet sharql c16
\newcommand{\rhabEnetSharqlClosePct}{$-$4.1}  % QLIKE change % vs har
\newcommand{\rhabEnetSharqlCloseDm}{$-$3.46}  % DM vs the reference arm
\newcommand{\rhabEnetSharqlCloseP}{\ensuremath{<}0.001}  % p-value
\newcommand{\rhabEnetSharqlCloseLateQlike}{0.1121}  % QLIKE enet sharql c16_2020
\newcommand{\rhabEnetSharqlCloseLateMse}{5.90}  % MSE x 1e10 enet sharql c16_2020
\newcommand{\rhabEnetSharqlCloseLatePct}{$-$1.4}  % QLIKE change % vs har
\newcommand{\rhabEnetSharqlCloseLateDm}{$-$0.64}  % DM vs the reference arm
\newcommand{\rhabEnetSharqlCloseLateP}{0.524}  % p-value
\newcommand{\rhabLassoAddEwAllQlike}{0.1877}  % QLIKE lasso add_ew all
\newcommand{\rhabLassoAddEwAllMse}{1.47}  % MSE x 1e10 lasso add_ew all
\newcommand{\rhabLassoAddEwAllPct}{$-$0.8}  % QLIKE change % vs har
\newcommand{\rhabLassoAddEwAllDm}{$-$2.01}  % DM vs the reference arm
\newcommand{\rhabLassoAddEwAllP}{0.045}  % p-value
\newcommand{\rhabLassoAddEwCloseQlike}{0.1257}  % QLIKE lasso add_ew c16
\newcommand{\rhabLassoAddEwCloseMse}{2.13}  % MSE x 1e10 lasso add_ew c16
\newcommand{\rhabLassoAddEwClosePct}{$-$1.0}  % QLIKE change % vs har
\newcommand{\rhabLassoAddEwCloseDm}{$-$0.81}  % DM vs the reference arm
\newcommand{\rhabLassoAddEwCloseP}{0.417}  % p-value
\newcommand{\rhabLassoAddEwCloseLateQlike}{0.1123}  % QLIKE lasso add_ew c16_2020
\newcommand{\rhabLassoAddEwCloseLateMse}{4.65}  % MSE x 1e10 lasso add_ew c16_2020
\newcommand{\rhabLassoAddEwCloseLatePct}{$-$1.2}  % QLIKE change % vs har
\newcommand{\rhabLassoAddEwCloseLateDm}{$-$0.70}  % DM vs the reference arm
\newcommand{\rhabLassoAddEwCloseLateP}{0.485}  % p-value
\newcommand{\rhabLassoAddFomcAllQlike}{0.1820}  % QLIKE lasso add_fomc all
\newcommand{\rhabLassoAddFomcAllMse}{1.13}  % MSE x 1e10 lasso add_fomc all
\newcommand{\rhabLassoAddFomcAllPct}{$-$3.8}  % QLIKE change % vs har
\newcommand{\rhabLassoAddFomcAllDm}{$-$5.76}  % DM vs the reference arm
\newcommand{\rhabLassoAddFomcAllP}{\ensuremath{<}0.001}  % p-value
\newcommand{\rhabLassoAddFomcCloseQlike}{0.1296}  % QLIKE lasso add_fomc c16
\newcommand{\rhabLassoAddFomcCloseMse}{1.83}  % MSE x 1e10 lasso add_fomc c16
\newcommand{\rhabLassoAddFomcClosePct}{2.2}  % QLIKE change % vs har
\newcommand{\rhabLassoAddFomcCloseDm}{4.33}  % DM vs the reference arm
\newcommand{\rhabLassoAddFomcCloseP}{\ensuremath{<}0.001}  % p-value
\newcommand{\rhabLassoAddFomcCloseLateQlike}{0.1135}  % QLIKE lasso add_fomc c16_2020
\newcommand{\rhabLassoAddFomcCloseLateMse}{3.92}  % MSE x 1e10 lasso add_fomc c16_2020
\newcommand{\rhabLassoAddFomcCloseLatePct}{$-$0.1}  % QLIKE change % vs har
\newcommand{\rhabLassoAddFomcCloseLateDm}{$-$0.16}  % DM vs the reference arm
\newcommand{\rhabLassoAddFomcCloseLateP}{0.875}  % p-value
\newcommand{\rhabLassoAddIvAllQlike}{0.1887}  % QLIKE lasso add_iv all
\newcommand{\rhabLassoAddIvAllMse}{1.91}  % MSE x 1e10 lasso add_iv all
\newcommand{\rhabLassoAddIvAllPct}{$-$0.3}  % QLIKE change % vs har
\newcommand{\rhabLassoAddIvAllDm}{$-$0.79}  % DM vs the reference arm
\newcommand{\rhabLassoAddIvAllP}{0.430}  % p-value
\newcommand{\rhabLassoAddIvCloseQlike}{0.1263}  % QLIKE lasso add_iv c16
\newcommand{\rhabLassoAddIvCloseMse}{1.82}  % MSE x 1e10 lasso add_iv c16
\newcommand{\rhabLassoAddIvClosePct}{$-$0.4}  % QLIKE change % vs har
\newcommand{\rhabLassoAddIvCloseDm}{$-$0.38}  % DM vs the reference arm
\newcommand{\rhabLassoAddIvCloseP}{0.705}  % p-value
\newcommand{\rhabLassoAddIvCloseLateQlike}{0.1144}  % QLIKE lasso add_iv c16_2020
\newcommand{\rhabLassoAddIvCloseLateMse}{3.92}  % MSE x 1e10 lasso add_iv c16_2020
\newcommand{\rhabLassoAddIvCloseLatePct}{0.6}  % QLIKE change % vs har
\newcommand{\rhabLassoAddIvCloseLateDm}{0.29}  % DM vs the reference arm
\newcommand{\rhabLassoAddIvCloseLateP}{0.768}  % p-value
\newcommand{\rhabLassoAddLiqAllQlike}{0.1888}  % QLIKE lasso add_liq all
\newcommand{\rhabLassoAddLiqAllMse}{1.17}  % MSE x 1e10 lasso add_liq all
\newcommand{\rhabLassoAddLiqAllPct}{$-$0.3}  % QLIKE change % vs har
\newcommand{\rhabLassoAddLiqAllDm}{$-$0.55}  % DM vs the reference arm
\newcommand{\rhabLassoAddLiqAllP}{0.583}  % p-value
\newcommand{\rhabLassoAddLiqCloseQlike}{0.1305}  % QLIKE lasso add_liq c16
\newcommand{\rhabLassoAddLiqCloseMse}{2.02}  % MSE x 1e10 lasso add_liq c16
\newcommand{\rhabLassoAddLiqClosePct}{2.9}  % QLIKE change % vs har
\newcommand{\rhabLassoAddLiqCloseDm}{1.86}  % DM vs the reference arm
\newcommand{\rhabLassoAddLiqCloseP}{0.063}  % p-value
\newcommand{\rhabLassoAddLiqCloseLateQlike}{0.1206}  % QLIKE lasso add_liq c16_2020
\newcommand{\rhabLassoAddLiqCloseLateMse}{4.37}  % MSE x 1e10 lasso add_liq c16_2020
\newcommand{\rhabLassoAddLiqCloseLatePct}{6.1}  % QLIKE change % vs har
\newcommand{\rhabLassoAddLiqCloseLateDm}{2.00}  % DM vs the reference arm
\newcommand{\rhabLassoAddLiqCloseLateP}{0.045}  % p-value
\newcommand{\rhabLassoAddRmAllQlike}{0.1852}  % QLIKE lasso add_rm all
\newcommand{\rhabLassoAddRmAllMse}{1.81}  % MSE x 1e10 lasso add_rm all
\newcommand{\rhabLassoAddRmAllPct}{$-$2.1}  % QLIKE change % vs har
\newcommand{\rhabLassoAddRmAllDm}{$-$2.64}  % DM vs the reference arm
\newcommand{\rhabLassoAddRmAllP}{0.008}  % p-value
\newcommand{\rhabLassoAddRmCloseQlike}{0.1234}  % QLIKE lasso add_rm c16
\newcommand{\rhabLassoAddRmCloseMse}{3.11}  % MSE x 1e10 lasso add_rm c16
\newcommand{\rhabLassoAddRmClosePct}{$-$2.8}  % QLIKE change % vs har
\newcommand{\rhabLassoAddRmCloseDm}{$-$1.41}  % DM vs the reference arm
\newcommand{\rhabLassoAddRmCloseP}{0.160}  % p-value
\newcommand{\rhabLassoAddRmCloseLateQlike}{0.1176}  % QLIKE lasso add_rm c16_2020
\newcommand{\rhabLassoAddRmCloseLateMse}{7.01}  % MSE x 1e10 lasso add_rm c16_2020
\newcommand{\rhabLassoAddRmCloseLatePct}{3.5}  % QLIKE change % vs har
\newcommand{\rhabLassoAddRmCloseLateDm}{0.90}  % DM vs the reference arm
\newcommand{\rhabLassoAddRmCloseLateP}{0.368}  % p-value
\newcommand{\rhabLassoAddSentAllQlike}{0.1895}  % QLIKE lasso add_sent all
\newcommand{\rhabLassoAddSentAllMse}{1.11}  % MSE x 1e10 lasso add_sent all
\newcommand{\rhabLassoAddSentAllPct}{0.1}  % QLIKE change % vs har
\newcommand{\rhabLassoAddSentAllDm}{0.60}  % DM vs the reference arm
\newcommand{\rhabLassoAddSentAllP}{0.550}  % p-value
\newcommand{\rhabLassoAddSentCloseQlike}{0.1310}  % QLIKE lasso add_sent c16
\newcommand{\rhabLassoAddSentCloseMse}{1.98}  % MSE x 1e10 lasso add_sent c16
\newcommand{\rhabLassoAddSentClosePct}{3.2}  % QLIKE change % vs har
\newcommand{\rhabLassoAddSentCloseDm}{4.34}  % DM vs the reference arm
\newcommand{\rhabLassoAddSentCloseP}{\ensuremath{<}0.001}  % p-value
\newcommand{\rhabLassoAddSentCloseLateQlike}{0.1214}  % QLIKE lasso add_sent c16_2020
\newcommand{\rhabLassoAddSentCloseLateMse}{4.27}  % MSE x 1e10 lasso add_sent c16_2020
\newcommand{\rhabLassoAddSentCloseLatePct}{6.8}  % QLIKE change % vs har
\newcommand{\rhabLassoAddSentCloseLateDm}{3.64}  % DM vs the reference arm
\newcommand{\rhabLassoAddSentCloseLateP}{\ensuremath{<}0.001}  % p-value
\newcommand{\rhabLassoAddVdAllQlike}{0.1903}  % QLIKE lasso add_vd all
\newcommand{\rhabLassoAddVdAllMse}{1.12}  % MSE x 1e10 lasso add_vd all
\newcommand{\rhabLassoAddVdAllPct}{0.6}  % QLIKE change % vs har
\newcommand{\rhabLassoAddVdAllDm}{3.05}  % DM vs the reference arm
\newcommand{\rhabLassoAddVdAllP}{0.002}  % p-value
\newcommand{\rhabLassoAddVdCloseQlike}{0.1272}  % QLIKE lasso add_vd c16
\newcommand{\rhabLassoAddVdCloseMse}{1.95}  % MSE x 1e10 lasso add_vd c16
\newcommand{\rhabLassoAddVdClosePct}{0.3}  % QLIKE change % vs har
\newcommand{\rhabLassoAddVdCloseDm}{0.65}  % DM vs the reference arm
\newcommand{\rhabLassoAddVdCloseP}{0.516}  % p-value
\newcommand{\rhabLassoAddVdCloseLateQlike}{0.1169}  % QLIKE lasso add_vd c16_2020
\newcommand{\rhabLassoAddVdCloseLateMse}{4.20}  % MSE x 1e10 lasso add_vd c16_2020
\newcommand{\rhabLassoAddVdCloseLatePct}{2.9}  % QLIKE change % vs har
\newcommand{\rhabLassoAddVdCloseLateDm}{3.03}  % DM vs the reference arm
\newcommand{\rhabLassoAddVdCloseLateP}{0.002}  % p-value
\newcommand{\rhabLassoAddVwAllQlike}{0.1881}  % QLIKE lasso add_vw all
\newcommand{\rhabLassoAddVwAllMse}{2.04}  % MSE x 1e10 lasso add_vw all
\newcommand{\rhabLassoAddVwAllPct}{$-$0.6}  % QLIKE change % vs har
\newcommand{\rhabLassoAddVwAllDm}{$-$1.02}  % DM vs the reference arm
\newcommand{\rhabLassoAddVwAllP}{0.309}  % p-value
\newcommand{\rhabLassoAddVwCloseQlike}{0.1262}  % QLIKE lasso add_vw c16
\newcommand{\rhabLassoAddVwCloseMse}{2.38}  % MSE x 1e10 lasso add_vw c16
\newcommand{\rhabLassoAddVwClosePct}{$-$0.5}  % QLIKE change % vs har
\newcommand{\rhabLassoAddVwCloseDm}{$-$0.35}  % DM vs the reference arm
\newcommand{\rhabLassoAddVwCloseP}{0.726}  % p-value
\newcommand{\rhabLassoAddVwCloseLateQlike}{0.1136}  % QLIKE lasso add_vw c16_2020
\newcommand{\rhabLassoAddVwCloseLateMse}{5.24}  % MSE x 1e10 lasso add_vw c16_2020
\newcommand{\rhabLassoAddVwCloseLatePct}{$-$0.1}  % QLIKE change % vs har
\newcommand{\rhabLassoAddVwCloseLateDm}{$-$0.04}  % DM vs the reference arm
\newcommand{\rhabLassoAddVwCloseLateP}{0.968}  % p-value
\newcommand{\rhabLassoFullAllQlike}{0.1775}  % QLIKE lasso full all
\newcommand{\rhabLassoFullAllMse}{2.37}  % MSE x 1e10 lasso full all
\newcommand{\rhabLassoFullAllPct}{$-$6.2}  % QLIKE change % vs har
\newcommand{\rhabLassoFullAllDm}{$-$4.72}  % DM vs the reference arm
\newcommand{\rhabLassoFullAllP}{\ensuremath{<}0.001}  % p-value
\newcommand{\rhabLassoFullCloseQlike}{0.1310}  % QLIKE lasso full c16
\newcommand{\rhabLassoFullCloseMse}{3.57}  % MSE x 1e10 lasso full c16
\newcommand{\rhabLassoFullClosePct}{3.3}  % QLIKE change % vs har
\newcommand{\rhabLassoFullCloseDm}{1.04}  % DM vs the reference arm
\newcommand{\rhabLassoFullCloseP}{0.300}  % p-value
\newcommand{\rhabLassoFullCloseLateQlike}{0.1301}  % QLIKE lasso full c16_2020
\newcommand{\rhabLassoFullCloseLateMse}{8.13}  % MSE x 1e10 lasso full c16_2020
\newcommand{\rhabLassoFullCloseLatePct}{14.4}  % QLIKE change % vs har
\newcommand{\rhabLassoFullCloseLateDm}{2.51}  % DM vs the reference arm
\newcommand{\rhabLassoFullCloseLateP}{0.012}  % p-value
\newcommand{\rhabLassoHarAllQlike}{0.1893}  % QLIKE lasso har all
\newcommand{\rhabLassoHarAllMse}{1.11}  % MSE x 1e10 lasso har all
\newcommand{\rhabLassoHarCloseQlike}{0.1269}  % QLIKE lasso har c16
\newcommand{\rhabLassoHarCloseMse}{1.97}  % MSE x 1e10 lasso har c16
\newcommand{\rhabLassoHarCloseLateQlike}{0.1137}  % QLIKE lasso har c16_2020
\newcommand{\rhabLassoHarCloseLateMse}{4.24}  % MSE x 1e10 lasso har c16_2020
\newcommand{\rhabLassoHarqAllQlike}{0.1893}  % QLIKE lasso harq all
\newcommand{\rhabLassoHarqAllMse}{1.17}  % MSE x 1e10 lasso harq all
\newcommand{\rhabLassoHarqAllPct}{0.0}  % QLIKE change % vs har
\newcommand{\rhabLassoHarqAllDm}{0.61}  % DM vs the reference arm
\newcommand{\rhabLassoHarqAllP}{0.544}  % p-value
\newcommand{\rhabLassoHarqCloseQlike}{0.1267}  % QLIKE lasso harq c16
\newcommand{\rhabLassoHarqCloseMse}{1.78}  % MSE x 1e10 lasso harq c16
\newcommand{\rhabLassoHarqClosePct}{$-$0.2}  % QLIKE change % vs har
\newcommand{\rhabLassoHarqCloseDm}{$-$1.37}  % DM vs the reference arm
\newcommand{\rhabLassoHarqCloseP}{0.171}  % p-value
\newcommand{\rhabLassoHarqCloseLateQlike}{0.1133}  % QLIKE lasso harq c16_2020
\newcommand{\rhabLassoHarqCloseLateMse}{3.80}  % MSE x 1e10 lasso harq c16_2020
\newcommand{\rhabLassoHarqCloseLatePct}{$-$0.3}  % QLIKE change % vs har
\newcommand{\rhabLassoHarqCloseLateDm}{$-$1.18}  % DM vs the reference arm
\newcommand{\rhabLassoHarqCloseLateP}{0.238}  % p-value
\newcommand{\rhabLassoLharAllQlike}{0.1842}  % QLIKE lasso lhar all
\newcommand{\rhabLassoLharAllMse}{1.71}  % MSE x 1e10 lasso lhar all
\newcommand{\rhabLassoLharAllPct}{$-$2.7}  % QLIKE change % vs har
\newcommand{\rhabLassoLharAllDm}{$-$6.30}  % DM vs the reference arm
\newcommand{\rhabLassoLharAllP}{\ensuremath{<}0.001}  % p-value
\newcommand{\rhabLassoLharCloseQlike}{0.1219}  % QLIKE lasso lhar c16
\newcommand{\rhabLassoLharCloseMse}{2.80}  % MSE x 1e10 lasso lhar c16
\newcommand{\rhabLassoLharClosePct}{$-$3.9}  % QLIKE change % vs har
\newcommand{\rhabLassoLharCloseDm}{$-$2.77}  % DM vs the reference arm
\newcommand{\rhabLassoLharCloseP}{0.006}  % p-value
\newcommand{\rhabLassoLharCloseLateQlike}{0.1130}  % QLIKE lasso lhar c16_2020
\newcommand{\rhabLassoLharCloseLateMse}{6.28}  % MSE x 1e10 lasso lhar c16_2020
\newcommand{\rhabLassoLharCloseLatePct}{$-$0.6}  % QLIKE change % vs har
\newcommand{\rhabLassoLharCloseLateDm}{$-$0.21}  % DM vs the reference arm
\newcommand{\rhabLassoLharCloseLateP}{0.836}  % p-value
\newcommand{\rhabLassoLoEwAllQlike}{0.1770}  % QLIKE lasso lo_ew all
\newcommand{\rhabLassoLoEwAllMse}{2.54}  % MSE x 1e10 lasso lo_ew all
\newcommand{\rhabLassoLoEwAllPct}{$-$0.3}  % QLIKE change % vs full
\newcommand{\rhabLassoLoEwAllDm}{$-$1.46}  % DM vs the reference arm
\newcommand{\rhabLassoLoEwAllP}{0.144}  % p-value
\newcommand{\rhabLassoLoEwCloseQlike}{0.1310}  % QLIKE lasso lo_ew c16
\newcommand{\rhabLassoLoEwCloseMse}{3.82}  % MSE x 1e10 lasso lo_ew c16
\newcommand{\rhabLassoLoEwClosePct}{0.0}  % QLIKE change % vs full
\newcommand{\rhabLassoLoEwCloseDm}{0.00}  % DM vs the reference arm
\newcommand{\rhabLassoLoEwCloseP}{0.999}  % p-value
\newcommand{\rhabLassoLoEwCloseLateQlike}{0.1289}  % QLIKE lasso lo_ew c16_2020
\newcommand{\rhabLassoLoEwCloseLateMse}{8.72}  % MSE x 1e10 lasso lo_ew c16_2020
\newcommand{\rhabLassoLoEwCloseLatePct}{$-$0.9}  % QLIKE change % vs full
\newcommand{\rhabLassoLoEwCloseLateDm}{$-$0.76}  % DM vs the reference arm
\newcommand{\rhabLassoLoEwCloseLateP}{0.446}  % p-value
\newcommand{\rhabLassoLoFomcAllQlike}{0.1839}  % QLIKE lasso lo_fomc all
\newcommand{\rhabLassoLoFomcAllMse}{2.39}  % MSE x 1e10 lasso lo_fomc all
\newcommand{\rhabLassoLoFomcAllPct}{3.6}  % QLIKE change % vs full
\newcommand{\rhabLassoLoFomcAllDm}{5.10}  % DM vs the reference arm
\newcommand{\rhabLassoLoFomcAllP}{\ensuremath{<}0.001}  % p-value
\newcommand{\rhabLassoLoFomcCloseQlike}{0.1293}  % QLIKE lasso lo_fomc c16
\newcommand{\rhabLassoLoFomcCloseMse}{3.62}  % MSE x 1e10 lasso lo_fomc c16
\newcommand{\rhabLassoLoFomcClosePct}{$-$1.4}  % QLIKE change % vs full
\newcommand{\rhabLassoLoFomcCloseDm}{$-$3.21}  % DM vs the reference arm
\newcommand{\rhabLassoLoFomcCloseP}{0.001}  % p-value
\newcommand{\rhabLassoLoFomcCloseLateQlike}{0.1293}  % QLIKE lasso lo_fomc c16_2020
\newcommand{\rhabLassoLoFomcCloseLateMse}{8.25}  % MSE x 1e10 lasso lo_fomc c16_2020
\newcommand{\rhabLassoLoFomcCloseLatePct}{$-$0.6}  % QLIKE change % vs full
\newcommand{\rhabLassoLoFomcCloseLateDm}{$-$0.92}  % DM vs the reference arm
\newcommand{\rhabLassoLoFomcCloseLateP}{0.358}  % p-value
\newcommand{\rhabLassoLoIvAllQlike}{0.1778}  % QLIKE lasso lo_iv all
\newcommand{\rhabLassoLoIvAllMse}{2.43}  % MSE x 1e10 lasso lo_iv all
\newcommand{\rhabLassoLoIvAllPct}{0.2}  % QLIKE change % vs full
\newcommand{\rhabLassoLoIvAllDm}{0.34}  % DM vs the reference arm
\newcommand{\rhabLassoLoIvAllP}{0.731}  % p-value
\newcommand{\rhabLassoLoIvCloseQlike}{0.1284}  % QLIKE lasso lo_iv c16
\newcommand{\rhabLassoLoIvCloseMse}{3.90}  % MSE x 1e10 lasso lo_iv c16
\newcommand{\rhabLassoLoIvClosePct}{$-$2.0}  % QLIKE change % vs full
\newcommand{\rhabLassoLoIvCloseDm}{$-$1.66}  % DM vs the reference arm
\newcommand{\rhabLassoLoIvCloseP}{0.096}  % p-value
\newcommand{\rhabLassoLoIvCloseLateQlike}{0.1251}  % QLIKE lasso lo_iv c16_2020
\newcommand{\rhabLassoLoIvCloseLateMse}{8.92}  % MSE x 1e10 lasso lo_iv c16_2020
\newcommand{\rhabLassoLoIvCloseLatePct}{$-$3.8}  % QLIKE change % vs full
\newcommand{\rhabLassoLoIvCloseLateDm}{$-$1.77}  % DM vs the reference arm
\newcommand{\rhabLassoLoIvCloseLateP}{0.076}  % p-value
\newcommand{\rhabLassoLoLiqAllQlike}{0.1774}  % QLIKE lasso lo_liq all
\newcommand{\rhabLassoLoLiqAllMse}{2.39}  % MSE x 1e10 lasso lo_liq all
\newcommand{\rhabLassoLoLiqAllPct}{$-$0.1}  % QLIKE change % vs full
\newcommand{\rhabLassoLoLiqAllDm}{$-$0.21}  % DM vs the reference arm
\newcommand{\rhabLassoLoLiqAllP}{0.831}  % p-value
\newcommand{\rhabLassoLoLiqCloseQlike}{0.1285}  % QLIKE lasso lo_liq c16
\newcommand{\rhabLassoLoLiqCloseMse}{4.03}  % MSE x 1e10 lasso lo_liq c16
\newcommand{\rhabLassoLoLiqClosePct}{$-$1.9}  % QLIKE change % vs full
\newcommand{\rhabLassoLoLiqCloseDm}{$-$2.17}  % DM vs the reference arm
\newcommand{\rhabLassoLoLiqCloseP}{0.030}  % p-value
\newcommand{\rhabLassoLoLiqCloseLateQlike}{0.1272}  % QLIKE lasso lo_liq c16_2020
\newcommand{\rhabLassoLoLiqCloseLateMse}{9.23}  % MSE x 1e10 lasso lo_liq c16_2020
\newcommand{\rhabLassoLoLiqCloseLatePct}{$-$2.2}  % QLIKE change % vs full
\newcommand{\rhabLassoLoLiqCloseLateDm}{$-$1.99}  % DM vs the reference arm
\newcommand{\rhabLassoLoLiqCloseLateP}{0.047}  % p-value
\newcommand{\rhabLassoLoRmAllQlike}{0.1798}  % QLIKE lasso lo_rm all
\newcommand{\rhabLassoLoRmAllMse}{2.07}  % MSE x 1e10 lasso lo_rm all
\newcommand{\rhabLassoLoRmAllPct}{1.3}  % QLIKE change % vs full
\newcommand{\rhabLassoLoRmAllDm}{2.07}  % DM vs the reference arm
\newcommand{\rhabLassoLoRmAllP}{0.039}  % p-value
\newcommand{\rhabLassoLoRmCloseQlike}{0.1311}  % QLIKE lasso lo_rm c16
\newcommand{\rhabLassoLoRmCloseMse}{2.27}  % MSE x 1e10 lasso lo_rm c16
\newcommand{\rhabLassoLoRmClosePct}{0.0}  % QLIKE change % vs full
\newcommand{\rhabLassoLoRmCloseDm}{0.01}  % DM vs the reference arm
\newcommand{\rhabLassoLoRmCloseP}{0.990}  % p-value
\newcommand{\rhabLassoLoRmCloseLateQlike}{0.1255}  % QLIKE lasso lo_rm c16_2020
\newcommand{\rhabLassoLoRmCloseLateMse}{4.99}  % MSE x 1e10 lasso lo_rm c16_2020
\newcommand{\rhabLassoLoRmCloseLatePct}{$-$3.5}  % QLIKE change % vs full
\newcommand{\rhabLassoLoRmCloseLateDm}{$-$1.28}  % DM vs the reference arm
\newcommand{\rhabLassoLoRmCloseLateP}{0.200}  % p-value
\newcommand{\rhabLassoLoSentAllQlike}{0.1776}  % QLIKE lasso lo_sent all
\newcommand{\rhabLassoLoSentAllMse}{2.37}  % MSE x 1e10 lasso lo_sent all
\newcommand{\rhabLassoLoSentAllPct}{0.1}  % QLIKE change % vs full
\newcommand{\rhabLassoLoSentAllDm}{0.41}  % DM vs the reference arm
\newcommand{\rhabLassoLoSentAllP}{0.679}  % p-value
\newcommand{\rhabLassoLoSentCloseQlike}{0.1294}  % QLIKE lasso lo_sent c16
\newcommand{\rhabLassoLoSentCloseMse}{3.60}  % MSE x 1e10 lasso lo_sent c16
\newcommand{\rhabLassoLoSentClosePct}{$-$1.3}  % QLIKE change % vs full
\newcommand{\rhabLassoLoSentCloseDm}{$-$2.54}  % DM vs the reference arm
\newcommand{\rhabLassoLoSentCloseP}{0.011}  % p-value
\newcommand{\rhabLassoLoSentCloseLateQlike}{0.1262}  % QLIKE lasso lo_sent c16_2020
\newcommand{\rhabLassoLoSentCloseLateMse}{8.21}  % MSE x 1e10 lasso lo_sent c16_2020
\newcommand{\rhabLassoLoSentCloseLatePct}{$-$3.0}  % QLIKE change % vs full
\newcommand{\rhabLassoLoSentCloseLateDm}{$-$2.62}  % DM vs the reference arm
\newcommand{\rhabLassoLoSentCloseLateP}{0.009}  % p-value
\newcommand{\rhabLassoLoVdAllQlike}{0.1773}  % QLIKE lasso lo_vd all
\newcommand{\rhabLassoLoVdAllMse}{2.35}  % MSE x 1e10 lasso lo_vd all
\newcommand{\rhabLassoLoVdAllPct}{$-$0.1}  % QLIKE change % vs full
\newcommand{\rhabLassoLoVdAllDm}{$-$0.63}  % DM vs the reference arm
\newcommand{\rhabLassoLoVdAllP}{0.531}  % p-value
\newcommand{\rhabLassoLoVdCloseQlike}{0.1299}  % QLIKE lasso lo_vd c16
\newcommand{\rhabLassoLoVdCloseMse}{3.56}  % MSE x 1e10 lasso lo_vd c16
\newcommand{\rhabLassoLoVdClosePct}{$-$0.9}  % QLIKE change % vs full
\newcommand{\rhabLassoLoVdCloseDm}{$-$2.10}  % DM vs the reference arm
\newcommand{\rhabLassoLoVdCloseP}{0.036}  % p-value
\newcommand{\rhabLassoLoVdCloseLateQlike}{0.1275}  % QLIKE lasso lo_vd c16_2020
\newcommand{\rhabLassoLoVdCloseLateMse}{8.12}  % MSE x 1e10 lasso lo_vd c16_2020
\newcommand{\rhabLassoLoVdCloseLatePct}{$-$2.0}  % QLIKE change % vs full
\newcommand{\rhabLassoLoVdCloseLateDm}{$-$2.12}  % DM vs the reference arm
\newcommand{\rhabLassoLoVdCloseLateP}{0.034}  % p-value
\newcommand{\rhabLassoLoVwAllQlike}{0.1769}  % QLIKE lasso lo_vw all
\newcommand{\rhabLassoLoVwAllMse}{2.07}  % MSE x 1e10 lasso lo_vw all
\newcommand{\rhabLassoLoVwAllPct}{$-$0.3}  % QLIKE change % vs full
\newcommand{\rhabLassoLoVwAllDm}{$-$1.09}  % DM vs the reference arm
\newcommand{\rhabLassoLoVwAllP}{0.274}  % p-value
\newcommand{\rhabLassoLoVwCloseQlike}{0.1302}  % QLIKE lasso lo_vw c16
\newcommand{\rhabLassoLoVwCloseMse}{3.61}  % MSE x 1e10 lasso lo_vw c16
\newcommand{\rhabLassoLoVwClosePct}{$-$0.6}  % QLIKE change % vs full
\newcommand{\rhabLassoLoVwCloseDm}{$-$1.06}  % DM vs the reference arm
\newcommand{\rhabLassoLoVwCloseP}{0.289}  % p-value
\newcommand{\rhabLassoLoVwCloseLateQlike}{0.1277}  % QLIKE lasso lo_vw c16_2020
\newcommand{\rhabLassoLoVwCloseLateMse}{8.22}  % MSE x 1e10 lasso lo_vw c16_2020
\newcommand{\rhabLassoLoVwCloseLatePct}{$-$1.8}  % QLIKE change % vs full
\newcommand{\rhabLassoLoVwCloseLateDm}{$-$1.96}  % DM vs the reference arm
\newcommand{\rhabLassoLoVwCloseLateP}{0.050}  % p-value
\newcommand{\rhabLassoSharAllQlike}{0.1877}  % QLIKE lasso shar all
\newcommand{\rhabLassoSharAllMse}{1.39}  % MSE x 1e10 lasso shar all
\newcommand{\rhabLassoSharAllPct}{$-$0.8}  % QLIKE change % vs har
\newcommand{\rhabLassoSharAllDm}{$-$3.52}  % DM vs the reference arm
\newcommand{\rhabLassoSharAllP}{\ensuremath{<}0.001}  % p-value
\newcommand{\rhabLassoSharCloseQlike}{0.1265}  % QLIKE lasso shar c16
\newcommand{\rhabLassoSharCloseMse}{1.84}  % MSE x 1e10 lasso shar c16
\newcommand{\rhabLassoSharClosePct}{$-$0.3}  % QLIKE change % vs har
\newcommand{\rhabLassoSharCloseDm}{$-$0.44}  % DM vs the reference arm
\newcommand{\rhabLassoSharCloseP}{0.657}  % p-value
\newcommand{\rhabLassoSharCloseLateQlike}{0.1142}  % QLIKE lasso shar c16_2020
\newcommand{\rhabLassoSharCloseLateMse}{3.95}  % MSE x 1e10 lasso shar c16_2020
\newcommand{\rhabLassoSharCloseLatePct}{0.4}  % QLIKE change % vs har
\newcommand{\rhabLassoSharCloseLateDm}{0.42}  % DM vs the reference arm
\newcommand{\rhabLassoSharCloseLateP}{0.673}  % p-value
\newcommand{\rhabLassoSharqlAllQlike}{0.1840}  % QLIKE lasso sharql all
\newcommand{\rhabLassoSharqlAllMse}{1.69}  % MSE x 1e10 lasso sharql all
\newcommand{\rhabLassoSharqlAllPct}{$-$2.8}  % QLIKE change % vs har
\newcommand{\rhabLassoSharqlAllDm}{$-$6.63}  % DM vs the reference arm
\newcommand{\rhabLassoSharqlAllP}{\ensuremath{<}0.001}  % p-value
\newcommand{\rhabLassoSharqlCloseQlike}{0.1218}  % QLIKE lasso sharql c16
\newcommand{\rhabLassoSharqlCloseMse}{2.68}  % MSE x 1e10 lasso sharql c16
\newcommand{\rhabLassoSharqlClosePct}{$-$4.0}  % QLIKE change % vs har
\newcommand{\rhabLassoSharqlCloseDm}{$-$3.14}  % DM vs the reference arm
\newcommand{\rhabLassoSharqlCloseP}{0.002}  % p-value
\newcommand{\rhabLassoSharqlCloseLateQlike}{0.1125}  % QLIKE lasso sharql c16_2020
\newcommand{\rhabLassoSharqlCloseLateMse}{5.98}  % MSE x 1e10 lasso sharql c16_2020
\newcommand{\rhabLassoSharqlCloseLatePct}{$-$1.0}  % QLIKE change % vs har
\newcommand{\rhabLassoSharqlCloseLateDm}{$-$0.44}  % DM vs the reference arm
\newcommand{\rhabLassoSharqlCloseLateP}{0.662}  % p-value
\newcommand{\rhabLgbmFullAllQlike}{0.1613}  % QLIKE lgbm full all
\newcommand{\rhabLgbmFullAllMse}{1.41}  % MSE x 1e10 lgbm full all
\newcommand{\rhabLgbmFullAllPct}{$-$14.0}  % QLIKE change % vs har
\newcommand{\rhabLgbmFullAllDm}{$-$10.32}  % DM vs the reference arm
\newcommand{\rhabLgbmFullAllP}{\ensuremath{<}0.001}  % p-value
\newcommand{\rhabLgbmFullCloseQlike}{0.1201}  % QLIKE lgbm full c16
\newcommand{\rhabLgbmFullCloseMse}{2.34}  % MSE x 1e10 lgbm full c16
\newcommand{\rhabLgbmFullClosePct}{$-$2.7}  % QLIKE change % vs har
\newcommand{\rhabLgbmFullCloseDm}{$-$1.81}  % DM vs the reference arm
\newcommand{\rhabLgbmFullCloseP}{0.071}  % p-value
\newcommand{\rhabLgbmFullCloseLateQlike}{0.1089}  % QLIKE lgbm full c16_2020
\newcommand{\rhabLgbmFullCloseLateMse}{5.09}  % MSE x 1e10 lgbm full c16_2020
\newcommand{\rhabLgbmFullCloseLatePct}{$-$1.1}  % QLIKE change % vs har
\newcommand{\rhabLgbmFullCloseLateDm}{$-$0.39}  % DM vs the reference arm
\newcommand{\rhabLgbmFullCloseLateP}{0.697}  % p-value
\newcommand{\rhabLgbmHarAllQlike}{0.1875}  % QLIKE lgbm har all
\newcommand{\rhabLgbmHarAllMse}{1.18}  % MSE x 1e10 lgbm har all
\newcommand{\rhabLgbmHarCloseQlike}{0.1235}  % QLIKE lgbm har c16
\newcommand{\rhabLgbmHarCloseMse}{1.77}  % MSE x 1e10 lgbm har c16
\newcommand{\rhabLgbmHarCloseLateQlike}{0.1101}  % QLIKE lgbm har c16_2020
\newcommand{\rhabLgbmHarCloseLateMse}{3.68}  % MSE x 1e10 lgbm har c16_2020
\newcommand{\rhabOlsAddRmAllQlike}{0.1858}  % QLIKE ols add_rm all
\newcommand{\rhabOlsAddRmAllMse}{1.79}  % MSE x 1e10 ols add_rm all
\newcommand{\rhabOlsAddRmAllPct}{$-$1.9}  % QLIKE change % vs har
\newcommand{\rhabOlsAddRmAllDm}{$-$2.00}  % DM vs the reference arm
\newcommand{\rhabOlsAddRmAllP}{0.045}  % p-value
\newcommand{\rhabOlsAddRmCloseQlike}{0.1252}  % QLIKE ols add_rm c16
\newcommand{\rhabOlsAddRmCloseMse}{3.02}  % MSE x 1e10 ols add_rm c16
\newcommand{\rhabOlsAddRmClosePct}{$-$1.3}  % QLIKE change % vs har
\newcommand{\rhabOlsAddRmCloseDm}{$-$0.56}  % DM vs the reference arm
\newcommand{\rhabOlsAddRmCloseP}{0.576}  % p-value
\newcommand{\rhabOlsAddRmCloseLateQlike}{0.1221}  % QLIKE ols add_rm c16_2020
\newcommand{\rhabOlsAddRmCloseLateMse}{6.80}  % MSE x 1e10 ols add_rm c16_2020
\newcommand{\rhabOlsAddRmCloseLatePct}{7.5}  % QLIKE change % vs har
\newcommand{\rhabOlsAddRmCloseLateDm}{1.43}  % DM vs the reference arm
\newcommand{\rhabOlsAddRmCloseLateP}{0.152}  % p-value
\newcommand{\rhabOlsHarAllQlike}{0.1893}  % QLIKE ols har all
\newcommand{\rhabOlsHarAllMse}{1.11}  % MSE x 1e10 ols har all
\newcommand{\rhabOlsHarCloseQlike}{0.1269}  % QLIKE ols har c16
\newcommand{\rhabOlsHarCloseMse}{1.97}  % MSE x 1e10 ols har c16
\newcommand{\rhabOlsHarCloseLateQlike}{0.1137}  % QLIKE ols har c16_2020
\newcommand{\rhabOlsHarCloseLateMse}{4.24}  % MSE x 1e10 ols har c16_2020
\newcommand{\rhabOlsHarqAllQlike}{0.1893}  % QLIKE ols harq all
\newcommand{\rhabOlsHarqAllMse}{1.17}  % MSE x 1e10 ols harq all
\newcommand{\rhabOlsHarqAllPct}{0.0}  % QLIKE change % vs har
\newcommand{\rhabOlsHarqAllDm}{0.61}  % DM vs the reference arm
\newcommand{\rhabOlsHarqAllP}{0.543}  % p-value
\newcommand{\rhabOlsHarqCloseQlike}{0.1267}  % QLIKE ols harq c16
\newcommand{\rhabOlsHarqCloseMse}{1.78}  % MSE x 1e10 ols harq c16
\newcommand{\rhabOlsHarqClosePct}{$-$0.2}  % QLIKE change % vs har
\newcommand{\rhabOlsHarqCloseDm}{$-$1.38}  % DM vs the reference arm
\newcommand{\rhabOlsHarqCloseP}{0.166}  % p-value
\newcommand{\rhabOlsHarqCloseLateQlike}{0.1133}  % QLIKE ols harq c16_2020
\newcommand{\rhabOlsHarqCloseLateMse}{3.80}  % MSE x 1e10 ols harq c16_2020
\newcommand{\rhabOlsHarqCloseLatePct}{$-$0.3}  % QLIKE change % vs har
\newcommand{\rhabOlsHarqCloseLateDm}{$-$1.19}  % DM vs the reference arm
\newcommand{\rhabOlsHarqCloseLateP}{0.235}  % p-value
\newcommand{\rhabOlsLharAllQlike}{0.1842}  % QLIKE ols lhar all
\newcommand{\rhabOlsLharAllMse}{1.71}  % MSE x 1e10 ols lhar all
\newcommand{\rhabOlsLharAllPct}{$-$2.7}  % QLIKE change % vs har
\newcommand{\rhabOlsLharAllDm}{$-$6.19}  % DM vs the reference arm
\newcommand{\rhabOlsLharAllP}{\ensuremath{<}0.001}  % p-value
\newcommand{\rhabOlsLharCloseQlike}{0.1220}  % QLIKE ols lhar c16
\newcommand{\rhabOlsLharCloseMse}{2.81}  % MSE x 1e10 ols lhar c16
\newcommand{\rhabOlsLharClosePct}{$-$3.9}  % QLIKE change % vs har
\newcommand{\rhabOlsLharCloseDm}{$-$2.66}  % DM vs the reference arm
\newcommand{\rhabOlsLharCloseP}{0.008}  % p-value
\newcommand{\rhabOlsLharCloseLateQlike}{0.1132}  % QLIKE ols lhar c16_2020
\newcommand{\rhabOlsLharCloseLateMse}{6.30}  % MSE x 1e10 ols lhar c16_2020
\newcommand{\rhabOlsLharCloseLatePct}{$-$0.4}  % QLIKE change % vs har
\newcommand{\rhabOlsLharCloseLateDm}{$-$0.15}  % DM vs the reference arm
\newcommand{\rhabOlsLharCloseLateP}{0.883}  % p-value
\newcommand{\rhabOlsSharAllQlike}{0.1877}  % QLIKE ols shar all
\newcommand{\rhabOlsSharAllMse}{1.40}  % MSE x 1e10 ols shar all
\newcommand{\rhabOlsSharAllPct}{$-$0.8}  % QLIKE change % vs har
\newcommand{\rhabOlsSharAllDm}{$-$3.48}  % DM vs the reference arm
\newcommand{\rhabOlsSharAllP}{\ensuremath{<}0.001}  % p-value
\newcommand{\rhabOlsSharCloseQlike}{0.1265}  % QLIKE ols shar c16
\newcommand{\rhabOlsSharCloseMse}{1.84}  % MSE x 1e10 ols shar c16
\newcommand{\rhabOlsSharClosePct}{$-$0.3}  % QLIKE change % vs har
\newcommand{\rhabOlsSharCloseDm}{$-$0.42}  % DM vs the reference arm
\newcommand{\rhabOlsSharCloseP}{0.674}  % p-value
\newcommand{\rhabOlsSharCloseLateQlike}{0.1142}  % QLIKE ols shar c16_2020
\newcommand{\rhabOlsSharCloseLateMse}{3.95}  % MSE x 1e10 ols shar c16_2020
\newcommand{\rhabOlsSharCloseLatePct}{0.5}  % QLIKE change % vs har
\newcommand{\rhabOlsSharCloseLateDm}{0.44}  % DM vs the reference arm
\newcommand{\rhabOlsSharCloseLateP}{0.659}  % p-value
\newcommand{\rhabOlsSharqlAllQlike}{0.1841}  % QLIKE ols sharql all
\newcommand{\rhabOlsSharqlAllMse}{1.69}  % MSE x 1e10 ols sharql all
\newcommand{\rhabOlsSharqlAllPct}{$-$2.7}  % QLIKE change % vs har
\newcommand{\rhabOlsSharqlAllDm}{$-$6.53}  % DM vs the reference arm
\newcommand{\rhabOlsSharqlAllP}{\ensuremath{<}0.001}  % p-value
\newcommand{\rhabOlsSharqlCloseQlike}{0.1218}  % QLIKE ols sharql c16
\newcommand{\rhabOlsSharqlCloseMse}{2.68}  % MSE x 1e10 ols sharql c16
\newcommand{\rhabOlsSharqlClosePct}{$-$4.0}  % QLIKE change % vs har
\newcommand{\rhabOlsSharqlCloseDm}{$-$3.05}  % DM vs the reference arm
\newcommand{\rhabOlsSharqlCloseP}{0.002}  % p-value
\newcommand{\rhabOlsSharqlCloseLateQlike}{0.1126}  % QLIKE ols sharql c16_2020
\newcommand{\rhabOlsSharqlCloseLateMse}{5.99}  % MSE x 1e10 ols sharql c16_2020
\newcommand{\rhabOlsSharqlCloseLatePct}{$-$0.9}  % QLIKE change % vs har
\newcommand{\rhabOlsSharqlCloseLateDm}{$-$0.38}  % DM vs the reference arm
\newcommand{\rhabOlsSharqlCloseLateP}{0.703}  % p-value
\newcommand{\rhabRfFullAllQlike}{0.1685}  % QLIKE rf full all
\newcommand{\rhabRfFullAllMse}{1.32}  % MSE x 1e10 rf full all
\newcommand{\rhabRfFullAllPct}{$-$11.1}  % QLIKE change % vs har
\newcommand{\rhabRfFullAllDm}{$-$9.85}  % DM vs the reference arm
\newcommand{\rhabRfFullAllP}{\ensuremath{<}0.001}  % p-value
\newcommand{\rhabRfFullCloseQlike}{0.1248}  % QLIKE rf full c16
\newcommand{\rhabRfFullCloseMse}{1.99}  % MSE x 1e10 rf full c16
\newcommand{\rhabRfFullClosePct}{$-$1.4}  % QLIKE change % vs har
\newcommand{\rhabRfFullCloseDm}{$-$0.92}  % DM vs the reference arm
\newcommand{\rhabRfFullCloseP}{0.357}  % p-value
\newcommand{\rhabRfFullCloseLateQlike}{0.1123}  % QLIKE rf full c16_2020
\newcommand{\rhabRfFullCloseLateMse}{4.26}  % MSE x 1e10 rf full c16_2020
\newcommand{\rhabRfFullCloseLatePct}{$-$0.7}  % QLIKE change % vs har
\newcommand{\rhabRfFullCloseLateDm}{$-$0.26}  % DM vs the reference arm
\newcommand{\rhabRfFullCloseLateP}{0.794}  % p-value
\newcommand{\rhabRfHarAllQlike}{0.1895}  % QLIKE rf har all
\newcommand{\rhabRfHarAllMse}{1.19}  % MSE x 1e10 rf har all
\newcommand{\rhabRfHarCloseQlike}{0.1266}  % QLIKE rf har c16
\newcommand{\rhabRfHarCloseMse}{1.63}  % MSE x 1e10 rf har c16
\newcommand{\rhabRfHarCloseLateQlike}{0.1131}  % QLIKE rf har c16_2020
\newcommand{\rhabRfHarCloseLateMse}{3.39}  % MSE x 1e10 rf har c16_2020
\newcommand{\rhabRidgeAddEwAllQlike}{0.1876}  % QLIKE ridge add_ew all
\newcommand{\rhabRidgeAddEwAllMse}{1.42}  % MSE x 1e10 ridge add_ew all
\newcommand{\rhabRidgeAddEwAllPct}{$-$0.9}  % QLIKE change % vs har
\newcommand{\rhabRidgeAddEwAllDm}{$-$2.26}  % DM vs the reference arm
\newcommand{\rhabRidgeAddEwAllP}{0.024}  % p-value
\newcommand{\rhabRidgeAddEwCloseQlike}{0.1256}  % QLIKE ridge add_ew c16
\newcommand{\rhabRidgeAddEwCloseMse}{2.08}  % MSE x 1e10 ridge add_ew c16
\newcommand{\rhabRidgeAddEwClosePct}{$-$1.0}  % QLIKE change % vs har
\newcommand{\rhabRidgeAddEwCloseDm}{$-$0.89}  % DM vs the reference arm
\newcommand{\rhabRidgeAddEwCloseP}{0.374}  % p-value
\newcommand{\rhabRidgeAddEwCloseLateQlike}{0.1125}  % QLIKE ridge add_ew c16_2020
\newcommand{\rhabRidgeAddEwCloseLateMse}{4.51}  % MSE x 1e10 ridge add_ew c16_2020
\newcommand{\rhabRidgeAddEwCloseLatePct}{$-$1.0}  % QLIKE change % vs har
\newcommand{\rhabRidgeAddEwCloseLateDm}{$-$0.62}  % DM vs the reference arm
\newcommand{\rhabRidgeAddEwCloseLateP}{0.537}  % p-value
\newcommand{\rhabRidgeAddFomcAllQlike}{0.1847}  % QLIKE ridge add_fomc all
\newcommand{\rhabRidgeAddFomcAllMse}{1.13}  % MSE x 1e10 ridge add_fomc all
\newcommand{\rhabRidgeAddFomcAllPct}{$-$2.4}  % QLIKE change % vs har
\newcommand{\rhabRidgeAddFomcAllDm}{$-$5.79}  % DM vs the reference arm
\newcommand{\rhabRidgeAddFomcAllP}{\ensuremath{<}0.001}  % p-value
\newcommand{\rhabRidgeAddFomcCloseQlike}{0.1276}  % QLIKE ridge add_fomc c16
\newcommand{\rhabRidgeAddFomcCloseMse}{1.84}  % MSE x 1e10 ridge add_fomc c16
\newcommand{\rhabRidgeAddFomcClosePct}{0.6}  % QLIKE change % vs har
\newcommand{\rhabRidgeAddFomcCloseDm}{2.17}  % DM vs the reference arm
\newcommand{\rhabRidgeAddFomcCloseP}{0.030}  % p-value
\newcommand{\rhabRidgeAddFomcCloseLateQlike}{0.1135}  % QLIKE ridge add_fomc c16_2020
\newcommand{\rhabRidgeAddFomcCloseLateMse}{3.94}  % MSE x 1e10 ridge add_fomc c16_2020
\newcommand{\rhabRidgeAddFomcCloseLatePct}{$-$0.2}  % QLIKE change % vs har
\newcommand{\rhabRidgeAddFomcCloseLateDm}{$-$0.35}  % DM vs the reference arm
\newcommand{\rhabRidgeAddFomcCloseLateP}{0.729}  % p-value
\newcommand{\rhabRidgeAddIvAllQlike}{0.1897}  % QLIKE ridge add_iv all
\newcommand{\rhabRidgeAddIvAllMse}{1.28}  % MSE x 1e10 ridge add_iv all
\newcommand{\rhabRidgeAddIvAllPct}{0.2}  % QLIKE change % vs har
\newcommand{\rhabRidgeAddIvAllDm}{0.67}  % DM vs the reference arm
\newcommand{\rhabRidgeAddIvAllP}{0.506}  % p-value
\newcommand{\rhabRidgeAddIvCloseQlike}{0.1266}  % QLIKE ridge add_iv c16
\newcommand{\rhabRidgeAddIvCloseMse}{1.79}  % MSE x 1e10 ridge add_iv c16
\newcommand{\rhabRidgeAddIvClosePct}{$-$0.2}  % QLIKE change % vs har
\newcommand{\rhabRidgeAddIvCloseDm}{$-$0.20}  % DM vs the reference arm
\newcommand{\rhabRidgeAddIvCloseP}{0.845}  % p-value
\newcommand{\rhabRidgeAddIvCloseLateQlike}{0.1153}  % QLIKE ridge add_iv c16_2020
\newcommand{\rhabRidgeAddIvCloseLateMse}{3.82}  % MSE x 1e10 ridge add_iv c16_2020
\newcommand{\rhabRidgeAddIvCloseLatePct}{1.4}  % QLIKE change % vs har
\newcommand{\rhabRidgeAddIvCloseLateDm}{0.70}  % DM vs the reference arm
\newcommand{\rhabRidgeAddIvCloseLateP}{0.482}  % p-value
\newcommand{\rhabRidgeAddLiqAllQlike}{0.1885}  % QLIKE ridge add_liq all
\newcommand{\rhabRidgeAddLiqAllMse}{1.18}  % MSE x 1e10 ridge add_liq all
\newcommand{\rhabRidgeAddLiqAllPct}{$-$0.4}  % QLIKE change % vs har
\newcommand{\rhabRidgeAddLiqAllDm}{$-$0.98}  % DM vs the reference arm
\newcommand{\rhabRidgeAddLiqAllP}{0.326}  % p-value
\newcommand{\rhabRidgeAddLiqCloseQlike}{0.1304}  % QLIKE ridge add_liq c16
\newcommand{\rhabRidgeAddLiqCloseMse}{2.01}  % MSE x 1e10 ridge add_liq c16
\newcommand{\rhabRidgeAddLiqClosePct}{2.8}  % QLIKE change % vs har
\newcommand{\rhabRidgeAddLiqCloseDm}{1.91}  % DM vs the reference arm
\newcommand{\rhabRidgeAddLiqCloseP}{0.056}  % p-value
\newcommand{\rhabRidgeAddLiqCloseLateQlike}{0.1204}  % QLIKE ridge add_liq c16_2020
\newcommand{\rhabRidgeAddLiqCloseLateMse}{4.35}  % MSE x 1e10 ridge add_liq c16_2020
\newcommand{\rhabRidgeAddLiqCloseLatePct}{6.0}  % QLIKE change % vs har
\newcommand{\rhabRidgeAddLiqCloseLateDm}{2.09}  % DM vs the reference arm
\newcommand{\rhabRidgeAddLiqCloseLateP}{0.036}  % p-value
\newcommand{\rhabRidgeAddRmAllQlike}{0.1843}  % QLIKE ridge add_rm all
\newcommand{\rhabRidgeAddRmAllMse}{1.69}  % MSE x 1e10 ridge add_rm all
\newcommand{\rhabRidgeAddRmAllPct}{$-$2.6}  % QLIKE change % vs har
\newcommand{\rhabRidgeAddRmAllDm}{$-$4.14}  % DM vs the reference arm
\newcommand{\rhabRidgeAddRmAllP}{\ensuremath{<}0.001}  % p-value
\newcommand{\rhabRidgeAddRmCloseQlike}{0.1230}  % QLIKE ridge add_rm c16
\newcommand{\rhabRidgeAddRmCloseMse}{2.81}  % MSE x 1e10 ridge add_rm c16
\newcommand{\rhabRidgeAddRmClosePct}{$-$3.1}  % QLIKE change % vs har
\newcommand{\rhabRidgeAddRmCloseDm}{$-$1.80}  % DM vs the reference arm
\newcommand{\rhabRidgeAddRmCloseP}{0.071}  % p-value
\newcommand{\rhabRidgeAddRmCloseLateQlike}{0.1156}  % QLIKE ridge add_rm c16_2020
\newcommand{\rhabRidgeAddRmCloseLateMse}{6.29}  % MSE x 1e10 ridge add_rm c16_2020
\newcommand{\rhabRidgeAddRmCloseLatePct}{1.7}  % QLIKE change % vs har
\newcommand{\rhabRidgeAddRmCloseLateDm}{0.56}  % DM vs the reference arm
\newcommand{\rhabRidgeAddRmCloseLateP}{0.577}  % p-value
\newcommand{\rhabRidgeAddSentAllQlike}{0.1897}  % QLIKE ridge add_sent all
\newcommand{\rhabRidgeAddSentAllMse}{1.11}  % MSE x 1e10 ridge add_sent all
\newcommand{\rhabRidgeAddSentAllPct}{0.2}  % QLIKE change % vs har
\newcommand{\rhabRidgeAddSentAllDm}{0.82}  % DM vs the reference arm
\newcommand{\rhabRidgeAddSentAllP}{0.411}  % p-value
\newcommand{\rhabRidgeAddSentCloseQlike}{0.1310}  % QLIKE ridge add_sent c16
\newcommand{\rhabRidgeAddSentCloseMse}{1.98}  % MSE x 1e10 ridge add_sent c16
\newcommand{\rhabRidgeAddSentClosePct}{3.2}  % QLIKE change % vs har
\newcommand{\rhabRidgeAddSentCloseDm}{4.41}  % DM vs the reference arm
\newcommand{\rhabRidgeAddSentCloseP}{\ensuremath{<}0.001}  % p-value
\newcommand{\rhabRidgeAddSentCloseLateQlike}{0.1214}  % QLIKE ridge add_sent c16_2020
\newcommand{\rhabRidgeAddSentCloseLateMse}{4.27}  % MSE x 1e10 ridge add_sent c16_2020
\newcommand{\rhabRidgeAddSentCloseLatePct}{6.8}  % QLIKE change % vs har
\newcommand{\rhabRidgeAddSentCloseLateDm}{3.65}  % DM vs the reference arm
\newcommand{\rhabRidgeAddSentCloseLateP}{\ensuremath{<}0.001}  % p-value
\newcommand{\rhabRidgeAddVdAllQlike}{0.1903}  % QLIKE ridge add_vd all
\newcommand{\rhabRidgeAddVdAllMse}{1.11}  % MSE x 1e10 ridge add_vd all
\newcommand{\rhabRidgeAddVdAllPct}{0.5}  % QLIKE change % vs har
\newcommand{\rhabRidgeAddVdAllDm}{3.10}  % DM vs the reference arm
\newcommand{\rhabRidgeAddVdAllP}{0.002}  % p-value
\newcommand{\rhabRidgeAddVdCloseQlike}{0.1276}  % QLIKE ridge add_vd c16
\newcommand{\rhabRidgeAddVdCloseMse}{1.95}  % MSE x 1e10 ridge add_vd c16
\newcommand{\rhabRidgeAddVdClosePct}{0.6}  % QLIKE change % vs har
\newcommand{\rhabRidgeAddVdCloseDm}{1.55}  % DM vs the reference arm
\newcommand{\rhabRidgeAddVdCloseP}{0.120}  % p-value
\newcommand{\rhabRidgeAddVdCloseLateQlike}{0.1169}  % QLIKE ridge add_vd c16_2020
\newcommand{\rhabRidgeAddVdCloseLateMse}{4.21}  % MSE x 1e10 ridge add_vd c16_2020
\newcommand{\rhabRidgeAddVdCloseLatePct}{2.8}  % QLIKE change % vs har
\newcommand{\rhabRidgeAddVdCloseLateDm}{3.10}  % DM vs the reference arm
\newcommand{\rhabRidgeAddVdCloseLateP}{0.002}  % p-value
\newcommand{\rhabRidgeAddVwAllQlike}{0.1876}  % QLIKE ridge add_vw all
\newcommand{\rhabRidgeAddVwAllMse}{1.87}  % MSE x 1e10 ridge add_vw all
\newcommand{\rhabRidgeAddVwAllPct}{$-$0.9}  % QLIKE change % vs har
\newcommand{\rhabRidgeAddVwAllDm}{$-$1.58}  % DM vs the reference arm
\newcommand{\rhabRidgeAddVwAllP}{0.114}  % p-value
\newcommand{\rhabRidgeAddVwCloseQlike}{0.1264}  % QLIKE ridge add_vw c16
\newcommand{\rhabRidgeAddVwCloseMse}{2.27}  % MSE x 1e10 ridge add_vw c16
\newcommand{\rhabRidgeAddVwClosePct}{$-$0.4}  % QLIKE change % vs har
\newcommand{\rhabRidgeAddVwCloseDm}{$-$0.28}  % DM vs the reference arm
\newcommand{\rhabRidgeAddVwCloseP}{0.781}  % p-value
\newcommand{\rhabRidgeAddVwCloseLateQlike}{0.1139}  % QLIKE ridge add_vw c16_2020
\newcommand{\rhabRidgeAddVwCloseLateMse}{5.00}  % MSE x 1e10 ridge add_vw c16_2020
\newcommand{\rhabRidgeAddVwCloseLatePct}{0.2}  % QLIKE change % vs har
\newcommand{\rhabRidgeAddVwCloseLateDm}{0.12}  % DM vs the reference arm
\newcommand{\rhabRidgeAddVwCloseLateP}{0.907}  % p-value
\newcommand{\rhabRidgeFullAllQlike}{0.1791}  % QLIKE ridge full all
\newcommand{\rhabRidgeFullAllMse}{2.24}  % MSE x 1e10 ridge full all
\newcommand{\rhabRidgeFullAllPct}{$-$5.4}  % QLIKE change % vs har
\newcommand{\rhabRidgeFullAllDm}{$-$5.63}  % DM vs the reference arm
\newcommand{\rhabRidgeFullAllP}{\ensuremath{<}0.001}  % p-value
\newcommand{\rhabRidgeFullCloseQlike}{0.1279}  % QLIKE ridge full c16
\newcommand{\rhabRidgeFullCloseMse}{3.52}  % MSE x 1e10 ridge full c16
\newcommand{\rhabRidgeFullClosePct}{0.8}  % QLIKE change % vs har
\newcommand{\rhabRidgeFullCloseDm}{0.31}  % DM vs the reference arm
\newcommand{\rhabRidgeFullCloseP}{0.759}  % p-value
\newcommand{\rhabRidgeFullCloseLateQlike}{0.1260}  % QLIKE ridge full c16_2020
\newcommand{\rhabRidgeFullCloseLateMse}{7.99}  % MSE x 1e10 ridge full c16_2020
\newcommand{\rhabRidgeFullCloseLatePct}{10.9}  % QLIKE change % vs har
\newcommand{\rhabRidgeFullCloseLateDm}{2.31}  % DM vs the reference arm
\newcommand{\rhabRidgeFullCloseLateP}{0.021}  % p-value
\newcommand{\rhabRidgeHarAllQlike}{0.1893}  % QLIKE ridge har all
\newcommand{\rhabRidgeHarAllMse}{1.11}  % MSE x 1e10 ridge har all
\newcommand{\rhabRidgeHarCloseQlike}{0.1269}  % QLIKE ridge har c16
\newcommand{\rhabRidgeHarCloseMse}{1.97}  % MSE x 1e10 ridge har c16
\newcommand{\rhabRidgeHarCloseLateQlike}{0.1137}  % QLIKE ridge har c16_2020
\newcommand{\rhabRidgeHarCloseLateMse}{4.24}  % MSE x 1e10 ridge har c16_2020
\newcommand{\rhabRidgeHarqAllQlike}{0.1893}  % QLIKE ridge harq all
\newcommand{\rhabRidgeHarqAllMse}{1.16}  % MSE x 1e10 ridge harq all
\newcommand{\rhabRidgeHarqAllPct}{0.0}  % QLIKE change % vs har
\newcommand{\rhabRidgeHarqAllDm}{0.76}  % DM vs the reference arm
\newcommand{\rhabRidgeHarqAllP}{0.446}  % p-value
\newcommand{\rhabRidgeHarqCloseQlike}{0.1267}  % QLIKE ridge harq c16
\newcommand{\rhabRidgeHarqCloseMse}{1.79}  % MSE x 1e10 ridge harq c16
\newcommand{\rhabRidgeHarqClosePct}{$-$0.2}  % QLIKE change % vs har
\newcommand{\rhabRidgeHarqCloseDm}{$-$1.42}  % DM vs the reference arm
\newcommand{\rhabRidgeHarqCloseP}{0.157}  % p-value
\newcommand{\rhabRidgeHarqCloseLateQlike}{0.1133}  % QLIKE ridge harq c16_2020
\newcommand{\rhabRidgeHarqCloseLateMse}{3.81}  % MSE x 1e10 ridge harq c16_2020
\newcommand{\rhabRidgeHarqCloseLatePct}{$-$0.3}  % QLIKE change % vs har
\newcommand{\rhabRidgeHarqCloseLateDm}{$-$1.20}  % DM vs the reference arm
\newcommand{\rhabRidgeHarqCloseLateP}{0.228}  % p-value
\newcommand{\rhabRidgeLharAllQlike}{0.1841}  % QLIKE ridge lhar all
\newcommand{\rhabRidgeLharAllMse}{1.72}  % MSE x 1e10 ridge lhar all
\newcommand{\rhabRidgeLharAllPct}{$-$2.7}  % QLIKE change % vs har
\newcommand{\rhabRidgeLharAllDm}{$-$6.56}  % DM vs the reference arm
\newcommand{\rhabRidgeLharAllP}{\ensuremath{<}0.001}  % p-value
\newcommand{\rhabRidgeLharCloseQlike}{0.1216}  % QLIKE ridge lhar c16
\newcommand{\rhabRidgeLharCloseMse}{2.78}  % MSE x 1e10 ridge lhar c16
\newcommand{\rhabRidgeLharClosePct}{$-$4.2}  % QLIKE change % vs har
\newcommand{\rhabRidgeLharCloseDm}{$-$3.28}  % DM vs the reference arm
\newcommand{\rhabRidgeLharCloseP}{0.001}  % p-value
\newcommand{\rhabRidgeLharCloseLateQlike}{0.1121}  % QLIKE ridge lhar c16_2020
\newcommand{\rhabRidgeLharCloseLateMse}{6.23}  % MSE x 1e10 ridge lhar c16_2020
\newcommand{\rhabRidgeLharCloseLatePct}{$-$1.4}  % QLIKE change % vs har
\newcommand{\rhabRidgeLharCloseLateDm}{$-$0.59}  % DM vs the reference arm
\newcommand{\rhabRidgeLharCloseLateP}{0.554}  % p-value
\newcommand{\rhabRidgeLoEwAllQlike}{0.1789}  % QLIKE ridge lo_ew all
\newcommand{\rhabRidgeLoEwAllMse}{2.38}  % MSE x 1e10 ridge lo_ew all
\newcommand{\rhabRidgeLoEwAllPct}{$-$0.1}  % QLIKE change % vs full
\newcommand{\rhabRidgeLoEwAllDm}{$-$1.20}  % DM vs the reference arm
\newcommand{\rhabRidgeLoEwAllP}{0.230}  % p-value
\newcommand{\rhabRidgeLoEwCloseQlike}{0.1285}  % QLIKE ridge lo_ew c16
\newcommand{\rhabRidgeLoEwCloseMse}{3.74}  % MSE x 1e10 ridge lo_ew c16
\newcommand{\rhabRidgeLoEwClosePct}{0.5}  % QLIKE change % vs full
\newcommand{\rhabRidgeLoEwCloseDm}{1.23}  % DM vs the reference arm
\newcommand{\rhabRidgeLoEwCloseP}{0.218}  % p-value
\newcommand{\rhabRidgeLoEwCloseLateQlike}{0.1262}  % QLIKE ridge lo_ew c16_2020
\newcommand{\rhabRidgeLoEwCloseLateMse}{8.53}  % MSE x 1e10 ridge lo_ew c16_2020
\newcommand{\rhabRidgeLoEwCloseLatePct}{0.2}  % QLIKE change % vs full
\newcommand{\rhabRidgeLoEwCloseLateDm}{0.22}  % DM vs the reference arm
\newcommand{\rhabRidgeLoEwCloseLateP}{0.828}  % p-value
\newcommand{\rhabRidgeLoFomcAllQlike}{0.1833}  % QLIKE ridge lo_fomc all
\newcommand{\rhabRidgeLoFomcAllMse}{2.25}  % MSE x 1e10 ridge lo_fomc all
\newcommand{\rhabRidgeLoFomcAllPct}{2.3}  % QLIKE change % vs full
\newcommand{\rhabRidgeLoFomcAllDm}{5.55}  % DM vs the reference arm
\newcommand{\rhabRidgeLoFomcAllP}{\ensuremath{<}0.001}  % p-value
\newcommand{\rhabRidgeLoFomcCloseQlike}{0.1275}  % QLIKE ridge lo_fomc c16
\newcommand{\rhabRidgeLoFomcCloseMse}{3.58}  % MSE x 1e10 ridge lo_fomc c16
\newcommand{\rhabRidgeLoFomcClosePct}{$-$0.3}  % QLIKE change % vs full
\newcommand{\rhabRidgeLoFomcCloseDm}{$-$1.23}  % DM vs the reference arm
\newcommand{\rhabRidgeLoFomcCloseP}{0.219}  % p-value
\newcommand{\rhabRidgeLoFomcCloseLateQlike}{0.1257}  % QLIKE ridge lo_fomc c16_2020
\newcommand{\rhabRidgeLoFomcCloseLateMse}{8.15}  % MSE x 1e10 ridge lo_fomc c16_2020
\newcommand{\rhabRidgeLoFomcCloseLatePct}{$-$0.2}  % QLIKE change % vs full
\newcommand{\rhabRidgeLoFomcCloseLateDm}{$-$0.49}  % DM vs the reference arm
\newcommand{\rhabRidgeLoFomcCloseLateP}{0.622}  % p-value
\newcommand{\rhabRidgeLoIvAllQlike}{0.1791}  % QLIKE ridge lo_iv all
\newcommand{\rhabRidgeLoIvAllMse}{2.30}  % MSE x 1e10 ridge lo_iv all
\newcommand{\rhabRidgeLoIvAllPct}{0.0}  % QLIKE change % vs full
\newcommand{\rhabRidgeLoIvAllDm}{0.03}  % DM vs the reference arm
\newcommand{\rhabRidgeLoIvAllP}{0.980}  % p-value
\newcommand{\rhabRidgeLoIvCloseQlike}{0.1267}  % QLIKE ridge lo_iv c16
\newcommand{\rhabRidgeLoIvCloseMse}{3.61}  % MSE x 1e10 ridge lo_iv c16
\newcommand{\rhabRidgeLoIvClosePct}{$-$1.0}  % QLIKE change % vs full
\newcommand{\rhabRidgeLoIvCloseDm}{$-$1.02}  % DM vs the reference arm
\newcommand{\rhabRidgeLoIvCloseP}{0.309}  % p-value
\newcommand{\rhabRidgeLoIvCloseLateQlike}{0.1228}  % QLIKE ridge lo_iv c16_2020
\newcommand{\rhabRidgeLoIvCloseLateMse}{8.20}  % MSE x 1e10 ridge lo_iv c16_2020
\newcommand{\rhabRidgeLoIvCloseLatePct}{$-$2.5}  % QLIKE change % vs full
\newcommand{\rhabRidgeLoIvCloseLateDm}{$-$1.41}  % DM vs the reference arm
\newcommand{\rhabRidgeLoIvCloseLateP}{0.159}  % p-value
\newcommand{\rhabRidgeLoLiqAllQlike}{0.1794}  % QLIKE ridge lo_liq all
\newcommand{\rhabRidgeLoLiqAllMse}{2.29}  % MSE x 1e10 ridge lo_liq all
\newcommand{\rhabRidgeLoLiqAllPct}{0.2}  % QLIKE change % vs full
\newcommand{\rhabRidgeLoLiqAllDm}{0.75}  % DM vs the reference arm
\newcommand{\rhabRidgeLoLiqAllP}{0.454}  % p-value
\newcommand{\rhabRidgeLoLiqCloseQlike}{0.1269}  % QLIKE ridge lo_liq c16
\newcommand{\rhabRidgeLoLiqCloseMse}{3.93}  % MSE x 1e10 ridge lo_liq c16
\newcommand{\rhabRidgeLoLiqClosePct}{$-$0.8}  % QLIKE change % vs full
\newcommand{\rhabRidgeLoLiqCloseDm}{$-$1.14}  % DM vs the reference arm
\newcommand{\rhabRidgeLoLiqCloseP}{0.256}  % p-value
\newcommand{\rhabRidgeLoLiqCloseLateQlike}{0.1251}  % QLIKE ridge lo_liq c16_2020
\newcommand{\rhabRidgeLoLiqCloseLateMse}{8.99}  % MSE x 1e10 ridge lo_liq c16_2020
\newcommand{\rhabRidgeLoLiqCloseLatePct}{$-$0.7}  % QLIKE change % vs full
\newcommand{\rhabRidgeLoLiqCloseLateDm}{$-$0.77}  % DM vs the reference arm
\newcommand{\rhabRidgeLoLiqCloseLateP}{0.442}  % p-value
\newcommand{\rhabRidgeLoRmAllQlike}{0.1822}  % QLIKE ridge lo_rm all
\newcommand{\rhabRidgeLoRmAllMse}{1.82}  % MSE x 1e10 ridge lo_rm all
\newcommand{\rhabRidgeLoRmAllPct}{1.7}  % QLIKE change % vs full
\newcommand{\rhabRidgeLoRmAllDm}{3.75}  % DM vs the reference arm
\newcommand{\rhabRidgeLoRmAllP}{\ensuremath{<}0.001}  % p-value
\newcommand{\rhabRidgeLoRmCloseQlike}{0.1292}  % QLIKE ridge lo_rm c16
\newcommand{\rhabRidgeLoRmCloseMse}{2.26}  % MSE x 1e10 ridge lo_rm c16
\newcommand{\rhabRidgeLoRmClosePct}{1.0}  % QLIKE change % vs full
\newcommand{\rhabRidgeLoRmCloseDm}{0.73}  % DM vs the reference arm
\newcommand{\rhabRidgeLoRmCloseP}{0.466}  % p-value
\newcommand{\rhabRidgeLoRmCloseLateQlike}{0.1245}  % QLIKE ridge lo_rm c16_2020
\newcommand{\rhabRidgeLoRmCloseLateMse}{4.97}  % MSE x 1e10 ridge lo_rm c16_2020
\newcommand{\rhabRidgeLoRmCloseLatePct}{$-$1.2}  % QLIKE change % vs full
\newcommand{\rhabRidgeLoRmCloseLateDm}{$-$0.55}  % DM vs the reference arm
\newcommand{\rhabRidgeLoRmCloseLateP}{0.579}  % p-value
\newcommand{\rhabRidgeLoSentAllQlike}{0.1793}  % QLIKE ridge lo_sent all
\newcommand{\rhabRidgeLoSentAllMse}{2.24}  % MSE x 1e10 ridge lo_sent all
\newcommand{\rhabRidgeLoSentAllPct}{0.1}  % QLIKE change % vs full
\newcommand{\rhabRidgeLoSentAllDm}{0.55}  % DM vs the reference arm
\newcommand{\rhabRidgeLoSentAllP}{0.584}  % p-value
\newcommand{\rhabRidgeLoSentCloseQlike}{0.1261}  % QLIKE ridge lo_sent c16
\newcommand{\rhabRidgeLoSentCloseMse}{3.51}  % MSE x 1e10 ridge lo_sent c16
\newcommand{\rhabRidgeLoSentClosePct}{$-$1.4}  % QLIKE change % vs full
\newcommand{\rhabRidgeLoSentCloseDm}{$-$3.16}  % DM vs the reference arm
\newcommand{\rhabRidgeLoSentCloseP}{0.002}  % p-value
\newcommand{\rhabRidgeLoSentCloseLateQlike}{0.1218}  % QLIKE ridge lo_sent c16_2020
\newcommand{\rhabRidgeLoSentCloseLateMse}{7.97}  % MSE x 1e10 ridge lo_sent c16_2020
\newcommand{\rhabRidgeLoSentCloseLatePct}{$-$3.3}  % QLIKE change % vs full
\newcommand{\rhabRidgeLoSentCloseLateDm}{$-$3.21}  % DM vs the reference arm
\newcommand{\rhabRidgeLoSentCloseLateP}{0.001}  % p-value
\newcommand{\rhabRidgeLoVdAllQlike}{0.1789}  % QLIKE ridge lo_vd all
\newcommand{\rhabRidgeLoVdAllMse}{2.24}  % MSE x 1e10 ridge lo_vd all
\newcommand{\rhabRidgeLoVdAllPct}{$-$0.1}  % QLIKE change % vs full
\newcommand{\rhabRidgeLoVdAllDm}{$-$0.98}  % DM vs the reference arm
\newcommand{\rhabRidgeLoVdAllP}{0.326}  % p-value
\newcommand{\rhabRidgeLoVdCloseQlike}{0.1266}  % QLIKE ridge lo_vd c16
\newcommand{\rhabRidgeLoVdCloseMse}{3.51}  % MSE x 1e10 ridge lo_vd c16
\newcommand{\rhabRidgeLoVdClosePct}{$-$1.0}  % QLIKE change % vs full
\newcommand{\rhabRidgeLoVdCloseDm}{$-$2.17}  % DM vs the reference arm
\newcommand{\rhabRidgeLoVdCloseP}{0.030}  % p-value
\newcommand{\rhabRidgeLoVdCloseLateQlike}{0.1232}  % QLIKE ridge lo_vd c16_2020
\newcommand{\rhabRidgeLoVdCloseLateMse}{7.98}  % MSE x 1e10 ridge lo_vd c16_2020
\newcommand{\rhabRidgeLoVdCloseLatePct}{$-$2.2}  % QLIKE change % vs full
\newcommand{\rhabRidgeLoVdCloseLateDm}{$-$2.26}  % DM vs the reference arm
\newcommand{\rhabRidgeLoVdCloseLateP}{0.024}  % p-value
\newcommand{\rhabRidgeLoVwAllQlike}{0.1789}  % QLIKE ridge lo_vw all
\newcommand{\rhabRidgeLoVwAllMse}{1.97}  % MSE x 1e10 ridge lo_vw all
\newcommand{\rhabRidgeLoVwAllPct}{$-$0.1}  % QLIKE change % vs full
\newcommand{\rhabRidgeLoVwAllDm}{$-$0.70}  % DM vs the reference arm
\newcommand{\rhabRidgeLoVwAllP}{0.484}  % p-value
\newcommand{\rhabRidgeLoVwCloseQlike}{0.1274}  % QLIKE ridge lo_vw c16
\newcommand{\rhabRidgeLoVwCloseMse}{3.50}  % MSE x 1e10 ridge lo_vw c16
\newcommand{\rhabRidgeLoVwClosePct}{$-$0.4}  % QLIKE change % vs full
\newcommand{\rhabRidgeLoVwCloseDm}{$-$0.90}  % DM vs the reference arm
\newcommand{\rhabRidgeLoVwCloseP}{0.366}  % p-value
\newcommand{\rhabRidgeLoVwCloseLateQlike}{0.1246}  % QLIKE ridge lo_vw c16_2020
\newcommand{\rhabRidgeLoVwCloseLateMse}{7.95}  % MSE x 1e10 ridge lo_vw c16_2020
\newcommand{\rhabRidgeLoVwCloseLatePct}{$-$1.1}  % QLIKE change % vs full
\newcommand{\rhabRidgeLoVwCloseLateDm}{$-$1.96}  % DM vs the reference arm
\newcommand{\rhabRidgeLoVwCloseLateP}{0.050}  % p-value
\newcommand{\rhabRidgeSharAllQlike}{0.1876}  % QLIKE ridge shar all
\newcommand{\rhabRidgeSharAllMse}{1.35}  % MSE x 1e10 ridge shar all
\newcommand{\rhabRidgeSharAllPct}{$-$0.9}  % QLIKE change % vs har
\newcommand{\rhabRidgeSharAllDm}{$-$3.95}  % DM vs the reference arm
\newcommand{\rhabRidgeSharAllP}{\ensuremath{<}0.001}  % p-value
\newcommand{\rhabRidgeSharCloseQlike}{0.1264}  % QLIKE ridge shar c16
\newcommand{\rhabRidgeSharCloseMse}{1.81}  % MSE x 1e10 ridge shar c16
\newcommand{\rhabRidgeSharClosePct}{$-$0.4}  % QLIKE change % vs har
\newcommand{\rhabRidgeSharCloseDm}{$-$0.69}  % DM vs the reference arm
\newcommand{\rhabRidgeSharCloseP}{0.488}  % p-value
\newcommand{\rhabRidgeSharCloseLateQlike}{0.1139}  % QLIKE ridge shar c16_2020
\newcommand{\rhabRidgeSharCloseLateMse}{3.87}  % MSE x 1e10 ridge shar c16_2020
\newcommand{\rhabRidgeSharCloseLatePct}{0.2}  % QLIKE change % vs har
\newcommand{\rhabRidgeSharCloseLateDm}{0.23}  % DM vs the reference arm
\newcommand{\rhabRidgeSharCloseLateP}{0.821}  % p-value
\newcommand{\rhabRidgeSharqlAllQlike}{0.1840}  % QLIKE ridge sharql all
\newcommand{\rhabRidgeSharqlAllMse}{1.68}  % MSE x 1e10 ridge sharql all
\newcommand{\rhabRidgeSharqlAllPct}{$-$2.8}  % QLIKE change % vs har
\newcommand{\rhabRidgeSharqlAllDm}{$-$6.82}  % DM vs the reference arm
\newcommand{\rhabRidgeSharqlAllP}{\ensuremath{<}0.001}  % p-value
\newcommand{\rhabRidgeSharqlCloseQlike}{0.1216}  % QLIKE ridge sharql c16
\newcommand{\rhabRidgeSharqlCloseMse}{2.65}  % MSE x 1e10 ridge sharql c16
\newcommand{\rhabRidgeSharqlClosePct}{$-$4.2}  % QLIKE change % vs har
\newcommand{\rhabRidgeSharqlCloseDm}{$-$3.49}  % DM vs the reference arm
\newcommand{\rhabRidgeSharqlCloseP}{\ensuremath{<}0.001}  % p-value
\newcommand{\rhabRidgeSharqlCloseLateQlike}{0.1119}  % QLIKE ridge sharql c16_2020
\newcommand{\rhabRidgeSharqlCloseLateMse}{5.92}  % MSE x 1e10 ridge sharql c16_2020
\newcommand{\rhabRidgeSharqlCloseLatePct}{$-$1.6}  % QLIKE change % vs har
\newcommand{\rhabRidgeSharqlCloseLateDm}{$-$0.73}  % DM vs the reference arm
\newcommand{\rhabRidgeSharqlCloseLateP}{0.468}  % p-value
\newcommand{\rhabXgbFullAllQlike}{0.1627}  % QLIKE xgb full all
\newcommand{\rhabXgbFullAllMse}{1.34}  % MSE x 1e10 xgb full all
\newcommand{\rhabXgbFullAllPct}{$-$13.8}  % QLIKE change % vs har
\newcommand{\rhabXgbFullAllDm}{$-$9.94}  % DM vs the reference arm
\newcommand{\rhabXgbFullAllP}{\ensuremath{<}0.001}  % p-value
\newcommand{\rhabXgbFullCloseQlike}{0.1192}  % QLIKE xgb full c16
\newcommand{\rhabXgbFullCloseMse}{2.20}  % MSE x 1e10 xgb full c16
\newcommand{\rhabXgbFullClosePct}{$-$3.4}  % QLIKE change % vs har
\newcommand{\rhabXgbFullCloseDm}{$-$2.18}  % DM vs the reference arm
\newcommand{\rhabXgbFullCloseP}{0.030}  % p-value
\newcommand{\rhabXgbFullCloseLateQlike}{0.1079}  % QLIKE xgb full c16_2020
\newcommand{\rhabXgbFullCloseLateMse}{4.79}  % MSE x 1e10 xgb full c16_2020
\newcommand{\rhabXgbFullCloseLatePct}{$-$2.1}  % QLIKE change % vs har
\newcommand{\rhabXgbFullCloseLateDm}{$-$0.76}  % DM vs the reference arm
\newcommand{\rhabXgbFullCloseLateP}{0.444}  % p-value
\newcommand{\rhabXgbHarAllQlike}{0.1888}  % QLIKE xgb har all
\newcommand{\rhabXgbHarAllMse}{1.16}  % MSE x 1e10 xgb har all
\newcommand{\rhabXgbHarCloseQlike}{0.1234}  % QLIKE xgb har c16
\newcommand{\rhabXgbHarCloseMse}{1.73}  % MSE x 1e10 xgb har c16
\newcommand{\rhabXgbHarCloseLateQlike}{0.1102}  % QLIKE xgb har c16_2020
\newcommand{\rhabXgbHarCloseLateMse}{3.62}  % MSE x 1e10 xgb har c16_2020
\newcommand{\rhabRidgeFullLiveSentQlike}{0.1833}  % QLIKE ridge full live_sent (subsample)
\newcommand{\rhabRidgeFullLiveSentPct}{$-$1.9}  % QLIKE change % vs the reference (subsample)
\newcommand{\rhabRidgeFullLiveSentDm}{$-$0.98}  % DM vs the reference arm (subsample)
\newcommand{\rhabRidgeFullLiveSentP}{0.327}  % p-value
\newcommand{\rhabRidgeLoSentLiveSentQlike}{0.1842}  % QLIKE ridge lo_sent live_sent (subsample)
\newcommand{\rhabRidgeLoSentLiveSentPct}{0.5}  % QLIKE change % vs the reference (subsample)
\newcommand{\rhabRidgeLoSentLiveSentDm}{0.71}  % DM vs the reference arm (subsample)
\newcommand{\rhabRidgeLoSentLiveSentP}{0.477}  % p-value
\newcommand{\rhabRidgeAddSentLiveSentQlike}{0.1846}  % QLIKE ridge add_sent live_sent (subsample)
\newcommand{\rhabRidgeAddSentLiveSentPct}{$-$1.2}  % QLIKE change % vs the reference (subsample)
\newcommand{\rhabRidgeAddSentLiveSentDm}{$-$1.18}  % DM vs the reference arm (subsample)
\newcommand{\rhabRidgeAddSentLiveSentP}{0.238}  % p-value
\newcommand{\rhabRidgeHarLiveSentQlike}{0.1868}  % QLIKE ridge har live_sent (subsample)
\newcommand{\rhabLassoFullLiveSentQlike}{0.1780}  % QLIKE lasso full live_sent (subsample)
\newcommand{\rhabLassoFullLiveSentPct}{$-$4.7}  % QLIKE change % vs the reference (subsample)
\newcommand{\rhabLassoFullLiveSentDm}{$-$1.83}  % DM vs the reference arm (subsample)
\newcommand{\rhabLassoFullLiveSentP}{0.068}  % p-value
\newcommand{\rhabLassoLoSentLiveSentQlike}{0.1786}  % QLIKE lasso lo_sent live_sent (subsample)
\newcommand{\rhabLassoLoSentLiveSentPct}{0.3}  % QLIKE change % vs the reference (subsample)
\newcommand{\rhabLassoLoSentLiveSentDm}{0.41}  % DM vs the reference arm (subsample)
\newcommand{\rhabLassoLoSentLiveSentP}{0.685}  % p-value
\newcommand{\rhabLassoAddSentLiveSentQlike}{0.1850}  % QLIKE lasso add_sent live_sent (subsample)
\newcommand{\rhabLassoAddSentLiveSentPct}{$-$1.0}  % QLIKE change % vs the reference (subsample)
\newcommand{\rhabLassoAddSentLiveSentDm}{$-$0.96}  % DM vs the reference arm (subsample)
\newcommand{\rhabLassoAddSentLiveSentP}{0.339}  % p-value
\newcommand{\rhabLassoHarLiveSentQlike}{0.1868}  % QLIKE lasso har live_sent (subsample)
\newcommand{\rhabEnetFullLiveSentQlike}{0.1818}  % QLIKE enet full live_sent (subsample)
\newcommand{\rhabEnetFullLiveSentPct}{$-$2.7}  % QLIKE change % vs the reference (subsample)
\newcommand{\rhabEnetFullLiveSentDm}{$-$1.55}  % DM vs the reference arm (subsample)
\newcommand{\rhabEnetFullLiveSentP}{0.121}  % p-value
\newcommand{\rhabEnetLoSentLiveSentQlike}{0.1833}  % QLIKE enet lo_sent live_sent (subsample)
\newcommand{\rhabEnetLoSentLiveSentPct}{0.8}  % QLIKE change % vs the reference (subsample)
\newcommand{\rhabEnetLoSentLiveSentDm}{1.43}  % DM vs the reference arm (subsample)
\newcommand{\rhabEnetLoSentLiveSentP}{0.154}  % p-value
\newcommand{\rhabEnetAddSentLiveSentQlike}{0.1845}  % QLIKE enet add_sent live_sent (subsample)
\newcommand{\rhabEnetAddSentLiveSentPct}{$-$1.3}  % QLIKE change % vs the reference (subsample)
\newcommand{\rhabEnetAddSentLiveSentDm}{$-$1.48}  % DM vs the reference arm (subsample)
\newcommand{\rhabEnetAddSentLiveSentP}{0.139}  % p-value
\newcommand{\rhabEnetHarLiveSentQlike}{0.1868}  % QLIKE enet har live_sent (subsample)
\newcommand{\rhabLgbmFullLiveSentQlike}{0.1564}  % QLIKE lgbm full live_sent (subsample)
\newcommand{\rhabLgbmFullLiveSentPct}{$-$13.8}  % QLIKE change % vs the reference (subsample)
\newcommand{\rhabLgbmFullLiveSentDm}{$-$3.76}  % DM vs the reference arm (subsample)
\newcommand{\rhabLgbmFullLiveSentP}{\ensuremath{<}0.001}  % p-value
\newcommand{\rhabLgbmHarLiveSentQlike}{0.1814}  % QLIKE lgbm har live_sent (subsample)
\newcommand{\rhabXgbFullLiveSentQlike}{0.1559}  % QLIKE xgb full live_sent (subsample)
\newcommand{\rhabXgbFullLiveSentPct}{$-$15.6}  % QLIKE change % vs the reference (subsample)
\newcommand{\rhabXgbFullLiveSentDm}{$-$3.69}  % DM vs the reference arm (subsample)
\newcommand{\rhabXgbFullLiveSentP}{\ensuremath{<}0.001}  % p-value
\newcommand{\rhabXgbHarLiveSentQlike}{0.1846}  % QLIKE xgb har live_sent (subsample)
\newcommand{\rhabRfFullLiveSentQlike}{0.1676}  % QLIKE rf full live_sent (subsample)
\newcommand{\rhabRfFullLiveSentPct}{$-$10.8}  % QLIKE change % vs the reference (subsample)
\newcommand{\rhabRfFullLiveSentDm}{$-$4.09}  % DM vs the reference arm (subsample)
\newcommand{\rhabRfFullLiveSentP}{\ensuremath{<}0.001}  % p-value
\newcommand{\rhabRfHarLiveSentQlike}{0.1878}  % QLIKE rf har live_sent (subsample)
\newcommand{\rhabRidgeFullLiveSentCloseQlike}{0.1222}  % QLIKE ridge full live16_sent (subsample)
\newcommand{\rhabRidgeFullLiveSentClosePct}{12.8}  % QLIKE change % vs the reference (subsample)
\newcommand{\rhabRidgeFullLiveSentCloseDm}{2.03}  % DM vs the reference arm (subsample)
\newcommand{\rhabRidgeFullLiveSentCloseP}{0.042}  % p-value
\newcommand{\rhabRidgeLoSentLiveSentCloseQlike}{0.1141}  % QLIKE ridge lo_sent live16_sent (subsample)
\newcommand{\rhabRidgeLoSentLiveSentClosePct}{$-$6.6}  % QLIKE change % vs the reference (subsample)
\newcommand{\rhabRidgeLoSentLiveSentCloseDm}{$-$3.48}  % DM vs the reference arm (subsample)
\newcommand{\rhabRidgeLoSentLiveSentCloseP}{\ensuremath{<}0.001}  % p-value
\newcommand{\rhabRidgeAddSentLiveSentCloseQlike}{0.1188}  % QLIKE ridge add_sent live16_sent (subsample)
\newcommand{\rhabRidgeAddSentLiveSentClosePct}{9.7}  % QLIKE change % vs the reference (subsample)
\newcommand{\rhabRidgeAddSentLiveSentCloseDm}{2.63}  % DM vs the reference arm (subsample)
\newcommand{\rhabRidgeAddSentLiveSentCloseP}{0.009}  % p-value
\newcommand{\rhabRidgeHarLiveSentCloseQlike}{0.1083}  % QLIKE ridge har live16_sent (subsample)
\newcommand{\rhabLassoFullLiveSentCloseQlike}{0.1234}  % QLIKE lasso full live16_sent (subsample)
\newcommand{\rhabLassoFullLiveSentClosePct}{13.9}  % QLIKE change % vs the reference (subsample)
\newcommand{\rhabLassoFullLiveSentCloseDm}{2.17}  % DM vs the reference arm (subsample)
\newcommand{\rhabLassoFullLiveSentCloseP}{0.030}  % p-value
\newcommand{\rhabLassoLoSentLiveSentCloseQlike}{0.1155}  % QLIKE lasso lo_sent live16_sent (subsample)
\newcommand{\rhabLassoLoSentLiveSentClosePct}{$-$6.3}  % QLIKE change % vs the reference (subsample)
\newcommand{\rhabLassoLoSentLiveSentCloseDm}{$-$2.92}  % DM vs the reference arm (subsample)
\newcommand{\rhabLassoLoSentLiveSentCloseP}{0.003}  % p-value
\newcommand{\rhabLassoAddSentLiveSentCloseQlike}{0.1195}  % QLIKE lasso add_sent live16_sent (subsample)
\newcommand{\rhabLassoAddSentLiveSentClosePct}{10.3}  % QLIKE change % vs the reference (subsample)
\newcommand{\rhabLassoAddSentLiveSentCloseDm}{2.76}  % DM vs the reference arm (subsample)
\newcommand{\rhabLassoAddSentLiveSentCloseP}{0.006}  % p-value
\newcommand{\rhabLassoHarLiveSentCloseQlike}{0.1083}  % QLIKE lasso har live16_sent (subsample)
\newcommand{\rhabEnetFullLiveSentCloseQlike}{0.1171}  % QLIKE enet full live16_sent (subsample)
\newcommand{\rhabEnetFullLiveSentClosePct}{8.1}  % QLIKE change % vs the reference (subsample)
\newcommand{\rhabEnetFullLiveSentCloseDm}{1.35}  % DM vs the reference arm (subsample)
\newcommand{\rhabEnetFullLiveSentCloseP}{0.176}  % p-value
\newcommand{\rhabEnetLoSentLiveSentCloseQlike}{0.1134}  % QLIKE enet lo_sent live16_sent (subsample)
\newcommand{\rhabEnetLoSentLiveSentClosePct}{$-$3.2}  % QLIKE change % vs the reference (subsample)
\newcommand{\rhabEnetLoSentLiveSentCloseDm}{$-$2.23}  % DM vs the reference arm (subsample)
\newcommand{\rhabEnetLoSentLiveSentCloseP}{0.026}  % p-value
\newcommand{\rhabEnetAddSentLiveSentCloseQlike}{0.1156}  % QLIKE enet add_sent live16_sent (subsample)
\newcommand{\rhabEnetAddSentLiveSentClosePct}{6.7}  % QLIKE change % vs the reference (subsample)
\newcommand{\rhabEnetAddSentLiveSentCloseDm}{1.91}  % DM vs the reference arm (subsample)
\newcommand{\rhabEnetAddSentLiveSentCloseP}{0.056}  % p-value
\newcommand{\rhabEnetHarLiveSentCloseQlike}{0.1083}  % QLIKE enet har live16_sent (subsample)
\newcommand{\rhabLgbmFullLiveSentCloseQlike}{0.0995}  % QLIKE lgbm full live16_sent (subsample)
\newcommand{\rhabLgbmFullLiveSentClosePct}{$-$2.5}  % QLIKE change % vs the reference (subsample)
\newcommand{\rhabLgbmFullLiveSentCloseDm}{$-$0.57}  % DM vs the reference arm (subsample)
\newcommand{\rhabLgbmFullLiveSentCloseP}{0.565}  % p-value
\newcommand{\rhabLgbmHarLiveSentCloseQlike}{0.1021}  % QLIKE lgbm har live16_sent (subsample)
\newcommand{\rhabXgbFullLiveSentCloseQlike}{0.0997}  % QLIKE xgb full live16_sent (subsample)
\newcommand{\rhabXgbFullLiveSentClosePct}{$-$2.8}  % QLIKE change % vs the reference (subsample)
\newcommand{\rhabXgbFullLiveSentCloseDm}{$-$0.61}  % DM vs the reference arm (subsample)
\newcommand{\rhabXgbFullLiveSentCloseP}{0.541}  % p-value
\newcommand{\rhabXgbHarLiveSentCloseQlike}{0.1025}  % QLIKE xgb har live16_sent (subsample)
\newcommand{\rhabRfFullLiveSentCloseQlike}{0.1026}  % QLIKE rf full live16_sent (subsample)
\newcommand{\rhabRfFullLiveSentClosePct}{$-$2.6}  % QLIKE change % vs the reference (subsample)
\newcommand{\rhabRfFullLiveSentCloseDm}{$-$0.65}  % DM vs the reference arm (subsample)
\newcommand{\rhabRfFullLiveSentCloseP}{0.513}  % p-value
\newcommand{\rhabRfHarLiveSentCloseQlike}{0.1053}  % QLIKE rf har live16_sent (subsample)
\newcommand{\rhabRidgeFullLiveIvQlike}{0.1940}  % QLIKE ridge full live_iv (subsample)
\newcommand{\rhabRidgeFullLiveIvPct}{$-$4.7}  % QLIKE change % vs the reference (subsample)
\newcommand{\rhabRidgeFullLiveIvDm}{$-$3.62}  % DM vs the reference arm (subsample)
\newcommand{\rhabRidgeFullLiveIvP}{\ensuremath{<}0.001}  % p-value
\newcommand{\rhabRidgeLoIvLiveIvQlike}{0.1942}  % QLIKE ridge lo_iv live_iv (subsample)
\newcommand{\rhabRidgeLoIvLiveIvPct}{0.1}  % QLIKE change % vs the reference (subsample)
\newcommand{\rhabRidgeLoIvLiveIvDm}{0.25}  % DM vs the reference arm (subsample)
\newcommand{\rhabRidgeLoIvLiveIvP}{0.806}  % p-value
\newcommand{\rhabRidgeAddIvLiveIvQlike}{0.2037}  % QLIKE ridge add_iv live_iv (subsample)
\newcommand{\rhabRidgeAddIvLiveIvPct}{0.0}  % QLIKE change % vs the reference (subsample)
\newcommand{\rhabRidgeAddIvLiveIvDm}{0.02}  % DM vs the reference arm (subsample)
\newcommand{\rhabRidgeAddIvLiveIvP}{0.986}  % p-value
\newcommand{\rhabRidgeHarLiveIvQlike}{0.2037}  % QLIKE ridge har live_iv (subsample)
\newcommand{\rhabLassoFullLiveIvQlike}{0.1924}  % QLIKE lasso full live_iv (subsample)
\newcommand{\rhabLassoFullLiveIvPct}{$-$5.5}  % QLIKE change % vs the reference (subsample)
\newcommand{\rhabLassoFullLiveIvDm}{$-$2.98}  % DM vs the reference arm (subsample)
\newcommand{\rhabLassoFullLiveIvP}{0.003}  % p-value
\newcommand{\rhabLassoLoIvLiveIvQlike}{0.1931}  % QLIKE lasso lo_iv live_iv (subsample)
\newcommand{\rhabLassoLoIvLiveIvPct}{0.4}  % QLIKE change % vs the reference (subsample)
\newcommand{\rhabLassoLoIvLiveIvDm}{0.47}  % DM vs the reference arm (subsample)
\newcommand{\rhabLassoLoIvLiveIvP}{0.641}  % p-value
\newcommand{\rhabLassoAddIvLiveIvQlike}{0.2019}  % QLIKE lasso add_iv live_iv (subsample)
\newcommand{\rhabLassoAddIvLiveIvPct}{$-$0.9}  % QLIKE change % vs the reference (subsample)
\newcommand{\rhabLassoAddIvLiveIvDm}{$-$1.39}  % DM vs the reference arm (subsample)
\newcommand{\rhabLassoAddIvLiveIvP}{0.165}  % p-value
\newcommand{\rhabLassoHarLiveIvQlike}{0.2037}  % QLIKE lasso har live_iv (subsample)
\newcommand{\rhabEnetFullLiveIvQlike}{0.1921}  % QLIKE enet full live_iv (subsample)
\newcommand{\rhabEnetFullLiveIvPct}{$-$5.7}  % QLIKE change % vs the reference (subsample)
\newcommand{\rhabEnetFullLiveIvDm}{$-$4.87}  % DM vs the reference arm (subsample)
\newcommand{\rhabEnetFullLiveIvP}{\ensuremath{<}0.001}  % p-value
\newcommand{\rhabEnetLoIvLiveIvQlike}{0.1924}  % QLIKE enet lo_iv live_iv (subsample)
\newcommand{\rhabEnetLoIvLiveIvPct}{0.1}  % QLIKE change % vs the reference (subsample)
\newcommand{\rhabEnetLoIvLiveIvDm}{0.23}  % DM vs the reference arm (subsample)
\newcommand{\rhabEnetLoIvLiveIvP}{0.816}  % p-value
\newcommand{\rhabEnetAddIvLiveIvQlike}{0.2033}  % QLIKE enet add_iv live_iv (subsample)
\newcommand{\rhabEnetAddIvLiveIvPct}{$-$0.2}  % QLIKE change % vs the reference (subsample)
\newcommand{\rhabEnetAddIvLiveIvDm}{$-$0.36}  % DM vs the reference arm (subsample)
\newcommand{\rhabEnetAddIvLiveIvP}{0.722}  % p-value
\newcommand{\rhabEnetHarLiveIvQlike}{0.2037}  % QLIKE enet har live_iv (subsample)
\newcommand{\rhabLgbmFullLiveIvQlike}{0.1739}  % QLIKE lgbm full live_iv (subsample)
\newcommand{\rhabLgbmFullLiveIvPct}{$-$13.3}  % QLIKE change % vs the reference (subsample)
\newcommand{\rhabLgbmFullLiveIvDm}{$-$7.31}  % DM vs the reference arm (subsample)
\newcommand{\rhabLgbmFullLiveIvP}{\ensuremath{<}0.001}  % p-value
\newcommand{\rhabLgbmHarLiveIvQlike}{0.2006}  % QLIKE lgbm har live_iv (subsample)
\newcommand{\rhabXgbFullLiveIvQlike}{0.1739}  % QLIKE xgb full live_iv (subsample)
\newcommand{\rhabXgbFullLiveIvPct}{$-$14.1}  % QLIKE change % vs the reference (subsample)
\newcommand{\rhabXgbFullLiveIvDm}{$-$7.16}  % DM vs the reference arm (subsample)
\newcommand{\rhabXgbFullLiveIvP}{\ensuremath{<}0.001}  % p-value
\newcommand{\rhabXgbHarLiveIvQlike}{0.2024}  % QLIKE xgb har live_iv (subsample)
\newcommand{\rhabRfFullLiveIvQlike}{0.1818}  % QLIKE rf full live_iv (subsample)
\newcommand{\rhabRfFullLiveIvPct}{$-$10.4}  % QLIKE change % vs the reference (subsample)
\newcommand{\rhabRfFullLiveIvDm}{$-$7.47}  % DM vs the reference arm (subsample)
\newcommand{\rhabRfFullLiveIvP}{\ensuremath{<}0.001}  % p-value
\newcommand{\rhabRfHarLiveIvQlike}{0.2029}  % QLIKE rf har live_iv (subsample)
\newcommand{\rhabRidgeFullLiveIvCloseQlike}{0.1380}  % QLIKE ridge full live16_iv (subsample)
\newcommand{\rhabRidgeFullLiveIvClosePct}{6.0}  % QLIKE change % vs the reference (subsample)
\newcommand{\rhabRidgeFullLiveIvCloseDm}{1.50}  % DM vs the reference arm (subsample)
\newcommand{\rhabRidgeFullLiveIvCloseP}{0.134}  % p-value
\newcommand{\rhabRidgeLoIvLiveIvCloseQlike}{0.1355}  % QLIKE ridge lo_iv live16_iv (subsample)
\newcommand{\rhabRidgeLoIvLiveIvClosePct}{$-$1.8}  % QLIKE change % vs the reference (subsample)
\newcommand{\rhabRidgeLoIvLiveIvCloseDm}{$-$1.22}  % DM vs the reference arm (subsample)
\newcommand{\rhabRidgeLoIvLiveIvCloseP}{0.222}  % p-value
\newcommand{\rhabRidgeAddIvLiveIvCloseQlike}{0.1306}  % QLIKE ridge add_iv live16_iv (subsample)
\newcommand{\rhabRidgeAddIvLiveIvClosePct}{0.3}  % QLIKE change % vs the reference (subsample)
\newcommand{\rhabRidgeAddIvLiveIvCloseDm}{0.19}  % DM vs the reference arm (subsample)
\newcommand{\rhabRidgeAddIvLiveIvCloseP}{0.853}  % p-value
\newcommand{\rhabRidgeHarLiveIvCloseQlike}{0.1302}  % QLIKE ridge har live16_iv (subsample)
\newcommand{\rhabLassoFullLiveIvCloseQlike}{0.1425}  % QLIKE lasso full live16_iv (subsample)
\newcommand{\rhabLassoFullLiveIvClosePct}{9.4}  % QLIKE change % vs the reference (subsample)
\newcommand{\rhabLassoFullLiveIvCloseDm}{1.93}  % DM vs the reference arm (subsample)
\newcommand{\rhabLassoFullLiveIvCloseP}{0.054}  % p-value
\newcommand{\rhabLassoLoIvLiveIvCloseQlike}{0.1378}  % QLIKE lasso lo_iv live16_iv (subsample)
\newcommand{\rhabLassoLoIvLiveIvClosePct}{$-$3.3}  % QLIKE change % vs the reference (subsample)
\newcommand{\rhabLassoLoIvLiveIvCloseDm}{$-$1.78}  % DM vs the reference arm (subsample)
\newcommand{\rhabLassoLoIvLiveIvCloseP}{0.075}  % p-value
\newcommand{\rhabLassoAddIvLiveIvCloseQlike}{0.1300}  % QLIKE lasso add_iv live16_iv (subsample)
\newcommand{\rhabLassoAddIvLiveIvClosePct}{$-$0.1}  % QLIKE change % vs the reference (subsample)
\newcommand{\rhabLassoAddIvLiveIvCloseDm}{$-$0.08}  % DM vs the reference arm (subsample)
\newcommand{\rhabLassoAddIvLiveIvCloseP}{0.940}  % p-value
\newcommand{\rhabLassoHarLiveIvCloseQlike}{0.1302}  % QLIKE lasso har live16_iv (subsample)
\newcommand{\rhabEnetFullLiveIvCloseQlike}{0.1355}  % QLIKE enet full live16_iv (subsample)
\newcommand{\rhabEnetFullLiveIvClosePct}{4.1}  % QLIKE change % vs the reference (subsample)
\newcommand{\rhabEnetFullLiveIvCloseDm}{1.12}  % DM vs the reference arm (subsample)
\newcommand{\rhabEnetFullLiveIvCloseP}{0.262}  % p-value
\newcommand{\rhabEnetLoIvLiveIvCloseQlike}{0.1321}  % QLIKE enet lo_iv live16_iv (subsample)
\newcommand{\rhabEnetLoIvLiveIvClosePct}{$-$2.5}  % QLIKE change % vs the reference (subsample)
\newcommand{\rhabEnetLoIvLiveIvCloseDm}{$-$1.58}  % DM vs the reference arm (subsample)
\newcommand{\rhabEnetLoIvLiveIvCloseP}{0.115}  % p-value
\newcommand{\rhabEnetAddIvLiveIvCloseQlike}{0.1308}  % QLIKE enet add_iv live16_iv (subsample)
\newcommand{\rhabEnetAddIvLiveIvClosePct}{0.5}  % QLIKE change % vs the reference (subsample)
\newcommand{\rhabEnetAddIvLiveIvCloseDm}{0.31}  % DM vs the reference arm (subsample)
\newcommand{\rhabEnetAddIvLiveIvCloseP}{0.756}  % p-value
\newcommand{\rhabEnetHarLiveIvCloseQlike}{0.1302}  % QLIKE enet har live16_iv (subsample)
\newcommand{\rhabLgbmFullLiveIvCloseQlike}{0.1242}  % QLIKE lgbm full live16_iv (subsample)
\newcommand{\rhabLgbmFullLiveIvClosePct}{$-$1.6}  % QLIKE change % vs the reference (subsample)
\newcommand{\rhabLgbmFullLiveIvCloseDm}{$-$0.74}  % DM vs the reference arm (subsample)
\newcommand{\rhabLgbmFullLiveIvCloseP}{0.461}  % p-value
\newcommand{\rhabLgbmHarLiveIvCloseQlike}{0.1263}  % QLIKE lgbm har live16_iv (subsample)
\newcommand{\rhabXgbFullLiveIvCloseQlike}{0.1234}  % QLIKE xgb full live16_iv (subsample)
\newcommand{\rhabXgbFullLiveIvClosePct}{$-$2.3}  % QLIKE change % vs the reference (subsample)
\newcommand{\rhabXgbFullLiveIvCloseDm}{$-$1.01}  % DM vs the reference arm (subsample)
\newcommand{\rhabXgbFullLiveIvCloseP}{0.312}  % p-value
\newcommand{\rhabXgbHarLiveIvCloseQlike}{0.1263}  % QLIKE xgb har live16_iv (subsample)
\newcommand{\rhabRfFullLiveIvCloseQlike}{0.1304}  % QLIKE rf full live16_iv (subsample)
\newcommand{\rhabRfFullLiveIvClosePct}{$-$0.4}  % QLIKE change % vs the reference (subsample)
\newcommand{\rhabRfFullLiveIvCloseDm}{$-$0.16}  % DM vs the reference arm (subsample)
\newcommand{\rhabRfFullLiveIvCloseP}{0.870}  % p-value
\newcommand{\rhabRfHarLiveIvCloseQlike}{0.1309}  % QLIKE rf har live16_iv (subsample)
\newcommand{\rhabRidgeFullLiveIvbQlike}{0.1757}  % QLIKE ridge full live_ivb (subsample)
\newcommand{\rhabRidgeFullLiveIvbPct}{$-$3.6}  % QLIKE change % vs the reference (subsample)
\newcommand{\rhabRidgeFullLiveIvbDm}{$-$2.12}  % DM vs the reference arm (subsample)
\newcommand{\rhabRidgeFullLiveIvbP}{0.034}  % p-value
\newcommand{\rhabRidgeLoIvLiveIvbQlike}{0.1746}  % QLIKE ridge lo_iv live_ivb (subsample)
\newcommand{\rhabRidgeLoIvLiveIvbPct}{$-$0.6}  % QLIKE change % vs the reference (subsample)
\newcommand{\rhabRidgeLoIvLiveIvbDm}{$-$0.79}  % DM vs the reference arm (subsample)
\newcommand{\rhabRidgeLoIvLiveIvbP}{0.431}  % p-value
\newcommand{\rhabRidgeAddIvLiveIvbQlike}{0.1850}  % QLIKE ridge add_iv live_ivb (subsample)
\newcommand{\rhabRidgeAddIvLiveIvbPct}{1.4}  % QLIKE change % vs the reference (subsample)
\newcommand{\rhabRidgeAddIvLiveIvbDm}{1.50}  % DM vs the reference arm (subsample)
\newcommand{\rhabRidgeAddIvLiveIvbP}{0.134}  % p-value
\newcommand{\rhabRidgeHarLiveIvbQlike}{0.1824}  % QLIKE ridge har live_ivb (subsample)
\newcommand{\rhabLassoFullLiveIvbQlike}{0.1718}  % QLIKE lasso full live_ivb (subsample)
\newcommand{\rhabLassoFullLiveIvbPct}{$-$5.8}  % QLIKE change % vs the reference (subsample)
\newcommand{\rhabLassoFullLiveIvbDm}{$-$2.58}  % DM vs the reference arm (subsample)
\newcommand{\rhabLassoFullLiveIvbP}{0.010}  % p-value
\newcommand{\rhabLassoLoIvLiveIvbQlike}{0.1700}  % QLIKE lasso lo_iv live_ivb (subsample)
\newcommand{\rhabLassoLoIvLiveIvbPct}{$-$1.0}  % QLIKE change % vs the reference (subsample)
\newcommand{\rhabLassoLoIvLiveIvbDm}{$-$0.96}  % DM vs the reference arm (subsample)
\newcommand{\rhabLassoLoIvLiveIvbP}{0.339}  % p-value
\newcommand{\rhabLassoAddIvLiveIvbQlike}{0.1844}  % QLIKE lasso add_iv live_ivb (subsample)
\newcommand{\rhabLassoAddIvLiveIvbPct}{1.1}  % QLIKE change % vs the reference (subsample)
\newcommand{\rhabLassoAddIvLiveIvbDm}{1.04}  % DM vs the reference arm (subsample)
\newcommand{\rhabLassoAddIvLiveIvbP}{0.300}  % p-value
\newcommand{\rhabLassoHarLiveIvbQlike}{0.1824}  % QLIKE lasso har live_ivb (subsample)
\newcommand{\rhabEnetFullLiveIvbQlike}{0.1745}  % QLIKE enet full live_ivb (subsample)
\newcommand{\rhabEnetFullLiveIvbPct}{$-$4.3}  % QLIKE change % vs the reference (subsample)
\newcommand{\rhabEnetFullLiveIvbDm}{$-$2.67}  % DM vs the reference arm (subsample)
\newcommand{\rhabEnetFullLiveIvbP}{0.008}  % p-value
\newcommand{\rhabEnetLoIvLiveIvbQlike}{0.1732}  % QLIKE enet lo_iv live_ivb (subsample)
\newcommand{\rhabEnetLoIvLiveIvbPct}{$-$0.7}  % QLIKE change % vs the reference (subsample)
\newcommand{\rhabEnetLoIvLiveIvbDm}{$-$0.88}  % DM vs the reference arm (subsample)
\newcommand{\rhabEnetLoIvLiveIvbP}{0.377}  % p-value
\newcommand{\rhabEnetAddIvLiveIvbQlike}{0.1844}  % QLIKE enet add_iv live_ivb (subsample)
\newcommand{\rhabEnetAddIvLiveIvbPct}{1.1}  % QLIKE change % vs the reference (subsample)
\newcommand{\rhabEnetAddIvLiveIvbDm}{1.30}  % DM vs the reference arm (subsample)
\newcommand{\rhabEnetAddIvLiveIvbP}{0.195}  % p-value
\newcommand{\rhabEnetHarLiveIvbQlike}{0.1824}  % QLIKE enet har live_ivb (subsample)
\newcommand{\rhabLgbmFullLiveIvbQlike}{0.1518}  % QLIKE lgbm full live_ivb (subsample)
\newcommand{\rhabLgbmFullLiveIvbPct}{$-$14.6}  % QLIKE change % vs the reference (subsample)
\newcommand{\rhabLgbmFullLiveIvbDm}{$-$5.14}  % DM vs the reference arm (subsample)
\newcommand{\rhabLgbmFullLiveIvbP}{\ensuremath{<}0.001}  % p-value
\newcommand{\rhabLgbmHarLiveIvbQlike}{0.1778}  % QLIKE lgbm har live_ivb (subsample)
\newcommand{\rhabXgbFullLiveIvbQlike}{0.1511}  % QLIKE xgb full live_ivb (subsample)
\newcommand{\rhabXgbFullLiveIvbPct}{$-$16.3}  % QLIKE change % vs the reference (subsample)
\newcommand{\rhabXgbFullLiveIvbDm}{$-$4.98}  % DM vs the reference arm (subsample)
\newcommand{\rhabXgbFullLiveIvbP}{\ensuremath{<}0.001}  % p-value
\newcommand{\rhabXgbHarLiveIvbQlike}{0.1805}  % QLIKE xgb har live_ivb (subsample)
\newcommand{\rhabRfFullLiveIvbQlike}{0.1621}  % QLIKE rf full live_ivb (subsample)
\newcommand{\rhabRfFullLiveIvbPct}{$-$11.3}  % QLIKE change % vs the reference (subsample)
\newcommand{\rhabRfFullLiveIvbDm}{$-$5.43}  % DM vs the reference arm (subsample)
\newcommand{\rhabRfFullLiveIvbP}{\ensuremath{<}0.001}  % p-value
\newcommand{\rhabRfHarLiveIvbQlike}{0.1829}  % QLIKE rf har live_ivb (subsample)
\newcommand{\rhabRidgeFullLiveIvbCloseQlike}{0.1180}  % QLIKE ridge full live16_ivb (subsample)
\newcommand{\rhabRidgeFullLiveIvbClosePct}{15.1}  % QLIKE change % vs the reference (subsample)
\newcommand{\rhabRidgeFullLiveIvbCloseDm}{2.54}  % DM vs the reference arm (subsample)
\newcommand{\rhabRidgeFullLiveIvbCloseP}{0.011}  % p-value
\newcommand{\rhabRidgeLoIvLiveIvbCloseQlike}{0.1123}  % QLIKE ridge lo_iv live16_ivb (subsample)
\newcommand{\rhabRidgeLoIvLiveIvbClosePct}{$-$4.9}  % QLIKE change % vs the reference (subsample)
\newcommand{\rhabRidgeLoIvLiveIvbCloseDm}{$-$1.73}  % DM vs the reference arm (subsample)
\newcommand{\rhabRidgeLoIvLiveIvbCloseP}{0.083}  % p-value
\newcommand{\rhabRidgeAddIvLiveIvbCloseQlike}{0.1073}  % QLIKE ridge add_iv live16_ivb (subsample)
\newcommand{\rhabRidgeAddIvLiveIvbClosePct}{4.6}  % QLIKE change % vs the reference (subsample)
\newcommand{\rhabRidgeAddIvLiveIvbCloseDm}{1.47}  % DM vs the reference arm (subsample)
\newcommand{\rhabRidgeAddIvLiveIvbCloseP}{0.142}  % p-value
\newcommand{\rhabRidgeHarLiveIvbCloseQlike}{0.1026}  % QLIKE ridge har live16_ivb (subsample)
\newcommand{\rhabLassoFullLiveIvbCloseQlike}{0.1202}  % QLIKE lasso full live16_ivb (subsample)
\newcommand{\rhabLassoFullLiveIvbClosePct}{17.2}  % QLIKE change % vs the reference (subsample)
\newcommand{\rhabLassoFullLiveIvbCloseDm}{2.74}  % DM vs the reference arm (subsample)
\newcommand{\rhabLassoFullLiveIvbCloseP}{0.006}  % p-value
\newcommand{\rhabLassoLoIvLiveIvbCloseQlike}{0.1120}  % QLIKE lasso lo_iv live16_ivb (subsample)
\newcommand{\rhabLassoLoIvLiveIvbClosePct}{$-$6.8}  % QLIKE change % vs the reference (subsample)
\newcommand{\rhabLassoLoIvLiveIvbCloseDm}{$-$2.01}  % DM vs the reference arm (subsample)
\newcommand{\rhabLassoLoIvLiveIvbCloseP}{0.044}  % p-value
\newcommand{\rhabLassoAddIvLiveIvbCloseQlike}{0.1066}  % QLIKE lasso add_iv live16_ivb (subsample)
\newcommand{\rhabLassoAddIvLiveIvbClosePct}{3.9}  % QLIKE change % vs the reference (subsample)
\newcommand{\rhabLassoAddIvLiveIvbCloseDm}{1.22}  % DM vs the reference arm (subsample)
\newcommand{\rhabLassoAddIvLiveIvbCloseP}{0.221}  % p-value
\newcommand{\rhabLassoHarLiveIvbCloseQlike}{0.1026}  % QLIKE lasso har live16_ivb (subsample)
\newcommand{\rhabEnetFullLiveIvbCloseQlike}{0.1147}  % QLIKE enet full live16_ivb (subsample)
\newcommand{\rhabEnetFullLiveIvbClosePct}{11.8}  % QLIKE change % vs the reference (subsample)
\newcommand{\rhabEnetFullLiveIvbCloseDm}{2.03}  % DM vs the reference arm (subsample)
\newcommand{\rhabEnetFullLiveIvbCloseP}{0.043}  % p-value
\newcommand{\rhabEnetLoIvLiveIvbCloseQlike}{0.1073}  % QLIKE enet lo_iv live16_ivb (subsample)
\newcommand{\rhabEnetLoIvLiveIvbClosePct}{$-$6.5}  % QLIKE change % vs the reference (subsample)
\newcommand{\rhabEnetLoIvLiveIvbCloseDm}{$-$2.14}  % DM vs the reference arm (subsample)
\newcommand{\rhabEnetLoIvLiveIvbCloseP}{0.033}  % p-value
\newcommand{\rhabEnetAddIvLiveIvbCloseQlike}{0.1076}  % QLIKE enet add_iv live16_ivb (subsample)
\newcommand{\rhabEnetAddIvLiveIvbClosePct}{4.9}  % QLIKE change % vs the reference (subsample)
\newcommand{\rhabEnetAddIvLiveIvbCloseDm}{1.60}  % DM vs the reference arm (subsample)
\newcommand{\rhabEnetAddIvLiveIvbCloseP}{0.110}  % p-value
\newcommand{\rhabEnetHarLiveIvbCloseQlike}{0.1026}  % QLIKE enet har live16_ivb (subsample)
\newcommand{\rhabLgbmFullLiveIvbCloseQlike}{0.0964}  % QLIKE lgbm full live16_ivb (subsample)
\newcommand{\rhabLgbmFullLiveIvbClosePct}{$-$1.0}  % QLIKE change % vs the reference (subsample)
\newcommand{\rhabLgbmFullLiveIvbCloseDm}{$-$0.29}  % DM vs the reference arm (subsample)
\newcommand{\rhabLgbmFullLiveIvbCloseP}{0.775}  % p-value
\newcommand{\rhabLgbmHarLiveIvbCloseQlike}{0.0974}  % QLIKE lgbm har live16_ivb (subsample)
\newcommand{\rhabXgbFullLiveIvbCloseQlike}{0.0962}  % QLIKE xgb full live16_ivb (subsample)
\newcommand{\rhabXgbFullLiveIvbClosePct}{$-$1.5}  % QLIKE change % vs the reference (subsample)
\newcommand{\rhabXgbFullLiveIvbCloseDm}{$-$0.40}  % DM vs the reference arm (subsample)
\newcommand{\rhabXgbFullLiveIvbCloseP}{0.689}  % p-value
\newcommand{\rhabXgbHarLiveIvbCloseQlike}{0.0976}  % QLIKE xgb har live16_ivb (subsample)
\newcommand{\rhabRfFullLiveIvbCloseQlike}{0.1005}  % QLIKE rf full live16_ivb (subsample)
\newcommand{\rhabRfFullLiveIvbClosePct}{0.2}  % QLIKE change % vs the reference (subsample)
\newcommand{\rhabRfFullLiveIvbCloseDm}{0.05}  % DM vs the reference arm (subsample)
\newcommand{\rhabRfFullLiveIvbCloseP}{0.958}  % p-value
\newcommand{\rhabRfHarLiveIvbCloseQlike}{0.1003}  % QLIKE rf har live16_ivb (subsample)
\newcommand{\rhabRidgeFullLiveVdQlike}{0.1942}  % QLIKE ridge full live_vd (subsample)
\newcommand{\rhabRidgeFullLiveVdPct}{4.2}  % QLIKE change % vs the reference (subsample)
\newcommand{\rhabRidgeFullLiveVdDm}{1.33}  % DM vs the reference arm (subsample)
\newcommand{\rhabRidgeFullLiveVdP}{0.183}  % p-value
\newcommand{\rhabRidgeLoVdLiveVdQlike}{0.1887}  % QLIKE ridge lo_vd live_vd (subsample)
\newcommand{\rhabRidgeLoVdLiveVdPct}{$-$2.8}  % QLIKE change % vs the reference (subsample)
\newcommand{\rhabRidgeLoVdLiveVdDm}{$-$3.17}  % DM vs the reference arm (subsample)
\newcommand{\rhabRidgeLoVdLiveVdP}{0.002}  % p-value
\newcommand{\rhabRidgeAddVdLiveVdQlike}{0.1944}  % QLIKE ridge add_vd live_vd (subsample)
\newcommand{\rhabRidgeAddVdLiveVdPct}{4.3}  % QLIKE change % vs the reference (subsample)
\newcommand{\rhabRidgeAddVdLiveVdDm}{2.52}  % DM vs the reference arm (subsample)
\newcommand{\rhabRidgeAddVdLiveVdP}{0.012}  % p-value
\newcommand{\rhabRidgeHarLiveVdQlike}{0.1864}  % QLIKE ridge har live_vd (subsample)
\newcommand{\rhabLassoFullLiveVdQlike}{0.1898}  % QLIKE lasso full live_vd (subsample)
\newcommand{\rhabLassoFullLiveVdPct}{1.9}  % QLIKE change % vs the reference (subsample)
\newcommand{\rhabLassoFullLiveVdDm}{0.45}  % DM vs the reference arm (subsample)
\newcommand{\rhabLassoFullLiveVdP}{0.651}  % p-value
\newcommand{\rhabLassoLoVdLiveVdQlike}{0.1846}  % QLIKE lasso lo_vd live_vd (subsample)
\newcommand{\rhabLassoLoVdLiveVdPct}{$-$2.7}  % QLIKE change % vs the reference (subsample)
\newcommand{\rhabLassoLoVdLiveVdDm}{$-$2.51}  % DM vs the reference arm (subsample)
\newcommand{\rhabLassoLoVdLiveVdP}{0.012}  % p-value
\newcommand{\rhabLassoAddVdLiveVdQlike}{0.1941}  % QLIKE lasso add_vd live_vd (subsample)
\newcommand{\rhabLassoAddVdLiveVdPct}{4.2}  % QLIKE change % vs the reference (subsample)
\newcommand{\rhabLassoAddVdLiveVdDm}{2.49}  % DM vs the reference arm (subsample)
\newcommand{\rhabLassoAddVdLiveVdP}{0.013}  % p-value
\newcommand{\rhabLassoHarLiveVdQlike}{0.1864}  % QLIKE lasso har live_vd (subsample)
\newcommand{\rhabEnetFullLiveVdQlike}{0.1908}  % QLIKE enet full live_vd (subsample)
\newcommand{\rhabEnetFullLiveVdPct}{2.4}  % QLIKE change % vs the reference (subsample)
\newcommand{\rhabEnetFullLiveVdDm}{0.85}  % DM vs the reference arm (subsample)
\newcommand{\rhabEnetFullLiveVdP}{0.397}  % p-value
\newcommand{\rhabEnetLoVdLiveVdQlike}{0.1869}  % QLIKE enet lo_vd live_vd (subsample)
\newcommand{\rhabEnetLoVdLiveVdPct}{$-$2.1}  % QLIKE change % vs the reference (subsample)
\newcommand{\rhabEnetLoVdLiveVdDm}{$-$3.47}  % DM vs the reference arm (subsample)
\newcommand{\rhabEnetLoVdLiveVdP}{\ensuremath{<}0.001}  % p-value
\newcommand{\rhabEnetAddVdLiveVdQlike}{0.1916}  % QLIKE enet add_vd live_vd (subsample)
\newcommand{\rhabEnetAddVdLiveVdPct}{2.8}  % QLIKE change % vs the reference (subsample)
\newcommand{\rhabEnetAddVdLiveVdDm}{2.20}  % DM vs the reference arm (subsample)
\newcommand{\rhabEnetAddVdLiveVdP}{0.028}  % p-value
\newcommand{\rhabEnetHarLiveVdQlike}{0.1864}  % QLIKE enet har live_vd (subsample)
\newcommand{\rhabLgbmFullLiveVdQlike}{0.1612}  % QLIKE lgbm full live_vd (subsample)
\newcommand{\rhabLgbmFullLiveVdPct}{$-$10.9}  % QLIKE change % vs the reference (subsample)
\newcommand{\rhabLgbmFullLiveVdDm}{$-$2.01}  % DM vs the reference arm (subsample)
\newcommand{\rhabLgbmFullLiveVdP}{0.044}  % p-value
\newcommand{\rhabLgbmHarLiveVdQlike}{0.1809}  % QLIKE lgbm har live_vd (subsample)
\newcommand{\rhabXgbFullLiveVdQlike}{0.1616}  % QLIKE xgb full live_vd (subsample)
\newcommand{\rhabXgbFullLiveVdPct}{$-$12.1}  % QLIKE change % vs the reference (subsample)
\newcommand{\rhabXgbFullLiveVdDm}{$-$2.10}  % DM vs the reference arm (subsample)
\newcommand{\rhabXgbFullLiveVdP}{0.036}  % p-value
\newcommand{\rhabXgbHarLiveVdQlike}{0.1839}  % QLIKE xgb har live_vd (subsample)
\newcommand{\rhabRfFullLiveVdQlike}{0.1704}  % QLIKE rf full live_vd (subsample)
\newcommand{\rhabRfFullLiveVdPct}{$-$9.1}  % QLIKE change % vs the reference (subsample)
\newcommand{\rhabRfFullLiveVdDm}{$-$2.40}  % DM vs the reference arm (subsample)
\newcommand{\rhabRfFullLiveVdP}{0.016}  % p-value
\newcommand{\rhabRfHarLiveVdQlike}{0.1874}  % QLIKE rf har live_vd (subsample)
\newcommand{\rhabRidgeFullLiveVdCloseQlike}{0.1269}  % QLIKE ridge full live16_vd (subsample)
\newcommand{\rhabRidgeFullLiveVdClosePct}{45.8}  % QLIKE change % vs the reference (subsample)
\newcommand{\rhabRidgeFullLiveVdCloseDm}{3.73}  % DM vs the reference arm (subsample)
\newcommand{\rhabRidgeFullLiveVdCloseP}{\ensuremath{<}0.001}  % p-value
\newcommand{\rhabRidgeLoVdLiveVdCloseQlike}{0.1048}  % QLIKE ridge lo_vd live16_vd (subsample)
\newcommand{\rhabRidgeLoVdLiveVdClosePct}{$-$17.4}  % QLIKE change % vs the reference (subsample)
\newcommand{\rhabRidgeLoVdLiveVdCloseDm}{$-$4.97}  % DM vs the reference arm (subsample)
\newcommand{\rhabRidgeLoVdLiveVdCloseP}{\ensuremath{<}0.001}  % p-value
\newcommand{\rhabRidgeAddVdLiveVdCloseQlike}{0.1018}  % QLIKE ridge add_vd live16_vd (subsample)
\newcommand{\rhabRidgeAddVdLiveVdClosePct}{17.0}  % QLIKE change % vs the reference (subsample)
\newcommand{\rhabRidgeAddVdLiveVdCloseDm}{3.66}  % DM vs the reference arm (subsample)
\newcommand{\rhabRidgeAddVdLiveVdCloseP}{\ensuremath{<}0.001}  % p-value
\newcommand{\rhabRidgeHarLiveVdCloseQlike}{0.0870}  % QLIKE ridge har live16_vd (subsample)
\newcommand{\rhabLassoFullLiveVdCloseQlike}{0.1271}  % QLIKE lasso full live16_vd (subsample)
\newcommand{\rhabLassoFullLiveVdClosePct}{46.0}  % QLIKE change % vs the reference (subsample)
\newcommand{\rhabLassoFullLiveVdCloseDm}{3.87}  % DM vs the reference arm (subsample)
\newcommand{\rhabLassoFullLiveVdCloseP}{\ensuremath{<}0.001}  % p-value
\newcommand{\rhabLassoLoVdLiveVdCloseQlike}{0.1097}  % QLIKE lasso lo_vd live16_vd (subsample)
\newcommand{\rhabLassoLoVdLiveVdClosePct}{$-$13.7}  % QLIKE change % vs the reference (subsample)
\newcommand{\rhabLassoLoVdLiveVdCloseDm}{$-$4.18}  % DM vs the reference arm (subsample)
\newcommand{\rhabLassoLoVdLiveVdCloseP}{\ensuremath{<}0.001}  % p-value
\newcommand{\rhabLassoAddVdLiveVdCloseQlike}{0.1018}  % QLIKE lasso add_vd live16_vd (subsample)
\newcommand{\rhabLassoAddVdLiveVdClosePct}{17.0}  % QLIKE change % vs the reference (subsample)
\newcommand{\rhabLassoAddVdLiveVdCloseDm}{3.75}  % DM vs the reference arm (subsample)
\newcommand{\rhabLassoAddVdLiveVdCloseP}{\ensuremath{<}0.001}  % p-value
\newcommand{\rhabLassoHarLiveVdCloseQlike}{0.0870}  % QLIKE lasso har live16_vd (subsample)
\newcommand{\rhabEnetFullLiveVdCloseQlike}{0.1227}  % QLIKE enet full live16_vd (subsample)
\newcommand{\rhabEnetFullLiveVdClosePct}{41.0}  % QLIKE change % vs the reference (subsample)
\newcommand{\rhabEnetFullLiveVdCloseDm}{3.59}  % DM vs the reference arm (subsample)
\newcommand{\rhabEnetFullLiveVdCloseP}{\ensuremath{<}0.001}  % p-value
\newcommand{\rhabEnetLoVdLiveVdCloseQlike}{0.1029}  % QLIKE enet lo_vd live16_vd (subsample)
\newcommand{\rhabEnetLoVdLiveVdClosePct}{$-$16.2}  % QLIKE change % vs the reference (subsample)
\newcommand{\rhabEnetLoVdLiveVdCloseDm}{$-$4.00}  % DM vs the reference arm (subsample)
\newcommand{\rhabEnetLoVdLiveVdCloseP}{\ensuremath{<}0.001}  % p-value
\newcommand{\rhabEnetAddVdLiveVdCloseQlike}{0.1005}  % QLIKE enet add_vd live16_vd (subsample)
\newcommand{\rhabEnetAddVdLiveVdClosePct}{15.5}  % QLIKE change % vs the reference (subsample)
\newcommand{\rhabEnetAddVdLiveVdCloseDm}{3.55}  % DM vs the reference arm (subsample)
\newcommand{\rhabEnetAddVdLiveVdCloseP}{\ensuremath{<}0.001}  % p-value
\newcommand{\rhabEnetHarLiveVdCloseQlike}{0.0870}  % QLIKE enet har live16_vd (subsample)
\newcommand{\rhabLgbmFullLiveVdCloseQlike}{0.0802}  % QLIKE lgbm full live16_vd (subsample)
\newcommand{\rhabLgbmFullLiveVdClosePct}{$-$1.5}  % QLIKE change % vs the reference (subsample)
\newcommand{\rhabLgbmFullLiveVdCloseDm}{$-$0.29}  % DM vs the reference arm (subsample)
\newcommand{\rhabLgbmFullLiveVdCloseP}{0.769}  % p-value
\newcommand{\rhabLgbmHarLiveVdCloseQlike}{0.0814}  % QLIKE lgbm har live16_vd (subsample)
\newcommand{\rhabXgbFullLiveVdCloseQlike}{0.0785}  % QLIKE xgb full live16_vd (subsample)
\newcommand{\rhabXgbFullLiveVdClosePct}{$-$4.3}  % QLIKE change % vs the reference (subsample)
\newcommand{\rhabXgbFullLiveVdCloseDm}{$-$0.89}  % DM vs the reference arm (subsample)
\newcommand{\rhabXgbFullLiveVdCloseP}{0.375}  % p-value
\newcommand{\rhabXgbHarLiveVdCloseQlike}{0.0820}  % QLIKE xgb har live16_vd (subsample)
\newcommand{\rhabRfFullLiveVdCloseQlike}{0.0822}  % QLIKE rf full live16_vd (subsample)
\newcommand{\rhabRfFullLiveVdClosePct}{$-$1.0}  % QLIKE change % vs the reference (subsample)
\newcommand{\rhabRfFullLiveVdCloseDm}{$-$0.21}  % DM vs the reference arm (subsample)
\newcommand{\rhabRfFullLiveVdCloseP}{0.833}  % p-value
\newcommand{\rhabRfHarLiveVdCloseQlike}{0.0831}  % QLIKE rf har live16_vd (subsample)
\newcommand{\rhabLiveSentStart}{19~April 2022}  % first bar of the live window sent
\newcommand{\rhabLiveSentEnd}{30~April 2024}  % last bar of the live window sent
\newcommand{\rhabLiveSentDays}{508}  % trading days of the live window sent
\newcommand{\rhabLiveIvStart}{19~April 2018}  % first bar of the live window iv
\newcommand{\rhabLiveIvEnd}{30~April 2024}  % last bar of the live window iv
\newcommand{\rhabLiveIvDays}{1{,}506}  % trading days of the live window iv
\newcommand{\rhabLiveIvbStart}{6~July 2021}  % first bar of the live window ivb
\newcommand{\rhabLiveIvbEnd}{30~April 2024}  % last bar of the live window ivb
\newcommand{\rhabLiveIvbDays}{706}  % trading days of the live window ivb
\newcommand{\rhabLiveVdStart}{16~November 2022}  % first bar of the live window vd
\newcommand{\rhabLiveVdEnd}{31~August 2023}  % last bar of the live window vd
\newcommand{\rhabLiveVdDays}{196}  % trading days of the live window vd
\newcommand{\rhabRidgeFullVsOlsAddRmAllDm}{$-$4.70}  % DM ridge full vs OLS add_rm
\newcommand{\rhabRidgeFullVsOlsAddRmAllP}{\ensuremath{<}0.001}  % p-value
\newcommand{\rhabRidgeFullVsOlsAddRmCloseDm}{1.16}  % DM ridge full vs OLS add_rm
\newcommand{\rhabRidgeFullVsOlsAddRmCloseP}{0.244}  % p-value
\newcommand{\rhabRidgeFullVsOlsAddRmCloseLateDm}{0.95}  % DM ridge full vs OLS add_rm
\newcommand{\rhabRidgeFullVsOlsAddRmCloseLateP}{0.341}  % p-value
\newcommand{\rhabRidgeFullVsOlsSharAllDm}{$-$4.77}  % DM ridge full vs OLS shar
\newcommand{\rhabRidgeFullVsOlsSharAllP}{\ensuremath{<}0.001}  % p-value
\newcommand{\rhabRidgeFullVsOlsSharCloseDm}{0.41}  % DM ridge full vs OLS shar
\newcommand{\rhabRidgeFullVsOlsSharCloseP}{0.682}  % p-value
\newcommand{\rhabRidgeFullVsOlsSharCloseLateDm}{2.34}  % DM ridge full vs OLS shar
\newcommand{\rhabRidgeFullVsOlsSharCloseLateP}{0.019}  % p-value
\newcommand{\rhabRidgeFullVsOlsHarqAllDm}{$-$5.66}  % DM ridge full vs OLS harq
\newcommand{\rhabRidgeFullVsOlsHarqAllP}{\ensuremath{<}0.001}  % p-value
\newcommand{\rhabRidgeFullVsOlsHarqCloseDm}{0.37}  % DM ridge full vs OLS harq
\newcommand{\rhabRidgeFullVsOlsHarqCloseP}{0.715}  % p-value
\newcommand{\rhabRidgeFullVsOlsHarqCloseLateDm}{2.35}  % DM ridge full vs OLS harq
\newcommand{\rhabRidgeFullVsOlsHarqCloseLateP}{0.019}  % p-value
\newcommand{\rhabRidgeFullVsOlsLharAllDm}{$-$3.23}  % DM ridge full vs OLS lhar
\newcommand{\rhabRidgeFullVsOlsLharAllP}{0.001}  % p-value
\newcommand{\rhabRidgeFullVsOlsLharCloseDm}{2.40}  % DM ridge full vs OLS lhar
\newcommand{\rhabRidgeFullVsOlsLharCloseP}{0.016}  % p-value
\newcommand{\rhabRidgeFullVsOlsLharCloseLateDm}{2.96}  % DM ridge full vs OLS lhar
\newcommand{\rhabRidgeFullVsOlsLharCloseLateP}{0.003}  % p-value
\newcommand{\rhabRidgeFullVsOlsSharqlAllDm}{$-$3.04}  % DM ridge full vs OLS sharql
\newcommand{\rhabRidgeFullVsOlsSharqlAllP}{0.002}  % p-value
\newcommand{\rhabRidgeFullVsOlsSharqlCloseDm}{2.44}  % DM ridge full vs OLS sharql
\newcommand{\rhabRidgeFullVsOlsSharqlCloseP}{0.015}  % p-value
\newcommand{\rhabRidgeFullVsOlsSharqlCloseLateDm}{3.13}  % DM ridge full vs OLS sharql
\newcommand{\rhabRidgeFullVsOlsSharqlCloseLateP}{0.002}  % p-value
\newcommand{\rhabRidgeFullVsOlsHarAllDm}{$-$5.63}  % DM ridge full vs OLS har
\newcommand{\rhabRidgeFullVsOlsHarAllP}{\ensuremath{<}0.001}  % p-value
\newcommand{\rhabRidgeFullVsOlsHarCloseDm}{0.31}  % DM ridge full vs OLS har
\newcommand{\rhabRidgeFullVsOlsHarCloseP}{0.759}  % p-value
\newcommand{\rhabRidgeFullVsOlsHarCloseLateDm}{2.31}  % DM ridge full vs OLS har
\newcommand{\rhabRidgeFullVsOlsHarCloseLateP}{0.021}  % p-value
\newcommand{\rhabLassoFullVsOlsAddRmAllDm}{$-$4.53}  % DM lasso full vs OLS add_rm
\newcommand{\rhabLassoFullVsOlsAddRmAllP}{\ensuremath{<}0.001}  % p-value
\newcommand{\rhabLassoFullVsOlsAddRmCloseDm}{2.40}  % DM lasso full vs OLS add_rm
\newcommand{\rhabLassoFullVsOlsAddRmCloseP}{0.016}  % p-value
\newcommand{\rhabLassoFullVsOlsAddRmCloseLateDm}{2.26}  % DM lasso full vs OLS add_rm
\newcommand{\rhabLassoFullVsOlsAddRmCloseLateP}{0.024}  % p-value
\newcommand{\rhabLassoFullVsOlsSharAllDm}{$-$4.07}  % DM lasso full vs OLS shar
\newcommand{\rhabLassoFullVsOlsSharAllP}{\ensuremath{<}0.001}  % p-value
\newcommand{\rhabLassoFullVsOlsSharCloseDm}{1.12}  % DM lasso full vs OLS shar
\newcommand{\rhabLassoFullVsOlsSharCloseP}{0.261}  % p-value
\newcommand{\rhabLassoFullVsOlsSharCloseLateDm}{2.54}  % DM lasso full vs OLS shar
\newcommand{\rhabLassoFullVsOlsSharCloseLateP}{0.011}  % p-value
\newcommand{\rhabLassoFullVsOlsHarqAllDm}{$-$4.75}  % DM lasso full vs OLS harq
\newcommand{\rhabLassoFullVsOlsHarqAllP}{\ensuremath{<}0.001}  % p-value
\newcommand{\rhabLassoFullVsOlsHarqCloseDm}{1.08}  % DM lasso full vs OLS harq
\newcommand{\rhabLassoFullVsOlsHarqCloseP}{0.279}  % p-value
\newcommand{\rhabLassoFullVsOlsHarqCloseLateDm}{2.55}  % DM lasso full vs OLS harq
\newcommand{\rhabLassoFullVsOlsHarqCloseLateP}{0.011}  % p-value
\newcommand{\rhabLassoFullVsOlsLharAllDm}{$-$2.95}  % DM lasso full vs OLS lhar
\newcommand{\rhabLassoFullVsOlsLharAllP}{0.003}  % p-value
\newcommand{\rhabLassoFullVsOlsLharCloseDm}{3.00}  % DM lasso full vs OLS lhar
\newcommand{\rhabLassoFullVsOlsLharCloseP}{0.003}  % p-value
\newcommand{\rhabLassoFullVsOlsLharCloseLateDm}{3.28}  % DM lasso full vs OLS lhar
\newcommand{\rhabLassoFullVsOlsLharCloseLateP}{0.001}  % p-value
\newcommand{\rhabLassoFullVsOlsSharqlAllDm}{$-$2.81}  % DM lasso full vs OLS sharql
\newcommand{\rhabLassoFullVsOlsSharqlAllP}{0.005}  % p-value
\newcommand{\rhabLassoFullVsOlsSharqlCloseDm}{2.96}  % DM lasso full vs OLS sharql
\newcommand{\rhabLassoFullVsOlsSharqlCloseP}{0.003}  % p-value
\newcommand{\rhabLassoFullVsOlsSharqlCloseLateDm}{3.33}  % DM lasso full vs OLS sharql
\newcommand{\rhabLassoFullVsOlsSharqlCloseLateP}{\ensuremath{<}0.001}  % p-value
\newcommand{\rhabLassoFullVsOlsHarAllDm}{$-$4.72}  % DM lasso full vs OLS har
\newcommand{\rhabLassoFullVsOlsHarAllP}{\ensuremath{<}0.001}  % p-value
\newcommand{\rhabLassoFullVsOlsHarCloseDm}{1.04}  % DM lasso full vs OLS har
\newcommand{\rhabLassoFullVsOlsHarCloseP}{0.300}  % p-value
\newcommand{\rhabLassoFullVsOlsHarCloseLateDm}{2.51}  % DM lasso full vs OLS har
\newcommand{\rhabLassoFullVsOlsHarCloseLateP}{0.012}  % p-value
\newcommand{\rhabEnetFullVsOlsAddRmAllDm}{$-$5.61}  % DM enet full vs OLS add_rm
\newcommand{\rhabEnetFullVsOlsAddRmAllP}{\ensuremath{<}0.001}  % p-value
\newcommand{\rhabEnetFullVsOlsAddRmCloseDm}{0.46}  % DM enet full vs OLS add_rm
\newcommand{\rhabEnetFullVsOlsAddRmCloseP}{0.649}  % p-value
\newcommand{\rhabEnetFullVsOlsAddRmCloseLateDm}{0.29}  % DM enet full vs OLS add_rm
\newcommand{\rhabEnetFullVsOlsAddRmCloseLateP}{0.770}  % p-value
\newcommand{\rhabEnetFullVsOlsSharAllDm}{$-$5.97}  % DM enet full vs OLS shar
\newcommand{\rhabEnetFullVsOlsSharAllP}{\ensuremath{<}0.001}  % p-value
\newcommand{\rhabEnetFullVsOlsSharCloseDm}{$-$0.11}  % DM enet full vs OLS shar
\newcommand{\rhabEnetFullVsOlsSharCloseP}{0.914}  % p-value
\newcommand{\rhabEnetFullVsOlsSharCloseLateDm}{1.92}  % DM enet full vs OLS shar
\newcommand{\rhabEnetFullVsOlsSharCloseLateP}{0.055}  % p-value
\newcommand{\rhabEnetFullVsOlsHarqAllDm}{$-$7.03}  % DM enet full vs OLS harq
\newcommand{\rhabEnetFullVsOlsHarqAllP}{\ensuremath{<}0.001}  % p-value
\newcommand{\rhabEnetFullVsOlsHarqCloseDm}{$-$0.15}  % DM enet full vs OLS harq
\newcommand{\rhabEnetFullVsOlsHarqCloseP}{0.883}  % p-value
\newcommand{\rhabEnetFullVsOlsHarqCloseLateDm}{1.96}  % DM enet full vs OLS harq
\newcommand{\rhabEnetFullVsOlsHarqCloseLateP}{0.050}  % p-value
\newcommand{\rhabEnetFullVsOlsLharAllDm}{$-$4.58}  % DM enet full vs OLS lhar
\newcommand{\rhabEnetFullVsOlsLharAllP}{\ensuremath{<}0.001}  % p-value
\newcommand{\rhabEnetFullVsOlsLharCloseDm}{1.86}  % DM enet full vs OLS lhar
\newcommand{\rhabEnetFullVsOlsLharCloseP}{0.063}  % p-value
\newcommand{\rhabEnetFullVsOlsLharCloseLateDm}{2.46}  % DM enet full vs OLS lhar
\newcommand{\rhabEnetFullVsOlsLharCloseLateP}{0.014}  % p-value
\newcommand{\rhabEnetFullVsOlsSharqlAllDm}{$-$4.34}  % DM enet full vs OLS sharql
\newcommand{\rhabEnetFullVsOlsSharqlAllP}{\ensuremath{<}0.001}  % p-value
\newcommand{\rhabEnetFullVsOlsSharqlCloseDm}{1.92}  % DM enet full vs OLS sharql
\newcommand{\rhabEnetFullVsOlsSharqlCloseP}{0.055}  % p-value
\newcommand{\rhabEnetFullVsOlsSharqlCloseLateDm}{2.65}  % DM enet full vs OLS sharql
\newcommand{\rhabEnetFullVsOlsSharqlCloseLateP}{0.008}  % p-value
\newcommand{\rhabEnetFullVsOlsHarAllDm}{$-$6.98}  % DM enet full vs OLS har
\newcommand{\rhabEnetFullVsOlsHarAllP}{\ensuremath{<}0.001}  % p-value
\newcommand{\rhabEnetFullVsOlsHarCloseDm}{$-$0.21}  % DM enet full vs OLS har
\newcommand{\rhabEnetFullVsOlsHarCloseP}{0.832}  % p-value
\newcommand{\rhabEnetFullVsOlsHarCloseLateDm}{1.91}  % DM enet full vs OLS har
\newcommand{\rhabEnetFullVsOlsHarCloseLateP}{0.056}  % p-value
\newcommand{\rhabLgbmFullVsOlsAddRmAllDm}{$-$7.94}  % DM lgbm full vs OLS add_rm
\newcommand{\rhabLgbmFullVsOlsAddRmAllP}{\ensuremath{<}0.001}  % p-value
\newcommand{\rhabLgbmFullVsOlsAddRmCloseDm}{$-$1.78}  % DM lgbm full vs OLS add_rm
\newcommand{\rhabLgbmFullVsOlsAddRmCloseP}{0.076}  % p-value
\newcommand{\rhabLgbmFullVsOlsAddRmCloseLateDm}{$-$2.44}  % DM lgbm full vs OLS add_rm
\newcommand{\rhabLgbmFullVsOlsAddRmCloseLateP}{0.015}  % p-value
\newcommand{\rhabLgbmFullVsOlsSharAllDm}{$-$9.26}  % DM lgbm full vs OLS shar
\newcommand{\rhabLgbmFullVsOlsSharAllP}{\ensuremath{<}0.001}  % p-value
\newcommand{\rhabLgbmFullVsOlsSharCloseDm}{$-$3.86}  % DM lgbm full vs OLS shar
\newcommand{\rhabLgbmFullVsOlsSharCloseP}{\ensuremath{<}0.001}  % p-value
\newcommand{\rhabLgbmFullVsOlsSharCloseLateDm}{$-$2.34}  % DM lgbm full vs OLS shar
\newcommand{\rhabLgbmFullVsOlsSharCloseLateP}{0.019}  % p-value
\newcommand{\rhabLgbmFullVsOlsHarqAllDm}{$-$10.44}  % DM lgbm full vs OLS harq
\newcommand{\rhabLgbmFullVsOlsHarqAllP}{\ensuremath{<}0.001}  % p-value
\newcommand{\rhabLgbmFullVsOlsHarqCloseDm}{$-$3.77}  % DM lgbm full vs OLS harq
\newcommand{\rhabLgbmFullVsOlsHarqCloseP}{\ensuremath{<}0.001}  % p-value
\newcommand{\rhabLgbmFullVsOlsHarqCloseLateDm}{$-$1.70}  % DM lgbm full vs OLS harq
\newcommand{\rhabLgbmFullVsOlsHarqCloseLateP}{0.088}  % p-value
\newcommand{\rhabLgbmFullVsOlsLharAllDm}{$-$8.32}  % DM lgbm full vs OLS lhar
\newcommand{\rhabLgbmFullVsOlsLharAllP}{\ensuremath{<}0.001}  % p-value
\newcommand{\rhabLgbmFullVsOlsLharCloseDm}{$-$1.02}  % DM lgbm full vs OLS lhar
\newcommand{\rhabLgbmFullVsOlsLharCloseP}{0.310}  % p-value
\newcommand{\rhabLgbmFullVsOlsLharCloseLateDm}{$-$1.42}  % DM lgbm full vs OLS lhar
\newcommand{\rhabLgbmFullVsOlsLharCloseLateP}{0.155}  % p-value
\newcommand{\rhabLgbmFullVsOlsSharqlAllDm}{$-$8.22}  % DM lgbm full vs OLS sharql
\newcommand{\rhabLgbmFullVsOlsSharqlAllP}{\ensuremath{<}0.001}  % p-value
\newcommand{\rhabLgbmFullVsOlsSharqlCloseDm}{$-$1.03}  % DM lgbm full vs OLS sharql
\newcommand{\rhabLgbmFullVsOlsSharqlCloseP}{0.302}  % p-value
\newcommand{\rhabLgbmFullVsOlsSharqlCloseLateDm}{$-$1.45}  % DM lgbm full vs OLS sharql
\newcommand{\rhabLgbmFullVsOlsSharqlCloseLateP}{0.148}  % p-value
\newcommand{\rhabLgbmFullVsOlsHarAllDm}{$-$10.41}  % DM lgbm full vs OLS har
\newcommand{\rhabLgbmFullVsOlsHarAllP}{\ensuremath{<}0.001}  % p-value
\newcommand{\rhabLgbmFullVsOlsHarCloseDm}{$-$3.89}  % DM lgbm full vs OLS har
\newcommand{\rhabLgbmFullVsOlsHarCloseP}{\ensuremath{<}0.001}  % p-value
\newcommand{\rhabLgbmFullVsOlsHarCloseLateDm}{$-$1.86}  % DM lgbm full vs OLS har
\newcommand{\rhabLgbmFullVsOlsHarCloseLateP}{0.063}  % p-value
\newcommand{\rhabXgbFullVsOlsAddRmAllDm}{$-$7.47}  % DM xgb full vs OLS add_rm
\newcommand{\rhabXgbFullVsOlsAddRmAllP}{\ensuremath{<}0.001}  % p-value
\newcommand{\rhabXgbFullVsOlsAddRmCloseDm}{$-$2.16}  % DM xgb full vs OLS add_rm
\newcommand{\rhabXgbFullVsOlsAddRmCloseP}{0.031}  % p-value
\newcommand{\rhabXgbFullVsOlsAddRmCloseLateDm}{$-$2.61}  % DM xgb full vs OLS add_rm
\newcommand{\rhabXgbFullVsOlsAddRmCloseLateP}{0.009}  % p-value
\newcommand{\rhabXgbFullVsOlsSharAllDm}{$-$8.81}  % DM xgb full vs OLS shar
\newcommand{\rhabXgbFullVsOlsSharAllP}{\ensuremath{<}0.001}  % p-value
\newcommand{\rhabXgbFullVsOlsSharCloseDm}{$-$4.87}  % DM xgb full vs OLS shar
\newcommand{\rhabXgbFullVsOlsSharCloseP}{\ensuremath{<}0.001}  % p-value
\newcommand{\rhabXgbFullVsOlsSharCloseLateDm}{$-$2.99}  % DM xgb full vs OLS shar
\newcommand{\rhabXgbFullVsOlsSharCloseLateP}{0.003}  % p-value
\newcommand{\rhabXgbFullVsOlsHarqAllDm}{$-$9.96}  % DM xgb full vs OLS harq
\newcommand{\rhabXgbFullVsOlsHarqAllP}{\ensuremath{<}0.001}  % p-value
\newcommand{\rhabXgbFullVsOlsHarqCloseDm}{$-$4.77}  % DM xgb full vs OLS harq
\newcommand{\rhabXgbFullVsOlsHarqCloseP}{\ensuremath{<}0.001}  % p-value
\newcommand{\rhabXgbFullVsOlsHarqCloseLateDm}{$-$2.27}  % DM xgb full vs OLS harq
\newcommand{\rhabXgbFullVsOlsHarqCloseLateP}{0.023}  % p-value
\newcommand{\rhabXgbFullVsOlsLharAllDm}{$-$7.85}  % DM xgb full vs OLS lhar
\newcommand{\rhabXgbFullVsOlsLharAllP}{\ensuremath{<}0.001}  % p-value
\newcommand{\rhabXgbFullVsOlsLharCloseDm}{$-$1.70}  % DM xgb full vs OLS lhar
\newcommand{\rhabXgbFullVsOlsLharCloseP}{0.090}  % p-value
\newcommand{\rhabXgbFullVsOlsLharCloseLateDm}{$-$1.85}  % DM xgb full vs OLS lhar
\newcommand{\rhabXgbFullVsOlsLharCloseLateP}{0.065}  % p-value
\newcommand{\rhabXgbFullVsOlsSharqlAllDm}{$-$7.76}  % DM xgb full vs OLS sharql
\newcommand{\rhabXgbFullVsOlsSharqlAllP}{\ensuremath{<}0.001}  % p-value
\newcommand{\rhabXgbFullVsOlsSharqlCloseDm}{$-$1.82}  % DM xgb full vs OLS sharql
\newcommand{\rhabXgbFullVsOlsSharqlCloseP}{0.069}  % p-value
\newcommand{\rhabXgbFullVsOlsSharqlCloseLateDm}{$-$1.95}  % DM xgb full vs OLS sharql
\newcommand{\rhabXgbFullVsOlsSharqlCloseLateP}{0.051}  % p-value
\newcommand{\rhabXgbFullVsOlsHarAllDm}{$-$9.93}  % DM xgb full vs OLS har
\newcommand{\rhabXgbFullVsOlsHarAllP}{\ensuremath{<}0.001}  % p-value
\newcommand{\rhabXgbFullVsOlsHarCloseDm}{$-$4.88}  % DM xgb full vs OLS har
\newcommand{\rhabXgbFullVsOlsHarCloseP}{\ensuremath{<}0.001}  % p-value
\newcommand{\rhabXgbFullVsOlsHarCloseLateDm}{$-$2.44}  % DM xgb full vs OLS har
\newcommand{\rhabXgbFullVsOlsHarCloseLateP}{0.015}  % p-value
\newcommand{\rhabRfFullVsOlsAddRmAllDm}{$-$6.68}  % DM rf full vs OLS add_rm
\newcommand{\rhabRfFullVsOlsAddRmAllP}{\ensuremath{<}0.001}  % p-value
\newcommand{\rhabRfFullVsOlsAddRmCloseDm}{$-$0.14}  % DM rf full vs OLS add_rm
\newcommand{\rhabRfFullVsOlsAddRmCloseP}{0.892}  % p-value
\newcommand{\rhabRfFullVsOlsAddRmCloseLateDm}{$-$1.75}  % DM rf full vs OLS add_rm
\newcommand{\rhabRfFullVsOlsAddRmCloseLateP}{0.080}  % p-value
\newcommand{\rhabRfFullVsOlsSharAllDm}{$-$8.45}  % DM rf full vs OLS shar
\newcommand{\rhabRfFullVsOlsSharAllP}{\ensuremath{<}0.001}  % p-value
\newcommand{\rhabRfFullVsOlsSharCloseDm}{$-$0.98}  % DM rf full vs OLS shar
\newcommand{\rhabRfFullVsOlsSharCloseP}{0.328}  % p-value
\newcommand{\rhabRfFullVsOlsSharCloseLateDm}{$-$0.76}  % DM rf full vs OLS shar
\newcommand{\rhabRfFullVsOlsSharCloseLateP}{0.444}  % p-value
\newcommand{\rhabRfFullVsOlsHarqAllDm}{$-$9.73}  % DM rf full vs OLS harq
\newcommand{\rhabRfFullVsOlsHarqAllP}{\ensuremath{<}0.001}  % p-value
\newcommand{\rhabRfFullVsOlsHarqCloseDm}{$-$1.05}  % DM rf full vs OLS harq
\newcommand{\rhabRfFullVsOlsHarqCloseP}{0.293}  % p-value
\newcommand{\rhabRfFullVsOlsHarqCloseLateDm}{$-$0.36}  % DM rf full vs OLS harq
\newcommand{\rhabRfFullVsOlsHarqCloseLateP}{0.715}  % p-value
\newcommand{\rhabRfFullVsOlsLharAllDm}{$-$7.08}  % DM rf full vs OLS lhar
\newcommand{\rhabRfFullVsOlsLharAllP}{\ensuremath{<}0.001}  % p-value
\newcommand{\rhabRfFullVsOlsLharCloseDm}{1.54}  % DM rf full vs OLS lhar
\newcommand{\rhabRfFullVsOlsLharCloseP}{0.124}  % p-value
\newcommand{\rhabRfFullVsOlsLharCloseLateDm}{$-$0.26}  % DM rf full vs OLS lhar
\newcommand{\rhabRfFullVsOlsLharCloseLateP}{0.794}  % p-value
\newcommand{\rhabRfFullVsOlsSharqlAllDm}{$-$6.99}  % DM rf full vs OLS sharql
\newcommand{\rhabRfFullVsOlsSharqlAllP}{\ensuremath{<}0.001}  % p-value
\newcommand{\rhabRfFullVsOlsSharqlCloseDm}{1.76}  % DM rf full vs OLS sharql
\newcommand{\rhabRfFullVsOlsSharqlCloseP}{0.078}  % p-value
\newcommand{\rhabRfFullVsOlsSharqlCloseLateDm}{$-$0.09}  % DM rf full vs OLS sharql
\newcommand{\rhabRfFullVsOlsSharqlCloseLateP}{0.926}  % p-value
\newcommand{\rhabRfFullVsOlsHarAllDm}{$-$9.69}  % DM rf full vs OLS har
\newcommand{\rhabRfFullVsOlsHarAllP}{\ensuremath{<}0.001}  % p-value
\newcommand{\rhabRfFullVsOlsHarCloseDm}{$-$1.16}  % DM rf full vs OLS har
\newcommand{\rhabRfFullVsOlsHarCloseP}{0.246}  % p-value
\newcommand{\rhabRfFullVsOlsHarCloseLateDm}{$-$0.49}  % DM rf full vs OLS har
\newcommand{\rhabRfFullVsOlsHarCloseLateP}{0.621}  % p-value
\newcommand{\rhabNClassesAddRm}{three}  % classes with an add_rm arm
\newcommand{\rhabNClassesBeatAddRm}{none}  % full set beats its own add_rm, DM p < 0.05
\newcommand{\rhabBestOlsBench}{SHAR + HARQ + LHAR}  % lowest QLIKE OLS augmented HAR
\newcommand{\rhabLiveAttnFirst}{23~March 2009}  % first print stocktwits_attention
\newcommand{\rhabLiveAttnLast}{30~April 2024}  % last print stocktwits_attention
\newcommand{\rhabLiveAttnFromKone}{11~February 2021}  % first non-zeroed value, k = 1
\newcommand{\rhabLiveAttnFromKlast}{13~May 2021}  % the same, k = 3125
\newcommand{\rhabLiveAttnSharePct}{31}  % % of evaluation bars live, k = 1
\newcommand{\rhabLiveNSentFirst}{1~June 2010}  % first print stocktwits_sentcount
\newcommand{\rhabLiveNSentLast}{30~April 2024}  % last print stocktwits_sentcount
\newcommand{\rhabLiveNSentFromKone}{19~April 2022}  % first non-zeroed value, k = 1
\newcommand{\rhabLiveNSentFromKlast}{19~July 2022}  % the same, k = 3125
\newcommand{\rhabLiveNSentSharePct}{20}  % % of evaluation bars live, k = 1
\newcommand{\rhabLiveSentFirst}{1~June 2010}  % first print stocktwits_sentiment
\newcommand{\rhabLiveSentLast}{30~April 2024}  % last print stocktwits_sentiment
\newcommand{\rhabLiveSentFromKone}{19~April 2022}  % first non-zeroed value, k = 1
\newcommand{\rhabLiveSentFromKlast}{19~July 2022}  % the same, k = 3125
\newcommand{\rhabLiveSentSharePct}{20}  % % of evaluation bars live, k = 1
\newcommand{\rhabLiveVixFirst}{22~May 2006}  % first print vix
\newcommand{\rhabLiveVixLast}{30~April 2024}  % last print vix
\newcommand{\rhabLiveVixFromKone}{19~April 2018}  % first non-zeroed value, k = 1
\newcommand{\rhabLiveVixFromKlast}{18~July 2018}  % the same, k = 3125
\newcommand{\rhabLiveVixSharePct}{58}  % % of evaluation bars live, k = 1
\newcommand{\rhabLiveVixthreemFirst}{12~August 2009}  % first print vix3m
\newcommand{\rhabLiveVixthreemLast}{30~April 2024}  % last print vix3m
\newcommand{\rhabLiveVixthreemFromKone}{6~July 2021}  % first non-zeroed value, k = 1
\newcommand{\rhabLiveVixthreemFromKlast}{1~October 2021}  % the same, k = 3125
\newcommand{\rhabLiveVixthreemSharePct}{27}  % % of evaluation bars live, k = 1
\newcommand{\rhabLiveVvixFirst}{26~March 2012}  % first print vvix
\newcommand{\rhabLiveVvixLast}{30~April 2024}  % last print vvix
\newcommand{\rhabLiveVvixFromKone}{12~February 2024}  % first non-zeroed value, k = 1
\newcommand{\rhabLiveVvixFromKlast}{never}  % the same, k = 3125
\newcommand{\rhabLiveVvixSharePct}{2}  % % of evaluation bars live, k = 1
\newcommand{\rhabLiveVdAllOcFirst}{3~January 2011}  % first print voldemand_all_open_and_close
\newcommand{\rhabLiveVdAllOcLast}{31~August 2023}  % last print voldemand_all_open_and_close
\newcommand{\rhabLiveVdAllOcFromKone}{16~November 2022}  % first non-zeroed value, k = 1
\newcommand{\rhabLiveVdAllOcFromKlast}{20~February 2023}  % the same, k = 3125
\newcommand{\rhabLiveVdAllOcSharePct}{8}  % % of evaluation bars live, k = 1
\newcommand{\rhabLiveVdAllOFirst}{3~January 2011}  % first print voldemand_all_open_only
\newcommand{\rhabLiveVdAllOLast}{31~August 2023}  % last print voldemand_all_open_only
\newcommand{\rhabLiveVdAllOFromKone}{16~November 2022}  % first non-zeroed value, k = 1
\newcommand{\rhabLiveVdAllOFromKlast}{20~February 2023}  % the same, k = 3125
\newcommand{\rhabLiveVdAllOSharePct}{8}  % % of evaluation bars live, k = 1
\newcommand{\rhabLiveVdSpxOcFirst}{3~January 2011}  % first print voldemand_spx_open_and_close
\newcommand{\rhabLiveVdSpxOcLast}{31~August 2023}  % last print voldemand_spx_open_and_close
\newcommand{\rhabLiveVdSpxOcFromKone}{16~November 2022}  % first non-zeroed value, k = 1
\newcommand{\rhabLiveVdSpxOcFromKlast}{20~February 2023}  % the same, k = 3125
\newcommand{\rhabLiveVdSpxOcSharePct}{8}  % % of evaluation bars live, k = 1
\newcommand{\rhabLiveVdSpxOFirst}{3~January 2011}  % first print voldemand_spx_open_only
\newcommand{\rhabLiveVdSpxOLast}{31~August 2023}  % last print voldemand_spx_open_only
\newcommand{\rhabLiveVdSpxOFromKone}{16~November 2022}  % first non-zeroed value, k = 1
\newcommand{\rhabLiveVdSpxOFromKlast}{20~February 2023}  % the same, k = 3125
\newcommand{\rhabLiveVdSpxOSharePct}{8}  % % of evaluation bars live, k = 1
\fi
\ifdefined\rhcaNdaysSpan\else% AUTO-GENERATED by writeup/make_rehaul_close_accuracy.py -- do not edit; re-run the script.
% \rhca... macros: closing-half-hour accuracy on both spans (tab:rh_close_accuracy) and the calibration
% diagnostics of the 12 pooled forecasts (tab:rh_calibration, fig:rh_calibration). Names: <Stat><Model><Set>
% <Sample><Span>; Model Ridge Lasso Enet Lgbm Xgb Rf Ols (pooled), Ci<Model> (trained on the closing half
% hour), Bm (closing-interval HAR benchmark); Set Har Full; Sample Close (16:00 bar) Bar (every bar);
% Span Span (2014-2024) Pre (2014-2019) Oos (T_oos, 2020-2024).
% --- spans and constants
\providecommand{\rhcaNdaysSpan}{}\renewcommand{\rhcaNdaysSpan}{2{,}579}  % trading days, 2014-01-02 .. 2024-04-30
\providecommand{\rhcaStartSpan}{}\renewcommand{\rhcaStartSpan}{2 January 2014}  % first day
\providecommand{\rhcaEndSpan}{}\renewcommand{\rhcaEndSpan}{30 April 2024}  % last day
\providecommand{\rhcaYearsSpan}{}\renewcommand{\rhcaYearsSpan}{2014--2024}  % years
\providecommand{\rhcaDMLagSpan}{}\renewcommand{\rhcaDMLagSpan}{8}  % Newey-West lag of the daily DM tests (dm_test automatic)
\providecommand{\rhcaHacLagBarSpan}{}\renewcommand{\rhcaHacLagBarSpan}{104}  % Newey-West lag of the every-bar MZ regressions, bars
\providecommand{\rhcaNdaysPre}{}\renewcommand{\rhcaNdaysPre}{1{,}496}  % trading days, 2014-01-02 .. 2019-12-31
\providecommand{\rhcaStartPre}{}\renewcommand{\rhcaStartPre}{2 January 2014}  % first day
\providecommand{\rhcaEndPre}{}\renewcommand{\rhcaEndPre}{31 December 2019}  % last day
\providecommand{\rhcaYearsPre}{}\renewcommand{\rhcaYearsPre}{2014--2019}  % years
\providecommand{\rhcaDMLagPre}{}\renewcommand{\rhcaDMLagPre}{7}  % Newey-West lag of the daily DM tests (dm_test automatic)
\providecommand{\rhcaHacLagBarPre}{}\renewcommand{\rhcaHacLagBarPre}{91}  % Newey-West lag of the every-bar MZ regressions, bars
\providecommand{\rhcaNdaysOos}{}\renewcommand{\rhcaNdaysOos}{1{,}082}  % trading days, 2020-01-03 .. 2024-04-30
\providecommand{\rhcaStartOos}{}\renewcommand{\rhcaStartOos}{3 January 2020}  % first day
\providecommand{\rhcaEndOos}{}\renewcommand{\rhcaEndOos}{30 April 2024}  % last day
\providecommand{\rhcaYearsOos}{}\renewcommand{\rhcaYearsOos}{2020--2024}  % years
\providecommand{\rhcaDMLagOos}{}\renewcommand{\rhcaDMLagOos}{6}  % Newey-West lag of the daily DM tests (dm_test automatic)
\providecommand{\rhcaHacLagBarOos}{}\renewcommand{\rhcaHacLagBarOos}{78}  % Newey-West lag of the every-bar MZ regressions, bars
\providecommand{\rhcaNbars}{}\renewcommand{\rhcaNbars}{33{,}527}  % scored bars of Table 1, 2014-2024
\providecommand{\rhcaCloseFirst}{}\renewcommand{\rhcaCloseFirst}{14 June 2006}  % first forecast of the closing-interval models
\providecommand{\rhcaMseExp}{}\renewcommand{\rhcaMseExp}{10}  % MSE printed x 10^this
\providecommand{\rhcaUnitExp}{}\renewcommand{\rhcaUnitExp}{6}  % MAE and the levels intercept printed x 10^this
\providecommand{\rhcaCrashStart}{}\renewcommand{\rhcaCrashStart}{24 February 2020}  % first day of the excluded window (make_rehaul_results.CRASH)
\providecommand{\rhcaCrashEnd}{}\renewcommand{\rhcaCrashEnd}{30 April 2020}  % last day of the excluded window
\providecommand{\rhcaTopDays}{}\renewcommand{\rhcaTopDays}{5}  % days in the extreme-day share
% --- closing half hour (16:00 bar): pooled forecasts, Span / Pre / Oos; Q QLIKE, Mse (x 10^10), Mae (x 10^6), Rsq R2_OOS %, Pct / Dm / P QLIKE full vs HAR, MsePct / MseDm / MseP, Gap / GapDm / GapP vs the closing-interval HAR benchmark (Oos only)
\providecommand{\rhcaQOlsHarCloseSpan}{}\renewcommand{\rhcaQOlsHarCloseSpan}{0.1269}
\providecommand{\rhcaMseOlsHarCloseSpan}{}\renewcommand{\rhcaMseOlsHarCloseSpan}{1.97}
\providecommand{\rhcaMaeOlsHarCloseSpan}{}\renewcommand{\rhcaMaeOlsHarCloseSpan}{2.92}
\providecommand{\rhcaRsqOlsHarCloseSpan}{}\renewcommand{\rhcaRsqOlsHarCloseSpan}{77.4}
\providecommand{\rhcaQRidgeHarCloseSpan}{}\renewcommand{\rhcaQRidgeHarCloseSpan}{0.1269}
\providecommand{\rhcaMseRidgeHarCloseSpan}{}\renewcommand{\rhcaMseRidgeHarCloseSpan}{1.97}
\providecommand{\rhcaMaeRidgeHarCloseSpan}{}\renewcommand{\rhcaMaeRidgeHarCloseSpan}{2.92}
\providecommand{\rhcaRsqRidgeHarCloseSpan}{}\renewcommand{\rhcaRsqRidgeHarCloseSpan}{77.4}
\providecommand{\rhcaQRidgeFullCloseSpan}{}\renewcommand{\rhcaQRidgeFullCloseSpan}{0.1291}
\providecommand{\rhcaMseRidgeFullCloseSpan}{}\renewcommand{\rhcaMseRidgeFullCloseSpan}{3.49}
\providecommand{\rhcaMaeRidgeFullCloseSpan}{}\renewcommand{\rhcaMaeRidgeFullCloseSpan}{3.22}
\providecommand{\rhcaRsqRidgeFullCloseSpan}{}\renewcommand{\rhcaRsqRidgeFullCloseSpan}{59.9}
\providecommand{\rhcaPctRidgeCloseSpan}{}\renewcommand{\rhcaPctRidgeCloseSpan}{1.8}  % QLIKE %, full vs HAR of the same class
\providecommand{\rhcaDmRidgeCloseSpan}{}\renewcommand{\rhcaDmRidgeCloseSpan}{0.67}
\providecommand{\rhcaPRidgeCloseSpan}{}\renewcommand{\rhcaPRidgeCloseSpan}{\ensuremath{p = 0.50}}
\providecommand{\rhcaMsePctRidgeCloseSpan}{}\renewcommand{\rhcaMsePctRidgeCloseSpan}{77.5}  % MSE %, full vs HAR
\providecommand{\rhcaMseDmRidgeCloseSpan}{}\renewcommand{\rhcaMseDmRidgeCloseSpan}{1.18}
\providecommand{\rhcaMsePRidgeCloseSpan}{}\renewcommand{\rhcaMsePRidgeCloseSpan}{\ensuremath{p = 0.24}}
\providecommand{\rhcaQLassoHarCloseSpan}{}\renewcommand{\rhcaQLassoHarCloseSpan}{0.1269}
\providecommand{\rhcaMseLassoHarCloseSpan}{}\renewcommand{\rhcaMseLassoHarCloseSpan}{1.97}
\providecommand{\rhcaMaeLassoHarCloseSpan}{}\renewcommand{\rhcaMaeLassoHarCloseSpan}{2.92}
\providecommand{\rhcaRsqLassoHarCloseSpan}{}\renewcommand{\rhcaRsqLassoHarCloseSpan}{77.4}
\providecommand{\rhcaQLassoFullCloseSpan}{}\renewcommand{\rhcaQLassoFullCloseSpan}{0.1322}
\providecommand{\rhcaMseLassoFullCloseSpan}{}\renewcommand{\rhcaMseLassoFullCloseSpan}{3.57}
\providecommand{\rhcaMaeLassoFullCloseSpan}{}\renewcommand{\rhcaMaeLassoFullCloseSpan}{3.24}
\providecommand{\rhcaRsqLassoFullCloseSpan}{}\renewcommand{\rhcaRsqLassoFullCloseSpan}{59.0}
\providecommand{\rhcaPctLassoCloseSpan}{}\renewcommand{\rhcaPctLassoCloseSpan}{4.2}  % QLIKE %, full vs HAR of the same class
\providecommand{\rhcaDmLassoCloseSpan}{}\renewcommand{\rhcaDmLassoCloseSpan}{1.35}
\providecommand{\rhcaPLassoCloseSpan}{}\renewcommand{\rhcaPLassoCloseSpan}{\ensuremath{p = 0.18}}
\providecommand{\rhcaMsePctLassoCloseSpan}{}\renewcommand{\rhcaMsePctLassoCloseSpan}{81.5}  % MSE %, full vs HAR
\providecommand{\rhcaMseDmLassoCloseSpan}{}\renewcommand{\rhcaMseDmLassoCloseSpan}{1.21}
\providecommand{\rhcaMsePLassoCloseSpan}{}\renewcommand{\rhcaMsePLassoCloseSpan}{\ensuremath{p = 0.23}}
\providecommand{\rhcaQEnetHarCloseSpan}{}\renewcommand{\rhcaQEnetHarCloseSpan}{0.1269}
\providecommand{\rhcaMseEnetHarCloseSpan}{}\renewcommand{\rhcaMseEnetHarCloseSpan}{1.97}
\providecommand{\rhcaMaeEnetHarCloseSpan}{}\renewcommand{\rhcaMaeEnetHarCloseSpan}{2.92}
\providecommand{\rhcaRsqEnetHarCloseSpan}{}\renewcommand{\rhcaRsqEnetHarCloseSpan}{77.4}
\providecommand{\rhcaQEnetFullCloseSpan}{}\renewcommand{\rhcaQEnetFullCloseSpan}{0.1276}
\providecommand{\rhcaMseEnetFullCloseSpan}{}\renewcommand{\rhcaMseEnetFullCloseSpan}{3.92}
\providecommand{\rhcaMaeEnetFullCloseSpan}{}\renewcommand{\rhcaMaeEnetFullCloseSpan}{3.28}
\providecommand{\rhcaRsqEnetFullCloseSpan}{}\renewcommand{\rhcaRsqEnetFullCloseSpan}{55.0}
\providecommand{\rhcaPctEnetCloseSpan}{}\renewcommand{\rhcaPctEnetCloseSpan}{0.6}  % QLIKE %, full vs HAR of the same class
\providecommand{\rhcaDmEnetCloseSpan}{}\renewcommand{\rhcaDmEnetCloseSpan}{0.24}
\providecommand{\rhcaPEnetCloseSpan}{}\renewcommand{\rhcaPEnetCloseSpan}{\ensuremath{p = 0.81}}
\providecommand{\rhcaMsePctEnetCloseSpan}{}\renewcommand{\rhcaMsePctEnetCloseSpan}{99.1}  % MSE %, full vs HAR
\providecommand{\rhcaMseDmEnetCloseSpan}{}\renewcommand{\rhcaMseDmEnetCloseSpan}{1.16}
\providecommand{\rhcaMsePEnetCloseSpan}{}\renewcommand{\rhcaMsePEnetCloseSpan}{\ensuremath{p = 0.25}}
\providecommand{\rhcaQLgbmHarCloseSpan}{}\renewcommand{\rhcaQLgbmHarCloseSpan}{0.1235}
\providecommand{\rhcaMseLgbmHarCloseSpan}{}\renewcommand{\rhcaMseLgbmHarCloseSpan}{1.77}
\providecommand{\rhcaMaeLgbmHarCloseSpan}{}\renewcommand{\rhcaMaeLgbmHarCloseSpan}{3.01}
\providecommand{\rhcaRsqLgbmHarCloseSpan}{}\renewcommand{\rhcaRsqLgbmHarCloseSpan}{79.6}
\providecommand{\rhcaQLgbmFullCloseSpan}{}\renewcommand{\rhcaQLgbmFullCloseSpan}{0.1202}
\providecommand{\rhcaMseLgbmFullCloseSpan}{}\renewcommand{\rhcaMseLgbmFullCloseSpan}{2.36}
\providecommand{\rhcaMaeLgbmFullCloseSpan}{}\renewcommand{\rhcaMaeLgbmFullCloseSpan}{3.02}
\providecommand{\rhcaRsqLgbmFullCloseSpan}{}\renewcommand{\rhcaRsqLgbmFullCloseSpan}{72.9}
\providecommand{\rhcaPctLgbmCloseSpan}{}\renewcommand{\rhcaPctLgbmCloseSpan}{\ensuremath{-}2.7}  % QLIKE %, full vs HAR of the same class
\providecommand{\rhcaDmLgbmCloseSpan}{}\renewcommand{\rhcaDmLgbmCloseSpan}{\ensuremath{-}1.76}
\providecommand{\rhcaPLgbmCloseSpan}{}\renewcommand{\rhcaPLgbmCloseSpan}{\ensuremath{p = 0.079}}
\providecommand{\rhcaMsePctLgbmCloseSpan}{}\renewcommand{\rhcaMsePctLgbmCloseSpan}{33.0}  % MSE %, full vs HAR
\providecommand{\rhcaMseDmLgbmCloseSpan}{}\renewcommand{\rhcaMseDmLgbmCloseSpan}{1.10}
\providecommand{\rhcaMsePLgbmCloseSpan}{}\renewcommand{\rhcaMsePLgbmCloseSpan}{\ensuremath{p = 0.27}}
\providecommand{\rhcaQXgbHarCloseSpan}{}\renewcommand{\rhcaQXgbHarCloseSpan}{0.1234}
\providecommand{\rhcaMseXgbHarCloseSpan}{}\renewcommand{\rhcaMseXgbHarCloseSpan}{1.73}
\providecommand{\rhcaMaeXgbHarCloseSpan}{}\renewcommand{\rhcaMaeXgbHarCloseSpan}{2.95}
\providecommand{\rhcaRsqXgbHarCloseSpan}{}\renewcommand{\rhcaRsqXgbHarCloseSpan}{80.1}
\providecommand{\rhcaQXgbFullCloseSpan}{}\renewcommand{\rhcaQXgbFullCloseSpan}{0.1195}
\providecommand{\rhcaMseXgbFullCloseSpan}{}\renewcommand{\rhcaMseXgbFullCloseSpan}{2.19}
\providecommand{\rhcaMaeXgbFullCloseSpan}{}\renewcommand{\rhcaMaeXgbFullCloseSpan}{3.01}
\providecommand{\rhcaRsqXgbFullCloseSpan}{}\renewcommand{\rhcaRsqXgbFullCloseSpan}{74.9}
\providecommand{\rhcaPctXgbCloseSpan}{}\renewcommand{\rhcaPctXgbCloseSpan}{\ensuremath{-}3.1}  % QLIKE %, full vs HAR of the same class
\providecommand{\rhcaDmXgbCloseSpan}{}\renewcommand{\rhcaDmXgbCloseSpan}{\ensuremath{-}2.01}
\providecommand{\rhcaPXgbCloseSpan}{}\renewcommand{\rhcaPXgbCloseSpan}{\ensuremath{p = 0.044}}
\providecommand{\rhcaMsePctXgbCloseSpan}{}\renewcommand{\rhcaMsePctXgbCloseSpan}{26.3}  % MSE %, full vs HAR
\providecommand{\rhcaMseDmXgbCloseSpan}{}\renewcommand{\rhcaMseDmXgbCloseSpan}{1.09}
\providecommand{\rhcaMsePXgbCloseSpan}{}\renewcommand{\rhcaMsePXgbCloseSpan}{\ensuremath{p = 0.28}}
\providecommand{\rhcaQRfHarCloseSpan}{}\renewcommand{\rhcaQRfHarCloseSpan}{0.1266}
\providecommand{\rhcaMseRfHarCloseSpan}{}\renewcommand{\rhcaMseRfHarCloseSpan}{1.63}
\providecommand{\rhcaMaeRfHarCloseSpan}{}\renewcommand{\rhcaMaeRfHarCloseSpan}{2.91}
\providecommand{\rhcaRsqRfHarCloseSpan}{}\renewcommand{\rhcaRsqRfHarCloseSpan}{81.3}
\providecommand{\rhcaQRfFullCloseSpan}{}\renewcommand{\rhcaQRfFullCloseSpan}{0.1246}
\providecommand{\rhcaMseRfFullCloseSpan}{}\renewcommand{\rhcaMseRfFullCloseSpan}{1.95}
\providecommand{\rhcaMaeRfFullCloseSpan}{}\renewcommand{\rhcaMaeRfFullCloseSpan}{2.97}
\providecommand{\rhcaRsqRfFullCloseSpan}{}\renewcommand{\rhcaRsqRfFullCloseSpan}{77.6}
\providecommand{\rhcaPctRfCloseSpan}{}\renewcommand{\rhcaPctRfCloseSpan}{\ensuremath{-}1.6}  % QLIKE %, full vs HAR of the same class
\providecommand{\rhcaDmRfCloseSpan}{}\renewcommand{\rhcaDmRfCloseSpan}{\ensuremath{-}1.03}
\providecommand{\rhcaPRfCloseSpan}{}\renewcommand{\rhcaPRfCloseSpan}{\ensuremath{p = 0.30}}
\providecommand{\rhcaMsePctRfCloseSpan}{}\renewcommand{\rhcaMsePctRfCloseSpan}{19.7}  % MSE %, full vs HAR
\providecommand{\rhcaMseDmRfCloseSpan}{}\renewcommand{\rhcaMseDmRfCloseSpan}{1.11}
\providecommand{\rhcaMsePRfCloseSpan}{}\renewcommand{\rhcaMsePRfCloseSpan}{\ensuremath{p = 0.27}}
\providecommand{\rhcaQOlsHarClosePre}{}\renewcommand{\rhcaQOlsHarClosePre}{0.1365}
\providecommand{\rhcaMseOlsHarClosePre}{}\renewcommand{\rhcaMseOlsHarClosePre}{0.32}
\providecommand{\rhcaMaeOlsHarClosePre}{}\renewcommand{\rhcaMaeOlsHarClosePre}{1.74}
\providecommand{\rhcaRsqOlsHarClosePre}{}\renewcommand{\rhcaRsqOlsHarClosePre}{73.9}
\providecommand{\rhcaQRidgeHarClosePre}{}\renewcommand{\rhcaQRidgeHarClosePre}{0.1365}
\providecommand{\rhcaMseRidgeHarClosePre}{}\renewcommand{\rhcaMseRidgeHarClosePre}{0.32}
\providecommand{\rhcaMaeRidgeHarClosePre}{}\renewcommand{\rhcaMaeRidgeHarClosePre}{1.74}
\providecommand{\rhcaRsqRidgeHarClosePre}{}\renewcommand{\rhcaRsqRidgeHarClosePre}{73.9}
\providecommand{\rhcaQRidgeFullClosePre}{}\renewcommand{\rhcaQRidgeFullClosePre}{0.1309}
\providecommand{\rhcaMseRidgeFullClosePre}{}\renewcommand{\rhcaMseRidgeFullClosePre}{0.28}
\providecommand{\rhcaMaeRidgeFullClosePre}{}\renewcommand{\rhcaMaeRidgeFullClosePre}{1.66}
\providecommand{\rhcaRsqRidgeFullClosePre}{}\renewcommand{\rhcaRsqRidgeFullClosePre}{77.3}
\providecommand{\rhcaPctRidgeClosePre}{}\renewcommand{\rhcaPctRidgeClosePre}{\ensuremath{-}4.1}  % QLIKE %, full vs HAR of the same class
\providecommand{\rhcaDmRidgeClosePre}{}\renewcommand{\rhcaDmRidgeClosePre}{\ensuremath{-}1.42}
\providecommand{\rhcaPRidgeClosePre}{}\renewcommand{\rhcaPRidgeClosePre}{\ensuremath{p = 0.16}}
\providecommand{\rhcaMsePctRidgeClosePre}{}\renewcommand{\rhcaMsePctRidgeClosePre}{\ensuremath{-}13.1}  % MSE %, full vs HAR
\providecommand{\rhcaMseDmRidgeClosePre}{}\renewcommand{\rhcaMseDmRidgeClosePre}{\ensuremath{-}1.79}
\providecommand{\rhcaMsePRidgeClosePre}{}\renewcommand{\rhcaMsePRidgeClosePre}{\ensuremath{p = 0.074}}
\providecommand{\rhcaQLassoHarClosePre}{}\renewcommand{\rhcaQLassoHarClosePre}{0.1365}
\providecommand{\rhcaMseLassoHarClosePre}{}\renewcommand{\rhcaMseLassoHarClosePre}{0.32}
\providecommand{\rhcaMaeLassoHarClosePre}{}\renewcommand{\rhcaMaeLassoHarClosePre}{1.74}
\providecommand{\rhcaRsqLassoHarClosePre}{}\renewcommand{\rhcaRsqLassoHarClosePre}{73.9}
\providecommand{\rhcaQLassoFullClosePre}{}\renewcommand{\rhcaQLassoFullClosePre}{0.1333}
\providecommand{\rhcaMseLassoFullClosePre}{}\renewcommand{\rhcaMseLassoFullClosePre}{0.27}
\providecommand{\rhcaMaeLassoFullClosePre}{}\renewcommand{\rhcaMaeLassoFullClosePre}{1.65}
\providecommand{\rhcaRsqLassoFullClosePre}{}\renewcommand{\rhcaRsqLassoFullClosePre}{78.0}
\providecommand{\rhcaPctLassoClosePre}{}\renewcommand{\rhcaPctLassoClosePre}{\ensuremath{-}2.3}  % QLIKE %, full vs HAR of the same class
\providecommand{\rhcaDmLassoClosePre}{}\renewcommand{\rhcaDmLassoClosePre}{\ensuremath{-}0.67}
\providecommand{\rhcaPLassoClosePre}{}\renewcommand{\rhcaPLassoClosePre}{\ensuremath{p = 0.50}}
\providecommand{\rhcaMsePctLassoClosePre}{}\renewcommand{\rhcaMsePctLassoClosePre}{\ensuremath{-}15.8}  % MSE %, full vs HAR
\providecommand{\rhcaMseDmLassoClosePre}{}\renewcommand{\rhcaMseDmLassoClosePre}{\ensuremath{-}1.78}
\providecommand{\rhcaMsePLassoClosePre}{}\renewcommand{\rhcaMsePLassoClosePre}{\ensuremath{p = 0.075}}
\providecommand{\rhcaQEnetHarClosePre}{}\renewcommand{\rhcaQEnetHarClosePre}{0.1365}
\providecommand{\rhcaMseEnetHarClosePre}{}\renewcommand{\rhcaMseEnetHarClosePre}{0.32}
\providecommand{\rhcaMaeEnetHarClosePre}{}\renewcommand{\rhcaMaeEnetHarClosePre}{1.74}
\providecommand{\rhcaRsqEnetHarClosePre}{}\renewcommand{\rhcaRsqEnetHarClosePre}{73.9}
\providecommand{\rhcaQEnetFullClosePre}{}\renewcommand{\rhcaQEnetFullClosePre}{0.1300}
\providecommand{\rhcaMseEnetFullClosePre}{}\renewcommand{\rhcaMseEnetFullClosePre}{0.28}
\providecommand{\rhcaMaeEnetFullClosePre}{}\renewcommand{\rhcaMaeEnetFullClosePre}{1.64}
\providecommand{\rhcaRsqEnetFullClosePre}{}\renewcommand{\rhcaRsqEnetFullClosePre}{77.6}
\providecommand{\rhcaPctEnetClosePre}{}\renewcommand{\rhcaPctEnetClosePre}{\ensuremath{-}4.8}  % QLIKE %, full vs HAR of the same class
\providecommand{\rhcaDmEnetClosePre}{}\renewcommand{\rhcaDmEnetClosePre}{\ensuremath{-}1.80}
\providecommand{\rhcaPEnetClosePre}{}\renewcommand{\rhcaPEnetClosePre}{\ensuremath{p = 0.071}}
\providecommand{\rhcaMsePctEnetClosePre}{}\renewcommand{\rhcaMsePctEnetClosePre}{\ensuremath{-}14.4}  % MSE %, full vs HAR
\providecommand{\rhcaMseDmEnetClosePre}{}\renewcommand{\rhcaMseDmEnetClosePre}{\ensuremath{-}1.90}
\providecommand{\rhcaMsePEnetClosePre}{}\renewcommand{\rhcaMsePEnetClosePre}{\ensuremath{p = 0.058}}
\providecommand{\rhcaQLgbmHarClosePre}{}\renewcommand{\rhcaQLgbmHarClosePre}{0.1333}
\providecommand{\rhcaMseLgbmHarClosePre}{}\renewcommand{\rhcaMseLgbmHarClosePre}{0.40}
\providecommand{\rhcaMaeLgbmHarClosePre}{}\renewcommand{\rhcaMaeLgbmHarClosePre}{1.84}
\providecommand{\rhcaRsqLgbmHarClosePre}{}\renewcommand{\rhcaRsqLgbmHarClosePre}{68.1}
\providecommand{\rhcaQLgbmFullClosePre}{}\renewcommand{\rhcaQLgbmFullClosePre}{0.1287}
\providecommand{\rhcaMseLgbmFullClosePre}{}\renewcommand{\rhcaMseLgbmFullClosePre}{0.37}
\providecommand{\rhcaMaeLgbmFullClosePre}{}\renewcommand{\rhcaMaeLgbmFullClosePre}{1.74}
\providecommand{\rhcaRsqLgbmFullClosePre}{}\renewcommand{\rhcaRsqLgbmFullClosePre}{70.5}
\providecommand{\rhcaPctLgbmClosePre}{}\renewcommand{\rhcaPctLgbmClosePre}{\ensuremath{-}3.4}  % QLIKE %, full vs HAR of the same class
\providecommand{\rhcaDmLgbmClosePre}{}\renewcommand{\rhcaDmLgbmClosePre}{\ensuremath{-}1.99}
\providecommand{\rhcaPLgbmClosePre}{}\renewcommand{\rhcaPLgbmClosePre}{\ensuremath{p = 0.046}}
\providecommand{\rhcaMsePctLgbmClosePre}{}\renewcommand{\rhcaMsePctLgbmClosePre}{\ensuremath{-}7.5}  % MSE %, full vs HAR
\providecommand{\rhcaMseDmLgbmClosePre}{}\renewcommand{\rhcaMseDmLgbmClosePre}{\ensuremath{-}1.78}
\providecommand{\rhcaMsePLgbmClosePre}{}\renewcommand{\rhcaMsePLgbmClosePre}{\ensuremath{p = 0.076}}
\providecommand{\rhcaQXgbHarClosePre}{}\renewcommand{\rhcaQXgbHarClosePre}{0.1331}
\providecommand{\rhcaMseXgbHarClosePre}{}\renewcommand{\rhcaMseXgbHarClosePre}{0.37}
\providecommand{\rhcaMaeXgbHarClosePre}{}\renewcommand{\rhcaMaeXgbHarClosePre}{1.82}
\providecommand{\rhcaRsqXgbHarClosePre}{}\renewcommand{\rhcaRsqXgbHarClosePre}{70.0}
\providecommand{\rhcaQXgbFullClosePre}{}\renewcommand{\rhcaQXgbFullClosePre}{0.1277}
\providecommand{\rhcaMseXgbFullClosePre}{}\renewcommand{\rhcaMseXgbFullClosePre}{0.33}
\providecommand{\rhcaMaeXgbFullClosePre}{}\renewcommand{\rhcaMaeXgbFullClosePre}{1.74}
\providecommand{\rhcaRsqXgbFullClosePre}{}\renewcommand{\rhcaRsqXgbFullClosePre}{73.0}
\providecommand{\rhcaPctXgbClosePre}{}\renewcommand{\rhcaPctXgbClosePre}{\ensuremath{-}4.0}  % QLIKE %, full vs HAR of the same class
\providecommand{\rhcaDmXgbClosePre}{}\renewcommand{\rhcaDmXgbClosePre}{\ensuremath{-}2.13}
\providecommand{\rhcaPXgbClosePre}{}\renewcommand{\rhcaPXgbClosePre}{\ensuremath{p = 0.033}}
\providecommand{\rhcaMsePctXgbClosePre}{}\renewcommand{\rhcaMsePctXgbClosePre}{\ensuremath{-}10.0}  % MSE %, full vs HAR
\providecommand{\rhcaMseDmXgbClosePre}{}\renewcommand{\rhcaMseDmXgbClosePre}{\ensuremath{-}1.66}
\providecommand{\rhcaMsePXgbClosePre}{}\renewcommand{\rhcaMsePXgbClosePre}{\ensuremath{p = 0.097}}
\providecommand{\rhcaQRfHarClosePre}{}\renewcommand{\rhcaQRfHarClosePre}{0.1364}
\providecommand{\rhcaMseRfHarClosePre}{}\renewcommand{\rhcaMseRfHarClosePre}{0.36}
\providecommand{\rhcaMaeRfHarClosePre}{}\renewcommand{\rhcaMaeRfHarClosePre}{1.81}
\providecommand{\rhcaRsqRfHarClosePre}{}\renewcommand{\rhcaRsqRfHarClosePre}{70.7}
\providecommand{\rhcaQRfFullClosePre}{}\renewcommand{\rhcaQRfFullClosePre}{0.1335}
\providecommand{\rhcaMseRfFullClosePre}{}\renewcommand{\rhcaMseRfFullClosePre}{0.34}
\providecommand{\rhcaMaeRfFullClosePre}{}\renewcommand{\rhcaMaeRfFullClosePre}{1.77}
\providecommand{\rhcaRsqRfFullClosePre}{}\renewcommand{\rhcaRsqRfFullClosePre}{72.3}
\providecommand{\rhcaPctRfClosePre}{}\renewcommand{\rhcaPctRfClosePre}{\ensuremath{-}2.1}  % QLIKE %, full vs HAR of the same class
\providecommand{\rhcaDmRfClosePre}{}\renewcommand{\rhcaDmRfClosePre}{\ensuremath{-}1.14}
\providecommand{\rhcaPRfClosePre}{}\renewcommand{\rhcaPRfClosePre}{\ensuremath{p = 0.25}}
\providecommand{\rhcaMsePctRfClosePre}{}\renewcommand{\rhcaMsePctRfClosePre}{\ensuremath{-}5.3}  % MSE %, full vs HAR
\providecommand{\rhcaMseDmRfClosePre}{}\renewcommand{\rhcaMseDmRfClosePre}{\ensuremath{-}1.39}
\providecommand{\rhcaMsePRfClosePre}{}\renewcommand{\rhcaMsePRfClosePre}{\ensuremath{p = 0.17}}
\providecommand{\rhcaQOlsHarCloseOos}{}\renewcommand{\rhcaQOlsHarCloseOos}{0.1137}
\providecommand{\rhcaMseOlsHarCloseOos}{}\renewcommand{\rhcaMseOlsHarCloseOos}{4.24}
\providecommand{\rhcaMaeOlsHarCloseOos}{}\renewcommand{\rhcaMaeOlsHarCloseOos}{4.55}
\providecommand{\rhcaRsqOlsHarCloseOos}{}\renewcommand{\rhcaRsqOlsHarCloseOos}{77.7}
\providecommand{\rhcaGapOlsHarCloseOos}{}\renewcommand{\rhcaGapOlsHarCloseOos}{10.9}  % QLIKE % vs the closing-interval HAR benchmark
\providecommand{\rhcaGapDmOlsHarCloseOos}{}\renewcommand{\rhcaGapDmOlsHarCloseOos}{2.84}
\providecommand{\rhcaGapPOlsHarCloseOos}{}\renewcommand{\rhcaGapPOlsHarCloseOos}{\ensuremath{p = 0.004}}
\providecommand{\rhcaQRidgeHarCloseOos}{}\renewcommand{\rhcaQRidgeHarCloseOos}{0.1137}
\providecommand{\rhcaMseRidgeHarCloseOos}{}\renewcommand{\rhcaMseRidgeHarCloseOos}{4.24}
\providecommand{\rhcaMaeRidgeHarCloseOos}{}\renewcommand{\rhcaMaeRidgeHarCloseOos}{4.55}
\providecommand{\rhcaRsqRidgeHarCloseOos}{}\renewcommand{\rhcaRsqRidgeHarCloseOos}{77.7}
\providecommand{\rhcaGapRidgeHarCloseOos}{}\renewcommand{\rhcaGapRidgeHarCloseOos}{10.9}  % QLIKE % vs the closing-interval HAR benchmark
\providecommand{\rhcaGapDmRidgeHarCloseOos}{}\renewcommand{\rhcaGapDmRidgeHarCloseOos}{2.84}
\providecommand{\rhcaGapPRidgeHarCloseOos}{}\renewcommand{\rhcaGapPRidgeHarCloseOos}{\ensuremath{p = 0.004}}
\providecommand{\rhcaQRidgeFullCloseOos}{}\renewcommand{\rhcaQRidgeFullCloseOos}{0.1269}
\providecommand{\rhcaMseRidgeFullCloseOos}{}\renewcommand{\rhcaMseRidgeFullCloseOos}{7.94}
\providecommand{\rhcaMaeRidgeFullCloseOos}{}\renewcommand{\rhcaMaeRidgeFullCloseOos}{5.39}
\providecommand{\rhcaRsqRidgeFullCloseOos}{}\renewcommand{\rhcaRsqRidgeFullCloseOos}{58.4}
\providecommand{\rhcaPctRidgeCloseOos}{}\renewcommand{\rhcaPctRidgeCloseOos}{11.6}  % QLIKE %, full vs HAR of the same class
\providecommand{\rhcaDmRidgeCloseOos}{}\renewcommand{\rhcaDmRidgeCloseOos}{2.44}
\providecommand{\rhcaPRidgeCloseOos}{}\renewcommand{\rhcaPRidgeCloseOos}{\ensuremath{p = 0.015}}
\providecommand{\rhcaMsePctRidgeCloseOos}{}\renewcommand{\rhcaMsePctRidgeCloseOos}{87.0}  % MSE %, full vs HAR
\providecommand{\rhcaMseDmRidgeCloseOos}{}\renewcommand{\rhcaMseDmRidgeCloseOos}{1.29}
\providecommand{\rhcaMsePRidgeCloseOos}{}\renewcommand{\rhcaMsePRidgeCloseOos}{\ensuremath{p = 0.20}}
\providecommand{\rhcaGapRidgeFullCloseOos}{}\renewcommand{\rhcaGapRidgeFullCloseOos}{23.8}  % QLIKE % vs the closing-interval HAR benchmark
\providecommand{\rhcaGapDmRidgeFullCloseOos}{}\renewcommand{\rhcaGapDmRidgeFullCloseOos}{3.72}
\providecommand{\rhcaGapPRidgeFullCloseOos}{}\renewcommand{\rhcaGapPRidgeFullCloseOos}{\ensuremath{p < 0.001}}
\providecommand{\rhcaQLassoHarCloseOos}{}\renewcommand{\rhcaQLassoHarCloseOos}{0.1137}
\providecommand{\rhcaMseLassoHarCloseOos}{}\renewcommand{\rhcaMseLassoHarCloseOos}{4.24}
\providecommand{\rhcaMaeLassoHarCloseOos}{}\renewcommand{\rhcaMaeLassoHarCloseOos}{4.55}
\providecommand{\rhcaRsqLassoHarCloseOos}{}\renewcommand{\rhcaRsqLassoHarCloseOos}{77.7}
\providecommand{\rhcaGapLassoHarCloseOos}{}\renewcommand{\rhcaGapLassoHarCloseOos}{10.9}  % QLIKE % vs the closing-interval HAR benchmark
\providecommand{\rhcaGapDmLassoHarCloseOos}{}\renewcommand{\rhcaGapDmLassoHarCloseOos}{2.84}
\providecommand{\rhcaGapPLassoHarCloseOos}{}\renewcommand{\rhcaGapPLassoHarCloseOos}{\ensuremath{p = 0.004}}
\providecommand{\rhcaQLassoFullCloseOos}{}\renewcommand{\rhcaQLassoFullCloseOos}{0.1309}
\providecommand{\rhcaMseLassoFullCloseOos}{}\renewcommand{\rhcaMseLassoFullCloseOos}{8.14}
\providecommand{\rhcaMaeLassoFullCloseOos}{}\renewcommand{\rhcaMaeLassoFullCloseOos}{5.45}
\providecommand{\rhcaRsqLassoFullCloseOos}{}\renewcommand{\rhcaRsqLassoFullCloseOos}{57.3}
\providecommand{\rhcaPctLassoCloseOos}{}\renewcommand{\rhcaPctLassoCloseOos}{15.2}  % QLIKE %, full vs HAR of the same class
\providecommand{\rhcaDmLassoCloseOos}{}\renewcommand{\rhcaDmLassoCloseOos}{2.68}
\providecommand{\rhcaPLassoCloseOos}{}\renewcommand{\rhcaPLassoCloseOos}{\ensuremath{p = 0.007}}
\providecommand{\rhcaMsePctLassoCloseOos}{}\renewcommand{\rhcaMsePctLassoCloseOos}{91.7}  % MSE %, full vs HAR
\providecommand{\rhcaMseDmLassoCloseOos}{}\renewcommand{\rhcaMseDmLassoCloseOos}{1.33}
\providecommand{\rhcaMsePLassoCloseOos}{}\renewcommand{\rhcaMsePLassoCloseOos}{\ensuremath{p = 0.18}}
\providecommand{\rhcaGapLassoFullCloseOos}{}\renewcommand{\rhcaGapLassoFullCloseOos}{27.7}  % QLIKE % vs the closing-interval HAR benchmark
\providecommand{\rhcaGapDmLassoFullCloseOos}{}\renewcommand{\rhcaGapDmLassoFullCloseOos}{3.80}
\providecommand{\rhcaGapPLassoFullCloseOos}{}\renewcommand{\rhcaGapPLassoFullCloseOos}{\ensuremath{p < 0.001}}
\providecommand{\rhcaQEnetHarCloseOos}{}\renewcommand{\rhcaQEnetHarCloseOos}{0.1137}
\providecommand{\rhcaMseEnetHarCloseOos}{}\renewcommand{\rhcaMseEnetHarCloseOos}{4.24}
\providecommand{\rhcaMaeEnetHarCloseOos}{}\renewcommand{\rhcaMaeEnetHarCloseOos}{4.55}
\providecommand{\rhcaRsqEnetHarCloseOos}{}\renewcommand{\rhcaRsqEnetHarCloseOos}{77.7}
\providecommand{\rhcaGapEnetHarCloseOos}{}\renewcommand{\rhcaGapEnetHarCloseOos}{10.9}  % QLIKE % vs the closing-interval HAR benchmark
\providecommand{\rhcaGapDmEnetHarCloseOos}{}\renewcommand{\rhcaGapDmEnetHarCloseOos}{2.84}
\providecommand{\rhcaGapPEnetHarCloseOos}{}\renewcommand{\rhcaGapPEnetHarCloseOos}{\ensuremath{p = 0.004}}
\providecommand{\rhcaQEnetFullCloseOos}{}\renewcommand{\rhcaQEnetFullCloseOos}{0.1245}
\providecommand{\rhcaMseEnetFullCloseOos}{}\renewcommand{\rhcaMseEnetFullCloseOos}{8.96}
\providecommand{\rhcaMaeEnetFullCloseOos}{}\renewcommand{\rhcaMaeEnetFullCloseOos}{5.54}
\providecommand{\rhcaRsqEnetFullCloseOos}{}\renewcommand{\rhcaRsqEnetFullCloseOos}{53.0}
\providecommand{\rhcaPctEnetCloseOos}{}\renewcommand{\rhcaPctEnetCloseOos}{9.5}  % QLIKE %, full vs HAR of the same class
\providecommand{\rhcaDmEnetCloseOos}{}\renewcommand{\rhcaDmEnetCloseOos}{2.09}
\providecommand{\rhcaPEnetCloseOos}{}\renewcommand{\rhcaPEnetCloseOos}{\ensuremath{p = 0.036}}
\providecommand{\rhcaMsePctEnetCloseOos}{}\renewcommand{\rhcaMsePctEnetCloseOos}{111.1}  % MSE %, full vs HAR
\providecommand{\rhcaMseDmEnetCloseOos}{}\renewcommand{\rhcaMseDmEnetCloseOos}{1.25}
\providecommand{\rhcaMsePEnetCloseOos}{}\renewcommand{\rhcaMsePEnetCloseOos}{\ensuremath{p = 0.21}}
\providecommand{\rhcaGapEnetFullCloseOos}{}\renewcommand{\rhcaGapEnetFullCloseOos}{21.4}  % QLIKE % vs the closing-interval HAR benchmark
\providecommand{\rhcaGapDmEnetFullCloseOos}{}\renewcommand{\rhcaGapDmEnetFullCloseOos}{3.48}
\providecommand{\rhcaGapPEnetFullCloseOos}{}\renewcommand{\rhcaGapPEnetFullCloseOos}{\ensuremath{p < 0.001}}
\providecommand{\rhcaQLgbmHarCloseOos}{}\renewcommand{\rhcaQLgbmHarCloseOos}{0.1101}
\providecommand{\rhcaMseLgbmHarCloseOos}{}\renewcommand{\rhcaMseLgbmHarCloseOos}{3.68}
\providecommand{\rhcaMaeLgbmHarCloseOos}{}\renewcommand{\rhcaMaeLgbmHarCloseOos}{4.63}
\providecommand{\rhcaRsqLgbmHarCloseOos}{}\renewcommand{\rhcaRsqLgbmHarCloseOos}{80.7}
\providecommand{\rhcaGapLgbmHarCloseOos}{}\renewcommand{\rhcaGapLgbmHarCloseOos}{7.4}  % QLIKE % vs the closing-interval HAR benchmark
\providecommand{\rhcaGapDmLgbmHarCloseOos}{}\renewcommand{\rhcaGapDmLgbmHarCloseOos}{2.60}
\providecommand{\rhcaGapPLgbmHarCloseOos}{}\renewcommand{\rhcaGapPLgbmHarCloseOos}{\ensuremath{p = 0.009}}
\providecommand{\rhcaQLgbmFullCloseOos}{}\renewcommand{\rhcaQLgbmFullCloseOos}{0.1086}
\providecommand{\rhcaMseLgbmFullCloseOos}{}\renewcommand{\rhcaMseLgbmFullCloseOos}{5.12}
\providecommand{\rhcaMaeLgbmFullCloseOos}{}\renewcommand{\rhcaMaeLgbmFullCloseOos}{4.80}
\providecommand{\rhcaRsqLgbmFullCloseOos}{}\renewcommand{\rhcaRsqLgbmFullCloseOos}{73.1}
\providecommand{\rhcaPctLgbmCloseOos}{}\renewcommand{\rhcaPctLgbmCloseOos}{\ensuremath{-}1.4}  % QLIKE %, full vs HAR of the same class
\providecommand{\rhcaDmLgbmCloseOos}{}\renewcommand{\rhcaDmLgbmCloseOos}{\ensuremath{-}0.49}
\providecommand{\rhcaPLgbmCloseOos}{}\renewcommand{\rhcaPLgbmCloseOos}{\ensuremath{p = 0.62}}
\providecommand{\rhcaMsePctLgbmCloseOos}{}\renewcommand{\rhcaMsePctLgbmCloseOos}{38.9}  % MSE %, full vs HAR
\providecommand{\rhcaMseDmLgbmCloseOos}{}\renewcommand{\rhcaMseDmLgbmCloseOos}{1.22}
\providecommand{\rhcaMsePLgbmCloseOos}{}\renewcommand{\rhcaMsePLgbmCloseOos}{\ensuremath{p = 0.22}}
\providecommand{\rhcaGapLgbmFullCloseOos}{}\renewcommand{\rhcaGapLgbmFullCloseOos}{6.0}  % QLIKE % vs the closing-interval HAR benchmark
\providecommand{\rhcaGapDmLgbmFullCloseOos}{}\renewcommand{\rhcaGapDmLgbmFullCloseOos}{1.41}
\providecommand{\rhcaGapPLgbmFullCloseOos}{}\renewcommand{\rhcaGapPLgbmFullCloseOos}{\ensuremath{p = 0.16}}
\providecommand{\rhcaQXgbHarCloseOos}{}\renewcommand{\rhcaQXgbHarCloseOos}{0.1102}
\providecommand{\rhcaMseXgbHarCloseOos}{}\renewcommand{\rhcaMseXgbHarCloseOos}{3.62}
\providecommand{\rhcaMaeXgbHarCloseOos}{}\renewcommand{\rhcaMaeXgbHarCloseOos}{4.51}
\providecommand{\rhcaRsqXgbHarCloseOos}{}\renewcommand{\rhcaRsqXgbHarCloseOos}{81.0}
\providecommand{\rhcaGapXgbHarCloseOos}{}\renewcommand{\rhcaGapXgbHarCloseOos}{7.5}  % QLIKE % vs the closing-interval HAR benchmark
\providecommand{\rhcaGapDmXgbHarCloseOos}{}\renewcommand{\rhcaGapDmXgbHarCloseOos}{2.70}
\providecommand{\rhcaGapPXgbHarCloseOos}{}\renewcommand{\rhcaGapPXgbHarCloseOos}{\ensuremath{p = 0.007}}
\providecommand{\rhcaQXgbFullCloseOos}{}\renewcommand{\rhcaQXgbFullCloseOos}{0.1083}
\providecommand{\rhcaMseXgbFullCloseOos}{}\renewcommand{\rhcaMseXgbFullCloseOos}{4.76}
\providecommand{\rhcaMaeXgbFullCloseOos}{}\renewcommand{\rhcaMaeXgbFullCloseOos}{4.76}
\providecommand{\rhcaRsqXgbFullCloseOos}{}\renewcommand{\rhcaRsqXgbFullCloseOos}{75.0}
\providecommand{\rhcaPctXgbCloseOos}{}\renewcommand{\rhcaPctXgbCloseOos}{\ensuremath{-}1.7}  % QLIKE %, full vs HAR of the same class
\providecommand{\rhcaDmXgbCloseOos}{}\renewcommand{\rhcaDmXgbCloseOos}{\ensuremath{-}0.63}
\providecommand{\rhcaPXgbCloseOos}{}\renewcommand{\rhcaPXgbCloseOos}{\ensuremath{p = 0.53}}
\providecommand{\rhcaMsePctXgbCloseOos}{}\renewcommand{\rhcaMsePctXgbCloseOos}{31.4}  % MSE %, full vs HAR
\providecommand{\rhcaMseDmXgbCloseOos}{}\renewcommand{\rhcaMseDmXgbCloseOos}{1.23}
\providecommand{\rhcaMsePXgbCloseOos}{}\renewcommand{\rhcaMsePXgbCloseOos}{\ensuremath{p = 0.22}}
\providecommand{\rhcaGapXgbFullCloseOos}{}\renewcommand{\rhcaGapXgbFullCloseOos}{5.7}  % QLIKE % vs the closing-interval HAR benchmark
\providecommand{\rhcaGapDmXgbFullCloseOos}{}\renewcommand{\rhcaGapDmXgbFullCloseOos}{1.40}
\providecommand{\rhcaGapPXgbFullCloseOos}{}\renewcommand{\rhcaGapPXgbFullCloseOos}{\ensuremath{p = 0.16}}
\providecommand{\rhcaQRfHarCloseOos}{}\renewcommand{\rhcaQRfHarCloseOos}{0.1131}
\providecommand{\rhcaMseRfHarCloseOos}{}\renewcommand{\rhcaMseRfHarCloseOos}{3.39}
\providecommand{\rhcaMaeRfHarCloseOos}{}\renewcommand{\rhcaMaeRfHarCloseOos}{4.42}
\providecommand{\rhcaRsqRfHarCloseOos}{}\renewcommand{\rhcaRsqRfHarCloseOos}{82.2}
\providecommand{\rhcaGapRfHarCloseOos}{}\renewcommand{\rhcaGapRfHarCloseOos}{10.4}  % QLIKE % vs the closing-interval HAR benchmark
\providecommand{\rhcaGapDmRfHarCloseOos}{}\renewcommand{\rhcaGapDmRfHarCloseOos}{3.12}
\providecommand{\rhcaGapPRfHarCloseOos}{}\renewcommand{\rhcaGapPRfHarCloseOos}{\ensuremath{p = 0.002}}
\providecommand{\rhcaQRfFullCloseOos}{}\renewcommand{\rhcaQRfFullCloseOos}{0.1123}
\providecommand{\rhcaMseRfFullCloseOos}{}\renewcommand{\rhcaMseRfFullCloseOos}{4.18}
\providecommand{\rhcaMaeRfFullCloseOos}{}\renewcommand{\rhcaMaeRfFullCloseOos}{4.63}
\providecommand{\rhcaRsqRfFullCloseOos}{}\renewcommand{\rhcaRsqRfFullCloseOos}{78.0}
\providecommand{\rhcaPctRfCloseOos}{}\renewcommand{\rhcaPctRfCloseOos}{\ensuremath{-}0.7}  % QLIKE %, full vs HAR of the same class
\providecommand{\rhcaDmRfCloseOos}{}\renewcommand{\rhcaDmRfCloseOos}{\ensuremath{-}0.27}
\providecommand{\rhcaPRfCloseOos}{}\renewcommand{\rhcaPRfCloseOos}{\ensuremath{p = 0.79}}
\providecommand{\rhcaMsePctRfCloseOos}{}\renewcommand{\rhcaMsePctRfCloseOos}{23.3}  % MSE %, full vs HAR
\providecommand{\rhcaMseDmRfCloseOos}{}\renewcommand{\rhcaMseDmRfCloseOos}{1.25}
\providecommand{\rhcaMsePRfCloseOos}{}\renewcommand{\rhcaMsePRfCloseOos}{\ensuremath{p = 0.21}}
\providecommand{\rhcaGapRfFullCloseOos}{}\renewcommand{\rhcaGapRfFullCloseOos}{9.6}  % QLIKE % vs the closing-interval HAR benchmark
\providecommand{\rhcaGapDmRfFullCloseOos}{}\renewcommand{\rhcaGapDmRfFullCloseOos}{2.13}
\providecommand{\rhcaGapPRfFullCloseOos}{}\renewcommand{\rhcaGapPRfFullCloseOos}{\ensuremath{p = 0.033}}
\providecommand{\rhcaQBmCloseOos}{}\renewcommand{\rhcaQBmCloseOos}{0.1025}
\providecommand{\rhcaMseBmCloseOos}{}\renewcommand{\rhcaMseBmCloseOos}{4.31}
\providecommand{\rhcaMaeBmCloseOos}{}\renewcommand{\rhcaMaeBmCloseOos}{4.78}
\providecommand{\rhcaRsqBmCloseOos}{}\renewcommand{\rhcaRsqBmCloseOos}{77.4}
\providecommand{\rhcaQCiRidgeHarCloseOos}{}\renewcommand{\rhcaQCiRidgeHarCloseOos}{0.1033}
\providecommand{\rhcaMseCiRidgeHarCloseOos}{}\renewcommand{\rhcaMseCiRidgeHarCloseOos}{4.15}
\providecommand{\rhcaMaeCiRidgeHarCloseOos}{}\renewcommand{\rhcaMaeCiRidgeHarCloseOos}{4.67}
\providecommand{\rhcaRsqCiRidgeHarCloseOos}{}\renewcommand{\rhcaRsqCiRidgeHarCloseOos}{78.2}
\providecommand{\rhcaGapCiRidgeHarCloseOos}{}\renewcommand{\rhcaGapCiRidgeHarCloseOos}{0.7}  % QLIKE % vs the closing-interval HAR benchmark
\providecommand{\rhcaGapDmCiRidgeHarCloseOos}{}\renewcommand{\rhcaGapDmCiRidgeHarCloseOos}{0.86}
\providecommand{\rhcaGapPCiRidgeHarCloseOos}{}\renewcommand{\rhcaGapPCiRidgeHarCloseOos}{\ensuremath{p = 0.39}}
\providecommand{\rhcaQCiRidgeFullCloseOos}{}\renewcommand{\rhcaQCiRidgeFullCloseOos}{0.1031}
\providecommand{\rhcaMseCiRidgeFullCloseOos}{}\renewcommand{\rhcaMseCiRidgeFullCloseOos}{7.30}
\providecommand{\rhcaMaeCiRidgeFullCloseOos}{}\renewcommand{\rhcaMaeCiRidgeFullCloseOos}{5.27}
\providecommand{\rhcaRsqCiRidgeFullCloseOos}{}\renewcommand{\rhcaRsqCiRidgeFullCloseOos}{61.7}
\providecommand{\rhcaPctCiRidgeCloseOos}{}\renewcommand{\rhcaPctCiRidgeCloseOos}{\ensuremath{-}0.1}  % QLIKE %, full vs HAR of the same class
\providecommand{\rhcaDmCiRidgeCloseOos}{}\renewcommand{\rhcaDmCiRidgeCloseOos}{\ensuremath{-}0.04}
\providecommand{\rhcaPCiRidgeCloseOos}{}\renewcommand{\rhcaPCiRidgeCloseOos}{\ensuremath{p = 0.97}}
\providecommand{\rhcaMsePctCiRidgeCloseOos}{}\renewcommand{\rhcaMsePctCiRidgeCloseOos}{75.8}  % MSE %, full vs HAR
\providecommand{\rhcaMseDmCiRidgeCloseOos}{}\renewcommand{\rhcaMseDmCiRidgeCloseOos}{1.45}
\providecommand{\rhcaMsePCiRidgeCloseOos}{}\renewcommand{\rhcaMsePCiRidgeCloseOos}{\ensuremath{p = 0.15}}
\providecommand{\rhcaGapCiRidgeFullCloseOos}{}\renewcommand{\rhcaGapCiRidgeFullCloseOos}{0.6}  % QLIKE % vs the closing-interval HAR benchmark
\providecommand{\rhcaGapDmCiRidgeFullCloseOos}{}\renewcommand{\rhcaGapDmCiRidgeFullCloseOos}{0.15}
\providecommand{\rhcaGapPCiRidgeFullCloseOos}{}\renewcommand{\rhcaGapPCiRidgeFullCloseOos}{\ensuremath{p = 0.88}}
\providecommand{\rhcaQCiLassoHarCloseOos}{}\renewcommand{\rhcaQCiLassoHarCloseOos}{0.1033}
\providecommand{\rhcaMseCiLassoHarCloseOos}{}\renewcommand{\rhcaMseCiLassoHarCloseOos}{4.15}
\providecommand{\rhcaMaeCiLassoHarCloseOos}{}\renewcommand{\rhcaMaeCiLassoHarCloseOos}{4.67}
\providecommand{\rhcaRsqCiLassoHarCloseOos}{}\renewcommand{\rhcaRsqCiLassoHarCloseOos}{78.2}
\providecommand{\rhcaGapCiLassoHarCloseOos}{}\renewcommand{\rhcaGapCiLassoHarCloseOos}{0.7}  % QLIKE % vs the closing-interval HAR benchmark
\providecommand{\rhcaGapDmCiLassoHarCloseOos}{}\renewcommand{\rhcaGapDmCiLassoHarCloseOos}{0.86}
\providecommand{\rhcaGapPCiLassoHarCloseOos}{}\renewcommand{\rhcaGapPCiLassoHarCloseOos}{\ensuremath{p = 0.39}}
\providecommand{\rhcaQCiLassoFullCloseOos}{}\renewcommand{\rhcaQCiLassoFullCloseOos}{0.0992}
\providecommand{\rhcaMseCiLassoFullCloseOos}{}\renewcommand{\rhcaMseCiLassoFullCloseOos}{5.94}
\providecommand{\rhcaMaeCiLassoFullCloseOos}{}\renewcommand{\rhcaMaeCiLassoFullCloseOos}{5.07}
\providecommand{\rhcaRsqCiLassoFullCloseOos}{}\renewcommand{\rhcaRsqCiLassoFullCloseOos}{68.8}
\providecommand{\rhcaPctCiLassoCloseOos}{}\renewcommand{\rhcaPctCiLassoCloseOos}{\ensuremath{-}4.0}  % QLIKE %, full vs HAR of the same class
\providecommand{\rhcaDmCiLassoCloseOos}{}\renewcommand{\rhcaDmCiLassoCloseOos}{\ensuremath{-}1.91}
\providecommand{\rhcaPCiLassoCloseOos}{}\renewcommand{\rhcaPCiLassoCloseOos}{\ensuremath{p = 0.056}}
\providecommand{\rhcaMsePctCiLassoCloseOos}{}\renewcommand{\rhcaMsePctCiLassoCloseOos}{43.1}  % MSE %, full vs HAR
\providecommand{\rhcaMseDmCiLassoCloseOos}{}\renewcommand{\rhcaMseDmCiLassoCloseOos}{1.34}
\providecommand{\rhcaMsePCiLassoCloseOos}{}\renewcommand{\rhcaMsePCiLassoCloseOos}{\ensuremath{p = 0.18}}
\providecommand{\rhcaGapCiLassoFullCloseOos}{}\renewcommand{\rhcaGapCiLassoFullCloseOos}{\ensuremath{-}3.3}  % QLIKE % vs the closing-interval HAR benchmark
\providecommand{\rhcaGapDmCiLassoFullCloseOos}{}\renewcommand{\rhcaGapDmCiLassoFullCloseOos}{\ensuremath{-}1.44}
\providecommand{\rhcaGapPCiLassoFullCloseOos}{}\renewcommand{\rhcaGapPCiLassoFullCloseOos}{\ensuremath{p = 0.15}}
\providecommand{\rhcaQCiEnetHarCloseOos}{}\renewcommand{\rhcaQCiEnetHarCloseOos}{0.1033}
\providecommand{\rhcaMseCiEnetHarCloseOos}{}\renewcommand{\rhcaMseCiEnetHarCloseOos}{4.15}
\providecommand{\rhcaMaeCiEnetHarCloseOos}{}\renewcommand{\rhcaMaeCiEnetHarCloseOos}{4.67}
\providecommand{\rhcaRsqCiEnetHarCloseOos}{}\renewcommand{\rhcaRsqCiEnetHarCloseOos}{78.2}
\providecommand{\rhcaGapCiEnetHarCloseOos}{}\renewcommand{\rhcaGapCiEnetHarCloseOos}{0.7}  % QLIKE % vs the closing-interval HAR benchmark
\providecommand{\rhcaGapDmCiEnetHarCloseOos}{}\renewcommand{\rhcaGapDmCiEnetHarCloseOos}{0.86}
\providecommand{\rhcaGapPCiEnetHarCloseOos}{}\renewcommand{\rhcaGapPCiEnetHarCloseOos}{\ensuremath{p = 0.39}}
\providecommand{\rhcaQCiEnetFullCloseOos}{}\renewcommand{\rhcaQCiEnetFullCloseOos}{0.1014}
\providecommand{\rhcaMseCiEnetFullCloseOos}{}\renewcommand{\rhcaMseCiEnetFullCloseOos}{6.89}
\providecommand{\rhcaMaeCiEnetFullCloseOos}{}\renewcommand{\rhcaMaeCiEnetFullCloseOos}{5.21}
\providecommand{\rhcaRsqCiEnetFullCloseOos}{}\renewcommand{\rhcaRsqCiEnetFullCloseOos}{63.9}
\providecommand{\rhcaPctCiEnetCloseOos}{}\renewcommand{\rhcaPctCiEnetCloseOos}{\ensuremath{-}1.8}  % QLIKE %, full vs HAR of the same class
\providecommand{\rhcaDmCiEnetCloseOos}{}\renewcommand{\rhcaDmCiEnetCloseOos}{\ensuremath{-}0.55}
\providecommand{\rhcaPCiEnetCloseOos}{}\renewcommand{\rhcaPCiEnetCloseOos}{\ensuremath{p = 0.58}}
\providecommand{\rhcaMsePctCiEnetCloseOos}{}\renewcommand{\rhcaMsePctCiEnetCloseOos}{65.9}  % MSE %, full vs HAR
\providecommand{\rhcaMseDmCiEnetCloseOos}{}\renewcommand{\rhcaMseDmCiEnetCloseOos}{1.45}
\providecommand{\rhcaMsePCiEnetCloseOos}{}\renewcommand{\rhcaMsePCiEnetCloseOos}{\ensuremath{p = 0.15}}
\providecommand{\rhcaGapCiEnetFullCloseOos}{}\renewcommand{\rhcaGapCiEnetFullCloseOos}{\ensuremath{-}1.1}  % QLIKE % vs the closing-interval HAR benchmark
\providecommand{\rhcaGapDmCiEnetFullCloseOos}{}\renewcommand{\rhcaGapDmCiEnetFullCloseOos}{\ensuremath{-}0.32}
\providecommand{\rhcaGapPCiEnetFullCloseOos}{}\renewcommand{\rhcaGapPCiEnetFullCloseOos}{\ensuremath{p = 0.75}}
\providecommand{\rhcaQCiLgbmHarCloseOos}{}\renewcommand{\rhcaQCiLgbmHarCloseOos}{0.1094}
\providecommand{\rhcaMseCiLgbmHarCloseOos}{}\renewcommand{\rhcaMseCiLgbmHarCloseOos}{5.75}
\providecommand{\rhcaMaeCiLgbmHarCloseOos}{}\renewcommand{\rhcaMaeCiLgbmHarCloseOos}{5.02}
\providecommand{\rhcaRsqCiLgbmHarCloseOos}{}\renewcommand{\rhcaRsqCiLgbmHarCloseOos}{69.8}
\providecommand{\rhcaGapCiLgbmHarCloseOos}{}\renewcommand{\rhcaGapCiLgbmHarCloseOos}{6.7}  % QLIKE % vs the closing-interval HAR benchmark
\providecommand{\rhcaGapDmCiLgbmHarCloseOos}{}\renewcommand{\rhcaGapDmCiLgbmHarCloseOos}{3.03}
\providecommand{\rhcaGapPCiLgbmHarCloseOos}{}\renewcommand{\rhcaGapPCiLgbmHarCloseOos}{\ensuremath{p = 0.002}}
\providecommand{\rhcaQCiLgbmFullCloseOos}{}\renewcommand{\rhcaQCiLgbmFullCloseOos}{0.1049}
\providecommand{\rhcaMseCiLgbmFullCloseOos}{}\renewcommand{\rhcaMseCiLgbmFullCloseOos}{17.19}
\providecommand{\rhcaMaeCiLgbmFullCloseOos}{}\renewcommand{\rhcaMaeCiLgbmFullCloseOos}{6.84}
\providecommand{\rhcaRsqCiLgbmFullCloseOos}{}\renewcommand{\rhcaRsqCiLgbmFullCloseOos}{9.8}
\providecommand{\rhcaPctCiLgbmCloseOos}{}\renewcommand{\rhcaPctCiLgbmCloseOos}{\ensuremath{-}4.1}  % QLIKE %, full vs HAR of the same class
\providecommand{\rhcaDmCiLgbmCloseOos}{}\renewcommand{\rhcaDmCiLgbmCloseOos}{\ensuremath{-}1.48}
\providecommand{\rhcaPCiLgbmCloseOos}{}\renewcommand{\rhcaPCiLgbmCloseOos}{\ensuremath{p = 0.14}}
\providecommand{\rhcaMsePctCiLgbmCloseOos}{}\renewcommand{\rhcaMsePctCiLgbmCloseOos}{199.1}  % MSE %, full vs HAR
\providecommand{\rhcaMseDmCiLgbmCloseOos}{}\renewcommand{\rhcaMseDmCiLgbmCloseOos}{1.21}
\providecommand{\rhcaMsePCiLgbmCloseOos}{}\renewcommand{\rhcaMsePCiLgbmCloseOos}{\ensuremath{p = 0.23}}
\providecommand{\rhcaGapCiLgbmFullCloseOos}{}\renewcommand{\rhcaGapCiLgbmFullCloseOos}{2.3}  % QLIKE % vs the closing-interval HAR benchmark
\providecommand{\rhcaGapDmCiLgbmFullCloseOos}{}\renewcommand{\rhcaGapDmCiLgbmFullCloseOos}{0.72}
\providecommand{\rhcaGapPCiLgbmFullCloseOos}{}\renewcommand{\rhcaGapPCiLgbmFullCloseOos}{\ensuremath{p = 0.47}}
\providecommand{\rhcaQCiXgbHarCloseOos}{}\renewcommand{\rhcaQCiXgbHarCloseOos}{0.1076}
\providecommand{\rhcaMseCiXgbHarCloseOos}{}\renewcommand{\rhcaMseCiXgbHarCloseOos}{5.27}
\providecommand{\rhcaMaeCiXgbHarCloseOos}{}\renewcommand{\rhcaMaeCiXgbHarCloseOos}{4.94}
\providecommand{\rhcaRsqCiXgbHarCloseOos}{}\renewcommand{\rhcaRsqCiXgbHarCloseOos}{72.4}
\providecommand{\rhcaGapCiXgbHarCloseOos}{}\renewcommand{\rhcaGapCiXgbHarCloseOos}{5.0}  % QLIKE % vs the closing-interval HAR benchmark
\providecommand{\rhcaGapDmCiXgbHarCloseOos}{}\renewcommand{\rhcaGapDmCiXgbHarCloseOos}{2.45}
\providecommand{\rhcaGapPCiXgbHarCloseOos}{}\renewcommand{\rhcaGapPCiXgbHarCloseOos}{\ensuremath{p = 0.014}}
\providecommand{\rhcaQCiXgbFullCloseOos}{}\renewcommand{\rhcaQCiXgbFullCloseOos}{0.1031}
\providecommand{\rhcaMseCiXgbFullCloseOos}{}\renewcommand{\rhcaMseCiXgbFullCloseOos}{20.94}
\providecommand{\rhcaMaeCiXgbFullCloseOos}{}\renewcommand{\rhcaMaeCiXgbFullCloseOos}{6.98}
\providecommand{\rhcaRsqCiXgbFullCloseOos}{}\renewcommand{\rhcaRsqCiXgbFullCloseOos}{\ensuremath{-}9.9}
\providecommand{\rhcaPctCiXgbCloseOos}{}\renewcommand{\rhcaPctCiXgbCloseOos}{\ensuremath{-}4.1}  % QLIKE %, full vs HAR of the same class
\providecommand{\rhcaDmCiXgbCloseOos}{}\renewcommand{\rhcaDmCiXgbCloseOos}{\ensuremath{-}1.63}
\providecommand{\rhcaPCiXgbCloseOos}{}\renewcommand{\rhcaPCiXgbCloseOos}{\ensuremath{p = 0.1039}}
\providecommand{\rhcaMsePctCiXgbCloseOos}{}\renewcommand{\rhcaMsePctCiXgbCloseOos}{297.6}  % MSE %, full vs HAR
\providecommand{\rhcaMseDmCiXgbCloseOos}{}\renewcommand{\rhcaMseDmCiXgbCloseOos}{1.15}
\providecommand{\rhcaMsePCiXgbCloseOos}{}\renewcommand{\rhcaMsePCiXgbCloseOos}{\ensuremath{p = 0.25}}
\providecommand{\rhcaGapCiXgbFullCloseOos}{}\renewcommand{\rhcaGapCiXgbFullCloseOos}{0.6}  % QLIKE % vs the closing-interval HAR benchmark
\providecommand{\rhcaGapDmCiXgbFullCloseOos}{}\renewcommand{\rhcaGapDmCiXgbFullCloseOos}{0.21}
\providecommand{\rhcaGapPCiXgbFullCloseOos}{}\renewcommand{\rhcaGapPCiXgbFullCloseOos}{\ensuremath{p = 0.83}}
\providecommand{\rhcaQCiRfHarCloseOos}{}\renewcommand{\rhcaQCiRfHarCloseOos}{0.1068}
\providecommand{\rhcaMseCiRfHarCloseOos}{}\renewcommand{\rhcaMseCiRfHarCloseOos}{4.62}
\providecommand{\rhcaMaeCiRfHarCloseOos}{}\renewcommand{\rhcaMaeCiRfHarCloseOos}{4.76}
\providecommand{\rhcaRsqCiRfHarCloseOos}{}\renewcommand{\rhcaRsqCiRfHarCloseOos}{75.8}
\providecommand{\rhcaGapCiRfHarCloseOos}{}\renewcommand{\rhcaGapCiRfHarCloseOos}{4.2}  % QLIKE % vs the closing-interval HAR benchmark
\providecommand{\rhcaGapDmCiRfHarCloseOos}{}\renewcommand{\rhcaGapDmCiRfHarCloseOos}{2.28}
\providecommand{\rhcaGapPCiRfHarCloseOos}{}\renewcommand{\rhcaGapPCiRfHarCloseOos}{\ensuremath{p = 0.022}}
\providecommand{\rhcaQCiRfFullCloseOos}{}\renewcommand{\rhcaQCiRfFullCloseOos}{0.1060}
\providecommand{\rhcaMseCiRfFullCloseOos}{}\renewcommand{\rhcaMseCiRfFullCloseOos}{14.15}
\providecommand{\rhcaMaeCiRfFullCloseOos}{}\renewcommand{\rhcaMaeCiRfFullCloseOos}{6.74}
\providecommand{\rhcaRsqCiRfFullCloseOos}{}\renewcommand{\rhcaRsqCiRfFullCloseOos}{25.8}
\providecommand{\rhcaPctCiRfCloseOos}{}\renewcommand{\rhcaPctCiRfCloseOos}{\ensuremath{-}0.8}  % QLIKE %, full vs HAR of the same class
\providecommand{\rhcaDmCiRfCloseOos}{}\renewcommand{\rhcaDmCiRfCloseOos}{\ensuremath{-}0.28}
\providecommand{\rhcaPCiRfCloseOos}{}\renewcommand{\rhcaPCiRfCloseOos}{\ensuremath{p = 0.78}}
\providecommand{\rhcaMsePctCiRfCloseOos}{}\renewcommand{\rhcaMsePctCiRfCloseOos}{206.6}  % MSE %, full vs HAR
\providecommand{\rhcaMseDmCiRfCloseOos}{}\renewcommand{\rhcaMseDmCiRfCloseOos}{1.22}
\providecommand{\rhcaMsePCiRfCloseOos}{}\renewcommand{\rhcaMsePCiRfCloseOos}{\ensuremath{p = 0.22}}
\providecommand{\rhcaGapCiRfFullCloseOos}{}\renewcommand{\rhcaGapCiRfFullCloseOos}{3.4}  % QLIKE % vs the closing-interval HAR benchmark
\providecommand{\rhcaGapDmCiRfFullCloseOos}{}\renewcommand{\rhcaGapDmCiRfFullCloseOos}{1.11}
\providecommand{\rhcaGapPCiRfFullCloseOos}{}\renewcommand{\rhcaGapPCiRfFullCloseOos}{\ensuremath{p = 0.27}}
% --- every bar: pooled forecasts, Span / Pre / Oos (Table 1 and its subperiods)
\providecommand{\rhcaQOlsHarBarSpan}{}\renewcommand{\rhcaQOlsHarBarSpan}{0.1893}
\providecommand{\rhcaMseOlsHarBarSpan}{}\renewcommand{\rhcaMseOlsHarBarSpan}{1.11}
\providecommand{\rhcaMaeOlsHarBarSpan}{}\renewcommand{\rhcaMaeOlsHarBarSpan}{2.25}
\providecommand{\rhcaRsqOlsHarBarSpan}{}\renewcommand{\rhcaRsqOlsHarBarSpan}{65.4}
\providecommand{\rhcaQRidgeHarBarSpan}{}\renewcommand{\rhcaQRidgeHarBarSpan}{0.1893}
\providecommand{\rhcaMseRidgeHarBarSpan}{}\renewcommand{\rhcaMseRidgeHarBarSpan}{1.11}
\providecommand{\rhcaMaeRidgeHarBarSpan}{}\renewcommand{\rhcaMaeRidgeHarBarSpan}{2.25}
\providecommand{\rhcaRsqRidgeHarBarSpan}{}\renewcommand{\rhcaRsqRidgeHarBarSpan}{65.4}
\providecommand{\rhcaQRidgeFullBarSpan}{}\renewcommand{\rhcaQRidgeFullBarSpan}{0.1778}
\providecommand{\rhcaMseRidgeFullBarSpan}{}\renewcommand{\rhcaMseRidgeFullBarSpan}{2.23}
\providecommand{\rhcaMaeRidgeFullBarSpan}{}\renewcommand{\rhcaMaeRidgeFullBarSpan}{2.38}
\providecommand{\rhcaRsqRidgeFullBarSpan}{}\renewcommand{\rhcaRsqRidgeFullBarSpan}{30.7}
\providecommand{\rhcaPctRidgeBarSpan}{}\renewcommand{\rhcaPctRidgeBarSpan}{\ensuremath{-}6.0}
\providecommand{\rhcaDmRidgeBarSpan}{}\renewcommand{\rhcaDmRidgeBarSpan}{\ensuremath{-}6.28}
\providecommand{\rhcaPRidgeBarSpan}{}\renewcommand{\rhcaPRidgeBarSpan}{\ensuremath{p < 0.001}}
\providecommand{\rhcaMsePctRidgeBarSpan}{}\renewcommand{\rhcaMsePctRidgeBarSpan}{100.6}
\providecommand{\rhcaMseDmRidgeBarSpan}{}\renewcommand{\rhcaMseDmRidgeBarSpan}{1.12}
\providecommand{\rhcaMsePRidgeBarSpan}{}\renewcommand{\rhcaMsePRidgeBarSpan}{\ensuremath{p = 0.26}}
\providecommand{\rhcaQLassoHarBarSpan}{}\renewcommand{\rhcaQLassoHarBarSpan}{0.1893}
\providecommand{\rhcaMseLassoHarBarSpan}{}\renewcommand{\rhcaMseLassoHarBarSpan}{1.11}
\providecommand{\rhcaMaeLassoHarBarSpan}{}\renewcommand{\rhcaMaeLassoHarBarSpan}{2.25}
\providecommand{\rhcaRsqLassoHarBarSpan}{}\renewcommand{\rhcaRsqLassoHarBarSpan}{65.4}
\providecommand{\rhcaQLassoFullBarSpan}{}\renewcommand{\rhcaQLassoFullBarSpan}{0.1723}
\providecommand{\rhcaMseLassoFullBarSpan}{}\renewcommand{\rhcaMseLassoFullBarSpan}{2.39}
\providecommand{\rhcaMaeLassoFullBarSpan}{}\renewcommand{\rhcaMaeLassoFullBarSpan}{2.40}
\providecommand{\rhcaRsqLassoFullBarSpan}{}\renewcommand{\rhcaRsqLassoFullBarSpan}{25.8}
\providecommand{\rhcaPctLassoBarSpan}{}\renewcommand{\rhcaPctLassoBarSpan}{\ensuremath{-}9.0}
\providecommand{\rhcaDmLassoBarSpan}{}\renewcommand{\rhcaDmLassoBarSpan}{\ensuremath{-}6.34}
\providecommand{\rhcaPLassoBarSpan}{}\renewcommand{\rhcaPLassoBarSpan}{\ensuremath{p < 0.001}}
\providecommand{\rhcaMsePctLassoBarSpan}{}\renewcommand{\rhcaMsePctLassoBarSpan}{114.7}
\providecommand{\rhcaMseDmLassoBarSpan}{}\renewcommand{\rhcaMseDmLassoBarSpan}{1.16}
\providecommand{\rhcaMsePLassoBarSpan}{}\renewcommand{\rhcaMsePLassoBarSpan}{\ensuremath{p = 0.25}}
\providecommand{\rhcaQEnetHarBarSpan}{}\renewcommand{\rhcaQEnetHarBarSpan}{0.1893}
\providecommand{\rhcaMseEnetHarBarSpan}{}\renewcommand{\rhcaMseEnetHarBarSpan}{1.11}
\providecommand{\rhcaMaeEnetHarBarSpan}{}\renewcommand{\rhcaMaeEnetHarBarSpan}{2.25}
\providecommand{\rhcaRsqEnetHarBarSpan}{}\renewcommand{\rhcaRsqEnetHarBarSpan}{65.4}
\providecommand{\rhcaQEnetFullBarSpan}{}\renewcommand{\rhcaQEnetFullBarSpan}{0.1766}
\providecommand{\rhcaMseEnetFullBarSpan}{}\renewcommand{\rhcaMseEnetFullBarSpan}{2.42}
\providecommand{\rhcaMaeEnetFullBarSpan}{}\renewcommand{\rhcaMaeEnetFullBarSpan}{2.40}
\providecommand{\rhcaRsqEnetFullBarSpan}{}\renewcommand{\rhcaRsqEnetFullBarSpan}{24.8}
\providecommand{\rhcaPctEnetBarSpan}{}\renewcommand{\rhcaPctEnetBarSpan}{\ensuremath{-}6.7}
\providecommand{\rhcaDmEnetBarSpan}{}\renewcommand{\rhcaDmEnetBarSpan}{\ensuremath{-}7.36}
\providecommand{\rhcaPEnetBarSpan}{}\renewcommand{\rhcaPEnetBarSpan}{\ensuremath{p < 0.001}}
\providecommand{\rhcaMsePctEnetBarSpan}{}\renewcommand{\rhcaMsePctEnetBarSpan}{117.7}
\providecommand{\rhcaMseDmEnetBarSpan}{}\renewcommand{\rhcaMseDmEnetBarSpan}{1.12}
\providecommand{\rhcaMsePEnetBarSpan}{}\renewcommand{\rhcaMsePEnetBarSpan}{\ensuremath{p = 0.26}}
\providecommand{\rhcaQLgbmHarBarSpan}{}\renewcommand{\rhcaQLgbmHarBarSpan}{0.1875}
\providecommand{\rhcaMseLgbmHarBarSpan}{}\renewcommand{\rhcaMseLgbmHarBarSpan}{1.18}
\providecommand{\rhcaMaeLgbmHarBarSpan}{}\renewcommand{\rhcaMaeLgbmHarBarSpan}{2.27}
\providecommand{\rhcaRsqLgbmHarBarSpan}{}\renewcommand{\rhcaRsqLgbmHarBarSpan}{63.3}
\providecommand{\rhcaQLgbmFullBarSpan}{}\renewcommand{\rhcaQLgbmFullBarSpan}{0.1611}
\providecommand{\rhcaMseLgbmFullBarSpan}{}\renewcommand{\rhcaMseLgbmFullBarSpan}{1.42}
\providecommand{\rhcaMaeLgbmFullBarSpan}{}\renewcommand{\rhcaMaeLgbmFullBarSpan}{2.28}
\providecommand{\rhcaRsqLgbmFullBarSpan}{}\renewcommand{\rhcaRsqLgbmFullBarSpan}{55.9}
\providecommand{\rhcaPctLgbmBarSpan}{}\renewcommand{\rhcaPctLgbmBarSpan}{\ensuremath{-}14.1}
\providecommand{\rhcaDmLgbmBarSpan}{}\renewcommand{\rhcaDmLgbmBarSpan}{\ensuremath{-}10.26}
\providecommand{\rhcaPLgbmBarSpan}{}\renewcommand{\rhcaPLgbmBarSpan}{\ensuremath{p < 0.001}}
\providecommand{\rhcaMsePctLgbmBarSpan}{}\renewcommand{\rhcaMsePctLgbmBarSpan}{20.0}
\providecommand{\rhcaMseDmLgbmBarSpan}{}\renewcommand{\rhcaMseDmLgbmBarSpan}{1.03}
\providecommand{\rhcaMsePLgbmBarSpan}{}\renewcommand{\rhcaMsePLgbmBarSpan}{\ensuremath{p = 0.30}}
\providecommand{\rhcaQXgbHarBarSpan}{}\renewcommand{\rhcaQXgbHarBarSpan}{0.1888}
\providecommand{\rhcaMseXgbHarBarSpan}{}\renewcommand{\rhcaMseXgbHarBarSpan}{1.16}
\providecommand{\rhcaMaeXgbHarBarSpan}{}\renewcommand{\rhcaMaeXgbHarBarSpan}{2.26}
\providecommand{\rhcaRsqXgbHarBarSpan}{}\renewcommand{\rhcaRsqXgbHarBarSpan}{63.9}
\providecommand{\rhcaQXgbFullBarSpan}{}\renewcommand{\rhcaQXgbFullBarSpan}{0.1612}
\providecommand{\rhcaMseXgbFullBarSpan}{}\renewcommand{\rhcaMseXgbFullBarSpan}{1.35}
\providecommand{\rhcaMaeXgbFullBarSpan}{}\renewcommand{\rhcaMaeXgbFullBarSpan}{2.27}
\providecommand{\rhcaRsqXgbFullBarSpan}{}\renewcommand{\rhcaRsqXgbFullBarSpan}{58.0}
\providecommand{\rhcaPctXgbBarSpan}{}\renewcommand{\rhcaPctXgbBarSpan}{\ensuremath{-}14.6}
\providecommand{\rhcaDmXgbBarSpan}{}\renewcommand{\rhcaDmXgbBarSpan}{\ensuremath{-}9.94}
\providecommand{\rhcaPXgbBarSpan}{}\renewcommand{\rhcaPXgbBarSpan}{\ensuremath{p < 0.001}}
\providecommand{\rhcaMsePctXgbBarSpan}{}\renewcommand{\rhcaMsePctXgbBarSpan}{16.3}
\providecommand{\rhcaMseDmXgbBarSpan}{}\renewcommand{\rhcaMseDmXgbBarSpan}{1.04}
\providecommand{\rhcaMsePXgbBarSpan}{}\renewcommand{\rhcaMsePXgbBarSpan}{\ensuremath{p = 0.30}}
\providecommand{\rhcaQRfHarBarSpan}{}\renewcommand{\rhcaQRfHarBarSpan}{0.1895}
\providecommand{\rhcaMseRfHarBarSpan}{}\renewcommand{\rhcaMseRfHarBarSpan}{1.19}
\providecommand{\rhcaMaeRfHarBarSpan}{}\renewcommand{\rhcaMaeRfHarBarSpan}{2.27}
\providecommand{\rhcaRsqRfHarBarSpan}{}\renewcommand{\rhcaRsqRfHarBarSpan}{63.0}
\providecommand{\rhcaQRfFullBarSpan}{}\renewcommand{\rhcaQRfFullBarSpan}{0.1681}
\providecommand{\rhcaMseRfFullBarSpan}{}\renewcommand{\rhcaMseRfFullBarSpan}{1.32}
\providecommand{\rhcaMaeRfFullBarSpan}{}\renewcommand{\rhcaMaeRfFullBarSpan}{2.27}
\providecommand{\rhcaRsqRfFullBarSpan}{}\renewcommand{\rhcaRsqRfFullBarSpan}{58.9}
\providecommand{\rhcaPctRfBarSpan}{}\renewcommand{\rhcaPctRfBarSpan}{\ensuremath{-}11.3}
\providecommand{\rhcaDmRfBarSpan}{}\renewcommand{\rhcaDmRfBarSpan}{\ensuremath{-}9.74}
\providecommand{\rhcaPRfBarSpan}{}\renewcommand{\rhcaPRfBarSpan}{\ensuremath{p < 0.001}}
\providecommand{\rhcaMsePctRfBarSpan}{}\renewcommand{\rhcaMsePctRfBarSpan}{10.9}
\providecommand{\rhcaMseDmRfBarSpan}{}\renewcommand{\rhcaMseDmRfBarSpan}{1.10}
\providecommand{\rhcaMsePRfBarSpan}{}\renewcommand{\rhcaMsePRfBarSpan}{\ensuremath{p = 0.27}}
\providecommand{\rhcaQOlsHarBarPre}{}\renewcommand{\rhcaQOlsHarBarPre}{0.1864}
\providecommand{\rhcaMseOlsHarBarPre}{}\renewcommand{\rhcaMseOlsHarBarPre}{0.31}
\providecommand{\rhcaMaeOlsHarBarPre}{}\renewcommand{\rhcaMaeOlsHarBarPre}{1.41}
\providecommand{\rhcaRsqOlsHarBarPre}{}\renewcommand{\rhcaRsqOlsHarBarPre}{61.4}
\providecommand{\rhcaQRidgeHarBarPre}{}\renewcommand{\rhcaQRidgeHarBarPre}{0.1864}
\providecommand{\rhcaMseRidgeHarBarPre}{}\renewcommand{\rhcaMseRidgeHarBarPre}{0.31}
\providecommand{\rhcaMaeRidgeHarBarPre}{}\renewcommand{\rhcaMaeRidgeHarBarPre}{1.41}
\providecommand{\rhcaRsqRidgeHarBarPre}{}\renewcommand{\rhcaRsqRidgeHarBarPre}{61.4}
\providecommand{\rhcaQRidgeFullBarPre}{}\renewcommand{\rhcaQRidgeFullBarPre}{0.1726}
\providecommand{\rhcaMseRidgeFullBarPre}{}\renewcommand{\rhcaMseRidgeFullBarPre}{0.28}
\providecommand{\rhcaMaeRidgeFullBarPre}{}\renewcommand{\rhcaMaeRidgeFullBarPre}{1.35}
\providecommand{\rhcaRsqRidgeFullBarPre}{}\renewcommand{\rhcaRsqRidgeFullBarPre}{64.9}
\providecommand{\rhcaPctRidgeBarPre}{}\renewcommand{\rhcaPctRidgeBarPre}{\ensuremath{-}7.4}
\providecommand{\rhcaDmRidgeBarPre}{}\renewcommand{\rhcaDmRidgeBarPre}{\ensuremath{-}6.82}
\providecommand{\rhcaPRidgeBarPre}{}\renewcommand{\rhcaPRidgeBarPre}{\ensuremath{p < 0.001}}
\providecommand{\rhcaMsePctRidgeBarPre}{}\renewcommand{\rhcaMsePctRidgeBarPre}{\ensuremath{-}8.9}
\providecommand{\rhcaMseDmRidgeBarPre}{}\renewcommand{\rhcaMseDmRidgeBarPre}{\ensuremath{-}1.79}
\providecommand{\rhcaMsePRidgeBarPre}{}\renewcommand{\rhcaMsePRidgeBarPre}{\ensuremath{p = 0.074}}
\providecommand{\rhcaQLassoHarBarPre}{}\renewcommand{\rhcaQLassoHarBarPre}{0.1864}
\providecommand{\rhcaMseLassoHarBarPre}{}\renewcommand{\rhcaMseLassoHarBarPre}{0.31}
\providecommand{\rhcaMaeLassoHarBarPre}{}\renewcommand{\rhcaMaeLassoHarBarPre}{1.41}
\providecommand{\rhcaRsqLassoHarBarPre}{}\renewcommand{\rhcaRsqLassoHarBarPre}{61.4}
\providecommand{\rhcaQLassoFullBarPre}{}\renewcommand{\rhcaQLassoFullBarPre}{0.1667}
\providecommand{\rhcaMseLassoFullBarPre}{}\renewcommand{\rhcaMseLassoFullBarPre}{0.28}
\providecommand{\rhcaMaeLassoFullBarPre}{}\renewcommand{\rhcaMaeLassoFullBarPre}{1.35}
\providecommand{\rhcaRsqLassoFullBarPre}{}\renewcommand{\rhcaRsqLassoFullBarPre}{65.6}
\providecommand{\rhcaPctLassoBarPre}{}\renewcommand{\rhcaPctLassoBarPre}{\ensuremath{-}10.6}
\providecommand{\rhcaDmLassoBarPre}{}\renewcommand{\rhcaDmLassoBarPre}{\ensuremath{-}6.88}
\providecommand{\rhcaPLassoBarPre}{}\renewcommand{\rhcaPLassoBarPre}{\ensuremath{p < 0.001}}
\providecommand{\rhcaMsePctLassoBarPre}{}\renewcommand{\rhcaMsePctLassoBarPre}{\ensuremath{-}10.8}
\providecommand{\rhcaMseDmLassoBarPre}{}\renewcommand{\rhcaMseDmLassoBarPre}{\ensuremath{-}1.81}
\providecommand{\rhcaMsePLassoBarPre}{}\renewcommand{\rhcaMsePLassoBarPre}{\ensuremath{p = 0.070}}
\providecommand{\rhcaQEnetHarBarPre}{}\renewcommand{\rhcaQEnetHarBarPre}{0.1864}
\providecommand{\rhcaMseEnetHarBarPre}{}\renewcommand{\rhcaMseEnetHarBarPre}{0.31}
\providecommand{\rhcaMaeEnetHarBarPre}{}\renewcommand{\rhcaMaeEnetHarBarPre}{1.41}
\providecommand{\rhcaRsqEnetHarBarPre}{}\renewcommand{\rhcaRsqEnetHarBarPre}{61.4}
\providecommand{\rhcaQEnetFullBarPre}{}\renewcommand{\rhcaQEnetFullBarPre}{0.1720}
\providecommand{\rhcaMseEnetFullBarPre}{}\renewcommand{\rhcaMseEnetFullBarPre}{0.28}
\providecommand{\rhcaMaeEnetFullBarPre}{}\renewcommand{\rhcaMaeEnetFullBarPre}{1.35}
\providecommand{\rhcaRsqEnetFullBarPre}{}\renewcommand{\rhcaRsqEnetFullBarPre}{65.2}
\providecommand{\rhcaPctEnetBarPre}{}\renewcommand{\rhcaPctEnetBarPre}{\ensuremath{-}7.7}
\providecommand{\rhcaDmEnetBarPre}{}\renewcommand{\rhcaDmEnetBarPre}{\ensuremath{-}6.97}
\providecommand{\rhcaPEnetBarPre}{}\renewcommand{\rhcaPEnetBarPre}{\ensuremath{p < 0.001}}
\providecommand{\rhcaMsePctEnetBarPre}{}\renewcommand{\rhcaMsePctEnetBarPre}{\ensuremath{-}9.7}
\providecommand{\rhcaMseDmEnetBarPre}{}\renewcommand{\rhcaMseDmEnetBarPre}{\ensuremath{-}1.75}
\providecommand{\rhcaMsePEnetBarPre}{}\renewcommand{\rhcaMsePEnetBarPre}{\ensuremath{p = 0.080}}
\providecommand{\rhcaQLgbmHarBarPre}{}\renewcommand{\rhcaQLgbmHarBarPre}{0.1858}
\providecommand{\rhcaMseLgbmHarBarPre}{}\renewcommand{\rhcaMseLgbmHarBarPre}{0.37}
\providecommand{\rhcaMaeLgbmHarBarPre}{}\renewcommand{\rhcaMaeLgbmHarBarPre}{1.43}
\providecommand{\rhcaRsqLgbmHarBarPre}{}\renewcommand{\rhcaRsqLgbmHarBarPre}{53.3}
\providecommand{\rhcaQLgbmFullBarPre}{}\renewcommand{\rhcaQLgbmFullBarPre}{0.1586}
\providecommand{\rhcaMseLgbmFullBarPre}{}\renewcommand{\rhcaMseLgbmFullBarPre}{0.36}
\providecommand{\rhcaMaeLgbmFullBarPre}{}\renewcommand{\rhcaMaeLgbmFullBarPre}{1.36}
\providecommand{\rhcaRsqLgbmFullBarPre}{}\renewcommand{\rhcaRsqLgbmFullBarPre}{55.5}
\providecommand{\rhcaPctLgbmBarPre}{}\renewcommand{\rhcaPctLgbmBarPre}{\ensuremath{-}14.6}
\providecommand{\rhcaDmLgbmBarPre}{}\renewcommand{\rhcaDmLgbmBarPre}{\ensuremath{-}7.91}
\providecommand{\rhcaPLgbmBarPre}{}\renewcommand{\rhcaPLgbmBarPre}{\ensuremath{p < 0.001}}
\providecommand{\rhcaMsePctLgbmBarPre}{}\renewcommand{\rhcaMsePctLgbmBarPre}{\ensuremath{-}4.7}
\providecommand{\rhcaMseDmLgbmBarPre}{}\renewcommand{\rhcaMseDmLgbmBarPre}{\ensuremath{-}2.59}
\providecommand{\rhcaMsePLgbmBarPre}{}\renewcommand{\rhcaMsePLgbmBarPre}{\ensuremath{p = 0.0097}}
\providecommand{\rhcaQXgbHarBarPre}{}\renewcommand{\rhcaQXgbHarBarPre}{0.1864}
\providecommand{\rhcaMseXgbHarBarPre}{}\renewcommand{\rhcaMseXgbHarBarPre}{0.37}
\providecommand{\rhcaMaeXgbHarBarPre}{}\renewcommand{\rhcaMaeXgbHarBarPre}{1.43}
\providecommand{\rhcaRsqXgbHarBarPre}{}\renewcommand{\rhcaRsqXgbHarBarPre}{53.9}
\providecommand{\rhcaQXgbFullBarPre}{}\renewcommand{\rhcaQXgbFullBarPre}{0.1597}
\providecommand{\rhcaMseXgbFullBarPre}{}\renewcommand{\rhcaMseXgbFullBarPre}{0.36}
\providecommand{\rhcaMaeXgbFullBarPre}{}\renewcommand{\rhcaMaeXgbFullBarPre}{1.37}
\providecommand{\rhcaRsqXgbFullBarPre}{}\renewcommand{\rhcaRsqXgbFullBarPre}{55.5}
\providecommand{\rhcaPctXgbBarPre}{}\renewcommand{\rhcaPctXgbBarPre}{\ensuremath{-}14.3}
\providecommand{\rhcaDmXgbBarPre}{}\renewcommand{\rhcaDmXgbBarPre}{\ensuremath{-}7.48}
\providecommand{\rhcaPXgbBarPre}{}\renewcommand{\rhcaPXgbBarPre}{\ensuremath{p < 0.001}}
\providecommand{\rhcaMsePctXgbBarPre}{}\renewcommand{\rhcaMsePctXgbBarPre}{\ensuremath{-}3.4}
\providecommand{\rhcaMseDmXgbBarPre}{}\renewcommand{\rhcaMseDmXgbBarPre}{\ensuremath{-}2.50}
\providecommand{\rhcaMsePXgbBarPre}{}\renewcommand{\rhcaMsePXgbBarPre}{\ensuremath{p = 0.012}}
\providecommand{\rhcaQRfHarBarPre}{}\renewcommand{\rhcaQRfHarBarPre}{0.1865}
\providecommand{\rhcaMseRfHarBarPre}{}\renewcommand{\rhcaMseRfHarBarPre}{0.37}
\providecommand{\rhcaMaeRfHarBarPre}{}\renewcommand{\rhcaMaeRfHarBarPre}{1.43}
\providecommand{\rhcaRsqRfHarBarPre}{}\renewcommand{\rhcaRsqRfHarBarPre}{53.9}
\providecommand{\rhcaQRfFullBarPre}{}\renewcommand{\rhcaQRfFullBarPre}{0.1646}
\providecommand{\rhcaMseRfFullBarPre}{}\renewcommand{\rhcaMseRfFullBarPre}{0.36}
\providecommand{\rhcaMaeRfFullBarPre}{}\renewcommand{\rhcaMaeRfFullBarPre}{1.37}
\providecommand{\rhcaRsqRfFullBarPre}{}\renewcommand{\rhcaRsqRfFullBarPre}{54.6}
\providecommand{\rhcaPctRfBarPre}{}\renewcommand{\rhcaPctRfBarPre}{\ensuremath{-}11.7}
\providecommand{\rhcaDmRfBarPre}{}\renewcommand{\rhcaDmRfBarPre}{\ensuremath{-}7.21}
\providecommand{\rhcaPRfBarPre}{}\renewcommand{\rhcaPRfBarPre}{\ensuremath{p < 0.001}}
\providecommand{\rhcaMsePctRfBarPre}{}\renewcommand{\rhcaMsePctRfBarPre}{\ensuremath{-}1.6}
\providecommand{\rhcaMseDmRfBarPre}{}\renewcommand{\rhcaMseDmRfBarPre}{\ensuremath{-}1.74}
\providecommand{\rhcaMsePRfBarPre}{}\renewcommand{\rhcaMsePRfBarPre}{\ensuremath{p = 0.083}}
\providecommand{\rhcaQOlsHarBarOos}{}\renewcommand{\rhcaQOlsHarBarOos}{0.1933}
\providecommand{\rhcaMseOlsHarBarOos}{}\renewcommand{\rhcaMseOlsHarBarOos}{2.22}
\providecommand{\rhcaMaeOlsHarBarOos}{}\renewcommand{\rhcaMaeOlsHarBarOos}{3.42}
\providecommand{\rhcaRsqOlsHarBarOos}{}\renewcommand{\rhcaRsqOlsHarBarOos}{66.1}
\providecommand{\rhcaQRidgeHarBarOos}{}\renewcommand{\rhcaQRidgeHarBarOos}{0.1933}
\providecommand{\rhcaMseRidgeHarBarOos}{}\renewcommand{\rhcaMseRidgeHarBarOos}{2.22}
\providecommand{\rhcaMaeRidgeHarBarOos}{}\renewcommand{\rhcaMaeRidgeHarBarOos}{3.42}
\providecommand{\rhcaRsqRidgeHarBarOos}{}\renewcommand{\rhcaRsqRidgeHarBarOos}{66.1}
\providecommand{\rhcaQRidgeFullBarOos}{}\renewcommand{\rhcaQRidgeFullBarOos}{0.1852}
\providecommand{\rhcaMseRidgeFullBarOos}{}\renewcommand{\rhcaMseRidgeFullBarOos}{4.93}
\providecommand{\rhcaMaeRidgeFullBarOos}{}\renewcommand{\rhcaMaeRidgeFullBarOos}{3.80}
\providecommand{\rhcaRsqRidgeFullBarOos}{}\renewcommand{\rhcaRsqRidgeFullBarOos}{24.9}
\providecommand{\rhcaPctRidgeBarOos}{}\renewcommand{\rhcaPctRidgeBarOos}{\ensuremath{-}4.2}
\providecommand{\rhcaDmRidgeBarOos}{}\renewcommand{\rhcaDmRidgeBarOos}{\ensuremath{-}2.52}
\providecommand{\rhcaPRidgeBarOos}{}\renewcommand{\rhcaPRidgeBarOos}{\ensuremath{p = 0.012}}
\providecommand{\rhcaMsePctRidgeBarOos}{}\renewcommand{\rhcaMsePctRidgeBarOos}{121.6}
\providecommand{\rhcaMseDmRidgeBarOos}{}\renewcommand{\rhcaMseDmRidgeBarOos}{1.19}
\providecommand{\rhcaMsePRidgeBarOos}{}\renewcommand{\rhcaMsePRidgeBarOos}{\ensuremath{p = 0.24}}
\providecommand{\rhcaQLassoHarBarOos}{}\renewcommand{\rhcaQLassoHarBarOos}{0.1933}
\providecommand{\rhcaMseLassoHarBarOos}{}\renewcommand{\rhcaMseLassoHarBarOos}{2.22}
\providecommand{\rhcaMaeLassoHarBarOos}{}\renewcommand{\rhcaMaeLassoHarBarOos}{3.42}
\providecommand{\rhcaRsqLassoHarBarOos}{}\renewcommand{\rhcaRsqLassoHarBarOos}{66.1}
\providecommand{\rhcaQLassoFullBarOos}{}\renewcommand{\rhcaQLassoFullBarOos}{0.1801}
\providecommand{\rhcaMseLassoFullBarOos}{}\renewcommand{\rhcaMseLassoFullBarOos}{5.31}
\providecommand{\rhcaMaeLassoFullBarOos}{}\renewcommand{\rhcaMaeLassoFullBarOos}{3.86}
\providecommand{\rhcaRsqLassoFullBarOos}{}\renewcommand{\rhcaRsqLassoFullBarOos}{19.1}
\providecommand{\rhcaPctLassoBarOos}{}\renewcommand{\rhcaPctLassoBarOos}{\ensuremath{-}6.8}
\providecommand{\rhcaDmLassoBarOos}{}\renewcommand{\rhcaDmLassoBarOos}{\ensuremath{-}2.70}
\providecommand{\rhcaPLassoBarOos}{}\renewcommand{\rhcaPLassoBarOos}{\ensuremath{p = 0.007}}
\providecommand{\rhcaMsePctLassoBarOos}{}\renewcommand{\rhcaMsePctLassoBarOos}{138.8}
\providecommand{\rhcaMseDmLassoBarOos}{}\renewcommand{\rhcaMseDmLassoBarOos}{1.23}
\providecommand{\rhcaMsePLassoBarOos}{}\renewcommand{\rhcaMsePLassoBarOos}{\ensuremath{p = 0.22}}
\providecommand{\rhcaQEnetHarBarOos}{}\renewcommand{\rhcaQEnetHarBarOos}{0.1933}
\providecommand{\rhcaMseEnetHarBarOos}{}\renewcommand{\rhcaMseEnetHarBarOos}{2.22}
\providecommand{\rhcaMaeEnetHarBarOos}{}\renewcommand{\rhcaMaeEnetHarBarOos}{3.42}
\providecommand{\rhcaRsqEnetHarBarOos}{}\renewcommand{\rhcaRsqEnetHarBarOos}{66.1}
\providecommand{\rhcaQEnetFullBarOos}{}\renewcommand{\rhcaQEnetFullBarOos}{0.1831}
\providecommand{\rhcaMseEnetFullBarOos}{}\renewcommand{\rhcaMseEnetFullBarOos}{5.38}
\providecommand{\rhcaMaeEnetFullBarOos}{}\renewcommand{\rhcaMaeEnetFullBarOos}{3.85}
\providecommand{\rhcaRsqEnetFullBarOos}{}\renewcommand{\rhcaRsqEnetFullBarOos}{18.0}
\providecommand{\rhcaPctEnetBarOos}{}\renewcommand{\rhcaPctEnetBarOos}{\ensuremath{-}5.3}
\providecommand{\rhcaDmEnetBarOos}{}\renewcommand{\rhcaDmEnetBarOos}{\ensuremath{-}3.58}
\providecommand{\rhcaPEnetBarOos}{}\renewcommand{\rhcaPEnetBarOos}{\ensuremath{p < 0.001}}
\providecommand{\rhcaMsePctEnetBarOos}{}\renewcommand{\rhcaMsePctEnetBarOos}{142.1}
\providecommand{\rhcaMseDmEnetBarOos}{}\renewcommand{\rhcaMseDmEnetBarOos}{1.18}
\providecommand{\rhcaMsePEnetBarOos}{}\renewcommand{\rhcaMsePEnetBarOos}{\ensuremath{p = 0.24}}
\providecommand{\rhcaQLgbmHarBarOos}{}\renewcommand{\rhcaQLgbmHarBarOos}{0.1900}
\providecommand{\rhcaMseLgbmHarBarOos}{}\renewcommand{\rhcaMseLgbmHarBarOos}{2.30}
\providecommand{\rhcaMaeLgbmHarBarOos}{}\renewcommand{\rhcaMaeLgbmHarBarOos}{3.43}
\providecommand{\rhcaRsqLgbmHarBarOos}{}\renewcommand{\rhcaRsqLgbmHarBarOos}{65.0}
\providecommand{\rhcaQLgbmFullBarOos}{}\renewcommand{\rhcaQLgbmFullBarOos}{0.1645}
\providecommand{\rhcaMseLgbmFullBarOos}{}\renewcommand{\rhcaMseLgbmFullBarOos}{2.89}
\providecommand{\rhcaMaeLgbmFullBarOos}{}\renewcommand{\rhcaMaeLgbmFullBarOos}{3.55}
\providecommand{\rhcaRsqLgbmFullBarOos}{}\renewcommand{\rhcaRsqLgbmFullBarOos}{56.0}
\providecommand{\rhcaPctLgbmBarOos}{}\renewcommand{\rhcaPctLgbmBarOos}{\ensuremath{-}13.4}
\providecommand{\rhcaDmLgbmBarOos}{}\renewcommand{\rhcaDmLgbmBarOos}{\ensuremath{-}6.42}
\providecommand{\rhcaPLgbmBarOos}{}\renewcommand{\rhcaPLgbmBarOos}{\ensuremath{p < 0.001}}
\providecommand{\rhcaMsePctLgbmBarOos}{}\renewcommand{\rhcaMsePctLgbmBarOos}{25.5}
\providecommand{\rhcaMseDmLgbmBarOos}{}\renewcommand{\rhcaMseDmLgbmBarOos}{1.13}
\providecommand{\rhcaMsePLgbmBarOos}{}\renewcommand{\rhcaMsePLgbmBarOos}{\ensuremath{p = 0.26}}
\providecommand{\rhcaQXgbHarBarOos}{}\renewcommand{\rhcaQXgbHarBarOos}{0.1921}
\providecommand{\rhcaMseXgbHarBarOos}{}\renewcommand{\rhcaMseXgbHarBarOos}{2.26}
\providecommand{\rhcaMaeXgbHarBarOos}{}\renewcommand{\rhcaMaeXgbHarBarOos}{3.42}
\providecommand{\rhcaRsqXgbHarBarOos}{}\renewcommand{\rhcaRsqXgbHarBarOos}{65.6}
\providecommand{\rhcaQXgbFullBarOos}{}\renewcommand{\rhcaQXgbFullBarOos}{0.1634}
\providecommand{\rhcaMseXgbFullBarOos}{}\renewcommand{\rhcaMseXgbFullBarOos}{2.72}
\providecommand{\rhcaMaeXgbFullBarOos}{}\renewcommand{\rhcaMaeXgbFullBarOos}{3.53}
\providecommand{\rhcaRsqXgbFullBarOos}{}\renewcommand{\rhcaRsqXgbFullBarOos}{58.5}
\providecommand{\rhcaPctXgbBarOos}{}\renewcommand{\rhcaPctXgbBarOos}{\ensuremath{-}14.9}
\providecommand{\rhcaDmXgbBarOos}{}\renewcommand{\rhcaDmXgbBarOos}{\ensuremath{-}6.42}
\providecommand{\rhcaPXgbBarOos}{}\renewcommand{\rhcaPXgbBarOos}{\ensuremath{p < 0.001}}
\providecommand{\rhcaMsePctXgbBarOos}{}\renewcommand{\rhcaMsePctXgbBarOos}{20.8}
\providecommand{\rhcaMseDmXgbBarOos}{}\renewcommand{\rhcaMseDmXgbBarOos}{1.15}
\providecommand{\rhcaMsePXgbBarOos}{}\renewcommand{\rhcaMsePXgbBarOos}{\ensuremath{p = 0.25}}
\providecommand{\rhcaQRfHarBarOos}{}\renewcommand{\rhcaQRfHarBarOos}{0.1938}
\providecommand{\rhcaMseRfHarBarOos}{}\renewcommand{\rhcaMseRfHarBarOos}{2.33}
\providecommand{\rhcaMaeRfHarBarOos}{}\renewcommand{\rhcaMaeRfHarBarOos}{3.44}
\providecommand{\rhcaRsqRfHarBarOos}{}\renewcommand{\rhcaRsqRfHarBarOos}{64.5}
\providecommand{\rhcaQRfFullBarOos}{}\renewcommand{\rhcaQRfFullBarOos}{0.1730}
\providecommand{\rhcaMseRfFullBarOos}{}\renewcommand{\rhcaMseRfFullBarOos}{2.65}
\providecommand{\rhcaMaeRfFullBarOos}{}\renewcommand{\rhcaMaeRfFullBarOos}{3.52}
\providecommand{\rhcaRsqRfFullBarOos}{}\renewcommand{\rhcaRsqRfFullBarOos}{59.6}
\providecommand{\rhcaPctRfBarOos}{}\renewcommand{\rhcaPctRfBarOos}{\ensuremath{-}10.7}
\providecommand{\rhcaDmRfBarOos}{}\renewcommand{\rhcaDmRfBarOos}{\ensuremath{-}6.47}
\providecommand{\rhcaPRfBarOos}{}\renewcommand{\rhcaPRfBarOos}{\ensuremath{p < 0.001}}
\providecommand{\rhcaMsePctRfBarOos}{}\renewcommand{\rhcaMsePctRfBarOos}{13.7}
\providecommand{\rhcaMseDmRfBarOos}{}\renewcommand{\rhcaMseDmRfBarOos}{1.19}
\providecommand{\rhcaMsePRfBarOos}{}\renewcommand{\rhcaMsePRfBarOos}{\ensuremath{p = 0.23}}
% --- calibration, Pre / Oos, every bar and 16:00: Mr mean of RV/sigma2, Rm ratio of means, Under %, MzA (x 10^6), MzB, MzBse, MzRsq (levels MZ), MzlA, MzlB, MzlBse, MzlRsq (logs MZ), ShBias / ShSlope / ShDist (Theil shares of the MSE, %), Ql / Qd (QLIKE level / dispersion terms), OverSe (% of SE from over-predictions)
\providecommand{\rhcaMrRidgeHarBarPre}{}\renewcommand{\rhcaMrRidgeHarBarPre}{1.078}
\providecommand{\rhcaRmRidgeHarBarPre}{}\renewcommand{\rhcaRmRidgeHarBarPre}{1.057}
\providecommand{\rhcaUnderRidgeHarBarPre}{}\renewcommand{\rhcaUnderRidgeHarBarPre}{39.9}
\providecommand{\rhcaMzARidgeHarBarPre}{}\renewcommand{\rhcaMzARidgeHarBarPre}{0.29}
\providecommand{\rhcaMzBRidgeHarBarPre}{}\renewcommand{\rhcaMzBRidgeHarBarPre}{0.97}
\providecommand{\rhcaMzBseRidgeHarBarPre}{}\renewcommand{\rhcaMzBseRidgeHarBarPre}{0.11}
\providecommand{\rhcaMzRsqRidgeHarBarPre}{}\renewcommand{\rhcaMzRsqRidgeHarBarPre}{0.48}
\providecommand{\rhcaMzlARidgeHarBarPre}{}\renewcommand{\rhcaMzlARidgeHarBarPre}{0.04}
\providecommand{\rhcaMzlBRidgeHarBarPre}{}\renewcommand{\rhcaMzlBRidgeHarBarPre}{1.01}
\providecommand{\rhcaMzlBseRidgeHarBarPre}{}\renewcommand{\rhcaMzlBseRidgeHarBarPre}{0.01}
\providecommand{\rhcaMzlRsqRidgeHarBarPre}{}\renewcommand{\rhcaMzlRsqRidgeHarBarPre}{0.76}
\providecommand{\rhcaShBiasRidgeHarBarPre}{}\renewcommand{\rhcaShBiasRidgeHarBarPre}{0.1}
\providecommand{\rhcaShSlopeRidgeHarBarPre}{}\renewcommand{\rhcaShSlopeRidgeHarBarPre}{0.1}
\providecommand{\rhcaShDistRidgeHarBarPre}{}\renewcommand{\rhcaShDistRidgeHarBarPre}{99.8}
\providecommand{\rhcaQlRidgeHarBarPre}{}\renewcommand{\rhcaQlRidgeHarBarPre}{0.0029}
\providecommand{\rhcaQdRidgeHarBarPre}{}\renewcommand{\rhcaQdRidgeHarBarPre}{0.1835}
\providecommand{\rhcaOverSeRidgeHarBarPre}{}\renewcommand{\rhcaOverSeRidgeHarBarPre}{20.0}
\providecommand{\rhcaMrRidgeFullBarPre}{}\renewcommand{\rhcaMrRidgeFullBarPre}{1.096}
\providecommand{\rhcaRmRidgeFullBarPre}{}\renewcommand{\rhcaRmRidgeFullBarPre}{1.049}
\providecommand{\rhcaUnderRidgeFullBarPre}{}\renewcommand{\rhcaUnderRidgeFullBarPre}{41.6}
\providecommand{\rhcaMzARidgeFullBarPre}{}\renewcommand{\rhcaMzARidgeFullBarPre}{0.17}
\providecommand{\rhcaMzBRidgeFullBarPre}{}\renewcommand{\rhcaMzBRidgeFullBarPre}{1.00}
\providecommand{\rhcaMzBseRidgeFullBarPre}{}\renewcommand{\rhcaMzBseRidgeFullBarPre}{0.10}
\providecommand{\rhcaMzRsqRidgeFullBarPre}{}\renewcommand{\rhcaMzRsqRidgeFullBarPre}{0.53}
\providecommand{\rhcaMzlARidgeFullBarPre}{}\renewcommand{\rhcaMzlARidgeFullBarPre}{\ensuremath{-}0.22}
\providecommand{\rhcaMzlBRidgeFullBarPre}{}\renewcommand{\rhcaMzlBRidgeFullBarPre}{0.99}
\providecommand{\rhcaMzlBseRidgeFullBarPre}{}\renewcommand{\rhcaMzlBseRidgeFullBarPre}{0.01}
\providecommand{\rhcaMzlRsqRidgeFullBarPre}{}\renewcommand{\rhcaMzlRsqRidgeFullBarPre}{0.78}
\providecommand{\rhcaShBiasRidgeFullBarPre}{}\renewcommand{\rhcaShBiasRidgeFullBarPre}{0.1}
\providecommand{\rhcaShSlopeRidgeFullBarPre}{}\renewcommand{\rhcaShSlopeRidgeFullBarPre}{0.0}
\providecommand{\rhcaShDistRidgeFullBarPre}{}\renewcommand{\rhcaShDistRidgeFullBarPre}{99.9}
\providecommand{\rhcaQlRidgeFullBarPre}{}\renewcommand{\rhcaQlRidgeFullBarPre}{0.0043}
\providecommand{\rhcaQdRidgeFullBarPre}{}\renewcommand{\rhcaQdRidgeFullBarPre}{0.1683}
\providecommand{\rhcaOverSeRidgeFullBarPre}{}\renewcommand{\rhcaOverSeRidgeFullBarPre}{19.6}
\providecommand{\rhcaMrLassoHarBarPre}{}\renewcommand{\rhcaMrLassoHarBarPre}{1.078}
\providecommand{\rhcaRmLassoHarBarPre}{}\renewcommand{\rhcaRmLassoHarBarPre}{1.057}
\providecommand{\rhcaUnderLassoHarBarPre}{}\renewcommand{\rhcaUnderLassoHarBarPre}{39.9}
\providecommand{\rhcaMzALassoHarBarPre}{}\renewcommand{\rhcaMzALassoHarBarPre}{0.29}
\providecommand{\rhcaMzBLassoHarBarPre}{}\renewcommand{\rhcaMzBLassoHarBarPre}{0.97}
\providecommand{\rhcaMzBseLassoHarBarPre}{}\renewcommand{\rhcaMzBseLassoHarBarPre}{0.11}
\providecommand{\rhcaMzRsqLassoHarBarPre}{}\renewcommand{\rhcaMzRsqLassoHarBarPre}{0.48}
\providecommand{\rhcaMzlALassoHarBarPre}{}\renewcommand{\rhcaMzlALassoHarBarPre}{0.04}
\providecommand{\rhcaMzlBLassoHarBarPre}{}\renewcommand{\rhcaMzlBLassoHarBarPre}{1.01}
\providecommand{\rhcaMzlBseLassoHarBarPre}{}\renewcommand{\rhcaMzlBseLassoHarBarPre}{0.01}
\providecommand{\rhcaMzlRsqLassoHarBarPre}{}\renewcommand{\rhcaMzlRsqLassoHarBarPre}{0.76}
\providecommand{\rhcaShBiasLassoHarBarPre}{}\renewcommand{\rhcaShBiasLassoHarBarPre}{0.1}
\providecommand{\rhcaShSlopeLassoHarBarPre}{}\renewcommand{\rhcaShSlopeLassoHarBarPre}{0.1}
\providecommand{\rhcaShDistLassoHarBarPre}{}\renewcommand{\rhcaShDistLassoHarBarPre}{99.8}
\providecommand{\rhcaQlLassoHarBarPre}{}\renewcommand{\rhcaQlLassoHarBarPre}{0.0029}
\providecommand{\rhcaQdLassoHarBarPre}{}\renewcommand{\rhcaQdLassoHarBarPre}{0.1835}
\providecommand{\rhcaOverSeLassoHarBarPre}{}\renewcommand{\rhcaOverSeLassoHarBarPre}{20.0}
\providecommand{\rhcaMrLassoFullBarPre}{}\renewcommand{\rhcaMrLassoFullBarPre}{1.092}
\providecommand{\rhcaRmLassoFullBarPre}{}\renewcommand{\rhcaRmLassoFullBarPre}{1.045}
\providecommand{\rhcaUnderLassoFullBarPre}{}\renewcommand{\rhcaUnderLassoFullBarPre}{42.1}
\providecommand{\rhcaMzALassoFullBarPre}{}\renewcommand{\rhcaMzALassoFullBarPre}{0.16}
\providecommand{\rhcaMzBLassoFullBarPre}{}\renewcommand{\rhcaMzBLassoFullBarPre}{0.99}
\providecommand{\rhcaMzBseLassoFullBarPre}{}\renewcommand{\rhcaMzBseLassoFullBarPre}{0.10}
\providecommand{\rhcaMzRsqLassoFullBarPre}{}\renewcommand{\rhcaMzRsqLassoFullBarPre}{0.54}
\providecommand{\rhcaMzlALassoFullBarPre}{}\renewcommand{\rhcaMzlALassoFullBarPre}{\ensuremath{-}0.25}
\providecommand{\rhcaMzlBLassoFullBarPre}{}\renewcommand{\rhcaMzlBLassoFullBarPre}{0.99}
\providecommand{\rhcaMzlBseLassoFullBarPre}{}\renewcommand{\rhcaMzlBseLassoFullBarPre}{0.01}
\providecommand{\rhcaMzlRsqLassoFullBarPre}{}\renewcommand{\rhcaMzlRsqLassoFullBarPre}{0.78}
\providecommand{\rhcaShBiasLassoFullBarPre}{}\renewcommand{\rhcaShBiasLassoFullBarPre}{0.1}
\providecommand{\rhcaShSlopeLassoFullBarPre}{}\renewcommand{\rhcaShSlopeLassoFullBarPre}{0.0}
\providecommand{\rhcaShDistLassoFullBarPre}{}\renewcommand{\rhcaShDistLassoFullBarPre}{99.9}
\providecommand{\rhcaQlLassoFullBarPre}{}\renewcommand{\rhcaQlLassoFullBarPre}{0.0040}
\providecommand{\rhcaQdLassoFullBarPre}{}\renewcommand{\rhcaQdLassoFullBarPre}{0.1627}
\providecommand{\rhcaOverSeLassoFullBarPre}{}\renewcommand{\rhcaOverSeLassoFullBarPre}{20.1}
\providecommand{\rhcaMrEnetHarBarPre}{}\renewcommand{\rhcaMrEnetHarBarPre}{1.078}
\providecommand{\rhcaRmEnetHarBarPre}{}\renewcommand{\rhcaRmEnetHarBarPre}{1.057}
\providecommand{\rhcaUnderEnetHarBarPre}{}\renewcommand{\rhcaUnderEnetHarBarPre}{39.9}
\providecommand{\rhcaMzAEnetHarBarPre}{}\renewcommand{\rhcaMzAEnetHarBarPre}{0.29}
\providecommand{\rhcaMzBEnetHarBarPre}{}\renewcommand{\rhcaMzBEnetHarBarPre}{0.97}
\providecommand{\rhcaMzBseEnetHarBarPre}{}\renewcommand{\rhcaMzBseEnetHarBarPre}{0.11}
\providecommand{\rhcaMzRsqEnetHarBarPre}{}\renewcommand{\rhcaMzRsqEnetHarBarPre}{0.48}
\providecommand{\rhcaMzlAEnetHarBarPre}{}\renewcommand{\rhcaMzlAEnetHarBarPre}{0.04}
\providecommand{\rhcaMzlBEnetHarBarPre}{}\renewcommand{\rhcaMzlBEnetHarBarPre}{1.01}
\providecommand{\rhcaMzlBseEnetHarBarPre}{}\renewcommand{\rhcaMzlBseEnetHarBarPre}{0.01}
\providecommand{\rhcaMzlRsqEnetHarBarPre}{}\renewcommand{\rhcaMzlRsqEnetHarBarPre}{0.76}
\providecommand{\rhcaShBiasEnetHarBarPre}{}\renewcommand{\rhcaShBiasEnetHarBarPre}{0.1}
\providecommand{\rhcaShSlopeEnetHarBarPre}{}\renewcommand{\rhcaShSlopeEnetHarBarPre}{0.1}
\providecommand{\rhcaShDistEnetHarBarPre}{}\renewcommand{\rhcaShDistEnetHarBarPre}{99.8}
\providecommand{\rhcaQlEnetHarBarPre}{}\renewcommand{\rhcaQlEnetHarBarPre}{0.0029}
\providecommand{\rhcaQdEnetHarBarPre}{}\renewcommand{\rhcaQdEnetHarBarPre}{0.1835}
\providecommand{\rhcaOverSeEnetHarBarPre}{}\renewcommand{\rhcaOverSeEnetHarBarPre}{20.0}
\providecommand{\rhcaMrEnetFullBarPre}{}\renewcommand{\rhcaMrEnetFullBarPre}{1.091}
\providecommand{\rhcaRmEnetFullBarPre}{}\renewcommand{\rhcaRmEnetFullBarPre}{1.049}
\providecommand{\rhcaUnderEnetFullBarPre}{}\renewcommand{\rhcaUnderEnetFullBarPre}{41.5}
\providecommand{\rhcaMzAEnetFullBarPre}{}\renewcommand{\rhcaMzAEnetFullBarPre}{0.16}
\providecommand{\rhcaMzBEnetFullBarPre}{}\renewcommand{\rhcaMzBEnetFullBarPre}{1.00}
\providecommand{\rhcaMzBseEnetFullBarPre}{}\renewcommand{\rhcaMzBseEnetFullBarPre}{0.10}
\providecommand{\rhcaMzRsqEnetFullBarPre}{}\renewcommand{\rhcaMzRsqEnetFullBarPre}{0.53}
\providecommand{\rhcaMzlAEnetFullBarPre}{}\renewcommand{\rhcaMzlAEnetFullBarPre}{\ensuremath{-}0.16}
\providecommand{\rhcaMzlBEnetFullBarPre}{}\renewcommand{\rhcaMzlBEnetFullBarPre}{0.99}
\providecommand{\rhcaMzlBseEnetFullBarPre}{}\renewcommand{\rhcaMzlBseEnetFullBarPre}{0.01}
\providecommand{\rhcaMzlRsqEnetFullBarPre}{}\renewcommand{\rhcaMzlRsqEnetFullBarPre}{0.78}
\providecommand{\rhcaShBiasEnetFullBarPre}{}\renewcommand{\rhcaShBiasEnetFullBarPre}{0.1}
\providecommand{\rhcaShSlopeEnetFullBarPre}{}\renewcommand{\rhcaShSlopeEnetFullBarPre}{0.0}
\providecommand{\rhcaShDistEnetFullBarPre}{}\renewcommand{\rhcaShDistEnetFullBarPre}{99.9}
\providecommand{\rhcaQlEnetFullBarPre}{}\renewcommand{\rhcaQlEnetFullBarPre}{0.0039}
\providecommand{\rhcaQdEnetFullBarPre}{}\renewcommand{\rhcaQdEnetFullBarPre}{0.1681}
\providecommand{\rhcaOverSeEnetFullBarPre}{}\renewcommand{\rhcaOverSeEnetFullBarPre}{19.7}
\providecommand{\rhcaMrLgbmHarBarPre}{}\renewcommand{\rhcaMrLgbmHarBarPre}{1.083}
\providecommand{\rhcaRmLgbmHarBarPre}{}\renewcommand{\rhcaRmLgbmHarBarPre}{1.075}
\providecommand{\rhcaUnderLgbmHarBarPre}{}\renewcommand{\rhcaUnderLgbmHarBarPre}{40.2}
\providecommand{\rhcaMzALgbmHarBarPre}{}\renewcommand{\rhcaMzALgbmHarBarPre}{0.27}
\providecommand{\rhcaMzBLgbmHarBarPre}{}\renewcommand{\rhcaMzBLgbmHarBarPre}{0.99}
\providecommand{\rhcaMzBseLgbmHarBarPre}{}\renewcommand{\rhcaMzBseLgbmHarBarPre}{0.07}
\providecommand{\rhcaMzRsqLgbmHarBarPre}{}\renewcommand{\rhcaMzRsqLgbmHarBarPre}{0.37}
\providecommand{\rhcaMzlALgbmHarBarPre}{}\renewcommand{\rhcaMzlALgbmHarBarPre}{0.04}
\providecommand{\rhcaMzlBLgbmHarBarPre}{}\renewcommand{\rhcaMzlBLgbmHarBarPre}{1.01}
\providecommand{\rhcaMzlBseLgbmHarBarPre}{}\renewcommand{\rhcaMzlBseLgbmHarBarPre}{0.01}
\providecommand{\rhcaMzlRsqLgbmHarBarPre}{}\renewcommand{\rhcaMzlRsqLgbmHarBarPre}{0.76}
\providecommand{\rhcaShBiasLgbmHarBarPre}{}\renewcommand{\rhcaShBiasLgbmHarBarPre}{0.1}
\providecommand{\rhcaShSlopeLgbmHarBarPre}{}\renewcommand{\rhcaShSlopeLgbmHarBarPre}{0.0}
\providecommand{\rhcaShDistLgbmHarBarPre}{}\renewcommand{\rhcaShDistLgbmHarBarPre}{99.8}
\providecommand{\rhcaQlLgbmHarBarPre}{}\renewcommand{\rhcaQlLgbmHarBarPre}{0.0032}
\providecommand{\rhcaQdLgbmHarBarPre}{}\renewcommand{\rhcaQdLgbmHarBarPre}{0.1826}
\providecommand{\rhcaOverSeLgbmHarBarPre}{}\renewcommand{\rhcaOverSeLgbmHarBarPre}{10.3}
\providecommand{\rhcaMrLgbmFullBarPre}{}\renewcommand{\rhcaMrLgbmFullBarPre}{1.069}
\providecommand{\rhcaRmLgbmFullBarPre}{}\renewcommand{\rhcaRmLgbmFullBarPre}{1.060}
\providecommand{\rhcaUnderLgbmFullBarPre}{}\renewcommand{\rhcaUnderLgbmFullBarPre}{41.0}
\providecommand{\rhcaMzALgbmFullBarPre}{}\renewcommand{\rhcaMzALgbmFullBarPre}{0.29}
\providecommand{\rhcaMzBLgbmFullBarPre}{}\renewcommand{\rhcaMzBLgbmFullBarPre}{0.97}
\providecommand{\rhcaMzBseLgbmFullBarPre}{}\renewcommand{\rhcaMzBseLgbmFullBarPre}{0.06}
\providecommand{\rhcaMzRsqLgbmFullBarPre}{}\renewcommand{\rhcaMzRsqLgbmFullBarPre}{0.40}
\providecommand{\rhcaMzlALgbmFullBarPre}{}\renewcommand{\rhcaMzlALgbmFullBarPre}{\ensuremath{-}0.03}
\providecommand{\rhcaMzlBLgbmFullBarPre}{}\renewcommand{\rhcaMzlBLgbmFullBarPre}{1.00}
\providecommand{\rhcaMzlBseLgbmFullBarPre}{}\renewcommand{\rhcaMzlBseLgbmFullBarPre}{0.01}
\providecommand{\rhcaMzlRsqLgbmFullBarPre}{}\renewcommand{\rhcaMzlRsqLgbmFullBarPre}{0.79}
\providecommand{\rhcaShBiasLgbmFullBarPre}{}\renewcommand{\rhcaShBiasLgbmFullBarPre}{0.1}
\providecommand{\rhcaShSlopeLgbmFullBarPre}{}\renewcommand{\rhcaShSlopeLgbmFullBarPre}{0.1}
\providecommand{\rhcaShDistLgbmFullBarPre}{}\renewcommand{\rhcaShDistLgbmFullBarPre}{99.8}
\providecommand{\rhcaQlLgbmFullBarPre}{}\renewcommand{\rhcaQlLgbmFullBarPre}{0.0023}
\providecommand{\rhcaQdLgbmFullBarPre}{}\renewcommand{\rhcaQdLgbmFullBarPre}{0.1563}
\providecommand{\rhcaOverSeLgbmFullBarPre}{}\renewcommand{\rhcaOverSeLgbmFullBarPre}{12.0}
\providecommand{\rhcaMrXgbHarBarPre}{}\renewcommand{\rhcaMrXgbHarBarPre}{1.083}
\providecommand{\rhcaRmXgbHarBarPre}{}\renewcommand{\rhcaRmXgbHarBarPre}{1.073}
\providecommand{\rhcaUnderXgbHarBarPre}{}\renewcommand{\rhcaUnderXgbHarBarPre}{40.2}
\providecommand{\rhcaMzAXgbHarBarPre}{}\renewcommand{\rhcaMzAXgbHarBarPre}{0.31}
\providecommand{\rhcaMzBXgbHarBarPre}{}\renewcommand{\rhcaMzBXgbHarBarPre}{0.97}
\providecommand{\rhcaMzBseXgbHarBarPre}{}\renewcommand{\rhcaMzBseXgbHarBarPre}{0.09}
\providecommand{\rhcaMzRsqXgbHarBarPre}{}\renewcommand{\rhcaMzRsqXgbHarBarPre}{0.38}
\providecommand{\rhcaMzlAXgbHarBarPre}{}\renewcommand{\rhcaMzlAXgbHarBarPre}{0.06}
\providecommand{\rhcaMzlBXgbHarBarPre}{}\renewcommand{\rhcaMzlBXgbHarBarPre}{1.01}
\providecommand{\rhcaMzlBseXgbHarBarPre}{}\renewcommand{\rhcaMzlBseXgbHarBarPre}{0.01}
\providecommand{\rhcaMzlRsqXgbHarBarPre}{}\renewcommand{\rhcaMzlRsqXgbHarBarPre}{0.76}
\providecommand{\rhcaShBiasXgbHarBarPre}{}\renewcommand{\rhcaShBiasXgbHarBarPre}{0.1}
\providecommand{\rhcaShSlopeXgbHarBarPre}{}\renewcommand{\rhcaShSlopeXgbHarBarPre}{0.0}
\providecommand{\rhcaShDistXgbHarBarPre}{}\renewcommand{\rhcaShDistXgbHarBarPre}{99.8}
\providecommand{\rhcaQlXgbHarBarPre}{}\renewcommand{\rhcaQlXgbHarBarPre}{0.0032}
\providecommand{\rhcaQdXgbHarBarPre}{}\renewcommand{\rhcaQdXgbHarBarPre}{0.1832}
\providecommand{\rhcaOverSeXgbHarBarPre}{}\renewcommand{\rhcaOverSeXgbHarBarPre}{12.2}
\providecommand{\rhcaMrXgbFullBarPre}{}\renewcommand{\rhcaMrXgbFullBarPre}{1.067}
\providecommand{\rhcaRmXgbFullBarPre}{}\renewcommand{\rhcaRmXgbFullBarPre}{1.055}
\providecommand{\rhcaUnderXgbFullBarPre}{}\renewcommand{\rhcaUnderXgbFullBarPre}{40.8}
\providecommand{\rhcaMzAXgbFullBarPre}{}\renewcommand{\rhcaMzAXgbFullBarPre}{0.40}
\providecommand{\rhcaMzBXgbFullBarPre}{}\renewcommand{\rhcaMzBXgbFullBarPre}{0.93}
\providecommand{\rhcaMzBseXgbFullBarPre}{}\renewcommand{\rhcaMzBseXgbFullBarPre}{0.09}
\providecommand{\rhcaMzRsqXgbFullBarPre}{}\renewcommand{\rhcaMzRsqXgbFullBarPre}{0.40}
\providecommand{\rhcaMzlAXgbFullBarPre}{}\renewcommand{\rhcaMzlAXgbFullBarPre}{0.00}
\providecommand{\rhcaMzlBXgbFullBarPre}{}\renewcommand{\rhcaMzlBXgbFullBarPre}{1.01}
\providecommand{\rhcaMzlBseXgbFullBarPre}{}\renewcommand{\rhcaMzlBseXgbFullBarPre}{0.01}
\providecommand{\rhcaMzlRsqXgbFullBarPre}{}\renewcommand{\rhcaMzlRsqXgbFullBarPre}{0.79}
\providecommand{\rhcaShBiasXgbFullBarPre}{}\renewcommand{\rhcaShBiasXgbFullBarPre}{0.1}
\providecommand{\rhcaShSlopeXgbFullBarPre}{}\renewcommand{\rhcaShSlopeXgbFullBarPre}{0.4}
\providecommand{\rhcaShDistXgbFullBarPre}{}\renewcommand{\rhcaShDistXgbFullBarPre}{99.5}
\providecommand{\rhcaQlXgbFullBarPre}{}\renewcommand{\rhcaQlXgbFullBarPre}{0.0021}
\providecommand{\rhcaQdXgbFullBarPre}{}\renewcommand{\rhcaQdXgbFullBarPre}{0.1575}
\providecommand{\rhcaOverSeXgbFullBarPre}{}\renewcommand{\rhcaOverSeXgbFullBarPre}{15.1}
\providecommand{\rhcaMrRfHarBarPre}{}\renewcommand{\rhcaMrRfHarBarPre}{1.085}
\providecommand{\rhcaRmRfHarBarPre}{}\renewcommand{\rhcaRmRfHarBarPre}{1.068}
\providecommand{\rhcaUnderRfHarBarPre}{}\renewcommand{\rhcaUnderRfHarBarPre}{40.1}
\providecommand{\rhcaMzARfHarBarPre}{}\renewcommand{\rhcaMzARfHarBarPre}{0.33}
\providecommand{\rhcaMzBRfHarBarPre}{}\renewcommand{\rhcaMzBRfHarBarPre}{0.96}
\providecommand{\rhcaMzBseRfHarBarPre}{}\renewcommand{\rhcaMzBseRfHarBarPre}{0.10}
\providecommand{\rhcaMzRsqRfHarBarPre}{}\renewcommand{\rhcaMzRsqRfHarBarPre}{0.38}
\providecommand{\rhcaMzlARfHarBarPre}{}\renewcommand{\rhcaMzlARfHarBarPre}{\ensuremath{-}0.04}
\providecommand{\rhcaMzlBRfHarBarPre}{}\renewcommand{\rhcaMzlBRfHarBarPre}{1.00}
\providecommand{\rhcaMzlBseRfHarBarPre}{}\renewcommand{\rhcaMzlBseRfHarBarPre}{0.01}
\providecommand{\rhcaMzlRsqRfHarBarPre}{}\renewcommand{\rhcaMzlRsqRfHarBarPre}{0.76}
\providecommand{\rhcaShBiasRfHarBarPre}{}\renewcommand{\rhcaShBiasRfHarBarPre}{0.1}
\providecommand{\rhcaShSlopeRfHarBarPre}{}\renewcommand{\rhcaShSlopeRfHarBarPre}{0.1}
\providecommand{\rhcaShDistRfHarBarPre}{}\renewcommand{\rhcaShDistRfHarBarPre}{99.8}
\providecommand{\rhcaQlRfHarBarPre}{}\renewcommand{\rhcaQlRfHarBarPre}{0.0034}
\providecommand{\rhcaQdRfHarBarPre}{}\renewcommand{\rhcaQdRfHarBarPre}{0.1831}
\providecommand{\rhcaOverSeRfHarBarPre}{}\renewcommand{\rhcaOverSeRfHarBarPre}{12.8}
\providecommand{\rhcaMrRfFullBarPre}{}\renewcommand{\rhcaMrRfFullBarPre}{1.078}
\providecommand{\rhcaRmRfFullBarPre}{}\renewcommand{\rhcaRmRfFullBarPre}{1.057}
\providecommand{\rhcaUnderRfFullBarPre}{}\renewcommand{\rhcaUnderRfFullBarPre}{41.1}
\providecommand{\rhcaMzARfFullBarPre}{}\renewcommand{\rhcaMzARfFullBarPre}{0.43}
\providecommand{\rhcaMzBRfFullBarPre}{}\renewcommand{\rhcaMzBRfFullBarPre}{0.92}
\providecommand{\rhcaMzBseRfFullBarPre}{}\renewcommand{\rhcaMzBseRfFullBarPre}{0.09}
\providecommand{\rhcaMzRsqRfFullBarPre}{}\renewcommand{\rhcaMzRsqRfFullBarPre}{0.39}
\providecommand{\rhcaMzlARfFullBarPre}{}\renewcommand{\rhcaMzlARfFullBarPre}{\ensuremath{-}0.11}
\providecommand{\rhcaMzlBRfFullBarPre}{}\renewcommand{\rhcaMzlBRfFullBarPre}{1.00}
\providecommand{\rhcaMzlBseRfFullBarPre}{}\renewcommand{\rhcaMzlBseRfFullBarPre}{0.01}
\providecommand{\rhcaMzlRsqRfFullBarPre}{}\renewcommand{\rhcaMzlRsqRfFullBarPre}{0.79}
\providecommand{\rhcaShBiasRfFullBarPre}{}\renewcommand{\rhcaShBiasRfFullBarPre}{0.1}
\providecommand{\rhcaShSlopeRfFullBarPre}{}\renewcommand{\rhcaShSlopeRfFullBarPre}{0.5}
\providecommand{\rhcaShDistRfFullBarPre}{}\renewcommand{\rhcaShDistRfFullBarPre}{99.5}
\providecommand{\rhcaQlRfFullBarPre}{}\renewcommand{\rhcaQlRfFullBarPre}{0.0029}
\providecommand{\rhcaQdRfFullBarPre}{}\renewcommand{\rhcaQdRfFullBarPre}{0.1617}
\providecommand{\rhcaOverSeRfFullBarPre}{}\renewcommand{\rhcaOverSeRfFullBarPre}{14.7}
\providecommand{\rhcaMrRidgeHarClosePre}{}\renewcommand{\rhcaMrRidgeHarClosePre}{1.121}
\providecommand{\rhcaRmRidgeHarClosePre}{}\renewcommand{\rhcaRmRidgeHarClosePre}{1.097}
\providecommand{\rhcaUnderRidgeHarClosePre}{}\renewcommand{\rhcaUnderRidgeHarClosePre}{46.7}
\providecommand{\rhcaMzARidgeHarClosePre}{}\renewcommand{\rhcaMzARidgeHarClosePre}{\ensuremath{-}0.16}
\providecommand{\rhcaMzBRidgeHarClosePre}{}\renewcommand{\rhcaMzBRidgeHarClosePre}{1.14}
\providecommand{\rhcaMzBseRidgeHarClosePre}{}\renewcommand{\rhcaMzBseRidgeHarClosePre}{0.10}
\providecommand{\rhcaMzRsqRidgeHarClosePre}{}\renewcommand{\rhcaMzRsqRidgeHarClosePre}{0.64}
\providecommand{\rhcaMzlARidgeHarClosePre}{}\renewcommand{\rhcaMzlARidgeHarClosePre}{\ensuremath{-}0.68}
\providecommand{\rhcaMzlBRidgeHarClosePre}{}\renewcommand{\rhcaMzlBRidgeHarClosePre}{0.95}
\providecommand{\rhcaMzlBseRidgeHarClosePre}{}\renewcommand{\rhcaMzlBseRidgeHarClosePre}{0.02}
\providecommand{\rhcaMzlRsqRidgeHarClosePre}{}\renewcommand{\rhcaMzlRsqRidgeHarClosePre}{0.78}
\providecommand{\rhcaShBiasRidgeHarClosePre}{}\renewcommand{\rhcaShBiasRidgeHarClosePre}{0.5}
\providecommand{\rhcaShSlopeRidgeHarClosePre}{}\renewcommand{\rhcaShSlopeRidgeHarClosePre}{2.5}
\providecommand{\rhcaShDistRidgeHarClosePre}{}\renewcommand{\rhcaShDistRidgeHarClosePre}{97.0}
\providecommand{\rhcaQlRidgeHarClosePre}{}\renewcommand{\rhcaQlRidgeHarClosePre}{0.0068}
\providecommand{\rhcaQdRidgeHarClosePre}{}\renewcommand{\rhcaQdRidgeHarClosePre}{0.1297}
\providecommand{\rhcaOverSeRidgeHarClosePre}{}\renewcommand{\rhcaOverSeRidgeHarClosePre}{13.1}
\providecommand{\rhcaMrRidgeFullClosePre}{}\renewcommand{\rhcaMrRidgeFullClosePre}{1.137}
\providecommand{\rhcaRmRidgeFullClosePre}{}\renewcommand{\rhcaRmRidgeFullClosePre}{1.061}
\providecommand{\rhcaUnderRidgeFullClosePre}{}\renewcommand{\rhcaUnderRidgeFullClosePre}{47.2}
\providecommand{\rhcaMzARidgeFullClosePre}{}\renewcommand{\rhcaMzARidgeFullClosePre}{\ensuremath{-}0.02}
\providecommand{\rhcaMzBRidgeFullClosePre}{}\renewcommand{\rhcaMzBRidgeFullClosePre}{1.07}
\providecommand{\rhcaMzBseRidgeFullClosePre}{}\renewcommand{\rhcaMzBseRidgeFullClosePre}{0.09}
\providecommand{\rhcaMzRsqRidgeFullClosePre}{}\renewcommand{\rhcaMzRsqRidgeFullClosePre}{0.68}
\providecommand{\rhcaMzlARidgeFullClosePre}{}\renewcommand{\rhcaMzlARidgeFullClosePre}{\ensuremath{-}0.93}
\providecommand{\rhcaMzlBRidgeFullClosePre}{}\renewcommand{\rhcaMzlBRidgeFullClosePre}{0.93}
\providecommand{\rhcaMzlBseRidgeFullClosePre}{}\renewcommand{\rhcaMzlBseRidgeFullClosePre}{0.02}
\providecommand{\rhcaMzlRsqRidgeFullClosePre}{}\renewcommand{\rhcaMzlRsqRidgeFullClosePre}{0.80}
\providecommand{\rhcaShBiasRidgeFullClosePre}{}\renewcommand{\rhcaShBiasRidgeFullClosePre}{0.2}
\providecommand{\rhcaShSlopeRidgeFullClosePre}{}\renewcommand{\rhcaShSlopeRidgeFullClosePre}{0.8}
\providecommand{\rhcaShDistRidgeFullClosePre}{}\renewcommand{\rhcaShDistRidgeFullClosePre}{99.0}
\providecommand{\rhcaQlRidgeFullClosePre}{}\renewcommand{\rhcaQlRidgeFullClosePre}{0.0086}
\providecommand{\rhcaQdRidgeFullClosePre}{}\renewcommand{\rhcaQdRidgeFullClosePre}{0.1222}
\providecommand{\rhcaOverSeRidgeFullClosePre}{}\renewcommand{\rhcaOverSeRidgeFullClosePre}{18.0}
\providecommand{\rhcaMrLassoHarClosePre}{}\renewcommand{\rhcaMrLassoHarClosePre}{1.121}
\providecommand{\rhcaRmLassoHarClosePre}{}\renewcommand{\rhcaRmLassoHarClosePre}{1.097}
\providecommand{\rhcaUnderLassoHarClosePre}{}\renewcommand{\rhcaUnderLassoHarClosePre}{46.7}
\providecommand{\rhcaMzALassoHarClosePre}{}\renewcommand{\rhcaMzALassoHarClosePre}{\ensuremath{-}0.16}
\providecommand{\rhcaMzBLassoHarClosePre}{}\renewcommand{\rhcaMzBLassoHarClosePre}{1.14}
\providecommand{\rhcaMzBseLassoHarClosePre}{}\renewcommand{\rhcaMzBseLassoHarClosePre}{0.10}
\providecommand{\rhcaMzRsqLassoHarClosePre}{}\renewcommand{\rhcaMzRsqLassoHarClosePre}{0.64}
\providecommand{\rhcaMzlALassoHarClosePre}{}\renewcommand{\rhcaMzlALassoHarClosePre}{\ensuremath{-}0.68}
\providecommand{\rhcaMzlBLassoHarClosePre}{}\renewcommand{\rhcaMzlBLassoHarClosePre}{0.95}
\providecommand{\rhcaMzlBseLassoHarClosePre}{}\renewcommand{\rhcaMzlBseLassoHarClosePre}{0.02}
\providecommand{\rhcaMzlRsqLassoHarClosePre}{}\renewcommand{\rhcaMzlRsqLassoHarClosePre}{0.78}
\providecommand{\rhcaShBiasLassoHarClosePre}{}\renewcommand{\rhcaShBiasLassoHarClosePre}{0.5}
\providecommand{\rhcaShSlopeLassoHarClosePre}{}\renewcommand{\rhcaShSlopeLassoHarClosePre}{2.5}
\providecommand{\rhcaShDistLassoHarClosePre}{}\renewcommand{\rhcaShDistLassoHarClosePre}{97.0}
\providecommand{\rhcaQlLassoHarClosePre}{}\renewcommand{\rhcaQlLassoHarClosePre}{0.0068}
\providecommand{\rhcaQdLassoHarClosePre}{}\renewcommand{\rhcaQdLassoHarClosePre}{0.1297}
\providecommand{\rhcaOverSeLassoHarClosePre}{}\renewcommand{\rhcaOverSeLassoHarClosePre}{13.1}
\providecommand{\rhcaMrLassoFullClosePre}{}\renewcommand{\rhcaMrLassoFullClosePre}{1.148}
\providecommand{\rhcaRmLassoFullClosePre}{}\renewcommand{\rhcaRmLassoFullClosePre}{1.059}
\providecommand{\rhcaUnderLassoFullClosePre}{}\renewcommand{\rhcaUnderLassoFullClosePre}{48.6}
\providecommand{\rhcaMzALassoFullClosePre}{}\renewcommand{\rhcaMzALassoFullClosePre}{\ensuremath{-}0.01}
\providecommand{\rhcaMzBLassoFullClosePre}{}\renewcommand{\rhcaMzBLassoFullClosePre}{1.06}
\providecommand{\rhcaMzBseLassoFullClosePre}{}\renewcommand{\rhcaMzBseLassoFullClosePre}{0.09}
\providecommand{\rhcaMzRsqLassoFullClosePre}{}\renewcommand{\rhcaMzRsqLassoFullClosePre}{0.69}
\providecommand{\rhcaMzlALassoFullClosePre}{}\renewcommand{\rhcaMzlALassoFullClosePre}{\ensuremath{-}1.01}
\providecommand{\rhcaMzlBLassoFullClosePre}{}\renewcommand{\rhcaMzlBLassoFullClosePre}{0.92}
\providecommand{\rhcaMzlBseLassoFullClosePre}{}\renewcommand{\rhcaMzlBseLassoFullClosePre}{0.02}
\providecommand{\rhcaMzlRsqLassoFullClosePre}{}\renewcommand{\rhcaMzlRsqLassoFullClosePre}{0.80}
\providecommand{\rhcaShBiasLassoFullClosePre}{}\renewcommand{\rhcaShBiasLassoFullClosePre}{0.2}
\providecommand{\rhcaShSlopeLassoFullClosePre}{}\renewcommand{\rhcaShSlopeLassoFullClosePre}{0.7}
\providecommand{\rhcaShDistLassoFullClosePre}{}\renewcommand{\rhcaShDistLassoFullClosePre}{99.1}
\providecommand{\rhcaQlLassoFullClosePre}{}\renewcommand{\rhcaQlLassoFullClosePre}{0.0100}
\providecommand{\rhcaQdLassoFullClosePre}{}\renewcommand{\rhcaQdLassoFullClosePre}{0.1233}
\providecommand{\rhcaOverSeLassoFullClosePre}{}\renewcommand{\rhcaOverSeLassoFullClosePre}{18.4}
\providecommand{\rhcaMrEnetHarClosePre}{}\renewcommand{\rhcaMrEnetHarClosePre}{1.121}
\providecommand{\rhcaRmEnetHarClosePre}{}\renewcommand{\rhcaRmEnetHarClosePre}{1.097}
\providecommand{\rhcaUnderEnetHarClosePre}{}\renewcommand{\rhcaUnderEnetHarClosePre}{46.7}
\providecommand{\rhcaMzAEnetHarClosePre}{}\renewcommand{\rhcaMzAEnetHarClosePre}{\ensuremath{-}0.16}
\providecommand{\rhcaMzBEnetHarClosePre}{}\renewcommand{\rhcaMzBEnetHarClosePre}{1.14}
\providecommand{\rhcaMzBseEnetHarClosePre}{}\renewcommand{\rhcaMzBseEnetHarClosePre}{0.10}
\providecommand{\rhcaMzRsqEnetHarClosePre}{}\renewcommand{\rhcaMzRsqEnetHarClosePre}{0.64}
\providecommand{\rhcaMzlAEnetHarClosePre}{}\renewcommand{\rhcaMzlAEnetHarClosePre}{\ensuremath{-}0.68}
\providecommand{\rhcaMzlBEnetHarClosePre}{}\renewcommand{\rhcaMzlBEnetHarClosePre}{0.95}
\providecommand{\rhcaMzlBseEnetHarClosePre}{}\renewcommand{\rhcaMzlBseEnetHarClosePre}{0.02}
\providecommand{\rhcaMzlRsqEnetHarClosePre}{}\renewcommand{\rhcaMzlRsqEnetHarClosePre}{0.78}
\providecommand{\rhcaShBiasEnetHarClosePre}{}\renewcommand{\rhcaShBiasEnetHarClosePre}{0.5}
\providecommand{\rhcaShSlopeEnetHarClosePre}{}\renewcommand{\rhcaShSlopeEnetHarClosePre}{2.5}
\providecommand{\rhcaShDistEnetHarClosePre}{}\renewcommand{\rhcaShDistEnetHarClosePre}{97.0}
\providecommand{\rhcaQlEnetHarClosePre}{}\renewcommand{\rhcaQlEnetHarClosePre}{0.0068}
\providecommand{\rhcaQdEnetHarClosePre}{}\renewcommand{\rhcaQdEnetHarClosePre}{0.1297}
\providecommand{\rhcaOverSeEnetHarClosePre}{}\renewcommand{\rhcaOverSeEnetHarClosePre}{13.1}
\providecommand{\rhcaMrEnetFullClosePre}{}\renewcommand{\rhcaMrEnetFullClosePre}{1.138}
\providecommand{\rhcaRmEnetFullClosePre}{}\renewcommand{\rhcaRmEnetFullClosePre}{1.067}
\providecommand{\rhcaUnderEnetFullClosePre}{}\renewcommand{\rhcaUnderEnetFullClosePre}{47.3}
\providecommand{\rhcaMzAEnetFullClosePre}{}\renewcommand{\rhcaMzAEnetFullClosePre}{\ensuremath{-}0.05}
\providecommand{\rhcaMzBEnetFullClosePre}{}\renewcommand{\rhcaMzBEnetFullClosePre}{1.08}
\providecommand{\rhcaMzBseEnetFullClosePre}{}\renewcommand{\rhcaMzBseEnetFullClosePre}{0.09}
\providecommand{\rhcaMzRsqEnetFullClosePre}{}\renewcommand{\rhcaMzRsqEnetFullClosePre}{0.69}
\providecommand{\rhcaMzlAEnetFullClosePre}{}\renewcommand{\rhcaMzlAEnetFullClosePre}{\ensuremath{-}0.90}
\providecommand{\rhcaMzlBEnetFullClosePre}{}\renewcommand{\rhcaMzlBEnetFullClosePre}{0.93}
\providecommand{\rhcaMzlBseEnetFullClosePre}{}\renewcommand{\rhcaMzlBseEnetFullClosePre}{0.02}
\providecommand{\rhcaMzlRsqEnetFullClosePre}{}\renewcommand{\rhcaMzlRsqEnetFullClosePre}{0.80}
\providecommand{\rhcaShBiasEnetFullClosePre}{}\renewcommand{\rhcaShBiasEnetFullClosePre}{0.3}
\providecommand{\rhcaShSlopeEnetFullClosePre}{}\renewcommand{\rhcaShSlopeEnetFullClosePre}{1.2}
\providecommand{\rhcaShDistEnetFullClosePre}{}\renewcommand{\rhcaShDistEnetFullClosePre}{98.5}
\providecommand{\rhcaQlEnetFullClosePre}{}\renewcommand{\rhcaQlEnetFullClosePre}{0.0087}
\providecommand{\rhcaQdEnetFullClosePre}{}\renewcommand{\rhcaQdEnetFullClosePre}{0.1213}
\providecommand{\rhcaOverSeEnetFullClosePre}{}\renewcommand{\rhcaOverSeEnetFullClosePre}{17.3}
\providecommand{\rhcaMrLgbmHarClosePre}{}\renewcommand{\rhcaMrLgbmHarClosePre}{1.068}
\providecommand{\rhcaRmLgbmHarClosePre}{}\renewcommand{\rhcaRmLgbmHarClosePre}{1.062}
\providecommand{\rhcaUnderLgbmHarClosePre}{}\renewcommand{\rhcaUnderLgbmHarClosePre}{42.9}
\providecommand{\rhcaMzALgbmHarClosePre}{}\renewcommand{\rhcaMzALgbmHarClosePre}{0.00}
\providecommand{\rhcaMzBLgbmHarClosePre}{}\renewcommand{\rhcaMzBLgbmHarClosePre}{1.06}
\providecommand{\rhcaMzBseLgbmHarClosePre}{}\renewcommand{\rhcaMzBseLgbmHarClosePre}{0.10}
\providecommand{\rhcaMzRsqLgbmHarClosePre}{}\renewcommand{\rhcaMzRsqLgbmHarClosePre}{0.55}
\providecommand{\rhcaMzlALgbmHarClosePre}{}\renewcommand{\rhcaMzlALgbmHarClosePre}{\ensuremath{-}0.59}
\providecommand{\rhcaMzlBLgbmHarClosePre}{}\renewcommand{\rhcaMzlBLgbmHarClosePre}{0.96}
\providecommand{\rhcaMzlBseLgbmHarClosePre}{}\renewcommand{\rhcaMzlBseLgbmHarClosePre}{0.02}
\providecommand{\rhcaMzlRsqLgbmHarClosePre}{}\renewcommand{\rhcaMzlRsqLgbmHarClosePre}{0.78}
\providecommand{\rhcaShBiasLgbmHarClosePre}{}\renewcommand{\rhcaShBiasLgbmHarClosePre}{0.2}
\providecommand{\rhcaShSlopeLgbmHarClosePre}{}\renewcommand{\rhcaShSlopeLgbmHarClosePre}{0.4}
\providecommand{\rhcaShDistLgbmHarClosePre}{}\renewcommand{\rhcaShDistLgbmHarClosePre}{99.4}
\providecommand{\rhcaQlLgbmHarClosePre}{}\renewcommand{\rhcaQlLgbmHarClosePre}{0.0022}
\providecommand{\rhcaQdLgbmHarClosePre}{}\renewcommand{\rhcaQdLgbmHarClosePre}{0.1310}
\providecommand{\rhcaOverSeLgbmHarClosePre}{}\renewcommand{\rhcaOverSeLgbmHarClosePre}{14.4}
\providecommand{\rhcaMrLgbmFullClosePre}{}\renewcommand{\rhcaMrLgbmFullClosePre}{1.101}
\providecommand{\rhcaRmLgbmFullClosePre}{}\renewcommand{\rhcaRmLgbmFullClosePre}{1.042}
\providecommand{\rhcaUnderLgbmFullClosePre}{}\renewcommand{\rhcaUnderLgbmFullClosePre}{45.9}
\providecommand{\rhcaMzALgbmFullClosePre}{}\renewcommand{\rhcaMzALgbmFullClosePre}{0.26}
\providecommand{\rhcaMzBLgbmFullClosePre}{}\renewcommand{\rhcaMzBLgbmFullClosePre}{0.98}
\providecommand{\rhcaMzBseLgbmFullClosePre}{}\renewcommand{\rhcaMzBseLgbmFullClosePre}{0.07}
\providecommand{\rhcaMzRsqLgbmFullClosePre}{}\renewcommand{\rhcaMzRsqLgbmFullClosePre}{0.58}
\providecommand{\rhcaMzlALgbmFullClosePre}{}\renewcommand{\rhcaMzlALgbmFullClosePre}{\ensuremath{-}0.89}
\providecommand{\rhcaMzlBLgbmFullClosePre}{}\renewcommand{\rhcaMzlBLgbmFullClosePre}{0.93}
\providecommand{\rhcaMzlBseLgbmFullClosePre}{}\renewcommand{\rhcaMzlBseLgbmFullClosePre}{0.02}
\providecommand{\rhcaMzlRsqLgbmFullClosePre}{}\renewcommand{\rhcaMzlRsqLgbmFullClosePre}{0.80}
\providecommand{\rhcaShBiasLgbmFullClosePre}{}\renewcommand{\rhcaShBiasLgbmFullClosePre}{0.1}
\providecommand{\rhcaShSlopeLgbmFullClosePre}{}\renewcommand{\rhcaShSlopeLgbmFullClosePre}{0.1}
\providecommand{\rhcaShDistLgbmFullClosePre}{}\renewcommand{\rhcaShDistLgbmFullClosePre}{99.9}
\providecommand{\rhcaQlLgbmFullClosePre}{}\renewcommand{\rhcaQlLgbmFullClosePre}{0.0048}
\providecommand{\rhcaQdLgbmFullClosePre}{}\renewcommand{\rhcaQdLgbmFullClosePre}{0.1239}
\providecommand{\rhcaOverSeLgbmFullClosePre}{}\renewcommand{\rhcaOverSeLgbmFullClosePre}{17.3}
\providecommand{\rhcaMrXgbHarClosePre}{}\renewcommand{\rhcaMrXgbHarClosePre}{1.064}
\providecommand{\rhcaRmXgbHarClosePre}{}\renewcommand{\rhcaRmXgbHarClosePre}{1.058}
\providecommand{\rhcaUnderXgbHarClosePre}{}\renewcommand{\rhcaUnderXgbHarClosePre}{42.6}
\providecommand{\rhcaMzAXgbHarClosePre}{}\renewcommand{\rhcaMzAXgbHarClosePre}{\ensuremath{-}0.12}
\providecommand{\rhcaMzBXgbHarClosePre}{}\renewcommand{\rhcaMzBXgbHarClosePre}{1.09}
\providecommand{\rhcaMzBseXgbHarClosePre}{}\renewcommand{\rhcaMzBseXgbHarClosePre}{0.09}
\providecommand{\rhcaMzRsqXgbHarClosePre}{}\renewcommand{\rhcaMzRsqXgbHarClosePre}{0.58}
\providecommand{\rhcaMzlAXgbHarClosePre}{}\renewcommand{\rhcaMzlAXgbHarClosePre}{\ensuremath{-}0.58}
\providecommand{\rhcaMzlBXgbHarClosePre}{}\renewcommand{\rhcaMzlBXgbHarClosePre}{0.96}
\providecommand{\rhcaMzlBseXgbHarClosePre}{}\renewcommand{\rhcaMzlBseXgbHarClosePre}{0.02}
\providecommand{\rhcaMzlRsqXgbHarClosePre}{}\renewcommand{\rhcaMzlRsqXgbHarClosePre}{0.78}
\providecommand{\rhcaShBiasXgbHarClosePre}{}\renewcommand{\rhcaShBiasXgbHarClosePre}{0.2}
\providecommand{\rhcaShSlopeXgbHarClosePre}{}\renewcommand{\rhcaShSlopeXgbHarClosePre}{0.9}
\providecommand{\rhcaShDistXgbHarClosePre}{}\renewcommand{\rhcaShDistXgbHarClosePre}{99.0}
\providecommand{\rhcaQlXgbHarClosePre}{}\renewcommand{\rhcaQlXgbHarClosePre}{0.0020}
\providecommand{\rhcaQdXgbHarClosePre}{}\renewcommand{\rhcaQdXgbHarClosePre}{0.1311}
\providecommand{\rhcaOverSeXgbHarClosePre}{}\renewcommand{\rhcaOverSeXgbHarClosePre}{14.2}
\providecommand{\rhcaMrXgbFullClosePre}{}\renewcommand{\rhcaMrXgbFullClosePre}{1.098}
\providecommand{\rhcaRmXgbFullClosePre}{}\renewcommand{\rhcaRmXgbFullClosePre}{1.037}
\providecommand{\rhcaUnderXgbFullClosePre}{}\renewcommand{\rhcaUnderXgbFullClosePre}{45.6}
\providecommand{\rhcaMzAXgbFullClosePre}{}\renewcommand{\rhcaMzAXgbFullClosePre}{0.19}
\providecommand{\rhcaMzBXgbFullClosePre}{}\renewcommand{\rhcaMzBXgbFullClosePre}{0.99}
\providecommand{\rhcaMzBseXgbFullClosePre}{}\renewcommand{\rhcaMzBseXgbFullClosePre}{0.08}
\providecommand{\rhcaMzRsqXgbFullClosePre}{}\renewcommand{\rhcaMzRsqXgbFullClosePre}{0.62}
\providecommand{\rhcaMzlAXgbFullClosePre}{}\renewcommand{\rhcaMzlAXgbFullClosePre}{\ensuremath{-}0.87}
\providecommand{\rhcaMzlBXgbFullClosePre}{}\renewcommand{\rhcaMzlBXgbFullClosePre}{0.94}
\providecommand{\rhcaMzlBseXgbFullClosePre}{}\renewcommand{\rhcaMzlBseXgbFullClosePre}{0.02}
\providecommand{\rhcaMzlRsqXgbFullClosePre}{}\renewcommand{\rhcaMzlRsqXgbFullClosePre}{0.80}
\providecommand{\rhcaShBiasXgbFullClosePre}{}\renewcommand{\rhcaShBiasXgbFullClosePre}{0.1}
\providecommand{\rhcaShSlopeXgbFullClosePre}{}\renewcommand{\rhcaShSlopeXgbFullClosePre}{0.0}
\providecommand{\rhcaShDistXgbFullClosePre}{}\renewcommand{\rhcaShDistXgbFullClosePre}{99.9}
\providecommand{\rhcaQlXgbFullClosePre}{}\renewcommand{\rhcaQlXgbFullClosePre}{0.0045}
\providecommand{\rhcaQdXgbFullClosePre}{}\renewcommand{\rhcaQdXgbFullClosePre}{0.1232}
\providecommand{\rhcaOverSeXgbFullClosePre}{}\renewcommand{\rhcaOverSeXgbFullClosePre}{18.8}
\providecommand{\rhcaMrRfHarClosePre}{}\renewcommand{\rhcaMrRfHarClosePre}{1.091}
\providecommand{\rhcaRmRfHarClosePre}{}\renewcommand{\rhcaRmRfHarClosePre}{1.073}
\providecommand{\rhcaUnderRfHarClosePre}{}\renewcommand{\rhcaUnderRfHarClosePre}{43.6}
\providecommand{\rhcaMzARfHarClosePre}{}\renewcommand{\rhcaMzARfHarClosePre}{\ensuremath{-}0.13}
\providecommand{\rhcaMzBRfHarClosePre}{}\renewcommand{\rhcaMzBRfHarClosePre}{1.11}
\providecommand{\rhcaMzBseRfHarClosePre}{}\renewcommand{\rhcaMzBseRfHarClosePre}{0.10}
\providecommand{\rhcaMzRsqRfHarClosePre}{}\renewcommand{\rhcaMzRsqRfHarClosePre}{0.59}
\providecommand{\rhcaMzlARfHarClosePre}{}\renewcommand{\rhcaMzlARfHarClosePre}{\ensuremath{-}0.68}
\providecommand{\rhcaMzlBRfHarClosePre}{}\renewcommand{\rhcaMzlBRfHarClosePre}{0.95}
\providecommand{\rhcaMzlBseRfHarClosePre}{}\renewcommand{\rhcaMzlBseRfHarClosePre}{0.02}
\providecommand{\rhcaMzlRsqRfHarClosePre}{}\renewcommand{\rhcaMzlRsqRfHarClosePre}{0.78}
\providecommand{\rhcaShBiasRfHarClosePre}{}\renewcommand{\rhcaShBiasRfHarClosePre}{0.2}
\providecommand{\rhcaShSlopeRfHarClosePre}{}\renewcommand{\rhcaShSlopeRfHarClosePre}{1.3}
\providecommand{\rhcaShDistRfHarClosePre}{}\renewcommand{\rhcaShDistRfHarClosePre}{98.5}
\providecommand{\rhcaQlRfHarClosePre}{}\renewcommand{\rhcaQlRfHarClosePre}{0.0039}
\providecommand{\rhcaQdRfHarClosePre}{}\renewcommand{\rhcaQdRfHarClosePre}{0.1325}
\providecommand{\rhcaOverSeRfHarClosePre}{}\renewcommand{\rhcaOverSeRfHarClosePre}{14.1}
\providecommand{\rhcaMrRfFullClosePre}{}\renewcommand{\rhcaMrRfFullClosePre}{1.119}
\providecommand{\rhcaRmRfFullClosePre}{}\renewcommand{\rhcaRmRfFullClosePre}{1.037}
\providecommand{\rhcaUnderRfFullClosePre}{}\renewcommand{\rhcaUnderRfFullClosePre}{46.1}
\providecommand{\rhcaMzARfFullClosePre}{}\renewcommand{\rhcaMzARfFullClosePre}{0.25}
\providecommand{\rhcaMzBRfFullClosePre}{}\renewcommand{\rhcaMzBRfFullClosePre}{0.98}
\providecommand{\rhcaMzBseRfFullClosePre}{}\renewcommand{\rhcaMzBseRfFullClosePre}{0.08}
\providecommand{\rhcaMzRsqRfFullClosePre}{}\renewcommand{\rhcaMzRsqRfFullClosePre}{0.61}
\providecommand{\rhcaMzlARfFullClosePre}{}\renewcommand{\rhcaMzlARfFullClosePre}{\ensuremath{-}1.04}
\providecommand{\rhcaMzlBRfFullClosePre}{}\renewcommand{\rhcaMzlBRfFullClosePre}{0.92}
\providecommand{\rhcaMzlBseRfFullClosePre}{}\renewcommand{\rhcaMzlBseRfFullClosePre}{0.02}
\providecommand{\rhcaMzlRsqRfFullClosePre}{}\renewcommand{\rhcaMzlRsqRfFullClosePre}{0.79}
\providecommand{\rhcaShBiasRfFullClosePre}{}\renewcommand{\rhcaShBiasRfFullClosePre}{0.1}
\providecommand{\rhcaShSlopeRfFullClosePre}{}\renewcommand{\rhcaShSlopeRfFullClosePre}{0.1}
\providecommand{\rhcaShDistRfFullClosePre}{}\renewcommand{\rhcaShDistRfFullClosePre}{99.8}
\providecommand{\rhcaQlRfFullClosePre}{}\renewcommand{\rhcaQlRfFullClosePre}{0.0065}
\providecommand{\rhcaQdRfFullClosePre}{}\renewcommand{\rhcaQdRfFullClosePre}{0.1270}
\providecommand{\rhcaOverSeRfFullClosePre}{}\renewcommand{\rhcaOverSeRfFullClosePre}{18.9}
\providecommand{\rhcaMrRidgeHarBarOos}{}\renewcommand{\rhcaMrRidgeHarBarOos}{1.061}
\providecommand{\rhcaRmRidgeHarBarOos}{}\renewcommand{\rhcaRmRidgeHarBarOos}{1.045}
\providecommand{\rhcaUnderRidgeHarBarOos}{}\renewcommand{\rhcaUnderRidgeHarBarOos}{38.0}
\providecommand{\rhcaMzARidgeHarBarOos}{}\renewcommand{\rhcaMzARidgeHarBarOos}{0.41}
\providecommand{\rhcaMzBRidgeHarBarOos}{}\renewcommand{\rhcaMzBRidgeHarBarOos}{0.99}
\providecommand{\rhcaMzBseRidgeHarBarOos}{}\renewcommand{\rhcaMzBseRidgeHarBarOos}{0.08}
\providecommand{\rhcaMzRsqRidgeHarBarOos}{}\renewcommand{\rhcaMzRsqRidgeHarBarOos}{0.67}
\providecommand{\rhcaMzlARidgeHarBarOos}{}\renewcommand{\rhcaMzlARidgeHarBarOos}{0.15}
\providecommand{\rhcaMzlBRidgeHarBarOos}{}\renewcommand{\rhcaMzlBRidgeHarBarOos}{1.02}
\providecommand{\rhcaMzlBseRidgeHarBarOos}{}\renewcommand{\rhcaMzlBseRidgeHarBarOos}{0.01}
\providecommand{\rhcaMzlRsqRidgeHarBarOos}{}\renewcommand{\rhcaMzlRsqRidgeHarBarOos}{0.79}
\providecommand{\rhcaShBiasRidgeHarBarOos}{}\renewcommand{\rhcaShBiasRidgeHarBarOos}{0.1}
\providecommand{\rhcaShSlopeRidgeHarBarOos}{}\renewcommand{\rhcaShSlopeRidgeHarBarOos}{0.0}
\providecommand{\rhcaShDistRidgeHarBarOos}{}\renewcommand{\rhcaShDistRidgeHarBarOos}{99.9}
\providecommand{\rhcaQlRidgeHarBarOos}{}\renewcommand{\rhcaQlRidgeHarBarOos}{0.0018}
\providecommand{\rhcaQdRidgeHarBarOos}{}\renewcommand{\rhcaQdRidgeHarBarOos}{0.1915}
\providecommand{\rhcaOverSeRidgeHarBarOos}{}\renewcommand{\rhcaOverSeRidgeHarBarOos}{21.5}
\providecommand{\rhcaMrRidgeFullBarOos}{}\renewcommand{\rhcaMrRidgeFullBarOos}{1.117}
\providecommand{\rhcaRmRidgeFullBarOos}{}\renewcommand{\rhcaRmRidgeFullBarOos}{0.977}
\providecommand{\rhcaUnderRidgeFullBarOos}{}\renewcommand{\rhcaUnderRidgeFullBarOos}{42.5}
\providecommand{\rhcaMzARidgeFullBarOos}{}\renewcommand{\rhcaMzARidgeFullBarOos}{3.27}
\providecommand{\rhcaMzBRidgeFullBarOos}{}\renewcommand{\rhcaMzBRidgeFullBarOos}{0.58}
\providecommand{\rhcaMzBseRidgeFullBarOos}{}\renewcommand{\rhcaMzBseRidgeFullBarOos}{0.10}
\providecommand{\rhcaMzRsqRidgeFullBarOos}{}\renewcommand{\rhcaMzRsqRidgeFullBarOos}{0.52}
\providecommand{\rhcaMzlARidgeFullBarOos}{}\renewcommand{\rhcaMzlARidgeFullBarOos}{\ensuremath{-}0.25}
\providecommand{\rhcaMzlBRidgeFullBarOos}{}\renewcommand{\rhcaMzlBRidgeFullBarOos}{0.99}
\providecommand{\rhcaMzlBseRidgeFullBarOos}{}\renewcommand{\rhcaMzlBseRidgeFullBarOos}{0.01}
\providecommand{\rhcaMzlRsqRidgeFullBarOos}{}\renewcommand{\rhcaMzlRsqRidgeFullBarOos}{0.81}
\providecommand{\rhcaShBiasRidgeFullBarOos}{}\renewcommand{\rhcaShBiasRidgeFullBarOos}{0.0}
\providecommand{\rhcaShSlopeRidgeFullBarOos}{}\renewcommand{\rhcaShSlopeRidgeFullBarOos}{35.9}
\providecommand{\rhcaShDistRidgeFullBarOos}{}\renewcommand{\rhcaShDistRidgeFullBarOos}{64.1}
\providecommand{\rhcaQlRidgeFullBarOos}{}\renewcommand{\rhcaQlRidgeFullBarOos}{0.0064}
\providecommand{\rhcaQdRidgeFullBarOos}{}\renewcommand{\rhcaQdRidgeFullBarOos}{0.1788}
\providecommand{\rhcaOverSeRidgeFullBarOos}{}\renewcommand{\rhcaOverSeRidgeFullBarOos}{63.1}
\providecommand{\rhcaMrLassoHarBarOos}{}\renewcommand{\rhcaMrLassoHarBarOos}{1.061}
\providecommand{\rhcaRmLassoHarBarOos}{}\renewcommand{\rhcaRmLassoHarBarOos}{1.045}
\providecommand{\rhcaUnderLassoHarBarOos}{}\renewcommand{\rhcaUnderLassoHarBarOos}{38.0}
\providecommand{\rhcaMzALassoHarBarOos}{}\renewcommand{\rhcaMzALassoHarBarOos}{0.41}
\providecommand{\rhcaMzBLassoHarBarOos}{}\renewcommand{\rhcaMzBLassoHarBarOos}{0.99}
\providecommand{\rhcaMzBseLassoHarBarOos}{}\renewcommand{\rhcaMzBseLassoHarBarOos}{0.08}
\providecommand{\rhcaMzRsqLassoHarBarOos}{}\renewcommand{\rhcaMzRsqLassoHarBarOos}{0.67}
\providecommand{\rhcaMzlALassoHarBarOos}{}\renewcommand{\rhcaMzlALassoHarBarOos}{0.15}
\providecommand{\rhcaMzlBLassoHarBarOos}{}\renewcommand{\rhcaMzlBLassoHarBarOos}{1.02}
\providecommand{\rhcaMzlBseLassoHarBarOos}{}\renewcommand{\rhcaMzlBseLassoHarBarOos}{0.01}
\providecommand{\rhcaMzlRsqLassoHarBarOos}{}\renewcommand{\rhcaMzlRsqLassoHarBarOos}{0.79}
\providecommand{\rhcaShBiasLassoHarBarOos}{}\renewcommand{\rhcaShBiasLassoHarBarOos}{0.1}
\providecommand{\rhcaShSlopeLassoHarBarOos}{}\renewcommand{\rhcaShSlopeLassoHarBarOos}{0.0}
\providecommand{\rhcaShDistLassoHarBarOos}{}\renewcommand{\rhcaShDistLassoHarBarOos}{99.9}
\providecommand{\rhcaQlLassoHarBarOos}{}\renewcommand{\rhcaQlLassoHarBarOos}{0.0018}
\providecommand{\rhcaQdLassoHarBarOos}{}\renewcommand{\rhcaQdLassoHarBarOos}{0.1915}
\providecommand{\rhcaOverSeLassoHarBarOos}{}\renewcommand{\rhcaOverSeLassoHarBarOos}{21.5}
\providecommand{\rhcaMrLassoFullBarOos}{}\renewcommand{\rhcaMrLassoFullBarOos}{1.116}
\providecommand{\rhcaRmLassoFullBarOos}{}\renewcommand{\rhcaRmLassoFullBarOos}{0.978}
\providecommand{\rhcaUnderLassoFullBarOos}{}\renewcommand{\rhcaUnderLassoFullBarOos}{43.2}
\providecommand{\rhcaMzALassoFullBarOos}{}\renewcommand{\rhcaMzALassoFullBarOos}{3.46}
\providecommand{\rhcaMzBLassoFullBarOos}{}\renewcommand{\rhcaMzBLassoFullBarOos}{0.56}
\providecommand{\rhcaMzBseLassoFullBarOos}{}\renewcommand{\rhcaMzBseLassoFullBarOos}{0.10}
\providecommand{\rhcaMzRsqLassoFullBarOos}{}\renewcommand{\rhcaMzRsqLassoFullBarOos}{0.51}
\providecommand{\rhcaMzlALassoFullBarOos}{}\renewcommand{\rhcaMzlALassoFullBarOos}{\ensuremath{-}0.26}
\providecommand{\rhcaMzlBLassoFullBarOos}{}\renewcommand{\rhcaMzlBLassoFullBarOos}{0.98}
\providecommand{\rhcaMzlBseLassoFullBarOos}{}\renewcommand{\rhcaMzlBseLassoFullBarOos}{0.01}
\providecommand{\rhcaMzlRsqLassoFullBarOos}{}\renewcommand{\rhcaMzlRsqLassoFullBarOos}{0.81}
\providecommand{\rhcaShBiasLassoFullBarOos}{}\renewcommand{\rhcaShBiasLassoFullBarOos}{0.0}
\providecommand{\rhcaShSlopeLassoFullBarOos}{}\renewcommand{\rhcaShSlopeLassoFullBarOos}{38.2}
\providecommand{\rhcaShDistLassoFullBarOos}{}\renewcommand{\rhcaShDistLassoFullBarOos}{61.8}
\providecommand{\rhcaQlLassoFullBarOos}{}\renewcommand{\rhcaQlLassoFullBarOos}{0.0062}
\providecommand{\rhcaQdLassoFullBarOos}{}\renewcommand{\rhcaQdLassoFullBarOos}{0.1739}
\providecommand{\rhcaOverSeLassoFullBarOos}{}\renewcommand{\rhcaOverSeLassoFullBarOos}{64.5}
\providecommand{\rhcaMrEnetHarBarOos}{}\renewcommand{\rhcaMrEnetHarBarOos}{1.061}
\providecommand{\rhcaRmEnetHarBarOos}{}\renewcommand{\rhcaRmEnetHarBarOos}{1.045}
\providecommand{\rhcaUnderEnetHarBarOos}{}\renewcommand{\rhcaUnderEnetHarBarOos}{38.0}
\providecommand{\rhcaMzAEnetHarBarOos}{}\renewcommand{\rhcaMzAEnetHarBarOos}{0.41}
\providecommand{\rhcaMzBEnetHarBarOos}{}\renewcommand{\rhcaMzBEnetHarBarOos}{0.99}
\providecommand{\rhcaMzBseEnetHarBarOos}{}\renewcommand{\rhcaMzBseEnetHarBarOos}{0.08}
\providecommand{\rhcaMzRsqEnetHarBarOos}{}\renewcommand{\rhcaMzRsqEnetHarBarOos}{0.67}
\providecommand{\rhcaMzlAEnetHarBarOos}{}\renewcommand{\rhcaMzlAEnetHarBarOos}{0.15}
\providecommand{\rhcaMzlBEnetHarBarOos}{}\renewcommand{\rhcaMzlBEnetHarBarOos}{1.02}
\providecommand{\rhcaMzlBseEnetHarBarOos}{}\renewcommand{\rhcaMzlBseEnetHarBarOos}{0.01}
\providecommand{\rhcaMzlRsqEnetHarBarOos}{}\renewcommand{\rhcaMzlRsqEnetHarBarOos}{0.79}
\providecommand{\rhcaShBiasEnetHarBarOos}{}\renewcommand{\rhcaShBiasEnetHarBarOos}{0.1}
\providecommand{\rhcaShSlopeEnetHarBarOos}{}\renewcommand{\rhcaShSlopeEnetHarBarOos}{0.0}
\providecommand{\rhcaShDistEnetHarBarOos}{}\renewcommand{\rhcaShDistEnetHarBarOos}{99.9}
\providecommand{\rhcaQlEnetHarBarOos}{}\renewcommand{\rhcaQlEnetHarBarOos}{0.0018}
\providecommand{\rhcaQdEnetHarBarOos}{}\renewcommand{\rhcaQdEnetHarBarOos}{0.1915}
\providecommand{\rhcaOverSeEnetHarBarOos}{}\renewcommand{\rhcaOverSeEnetHarBarOos}{21.5}
\providecommand{\rhcaMrEnetFullBarOos}{}\renewcommand{\rhcaMrEnetFullBarOos}{1.109}
\providecommand{\rhcaRmEnetFullBarOos}{}\renewcommand{\rhcaRmEnetFullBarOos}{0.966}
\providecommand{\rhcaUnderEnetFullBarOos}{}\renewcommand{\rhcaUnderEnetFullBarOos}{42.2}
\providecommand{\rhcaMzAEnetFullBarOos}{}\renewcommand{\rhcaMzAEnetFullBarOos}{3.45}
\providecommand{\rhcaMzBEnetFullBarOos}{}\renewcommand{\rhcaMzBEnetFullBarOos}{0.56}
\providecommand{\rhcaMzBseEnetFullBarOos}{}\renewcommand{\rhcaMzBseEnetFullBarOos}{0.10}
\providecommand{\rhcaMzRsqEnetFullBarOos}{}\renewcommand{\rhcaMzRsqEnetFullBarOos}{0.52}
\providecommand{\rhcaMzlAEnetFullBarOos}{}\renewcommand{\rhcaMzlAEnetFullBarOos}{\ensuremath{-}0.22}
\providecommand{\rhcaMzlBEnetFullBarOos}{}\renewcommand{\rhcaMzlBEnetFullBarOos}{0.99}
\providecommand{\rhcaMzlBseEnetFullBarOos}{}\renewcommand{\rhcaMzlBseEnetFullBarOos}{0.01}
\providecommand{\rhcaMzlRsqEnetFullBarOos}{}\renewcommand{\rhcaMzlRsqEnetFullBarOos}{0.81}
\providecommand{\rhcaShBiasEnetFullBarOos}{}\renewcommand{\rhcaShBiasEnetFullBarOos}{0.0}
\providecommand{\rhcaShSlopeEnetFullBarOos}{}\renewcommand{\rhcaShSlopeEnetFullBarOos}{40.8}
\providecommand{\rhcaShDistEnetFullBarOos}{}\renewcommand{\rhcaShDistEnetFullBarOos}{59.1}
\providecommand{\rhcaQlEnetFullBarOos}{}\renewcommand{\rhcaQlEnetFullBarOos}{0.0056}
\providecommand{\rhcaQdEnetFullBarOos}{}\renewcommand{\rhcaQdEnetFullBarOos}{0.1775}
\providecommand{\rhcaOverSeEnetFullBarOos}{}\renewcommand{\rhcaOverSeEnetFullBarOos}{67.0}
\providecommand{\rhcaMrLgbmHarBarOos}{}\renewcommand{\rhcaMrLgbmHarBarOos}{1.067}
\providecommand{\rhcaRmLgbmHarBarOos}{}\renewcommand{\rhcaRmLgbmHarBarOos}{1.028}
\providecommand{\rhcaUnderLgbmHarBarOos}{}\renewcommand{\rhcaUnderLgbmHarBarOos}{38.5}
\providecommand{\rhcaMzALgbmHarBarOos}{}\renewcommand{\rhcaMzALgbmHarBarOos}{0.72}
\providecommand{\rhcaMzBLgbmHarBarOos}{}\renewcommand{\rhcaMzBLgbmHarBarOos}{0.94}
\providecommand{\rhcaMzBseLgbmHarBarOos}{}\renewcommand{\rhcaMzBseLgbmHarBarOos}{0.06}
\providecommand{\rhcaMzRsqLgbmHarBarOos}{}\renewcommand{\rhcaMzRsqLgbmHarBarOos}{0.66}
\providecommand{\rhcaMzlALgbmHarBarOos}{}\renewcommand{\rhcaMzlALgbmHarBarOos}{0.03}
\providecommand{\rhcaMzlBLgbmHarBarOos}{}\renewcommand{\rhcaMzlBLgbmHarBarOos}{1.01}
\providecommand{\rhcaMzlBseLgbmHarBarOos}{}\renewcommand{\rhcaMzlBseLgbmHarBarOos}{0.01}
\providecommand{\rhcaMzlRsqLgbmHarBarOos}{}\renewcommand{\rhcaMzlRsqLgbmHarBarOos}{0.80}
\providecommand{\rhcaShBiasLgbmHarBarOos}{}\renewcommand{\rhcaShBiasLgbmHarBarOos}{0.0}
\providecommand{\rhcaShSlopeLgbmHarBarOos}{}\renewcommand{\rhcaShSlopeLgbmHarBarOos}{0.8}
\providecommand{\rhcaShDistLgbmHarBarOos}{}\renewcommand{\rhcaShDistLgbmHarBarOos}{99.1}
\providecommand{\rhcaQlLgbmHarBarOos}{}\renewcommand{\rhcaQlLgbmHarBarOos}{0.0022}
\providecommand{\rhcaQdLgbmHarBarOos}{}\renewcommand{\rhcaQdLgbmHarBarOos}{0.1878}
\providecommand{\rhcaOverSeLgbmHarBarOos}{}\renewcommand{\rhcaOverSeLgbmHarBarOos}{25.9}
\providecommand{\rhcaMrLgbmFullBarOos}{}\renewcommand{\rhcaMrLgbmFullBarOos}{1.072}
\providecommand{\rhcaRmLgbmFullBarOos}{}\renewcommand{\rhcaRmLgbmFullBarOos}{0.963}
\providecommand{\rhcaUnderLgbmFullBarOos}{}\renewcommand{\rhcaUnderLgbmFullBarOos}{40.6}
\providecommand{\rhcaMzALgbmFullBarOos}{}\renewcommand{\rhcaMzALgbmFullBarOos}{1.86}
\providecommand{\rhcaMzBLgbmFullBarOos}{}\renewcommand{\rhcaMzBLgbmFullBarOos}{0.74}
\providecommand{\rhcaMzBseLgbmFullBarOos}{}\renewcommand{\rhcaMzBseLgbmFullBarOos}{0.07}
\providecommand{\rhcaMzRsqLgbmFullBarOos}{}\renewcommand{\rhcaMzRsqLgbmFullBarOos}{0.64}
\providecommand{\rhcaMzlALgbmFullBarOos}{}\renewcommand{\rhcaMzlALgbmFullBarOos}{\ensuremath{-}0.20}
\providecommand{\rhcaMzlBLgbmFullBarOos}{}\renewcommand{\rhcaMzlBLgbmFullBarOos}{0.99}
\providecommand{\rhcaMzlBseLgbmFullBarOos}{}\renewcommand{\rhcaMzlBseLgbmFullBarOos}{0.01}
\providecommand{\rhcaMzlRsqLgbmFullBarOos}{}\renewcommand{\rhcaMzlRsqLgbmFullBarOos}{0.82}
\providecommand{\rhcaShBiasLgbmFullBarOos}{}\renewcommand{\rhcaShBiasLgbmFullBarOos}{0.0}
\providecommand{\rhcaShSlopeLgbmFullBarOos}{}\renewcommand{\rhcaShSlopeLgbmFullBarOos}{17.7}
\providecommand{\rhcaShDistLgbmFullBarOos}{}\renewcommand{\rhcaShDistLgbmFullBarOos}{82.2}
\providecommand{\rhcaQlLgbmFullBarOos}{}\renewcommand{\rhcaQlLgbmFullBarOos}{0.0025}
\providecommand{\rhcaQdLgbmFullBarOos}{}\renewcommand{\rhcaQdLgbmFullBarOos}{0.1621}
\providecommand{\rhcaOverSeLgbmFullBarOos}{}\renewcommand{\rhcaOverSeLgbmFullBarOos}{49.1}
\providecommand{\rhcaMrXgbHarBarOos}{}\renewcommand{\rhcaMrXgbHarBarOos}{1.069}
\providecommand{\rhcaRmXgbHarBarOos}{}\renewcommand{\rhcaRmXgbHarBarOos}{1.028}
\providecommand{\rhcaUnderXgbHarBarOos}{}\renewcommand{\rhcaUnderXgbHarBarOos}{38.4}
\providecommand{\rhcaMzAXgbHarBarOos}{}\renewcommand{\rhcaMzAXgbHarBarOos}{0.69}
\providecommand{\rhcaMzBXgbHarBarOos}{}\renewcommand{\rhcaMzBXgbHarBarOos}{0.94}
\providecommand{\rhcaMzBseXgbHarBarOos}{}\renewcommand{\rhcaMzBseXgbHarBarOos}{0.06}
\providecommand{\rhcaMzRsqXgbHarBarOos}{}\renewcommand{\rhcaMzRsqXgbHarBarOos}{0.66}
\providecommand{\rhcaMzlAXgbHarBarOos}{}\renewcommand{\rhcaMzlAXgbHarBarOos}{0.04}
\providecommand{\rhcaMzlBXgbHarBarOos}{}\renewcommand{\rhcaMzlBXgbHarBarOos}{1.01}
\providecommand{\rhcaMzlBseXgbHarBarOos}{}\renewcommand{\rhcaMzlBseXgbHarBarOos}{0.01}
\providecommand{\rhcaMzlRsqXgbHarBarOos}{}\renewcommand{\rhcaMzlRsqXgbHarBarOos}{0.80}
\providecommand{\rhcaShBiasXgbHarBarOos}{}\renewcommand{\rhcaShBiasXgbHarBarOos}{0.0}
\providecommand{\rhcaShSlopeXgbHarBarOos}{}\renewcommand{\rhcaShSlopeXgbHarBarOos}{0.8}
\providecommand{\rhcaShDistXgbHarBarOos}{}\renewcommand{\rhcaShDistXgbHarBarOos}{99.2}
\providecommand{\rhcaQlXgbHarBarOos}{}\renewcommand{\rhcaQlXgbHarBarOos}{0.0023}
\providecommand{\rhcaQdXgbHarBarOos}{}\renewcommand{\rhcaQdXgbHarBarOos}{0.1899}
\providecommand{\rhcaOverSeXgbHarBarOos}{}\renewcommand{\rhcaOverSeXgbHarBarOos}{25.9}
\providecommand{\rhcaMrXgbFullBarOos}{}\renewcommand{\rhcaMrXgbFullBarOos}{1.068}
\providecommand{\rhcaRmXgbFullBarOos}{}\renewcommand{\rhcaRmXgbFullBarOos}{0.966}
\providecommand{\rhcaUnderXgbFullBarOos}{}\renewcommand{\rhcaUnderXgbFullBarOos}{40.5}
\providecommand{\rhcaMzAXgbFullBarOos}{}\renewcommand{\rhcaMzAXgbFullBarOos}{1.74}
\providecommand{\rhcaMzBXgbFullBarOos}{}\renewcommand{\rhcaMzBXgbFullBarOos}{0.76}
\providecommand{\rhcaMzBseXgbFullBarOos}{}\renewcommand{\rhcaMzBseXgbFullBarOos}{0.06}
\providecommand{\rhcaMzRsqXgbFullBarOos}{}\renewcommand{\rhcaMzRsqXgbFullBarOos}{0.66}
\providecommand{\rhcaMzlAXgbFullBarOos}{}\renewcommand{\rhcaMzlAXgbFullBarOos}{\ensuremath{-}0.18}
\providecommand{\rhcaMzlBXgbFullBarOos}{}\renewcommand{\rhcaMzlBXgbFullBarOos}{0.99}
\providecommand{\rhcaMzlBseXgbFullBarOos}{}\renewcommand{\rhcaMzlBseXgbFullBarOos}{0.01}
\providecommand{\rhcaMzlRsqXgbFullBarOos}{}\renewcommand{\rhcaMzlRsqXgbFullBarOos}{0.82}
\providecommand{\rhcaShBiasXgbFullBarOos}{}\renewcommand{\rhcaShBiasXgbFullBarOos}{0.0}
\providecommand{\rhcaShSlopeXgbFullBarOos}{}\renewcommand{\rhcaShSlopeXgbFullBarOos}{16.1}
\providecommand{\rhcaShDistXgbFullBarOos}{}\renewcommand{\rhcaShDistXgbFullBarOos}{83.9}
\providecommand{\rhcaQlXgbFullBarOos}{}\renewcommand{\rhcaQlXgbFullBarOos}{0.0022}
\providecommand{\rhcaQdXgbFullBarOos}{}\renewcommand{\rhcaQdXgbFullBarOos}{0.1612}
\providecommand{\rhcaOverSeXgbFullBarOos}{}\renewcommand{\rhcaOverSeXgbFullBarOos}{48.3}
\providecommand{\rhcaMrRfHarBarOos}{}\renewcommand{\rhcaMrRfHarBarOos}{1.077}
\providecommand{\rhcaRmRfHarBarOos}{}\renewcommand{\rhcaRmRfHarBarOos}{1.024}
\providecommand{\rhcaUnderRfHarBarOos}{}\renewcommand{\rhcaUnderRfHarBarOos}{38.8}
\providecommand{\rhcaMzARfHarBarOos}{}\renewcommand{\rhcaMzARfHarBarOos}{0.82}
\providecommand{\rhcaMzBRfHarBarOos}{}\renewcommand{\rhcaMzBRfHarBarOos}{0.92}
\providecommand{\rhcaMzBseRfHarBarOos}{}\renewcommand{\rhcaMzBseRfHarBarOos}{0.07}
\providecommand{\rhcaMzRsqRfHarBarOos}{}\renewcommand{\rhcaMzRsqRfHarBarOos}{0.65}
\providecommand{\rhcaMzlARfHarBarOos}{}\renewcommand{\rhcaMzlARfHarBarOos}{\ensuremath{-}0.07}
\providecommand{\rhcaMzlBRfHarBarOos}{}\renewcommand{\rhcaMzlBRfHarBarOos}{1.00}
\providecommand{\rhcaMzlBseRfHarBarOos}{}\renewcommand{\rhcaMzlBseRfHarBarOos}{0.01}
\providecommand{\rhcaMzlRsqRfHarBarOos}{}\renewcommand{\rhcaMzlRsqRfHarBarOos}{0.80}
\providecommand{\rhcaShBiasRfHarBarOos}{}\renewcommand{\rhcaShBiasRfHarBarOos}{0.0}
\providecommand{\rhcaShSlopeRfHarBarOos}{}\renewcommand{\rhcaShSlopeRfHarBarOos}{1.4}
\providecommand{\rhcaShDistRfHarBarOos}{}\renewcommand{\rhcaShDistRfHarBarOos}{98.6}
\providecommand{\rhcaQlRfHarBarOos}{}\renewcommand{\rhcaQlRfHarBarOos}{0.0028}
\providecommand{\rhcaQdRfHarBarOos}{}\renewcommand{\rhcaQdRfHarBarOos}{0.1909}
\providecommand{\rhcaOverSeRfHarBarOos}{}\renewcommand{\rhcaOverSeRfHarBarOos}{26.7}
\providecommand{\rhcaMrRfFullBarOos}{}\renewcommand{\rhcaMrRfFullBarOos}{1.087}
\providecommand{\rhcaRmRfFullBarOos}{}\renewcommand{\rhcaRmRfFullBarOos}{0.974}
\providecommand{\rhcaUnderRfFullBarOos}{}\renewcommand{\rhcaUnderRfFullBarOos}{41.0}
\providecommand{\rhcaMzARfFullBarOos}{}\renewcommand{\rhcaMzARfFullBarOos}{1.61}
\providecommand{\rhcaMzBRfFullBarOos}{}\renewcommand{\rhcaMzBRfFullBarOos}{0.78}
\providecommand{\rhcaMzBseRfFullBarOos}{}\renewcommand{\rhcaMzBseRfFullBarOos}{0.07}
\providecommand{\rhcaMzRsqRfFullBarOos}{}\renewcommand{\rhcaMzRsqRfFullBarOos}{0.65}
\providecommand{\rhcaMzlARfFullBarOos}{}\renewcommand{\rhcaMzlARfFullBarOos}{\ensuremath{-}0.25}
\providecommand{\rhcaMzlBRfFullBarOos}{}\renewcommand{\rhcaMzlBRfFullBarOos}{0.99}
\providecommand{\rhcaMzlBseRfFullBarOos}{}\renewcommand{\rhcaMzlBseRfFullBarOos}{0.01}
\providecommand{\rhcaMzlRsqRfFullBarOos}{}\renewcommand{\rhcaMzlRsqRfFullBarOos}{0.82}
\providecommand{\rhcaShBiasRfFullBarOos}{}\renewcommand{\rhcaShBiasRfFullBarOos}{0.0}
\providecommand{\rhcaShSlopeRfFullBarOos}{}\renewcommand{\rhcaShSlopeRfFullBarOos}{12.7}
\providecommand{\rhcaShDistRfFullBarOos}{}\renewcommand{\rhcaShDistRfFullBarOos}{87.3}
\providecommand{\rhcaQlRfFullBarOos}{}\renewcommand{\rhcaQlRfFullBarOos}{0.0035}
\providecommand{\rhcaQdRfFullBarOos}{}\renewcommand{\rhcaQdRfFullBarOos}{0.1694}
\providecommand{\rhcaOverSeRfFullBarOos}{}\renewcommand{\rhcaOverSeRfFullBarOos}{44.4}
\providecommand{\rhcaMrRidgeHarCloseOos}{}\renewcommand{\rhcaMrRidgeHarCloseOos}{1.092}
\providecommand{\rhcaRmRidgeHarCloseOos}{}\renewcommand{\rhcaRmRidgeHarCloseOos}{1.103}
\providecommand{\rhcaUnderRidgeHarCloseOos}{}\renewcommand{\rhcaUnderRidgeHarCloseOos}{47.7}
\providecommand{\rhcaMzARidgeHarCloseOos}{}\renewcommand{\rhcaMzARidgeHarCloseOos}{\ensuremath{-}1.23}
\providecommand{\rhcaMzBRidgeHarCloseOos}{}\renewcommand{\rhcaMzBRidgeHarCloseOos}{1.21}
\providecommand{\rhcaMzBseRidgeHarCloseOos}{}\renewcommand{\rhcaMzBseRidgeHarCloseOos}{0.09}
\providecommand{\rhcaMzRsqRidgeHarCloseOos}{}\renewcommand{\rhcaMzRsqRidgeHarCloseOos}{0.80}
\providecommand{\rhcaMzlARidgeHarCloseOos}{}\renewcommand{\rhcaMzlARidgeHarCloseOos}{\ensuremath{-}0.38}
\providecommand{\rhcaMzlBRidgeHarCloseOos}{}\renewcommand{\rhcaMzlBRidgeHarCloseOos}{0.97}
\providecommand{\rhcaMzlBseRidgeHarCloseOos}{}\renewcommand{\rhcaMzlBseRidgeHarCloseOos}{0.02}
\providecommand{\rhcaMzlRsqRidgeHarCloseOos}{}\renewcommand{\rhcaMzlRsqRidgeHarCloseOos}{0.84}
\providecommand{\rhcaShBiasRidgeHarCloseOos}{}\renewcommand{\rhcaShBiasRidgeHarCloseOos}{0.3}
\providecommand{\rhcaShSlopeRidgeHarCloseOos}{}\renewcommand{\rhcaShSlopeRidgeHarCloseOos}{10.7}
\providecommand{\rhcaShDistRidgeHarCloseOos}{}\renewcommand{\rhcaShDistRidgeHarCloseOos}{88.9}
\providecommand{\rhcaQlRidgeHarCloseOos}{}\renewcommand{\rhcaQlRidgeHarCloseOos}{0.0040}
\providecommand{\rhcaQdRidgeHarCloseOos}{}\renewcommand{\rhcaQdRidgeHarCloseOos}{0.1097}
\providecommand{\rhcaOverSeRidgeHarCloseOos}{}\renewcommand{\rhcaOverSeRidgeHarCloseOos}{13.8}
\providecommand{\rhcaMrRidgeFullCloseOos}{}\renewcommand{\rhcaMrRidgeFullCloseOos}{1.169}
\providecommand{\rhcaRmRidgeFullCloseOos}{}\renewcommand{\rhcaRmRidgeFullCloseOos}{0.971}
\providecommand{\rhcaUnderRidgeFullCloseOos}{}\renewcommand{\rhcaUnderRidgeFullCloseOos}{52.1}
\providecommand{\rhcaMzARidgeFullCloseOos}{}\renewcommand{\rhcaMzARidgeFullCloseOos}{3.60}
\providecommand{\rhcaMzBRidgeFullCloseOos}{}\renewcommand{\rhcaMzBRidgeFullCloseOos}{0.70}
\providecommand{\rhcaMzBseRidgeFullCloseOos}{}\renewcommand{\rhcaMzBseRidgeFullCloseOos}{0.06}
\providecommand{\rhcaMzRsqRidgeFullCloseOos}{}\renewcommand{\rhcaMzRsqRidgeFullCloseOos}{0.72}
\providecommand{\rhcaMzlARidgeFullCloseOos}{}\renewcommand{\rhcaMzlARidgeFullCloseOos}{\ensuremath{-}0.95}
\providecommand{\rhcaMzlBRidgeFullCloseOos}{}\renewcommand{\rhcaMzlBRidgeFullCloseOos}{0.92}
\providecommand{\rhcaMzlBseRidgeFullCloseOos}{}\renewcommand{\rhcaMzlBseRidgeFullCloseOos}{0.01}
\providecommand{\rhcaMzlRsqRidgeFullCloseOos}{}\renewcommand{\rhcaMzlRsqRidgeFullCloseOos}{0.84}
\providecommand{\rhcaShBiasRidgeFullCloseOos}{}\renewcommand{\rhcaShBiasRidgeFullCloseOos}{0.0}
\providecommand{\rhcaShSlopeRidgeFullCloseOos}{}\renewcommand{\rhcaShSlopeRidgeFullCloseOos}{32.0}
\providecommand{\rhcaShDistRidgeFullCloseOos}{}\renewcommand{\rhcaShDistRidgeFullCloseOos}{68.0}
\providecommand{\rhcaQlRidgeFullCloseOos}{}\renewcommand{\rhcaQlRidgeFullCloseOos}{0.0128}
\providecommand{\rhcaQdRidgeFullCloseOos}{}\renewcommand{\rhcaQdRidgeFullCloseOos}{0.1140}
\providecommand{\rhcaOverSeRidgeFullCloseOos}{}\renewcommand{\rhcaOverSeRidgeFullCloseOos}{52.6}
\providecommand{\rhcaMrLassoHarCloseOos}{}\renewcommand{\rhcaMrLassoHarCloseOos}{1.092}
\providecommand{\rhcaRmLassoHarCloseOos}{}\renewcommand{\rhcaRmLassoHarCloseOos}{1.103}
\providecommand{\rhcaUnderLassoHarCloseOos}{}\renewcommand{\rhcaUnderLassoHarCloseOos}{47.7}
\providecommand{\rhcaMzALassoHarCloseOos}{}\renewcommand{\rhcaMzALassoHarCloseOos}{\ensuremath{-}1.23}
\providecommand{\rhcaMzBLassoHarCloseOos}{}\renewcommand{\rhcaMzBLassoHarCloseOos}{1.21}
\providecommand{\rhcaMzBseLassoHarCloseOos}{}\renewcommand{\rhcaMzBseLassoHarCloseOos}{0.09}
\providecommand{\rhcaMzRsqLassoHarCloseOos}{}\renewcommand{\rhcaMzRsqLassoHarCloseOos}{0.80}
\providecommand{\rhcaMzlALassoHarCloseOos}{}\renewcommand{\rhcaMzlALassoHarCloseOos}{\ensuremath{-}0.38}
\providecommand{\rhcaMzlBLassoHarCloseOos}{}\renewcommand{\rhcaMzlBLassoHarCloseOos}{0.97}
\providecommand{\rhcaMzlBseLassoHarCloseOos}{}\renewcommand{\rhcaMzlBseLassoHarCloseOos}{0.02}
\providecommand{\rhcaMzlRsqLassoHarCloseOos}{}\renewcommand{\rhcaMzlRsqLassoHarCloseOos}{0.84}
\providecommand{\rhcaShBiasLassoHarCloseOos}{}\renewcommand{\rhcaShBiasLassoHarCloseOos}{0.3}
\providecommand{\rhcaShSlopeLassoHarCloseOos}{}\renewcommand{\rhcaShSlopeLassoHarCloseOos}{10.7}
\providecommand{\rhcaShDistLassoHarCloseOos}{}\renewcommand{\rhcaShDistLassoHarCloseOos}{88.9}
\providecommand{\rhcaQlLassoHarCloseOos}{}\renewcommand{\rhcaQlLassoHarCloseOos}{0.0040}
\providecommand{\rhcaQdLassoHarCloseOos}{}\renewcommand{\rhcaQdLassoHarCloseOos}{0.1097}
\providecommand{\rhcaOverSeLassoHarCloseOos}{}\renewcommand{\rhcaOverSeLassoHarCloseOos}{13.8}
\providecommand{\rhcaMrLassoFullCloseOos}{}\renewcommand{\rhcaMrLassoFullCloseOos}{1.183}
\providecommand{\rhcaRmLassoFullCloseOos}{}\renewcommand{\rhcaRmLassoFullCloseOos}{0.982}
\providecommand{\rhcaUnderLassoFullCloseOos}{}\renewcommand{\rhcaUnderLassoFullCloseOos}{52.6}
\providecommand{\rhcaMzALassoFullCloseOos}{}\renewcommand{\rhcaMzALassoFullCloseOos}{3.74}
\providecommand{\rhcaMzBLassoFullCloseOos}{}\renewcommand{\rhcaMzBLassoFullCloseOos}{0.69}
\providecommand{\rhcaMzBseLassoFullCloseOos}{}\renewcommand{\rhcaMzBseLassoFullCloseOos}{0.05}
\providecommand{\rhcaMzRsqLassoFullCloseOos}{}\renewcommand{\rhcaMzRsqLassoFullCloseOos}{0.71}
\providecommand{\rhcaMzlALassoFullCloseOos}{}\renewcommand{\rhcaMzlALassoFullCloseOos}{\ensuremath{-}0.93}
\providecommand{\rhcaMzlBLassoFullCloseOos}{}\renewcommand{\rhcaMzlBLassoFullCloseOos}{0.92}
\providecommand{\rhcaMzlBseLassoFullCloseOos}{}\renewcommand{\rhcaMzlBseLassoFullCloseOos}{0.01}
\providecommand{\rhcaMzlRsqLassoFullCloseOos}{}\renewcommand{\rhcaMzlRsqLassoFullCloseOos}{0.84}
\providecommand{\rhcaShBiasLassoFullCloseOos}{}\renewcommand{\rhcaShBiasLassoFullCloseOos}{0.0}
\providecommand{\rhcaShSlopeLassoFullCloseOos}{}\renewcommand{\rhcaShSlopeLassoFullCloseOos}{32.1}
\providecommand{\rhcaShDistLassoFullCloseOos}{}\renewcommand{\rhcaShDistLassoFullCloseOos}{67.9}
\providecommand{\rhcaQlLassoFullCloseOos}{}\renewcommand{\rhcaQlLassoFullCloseOos}{0.0150}
\providecommand{\rhcaQdLassoFullCloseOos}{}\renewcommand{\rhcaQdLassoFullCloseOos}{0.1159}
\providecommand{\rhcaOverSeLassoFullCloseOos}{}\renewcommand{\rhcaOverSeLassoFullCloseOos}{51.5}
\providecommand{\rhcaMrEnetHarCloseOos}{}\renewcommand{\rhcaMrEnetHarCloseOos}{1.092}
\providecommand{\rhcaRmEnetHarCloseOos}{}\renewcommand{\rhcaRmEnetHarCloseOos}{1.103}
\providecommand{\rhcaUnderEnetHarCloseOos}{}\renewcommand{\rhcaUnderEnetHarCloseOos}{47.7}
\providecommand{\rhcaMzAEnetHarCloseOos}{}\renewcommand{\rhcaMzAEnetHarCloseOos}{\ensuremath{-}1.23}
\providecommand{\rhcaMzBEnetHarCloseOos}{}\renewcommand{\rhcaMzBEnetHarCloseOos}{1.21}
\providecommand{\rhcaMzBseEnetHarCloseOos}{}\renewcommand{\rhcaMzBseEnetHarCloseOos}{0.09}
\providecommand{\rhcaMzRsqEnetHarCloseOos}{}\renewcommand{\rhcaMzRsqEnetHarCloseOos}{0.80}
\providecommand{\rhcaMzlAEnetHarCloseOos}{}\renewcommand{\rhcaMzlAEnetHarCloseOos}{\ensuremath{-}0.38}
\providecommand{\rhcaMzlBEnetHarCloseOos}{}\renewcommand{\rhcaMzlBEnetHarCloseOos}{0.97}
\providecommand{\rhcaMzlBseEnetHarCloseOos}{}\renewcommand{\rhcaMzlBseEnetHarCloseOos}{0.02}
\providecommand{\rhcaMzlRsqEnetHarCloseOos}{}\renewcommand{\rhcaMzlRsqEnetHarCloseOos}{0.84}
\providecommand{\rhcaShBiasEnetHarCloseOos}{}\renewcommand{\rhcaShBiasEnetHarCloseOos}{0.3}
\providecommand{\rhcaShSlopeEnetHarCloseOos}{}\renewcommand{\rhcaShSlopeEnetHarCloseOos}{10.7}
\providecommand{\rhcaShDistEnetHarCloseOos}{}\renewcommand{\rhcaShDistEnetHarCloseOos}{88.9}
\providecommand{\rhcaQlEnetHarCloseOos}{}\renewcommand{\rhcaQlEnetHarCloseOos}{0.0040}
\providecommand{\rhcaQdEnetHarCloseOos}{}\renewcommand{\rhcaQdEnetHarCloseOos}{0.1097}
\providecommand{\rhcaOverSeEnetHarCloseOos}{}\renewcommand{\rhcaOverSeEnetHarCloseOos}{13.8}
\providecommand{\rhcaMrEnetFullCloseOos}{}\renewcommand{\rhcaMrEnetFullCloseOos}{1.165}
\providecommand{\rhcaRmEnetFullCloseOos}{}\renewcommand{\rhcaRmEnetFullCloseOos}{0.955}
\providecommand{\rhcaUnderEnetFullCloseOos}{}\renewcommand{\rhcaUnderEnetFullCloseOos}{52.9}
\providecommand{\rhcaMzAEnetFullCloseOos}{}\renewcommand{\rhcaMzAEnetFullCloseOos}{3.93}
\providecommand{\rhcaMzBEnetFullCloseOos}{}\renewcommand{\rhcaMzBEnetFullCloseOos}{0.66}
\providecommand{\rhcaMzBseEnetFullCloseOos}{}\renewcommand{\rhcaMzBseEnetFullCloseOos}{0.05}
\providecommand{\rhcaMzRsqEnetFullCloseOos}{}\renewcommand{\rhcaMzRsqEnetFullCloseOos}{0.72}
\providecommand{\rhcaMzlAEnetFullCloseOos}{}\renewcommand{\rhcaMzlAEnetFullCloseOos}{\ensuremath{-}0.94}
\providecommand{\rhcaMzlBEnetFullCloseOos}{}\renewcommand{\rhcaMzlBEnetFullCloseOos}{0.92}
\providecommand{\rhcaMzlBseEnetFullCloseOos}{}\renewcommand{\rhcaMzlBseEnetFullCloseOos}{0.01}
\providecommand{\rhcaMzlRsqEnetFullCloseOos}{}\renewcommand{\rhcaMzlRsqEnetFullCloseOos}{0.84}
\providecommand{\rhcaShBiasEnetFullCloseOos}{}\renewcommand{\rhcaShBiasEnetFullCloseOos}{0.0}
\providecommand{\rhcaShSlopeEnetFullCloseOos}{}\renewcommand{\rhcaShSlopeEnetFullCloseOos}{39.7}
\providecommand{\rhcaShDistEnetFullCloseOos}{}\renewcommand{\rhcaShDistEnetFullCloseOos}{60.3}
\providecommand{\rhcaQlEnetFullCloseOos}{}\renewcommand{\rhcaQlEnetFullCloseOos}{0.0123}
\providecommand{\rhcaQdEnetFullCloseOos}{}\renewcommand{\rhcaQdEnetFullCloseOos}{0.1122}
\providecommand{\rhcaOverSeEnetFullCloseOos}{}\renewcommand{\rhcaOverSeEnetFullCloseOos}{59.7}
\providecommand{\rhcaMrLgbmHarCloseOos}{}\renewcommand{\rhcaMrLgbmHarCloseOos}{1.035}
\providecommand{\rhcaRmLgbmHarCloseOos}{}\renewcommand{\rhcaRmLgbmHarCloseOos}{1.028}
\providecommand{\rhcaUnderLgbmHarCloseOos}{}\renewcommand{\rhcaUnderLgbmHarCloseOos}{42.9}
\providecommand{\rhcaMzALgbmHarCloseOos}{}\renewcommand{\rhcaMzALgbmHarCloseOos}{\ensuremath{-}0.68}
\providecommand{\rhcaMzBLgbmHarCloseOos}{}\renewcommand{\rhcaMzBLgbmHarCloseOos}{1.08}
\providecommand{\rhcaMzBseLgbmHarCloseOos}{}\renewcommand{\rhcaMzBseLgbmHarCloseOos}{0.06}
\providecommand{\rhcaMzRsqLgbmHarCloseOos}{}\renewcommand{\rhcaMzRsqLgbmHarCloseOos}{0.81}
\providecommand{\rhcaMzlALgbmHarCloseOos}{}\renewcommand{\rhcaMzlALgbmHarCloseOos}{\ensuremath{-}0.33}
\providecommand{\rhcaMzlBLgbmHarCloseOos}{}\renewcommand{\rhcaMzlBLgbmHarCloseOos}{0.98}
\providecommand{\rhcaMzlBseLgbmHarCloseOos}{}\renewcommand{\rhcaMzlBseLgbmHarCloseOos}{0.02}
\providecommand{\rhcaMzlRsqLgbmHarCloseOos}{}\renewcommand{\rhcaMzlRsqLgbmHarCloseOos}{0.84}
\providecommand{\rhcaShBiasLgbmHarCloseOos}{}\renewcommand{\rhcaShBiasLgbmHarCloseOos}{0.0}
\providecommand{\rhcaShSlopeLgbmHarCloseOos}{}\renewcommand{\rhcaShSlopeLgbmHarCloseOos}{2.4}
\providecommand{\rhcaShDistLgbmHarCloseOos}{}\renewcommand{\rhcaShDistLgbmHarCloseOos}{97.6}
\providecommand{\rhcaQlLgbmHarCloseOos}{}\renewcommand{\rhcaQlLgbmHarCloseOos}{0.0006}
\providecommand{\rhcaQdLgbmHarCloseOos}{}\renewcommand{\rhcaQdLgbmHarCloseOos}{0.1095}
\providecommand{\rhcaOverSeLgbmHarCloseOos}{}\renewcommand{\rhcaOverSeLgbmHarCloseOos}{22.8}
\providecommand{\rhcaMrLgbmFullCloseOos}{}\renewcommand{\rhcaMrLgbmFullCloseOos}{1.093}
\providecommand{\rhcaRmLgbmFullCloseOos}{}\renewcommand{\rhcaRmLgbmFullCloseOos}{0.969}
\providecommand{\rhcaUnderLgbmFullCloseOos}{}\renewcommand{\rhcaUnderLgbmFullCloseOos}{47.1}
\providecommand{\rhcaMzALgbmFullCloseOos}{}\renewcommand{\rhcaMzALgbmFullCloseOos}{2.03}
\providecommand{\rhcaMzBLgbmFullCloseOos}{}\renewcommand{\rhcaMzBLgbmFullCloseOos}{0.81}
\providecommand{\rhcaMzBseLgbmFullCloseOos}{}\renewcommand{\rhcaMzBseLgbmFullCloseOos}{0.05}
\providecommand{\rhcaMzRsqLgbmFullCloseOos}{}\renewcommand{\rhcaMzRsqLgbmFullCloseOos}{0.77}
\providecommand{\rhcaMzlALgbmFullCloseOos}{}\renewcommand{\rhcaMzlALgbmFullCloseOos}{\ensuremath{-}0.69}
\providecommand{\rhcaMzlBLgbmFullCloseOos}{}\renewcommand{\rhcaMzlBLgbmFullCloseOos}{0.95}
\providecommand{\rhcaMzlBseLgbmFullCloseOos}{}\renewcommand{\rhcaMzlBseLgbmFullCloseOos}{0.01}
\providecommand{\rhcaMzlRsqLgbmFullCloseOos}{}\renewcommand{\rhcaMzlRsqLgbmFullCloseOos}{0.85}
\providecommand{\rhcaShBiasLgbmFullCloseOos}{}\renewcommand{\rhcaShBiasLgbmFullCloseOos}{0.0}
\providecommand{\rhcaShSlopeLgbmFullCloseOos}{}\renewcommand{\rhcaShSlopeLgbmFullCloseOos}{14.8}
\providecommand{\rhcaShDistLgbmFullCloseOos}{}\renewcommand{\rhcaShDistLgbmFullCloseOos}{85.2}
\providecommand{\rhcaQlLgbmFullCloseOos}{}\renewcommand{\rhcaQlLgbmFullCloseOos}{0.0041}
\providecommand{\rhcaQdLgbmFullCloseOos}{}\renewcommand{\rhcaQdLgbmFullCloseOos}{0.1045}
\providecommand{\rhcaOverSeLgbmFullCloseOos}{}\renewcommand{\rhcaOverSeLgbmFullCloseOos}{44.8}
\providecommand{\rhcaMrXgbHarCloseOos}{}\renewcommand{\rhcaMrXgbHarCloseOos}{1.033}
\providecommand{\rhcaRmXgbHarCloseOos}{}\renewcommand{\rhcaRmXgbHarCloseOos}{1.029}
\providecommand{\rhcaUnderXgbHarCloseOos}{}\renewcommand{\rhcaUnderXgbHarCloseOos}{43.2}
\providecommand{\rhcaMzAXgbHarCloseOos}{}\renewcommand{\rhcaMzAXgbHarCloseOos}{\ensuremath{-}0.70}
\providecommand{\rhcaMzBXgbHarCloseOos}{}\renewcommand{\rhcaMzBXgbHarCloseOos}{1.08}
\providecommand{\rhcaMzBseXgbHarCloseOos}{}\renewcommand{\rhcaMzBseXgbHarCloseOos}{0.06}
\providecommand{\rhcaMzRsqXgbHarCloseOos}{}\renewcommand{\rhcaMzRsqXgbHarCloseOos}{0.81}
\providecommand{\rhcaMzlAXgbHarCloseOos}{}\renewcommand{\rhcaMzlAXgbHarCloseOos}{\ensuremath{-}0.30}
\providecommand{\rhcaMzlBXgbHarCloseOos}{}\renewcommand{\rhcaMzlBXgbHarCloseOos}{0.98}
\providecommand{\rhcaMzlBseXgbHarCloseOos}{}\renewcommand{\rhcaMzlBseXgbHarCloseOos}{0.02}
\providecommand{\rhcaMzlRsqXgbHarCloseOos}{}\renewcommand{\rhcaMzlRsqXgbHarCloseOos}{0.84}
\providecommand{\rhcaShBiasXgbHarCloseOos}{}\renewcommand{\rhcaShBiasXgbHarCloseOos}{0.0}
\providecommand{\rhcaShSlopeXgbHarCloseOos}{}\renewcommand{\rhcaShSlopeXgbHarCloseOos}{2.6}
\providecommand{\rhcaShDistXgbHarCloseOos}{}\renewcommand{\rhcaShDistXgbHarCloseOos}{97.4}
\providecommand{\rhcaQlXgbHarCloseOos}{}\renewcommand{\rhcaQlXgbHarCloseOos}{0.0005}
\providecommand{\rhcaQdXgbHarCloseOos}{}\renewcommand{\rhcaQdXgbHarCloseOos}{0.1096}
\providecommand{\rhcaOverSeXgbHarCloseOos}{}\renewcommand{\rhcaOverSeXgbHarCloseOos}{22.7}
\providecommand{\rhcaMrXgbFullCloseOos}{}\renewcommand{\rhcaMrXgbFullCloseOos}{1.091}
\providecommand{\rhcaRmXgbFullCloseOos}{}\renewcommand{\rhcaRmXgbFullCloseOos}{0.965}
\providecommand{\rhcaUnderXgbFullCloseOos}{}\renewcommand{\rhcaUnderXgbFullCloseOos}{47.9}
\providecommand{\rhcaMzAXgbFullCloseOos}{}\renewcommand{\rhcaMzAXgbFullCloseOos}{1.80}
\providecommand{\rhcaMzBXgbFullCloseOos}{}\renewcommand{\rhcaMzBXgbFullCloseOos}{0.83}
\providecommand{\rhcaMzBseXgbFullCloseOos}{}\renewcommand{\rhcaMzBseXgbFullCloseOos}{0.04}
\providecommand{\rhcaMzRsqXgbFullCloseOos}{}\renewcommand{\rhcaMzRsqXgbFullCloseOos}{0.78}
\providecommand{\rhcaMzlAXgbFullCloseOos}{}\renewcommand{\rhcaMzlAXgbFullCloseOos}{\ensuremath{-}0.68}
\providecommand{\rhcaMzlBXgbFullCloseOos}{}\renewcommand{\rhcaMzlBXgbFullCloseOos}{0.95}
\providecommand{\rhcaMzlBseXgbFullCloseOos}{}\renewcommand{\rhcaMzlBseXgbFullCloseOos}{0.01}
\providecommand{\rhcaMzlRsqXgbFullCloseOos}{}\renewcommand{\rhcaMzlRsqXgbFullCloseOos}{0.85}
\providecommand{\rhcaShBiasXgbFullCloseOos}{}\renewcommand{\rhcaShBiasXgbFullCloseOos}{0.0}
\providecommand{\rhcaShSlopeXgbFullCloseOos}{}\renewcommand{\rhcaShSlopeXgbFullCloseOos}{13.1}
\providecommand{\rhcaShDistXgbFullCloseOos}{}\renewcommand{\rhcaShDistXgbFullCloseOos}{86.8}
\providecommand{\rhcaQlXgbFullCloseOos}{}\renewcommand{\rhcaQlXgbFullCloseOos}{0.0039}
\providecommand{\rhcaQdXgbFullCloseOos}{}\renewcommand{\rhcaQdXgbFullCloseOos}{0.1044}
\providecommand{\rhcaOverSeXgbFullCloseOos}{}\renewcommand{\rhcaOverSeXgbFullCloseOos}{46.1}
\providecommand{\rhcaMrRfHarCloseOos}{}\renewcommand{\rhcaMrRfHarCloseOos}{1.063}
\providecommand{\rhcaRmRfHarCloseOos}{}\renewcommand{\rhcaRmRfHarCloseOos}{1.033}
\providecommand{\rhcaUnderRfHarCloseOos}{}\renewcommand{\rhcaUnderRfHarCloseOos}{45.5}
\providecommand{\rhcaMzARfHarCloseOos}{}\renewcommand{\rhcaMzARfHarCloseOos}{\ensuremath{-}0.38}
\providecommand{\rhcaMzBRfHarCloseOos}{}\renewcommand{\rhcaMzBRfHarCloseOos}{1.06}
\providecommand{\rhcaMzBseRfHarCloseOos}{}\renewcommand{\rhcaMzBseRfHarCloseOos}{0.07}
\providecommand{\rhcaMzRsqRfHarCloseOos}{}\renewcommand{\rhcaMzRsqRfHarCloseOos}{0.82}
\providecommand{\rhcaMzlARfHarCloseOos}{}\renewcommand{\rhcaMzlARfHarCloseOos}{\ensuremath{-}0.46}
\providecommand{\rhcaMzlBRfHarCloseOos}{}\renewcommand{\rhcaMzlBRfHarCloseOos}{0.97}
\providecommand{\rhcaMzlBseRfHarCloseOos}{}\renewcommand{\rhcaMzlBseRfHarCloseOos}{0.02}
\providecommand{\rhcaMzlRsqRfHarCloseOos}{}\renewcommand{\rhcaMzlRsqRfHarCloseOos}{0.83}
\providecommand{\rhcaShBiasRfHarCloseOos}{}\renewcommand{\rhcaShBiasRfHarCloseOos}{0.1}
\providecommand{\rhcaShSlopeRfHarCloseOos}{}\renewcommand{\rhcaShSlopeRfHarCloseOos}{1.7}
\providecommand{\rhcaShDistRfHarCloseOos}{}\renewcommand{\rhcaShDistRfHarCloseOos}{98.3}
\providecommand{\rhcaQlRfHarCloseOos}{}\renewcommand{\rhcaQlRfHarCloseOos}{0.0019}
\providecommand{\rhcaQdRfHarCloseOos}{}\renewcommand{\rhcaQdRfHarCloseOos}{0.1112}
\providecommand{\rhcaOverSeRfHarCloseOos}{}\renewcommand{\rhcaOverSeRfHarCloseOos}{25.7}
\providecommand{\rhcaMrRfFullCloseOos}{}\renewcommand{\rhcaMrRfFullCloseOos}{1.113}
\providecommand{\rhcaRmRfFullCloseOos}{}\renewcommand{\rhcaRmRfFullCloseOos}{0.978}
\providecommand{\rhcaUnderRfFullCloseOos}{}\renewcommand{\rhcaUnderRfFullCloseOos}{50.1}
\providecommand{\rhcaMzARfFullCloseOos}{}\renewcommand{\rhcaMzARfFullCloseOos}{1.15}
\providecommand{\rhcaMzBRfFullCloseOos}{}\renewcommand{\rhcaMzBRfFullCloseOos}{0.89}
\providecommand{\rhcaMzBseRfFullCloseOos}{}\renewcommand{\rhcaMzBseRfFullCloseOos}{0.06}
\providecommand{\rhcaMzRsqRfFullCloseOos}{}\renewcommand{\rhcaMzRsqRfFullCloseOos}{0.79}
\providecommand{\rhcaMzlARfFullCloseOos}{}\renewcommand{\rhcaMzlARfFullCloseOos}{\ensuremath{-}0.87}
\providecommand{\rhcaMzlBRfFullCloseOos}{}\renewcommand{\rhcaMzlBRfFullCloseOos}{0.93}
\providecommand{\rhcaMzlBseRfFullCloseOos}{}\renewcommand{\rhcaMzlBseRfFullCloseOos}{0.01}
\providecommand{\rhcaMzlRsqRfFullCloseOos}{}\renewcommand{\rhcaMzlRsqRfFullCloseOos}{0.85}
\providecommand{\rhcaShBiasRfFullCloseOos}{}\renewcommand{\rhcaShBiasRfFullCloseOos}{0.0}
\providecommand{\rhcaShSlopeRfFullCloseOos}{}\renewcommand{\rhcaShSlopeRfFullCloseOos}{5.5}
\providecommand{\rhcaShDistRfFullCloseOos}{}\renewcommand{\rhcaShDistRfFullCloseOos}{94.5}
\providecommand{\rhcaQlRfFullCloseOos}{}\renewcommand{\rhcaQlRfFullCloseOos}{0.0059}
\providecommand{\rhcaQdRfFullCloseOos}{}\renewcommand{\rhcaQdRfFullCloseOos}{0.1064}
\providecommand{\rhcaOverSeRfFullCloseOos}{}\renewcommand{\rhcaOverSeRfFullCloseOos}{43.1}
% --- aggregates for the prose
\providecommand{\rhcaNLowerCloseSpan}{}\renewcommand{\rhcaNLowerCloseSpan}{3}  % classes (of 6) with the lower 16:00 QLIKE on the full set, Span
\providecommand{\rhcaNLowerSigCloseSpan}{}\renewcommand{\rhcaNLowerSigCloseSpan}{1}  % ... and DM p < 0.05
\providecommand{\rhcaNHigherCloseSpan}{}\renewcommand{\rhcaNHigherCloseSpan}{3}  % classes with the higher 16:00 QLIKE on the full set
\providecommand{\rhcaNHigherSigCloseSpan}{}\renewcommand{\rhcaNHigherSigCloseSpan}{0}  % ... and DM p < 0.05
\providecommand{\rhcaNLowerCloseWordSpan}{}\renewcommand{\rhcaNLowerCloseWordSpan}{three}  % as a word
\providecommand{\rhcaNLinHigherSigCloseSpan}{}\renewcommand{\rhcaNLinHigherSigCloseSpan}{0}  % linear classes (of 3): full-set 16:00 QLIKE significantly higher
\providecommand{\rhcaNLinLowerCloseSpan}{}\renewcommand{\rhcaNLinLowerCloseSpan}{0}  % linear classes with a lower 16:00 QLIKE
\providecommand{\rhcaNTreeLowerSigCloseSpan}{}\renewcommand{\rhcaNTreeLowerSigCloseSpan}{1}  % tree classes (of 3): full-set 16:00 QLIKE significantly lower
\providecommand{\rhcaLinPctCloseSpanMin}{}\renewcommand{\rhcaLinPctCloseSpanMin}{0.6}  % linear 16:00 QLIKE %, full vs HAR, Span: min
\providecommand{\rhcaLinPctCloseSpanMax}{}\renewcommand{\rhcaLinPctCloseSpanMax}{4.2}  % ... max
\providecommand{\rhcaLinPctCloseSpanRng}{}\renewcommand{\rhcaLinPctCloseSpanRng}{0.6--4.2}  % ... the range as text
\providecommand{\rhcaTreePctCloseSpanMin}{}\renewcommand{\rhcaTreePctCloseSpanMin}{\ensuremath{-}3.1}  % tree 16:00 QLIKE %, full vs HAR, Span: min
\providecommand{\rhcaTreePctCloseSpanMax}{}\renewcommand{\rhcaTreePctCloseSpanMax}{\ensuremath{-}1.6}  % ... max
\providecommand{\rhcaTreePctCloseSpanRng}{}\renewcommand{\rhcaTreePctCloseSpanRng}{\ensuremath{-}3.1 to \ensuremath{-}1.6}  % ... the range as text
\providecommand{\rhcaLinMsePctCloseSpanMin}{}\renewcommand{\rhcaLinMsePctCloseSpanMin}{77.5}  % linear 16:00 MSE %, full vs HAR, Span: min
\providecommand{\rhcaLinMsePctCloseSpanMax}{}\renewcommand{\rhcaLinMsePctCloseSpanMax}{99.1}  % ... max
\providecommand{\rhcaLinMsePctCloseSpanRng}{}\renewcommand{\rhcaLinMsePctCloseSpanRng}{77.5--99.1}  % ... the range as text
\providecommand{\rhcaTreeMsePctCloseSpanMin}{}\renewcommand{\rhcaTreeMsePctCloseSpanMin}{19.7}  % tree 16:00 MSE %, full vs HAR, Span: min
\providecommand{\rhcaTreeMsePctCloseSpanMax}{}\renewcommand{\rhcaTreeMsePctCloseSpanMax}{33.0}  % ... max
\providecommand{\rhcaTreeMsePctCloseSpanRng}{}\renewcommand{\rhcaTreeMsePctCloseSpanRng}{19.7--33.0}  % ... the range as text
\providecommand{\rhcaNMseHigherCloseSpan}{}\renewcommand{\rhcaNMseHigherCloseSpan}{6}  % classes with a higher 16:00 MSE, full
\providecommand{\rhcaNMseHigherSigCloseSpan}{}\renewcommand{\rhcaNMseHigherSigCloseSpan}{0}  % ... significantly
\providecommand{\rhcaNLowerClosePre}{}\renewcommand{\rhcaNLowerClosePre}{6}  % classes (of 6) with the lower 16:00 QLIKE on the full set, Pre
\providecommand{\rhcaNLowerSigClosePre}{}\renewcommand{\rhcaNLowerSigClosePre}{2}  % ... and DM p < 0.05
\providecommand{\rhcaNHigherClosePre}{}\renewcommand{\rhcaNHigherClosePre}{0}  % classes with the higher 16:00 QLIKE on the full set
\providecommand{\rhcaNHigherSigClosePre}{}\renewcommand{\rhcaNHigherSigClosePre}{0}  % ... and DM p < 0.05
\providecommand{\rhcaNLowerCloseWordPre}{}\renewcommand{\rhcaNLowerCloseWordPre}{six}  % as a word
\providecommand{\rhcaNLinHigherSigClosePre}{}\renewcommand{\rhcaNLinHigherSigClosePre}{0}  % linear classes (of 3): full-set 16:00 QLIKE significantly higher
\providecommand{\rhcaNLinLowerClosePre}{}\renewcommand{\rhcaNLinLowerClosePre}{3}  % linear classes with a lower 16:00 QLIKE
\providecommand{\rhcaNTreeLowerSigClosePre}{}\renewcommand{\rhcaNTreeLowerSigClosePre}{2}  % tree classes (of 3): full-set 16:00 QLIKE significantly lower
\providecommand{\rhcaLinPctClosePreMin}{}\renewcommand{\rhcaLinPctClosePreMin}{\ensuremath{-}4.8}  % linear 16:00 QLIKE %, full vs HAR, Pre: min
\providecommand{\rhcaLinPctClosePreMax}{}\renewcommand{\rhcaLinPctClosePreMax}{\ensuremath{-}2.3}  % ... max
\providecommand{\rhcaLinPctClosePreRng}{}\renewcommand{\rhcaLinPctClosePreRng}{\ensuremath{-}4.8 to \ensuremath{-}2.3}  % ... the range as text
\providecommand{\rhcaTreePctClosePreMin}{}\renewcommand{\rhcaTreePctClosePreMin}{\ensuremath{-}4.0}  % tree 16:00 QLIKE %, full vs HAR, Pre: min
\providecommand{\rhcaTreePctClosePreMax}{}\renewcommand{\rhcaTreePctClosePreMax}{\ensuremath{-}2.1}  % ... max
\providecommand{\rhcaTreePctClosePreRng}{}\renewcommand{\rhcaTreePctClosePreRng}{\ensuremath{-}4.0 to \ensuremath{-}2.1}  % ... the range as text
\providecommand{\rhcaLinMsePctClosePreMin}{}\renewcommand{\rhcaLinMsePctClosePreMin}{\ensuremath{-}15.8}  % linear 16:00 MSE %, full vs HAR, Pre: min
\providecommand{\rhcaLinMsePctClosePreMax}{}\renewcommand{\rhcaLinMsePctClosePreMax}{\ensuremath{-}13.1}  % ... max
\providecommand{\rhcaLinMsePctClosePreRng}{}\renewcommand{\rhcaLinMsePctClosePreRng}{\ensuremath{-}15.8 to \ensuremath{-}13.1}  % ... the range as text
\providecommand{\rhcaTreeMsePctClosePreMin}{}\renewcommand{\rhcaTreeMsePctClosePreMin}{\ensuremath{-}10.0}  % tree 16:00 MSE %, full vs HAR, Pre: min
\providecommand{\rhcaTreeMsePctClosePreMax}{}\renewcommand{\rhcaTreeMsePctClosePreMax}{\ensuremath{-}5.3}  % ... max
\providecommand{\rhcaTreeMsePctClosePreRng}{}\renewcommand{\rhcaTreeMsePctClosePreRng}{\ensuremath{-}10.0 to \ensuremath{-}5.3}  % ... the range as text
\providecommand{\rhcaNMseHigherClosePre}{}\renewcommand{\rhcaNMseHigherClosePre}{0}  % classes with a higher 16:00 MSE, full
\providecommand{\rhcaNMseHigherSigClosePre}{}\renewcommand{\rhcaNMseHigherSigClosePre}{0}  % ... significantly
\providecommand{\rhcaNLowerCloseOos}{}\renewcommand{\rhcaNLowerCloseOos}{3}  % classes (of 6) with the lower 16:00 QLIKE on the full set, Oos
\providecommand{\rhcaNLowerSigCloseOos}{}\renewcommand{\rhcaNLowerSigCloseOos}{0}  % ... and DM p < 0.05
\providecommand{\rhcaNHigherCloseOos}{}\renewcommand{\rhcaNHigherCloseOos}{3}  % classes with the higher 16:00 QLIKE on the full set
\providecommand{\rhcaNHigherSigCloseOos}{}\renewcommand{\rhcaNHigherSigCloseOos}{3}  % ... and DM p < 0.05
\providecommand{\rhcaNLowerCloseWordOos}{}\renewcommand{\rhcaNLowerCloseWordOos}{three}  % as a word
\providecommand{\rhcaNLinHigherSigCloseOos}{}\renewcommand{\rhcaNLinHigherSigCloseOos}{3}  % linear classes (of 3): full-set 16:00 QLIKE significantly higher
\providecommand{\rhcaNLinLowerCloseOos}{}\renewcommand{\rhcaNLinLowerCloseOos}{0}  % linear classes with a lower 16:00 QLIKE
\providecommand{\rhcaNTreeLowerSigCloseOos}{}\renewcommand{\rhcaNTreeLowerSigCloseOos}{0}  % tree classes (of 3): full-set 16:00 QLIKE significantly lower
\providecommand{\rhcaLinPctCloseOosMin}{}\renewcommand{\rhcaLinPctCloseOosMin}{9.5}  % linear 16:00 QLIKE %, full vs HAR, Oos: min
\providecommand{\rhcaLinPctCloseOosMax}{}\renewcommand{\rhcaLinPctCloseOosMax}{15.2}  % ... max
\providecommand{\rhcaLinPctCloseOosRng}{}\renewcommand{\rhcaLinPctCloseOosRng}{9.5--15.2}  % ... the range as text
\providecommand{\rhcaTreePctCloseOosMin}{}\renewcommand{\rhcaTreePctCloseOosMin}{\ensuremath{-}1.7}  % tree 16:00 QLIKE %, full vs HAR, Oos: min
\providecommand{\rhcaTreePctCloseOosMax}{}\renewcommand{\rhcaTreePctCloseOosMax}{\ensuremath{-}0.7}  % ... max
\providecommand{\rhcaTreePctCloseOosRng}{}\renewcommand{\rhcaTreePctCloseOosRng}{\ensuremath{-}1.7 to \ensuremath{-}0.7}  % ... the range as text
\providecommand{\rhcaLinMsePctCloseOosMin}{}\renewcommand{\rhcaLinMsePctCloseOosMin}{87.0}  % linear 16:00 MSE %, full vs HAR, Oos: min
\providecommand{\rhcaLinMsePctCloseOosMax}{}\renewcommand{\rhcaLinMsePctCloseOosMax}{111.1}  % ... max
\providecommand{\rhcaLinMsePctCloseOosRng}{}\renewcommand{\rhcaLinMsePctCloseOosRng}{87.0--111.1}  % ... the range as text
\providecommand{\rhcaTreeMsePctCloseOosMin}{}\renewcommand{\rhcaTreeMsePctCloseOosMin}{23.3}  % tree 16:00 MSE %, full vs HAR, Oos: min
\providecommand{\rhcaTreeMsePctCloseOosMax}{}\renewcommand{\rhcaTreeMsePctCloseOosMax}{38.9}  % ... max
\providecommand{\rhcaTreeMsePctCloseOosRng}{}\renewcommand{\rhcaTreeMsePctCloseOosRng}{23.3--38.9}  % ... the range as text
\providecommand{\rhcaNMseHigherCloseOos}{}\renewcommand{\rhcaNMseHigherCloseOos}{6}  % classes with a higher 16:00 MSE, full
\providecommand{\rhcaNMseHigherSigCloseOos}{}\renewcommand{\rhcaNMseHigherSigCloseOos}{0}  % ... significantly
\providecommand{\rhcaNMseHigherBarSpan}{}\renewcommand{\rhcaNMseHigherBarSpan}{6}  % classes with a higher every-bar MSE, full, Span
\providecommand{\rhcaNMseHigherSigBarSpan}{}\renewcommand{\rhcaNMseHigherSigBarSpan}{0}  % ... significantly
\providecommand{\rhcaNMseLowerBarSpan}{}\renewcommand{\rhcaNMseLowerBarSpan}{0}  % classes with a lower every-bar MSE, full
\providecommand{\rhcaNMseLowerSigBarSpan}{}\renewcommand{\rhcaNMseLowerSigBarSpan}{0}  % ... significantly
\providecommand{\rhcaLinMsePctBarSpanMin}{}\renewcommand{\rhcaLinMsePctBarSpanMin}{100.6}  % linear every-bar MSE %, full vs HAR, Span: min
\providecommand{\rhcaLinMsePctBarSpanMax}{}\renewcommand{\rhcaLinMsePctBarSpanMax}{117.7}  % ... max
\providecommand{\rhcaLinMsePctBarSpanRng}{}\renewcommand{\rhcaLinMsePctBarSpanRng}{100.6--117.7}  % ... the range as text
\providecommand{\rhcaTreeMsePctBarSpanMin}{}\renewcommand{\rhcaTreeMsePctBarSpanMin}{10.9}  % tree every-bar MSE %, full vs HAR, Span: min
\providecommand{\rhcaTreeMsePctBarSpanMax}{}\renewcommand{\rhcaTreeMsePctBarSpanMax}{20.0}  % ... max
\providecommand{\rhcaTreeMsePctBarSpanRng}{}\renewcommand{\rhcaTreeMsePctBarSpanRng}{10.9--20.0}  % ... the range as text
\providecommand{\rhcaNMseHigherBarPre}{}\renewcommand{\rhcaNMseHigherBarPre}{0}  % classes with a higher every-bar MSE, full, Pre
\providecommand{\rhcaNMseHigherSigBarPre}{}\renewcommand{\rhcaNMseHigherSigBarPre}{0}  % ... significantly
\providecommand{\rhcaNMseLowerBarPre}{}\renewcommand{\rhcaNMseLowerBarPre}{6}  % classes with a lower every-bar MSE, full
\providecommand{\rhcaNMseLowerSigBarPre}{}\renewcommand{\rhcaNMseLowerSigBarPre}{2}  % ... significantly
\providecommand{\rhcaLinMsePctBarPreMin}{}\renewcommand{\rhcaLinMsePctBarPreMin}{\ensuremath{-}10.8}  % linear every-bar MSE %, full vs HAR, Pre: min
\providecommand{\rhcaLinMsePctBarPreMax}{}\renewcommand{\rhcaLinMsePctBarPreMax}{\ensuremath{-}8.9}  % ... max
\providecommand{\rhcaLinMsePctBarPreRng}{}\renewcommand{\rhcaLinMsePctBarPreRng}{\ensuremath{-}10.8 to \ensuremath{-}8.9}  % ... the range as text
\providecommand{\rhcaTreeMsePctBarPreMin}{}\renewcommand{\rhcaTreeMsePctBarPreMin}{\ensuremath{-}4.7}  % tree every-bar MSE %, full vs HAR, Pre: min
\providecommand{\rhcaTreeMsePctBarPreMax}{}\renewcommand{\rhcaTreeMsePctBarPreMax}{\ensuremath{-}1.6}  % ... max
\providecommand{\rhcaTreeMsePctBarPreRng}{}\renewcommand{\rhcaTreeMsePctBarPreRng}{\ensuremath{-}4.7 to \ensuremath{-}1.6}  % ... the range as text
\providecommand{\rhcaNMseHigherBarOos}{}\renewcommand{\rhcaNMseHigherBarOos}{6}  % classes with a higher every-bar MSE, full, Oos
\providecommand{\rhcaNMseHigherSigBarOos}{}\renewcommand{\rhcaNMseHigherSigBarOos}{0}  % ... significantly
\providecommand{\rhcaNMseLowerBarOos}{}\renewcommand{\rhcaNMseLowerBarOos}{0}  % classes with a lower every-bar MSE, full
\providecommand{\rhcaNMseLowerSigBarOos}{}\renewcommand{\rhcaNMseLowerSigBarOos}{0}  % ... significantly
\providecommand{\rhcaLinMsePctBarOosMin}{}\renewcommand{\rhcaLinMsePctBarOosMin}{121.6}  % linear every-bar MSE %, full vs HAR, Oos: min
\providecommand{\rhcaLinMsePctBarOosMax}{}\renewcommand{\rhcaLinMsePctBarOosMax}{142.1}  % ... max
\providecommand{\rhcaLinMsePctBarOosRng}{}\renewcommand{\rhcaLinMsePctBarOosRng}{121.6--142.1}  % ... the range as text
\providecommand{\rhcaTreeMsePctBarOosMin}{}\renewcommand{\rhcaTreeMsePctBarOosMin}{13.7}  % tree every-bar MSE %, full vs HAR, Oos: min
\providecommand{\rhcaTreeMsePctBarOosMax}{}\renewcommand{\rhcaTreeMsePctBarOosMax}{25.5}  % ... max
\providecommand{\rhcaTreeMsePctBarOosRng}{}\renewcommand{\rhcaTreeMsePctBarOosRng}{13.7--25.5}  % ... the range as text
\providecommand{\rhcaNCiLowerBm}{}\renewcommand{\rhcaNCiLowerBm}{2}  % closing-interval models (of 12) below the benchmark's QLIKE
\providecommand{\rhcaNCiLowerSigBm}{}\renewcommand{\rhcaNCiLowerSigBm}{0}  % ... significantly
\providecommand{\rhcaNCi}{}\renewcommand{\rhcaNCi}{12}  % closing-interval models other than the benchmark
\providecommand{\rhcaCiHarGapBound}{}\renewcommand{\rhcaCiHarGapBound}{0.8}  % |QLIKE %| of the closing-interval HAR-set linear models vs the benchmark is below
\providecommand{\rhcaLinFullMzBBarPreMin}{}\renewcommand{\rhcaLinFullMzBBarPreMin}{0.99}  % levels MZ slope, linear full, BarPre: min
\providecommand{\rhcaLinFullMzBBarPreMax}{}\renewcommand{\rhcaLinFullMzBBarPreMax}{1.00}  % ... max
\providecommand{\rhcaLinFullMzBBarPreRng}{}\renewcommand{\rhcaLinFullMzBBarPreRng}{0.99--1.00}  % ... the range as text
\providecommand{\rhcaTreeFullMzBBarPreMin}{}\renewcommand{\rhcaTreeFullMzBBarPreMin}{0.92}  % levels MZ slope, tree full, BarPre: min
\providecommand{\rhcaTreeFullMzBBarPreMax}{}\renewcommand{\rhcaTreeFullMzBBarPreMax}{0.97}  % ... max
\providecommand{\rhcaTreeFullMzBBarPreRng}{}\renewcommand{\rhcaTreeFullMzBBarPreRng}{0.92--0.97}  % ... the range as text
\providecommand{\rhcaTreeHarMzBBarPreMin}{}\renewcommand{\rhcaTreeHarMzBBarPreMin}{0.96}  % levels MZ slope, tree HAR, BarPre: min
\providecommand{\rhcaTreeHarMzBBarPreMax}{}\renewcommand{\rhcaTreeHarMzBBarPreMax}{0.99}  % ... max
\providecommand{\rhcaTreeHarMzBBarPreRng}{}\renewcommand{\rhcaTreeHarMzBBarPreRng}{0.96--0.99}  % ... the range as text
\providecommand{\rhcaLinHarMzBBarPre}{}\renewcommand{\rhcaLinHarMzBBarPre}{0.97}  % levels MZ slope, linear HAR (= OLS), BarPre
\providecommand{\rhcaLinFullMzlBBarPreMin}{}\renewcommand{\rhcaLinFullMzlBBarPreMin}{0.99}  % logs MZ slope, linear full, BarPre: min
\providecommand{\rhcaLinFullMzlBBarPreMax}{}\renewcommand{\rhcaLinFullMzlBBarPreMax}{0.99}  % ... max
\providecommand{\rhcaLinFullMzlBBarPreRng}{}\renewcommand{\rhcaLinFullMzlBBarPreRng}{0.99}  % ... the range as text
\providecommand{\rhcaTreeFullMzlBBarPreMin}{}\renewcommand{\rhcaTreeFullMzlBBarPreMin}{1.00}  % logs MZ slope, tree full, BarPre: min
\providecommand{\rhcaTreeFullMzlBBarPreMax}{}\renewcommand{\rhcaTreeFullMzlBBarPreMax}{1.01}  % ... max
\providecommand{\rhcaTreeFullMzlBBarPreRng}{}\renewcommand{\rhcaTreeFullMzlBBarPreRng}{1.00--1.01}  % ... the range as text
\providecommand{\rhcaAllHarMzlBBarPreMin}{}\renewcommand{\rhcaAllHarMzlBBarPreMin}{1.00}  % logs MZ slope, HAR forecasts, BarPre: min
\providecommand{\rhcaAllHarMzlBBarPreMax}{}\renewcommand{\rhcaAllHarMzlBBarPreMax}{1.01}  % ... max
\providecommand{\rhcaAllHarMzlBBarPreRng}{}\renewcommand{\rhcaAllHarMzlBBarPreRng}{1.00--1.01}  % ... the range as text
\providecommand{\rhcaLinFullMzlRsqBarPreMin}{}\renewcommand{\rhcaLinFullMzlRsqBarPreMin}{0.78}  % logs MZ R2, linear full, BarPre: min
\providecommand{\rhcaLinFullMzlRsqBarPreMax}{}\renewcommand{\rhcaLinFullMzlRsqBarPreMax}{0.78}  % ... max
\providecommand{\rhcaLinFullMzlRsqBarPreRng}{}\renewcommand{\rhcaLinFullMzlRsqBarPreRng}{0.78}  % ... the range as text
\providecommand{\rhcaTreeFullMzlRsqBarPreMin}{}\renewcommand{\rhcaTreeFullMzlRsqBarPreMin}{0.79}  % logs MZ R2, tree full, BarPre: min
\providecommand{\rhcaTreeFullMzlRsqBarPreMax}{}\renewcommand{\rhcaTreeFullMzlRsqBarPreMax}{0.79}  % ... max
\providecommand{\rhcaTreeFullMzlRsqBarPreRng}{}\renewcommand{\rhcaTreeFullMzlRsqBarPreRng}{0.79}  % ... the range as text
\providecommand{\rhcaAllHarMzlRsqBarPreMin}{}\renewcommand{\rhcaAllHarMzlRsqBarPreMin}{0.76}  % logs MZ R2, HAR forecasts, BarPre: min
\providecommand{\rhcaAllHarMzlRsqBarPreMax}{}\renewcommand{\rhcaAllHarMzlRsqBarPreMax}{0.76}  % ... max
\providecommand{\rhcaAllHarMzlRsqBarPreRng}{}\renewcommand{\rhcaAllHarMzlRsqBarPreRng}{0.76}  % ... the range as text
\providecommand{\rhcaLinFullMzRsqBarPreMin}{}\renewcommand{\rhcaLinFullMzRsqBarPreMin}{0.53}  % levels MZ R2, linear full, BarPre: min
\providecommand{\rhcaLinFullMzRsqBarPreMax}{}\renewcommand{\rhcaLinFullMzRsqBarPreMax}{0.54}  % ... max
\providecommand{\rhcaLinFullMzRsqBarPreRng}{}\renewcommand{\rhcaLinFullMzRsqBarPreRng}{0.53--0.54}  % ... the range as text
\providecommand{\rhcaAllHarMzRsqBarPreMin}{}\renewcommand{\rhcaAllHarMzRsqBarPreMin}{0.37}  % levels MZ R2, HAR forecasts, BarPre: min
\providecommand{\rhcaAllHarMzRsqBarPreMax}{}\renewcommand{\rhcaAllHarMzRsqBarPreMax}{0.48}  % ... max
\providecommand{\rhcaAllHarMzRsqBarPreRng}{}\renewcommand{\rhcaAllHarMzRsqBarPreRng}{0.37--0.48}  % ... the range as text
\providecommand{\rhcaLinFullShSlopeBarPreMin}{}\renewcommand{\rhcaLinFullShSlopeBarPreMin}{0.0}  % Theil slope share %, linear full, BarPre: min
\providecommand{\rhcaLinFullShSlopeBarPreMax}{}\renewcommand{\rhcaLinFullShSlopeBarPreMax}{0.0}  % ... max
\providecommand{\rhcaLinFullShSlopeBarPreRng}{}\renewcommand{\rhcaLinFullShSlopeBarPreRng}{0.0}  % ... the range as text
\providecommand{\rhcaTreeFullShSlopeBarPreMin}{}\renewcommand{\rhcaTreeFullShSlopeBarPreMin}{0.1}  % Theil slope share %, tree full, BarPre: min
\providecommand{\rhcaTreeFullShSlopeBarPreMax}{}\renewcommand{\rhcaTreeFullShSlopeBarPreMax}{0.5}  % ... max
\providecommand{\rhcaTreeFullShSlopeBarPreRng}{}\renewcommand{\rhcaTreeFullShSlopeBarPreRng}{0.1--0.5}  % ... the range as text
\providecommand{\rhcaAllHarShSlopeBarPreMin}{}\renewcommand{\rhcaAllHarShSlopeBarPreMin}{0.0}  % Theil slope share %, HAR forecasts, BarPre: min
\providecommand{\rhcaAllHarShSlopeBarPreMax}{}\renewcommand{\rhcaAllHarShSlopeBarPreMax}{0.1}  % ... max
\providecommand{\rhcaAllHarShSlopeBarPreRng}{}\renewcommand{\rhcaAllHarShSlopeBarPreRng}{0.0--0.1}  % ... the range as text
\providecommand{\rhcaAllShBiasBarPreMin}{}\renewcommand{\rhcaAllShBiasBarPreMin}{0.1}  % Theil bias share %, all 12, BarPre: min
\providecommand{\rhcaAllShBiasBarPreMax}{}\renewcommand{\rhcaAllShBiasBarPreMax}{0.1}  % ... max
\providecommand{\rhcaAllShBiasBarPreRng}{}\renewcommand{\rhcaAllShBiasBarPreRng}{0.1}  % ... the range as text
\providecommand{\rhcaAllRmBarPreMin}{}\renewcommand{\rhcaAllRmBarPreMin}{1.04}  % ratio of means, all 12, BarPre: min
\providecommand{\rhcaAllRmBarPreMax}{}\renewcommand{\rhcaAllRmBarPreMax}{1.07}  % ... max
\providecommand{\rhcaAllRmBarPreRng}{}\renewcommand{\rhcaAllRmBarPreRng}{1.04--1.07}  % ... the range as text
\providecommand{\rhcaFullRmBarPreMin}{}\renewcommand{\rhcaFullRmBarPreMin}{1.04}  % ratio of means, full set, BarPre: min
\providecommand{\rhcaFullRmBarPreMax}{}\renewcommand{\rhcaFullRmBarPreMax}{1.06}  % ... max
\providecommand{\rhcaFullRmBarPreRng}{}\renewcommand{\rhcaFullRmBarPreRng}{1.04--1.06}  % ... the range as text
\providecommand{\rhcaHarRmBarPreMin}{}\renewcommand{\rhcaHarRmBarPreMin}{1.06}  % ratio of means, HAR, BarPre: min
\providecommand{\rhcaHarRmBarPreMax}{}\renewcommand{\rhcaHarRmBarPreMax}{1.07}  % ... max
\providecommand{\rhcaHarRmBarPreRng}{}\renewcommand{\rhcaHarRmBarPreRng}{1.06--1.07}  % ... the range as text
\providecommand{\rhcaAllMrBarPreMin}{}\renewcommand{\rhcaAllMrBarPreMin}{1.07}  % mean of RV/sigma2, all 12, BarPre: min
\providecommand{\rhcaAllMrBarPreMax}{}\renewcommand{\rhcaAllMrBarPreMax}{1.10}  % ... max
\providecommand{\rhcaAllMrBarPreRng}{}\renewcommand{\rhcaAllMrBarPreRng}{1.07--1.10}  % ... the range as text
\providecommand{\rhcaAllUnderBarPreMin}{}\renewcommand{\rhcaAllUnderBarPreMin}{39.9}  % under-prediction %, all 12, BarPre: min
\providecommand{\rhcaAllUnderBarPreMax}{}\renewcommand{\rhcaAllUnderBarPreMax}{42.1}  % ... max
\providecommand{\rhcaAllUnderBarPreRng}{}\renewcommand{\rhcaAllUnderBarPreRng}{39.9--42.1}  % ... the range as text
\providecommand{\rhcaLinFullUnderBarPreMin}{}\renewcommand{\rhcaLinFullUnderBarPreMin}{41.5}  % under-prediction %, linear full, BarPre: min
\providecommand{\rhcaLinFullUnderBarPreMax}{}\renewcommand{\rhcaLinFullUnderBarPreMax}{42.1}  % ... max
\providecommand{\rhcaLinFullUnderBarPreRng}{}\renewcommand{\rhcaLinFullUnderBarPreRng}{41.5--42.1}  % ... the range as text
\providecommand{\rhcaLinHarUnderBarPre}{}\renewcommand{\rhcaLinHarUnderBarPre}{39.9}  % under-prediction %, linear HAR, BarPre
\providecommand{\rhcaLinFullOverSeBarPreMin}{}\renewcommand{\rhcaLinFullOverSeBarPreMin}{19.6}  % % of SE from over-predictions, linear full, BarPre: min
\providecommand{\rhcaLinFullOverSeBarPreMax}{}\renewcommand{\rhcaLinFullOverSeBarPreMax}{20.1}  % ... max
\providecommand{\rhcaLinFullOverSeBarPreRng}{}\renewcommand{\rhcaLinFullOverSeBarPreRng}{19.6--20.1}  % ... the range as text
\providecommand{\rhcaTreeFullOverSeBarPreMin}{}\renewcommand{\rhcaTreeFullOverSeBarPreMin}{12.0}  % % of SE from over-predictions, tree full, BarPre: min
\providecommand{\rhcaTreeFullOverSeBarPreMax}{}\renewcommand{\rhcaTreeFullOverSeBarPreMax}{15.1}  % ... max
\providecommand{\rhcaTreeFullOverSeBarPreRng}{}\renewcommand{\rhcaTreeFullOverSeBarPreRng}{12.0--15.1}  % ... the range as text
\providecommand{\rhcaAllHarOverSeBarPreMin}{}\renewcommand{\rhcaAllHarOverSeBarPreMin}{10.3}  % % of SE from over-predictions, HAR, BarPre: min
\providecommand{\rhcaAllHarOverSeBarPreMax}{}\renewcommand{\rhcaAllHarOverSeBarPreMax}{20.0}  % ... max
\providecommand{\rhcaAllHarOverSeBarPreRng}{}\renewcommand{\rhcaAllHarOverSeBarPreRng}{10.3--20.0}  % ... the range as text
\providecommand{\rhcaNMaeLowerBarPre}{}\renewcommand{\rhcaNMaeLowerBarPre}{6}  % classes (of 6) with the lower MAE on the full set, BarPre
\providecommand{\rhcaLinFullMrBarPreMin}{}\renewcommand{\rhcaLinFullMrBarPreMin}{1.091}  % mean of RV/sigma2, linear full, BarPre: min
\providecommand{\rhcaLinFullMrBarPreMax}{}\renewcommand{\rhcaLinFullMrBarPreMax}{1.096}  % ... max
\providecommand{\rhcaLinFullMrBarPreRng}{}\renewcommand{\rhcaLinFullMrBarPreRng}{1.091--1.096}  % ... the range as text
\providecommand{\rhcaLinHarMrBarPre}{}\renewcommand{\rhcaLinHarMrBarPre}{1.078}  % mean of RV/sigma2, linear HAR, BarPre
\providecommand{\rhcaLinFullQlBarPreMin}{}\renewcommand{\rhcaLinFullQlBarPreMin}{0.0039}  % QLIKE level term, linear full, BarPre: min
\providecommand{\rhcaLinFullQlBarPreMax}{}\renewcommand{\rhcaLinFullQlBarPreMax}{0.0043}  % ... max
\providecommand{\rhcaLinFullQlBarPreRng}{}\renewcommand{\rhcaLinFullQlBarPreRng}{0.0039--0.0043}  % ... the range as text
\providecommand{\rhcaLinHarQlBarPre}{}\renewcommand{\rhcaLinHarQlBarPre}{0.0029}  % QLIKE level term, linear HAR, BarPre
\providecommand{\rhcaLinFullQdBarPreMin}{}\renewcommand{\rhcaLinFullQdBarPreMin}{0.1627}  % QLIKE dispersion term, linear full, BarPre: min
\providecommand{\rhcaLinFullQdBarPreMax}{}\renewcommand{\rhcaLinFullQdBarPreMax}{0.1683}  % ... max
\providecommand{\rhcaLinFullQdBarPreRng}{}\renewcommand{\rhcaLinFullQdBarPreRng}{0.1627--0.1683}  % ... the range as text
\providecommand{\rhcaLinHarQdBarPre}{}\renewcommand{\rhcaLinHarQdBarPre}{0.1835}  % QLIKE dispersion term, linear HAR, BarPre
\providecommand{\rhcaTreeFullQlBarPreMin}{}\renewcommand{\rhcaTreeFullQlBarPreMin}{0.0021}  % QLIKE level term, tree full, BarPre: min
\providecommand{\rhcaTreeFullQlBarPreMax}{}\renewcommand{\rhcaTreeFullQlBarPreMax}{0.0029}  % ... max
\providecommand{\rhcaTreeFullQlBarPreRng}{}\renewcommand{\rhcaTreeFullQlBarPreRng}{0.0021--0.0029}  % ... the range as text
\providecommand{\rhcaTreeHarQlBarPreMin}{}\renewcommand{\rhcaTreeHarQlBarPreMin}{0.0032}  % QLIKE level term, tree HAR, BarPre: min
\providecommand{\rhcaTreeHarQlBarPreMax}{}\renewcommand{\rhcaTreeHarQlBarPreMax}{0.0034}  % ... max
\providecommand{\rhcaTreeHarQlBarPreRng}{}\renewcommand{\rhcaTreeHarQlBarPreRng}{0.0032--0.0034}  % ... the range as text
\providecommand{\rhcaTreeFullQdBarPreMin}{}\renewcommand{\rhcaTreeFullQdBarPreMin}{0.1563}  % QLIKE dispersion term, tree full, BarPre: min
\providecommand{\rhcaTreeFullQdBarPreMax}{}\renewcommand{\rhcaTreeFullQdBarPreMax}{0.1617}  % ... max
\providecommand{\rhcaTreeFullQdBarPreRng}{}\renewcommand{\rhcaTreeFullQdBarPreRng}{0.1563--0.1617}  % ... the range as text
\providecommand{\rhcaTreeHarQdBarPreMin}{}\renewcommand{\rhcaTreeHarQdBarPreMin}{0.1826}  % QLIKE dispersion term, tree HAR, BarPre: min
\providecommand{\rhcaTreeHarQdBarPreMax}{}\renewcommand{\rhcaTreeHarQdBarPreMax}{0.1832}  % ... max
\providecommand{\rhcaTreeHarQdBarPreRng}{}\renewcommand{\rhcaTreeHarQdBarPreRng}{0.1826--0.1832}  % ... the range as text
\providecommand{\rhcaTreeFullUnderBarPreMin}{}\renewcommand{\rhcaTreeFullUnderBarPreMin}{40.8}  % under-prediction %, tree full, BarPre: min
\providecommand{\rhcaTreeFullUnderBarPreMax}{}\renewcommand{\rhcaTreeFullUnderBarPreMax}{41.1}  % ... max
\providecommand{\rhcaTreeFullUnderBarPreRng}{}\renewcommand{\rhcaTreeFullUnderBarPreRng}{40.8--41.1}  % ... the range as text
\providecommand{\rhcaTreeHarUnderBarPreMin}{}\renewcommand{\rhcaTreeHarUnderBarPreMin}{40.1}  % under-prediction %, tree HAR, BarPre: min
\providecommand{\rhcaTreeHarUnderBarPreMax}{}\renewcommand{\rhcaTreeHarUnderBarPreMax}{40.2}  % ... max
\providecommand{\rhcaTreeHarUnderBarPreRng}{}\renewcommand{\rhcaTreeHarUnderBarPreRng}{40.1--40.2}  % ... the range as text
\providecommand{\rhcaLinFullMzBBarOosMin}{}\renewcommand{\rhcaLinFullMzBBarOosMin}{0.56}  % levels MZ slope, linear full, BarOos: min
\providecommand{\rhcaLinFullMzBBarOosMax}{}\renewcommand{\rhcaLinFullMzBBarOosMax}{0.58}  % ... max
\providecommand{\rhcaLinFullMzBBarOosRng}{}\renewcommand{\rhcaLinFullMzBBarOosRng}{0.56--0.58}  % ... the range as text
\providecommand{\rhcaTreeFullMzBBarOosMin}{}\renewcommand{\rhcaTreeFullMzBBarOosMin}{0.74}  % levels MZ slope, tree full, BarOos: min
\providecommand{\rhcaTreeFullMzBBarOosMax}{}\renewcommand{\rhcaTreeFullMzBBarOosMax}{0.78}  % ... max
\providecommand{\rhcaTreeFullMzBBarOosRng}{}\renewcommand{\rhcaTreeFullMzBBarOosRng}{0.74--0.78}  % ... the range as text
\providecommand{\rhcaTreeHarMzBBarOosMin}{}\renewcommand{\rhcaTreeHarMzBBarOosMin}{0.92}  % levels MZ slope, tree HAR, BarOos: min
\providecommand{\rhcaTreeHarMzBBarOosMax}{}\renewcommand{\rhcaTreeHarMzBBarOosMax}{0.94}  % ... max
\providecommand{\rhcaTreeHarMzBBarOosRng}{}\renewcommand{\rhcaTreeHarMzBBarOosRng}{0.92--0.94}  % ... the range as text
\providecommand{\rhcaLinHarMzBBarOos}{}\renewcommand{\rhcaLinHarMzBBarOos}{0.99}  % levels MZ slope, linear HAR (= OLS), BarOos
\providecommand{\rhcaLinFullMzlBBarOosMin}{}\renewcommand{\rhcaLinFullMzlBBarOosMin}{0.98}  % logs MZ slope, linear full, BarOos: min
\providecommand{\rhcaLinFullMzlBBarOosMax}{}\renewcommand{\rhcaLinFullMzlBBarOosMax}{0.99}  % ... max
\providecommand{\rhcaLinFullMzlBBarOosRng}{}\renewcommand{\rhcaLinFullMzlBBarOosRng}{0.98--0.99}  % ... the range as text
\providecommand{\rhcaTreeFullMzlBBarOosMin}{}\renewcommand{\rhcaTreeFullMzlBBarOosMin}{0.99}  % logs MZ slope, tree full, BarOos: min
\providecommand{\rhcaTreeFullMzlBBarOosMax}{}\renewcommand{\rhcaTreeFullMzlBBarOosMax}{0.99}  % ... max
\providecommand{\rhcaTreeFullMzlBBarOosRng}{}\renewcommand{\rhcaTreeFullMzlBBarOosRng}{0.99}  % ... the range as text
\providecommand{\rhcaAllHarMzlBBarOosMin}{}\renewcommand{\rhcaAllHarMzlBBarOosMin}{1.00}  % logs MZ slope, HAR forecasts, BarOos: min
\providecommand{\rhcaAllHarMzlBBarOosMax}{}\renewcommand{\rhcaAllHarMzlBBarOosMax}{1.02}  % ... max
\providecommand{\rhcaAllHarMzlBBarOosRng}{}\renewcommand{\rhcaAllHarMzlBBarOosRng}{1.00--1.02}  % ... the range as text
\providecommand{\rhcaLinFullMzlRsqBarOosMin}{}\renewcommand{\rhcaLinFullMzlRsqBarOosMin}{0.81}  % logs MZ R2, linear full, BarOos: min
\providecommand{\rhcaLinFullMzlRsqBarOosMax}{}\renewcommand{\rhcaLinFullMzlRsqBarOosMax}{0.81}  % ... max
\providecommand{\rhcaLinFullMzlRsqBarOosRng}{}\renewcommand{\rhcaLinFullMzlRsqBarOosRng}{0.81}  % ... the range as text
\providecommand{\rhcaTreeFullMzlRsqBarOosMin}{}\renewcommand{\rhcaTreeFullMzlRsqBarOosMin}{0.82}  % logs MZ R2, tree full, BarOos: min
\providecommand{\rhcaTreeFullMzlRsqBarOosMax}{}\renewcommand{\rhcaTreeFullMzlRsqBarOosMax}{0.82}  % ... max
\providecommand{\rhcaTreeFullMzlRsqBarOosRng}{}\renewcommand{\rhcaTreeFullMzlRsqBarOosRng}{0.82}  % ... the range as text
\providecommand{\rhcaAllHarMzlRsqBarOosMin}{}\renewcommand{\rhcaAllHarMzlRsqBarOosMin}{0.79}  % logs MZ R2, HAR forecasts, BarOos: min
\providecommand{\rhcaAllHarMzlRsqBarOosMax}{}\renewcommand{\rhcaAllHarMzlRsqBarOosMax}{0.80}  % ... max
\providecommand{\rhcaAllHarMzlRsqBarOosRng}{}\renewcommand{\rhcaAllHarMzlRsqBarOosRng}{0.79--0.80}  % ... the range as text
\providecommand{\rhcaLinFullMzRsqBarOosMin}{}\renewcommand{\rhcaLinFullMzRsqBarOosMin}{0.51}  % levels MZ R2, linear full, BarOos: min
\providecommand{\rhcaLinFullMzRsqBarOosMax}{}\renewcommand{\rhcaLinFullMzRsqBarOosMax}{0.52}  % ... max
\providecommand{\rhcaLinFullMzRsqBarOosRng}{}\renewcommand{\rhcaLinFullMzRsqBarOosRng}{0.51--0.52}  % ... the range as text
\providecommand{\rhcaAllHarMzRsqBarOosMin}{}\renewcommand{\rhcaAllHarMzRsqBarOosMin}{0.65}  % levels MZ R2, HAR forecasts, BarOos: min
\providecommand{\rhcaAllHarMzRsqBarOosMax}{}\renewcommand{\rhcaAllHarMzRsqBarOosMax}{0.67}  % ... max
\providecommand{\rhcaAllHarMzRsqBarOosRng}{}\renewcommand{\rhcaAllHarMzRsqBarOosRng}{0.65--0.67}  % ... the range as text
\providecommand{\rhcaLinFullShSlopeBarOosMin}{}\renewcommand{\rhcaLinFullShSlopeBarOosMin}{35.9}  % Theil slope share %, linear full, BarOos: min
\providecommand{\rhcaLinFullShSlopeBarOosMax}{}\renewcommand{\rhcaLinFullShSlopeBarOosMax}{40.8}  % ... max
\providecommand{\rhcaLinFullShSlopeBarOosRng}{}\renewcommand{\rhcaLinFullShSlopeBarOosRng}{35.9--40.8}  % ... the range as text
\providecommand{\rhcaTreeFullShSlopeBarOosMin}{}\renewcommand{\rhcaTreeFullShSlopeBarOosMin}{12.7}  % Theil slope share %, tree full, BarOos: min
\providecommand{\rhcaTreeFullShSlopeBarOosMax}{}\renewcommand{\rhcaTreeFullShSlopeBarOosMax}{17.7}  % ... max
\providecommand{\rhcaTreeFullShSlopeBarOosRng}{}\renewcommand{\rhcaTreeFullShSlopeBarOosRng}{12.7--17.7}  % ... the range as text
\providecommand{\rhcaAllHarShSlopeBarOosMin}{}\renewcommand{\rhcaAllHarShSlopeBarOosMin}{0.0}  % Theil slope share %, HAR forecasts, BarOos: min
\providecommand{\rhcaAllHarShSlopeBarOosMax}{}\renewcommand{\rhcaAllHarShSlopeBarOosMax}{1.4}  % ... max
\providecommand{\rhcaAllHarShSlopeBarOosRng}{}\renewcommand{\rhcaAllHarShSlopeBarOosRng}{0.0--1.4}  % ... the range as text
\providecommand{\rhcaAllShBiasBarOosMin}{}\renewcommand{\rhcaAllShBiasBarOosMin}{0.0}  % Theil bias share %, all 12, BarOos: min
\providecommand{\rhcaAllShBiasBarOosMax}{}\renewcommand{\rhcaAllShBiasBarOosMax}{0.1}  % ... max
\providecommand{\rhcaAllShBiasBarOosRng}{}\renewcommand{\rhcaAllShBiasBarOosRng}{0.0--0.1}  % ... the range as text
\providecommand{\rhcaAllRmBarOosMin}{}\renewcommand{\rhcaAllRmBarOosMin}{0.96}  % ratio of means, all 12, BarOos: min
\providecommand{\rhcaAllRmBarOosMax}{}\renewcommand{\rhcaAllRmBarOosMax}{1.05}  % ... max
\providecommand{\rhcaAllRmBarOosRng}{}\renewcommand{\rhcaAllRmBarOosRng}{0.96--1.05}  % ... the range as text
\providecommand{\rhcaFullRmBarOosMin}{}\renewcommand{\rhcaFullRmBarOosMin}{0.96}  % ratio of means, full set, BarOos: min
\providecommand{\rhcaFullRmBarOosMax}{}\renewcommand{\rhcaFullRmBarOosMax}{0.98}  % ... max
\providecommand{\rhcaFullRmBarOosRng}{}\renewcommand{\rhcaFullRmBarOosRng}{0.96--0.98}  % ... the range as text
\providecommand{\rhcaHarRmBarOosMin}{}\renewcommand{\rhcaHarRmBarOosMin}{1.02}  % ratio of means, HAR, BarOos: min
\providecommand{\rhcaHarRmBarOosMax}{}\renewcommand{\rhcaHarRmBarOosMax}{1.05}  % ... max
\providecommand{\rhcaHarRmBarOosRng}{}\renewcommand{\rhcaHarRmBarOosRng}{1.02--1.05}  % ... the range as text
\providecommand{\rhcaAllMrBarOosMin}{}\renewcommand{\rhcaAllMrBarOosMin}{1.06}  % mean of RV/sigma2, all 12, BarOos: min
\providecommand{\rhcaAllMrBarOosMax}{}\renewcommand{\rhcaAllMrBarOosMax}{1.12}  % ... max
\providecommand{\rhcaAllMrBarOosRng}{}\renewcommand{\rhcaAllMrBarOosRng}{1.06--1.12}  % ... the range as text
\providecommand{\rhcaAllUnderBarOosMin}{}\renewcommand{\rhcaAllUnderBarOosMin}{38.0}  % under-prediction %, all 12, BarOos: min
\providecommand{\rhcaAllUnderBarOosMax}{}\renewcommand{\rhcaAllUnderBarOosMax}{43.2}  % ... max
\providecommand{\rhcaAllUnderBarOosRng}{}\renewcommand{\rhcaAllUnderBarOosRng}{38.0--43.2}  % ... the range as text
\providecommand{\rhcaLinFullUnderBarOosMin}{}\renewcommand{\rhcaLinFullUnderBarOosMin}{42.2}  % under-prediction %, linear full, BarOos: min
\providecommand{\rhcaLinFullUnderBarOosMax}{}\renewcommand{\rhcaLinFullUnderBarOosMax}{43.2}  % ... max
\providecommand{\rhcaLinFullUnderBarOosRng}{}\renewcommand{\rhcaLinFullUnderBarOosRng}{42.2--43.2}  % ... the range as text
\providecommand{\rhcaLinHarUnderBarOos}{}\renewcommand{\rhcaLinHarUnderBarOos}{38.0}  % under-prediction %, linear HAR, BarOos
\providecommand{\rhcaLinFullOverSeBarOosMin}{}\renewcommand{\rhcaLinFullOverSeBarOosMin}{63.1}  % % of SE from over-predictions, linear full, BarOos: min
\providecommand{\rhcaLinFullOverSeBarOosMax}{}\renewcommand{\rhcaLinFullOverSeBarOosMax}{67.0}  % ... max
\providecommand{\rhcaLinFullOverSeBarOosRng}{}\renewcommand{\rhcaLinFullOverSeBarOosRng}{63.1--67.0}  % ... the range as text
\providecommand{\rhcaTreeFullOverSeBarOosMin}{}\renewcommand{\rhcaTreeFullOverSeBarOosMin}{44.4}  % % of SE from over-predictions, tree full, BarOos: min
\providecommand{\rhcaTreeFullOverSeBarOosMax}{}\renewcommand{\rhcaTreeFullOverSeBarOosMax}{49.1}  % ... max
\providecommand{\rhcaTreeFullOverSeBarOosRng}{}\renewcommand{\rhcaTreeFullOverSeBarOosRng}{44.4--49.1}  % ... the range as text
\providecommand{\rhcaAllHarOverSeBarOosMin}{}\renewcommand{\rhcaAllHarOverSeBarOosMin}{21.5}  % % of SE from over-predictions, HAR, BarOos: min
\providecommand{\rhcaAllHarOverSeBarOosMax}{}\renewcommand{\rhcaAllHarOverSeBarOosMax}{26.7}  % ... max
\providecommand{\rhcaAllHarOverSeBarOosRng}{}\renewcommand{\rhcaAllHarOverSeBarOosRng}{21.5--26.7}  % ... the range as text
\providecommand{\rhcaNMaeLowerBarOos}{}\renewcommand{\rhcaNMaeLowerBarOos}{0}  % classes (of 6) with the lower MAE on the full set, BarOos
\providecommand{\rhcaLinFullMrBarOosMin}{}\renewcommand{\rhcaLinFullMrBarOosMin}{1.109}  % mean of RV/sigma2, linear full, BarOos: min
\providecommand{\rhcaLinFullMrBarOosMax}{}\renewcommand{\rhcaLinFullMrBarOosMax}{1.117}  % ... max
\providecommand{\rhcaLinFullMrBarOosRng}{}\renewcommand{\rhcaLinFullMrBarOosRng}{1.109--1.117}  % ... the range as text
\providecommand{\rhcaLinHarMrBarOos}{}\renewcommand{\rhcaLinHarMrBarOos}{1.061}  % mean of RV/sigma2, linear HAR, BarOos
\providecommand{\rhcaLinFullQlBarOosMin}{}\renewcommand{\rhcaLinFullQlBarOosMin}{0.0056}  % QLIKE level term, linear full, BarOos: min
\providecommand{\rhcaLinFullQlBarOosMax}{}\renewcommand{\rhcaLinFullQlBarOosMax}{0.0064}  % ... max
\providecommand{\rhcaLinFullQlBarOosRng}{}\renewcommand{\rhcaLinFullQlBarOosRng}{0.0056--0.0064}  % ... the range as text
\providecommand{\rhcaLinHarQlBarOos}{}\renewcommand{\rhcaLinHarQlBarOos}{0.0018}  % QLIKE level term, linear HAR, BarOos
\providecommand{\rhcaLinFullQdBarOosMin}{}\renewcommand{\rhcaLinFullQdBarOosMin}{0.1739}  % QLIKE dispersion term, linear full, BarOos: min
\providecommand{\rhcaLinFullQdBarOosMax}{}\renewcommand{\rhcaLinFullQdBarOosMax}{0.1788}  % ... max
\providecommand{\rhcaLinFullQdBarOosRng}{}\renewcommand{\rhcaLinFullQdBarOosRng}{0.1739--0.1788}  % ... the range as text
\providecommand{\rhcaLinHarQdBarOos}{}\renewcommand{\rhcaLinHarQdBarOos}{0.1915}  % QLIKE dispersion term, linear HAR, BarOos
\providecommand{\rhcaTreeFullQlBarOosMin}{}\renewcommand{\rhcaTreeFullQlBarOosMin}{0.0022}  % QLIKE level term, tree full, BarOos: min
\providecommand{\rhcaTreeFullQlBarOosMax}{}\renewcommand{\rhcaTreeFullQlBarOosMax}{0.0035}  % ... max
\providecommand{\rhcaTreeFullQlBarOosRng}{}\renewcommand{\rhcaTreeFullQlBarOosRng}{0.0022--0.0035}  % ... the range as text
\providecommand{\rhcaTreeHarQlBarOosMin}{}\renewcommand{\rhcaTreeHarQlBarOosMin}{0.0022}  % QLIKE level term, tree HAR, BarOos: min
\providecommand{\rhcaTreeHarQlBarOosMax}{}\renewcommand{\rhcaTreeHarQlBarOosMax}{0.0028}  % ... max
\providecommand{\rhcaTreeHarQlBarOosRng}{}\renewcommand{\rhcaTreeHarQlBarOosRng}{0.0022--0.0028}  % ... the range as text
\providecommand{\rhcaTreeFullQdBarOosMin}{}\renewcommand{\rhcaTreeFullQdBarOosMin}{0.1612}  % QLIKE dispersion term, tree full, BarOos: min
\providecommand{\rhcaTreeFullQdBarOosMax}{}\renewcommand{\rhcaTreeFullQdBarOosMax}{0.1694}  % ... max
\providecommand{\rhcaTreeFullQdBarOosRng}{}\renewcommand{\rhcaTreeFullQdBarOosRng}{0.1612--0.1694}  % ... the range as text
\providecommand{\rhcaTreeHarQdBarOosMin}{}\renewcommand{\rhcaTreeHarQdBarOosMin}{0.1878}  % QLIKE dispersion term, tree HAR, BarOos: min
\providecommand{\rhcaTreeHarQdBarOosMax}{}\renewcommand{\rhcaTreeHarQdBarOosMax}{0.1909}  % ... max
\providecommand{\rhcaTreeHarQdBarOosRng}{}\renewcommand{\rhcaTreeHarQdBarOosRng}{0.1878--0.1909}  % ... the range as text
\providecommand{\rhcaTreeFullUnderBarOosMin}{}\renewcommand{\rhcaTreeFullUnderBarOosMin}{40.5}  % under-prediction %, tree full, BarOos: min
\providecommand{\rhcaTreeFullUnderBarOosMax}{}\renewcommand{\rhcaTreeFullUnderBarOosMax}{41.0}  % ... max
\providecommand{\rhcaTreeFullUnderBarOosRng}{}\renewcommand{\rhcaTreeFullUnderBarOosRng}{40.5--41.0}  % ... the range as text
\providecommand{\rhcaTreeHarUnderBarOosMin}{}\renewcommand{\rhcaTreeHarUnderBarOosMin}{38.4}  % under-prediction %, tree HAR, BarOos: min
\providecommand{\rhcaTreeHarUnderBarOosMax}{}\renewcommand{\rhcaTreeHarUnderBarOosMax}{38.8}  % ... max
\providecommand{\rhcaTreeHarUnderBarOosRng}{}\renewcommand{\rhcaTreeHarUnderBarOosRng}{38.4--38.8}  % ... the range as text
\providecommand{\rhcaLinFullMzBClosePreMin}{}\renewcommand{\rhcaLinFullMzBClosePreMin}{1.06}  % levels MZ slope, linear full, ClosePre: min
\providecommand{\rhcaLinFullMzBClosePreMax}{}\renewcommand{\rhcaLinFullMzBClosePreMax}{1.08}  % ... max
\providecommand{\rhcaLinFullMzBClosePreRng}{}\renewcommand{\rhcaLinFullMzBClosePreRng}{1.06--1.08}  % ... the range as text
\providecommand{\rhcaTreeFullMzBClosePreMin}{}\renewcommand{\rhcaTreeFullMzBClosePreMin}{0.98}  % levels MZ slope, tree full, ClosePre: min
\providecommand{\rhcaTreeFullMzBClosePreMax}{}\renewcommand{\rhcaTreeFullMzBClosePreMax}{0.99}  % ... max
\providecommand{\rhcaTreeFullMzBClosePreRng}{}\renewcommand{\rhcaTreeFullMzBClosePreRng}{0.98--0.99}  % ... the range as text
\providecommand{\rhcaTreeHarMzBClosePreMin}{}\renewcommand{\rhcaTreeHarMzBClosePreMin}{1.06}  % levels MZ slope, tree HAR, ClosePre: min
\providecommand{\rhcaTreeHarMzBClosePreMax}{}\renewcommand{\rhcaTreeHarMzBClosePreMax}{1.11}  % ... max
\providecommand{\rhcaTreeHarMzBClosePreRng}{}\renewcommand{\rhcaTreeHarMzBClosePreRng}{1.06--1.11}  % ... the range as text
\providecommand{\rhcaLinHarMzBClosePre}{}\renewcommand{\rhcaLinHarMzBClosePre}{1.14}  % levels MZ slope, linear HAR (= OLS), ClosePre
\providecommand{\rhcaLinFullMzlBClosePreMin}{}\renewcommand{\rhcaLinFullMzlBClosePreMin}{0.92}  % logs MZ slope, linear full, ClosePre: min
\providecommand{\rhcaLinFullMzlBClosePreMax}{}\renewcommand{\rhcaLinFullMzlBClosePreMax}{0.93}  % ... max
\providecommand{\rhcaLinFullMzlBClosePreRng}{}\renewcommand{\rhcaLinFullMzlBClosePreRng}{0.92--0.93}  % ... the range as text
\providecommand{\rhcaTreeFullMzlBClosePreMin}{}\renewcommand{\rhcaTreeFullMzlBClosePreMin}{0.92}  % logs MZ slope, tree full, ClosePre: min
\providecommand{\rhcaTreeFullMzlBClosePreMax}{}\renewcommand{\rhcaTreeFullMzlBClosePreMax}{0.94}  % ... max
\providecommand{\rhcaTreeFullMzlBClosePreRng}{}\renewcommand{\rhcaTreeFullMzlBClosePreRng}{0.92--0.94}  % ... the range as text
\providecommand{\rhcaAllHarMzlBClosePreMin}{}\renewcommand{\rhcaAllHarMzlBClosePreMin}{0.95}  % logs MZ slope, HAR forecasts, ClosePre: min
\providecommand{\rhcaAllHarMzlBClosePreMax}{}\renewcommand{\rhcaAllHarMzlBClosePreMax}{0.96}  % ... max
\providecommand{\rhcaAllHarMzlBClosePreRng}{}\renewcommand{\rhcaAllHarMzlBClosePreRng}{0.95--0.96}  % ... the range as text
\providecommand{\rhcaLinFullMzlRsqClosePreMin}{}\renewcommand{\rhcaLinFullMzlRsqClosePreMin}{0.80}  % logs MZ R2, linear full, ClosePre: min
\providecommand{\rhcaLinFullMzlRsqClosePreMax}{}\renewcommand{\rhcaLinFullMzlRsqClosePreMax}{0.80}  % ... max
\providecommand{\rhcaLinFullMzlRsqClosePreRng}{}\renewcommand{\rhcaLinFullMzlRsqClosePreRng}{0.80}  % ... the range as text
\providecommand{\rhcaTreeFullMzlRsqClosePreMin}{}\renewcommand{\rhcaTreeFullMzlRsqClosePreMin}{0.79}  % logs MZ R2, tree full, ClosePre: min
\providecommand{\rhcaTreeFullMzlRsqClosePreMax}{}\renewcommand{\rhcaTreeFullMzlRsqClosePreMax}{0.80}  % ... max
\providecommand{\rhcaTreeFullMzlRsqClosePreRng}{}\renewcommand{\rhcaTreeFullMzlRsqClosePreRng}{0.79--0.80}  % ... the range as text
\providecommand{\rhcaAllHarMzlRsqClosePreMin}{}\renewcommand{\rhcaAllHarMzlRsqClosePreMin}{0.78}  % logs MZ R2, HAR forecasts, ClosePre: min
\providecommand{\rhcaAllHarMzlRsqClosePreMax}{}\renewcommand{\rhcaAllHarMzlRsqClosePreMax}{0.78}  % ... max
\providecommand{\rhcaAllHarMzlRsqClosePreRng}{}\renewcommand{\rhcaAllHarMzlRsqClosePreRng}{0.78}  % ... the range as text
\providecommand{\rhcaLinFullMzRsqClosePreMin}{}\renewcommand{\rhcaLinFullMzRsqClosePreMin}{0.68}  % levels MZ R2, linear full, ClosePre: min
\providecommand{\rhcaLinFullMzRsqClosePreMax}{}\renewcommand{\rhcaLinFullMzRsqClosePreMax}{0.69}  % ... max
\providecommand{\rhcaLinFullMzRsqClosePreRng}{}\renewcommand{\rhcaLinFullMzRsqClosePreRng}{0.68--0.69}  % ... the range as text
\providecommand{\rhcaAllHarMzRsqClosePreMin}{}\renewcommand{\rhcaAllHarMzRsqClosePreMin}{0.55}  % levels MZ R2, HAR forecasts, ClosePre: min
\providecommand{\rhcaAllHarMzRsqClosePreMax}{}\renewcommand{\rhcaAllHarMzRsqClosePreMax}{0.64}  % ... max
\providecommand{\rhcaAllHarMzRsqClosePreRng}{}\renewcommand{\rhcaAllHarMzRsqClosePreRng}{0.55--0.64}  % ... the range as text
\providecommand{\rhcaLinFullShSlopeClosePreMin}{}\renewcommand{\rhcaLinFullShSlopeClosePreMin}{0.7}  % Theil slope share %, linear full, ClosePre: min
\providecommand{\rhcaLinFullShSlopeClosePreMax}{}\renewcommand{\rhcaLinFullShSlopeClosePreMax}{1.2}  % ... max
\providecommand{\rhcaLinFullShSlopeClosePreRng}{}\renewcommand{\rhcaLinFullShSlopeClosePreRng}{0.7--1.2}  % ... the range as text
\providecommand{\rhcaTreeFullShSlopeClosePreMin}{}\renewcommand{\rhcaTreeFullShSlopeClosePreMin}{0.0}  % Theil slope share %, tree full, ClosePre: min
\providecommand{\rhcaTreeFullShSlopeClosePreMax}{}\renewcommand{\rhcaTreeFullShSlopeClosePreMax}{0.1}  % ... max
\providecommand{\rhcaTreeFullShSlopeClosePreRng}{}\renewcommand{\rhcaTreeFullShSlopeClosePreRng}{0.0--0.1}  % ... the range as text
\providecommand{\rhcaAllHarShSlopeClosePreMin}{}\renewcommand{\rhcaAllHarShSlopeClosePreMin}{0.4}  % Theil slope share %, HAR forecasts, ClosePre: min
\providecommand{\rhcaAllHarShSlopeClosePreMax}{}\renewcommand{\rhcaAllHarShSlopeClosePreMax}{2.5}  % ... max
\providecommand{\rhcaAllHarShSlopeClosePreRng}{}\renewcommand{\rhcaAllHarShSlopeClosePreRng}{0.4--2.5}  % ... the range as text
\providecommand{\rhcaAllShBiasClosePreMin}{}\renewcommand{\rhcaAllShBiasClosePreMin}{0.1}  % Theil bias share %, all 12, ClosePre: min
\providecommand{\rhcaAllShBiasClosePreMax}{}\renewcommand{\rhcaAllShBiasClosePreMax}{0.5}  % ... max
\providecommand{\rhcaAllShBiasClosePreRng}{}\renewcommand{\rhcaAllShBiasClosePreRng}{0.1--0.5}  % ... the range as text
\providecommand{\rhcaAllRmClosePreMin}{}\renewcommand{\rhcaAllRmClosePreMin}{1.04}  % ratio of means, all 12, ClosePre: min
\providecommand{\rhcaAllRmClosePreMax}{}\renewcommand{\rhcaAllRmClosePreMax}{1.10}  % ... max
\providecommand{\rhcaAllRmClosePreRng}{}\renewcommand{\rhcaAllRmClosePreRng}{1.04--1.10}  % ... the range as text
\providecommand{\rhcaFullRmClosePreMin}{}\renewcommand{\rhcaFullRmClosePreMin}{1.04}  % ratio of means, full set, ClosePre: min
\providecommand{\rhcaFullRmClosePreMax}{}\renewcommand{\rhcaFullRmClosePreMax}{1.07}  % ... max
\providecommand{\rhcaFullRmClosePreRng}{}\renewcommand{\rhcaFullRmClosePreRng}{1.04--1.07}  % ... the range as text
\providecommand{\rhcaHarRmClosePreMin}{}\renewcommand{\rhcaHarRmClosePreMin}{1.06}  % ratio of means, HAR, ClosePre: min
\providecommand{\rhcaHarRmClosePreMax}{}\renewcommand{\rhcaHarRmClosePreMax}{1.10}  % ... max
\providecommand{\rhcaHarRmClosePreRng}{}\renewcommand{\rhcaHarRmClosePreRng}{1.06--1.10}  % ... the range as text
\providecommand{\rhcaAllMrClosePreMin}{}\renewcommand{\rhcaAllMrClosePreMin}{1.06}  % mean of RV/sigma2, all 12, ClosePre: min
\providecommand{\rhcaAllMrClosePreMax}{}\renewcommand{\rhcaAllMrClosePreMax}{1.15}  % ... max
\providecommand{\rhcaAllMrClosePreRng}{}\renewcommand{\rhcaAllMrClosePreRng}{1.06--1.15}  % ... the range as text
\providecommand{\rhcaAllUnderClosePreMin}{}\renewcommand{\rhcaAllUnderClosePreMin}{42.6}  % under-prediction %, all 12, ClosePre: min
\providecommand{\rhcaAllUnderClosePreMax}{}\renewcommand{\rhcaAllUnderClosePreMax}{48.6}  % ... max
\providecommand{\rhcaAllUnderClosePreRng}{}\renewcommand{\rhcaAllUnderClosePreRng}{42.6--48.6}  % ... the range as text
\providecommand{\rhcaLinFullUnderClosePreMin}{}\renewcommand{\rhcaLinFullUnderClosePreMin}{47.2}  % under-prediction %, linear full, ClosePre: min
\providecommand{\rhcaLinFullUnderClosePreMax}{}\renewcommand{\rhcaLinFullUnderClosePreMax}{48.6}  % ... max
\providecommand{\rhcaLinFullUnderClosePreRng}{}\renewcommand{\rhcaLinFullUnderClosePreRng}{47.2--48.6}  % ... the range as text
\providecommand{\rhcaLinHarUnderClosePre}{}\renewcommand{\rhcaLinHarUnderClosePre}{46.7}  % under-prediction %, linear HAR, ClosePre
\providecommand{\rhcaLinFullOverSeClosePreMin}{}\renewcommand{\rhcaLinFullOverSeClosePreMin}{17.3}  % % of SE from over-predictions, linear full, ClosePre: min
\providecommand{\rhcaLinFullOverSeClosePreMax}{}\renewcommand{\rhcaLinFullOverSeClosePreMax}{18.4}  % ... max
\providecommand{\rhcaLinFullOverSeClosePreRng}{}\renewcommand{\rhcaLinFullOverSeClosePreRng}{17.3--18.4}  % ... the range as text
\providecommand{\rhcaTreeFullOverSeClosePreMin}{}\renewcommand{\rhcaTreeFullOverSeClosePreMin}{17.3}  % % of SE from over-predictions, tree full, ClosePre: min
\providecommand{\rhcaTreeFullOverSeClosePreMax}{}\renewcommand{\rhcaTreeFullOverSeClosePreMax}{18.9}  % ... max
\providecommand{\rhcaTreeFullOverSeClosePreRng}{}\renewcommand{\rhcaTreeFullOverSeClosePreRng}{17.3--18.9}  % ... the range as text
\providecommand{\rhcaAllHarOverSeClosePreMin}{}\renewcommand{\rhcaAllHarOverSeClosePreMin}{13.1}  % % of SE from over-predictions, HAR, ClosePre: min
\providecommand{\rhcaAllHarOverSeClosePreMax}{}\renewcommand{\rhcaAllHarOverSeClosePreMax}{14.4}  % ... max
\providecommand{\rhcaAllHarOverSeClosePreRng}{}\renewcommand{\rhcaAllHarOverSeClosePreRng}{13.1--14.4}  % ... the range as text
\providecommand{\rhcaNMaeLowerClosePre}{}\renewcommand{\rhcaNMaeLowerClosePre}{6}  % classes (of 6) with the lower MAE on the full set, ClosePre
\providecommand{\rhcaLinFullMrClosePreMin}{}\renewcommand{\rhcaLinFullMrClosePreMin}{1.137}  % mean of RV/sigma2, linear full, ClosePre: min
\providecommand{\rhcaLinFullMrClosePreMax}{}\renewcommand{\rhcaLinFullMrClosePreMax}{1.148}  % ... max
\providecommand{\rhcaLinFullMrClosePreRng}{}\renewcommand{\rhcaLinFullMrClosePreRng}{1.137--1.148}  % ... the range as text
\providecommand{\rhcaLinHarMrClosePre}{}\renewcommand{\rhcaLinHarMrClosePre}{1.121}  % mean of RV/sigma2, linear HAR, ClosePre
\providecommand{\rhcaLinFullQlClosePreMin}{}\renewcommand{\rhcaLinFullQlClosePreMin}{0.0086}  % QLIKE level term, linear full, ClosePre: min
\providecommand{\rhcaLinFullQlClosePreMax}{}\renewcommand{\rhcaLinFullQlClosePreMax}{0.0100}  % ... max
\providecommand{\rhcaLinFullQlClosePreRng}{}\renewcommand{\rhcaLinFullQlClosePreRng}{0.0086--0.0100}  % ... the range as text
\providecommand{\rhcaLinHarQlClosePre}{}\renewcommand{\rhcaLinHarQlClosePre}{0.0068}  % QLIKE level term, linear HAR, ClosePre
\providecommand{\rhcaLinFullQdClosePreMin}{}\renewcommand{\rhcaLinFullQdClosePreMin}{0.1213}  % QLIKE dispersion term, linear full, ClosePre: min
\providecommand{\rhcaLinFullQdClosePreMax}{}\renewcommand{\rhcaLinFullQdClosePreMax}{0.1233}  % ... max
\providecommand{\rhcaLinFullQdClosePreRng}{}\renewcommand{\rhcaLinFullQdClosePreRng}{0.1213--0.1233}  % ... the range as text
\providecommand{\rhcaLinHarQdClosePre}{}\renewcommand{\rhcaLinHarQdClosePre}{0.1297}  % QLIKE dispersion term, linear HAR, ClosePre
\providecommand{\rhcaTreeFullQlClosePreMin}{}\renewcommand{\rhcaTreeFullQlClosePreMin}{0.0045}  % QLIKE level term, tree full, ClosePre: min
\providecommand{\rhcaTreeFullQlClosePreMax}{}\renewcommand{\rhcaTreeFullQlClosePreMax}{0.0065}  % ... max
\providecommand{\rhcaTreeFullQlClosePreRng}{}\renewcommand{\rhcaTreeFullQlClosePreRng}{0.0045--0.0065}  % ... the range as text
\providecommand{\rhcaTreeHarQlClosePreMin}{}\renewcommand{\rhcaTreeHarQlClosePreMin}{0.0020}  % QLIKE level term, tree HAR, ClosePre: min
\providecommand{\rhcaTreeHarQlClosePreMax}{}\renewcommand{\rhcaTreeHarQlClosePreMax}{0.0039}  % ... max
\providecommand{\rhcaTreeHarQlClosePreRng}{}\renewcommand{\rhcaTreeHarQlClosePreRng}{0.0020--0.0039}  % ... the range as text
\providecommand{\rhcaTreeFullQdClosePreMin}{}\renewcommand{\rhcaTreeFullQdClosePreMin}{0.1232}  % QLIKE dispersion term, tree full, ClosePre: min
\providecommand{\rhcaTreeFullQdClosePreMax}{}\renewcommand{\rhcaTreeFullQdClosePreMax}{0.1270}  % ... max
\providecommand{\rhcaTreeFullQdClosePreRng}{}\renewcommand{\rhcaTreeFullQdClosePreRng}{0.1232--0.1270}  % ... the range as text
\providecommand{\rhcaTreeHarQdClosePreMin}{}\renewcommand{\rhcaTreeHarQdClosePreMin}{0.1310}  % QLIKE dispersion term, tree HAR, ClosePre: min
\providecommand{\rhcaTreeHarQdClosePreMax}{}\renewcommand{\rhcaTreeHarQdClosePreMax}{0.1325}  % ... max
\providecommand{\rhcaTreeHarQdClosePreRng}{}\renewcommand{\rhcaTreeHarQdClosePreRng}{0.1310--0.1325}  % ... the range as text
\providecommand{\rhcaTreeFullUnderClosePreMin}{}\renewcommand{\rhcaTreeFullUnderClosePreMin}{45.6}  % under-prediction %, tree full, ClosePre: min
\providecommand{\rhcaTreeFullUnderClosePreMax}{}\renewcommand{\rhcaTreeFullUnderClosePreMax}{46.1}  % ... max
\providecommand{\rhcaTreeFullUnderClosePreRng}{}\renewcommand{\rhcaTreeFullUnderClosePreRng}{45.6--46.1}  % ... the range as text
\providecommand{\rhcaTreeHarUnderClosePreMin}{}\renewcommand{\rhcaTreeHarUnderClosePreMin}{42.6}  % under-prediction %, tree HAR, ClosePre: min
\providecommand{\rhcaTreeHarUnderClosePreMax}{}\renewcommand{\rhcaTreeHarUnderClosePreMax}{43.6}  % ... max
\providecommand{\rhcaTreeHarUnderClosePreRng}{}\renewcommand{\rhcaTreeHarUnderClosePreRng}{42.6--43.6}  % ... the range as text
\providecommand{\rhcaLinFullMzBCloseOosMin}{}\renewcommand{\rhcaLinFullMzBCloseOosMin}{0.66}  % levels MZ slope, linear full, CloseOos: min
\providecommand{\rhcaLinFullMzBCloseOosMax}{}\renewcommand{\rhcaLinFullMzBCloseOosMax}{0.70}  % ... max
\providecommand{\rhcaLinFullMzBCloseOosRng}{}\renewcommand{\rhcaLinFullMzBCloseOosRng}{0.66--0.70}  % ... the range as text
\providecommand{\rhcaTreeFullMzBCloseOosMin}{}\renewcommand{\rhcaTreeFullMzBCloseOosMin}{0.81}  % levels MZ slope, tree full, CloseOos: min
\providecommand{\rhcaTreeFullMzBCloseOosMax}{}\renewcommand{\rhcaTreeFullMzBCloseOosMax}{0.89}  % ... max
\providecommand{\rhcaTreeFullMzBCloseOosRng}{}\renewcommand{\rhcaTreeFullMzBCloseOosRng}{0.81--0.89}  % ... the range as text
\providecommand{\rhcaTreeHarMzBCloseOosMin}{}\renewcommand{\rhcaTreeHarMzBCloseOosMin}{1.06}  % levels MZ slope, tree HAR, CloseOos: min
\providecommand{\rhcaTreeHarMzBCloseOosMax}{}\renewcommand{\rhcaTreeHarMzBCloseOosMax}{1.08}  % ... max
\providecommand{\rhcaTreeHarMzBCloseOosRng}{}\renewcommand{\rhcaTreeHarMzBCloseOosRng}{1.06--1.08}  % ... the range as text
\providecommand{\rhcaLinHarMzBCloseOos}{}\renewcommand{\rhcaLinHarMzBCloseOos}{1.21}  % levels MZ slope, linear HAR (= OLS), CloseOos
\providecommand{\rhcaLinFullMzlBCloseOosMin}{}\renewcommand{\rhcaLinFullMzlBCloseOosMin}{0.92}  % logs MZ slope, linear full, CloseOos: min
\providecommand{\rhcaLinFullMzlBCloseOosMax}{}\renewcommand{\rhcaLinFullMzlBCloseOosMax}{0.92}  % ... max
\providecommand{\rhcaLinFullMzlBCloseOosRng}{}\renewcommand{\rhcaLinFullMzlBCloseOosRng}{0.92}  % ... the range as text
\providecommand{\rhcaTreeFullMzlBCloseOosMin}{}\renewcommand{\rhcaTreeFullMzlBCloseOosMin}{0.93}  % logs MZ slope, tree full, CloseOos: min
\providecommand{\rhcaTreeFullMzlBCloseOosMax}{}\renewcommand{\rhcaTreeFullMzlBCloseOosMax}{0.95}  % ... max
\providecommand{\rhcaTreeFullMzlBCloseOosRng}{}\renewcommand{\rhcaTreeFullMzlBCloseOosRng}{0.93--0.95}  % ... the range as text
\providecommand{\rhcaAllHarMzlBCloseOosMin}{}\renewcommand{\rhcaAllHarMzlBCloseOosMin}{0.97}  % logs MZ slope, HAR forecasts, CloseOos: min
\providecommand{\rhcaAllHarMzlBCloseOosMax}{}\renewcommand{\rhcaAllHarMzlBCloseOosMax}{0.98}  % ... max
\providecommand{\rhcaAllHarMzlBCloseOosRng}{}\renewcommand{\rhcaAllHarMzlBCloseOosRng}{0.97--0.98}  % ... the range as text
\providecommand{\rhcaLinFullMzlRsqCloseOosMin}{}\renewcommand{\rhcaLinFullMzlRsqCloseOosMin}{0.84}  % logs MZ R2, linear full, CloseOos: min
\providecommand{\rhcaLinFullMzlRsqCloseOosMax}{}\renewcommand{\rhcaLinFullMzlRsqCloseOosMax}{0.84}  % ... max
\providecommand{\rhcaLinFullMzlRsqCloseOosRng}{}\renewcommand{\rhcaLinFullMzlRsqCloseOosRng}{0.84}  % ... the range as text
\providecommand{\rhcaTreeFullMzlRsqCloseOosMin}{}\renewcommand{\rhcaTreeFullMzlRsqCloseOosMin}{0.85}  % logs MZ R2, tree full, CloseOos: min
\providecommand{\rhcaTreeFullMzlRsqCloseOosMax}{}\renewcommand{\rhcaTreeFullMzlRsqCloseOosMax}{0.85}  % ... max
\providecommand{\rhcaTreeFullMzlRsqCloseOosRng}{}\renewcommand{\rhcaTreeFullMzlRsqCloseOosRng}{0.85}  % ... the range as text
\providecommand{\rhcaAllHarMzlRsqCloseOosMin}{}\renewcommand{\rhcaAllHarMzlRsqCloseOosMin}{0.83}  % logs MZ R2, HAR forecasts, CloseOos: min
\providecommand{\rhcaAllHarMzlRsqCloseOosMax}{}\renewcommand{\rhcaAllHarMzlRsqCloseOosMax}{0.84}  % ... max
\providecommand{\rhcaAllHarMzlRsqCloseOosRng}{}\renewcommand{\rhcaAllHarMzlRsqCloseOosRng}{0.83--0.84}  % ... the range as text
\providecommand{\rhcaLinFullMzRsqCloseOosMin}{}\renewcommand{\rhcaLinFullMzRsqCloseOosMin}{0.71}  % levels MZ R2, linear full, CloseOos: min
\providecommand{\rhcaLinFullMzRsqCloseOosMax}{}\renewcommand{\rhcaLinFullMzRsqCloseOosMax}{0.72}  % ... max
\providecommand{\rhcaLinFullMzRsqCloseOosRng}{}\renewcommand{\rhcaLinFullMzRsqCloseOosRng}{0.71--0.72}  % ... the range as text
\providecommand{\rhcaAllHarMzRsqCloseOosMin}{}\renewcommand{\rhcaAllHarMzRsqCloseOosMin}{0.80}  % levels MZ R2, HAR forecasts, CloseOos: min
\providecommand{\rhcaAllHarMzRsqCloseOosMax}{}\renewcommand{\rhcaAllHarMzRsqCloseOosMax}{0.82}  % ... max
\providecommand{\rhcaAllHarMzRsqCloseOosRng}{}\renewcommand{\rhcaAllHarMzRsqCloseOosRng}{0.80--0.82}  % ... the range as text
\providecommand{\rhcaLinFullShSlopeCloseOosMin}{}\renewcommand{\rhcaLinFullShSlopeCloseOosMin}{32.0}  % Theil slope share %, linear full, CloseOos: min
\providecommand{\rhcaLinFullShSlopeCloseOosMax}{}\renewcommand{\rhcaLinFullShSlopeCloseOosMax}{39.7}  % ... max
\providecommand{\rhcaLinFullShSlopeCloseOosRng}{}\renewcommand{\rhcaLinFullShSlopeCloseOosRng}{32.0--39.7}  % ... the range as text
\providecommand{\rhcaTreeFullShSlopeCloseOosMin}{}\renewcommand{\rhcaTreeFullShSlopeCloseOosMin}{5.5}  % Theil slope share %, tree full, CloseOos: min
\providecommand{\rhcaTreeFullShSlopeCloseOosMax}{}\renewcommand{\rhcaTreeFullShSlopeCloseOosMax}{14.8}  % ... max
\providecommand{\rhcaTreeFullShSlopeCloseOosRng}{}\renewcommand{\rhcaTreeFullShSlopeCloseOosRng}{5.5--14.8}  % ... the range as text
\providecommand{\rhcaAllHarShSlopeCloseOosMin}{}\renewcommand{\rhcaAllHarShSlopeCloseOosMin}{1.7}  % Theil slope share %, HAR forecasts, CloseOos: min
\providecommand{\rhcaAllHarShSlopeCloseOosMax}{}\renewcommand{\rhcaAllHarShSlopeCloseOosMax}{10.7}  % ... max
\providecommand{\rhcaAllHarShSlopeCloseOosRng}{}\renewcommand{\rhcaAllHarShSlopeCloseOosRng}{1.7--10.7}  % ... the range as text
\providecommand{\rhcaAllShBiasCloseOosMin}{}\renewcommand{\rhcaAllShBiasCloseOosMin}{0.0}  % Theil bias share %, all 12, CloseOos: min
\providecommand{\rhcaAllShBiasCloseOosMax}{}\renewcommand{\rhcaAllShBiasCloseOosMax}{0.3}  % ... max
\providecommand{\rhcaAllShBiasCloseOosRng}{}\renewcommand{\rhcaAllShBiasCloseOosRng}{0.0--0.3}  % ... the range as text
\providecommand{\rhcaAllRmCloseOosMin}{}\renewcommand{\rhcaAllRmCloseOosMin}{0.96}  % ratio of means, all 12, CloseOos: min
\providecommand{\rhcaAllRmCloseOosMax}{}\renewcommand{\rhcaAllRmCloseOosMax}{1.10}  % ... max
\providecommand{\rhcaAllRmCloseOosRng}{}\renewcommand{\rhcaAllRmCloseOosRng}{0.96--1.10}  % ... the range as text
\providecommand{\rhcaFullRmCloseOosMin}{}\renewcommand{\rhcaFullRmCloseOosMin}{0.96}  % ratio of means, full set, CloseOos: min
\providecommand{\rhcaFullRmCloseOosMax}{}\renewcommand{\rhcaFullRmCloseOosMax}{0.98}  % ... max
\providecommand{\rhcaFullRmCloseOosRng}{}\renewcommand{\rhcaFullRmCloseOosRng}{0.96--0.98}  % ... the range as text
\providecommand{\rhcaHarRmCloseOosMin}{}\renewcommand{\rhcaHarRmCloseOosMin}{1.03}  % ratio of means, HAR, CloseOos: min
\providecommand{\rhcaHarRmCloseOosMax}{}\renewcommand{\rhcaHarRmCloseOosMax}{1.10}  % ... max
\providecommand{\rhcaHarRmCloseOosRng}{}\renewcommand{\rhcaHarRmCloseOosRng}{1.03--1.10}  % ... the range as text
\providecommand{\rhcaAllMrCloseOosMin}{}\renewcommand{\rhcaAllMrCloseOosMin}{1.03}  % mean of RV/sigma2, all 12, CloseOos: min
\providecommand{\rhcaAllMrCloseOosMax}{}\renewcommand{\rhcaAllMrCloseOosMax}{1.18}  % ... max
\providecommand{\rhcaAllMrCloseOosRng}{}\renewcommand{\rhcaAllMrCloseOosRng}{1.03--1.18}  % ... the range as text
\providecommand{\rhcaAllUnderCloseOosMin}{}\renewcommand{\rhcaAllUnderCloseOosMin}{42.9}  % under-prediction %, all 12, CloseOos: min
\providecommand{\rhcaAllUnderCloseOosMax}{}\renewcommand{\rhcaAllUnderCloseOosMax}{52.9}  % ... max
\providecommand{\rhcaAllUnderCloseOosRng}{}\renewcommand{\rhcaAllUnderCloseOosRng}{42.9--52.9}  % ... the range as text
\providecommand{\rhcaLinFullUnderCloseOosMin}{}\renewcommand{\rhcaLinFullUnderCloseOosMin}{52.1}  % under-prediction %, linear full, CloseOos: min
\providecommand{\rhcaLinFullUnderCloseOosMax}{}\renewcommand{\rhcaLinFullUnderCloseOosMax}{52.9}  % ... max
\providecommand{\rhcaLinFullUnderCloseOosRng}{}\renewcommand{\rhcaLinFullUnderCloseOosRng}{52.1--52.9}  % ... the range as text
\providecommand{\rhcaLinHarUnderCloseOos}{}\renewcommand{\rhcaLinHarUnderCloseOos}{47.7}  % under-prediction %, linear HAR, CloseOos
\providecommand{\rhcaLinFullOverSeCloseOosMin}{}\renewcommand{\rhcaLinFullOverSeCloseOosMin}{51.5}  % % of SE from over-predictions, linear full, CloseOos: min
\providecommand{\rhcaLinFullOverSeCloseOosMax}{}\renewcommand{\rhcaLinFullOverSeCloseOosMax}{59.7}  % ... max
\providecommand{\rhcaLinFullOverSeCloseOosRng}{}\renewcommand{\rhcaLinFullOverSeCloseOosRng}{51.5--59.7}  % ... the range as text
\providecommand{\rhcaTreeFullOverSeCloseOosMin}{}\renewcommand{\rhcaTreeFullOverSeCloseOosMin}{43.1}  % % of SE from over-predictions, tree full, CloseOos: min
\providecommand{\rhcaTreeFullOverSeCloseOosMax}{}\renewcommand{\rhcaTreeFullOverSeCloseOosMax}{46.1}  % ... max
\providecommand{\rhcaTreeFullOverSeCloseOosRng}{}\renewcommand{\rhcaTreeFullOverSeCloseOosRng}{43.1--46.1}  % ... the range as text
\providecommand{\rhcaAllHarOverSeCloseOosMin}{}\renewcommand{\rhcaAllHarOverSeCloseOosMin}{13.8}  % % of SE from over-predictions, HAR, CloseOos: min
\providecommand{\rhcaAllHarOverSeCloseOosMax}{}\renewcommand{\rhcaAllHarOverSeCloseOosMax}{25.7}  % ... max
\providecommand{\rhcaAllHarOverSeCloseOosRng}{}\renewcommand{\rhcaAllHarOverSeCloseOosRng}{13.8--25.7}  % ... the range as text
\providecommand{\rhcaNMaeLowerCloseOos}{}\renewcommand{\rhcaNMaeLowerCloseOos}{0}  % classes (of 6) with the lower MAE on the full set, CloseOos
\providecommand{\rhcaLinQlShareCloseOosMin}{}\renewcommand{\rhcaLinQlShareCloseOosMin}{64}  % % of the linear full-set QLIKE increase from the level term, CloseOos: min
\providecommand{\rhcaLinQlShareCloseOosMax}{}\renewcommand{\rhcaLinQlShareCloseOosMax}{76}  % ... max
\providecommand{\rhcaLinQlShareCloseOosRng}{}\renewcommand{\rhcaLinQlShareCloseOosRng}{64--76}  % ... the range as text
\providecommand{\rhcaLinFullMrCloseOosMin}{}\renewcommand{\rhcaLinFullMrCloseOosMin}{1.165}  % mean of RV/sigma2, linear full, CloseOos: min
\providecommand{\rhcaLinFullMrCloseOosMax}{}\renewcommand{\rhcaLinFullMrCloseOosMax}{1.183}  % ... max
\providecommand{\rhcaLinFullMrCloseOosRng}{}\renewcommand{\rhcaLinFullMrCloseOosRng}{1.165--1.183}  % ... the range as text
\providecommand{\rhcaLinHarMrCloseOos}{}\renewcommand{\rhcaLinHarMrCloseOos}{1.092}  % mean of RV/sigma2, linear HAR, CloseOos
\providecommand{\rhcaLinFullQlCloseOosMin}{}\renewcommand{\rhcaLinFullQlCloseOosMin}{0.0123}  % QLIKE level term, linear full, CloseOos: min
\providecommand{\rhcaLinFullQlCloseOosMax}{}\renewcommand{\rhcaLinFullQlCloseOosMax}{0.0150}  % ... max
\providecommand{\rhcaLinFullQlCloseOosRng}{}\renewcommand{\rhcaLinFullQlCloseOosRng}{0.0123--0.0150}  % ... the range as text
\providecommand{\rhcaLinHarQlCloseOos}{}\renewcommand{\rhcaLinHarQlCloseOos}{0.0040}  % QLIKE level term, linear HAR, CloseOos
\providecommand{\rhcaLinFullQdCloseOosMin}{}\renewcommand{\rhcaLinFullQdCloseOosMin}{0.1122}  % QLIKE dispersion term, linear full, CloseOos: min
\providecommand{\rhcaLinFullQdCloseOosMax}{}\renewcommand{\rhcaLinFullQdCloseOosMax}{0.1159}  % ... max
\providecommand{\rhcaLinFullQdCloseOosRng}{}\renewcommand{\rhcaLinFullQdCloseOosRng}{0.1122--0.1159}  % ... the range as text
\providecommand{\rhcaLinHarQdCloseOos}{}\renewcommand{\rhcaLinHarQdCloseOos}{0.1097}  % QLIKE dispersion term, linear HAR, CloseOos
\providecommand{\rhcaTreeFullQlCloseOosMin}{}\renewcommand{\rhcaTreeFullQlCloseOosMin}{0.0039}  % QLIKE level term, tree full, CloseOos: min
\providecommand{\rhcaTreeFullQlCloseOosMax}{}\renewcommand{\rhcaTreeFullQlCloseOosMax}{0.0059}  % ... max
\providecommand{\rhcaTreeFullQlCloseOosRng}{}\renewcommand{\rhcaTreeFullQlCloseOosRng}{0.0039--0.0059}  % ... the range as text
\providecommand{\rhcaTreeHarQlCloseOosMin}{}\renewcommand{\rhcaTreeHarQlCloseOosMin}{0.0005}  % QLIKE level term, tree HAR, CloseOos: min
\providecommand{\rhcaTreeHarQlCloseOosMax}{}\renewcommand{\rhcaTreeHarQlCloseOosMax}{0.0019}  % ... max
\providecommand{\rhcaTreeHarQlCloseOosRng}{}\renewcommand{\rhcaTreeHarQlCloseOosRng}{0.0005--0.0019}  % ... the range as text
\providecommand{\rhcaTreeFullQdCloseOosMin}{}\renewcommand{\rhcaTreeFullQdCloseOosMin}{0.1044}  % QLIKE dispersion term, tree full, CloseOos: min
\providecommand{\rhcaTreeFullQdCloseOosMax}{}\renewcommand{\rhcaTreeFullQdCloseOosMax}{0.1064}  % ... max
\providecommand{\rhcaTreeFullQdCloseOosRng}{}\renewcommand{\rhcaTreeFullQdCloseOosRng}{0.1044--0.1064}  % ... the range as text
\providecommand{\rhcaTreeHarQdCloseOosMin}{}\renewcommand{\rhcaTreeHarQdCloseOosMin}{0.1095}  % QLIKE dispersion term, tree HAR, CloseOos: min
\providecommand{\rhcaTreeHarQdCloseOosMax}{}\renewcommand{\rhcaTreeHarQdCloseOosMax}{0.1112}  % ... max
\providecommand{\rhcaTreeHarQdCloseOosRng}{}\renewcommand{\rhcaTreeHarQdCloseOosRng}{0.1095--0.1112}  % ... the range as text
\providecommand{\rhcaTreeFullUnderCloseOosMin}{}\renewcommand{\rhcaTreeFullUnderCloseOosMin}{47.1}  % under-prediction %, tree full, CloseOos: min
\providecommand{\rhcaTreeFullUnderCloseOosMax}{}\renewcommand{\rhcaTreeFullUnderCloseOosMax}{50.1}  % ... max
\providecommand{\rhcaTreeFullUnderCloseOosRng}{}\renewcommand{\rhcaTreeFullUnderCloseOosRng}{47.1--50.1}  % ... the range as text
\providecommand{\rhcaTreeHarUnderCloseOosMin}{}\renewcommand{\rhcaTreeHarUnderCloseOosMin}{42.9}  % under-prediction %, tree HAR, CloseOos: min
\providecommand{\rhcaTreeHarUnderCloseOosMax}{}\renewcommand{\rhcaTreeHarUnderCloseOosMax}{45.5}  % ... max
\providecommand{\rhcaTreeHarUnderCloseOosRng}{}\renewcommand{\rhcaTreeHarUnderCloseOosRng}{42.9--45.5}  % ... the range as text
\providecommand{\rhcaLinTopShareBarOosMin}{}\renewcommand{\rhcaLinTopShareBarOosMin}{92}  % % of the full-minus-HAR squared-error sum (every bar, Oos) on the 5 days of largest difference, linear: min
\providecommand{\rhcaLinTopShareBarOosMax}{}\renewcommand{\rhcaLinTopShareBarOosMax}{94}  % ... max
\providecommand{\rhcaLinTopShareBarOosRng}{}\renewcommand{\rhcaLinTopShareBarOosRng}{92--94}  % ... the range as text
\providecommand{\rhcaTreeTopShareBarOosMin}{}\renewcommand{\rhcaTreeTopShareBarOosMin}{92}  % ... tree ensembles: min
\providecommand{\rhcaTreeTopShareBarOosMax}{}\renewcommand{\rhcaTreeTopShareBarOosMax}{101}  % ... max
\providecommand{\rhcaTreeTopShareBarOosRng}{}\renewcommand{\rhcaTreeTopShareBarOosRng}{92--101}  % ... the range as text
\providecommand{\rhcaTopFirstBarOos}{}\renewcommand{\rhcaTopFirstBarOos}{12 March 2020}  % earliest of those days, any class
\providecommand{\rhcaTopLastBarOos}{}\renewcommand{\rhcaTopLastBarOos}{2 April 2020}  % latest of those days, any class
\providecommand{\rhcaLinMsePctBarOosExCrashMin}{}\renewcommand{\rhcaLinMsePctBarOosExCrashMin}{\ensuremath{-}7.3}  % linear every-bar MSE %, full vs HAR, Oos without the crash window: min
\providecommand{\rhcaLinMsePctBarOosExCrashMax}{}\renewcommand{\rhcaLinMsePctBarOosExCrashMax}{\ensuremath{-}5.2}  % ... max
\providecommand{\rhcaLinMsePctBarOosExCrashRng}{}\renewcommand{\rhcaLinMsePctBarOosExCrashRng}{\ensuremath{-}7.3 to \ensuremath{-}5.2}  % ... the range as text
\providecommand{\rhcaTreeMsePctBarOosExCrashMin}{}\renewcommand{\rhcaTreeMsePctBarOosExCrashMin}{\ensuremath{-}9.0}  % ... tree ensembles: min
\providecommand{\rhcaTreeMsePctBarOosExCrashMax}{}\renewcommand{\rhcaTreeMsePctBarOosExCrashMax}{\ensuremath{-}6.0}  % ... max
\providecommand{\rhcaTreeMsePctBarOosExCrashRng}{}\renewcommand{\rhcaTreeMsePctBarOosExCrashRng}{\ensuremath{-}9.0 to \ensuremath{-}6.0}  % ... the range as text
\providecommand{\rhcaNMseLowerBarOosExCrash}{}\renewcommand{\rhcaNMseLowerBarOosExCrash}{6}  % classes (of 6) with the lower every-bar MSE on the full set, Oos without the crash window
\providecommand{\rhcaNMseLowerSigBarOosExCrash}{}\renewcommand{\rhcaNMseLowerSigBarOosExCrash}{6}  % ... significantly (DM on the daily mean squared error)
\providecommand{\rhcaLinMsePctCloseOosExCrashMin}{}\renewcommand{\rhcaLinMsePctCloseOosExCrashMin}{6.6}  % linear 16:00 MSE %, full vs HAR, Oos without the crash window: min
\providecommand{\rhcaLinMsePctCloseOosExCrashMax}{}\renewcommand{\rhcaLinMsePctCloseOosExCrashMax}{10.0}  % ... max
\providecommand{\rhcaLinMsePctCloseOosExCrashRng}{}\renewcommand{\rhcaLinMsePctCloseOosExCrashRng}{6.6--10.0}  % ... the range as text
\providecommand{\rhcaNMseLowerCloseOosExCrash}{}\renewcommand{\rhcaNMseLowerCloseOosExCrash}{1}  % classes with the lower 16:00 MSE on the full set, Oos without the crash window
\providecommand{\rhcaNMseHigherSigCloseOosExCrash}{}\renewcommand{\rhcaNMseHigherSigCloseOosExCrash}{0}  % ... classes with a significantly higher 16:00 MSE there
\providecommand{\rhcaTreeMsePctCloseOosExCrashMin}{}\renewcommand{\rhcaTreeMsePctCloseOosExCrashMin}{\ensuremath{-}0.7}  % tree 16:00 MSE %, full vs HAR, Oos without the crash window: min
\providecommand{\rhcaTreeMsePctCloseOosExCrashMax}{}\renewcommand{\rhcaTreeMsePctCloseOosExCrashMax}{3.4}  % ... max
\providecommand{\rhcaTreeMsePctCloseOosExCrashRng}{}\renewcommand{\rhcaTreeMsePctCloseOosExCrashRng}{\ensuremath{-}0.7 to 3.4}  % ... the range as text
\providecommand{\rhcaAllMaePctBarOosMin}{}\renewcommand{\rhcaAllMaePctBarOosMin}{2.3}  % every-bar MAE %, full vs HAR, Oos, all six classes: min
\providecommand{\rhcaAllMaePctBarOosMax}{}\renewcommand{\rhcaAllMaePctBarOosMax}{12.8}  % ... max
\providecommand{\rhcaAllMaePctBarOosRng}{}\renewcommand{\rhcaAllMaePctBarOosRng}{2.3--12.8}  % ... the range as text
\providecommand{\rhcaAllMaePctBarOosExCrashMin}{}\renewcommand{\rhcaAllMaePctBarOosExCrashMin}{\ensuremath{-}6.3}  % ... without the crash window: min
\providecommand{\rhcaAllMaePctBarOosExCrashMax}{}\renewcommand{\rhcaAllMaePctBarOosExCrashMax}{\ensuremath{-}4.4}  % ... max
\providecommand{\rhcaAllMaePctBarOosExCrashRng}{}\renewcommand{\rhcaAllMaePctBarOosExCrashRng}{\ensuremath{-}6.3 to \ensuremath{-}4.4}  % ... the range as text
\providecommand{\rhcaNMaeLowerBarOosExCrash}{}\renewcommand{\rhcaNMaeLowerBarOosExCrash}{6}  % classes with the lower every-bar MAE on the full set, Oos without the crash window
\providecommand{\rhcaNMaeLowerSigBarOosExCrash}{}\renewcommand{\rhcaNMaeLowerSigBarOosExCrash}{6}  % ... significantly
\providecommand{\rhcaTopShareRidgeBarOos}{}\renewcommand{\rhcaTopShareRidgeBarOos}{94}  % top-day share, %
\providecommand{\rhcaMsePctRidgeBarOosExCrash}{}\renewcommand{\rhcaMsePctRidgeBarOosExCrash}{\ensuremath{-}5.2}  % every-bar MSE %, ex crash
\providecommand{\rhcaMseDmRidgeBarOosExCrash}{}\renewcommand{\rhcaMseDmRidgeBarOosExCrash}{\ensuremath{-}3.10}
\providecommand{\rhcaMsePRidgeBarOosExCrash}{}\renewcommand{\rhcaMsePRidgeBarOosExCrash}{\ensuremath{p = 0.002}}
\providecommand{\rhcaTopShareLassoBarOos}{}\renewcommand{\rhcaTopShareLassoBarOos}{92}  % top-day share, %
\providecommand{\rhcaMsePctLassoBarOosExCrash}{}\renewcommand{\rhcaMsePctLassoBarOosExCrash}{\ensuremath{-}7.3}  % every-bar MSE %, ex crash
\providecommand{\rhcaMseDmLassoBarOosExCrash}{}\renewcommand{\rhcaMseDmLassoBarOosExCrash}{\ensuremath{-}3.29}
\providecommand{\rhcaMsePLassoBarOosExCrash}{}\renewcommand{\rhcaMsePLassoBarOosExCrash}{\ensuremath{p = 0.001}}
\providecommand{\rhcaTopShareEnetBarOos}{}\renewcommand{\rhcaTopShareEnetBarOos}{94}  % top-day share, %
\providecommand{\rhcaMsePctEnetBarOosExCrash}{}\renewcommand{\rhcaMsePctEnetBarOosExCrash}{\ensuremath{-}6.0}  % every-bar MSE %, ex crash
\providecommand{\rhcaMseDmEnetBarOosExCrash}{}\renewcommand{\rhcaMseDmEnetBarOosExCrash}{\ensuremath{-}3.43}
\providecommand{\rhcaMsePEnetBarOosExCrash}{}\renewcommand{\rhcaMsePEnetBarOosExCrash}{\ensuremath{p < 0.001}}
\providecommand{\rhcaTopShareLgbmBarOos}{}\renewcommand{\rhcaTopShareLgbmBarOos}{101}  % top-day share, %
\providecommand{\rhcaMsePctLgbmBarOosExCrash}{}\renewcommand{\rhcaMsePctLgbmBarOosExCrash}{\ensuremath{-}8.8}  % every-bar MSE %, ex crash
\providecommand{\rhcaMseDmLgbmBarOosExCrash}{}\renewcommand{\rhcaMseDmLgbmBarOosExCrash}{\ensuremath{-}4.30}
\providecommand{\rhcaMsePLgbmBarOosExCrash}{}\renewcommand{\rhcaMsePLgbmBarOosExCrash}{\ensuremath{p < 0.001}}
\providecommand{\rhcaTopShareXgbBarOos}{}\renewcommand{\rhcaTopShareXgbBarOos}{99}  % top-day share, %
\providecommand{\rhcaMsePctXgbBarOosExCrash}{}\renewcommand{\rhcaMsePctXgbBarOosExCrash}{\ensuremath{-}9.0}  % every-bar MSE %, ex crash
\providecommand{\rhcaMseDmXgbBarOosExCrash}{}\renewcommand{\rhcaMseDmXgbBarOosExCrash}{\ensuremath{-}3.67}
\providecommand{\rhcaMsePXgbBarOosExCrash}{}\renewcommand{\rhcaMsePXgbBarOosExCrash}{\ensuremath{p < 0.001}}
\providecommand{\rhcaTopShareRfBarOos}{}\renewcommand{\rhcaTopShareRfBarOos}{92}  % top-day share, %
\providecommand{\rhcaMsePctRfBarOosExCrash}{}\renewcommand{\rhcaMsePctRfBarOosExCrash}{\ensuremath{-}6.0}  % every-bar MSE %, ex crash
\providecommand{\rhcaMseDmRfBarOosExCrash}{}\renewcommand{\rhcaMseDmRfBarOosExCrash}{\ensuremath{-}3.97}
\providecommand{\rhcaMsePRfBarOosExCrash}{}\renewcommand{\rhcaMsePRfBarOosExCrash}{\ensuremath{p < 0.001}}
\providecommand{\rhcaMaxFcRidge}{}\renewcommand{\rhcaMaxFcRidge}{1947}  % largest full-set forecast on a bar, Oos, x 10^6
\providecommand{\rhcaMaxFcRvRidge}{}\renewcommand{\rhcaMaxFcRvRidge}{501}  % the realized variance of that bar, x 10^6
\providecommand{\rhcaMaxFcHarRidge}{}\renewcommand{\rhcaMaxFcHarRidge}{627}  % the HAR forecast of that bar, x 10^6
\providecommand{\rhcaMaxFcRatioRidge}{}\renewcommand{\rhcaMaxFcRatioRidge}{3.9}  % forecast / realized variance on that bar
\providecommand{\rhcaMaxFcLasso}{}\renewcommand{\rhcaMaxFcLasso}{2035}  % largest full-set forecast on a bar, Oos, x 10^6
\providecommand{\rhcaMaxFcRvLasso}{}\renewcommand{\rhcaMaxFcRvLasso}{501}  % the realized variance of that bar, x 10^6
\providecommand{\rhcaMaxFcHarLasso}{}\renewcommand{\rhcaMaxFcHarLasso}{627}  % the HAR forecast of that bar, x 10^6
\providecommand{\rhcaMaxFcRatioLasso}{}\renewcommand{\rhcaMaxFcRatioLasso}{4.1}  % forecast / realized variance on that bar
\providecommand{\rhcaMaxFcEnet}{}\renewcommand{\rhcaMaxFcEnet}{2064}  % largest full-set forecast on a bar, Oos, x 10^6
\providecommand{\rhcaMaxFcRvEnet}{}\renewcommand{\rhcaMaxFcRvEnet}{501}  % the realized variance of that bar, x 10^6
\providecommand{\rhcaMaxFcHarEnet}{}\renewcommand{\rhcaMaxFcHarEnet}{627}  % the HAR forecast of that bar, x 10^6
\providecommand{\rhcaMaxFcRatioEnet}{}\renewcommand{\rhcaMaxFcRatioEnet}{4.1}  % forecast / realized variance on that bar
\providecommand{\rhcaMaxFcDate}{}\renewcommand{\rhcaMaxFcDate}{23 March 2020}  % the bar of the largest linear full-set forecast (all three)
\providecommand{\rhcaMaxFcClock}{}\renewcommand{\rhcaMaxFcClock}{10:00}  % its bar end, ET
\providecommand{\rhcaMaxRvOos}{}\renewcommand{\rhcaMaxRvOos}{1150}  % largest bar RV on Oos, x 10^6
\providecommand{\rhcaYrsUpRidgeCloseOos}{}\renewcommand{\rhcaYrsUpRidgeCloseOos}{4}  % calendar years of Oos (of 5) with a higher 16:00 QLIKE on the full set
\providecommand{\rhcaYrsDownRidgeClosePre}{}\renewcommand{\rhcaYrsDownRidgeClosePre}{4}  % calendar years of Pre (of 6) with a lower 16:00 QLIKE
\providecommand{\rhcaWorstYrRidgeCloseOos}{}\renewcommand{\rhcaWorstYrRidgeCloseOos}{2023}  % the Oos year of the largest full-set increase
\providecommand{\rhcaWorstYrPctRidgeCloseOos}{}\renewcommand{\rhcaWorstYrPctRidgeCloseOos}{32.0}  % ... its QLIKE %
\providecommand{\rhcaYrsUpLassoCloseOos}{}\renewcommand{\rhcaYrsUpLassoCloseOos}{4}  % calendar years of Oos (of 5) with a higher 16:00 QLIKE on the full set
\providecommand{\rhcaYrsDownLassoClosePre}{}\renewcommand{\rhcaYrsDownLassoClosePre}{3}  % calendar years of Pre (of 6) with a lower 16:00 QLIKE
\providecommand{\rhcaWorstYrLassoCloseOos}{}\renewcommand{\rhcaWorstYrLassoCloseOos}{2023}  % the Oos year of the largest full-set increase
\providecommand{\rhcaWorstYrPctLassoCloseOos}{}\renewcommand{\rhcaWorstYrPctLassoCloseOos}{33.0}  % ... its QLIKE %
\providecommand{\rhcaYrsUpEnetCloseOos}{}\renewcommand{\rhcaYrsUpEnetCloseOos}{4}  % calendar years of Oos (of 5) with a higher 16:00 QLIKE on the full set
\providecommand{\rhcaYrsDownEnetClosePre}{}\renewcommand{\rhcaYrsDownEnetClosePre}{5}  % calendar years of Pre (of 6) with a lower 16:00 QLIKE
\providecommand{\rhcaWorstYrEnetCloseOos}{}\renewcommand{\rhcaWorstYrEnetCloseOos}{2023}  % the Oos year of the largest full-set increase
\providecommand{\rhcaWorstYrPctEnetCloseOos}{}\renewcommand{\rhcaWorstYrPctEnetCloseOos}{30.7}  % ... its QLIKE %
\providecommand{\rhcaYrsUpLgbmCloseOos}{}\renewcommand{\rhcaYrsUpLgbmCloseOos}{2}  % calendar years of Oos (of 5) with a higher 16:00 QLIKE on the full set
\providecommand{\rhcaYrsDownLgbmClosePre}{}\renewcommand{\rhcaYrsDownLgbmClosePre}{4}  % calendar years of Pre (of 6) with a lower 16:00 QLIKE
\providecommand{\rhcaWorstYrLgbmCloseOos}{}\renewcommand{\rhcaWorstYrLgbmCloseOos}{2024}  % the Oos year of the largest full-set increase
\providecommand{\rhcaWorstYrPctLgbmCloseOos}{}\renewcommand{\rhcaWorstYrPctLgbmCloseOos}{10.5}  % ... its QLIKE %
\providecommand{\rhcaYrsUpXgbCloseOos}{}\renewcommand{\rhcaYrsUpXgbCloseOos}{2}  % calendar years of Oos (of 5) with a higher 16:00 QLIKE on the full set
\providecommand{\rhcaYrsDownXgbClosePre}{}\renewcommand{\rhcaYrsDownXgbClosePre}{4}  % calendar years of Pre (of 6) with a lower 16:00 QLIKE
\providecommand{\rhcaWorstYrXgbCloseOos}{}\renewcommand{\rhcaWorstYrXgbCloseOos}{2024}  % the Oos year of the largest full-set increase
\providecommand{\rhcaWorstYrPctXgbCloseOos}{}\renewcommand{\rhcaWorstYrPctXgbCloseOos}{11.9}  % ... its QLIKE %
\providecommand{\rhcaYrsUpRfCloseOos}{}\renewcommand{\rhcaYrsUpRfCloseOos}{2}  % calendar years of Oos (of 5) with a higher 16:00 QLIKE on the full set
\providecommand{\rhcaYrsDownRfClosePre}{}\renewcommand{\rhcaYrsDownRfClosePre}{4}  % calendar years of Pre (of 6) with a lower 16:00 QLIKE
\providecommand{\rhcaWorstYrRfCloseOos}{}\renewcommand{\rhcaWorstYrRfCloseOos}{2021}  % the Oos year of the largest full-set increase
\providecommand{\rhcaWorstYrPctRfCloseOos}{}\renewcommand{\rhcaWorstYrPctRfCloseOos}{7.3}  % ... its QLIKE %
\providecommand{\rhcaNYrsOos}{}\renewcommand{\rhcaNYrsOos}{5}  % calendar years of Oos
\providecommand{\rhcaNYrsPre}{}\renewcommand{\rhcaNYrsPre}{6}  % calendar years of Pre
\providecommand{\rhcaLinMedLogRatioClosePreMin}{}\renewcommand{\rhcaLinMedLogRatioClosePreMin}{\ensuremath{-}3.1}  % median 100 ln(full / HAR forecast), linear, 16:00, Pre: min
\providecommand{\rhcaLinMedLogRatioClosePreMax}{}\renewcommand{\rhcaLinMedLogRatioClosePreMax}{\ensuremath{-}2.1}  % ... max
\providecommand{\rhcaLinMedLogRatioClosePreRng}{}\renewcommand{\rhcaLinMedLogRatioClosePreRng}{\ensuremath{-}3.1 to \ensuremath{-}2.1}  % ... the range as text
\providecommand{\rhcaTreeMedLogRatioClosePreMin}{}\renewcommand{\rhcaTreeMedLogRatioClosePreMin}{\ensuremath{-}5.1}  % median 100 ln(full / HAR forecast), trees, 16:00, Pre: min
\providecommand{\rhcaTreeMedLogRatioClosePreMax}{}\renewcommand{\rhcaTreeMedLogRatioClosePreMax}{\ensuremath{-}3.6}  % ... max
\providecommand{\rhcaTreeMedLogRatioClosePreRng}{}\renewcommand{\rhcaTreeMedLogRatioClosePreRng}{\ensuremath{-}5.1 to \ensuremath{-}3.6}  % ... the range as text
\providecommand{\rhcaLinMedLogRatioCloseOosMin}{}\renewcommand{\rhcaLinMedLogRatioCloseOosMin}{\ensuremath{-}6.9}  % median 100 ln(full / HAR forecast), linear, 16:00, Oos: min
\providecommand{\rhcaLinMedLogRatioCloseOosMax}{}\renewcommand{\rhcaLinMedLogRatioCloseOosMax}{\ensuremath{-}5.6}  % ... max
\providecommand{\rhcaLinMedLogRatioCloseOosRng}{}\renewcommand{\rhcaLinMedLogRatioCloseOosRng}{\ensuremath{-}6.9 to \ensuremath{-}5.6}  % ... the range as text
\providecommand{\rhcaTreeMedLogRatioCloseOosMin}{}\renewcommand{\rhcaTreeMedLogRatioCloseOosMin}{\ensuremath{-}8.6}  % median 100 ln(full / HAR forecast), trees, 16:00, Oos: min
\providecommand{\rhcaTreeMedLogRatioCloseOosMax}{}\renewcommand{\rhcaTreeMedLogRatioCloseOosMax}{\ensuremath{-}6.1}  % ... max
\providecommand{\rhcaTreeMedLogRatioCloseOosRng}{}\renewcommand{\rhcaTreeMedLogRatioCloseOosRng}{\ensuremath{-}8.6 to \ensuremath{-}6.1}  % ... the range as text
\providecommand{\rhcaTopDecRatioEnetHarBarPre}{}\renewcommand{\rhcaTopDecRatioEnetHarBarPre}{1.04}  % top forecast decile: mean RV / mean forecast
\providecommand{\rhcaTopDecSeEnetHarBarPre}{}\renewcommand{\rhcaTopDecSeEnetHarBarPre}{85}  % top forecast decile: % of the squared error
\providecommand{\rhcaBotDecRatioEnetHarBarPre}{}\renewcommand{\rhcaBotDecRatioEnetHarBarPre}{1.12}  % bottom forecast decile: mean RV / mean forecast
\providecommand{\rhcaTopDecRatioEnetFullBarPre}{}\renewcommand{\rhcaTopDecRatioEnetFullBarPre}{1.01}  % top forecast decile: mean RV / mean forecast
\providecommand{\rhcaTopDecSeEnetFullBarPre}{}\renewcommand{\rhcaTopDecSeEnetFullBarPre}{84}  % top forecast decile: % of the squared error
\providecommand{\rhcaBotDecRatioEnetFullBarPre}{}\renewcommand{\rhcaBotDecRatioEnetFullBarPre}{1.16}  % bottom forecast decile: mean RV / mean forecast
\providecommand{\rhcaTopDecRatioLgbmHarBarPre}{}\renewcommand{\rhcaTopDecRatioLgbmHarBarPre}{1.08}  % top forecast decile: mean RV / mean forecast
\providecommand{\rhcaTopDecSeLgbmHarBarPre}{}\renewcommand{\rhcaTopDecSeLgbmHarBarPre}{88}  % top forecast decile: % of the squared error
\providecommand{\rhcaBotDecRatioLgbmHarBarPre}{}\renewcommand{\rhcaBotDecRatioLgbmHarBarPre}{1.14}  % bottom forecast decile: mean RV / mean forecast
\providecommand{\rhcaTopDecRatioLgbmFullBarPre}{}\renewcommand{\rhcaTopDecRatioLgbmFullBarPre}{1.06}  % top forecast decile: mean RV / mean forecast
\providecommand{\rhcaTopDecSeLgbmFullBarPre}{}\renewcommand{\rhcaTopDecSeLgbmFullBarPre}{91}  % top forecast decile: % of the squared error
\providecommand{\rhcaBotDecRatioLgbmFullBarPre}{}\renewcommand{\rhcaBotDecRatioLgbmFullBarPre}{1.09}  % bottom forecast decile: mean RV / mean forecast
\providecommand{\rhcaTopDecRatioEnetHarBarOos}{}\renewcommand{\rhcaTopDecRatioEnetHarBarOos}{1.03}  % top forecast decile: mean RV / mean forecast
\providecommand{\rhcaTopDecSeEnetHarBarOos}{}\renewcommand{\rhcaTopDecSeEnetHarBarOos}{91}  % top forecast decile: % of the squared error
\providecommand{\rhcaBotDecRatioEnetHarBarOos}{}\renewcommand{\rhcaBotDecRatioEnetHarBarOos}{1.06}  % bottom forecast decile: mean RV / mean forecast
\providecommand{\rhcaTopDecRatioEnetFullBarOos}{}\renewcommand{\rhcaTopDecRatioEnetFullBarOos}{0.88}  % top forecast decile: mean RV / mean forecast
\providecommand{\rhcaTopDecSeEnetFullBarOos}{}\renewcommand{\rhcaTopDecSeEnetFullBarOos}{97}  % top forecast decile: % of the squared error
\providecommand{\rhcaBotDecRatioEnetFullBarOos}{}\renewcommand{\rhcaBotDecRatioEnetFullBarOos}{1.14}  % bottom forecast decile: mean RV / mean forecast
\providecommand{\rhcaTopDecRatioLgbmHarBarOos}{}\renewcommand{\rhcaTopDecRatioLgbmHarBarOos}{1.00}  % top forecast decile: mean RV / mean forecast
\providecommand{\rhcaTopDecSeLgbmHarBarOos}{}\renewcommand{\rhcaTopDecSeLgbmHarBarOos}{92}  % top forecast decile: % of the squared error
\providecommand{\rhcaBotDecRatioLgbmHarBarOos}{}\renewcommand{\rhcaBotDecRatioLgbmHarBarOos}{1.07}  % bottom forecast decile: mean RV / mean forecast
\providecommand{\rhcaTopDecRatioLgbmFullBarOos}{}\renewcommand{\rhcaTopDecRatioLgbmFullBarOos}{0.89}  % top forecast decile: mean RV / mean forecast
\providecommand{\rhcaTopDecSeLgbmFullBarOos}{}\renewcommand{\rhcaTopDecSeLgbmFullBarOos}{95}  % top forecast decile: % of the squared error
\providecommand{\rhcaBotDecRatioLgbmFullBarOos}{}\renewcommand{\rhcaBotDecRatioLgbmFullBarOos}{1.12}  % bottom forecast decile: mean RV / mean forecast
\providecommand{\rhcaTopDecRatioEnetHarClosePre}{}\renewcommand{\rhcaTopDecRatioEnetHarClosePre}{1.11}  % top forecast decile: mean RV / mean forecast
\providecommand{\rhcaTopDecSeEnetHarClosePre}{}\renewcommand{\rhcaTopDecSeEnetHarClosePre}{89}  % top forecast decile: % of the squared error
\providecommand{\rhcaBotDecRatioEnetHarClosePre}{}\renewcommand{\rhcaBotDecRatioEnetHarClosePre}{1.31}  % bottom forecast decile: mean RV / mean forecast
\providecommand{\rhcaTopDecRatioEnetFullClosePre}{}\renewcommand{\rhcaTopDecRatioEnetFullClosePre}{1.03}  % top forecast decile: mean RV / mean forecast
\providecommand{\rhcaTopDecSeEnetFullClosePre}{}\renewcommand{\rhcaTopDecSeEnetFullClosePre}{88}  % top forecast decile: % of the squared error
\providecommand{\rhcaBotDecRatioEnetFullClosePre}{}\renewcommand{\rhcaBotDecRatioEnetFullClosePre}{1.31}  % bottom forecast decile: mean RV / mean forecast
\providecommand{\rhcaTopDecRatioLgbmHarClosePre}{}\renewcommand{\rhcaTopDecRatioLgbmHarClosePre}{1.08}  % top forecast decile: mean RV / mean forecast
\providecommand{\rhcaTopDecSeLgbmHarClosePre}{}\renewcommand{\rhcaTopDecSeLgbmHarClosePre}{91}  % top forecast decile: % of the squared error
\providecommand{\rhcaBotDecRatioLgbmHarClosePre}{}\renewcommand{\rhcaBotDecRatioLgbmHarClosePre}{1.19}  % bottom forecast decile: mean RV / mean forecast
\providecommand{\rhcaTopDecRatioLgbmFullClosePre}{}\renewcommand{\rhcaTopDecRatioLgbmFullClosePre}{1.01}  % top forecast decile: mean RV / mean forecast
\providecommand{\rhcaTopDecSeLgbmFullClosePre}{}\renewcommand{\rhcaTopDecSeLgbmFullClosePre}{91}  % top forecast decile: % of the squared error
\providecommand{\rhcaBotDecRatioLgbmFullClosePre}{}\renewcommand{\rhcaBotDecRatioLgbmFullClosePre}{1.23}  % bottom forecast decile: mean RV / mean forecast
\providecommand{\rhcaTopDecRatioEnetHarCloseOos}{}\renewcommand{\rhcaTopDecRatioEnetHarCloseOos}{1.13}  % top forecast decile: mean RV / mean forecast
\providecommand{\rhcaTopDecSeEnetHarCloseOos}{}\renewcommand{\rhcaTopDecSeEnetHarCloseOos}{98}  % top forecast decile: % of the squared error
\providecommand{\rhcaBotDecRatioEnetHarCloseOos}{}\renewcommand{\rhcaBotDecRatioEnetHarCloseOos}{1.15}  % bottom forecast decile: mean RV / mean forecast
\providecommand{\rhcaTopDecRatioEnetFullCloseOos}{}\renewcommand{\rhcaTopDecRatioEnetFullCloseOos}{0.87}  % top forecast decile: mean RV / mean forecast
\providecommand{\rhcaTopDecSeEnetFullCloseOos}{}\renewcommand{\rhcaTopDecSeEnetFullCloseOos}{99}  % top forecast decile: % of the squared error
\providecommand{\rhcaBotDecRatioEnetFullCloseOos}{}\renewcommand{\rhcaBotDecRatioEnetFullCloseOos}{1.31}  % bottom forecast decile: mean RV / mean forecast
\providecommand{\rhcaTopDecRatioLgbmHarCloseOos}{}\renewcommand{\rhcaTopDecRatioLgbmHarCloseOos}{1.02}  % top forecast decile: mean RV / mean forecast
\providecommand{\rhcaTopDecSeLgbmHarCloseOos}{}\renewcommand{\rhcaTopDecSeLgbmHarCloseOos}{97}  % top forecast decile: % of the squared error
\providecommand{\rhcaBotDecRatioLgbmHarCloseOos}{}\renewcommand{\rhcaBotDecRatioLgbmHarCloseOos}{1.11}  % bottom forecast decile: mean RV / mean forecast
\providecommand{\rhcaTopDecRatioLgbmFullCloseOos}{}\renewcommand{\rhcaTopDecRatioLgbmFullCloseOos}{0.91}  % top forecast decile: mean RV / mean forecast
\providecommand{\rhcaTopDecSeLgbmFullCloseOos}{}\renewcommand{\rhcaTopDecSeLgbmFullCloseOos}{98}  % top forecast decile: % of the squared error
\providecommand{\rhcaBotDecRatioLgbmFullCloseOos}{}\renewcommand{\rhcaBotDecRatioLgbmFullCloseOos}{1.22}  % bottom forecast decile: mean RV / mean forecast
\providecommand{\rhcaHarLinGap}{}\renewcommand{\rhcaHarLinGap}{5 \times 10^{-6}}  % max relative gap, HAR linear vs OLS forecasts (math mode)
\fi
\ifdefined\rhmcsNModels\else% AUTO-GENERATED by writeup/make_rehaul_mcs.py -- do not edit; re-run the script.
% \rhmcs... macros: model confidence set (Hansen, Lunde and Nason 2011), SPA (Hansen 2005) and StepM (Romano and
% Wolf 2005) on the 12 pooled forecasts of the final rerun (the three linear HAR forecasts as one) and the
% closing-interval HAR benchmark. Samples: Bar (every bar, 2014-2024, daily-averaged losses), All (16:00, 2014-2024),
% Oos (16:00, T_oos), Ch (16:00, the benchmark's span). Losses: Qlike, Mse. MCS statistics: Tmax, Trange (T_R).
% SPA / StepM families: <Sample><Q|M>Af (Full vs own HAR), <Sample><Q|M>Bm (Full vs the closing-interval HAR),
% <Mid|Ba><Short|Twin><Mean|Sr> (sign strategies on the 614 post-warmup days, vs the short straddle or the HAR twin,
% mean return or the paper's pooled Sharpe ratio, D19). p-values of StepM and SPA are one-sided (H_k: no better than the benchmark).
\providecommand{\rhmcsB}{}\renewcommand{\rhmcsB}{4{,}999}  % bootstrap replications, every MCS / SPA / StepM (accuracy and strategies)
\providecommand{\rhmcsBlock}{}\renewcommand{\rhmcsBlock}{10}  % block length in trading days (circular moving-block bootstrap), main
\providecommand{\rhmcsBlockLo}{}\renewcommand{\rhmcsBlockLo}{5}  % block length, sensitivity (shorter)
\providecommand{\rhmcsBlockHi}{}\renewcommand{\rhmcsBlockHi}{20}  % block length, sensitivity (longer)
\providecommand{\rhmcsSeed}{}\renewcommand{\rhmcsSeed}{20261025}  % random seed of the accuracy bootstrap (MCS, SPA, StepM on losses)
\providecommand{\rhmcsSeedEcon}{}\renewcommand{\rhmcsSeedEcon}{20261022}  % random seed of the strategy bootstrap (= the paper's D6 bootstrap seed)
\providecommand{\rhmcsNModels}{}\renewcommand{\rhmcsNModels}{10}  % distinct pooled forecasts in the MCS (the three linear HAR forecasts as one)
\providecommand{\rhmcsNModelsWord}{}\renewcommand{\rhmcsNModelsWord}{ten}  % rhmcsNModels as a word
\providecommand{\rhmcsNModelsBench}{}\renewcommand{\rhmcsNModelsBench}{11}  % ... with the closing-interval HAR benchmark
\providecommand{\rhmcsNModelsBenchWord}{}\renewcommand{\rhmcsNModelsBenchWord}{eleven}  % rhmcsNModelsBench as a word
\providecommand{\rhmcsNPool}{}\renewcommand{\rhmcsNPool}{12}  % pooled forecasts (6 classes x HAR, Full)
\providecommand{\rhmcsNAlt}{}\renewcommand{\rhmcsNAlt}{6}  % alternatives in every SPA / StepM family (the six full-set forecasts or strategies)
\providecommand{\rhmcsNAltWord}{}\renewcommand{\rhmcsNAltWord}{six}  % rhmcsNAlt as a word
\providecommand{\rhmcsLinGap}{}\renewcommand{\rhmcsLinGap}{5 \times 10^{-6}}  % largest relative gap between the daily every-bar QLIKE of a linear HAR forecast and the pooled OLS (math mode)
\providecommand{\rhmcsDaysAll}{}\renewcommand{\rhmcsDaysAll}{2{,}579}  % trading days 2014-2024 (Table 1's span)
\providecommand{\rhmcsDaysOos}{}\renewcommand{\rhmcsDaysOos}{1{,}082}  % trading days of T_oos
\providecommand{\rhmcsDaysCh}{}\renewcommand{\rhmcsDaysCh}{2{,}579}  % trading days of the closing-interval benchmark's span within 2014-2024
\providecommand{\rhmcsChStart}{}\renewcommand{\rhmcsChStart}{2 January 2014}  % first day of the closing-interval HAR benchmark's forecasts
\providecommand{\rhmcsNBars}{}\renewcommand{\rhmcsNBars}{33{,}527}  % bars of the every-bar sample (13 a day)
\providecommand{\rhmcsNPw}{}\renewcommand{\rhmcsNPw}{614}  % post-warmup option days of the strategy tests
\providecommand{\rhmcsLinSameTrades}{}\renewcommand{\rhmcsLinSameTrades}{yes}  % the three linear HAR sign strategies have identical daily returns (mid and bid/ask)
\providecommand{\rhmcsBarQlikeTmaxPLinHar}{}\renewcommand{\rhmcsBarQlikeTmaxPLinHar}{\ensuremath{<}0.001}  % MCS p-value, every bar, 2014--2024, Qlike, Tmax: Linear, HAR
\providecommand{\rhmcsBarQlikeTmaxPRidgeFull}{}\renewcommand{\rhmcsBarQlikeTmaxPRidgeFull}{\ensuremath{<}0.001}  % MCS p-value, every bar, 2014--2024, Qlike, Tmax: Ridge, Full
\providecommand{\rhmcsBarQlikeTmaxPLassoFull}{}\renewcommand{\rhmcsBarQlikeTmaxPLassoFull}{\ensuremath{<}0.001}  % MCS p-value, every bar, 2014--2024, Qlike, Tmax: Lasso, Full
\providecommand{\rhmcsBarQlikeTmaxPEnetFull}{}\renewcommand{\rhmcsBarQlikeTmaxPEnetFull}{\ensuremath{<}0.001}  % MCS p-value, every bar, 2014--2024, Qlike, Tmax: Elastic net, Full
\providecommand{\rhmcsBarQlikeTmaxPLgbmHar}{}\renewcommand{\rhmcsBarQlikeTmaxPLgbmHar}{\ensuremath{<}0.001}  % MCS p-value, every bar, 2014--2024, Qlike, Tmax: LightGBM, HAR
\providecommand{\rhmcsBarQlikeTmaxPLgbmFull}{}\renewcommand{\rhmcsBarQlikeTmaxPLgbmFull}{1.00}  % MCS p-value, every bar, 2014--2024, Qlike, Tmax: LightGBM, Full
\providecommand{\rhmcsBarQlikeTmaxPXgbHar}{}\renewcommand{\rhmcsBarQlikeTmaxPXgbHar}{\ensuremath{<}0.001}  % MCS p-value, every bar, 2014--2024, Qlike, Tmax: XGBoost, HAR
\providecommand{\rhmcsBarQlikeTmaxPXgbFull}{}\renewcommand{\rhmcsBarQlikeTmaxPXgbFull}{0.69}  % MCS p-value, every bar, 2014--2024, Qlike, Tmax: XGBoost, Full
\providecommand{\rhmcsBarQlikeTmaxPRfHar}{}\renewcommand{\rhmcsBarQlikeTmaxPRfHar}{\ensuremath{<}0.001}  % MCS p-value, every bar, 2014--2024, Qlike, Tmax: Random forest, HAR
\providecommand{\rhmcsBarQlikeTmaxPRfFull}{}\renewcommand{\rhmcsBarQlikeTmaxPRfFull}{\ensuremath{<}0.001}  % MCS p-value, every bar, 2014--2024, Qlike, Tmax: Random forest, Full
\providecommand{\rhmcsBarQlikeTmaxNNinety}{}\renewcommand{\rhmcsBarQlikeTmaxNNinety}{2}  % models in the 90 % MCS (of 10)
\providecommand{\rhmcsBarQlikeTmaxSetNinety}{}\renewcommand{\rhmcsBarQlikeTmaxSetNinety}{LightGBM on the full predictor set and XGBoost on the full predictor set}  % the 90 % MCS, best (last eliminated) first
\providecommand{\rhmcsBarQlikeTmaxStableNinety}{}\renewcommand{\rhmcsBarQlikeTmaxStableNinety}{yes}  % the 90 % MCS is the same at block lengths (5, 20)
\providecommand{\rhmcsBarQlikeTmaxPMaxOutNinety}{}\renewcommand{\rhmcsBarQlikeTmaxPMaxOutNinety}{\ensuremath{<}0.001}  % largest MCS p-value outside the 90 % set
\providecommand{\rhmcsBarQlikeTmaxNNinetyLo}{}\renewcommand{\rhmcsBarQlikeTmaxNNinetyLo}{2}  % ... at block length 5
\providecommand{\rhmcsBarQlikeTmaxSetNinetyLo}{}\renewcommand{\rhmcsBarQlikeTmaxSetNinetyLo}{LightGBM on the full predictor set and XGBoost on the full predictor set}  % ... at block length 5
\providecommand{\rhmcsBarQlikeTmaxNNinetyHi}{}\renewcommand{\rhmcsBarQlikeTmaxNNinetyHi}{2}  % ... at block length 20
\providecommand{\rhmcsBarQlikeTmaxSetNinetyHi}{}\renewcommand{\rhmcsBarQlikeTmaxSetNinetyHi}{LightGBM on the full predictor set and XGBoost on the full predictor set}  % ... at block length 20
\providecommand{\rhmcsBarQlikeTmaxNSevenFive}{}\renewcommand{\rhmcsBarQlikeTmaxNSevenFive}{2}  % models in the 75 % MCS (of 10)
\providecommand{\rhmcsBarQlikeTmaxSetSevenFive}{}\renewcommand{\rhmcsBarQlikeTmaxSetSevenFive}{LightGBM on the full predictor set and XGBoost on the full predictor set}  % the 75 % MCS, best (last eliminated) first
\providecommand{\rhmcsBarQlikeTmaxStableSevenFive}{}\renewcommand{\rhmcsBarQlikeTmaxStableSevenFive}{yes}  % the 75 % MCS is the same at block lengths (5, 20)
\providecommand{\rhmcsBarQlikeTmaxPMaxOutSevenFive}{}\renewcommand{\rhmcsBarQlikeTmaxPMaxOutSevenFive}{\ensuremath{<}0.001}  % largest MCS p-value outside the 75 % set
\providecommand{\rhmcsBarQlikeTmaxNSevenFiveLo}{}\renewcommand{\rhmcsBarQlikeTmaxNSevenFiveLo}{2}  % ... at block length 5
\providecommand{\rhmcsBarQlikeTmaxSetSevenFiveLo}{}\renewcommand{\rhmcsBarQlikeTmaxSetSevenFiveLo}{LightGBM on the full predictor set and XGBoost on the full predictor set}  % ... at block length 5
\providecommand{\rhmcsBarQlikeTmaxNSevenFiveHi}{}\renewcommand{\rhmcsBarQlikeTmaxNSevenFiveHi}{2}  % ... at block length 20
\providecommand{\rhmcsBarQlikeTmaxSetSevenFiveHi}{}\renewcommand{\rhmcsBarQlikeTmaxSetSevenFiveHi}{LightGBM on the full predictor set and XGBoost on the full predictor set}  % ... at block length 20
\providecommand{\rhmcsBarQlikeTmaxBest}{}\renewcommand{\rhmcsBarQlikeTmaxBest}{LightGBM on the full predictor set}  % the last model left (MCS p-value 1)
\providecommand{\rhmcsBarQlikeTrangePLinHar}{}\renewcommand{\rhmcsBarQlikeTrangePLinHar}{\ensuremath{<}0.001}  % MCS p-value, every bar, 2014--2024, Qlike, Trange: Linear, HAR
\providecommand{\rhmcsBarQlikeTrangePRidgeFull}{}\renewcommand{\rhmcsBarQlikeTrangePRidgeFull}{\ensuremath{<}0.001}  % MCS p-value, every bar, 2014--2024, Qlike, Trange: Ridge, Full
\providecommand{\rhmcsBarQlikeTrangePLassoFull}{}\renewcommand{\rhmcsBarQlikeTrangePLassoFull}{\ensuremath{<}0.001}  % MCS p-value, every bar, 2014--2024, Qlike, Trange: Lasso, Full
\providecommand{\rhmcsBarQlikeTrangePEnetFull}{}\renewcommand{\rhmcsBarQlikeTrangePEnetFull}{\ensuremath{<}0.001}  % MCS p-value, every bar, 2014--2024, Qlike, Trange: Elastic net, Full
\providecommand{\rhmcsBarQlikeTrangePLgbmHar}{}\renewcommand{\rhmcsBarQlikeTrangePLgbmHar}{\ensuremath{<}0.001}  % MCS p-value, every bar, 2014--2024, Qlike, Trange: LightGBM, HAR
\providecommand{\rhmcsBarQlikeTrangePLgbmFull}{}\renewcommand{\rhmcsBarQlikeTrangePLgbmFull}{1.00}  % MCS p-value, every bar, 2014--2024, Qlike, Trange: LightGBM, Full
\providecommand{\rhmcsBarQlikeTrangePXgbHar}{}\renewcommand{\rhmcsBarQlikeTrangePXgbHar}{\ensuremath{<}0.001}  % MCS p-value, every bar, 2014--2024, Qlike, Trange: XGBoost, HAR
\providecommand{\rhmcsBarQlikeTrangePXgbFull}{}\renewcommand{\rhmcsBarQlikeTrangePXgbFull}{0.69}  % MCS p-value, every bar, 2014--2024, Qlike, Trange: XGBoost, Full
\providecommand{\rhmcsBarQlikeTrangePRfHar}{}\renewcommand{\rhmcsBarQlikeTrangePRfHar}{\ensuremath{<}0.001}  % MCS p-value, every bar, 2014--2024, Qlike, Trange: Random forest, HAR
\providecommand{\rhmcsBarQlikeTrangePRfFull}{}\renewcommand{\rhmcsBarQlikeTrangePRfFull}{\ensuremath{<}0.001}  % MCS p-value, every bar, 2014--2024, Qlike, Trange: Random forest, Full
\providecommand{\rhmcsBarQlikeTrangeNNinety}{}\renewcommand{\rhmcsBarQlikeTrangeNNinety}{2}  % models in the 90 % MCS (of 10)
\providecommand{\rhmcsBarQlikeTrangeSetNinety}{}\renewcommand{\rhmcsBarQlikeTrangeSetNinety}{LightGBM on the full predictor set and XGBoost on the full predictor set}  % the 90 % MCS, best (last eliminated) first
\providecommand{\rhmcsBarQlikeTrangeStableNinety}{}\renewcommand{\rhmcsBarQlikeTrangeStableNinety}{yes}  % the 90 % MCS is the same at block lengths (5, 20)
\providecommand{\rhmcsBarQlikeTrangePMaxOutNinety}{}\renewcommand{\rhmcsBarQlikeTrangePMaxOutNinety}{\ensuremath{<}0.001}  % largest MCS p-value outside the 90 % set
\providecommand{\rhmcsBarQlikeTrangeNNinetyLo}{}\renewcommand{\rhmcsBarQlikeTrangeNNinetyLo}{2}  % ... at block length 5
\providecommand{\rhmcsBarQlikeTrangeSetNinetyLo}{}\renewcommand{\rhmcsBarQlikeTrangeSetNinetyLo}{LightGBM on the full predictor set and XGBoost on the full predictor set}  % ... at block length 5
\providecommand{\rhmcsBarQlikeTrangeNNinetyHi}{}\renewcommand{\rhmcsBarQlikeTrangeNNinetyHi}{2}  % ... at block length 20
\providecommand{\rhmcsBarQlikeTrangeSetNinetyHi}{}\renewcommand{\rhmcsBarQlikeTrangeSetNinetyHi}{LightGBM on the full predictor set and XGBoost on the full predictor set}  % ... at block length 20
\providecommand{\rhmcsBarQlikeTrangeNSevenFive}{}\renewcommand{\rhmcsBarQlikeTrangeNSevenFive}{2}  % models in the 75 % MCS (of 10)
\providecommand{\rhmcsBarQlikeTrangeSetSevenFive}{}\renewcommand{\rhmcsBarQlikeTrangeSetSevenFive}{LightGBM on the full predictor set and XGBoost on the full predictor set}  % the 75 % MCS, best (last eliminated) first
\providecommand{\rhmcsBarQlikeTrangeStableSevenFive}{}\renewcommand{\rhmcsBarQlikeTrangeStableSevenFive}{yes}  % the 75 % MCS is the same at block lengths (5, 20)
\providecommand{\rhmcsBarQlikeTrangePMaxOutSevenFive}{}\renewcommand{\rhmcsBarQlikeTrangePMaxOutSevenFive}{\ensuremath{<}0.001}  % largest MCS p-value outside the 75 % set
\providecommand{\rhmcsBarQlikeTrangeNSevenFiveLo}{}\renewcommand{\rhmcsBarQlikeTrangeNSevenFiveLo}{2}  % ... at block length 5
\providecommand{\rhmcsBarQlikeTrangeSetSevenFiveLo}{}\renewcommand{\rhmcsBarQlikeTrangeSetSevenFiveLo}{LightGBM on the full predictor set and XGBoost on the full predictor set}  % ... at block length 5
\providecommand{\rhmcsBarQlikeTrangeNSevenFiveHi}{}\renewcommand{\rhmcsBarQlikeTrangeNSevenFiveHi}{2}  % ... at block length 20
\providecommand{\rhmcsBarQlikeTrangeSetSevenFiveHi}{}\renewcommand{\rhmcsBarQlikeTrangeSetSevenFiveHi}{LightGBM on the full predictor set and XGBoost on the full predictor set}  % ... at block length 20
\providecommand{\rhmcsBarQlikeTrangeBest}{}\renewcommand{\rhmcsBarQlikeTrangeBest}{LightGBM on the full predictor set}  % the last model left (MCS p-value 1)
\providecommand{\rhmcsBarMseTmaxPLinHar}{}\renewcommand{\rhmcsBarMseTmaxPLinHar}{1.00}  % MCS p-value, every bar, 2014--2024, Mse, Tmax: Linear, HAR
\providecommand{\rhmcsBarMseTmaxPRidgeFull}{}\renewcommand{\rhmcsBarMseTmaxPRidgeFull}{0.49}  % MCS p-value, every bar, 2014--2024, Mse, Tmax: Ridge, Full
\providecommand{\rhmcsBarMseTmaxPLassoFull}{}\renewcommand{\rhmcsBarMseTmaxPLassoFull}{0.32}  % MCS p-value, every bar, 2014--2024, Mse, Tmax: Lasso, Full
\providecommand{\rhmcsBarMseTmaxPEnetFull}{}\renewcommand{\rhmcsBarMseTmaxPEnetFull}{0.43}  % MCS p-value, every bar, 2014--2024, Mse, Tmax: Elastic net, Full
\providecommand{\rhmcsBarMseTmaxPLgbmHar}{}\renewcommand{\rhmcsBarMseTmaxPLgbmHar}{0.49}  % MCS p-value, every bar, 2014--2024, Mse, Tmax: LightGBM, HAR
\providecommand{\rhmcsBarMseTmaxPLgbmFull}{}\renewcommand{\rhmcsBarMseTmaxPLgbmFull}{0.49}  % MCS p-value, every bar, 2014--2024, Mse, Tmax: LightGBM, Full
\providecommand{\rhmcsBarMseTmaxPXgbHar}{}\renewcommand{\rhmcsBarMseTmaxPXgbHar}{0.49}  % MCS p-value, every bar, 2014--2024, Mse, Tmax: XGBoost, HAR
\providecommand{\rhmcsBarMseTmaxPXgbFull}{}\renewcommand{\rhmcsBarMseTmaxPXgbFull}{0.49}  % MCS p-value, every bar, 2014--2024, Mse, Tmax: XGBoost, Full
\providecommand{\rhmcsBarMseTmaxPRfHar}{}\renewcommand{\rhmcsBarMseTmaxPRfHar}{0.49}  % MCS p-value, every bar, 2014--2024, Mse, Tmax: Random forest, HAR
\providecommand{\rhmcsBarMseTmaxPRfFull}{}\renewcommand{\rhmcsBarMseTmaxPRfFull}{0.49}  % MCS p-value, every bar, 2014--2024, Mse, Tmax: Random forest, Full
\providecommand{\rhmcsBarMseTmaxNNinety}{}\renewcommand{\rhmcsBarMseTmaxNNinety}{10}  % models in the 90 % MCS (of 10)
\providecommand{\rhmcsBarMseTmaxSetNinety}{}\renewcommand{\rhmcsBarMseTmaxSetNinety}{the linear models on the HAR benchmark, XGBoost on the HAR benchmark, LightGBM on the HAR benchmark, the random forest on the HAR benchmark, XGBoost on the full predictor set, LightGBM on the full predictor set, the random forest on the full predictor set, ridge regression on the full predictor set, the elastic net on the full predictor set, and the lasso on the full predictor set}  % the 90 % MCS, best (last eliminated) first
\providecommand{\rhmcsBarMseTmaxStableNinety}{}\renewcommand{\rhmcsBarMseTmaxStableNinety}{yes}  % the 90 % MCS is the same at block lengths (5, 20)
\providecommand{\rhmcsBarMseTmaxPMaxOutNinety}{}\renewcommand{\rhmcsBarMseTmaxPMaxOutNinety}{--}  % largest MCS p-value outside the 90 % set
\providecommand{\rhmcsBarMseTmaxNNinetyLo}{}\renewcommand{\rhmcsBarMseTmaxNNinetyLo}{10}  % ... at block length 5
\providecommand{\rhmcsBarMseTmaxSetNinetyLo}{}\renewcommand{\rhmcsBarMseTmaxSetNinetyLo}{the linear models on the HAR benchmark, XGBoost on the HAR benchmark, LightGBM on the HAR benchmark, the random forest on the HAR benchmark, XGBoost on the full predictor set, LightGBM on the full predictor set, the random forest on the full predictor set, the elastic net on the full predictor set, ridge regression on the full predictor set, and the lasso on the full predictor set}  % ... at block length 5
\providecommand{\rhmcsBarMseTmaxNNinetyHi}{}\renewcommand{\rhmcsBarMseTmaxNNinetyHi}{10}  % ... at block length 20
\providecommand{\rhmcsBarMseTmaxSetNinetyHi}{}\renewcommand{\rhmcsBarMseTmaxSetNinetyHi}{the linear models on the HAR benchmark, XGBoost on the HAR benchmark, LightGBM on the HAR benchmark, the random forest on the HAR benchmark, XGBoost on the full predictor set, LightGBM on the full predictor set, the random forest on the full predictor set, ridge regression on the full predictor set, the elastic net on the full predictor set, and the lasso on the full predictor set}  % ... at block length 20
\providecommand{\rhmcsBarMseTmaxNSevenFive}{}\renewcommand{\rhmcsBarMseTmaxNSevenFive}{10}  % models in the 75 % MCS (of 10)
\providecommand{\rhmcsBarMseTmaxSetSevenFive}{}\renewcommand{\rhmcsBarMseTmaxSetSevenFive}{the linear models on the HAR benchmark, XGBoost on the HAR benchmark, LightGBM on the HAR benchmark, the random forest on the HAR benchmark, XGBoost on the full predictor set, LightGBM on the full predictor set, the random forest on the full predictor set, ridge regression on the full predictor set, the elastic net on the full predictor set, and the lasso on the full predictor set}  % the 75 % MCS, best (last eliminated) first
\providecommand{\rhmcsBarMseTmaxStableSevenFive}{}\renewcommand{\rhmcsBarMseTmaxStableSevenFive}{no}  % the 75 % MCS is the same at block lengths (5, 20)
\providecommand{\rhmcsBarMseTmaxPMaxOutSevenFive}{}\renewcommand{\rhmcsBarMseTmaxPMaxOutSevenFive}{--}  % largest MCS p-value outside the 75 % set
\providecommand{\rhmcsBarMseTmaxNSevenFiveLo}{}\renewcommand{\rhmcsBarMseTmaxNSevenFiveLo}{9}  % ... at block length 5
\providecommand{\rhmcsBarMseTmaxSetSevenFiveLo}{}\renewcommand{\rhmcsBarMseTmaxSetSevenFiveLo}{the linear models on the HAR benchmark, XGBoost on the HAR benchmark, LightGBM on the HAR benchmark, the random forest on the HAR benchmark, XGBoost on the full predictor set, LightGBM on the full predictor set, the random forest on the full predictor set, the elastic net on the full predictor set, and ridge regression on the full predictor set}  % ... at block length 5
\providecommand{\rhmcsBarMseTmaxNSevenFiveHi}{}\renewcommand{\rhmcsBarMseTmaxNSevenFiveHi}{10}  % ... at block length 20
\providecommand{\rhmcsBarMseTmaxSetSevenFiveHi}{}\renewcommand{\rhmcsBarMseTmaxSetSevenFiveHi}{the linear models on the HAR benchmark, XGBoost on the HAR benchmark, LightGBM on the HAR benchmark, the random forest on the HAR benchmark, XGBoost on the full predictor set, LightGBM on the full predictor set, the random forest on the full predictor set, ridge regression on the full predictor set, the elastic net on the full predictor set, and the lasso on the full predictor set}  % ... at block length 20
\providecommand{\rhmcsBarMseTmaxBest}{}\renewcommand{\rhmcsBarMseTmaxBest}{the linear models on the HAR benchmark}  % the last model left (MCS p-value 1)
\providecommand{\rhmcsBarMseTrangePLinHar}{}\renewcommand{\rhmcsBarMseTrangePLinHar}{1.00}  % MCS p-value, every bar, 2014--2024, Mse, Trange: Linear, HAR
\providecommand{\rhmcsBarMseTrangePRidgeFull}{}\renewcommand{\rhmcsBarMseTrangePRidgeFull}{0.42}  % MCS p-value, every bar, 2014--2024, Mse, Trange: Ridge, Full
\providecommand{\rhmcsBarMseTrangePLassoFull}{}\renewcommand{\rhmcsBarMseTrangePLassoFull}{0.42}  % MCS p-value, every bar, 2014--2024, Mse, Trange: Lasso, Full
\providecommand{\rhmcsBarMseTrangePEnetFull}{}\renewcommand{\rhmcsBarMseTrangePEnetFull}{0.42}  % MCS p-value, every bar, 2014--2024, Mse, Trange: Elastic net, Full
\providecommand{\rhmcsBarMseTrangePLgbmHar}{}\renewcommand{\rhmcsBarMseTrangePLgbmHar}{0.31}  % MCS p-value, every bar, 2014--2024, Mse, Trange: LightGBM, HAR
\providecommand{\rhmcsBarMseTrangePLgbmFull}{}\renewcommand{\rhmcsBarMseTrangePLgbmFull}{0.42}  % MCS p-value, every bar, 2014--2024, Mse, Trange: LightGBM, Full
\providecommand{\rhmcsBarMseTrangePXgbHar}{}\renewcommand{\rhmcsBarMseTrangePXgbHar}{0.42}  % MCS p-value, every bar, 2014--2024, Mse, Trange: XGBoost, HAR
\providecommand{\rhmcsBarMseTrangePXgbFull}{}\renewcommand{\rhmcsBarMseTrangePXgbFull}{0.42}  % MCS p-value, every bar, 2014--2024, Mse, Trange: XGBoost, Full
\providecommand{\rhmcsBarMseTrangePRfHar}{}\renewcommand{\rhmcsBarMseTrangePRfHar}{0.31}  % MCS p-value, every bar, 2014--2024, Mse, Trange: Random forest, HAR
\providecommand{\rhmcsBarMseTrangePRfFull}{}\renewcommand{\rhmcsBarMseTrangePRfFull}{0.42}  % MCS p-value, every bar, 2014--2024, Mse, Trange: Random forest, Full
\providecommand{\rhmcsBarMseTrangeNNinety}{}\renewcommand{\rhmcsBarMseTrangeNNinety}{10}  % models in the 90 % MCS (of 10)
\providecommand{\rhmcsBarMseTrangeSetNinety}{}\renewcommand{\rhmcsBarMseTrangeSetNinety}{the linear models on the HAR benchmark, the elastic net on the full predictor set, ridge regression on the full predictor set, XGBoost on the full predictor set, LightGBM on the full predictor set, the random forest on the full predictor set, the lasso on the full predictor set, XGBoost on the HAR benchmark, LightGBM on the HAR benchmark, and the random forest on the HAR benchmark}  % the 90 % MCS, best (last eliminated) first
\providecommand{\rhmcsBarMseTrangeStableNinety}{}\renewcommand{\rhmcsBarMseTrangeStableNinety}{yes}  % the 90 % MCS is the same at block lengths (5, 20)
\providecommand{\rhmcsBarMseTrangePMaxOutNinety}{}\renewcommand{\rhmcsBarMseTrangePMaxOutNinety}{--}  % largest MCS p-value outside the 90 % set
\providecommand{\rhmcsBarMseTrangeNNinetyLo}{}\renewcommand{\rhmcsBarMseTrangeNNinetyLo}{10}  % ... at block length 5
\providecommand{\rhmcsBarMseTrangeSetNinetyLo}{}\renewcommand{\rhmcsBarMseTrangeSetNinetyLo}{the linear models on the HAR benchmark, the elastic net on the full predictor set, ridge regression on the full predictor set, XGBoost on the HAR benchmark, XGBoost on the full predictor set, LightGBM on the full predictor set, the random forest on the full predictor set, LightGBM on the HAR benchmark, the lasso on the full predictor set, and the random forest on the HAR benchmark}  % ... at block length 5
\providecommand{\rhmcsBarMseTrangeNNinetyHi}{}\renewcommand{\rhmcsBarMseTrangeNNinetyHi}{10}  % ... at block length 20
\providecommand{\rhmcsBarMseTrangeSetNinetyHi}{}\renewcommand{\rhmcsBarMseTrangeSetNinetyHi}{the linear models on the HAR benchmark, the elastic net on the full predictor set, ridge regression on the full predictor set, XGBoost on the full predictor set, the lasso on the full predictor set, LightGBM on the full predictor set, the random forest on the full predictor set, XGBoost on the HAR benchmark, LightGBM on the HAR benchmark, and the random forest on the HAR benchmark}  % ... at block length 20
\providecommand{\rhmcsBarMseTrangeNSevenFive}{}\renewcommand{\rhmcsBarMseTrangeNSevenFive}{10}  % models in the 75 % MCS (of 10)
\providecommand{\rhmcsBarMseTrangeSetSevenFive}{}\renewcommand{\rhmcsBarMseTrangeSetSevenFive}{the linear models on the HAR benchmark, the elastic net on the full predictor set, ridge regression on the full predictor set, XGBoost on the full predictor set, LightGBM on the full predictor set, the random forest on the full predictor set, the lasso on the full predictor set, XGBoost on the HAR benchmark, LightGBM on the HAR benchmark, and the random forest on the HAR benchmark}  % the 75 % MCS, best (last eliminated) first
\providecommand{\rhmcsBarMseTrangeStableSevenFive}{}\renewcommand{\rhmcsBarMseTrangeStableSevenFive}{yes}  % the 75 % MCS is the same at block lengths (5, 20)
\providecommand{\rhmcsBarMseTrangePMaxOutSevenFive}{}\renewcommand{\rhmcsBarMseTrangePMaxOutSevenFive}{--}  % largest MCS p-value outside the 75 % set
\providecommand{\rhmcsBarMseTrangeNSevenFiveLo}{}\renewcommand{\rhmcsBarMseTrangeNSevenFiveLo}{10}  % ... at block length 5
\providecommand{\rhmcsBarMseTrangeSetSevenFiveLo}{}\renewcommand{\rhmcsBarMseTrangeSetSevenFiveLo}{the linear models on the HAR benchmark, the elastic net on the full predictor set, ridge regression on the full predictor set, XGBoost on the HAR benchmark, XGBoost on the full predictor set, LightGBM on the full predictor set, the random forest on the full predictor set, LightGBM on the HAR benchmark, the lasso on the full predictor set, and the random forest on the HAR benchmark}  % ... at block length 5
\providecommand{\rhmcsBarMseTrangeNSevenFiveHi}{}\renewcommand{\rhmcsBarMseTrangeNSevenFiveHi}{10}  % ... at block length 20
\providecommand{\rhmcsBarMseTrangeSetSevenFiveHi}{}\renewcommand{\rhmcsBarMseTrangeSetSevenFiveHi}{the linear models on the HAR benchmark, the elastic net on the full predictor set, ridge regression on the full predictor set, XGBoost on the full predictor set, the lasso on the full predictor set, LightGBM on the full predictor set, the random forest on the full predictor set, XGBoost on the HAR benchmark, LightGBM on the HAR benchmark, and the random forest on the HAR benchmark}  % ... at block length 20
\providecommand{\rhmcsBarMseTrangeBest}{}\renewcommand{\rhmcsBarMseTrangeBest}{the linear models on the HAR benchmark}  % the last model left (MCS p-value 1)
\providecommand{\rhmcsAllQlikeTmaxPLinHar}{}\renewcommand{\rhmcsAllQlikeTmaxPLinHar}{0.19}  % MCS p-value, 16:00 bar, 2014--2024, Qlike, Tmax: Linear, HAR
\providecommand{\rhmcsAllQlikeTmaxPRidgeFull}{}\renewcommand{\rhmcsAllQlikeTmaxPRidgeFull}{0.19}  % MCS p-value, 16:00 bar, 2014--2024, Qlike, Tmax: Ridge, Full
\providecommand{\rhmcsAllQlikeTmaxPLassoFull}{}\renewcommand{\rhmcsAllQlikeTmaxPLassoFull}{0.071}  % MCS p-value, 16:00 bar, 2014--2024, Qlike, Tmax: Lasso, Full
\providecommand{\rhmcsAllQlikeTmaxPEnetFull}{}\renewcommand{\rhmcsAllQlikeTmaxPEnetFull}{0.19}  % MCS p-value, 16:00 bar, 2014--2024, Qlike, Tmax: Elastic net, Full
\providecommand{\rhmcsAllQlikeTmaxPLgbmHar}{}\renewcommand{\rhmcsAllQlikeTmaxPLgbmHar}{0.19}  % MCS p-value, 16:00 bar, 2014--2024, Qlike, Tmax: LightGBM, HAR
\providecommand{\rhmcsAllQlikeTmaxPLgbmFull}{}\renewcommand{\rhmcsAllQlikeTmaxPLgbmFull}{0.29}  % MCS p-value, 16:00 bar, 2014--2024, Qlike, Tmax: LightGBM, Full
\providecommand{\rhmcsAllQlikeTmaxPXgbHar}{}\renewcommand{\rhmcsAllQlikeTmaxPXgbHar}{0.19}  % MCS p-value, 16:00 bar, 2014--2024, Qlike, Tmax: XGBoost, HAR
\providecommand{\rhmcsAllQlikeTmaxPXgbFull}{}\renewcommand{\rhmcsAllQlikeTmaxPXgbFull}{1.00}  % MCS p-value, 16:00 bar, 2014--2024, Qlike, Tmax: XGBoost, Full
\providecommand{\rhmcsAllQlikeTmaxPRfHar}{}\renewcommand{\rhmcsAllQlikeTmaxPRfHar}{0.19}  % MCS p-value, 16:00 bar, 2014--2024, Qlike, Tmax: Random forest, HAR
\providecommand{\rhmcsAllQlikeTmaxPRfFull}{}\renewcommand{\rhmcsAllQlikeTmaxPRfFull}{0.19}  % MCS p-value, 16:00 bar, 2014--2024, Qlike, Tmax: Random forest, Full
\providecommand{\rhmcsAllQlikeTmaxNNinety}{}\renewcommand{\rhmcsAllQlikeTmaxNNinety}{9}  % models in the 90 % MCS (of 10)
\providecommand{\rhmcsAllQlikeTmaxSetNinety}{}\renewcommand{\rhmcsAllQlikeTmaxSetNinety}{XGBoost on the full predictor set, LightGBM on the full predictor set, XGBoost on the HAR benchmark, LightGBM on the HAR benchmark, the random forest on the full predictor set, the elastic net on the full predictor set, the random forest on the HAR benchmark, ridge regression on the full predictor set, and the linear models on the HAR benchmark}  % the 90 % MCS, best (last eliminated) first
\providecommand{\rhmcsAllQlikeTmaxStableNinety}{}\renewcommand{\rhmcsAllQlikeTmaxStableNinety}{yes}  % the 90 % MCS is the same at block lengths (5, 20)
\providecommand{\rhmcsAllQlikeTmaxPMaxOutNinety}{}\renewcommand{\rhmcsAllQlikeTmaxPMaxOutNinety}{0.071}  % largest MCS p-value outside the 90 % set
\providecommand{\rhmcsAllQlikeTmaxNNinetyLo}{}\renewcommand{\rhmcsAllQlikeTmaxNNinetyLo}{9}  % ... at block length 5
\providecommand{\rhmcsAllQlikeTmaxSetNinetyLo}{}\renewcommand{\rhmcsAllQlikeTmaxSetNinetyLo}{XGBoost on the full predictor set, LightGBM on the full predictor set, XGBoost on the HAR benchmark, LightGBM on the HAR benchmark, the random forest on the full predictor set, the elastic net on the full predictor set, the random forest on the HAR benchmark, ridge regression on the full predictor set, and the linear models on the HAR benchmark}  % ... at block length 5
\providecommand{\rhmcsAllQlikeTmaxNNinetyHi}{}\renewcommand{\rhmcsAllQlikeTmaxNNinetyHi}{9}  % ... at block length 20
\providecommand{\rhmcsAllQlikeTmaxSetNinetyHi}{}\renewcommand{\rhmcsAllQlikeTmaxSetNinetyHi}{XGBoost on the full predictor set, LightGBM on the full predictor set, XGBoost on the HAR benchmark, LightGBM on the HAR benchmark, the random forest on the full predictor set, the elastic net on the full predictor set, the random forest on the HAR benchmark, ridge regression on the full predictor set, and the linear models on the HAR benchmark}  % ... at block length 20
\providecommand{\rhmcsAllQlikeTmaxNSevenFive}{}\renewcommand{\rhmcsAllQlikeTmaxNSevenFive}{2}  % models in the 75 % MCS (of 10)
\providecommand{\rhmcsAllQlikeTmaxSetSevenFive}{}\renewcommand{\rhmcsAllQlikeTmaxSetSevenFive}{XGBoost on the full predictor set and LightGBM on the full predictor set}  % the 75 % MCS, best (last eliminated) first
\providecommand{\rhmcsAllQlikeTmaxStableSevenFive}{}\renewcommand{\rhmcsAllQlikeTmaxStableSevenFive}{yes}  % the 75 % MCS is the same at block lengths (5, 20)
\providecommand{\rhmcsAllQlikeTmaxPMaxOutSevenFive}{}\renewcommand{\rhmcsAllQlikeTmaxPMaxOutSevenFive}{0.19}  % largest MCS p-value outside the 75 % set
\providecommand{\rhmcsAllQlikeTmaxNSevenFiveLo}{}\renewcommand{\rhmcsAllQlikeTmaxNSevenFiveLo}{2}  % ... at block length 5
\providecommand{\rhmcsAllQlikeTmaxSetSevenFiveLo}{}\renewcommand{\rhmcsAllQlikeTmaxSetSevenFiveLo}{XGBoost on the full predictor set and LightGBM on the full predictor set}  % ... at block length 5
\providecommand{\rhmcsAllQlikeTmaxNSevenFiveHi}{}\renewcommand{\rhmcsAllQlikeTmaxNSevenFiveHi}{2}  % ... at block length 20
\providecommand{\rhmcsAllQlikeTmaxSetSevenFiveHi}{}\renewcommand{\rhmcsAllQlikeTmaxSetSevenFiveHi}{XGBoost on the full predictor set and LightGBM on the full predictor set}  % ... at block length 20
\providecommand{\rhmcsAllQlikeTmaxBest}{}\renewcommand{\rhmcsAllQlikeTmaxBest}{XGBoost on the full predictor set}  % the last model left (MCS p-value 1)
\providecommand{\rhmcsAllQlikeTrangePLinHar}{}\renewcommand{\rhmcsAllQlikeTrangePLinHar}{\ensuremath{<}0.001}  % MCS p-value, 16:00 bar, 2014--2024, Qlike, Trange: Linear, HAR
\providecommand{\rhmcsAllQlikeTrangePRidgeFull}{}\renewcommand{\rhmcsAllQlikeTrangePRidgeFull}{0.008}  % MCS p-value, 16:00 bar, 2014--2024, Qlike, Trange: Ridge, Full
\providecommand{\rhmcsAllQlikeTrangePLassoFull}{}\renewcommand{\rhmcsAllQlikeTrangePLassoFull}{0.006}  % MCS p-value, 16:00 bar, 2014--2024, Qlike, Trange: Lasso, Full
\providecommand{\rhmcsAllQlikeTrangePEnetFull}{}\renewcommand{\rhmcsAllQlikeTrangePEnetFull}{0.010}  % MCS p-value, 16:00 bar, 2014--2024, Qlike, Trange: Elastic net, Full
\providecommand{\rhmcsAllQlikeTrangePLgbmHar}{}\renewcommand{\rhmcsAllQlikeTrangePLgbmHar}{0.14}  % MCS p-value, 16:00 bar, 2014--2024, Qlike, Trange: LightGBM, HAR
\providecommand{\rhmcsAllQlikeTrangePLgbmFull}{}\renewcommand{\rhmcsAllQlikeTrangePLgbmFull}{0.29}  % MCS p-value, 16:00 bar, 2014--2024, Qlike, Trange: LightGBM, Full
\providecommand{\rhmcsAllQlikeTrangePXgbHar}{}\renewcommand{\rhmcsAllQlikeTrangePXgbHar}{0.14}  % MCS p-value, 16:00 bar, 2014--2024, Qlike, Trange: XGBoost, HAR
\providecommand{\rhmcsAllQlikeTrangePXgbFull}{}\renewcommand{\rhmcsAllQlikeTrangePXgbFull}{1.00}  % MCS p-value, 16:00 bar, 2014--2024, Qlike, Trange: XGBoost, Full
\providecommand{\rhmcsAllQlikeTrangePRfHar}{}\renewcommand{\rhmcsAllQlikeTrangePRfHar}{0.001}  % MCS p-value, 16:00 bar, 2014--2024, Qlike, Trange: Random forest, HAR
\providecommand{\rhmcsAllQlikeTrangePRfFull}{}\renewcommand{\rhmcsAllQlikeTrangePRfFull}{\ensuremath{<}0.001}  % MCS p-value, 16:00 bar, 2014--2024, Qlike, Trange: Random forest, Full
\providecommand{\rhmcsAllQlikeTrangeNNinety}{}\renewcommand{\rhmcsAllQlikeTrangeNNinety}{4}  % models in the 90 % MCS (of 10)
\providecommand{\rhmcsAllQlikeTrangeSetNinety}{}\renewcommand{\rhmcsAllQlikeTrangeSetNinety}{XGBoost on the full predictor set, LightGBM on the full predictor set, XGBoost on the HAR benchmark, and LightGBM on the HAR benchmark}  % the 90 % MCS, best (last eliminated) first
\providecommand{\rhmcsAllQlikeTrangeStableNinety}{}\renewcommand{\rhmcsAllQlikeTrangeStableNinety}{yes}  % the 90 % MCS is the same at block lengths (5, 20)
\providecommand{\rhmcsAllQlikeTrangePMaxOutNinety}{}\renewcommand{\rhmcsAllQlikeTrangePMaxOutNinety}{0.010}  % largest MCS p-value outside the 90 % set
\providecommand{\rhmcsAllQlikeTrangeNNinetyLo}{}\renewcommand{\rhmcsAllQlikeTrangeNNinetyLo}{4}  % ... at block length 5
\providecommand{\rhmcsAllQlikeTrangeSetNinetyLo}{}\renewcommand{\rhmcsAllQlikeTrangeSetNinetyLo}{XGBoost on the full predictor set, LightGBM on the full predictor set, LightGBM on the HAR benchmark, and XGBoost on the HAR benchmark}  % ... at block length 5
\providecommand{\rhmcsAllQlikeTrangeNNinetyHi}{}\renewcommand{\rhmcsAllQlikeTrangeNNinetyHi}{4}  % ... at block length 20
\providecommand{\rhmcsAllQlikeTrangeSetNinetyHi}{}\renewcommand{\rhmcsAllQlikeTrangeSetNinetyHi}{XGBoost on the full predictor set, LightGBM on the full predictor set, LightGBM on the HAR benchmark, and XGBoost on the HAR benchmark}  % ... at block length 20
\providecommand{\rhmcsAllQlikeTrangeNSevenFive}{}\renewcommand{\rhmcsAllQlikeTrangeNSevenFive}{2}  % models in the 75 % MCS (of 10)
\providecommand{\rhmcsAllQlikeTrangeSetSevenFive}{}\renewcommand{\rhmcsAllQlikeTrangeSetSevenFive}{XGBoost on the full predictor set and LightGBM on the full predictor set}  % the 75 % MCS, best (last eliminated) first
\providecommand{\rhmcsAllQlikeTrangeStableSevenFive}{}\renewcommand{\rhmcsAllQlikeTrangeStableSevenFive}{yes}  % the 75 % MCS is the same at block lengths (5, 20)
\providecommand{\rhmcsAllQlikeTrangePMaxOutSevenFive}{}\renewcommand{\rhmcsAllQlikeTrangePMaxOutSevenFive}{0.14}  % largest MCS p-value outside the 75 % set
\providecommand{\rhmcsAllQlikeTrangeNSevenFiveLo}{}\renewcommand{\rhmcsAllQlikeTrangeNSevenFiveLo}{2}  % ... at block length 5
\providecommand{\rhmcsAllQlikeTrangeSetSevenFiveLo}{}\renewcommand{\rhmcsAllQlikeTrangeSetSevenFiveLo}{XGBoost on the full predictor set and LightGBM on the full predictor set}  % ... at block length 5
\providecommand{\rhmcsAllQlikeTrangeNSevenFiveHi}{}\renewcommand{\rhmcsAllQlikeTrangeNSevenFiveHi}{2}  % ... at block length 20
\providecommand{\rhmcsAllQlikeTrangeSetSevenFiveHi}{}\renewcommand{\rhmcsAllQlikeTrangeSetSevenFiveHi}{XGBoost on the full predictor set and LightGBM on the full predictor set}  % ... at block length 20
\providecommand{\rhmcsAllQlikeTrangeBest}{}\renewcommand{\rhmcsAllQlikeTrangeBest}{XGBoost on the full predictor set}  % the last model left (MCS p-value 1)
\providecommand{\rhmcsAllMseTmaxPLinHar}{}\renewcommand{\rhmcsAllMseTmaxPLinHar}{0.59}  % MCS p-value, 16:00 bar, 2014--2024, Mse, Tmax: Linear, HAR
\providecommand{\rhmcsAllMseTmaxPRidgeFull}{}\renewcommand{\rhmcsAllMseTmaxPRidgeFull}{0.43}  % MCS p-value, 16:00 bar, 2014--2024, Mse, Tmax: Ridge, Full
\providecommand{\rhmcsAllMseTmaxPLassoFull}{}\renewcommand{\rhmcsAllMseTmaxPLassoFull}{0.38}  % MCS p-value, 16:00 bar, 2014--2024, Mse, Tmax: Lasso, Full
\providecommand{\rhmcsAllMseTmaxPEnetFull}{}\renewcommand{\rhmcsAllMseTmaxPEnetFull}{0.48}  % MCS p-value, 16:00 bar, 2014--2024, Mse, Tmax: Elastic net, Full
\providecommand{\rhmcsAllMseTmaxPLgbmHar}{}\renewcommand{\rhmcsAllMseTmaxPLgbmHar}{0.59}  % MCS p-value, 16:00 bar, 2014--2024, Mse, Tmax: LightGBM, HAR
\providecommand{\rhmcsAllMseTmaxPLgbmFull}{}\renewcommand{\rhmcsAllMseTmaxPLgbmFull}{0.48}  % MCS p-value, 16:00 bar, 2014--2024, Mse, Tmax: LightGBM, Full
\providecommand{\rhmcsAllMseTmaxPXgbHar}{}\renewcommand{\rhmcsAllMseTmaxPXgbHar}{0.59}  % MCS p-value, 16:00 bar, 2014--2024, Mse, Tmax: XGBoost, HAR
\providecommand{\rhmcsAllMseTmaxPXgbFull}{}\renewcommand{\rhmcsAllMseTmaxPXgbFull}{0.48}  % MCS p-value, 16:00 bar, 2014--2024, Mse, Tmax: XGBoost, Full
\providecommand{\rhmcsAllMseTmaxPRfHar}{}\renewcommand{\rhmcsAllMseTmaxPRfHar}{1.00}  % MCS p-value, 16:00 bar, 2014--2024, Mse, Tmax: Random forest, HAR
\providecommand{\rhmcsAllMseTmaxPRfFull}{}\renewcommand{\rhmcsAllMseTmaxPRfFull}{0.59}  % MCS p-value, 16:00 bar, 2014--2024, Mse, Tmax: Random forest, Full
\providecommand{\rhmcsAllMseTmaxNNinety}{}\renewcommand{\rhmcsAllMseTmaxNNinety}{10}  % models in the 90 % MCS (of 10)
\providecommand{\rhmcsAllMseTmaxSetNinety}{}\renewcommand{\rhmcsAllMseTmaxSetNinety}{the random forest on the HAR benchmark, XGBoost on the HAR benchmark, LightGBM on the HAR benchmark, the linear models on the HAR benchmark, the random forest on the full predictor set, XGBoost on the full predictor set, LightGBM on the full predictor set, the elastic net on the full predictor set, ridge regression on the full predictor set, and the lasso on the full predictor set}  % the 90 % MCS, best (last eliminated) first
\providecommand{\rhmcsAllMseTmaxStableNinety}{}\renewcommand{\rhmcsAllMseTmaxStableNinety}{yes}  % the 90 % MCS is the same at block lengths (5, 20)
\providecommand{\rhmcsAllMseTmaxPMaxOutNinety}{}\renewcommand{\rhmcsAllMseTmaxPMaxOutNinety}{--}  % largest MCS p-value outside the 90 % set
\providecommand{\rhmcsAllMseTmaxNNinetyLo}{}\renewcommand{\rhmcsAllMseTmaxNNinetyLo}{10}  % ... at block length 5
\providecommand{\rhmcsAllMseTmaxSetNinetyLo}{}\renewcommand{\rhmcsAllMseTmaxSetNinetyLo}{the random forest on the HAR benchmark, XGBoost on the HAR benchmark, LightGBM on the HAR benchmark, the linear models on the HAR benchmark, the random forest on the full predictor set, XGBoost on the full predictor set, LightGBM on the full predictor set, the elastic net on the full predictor set, ridge regression on the full predictor set, and the lasso on the full predictor set}  % ... at block length 5
\providecommand{\rhmcsAllMseTmaxNNinetyHi}{}\renewcommand{\rhmcsAllMseTmaxNNinetyHi}{10}  % ... at block length 20
\providecommand{\rhmcsAllMseTmaxSetNinetyHi}{}\renewcommand{\rhmcsAllMseTmaxSetNinetyHi}{the random forest on the HAR benchmark, XGBoost on the HAR benchmark, LightGBM on the HAR benchmark, the linear models on the HAR benchmark, the random forest on the full predictor set, XGBoost on the full predictor set, LightGBM on the full predictor set, the elastic net on the full predictor set, ridge regression on the full predictor set, and the lasso on the full predictor set}  % ... at block length 20
\providecommand{\rhmcsAllMseTmaxNSevenFive}{}\renewcommand{\rhmcsAllMseTmaxNSevenFive}{10}  % models in the 75 % MCS (of 10)
\providecommand{\rhmcsAllMseTmaxSetSevenFive}{}\renewcommand{\rhmcsAllMseTmaxSetSevenFive}{the random forest on the HAR benchmark, XGBoost on the HAR benchmark, LightGBM on the HAR benchmark, the linear models on the HAR benchmark, the random forest on the full predictor set, XGBoost on the full predictor set, LightGBM on the full predictor set, the elastic net on the full predictor set, ridge regression on the full predictor set, and the lasso on the full predictor set}  % the 75 % MCS, best (last eliminated) first
\providecommand{\rhmcsAllMseTmaxStableSevenFive}{}\renewcommand{\rhmcsAllMseTmaxStableSevenFive}{no}  % the 75 % MCS is the same at block lengths (5, 20)
\providecommand{\rhmcsAllMseTmaxPMaxOutSevenFive}{}\renewcommand{\rhmcsAllMseTmaxPMaxOutSevenFive}{--}  % largest MCS p-value outside the 75 % set
\providecommand{\rhmcsAllMseTmaxNSevenFiveLo}{}\renewcommand{\rhmcsAllMseTmaxNSevenFiveLo}{9}  % ... at block length 5
\providecommand{\rhmcsAllMseTmaxSetSevenFiveLo}{}\renewcommand{\rhmcsAllMseTmaxSetSevenFiveLo}{the random forest on the HAR benchmark, XGBoost on the HAR benchmark, LightGBM on the HAR benchmark, the linear models on the HAR benchmark, the random forest on the full predictor set, XGBoost on the full predictor set, LightGBM on the full predictor set, the elastic net on the full predictor set, and ridge regression on the full predictor set}  % ... at block length 5
\providecommand{\rhmcsAllMseTmaxNSevenFiveHi}{}\renewcommand{\rhmcsAllMseTmaxNSevenFiveHi}{10}  % ... at block length 20
\providecommand{\rhmcsAllMseTmaxSetSevenFiveHi}{}\renewcommand{\rhmcsAllMseTmaxSetSevenFiveHi}{the random forest on the HAR benchmark, XGBoost on the HAR benchmark, LightGBM on the HAR benchmark, the linear models on the HAR benchmark, the random forest on the full predictor set, XGBoost on the full predictor set, LightGBM on the full predictor set, the elastic net on the full predictor set, ridge regression on the full predictor set, and the lasso on the full predictor set}  % ... at block length 20
\providecommand{\rhmcsAllMseTmaxBest}{}\renewcommand{\rhmcsAllMseTmaxBest}{the random forest on the HAR benchmark}  % the last model left (MCS p-value 1)
\providecommand{\rhmcsAllMseTrangePLinHar}{}\renewcommand{\rhmcsAllMseTrangePLinHar}{0.60}  % MCS p-value, 16:00 bar, 2014--2024, Mse, Trange: Linear, HAR
\providecommand{\rhmcsAllMseTrangePRidgeFull}{}\renewcommand{\rhmcsAllMseTrangePRidgeFull}{0.60}  % MCS p-value, 16:00 bar, 2014--2024, Mse, Trange: Ridge, Full
\providecommand{\rhmcsAllMseTrangePLassoFull}{}\renewcommand{\rhmcsAllMseTrangePLassoFull}{0.60}  % MCS p-value, 16:00 bar, 2014--2024, Mse, Trange: Lasso, Full
\providecommand{\rhmcsAllMseTrangePEnetFull}{}\renewcommand{\rhmcsAllMseTrangePEnetFull}{0.60}  % MCS p-value, 16:00 bar, 2014--2024, Mse, Trange: Elastic net, Full
\providecommand{\rhmcsAllMseTrangePLgbmHar}{}\renewcommand{\rhmcsAllMseTrangePLgbmHar}{0.35}  % MCS p-value, 16:00 bar, 2014--2024, Mse, Trange: LightGBM, HAR
\providecommand{\rhmcsAllMseTrangePLgbmFull}{}\renewcommand{\rhmcsAllMseTrangePLgbmFull}{0.60}  % MCS p-value, 16:00 bar, 2014--2024, Mse, Trange: LightGBM, Full
\providecommand{\rhmcsAllMseTrangePXgbHar}{}\renewcommand{\rhmcsAllMseTrangePXgbHar}{0.60}  % MCS p-value, 16:00 bar, 2014--2024, Mse, Trange: XGBoost, HAR
\providecommand{\rhmcsAllMseTrangePXgbFull}{}\renewcommand{\rhmcsAllMseTrangePXgbFull}{0.60}  % MCS p-value, 16:00 bar, 2014--2024, Mse, Trange: XGBoost, Full
\providecommand{\rhmcsAllMseTrangePRfHar}{}\renewcommand{\rhmcsAllMseTrangePRfHar}{1.00}  % MCS p-value, 16:00 bar, 2014--2024, Mse, Trange: Random forest, HAR
\providecommand{\rhmcsAllMseTrangePRfFull}{}\renewcommand{\rhmcsAllMseTrangePRfFull}{0.60}  % MCS p-value, 16:00 bar, 2014--2024, Mse, Trange: Random forest, Full
\providecommand{\rhmcsAllMseTrangeNNinety}{}\renewcommand{\rhmcsAllMseTrangeNNinety}{10}  % models in the 90 % MCS (of 10)
\providecommand{\rhmcsAllMseTrangeSetNinety}{}\renewcommand{\rhmcsAllMseTrangeSetNinety}{the random forest on the HAR benchmark, the linear models on the HAR benchmark, XGBoost on the HAR benchmark, the random forest on the full predictor set, XGBoost on the full predictor set, the elastic net on the full predictor set, ridge regression on the full predictor set, the lasso on the full predictor set, LightGBM on the full predictor set, and LightGBM on the HAR benchmark}  % the 90 % MCS, best (last eliminated) first
\providecommand{\rhmcsAllMseTrangeStableNinety}{}\renewcommand{\rhmcsAllMseTrangeStableNinety}{yes}  % the 90 % MCS is the same at block lengths (5, 20)
\providecommand{\rhmcsAllMseTrangePMaxOutNinety}{}\renewcommand{\rhmcsAllMseTrangePMaxOutNinety}{--}  % largest MCS p-value outside the 90 % set
\providecommand{\rhmcsAllMseTrangeNNinetyLo}{}\renewcommand{\rhmcsAllMseTrangeNNinetyLo}{10}  % ... at block length 5
\providecommand{\rhmcsAllMseTrangeSetNinetyLo}{}\renewcommand{\rhmcsAllMseTrangeSetNinetyLo}{the random forest on the HAR benchmark, the linear models on the HAR benchmark, XGBoost on the HAR benchmark, XGBoost on the full predictor set, the random forest on the full predictor set, the elastic net on the full predictor set, ridge regression on the full predictor set, LightGBM on the full predictor set, the lasso on the full predictor set, and LightGBM on the HAR benchmark}  % ... at block length 5
\providecommand{\rhmcsAllMseTrangeNNinetyHi}{}\renewcommand{\rhmcsAllMseTrangeNNinetyHi}{10}  % ... at block length 20
\providecommand{\rhmcsAllMseTrangeSetNinetyHi}{}\renewcommand{\rhmcsAllMseTrangeSetNinetyHi}{the random forest on the HAR benchmark, the linear models on the HAR benchmark, the random forest on the full predictor set, XGBoost on the full predictor set, the elastic net on the full predictor set, ridge regression on the full predictor set, the lasso on the full predictor set, XGBoost on the HAR benchmark, LightGBM on the full predictor set, and LightGBM on the HAR benchmark}  % ... at block length 20
\providecommand{\rhmcsAllMseTrangeNSevenFive}{}\renewcommand{\rhmcsAllMseTrangeNSevenFive}{10}  % models in the 75 % MCS (of 10)
\providecommand{\rhmcsAllMseTrangeSetSevenFive}{}\renewcommand{\rhmcsAllMseTrangeSetSevenFive}{the random forest on the HAR benchmark, the linear models on the HAR benchmark, XGBoost on the HAR benchmark, the random forest on the full predictor set, XGBoost on the full predictor set, the elastic net on the full predictor set, ridge regression on the full predictor set, the lasso on the full predictor set, LightGBM on the full predictor set, and LightGBM on the HAR benchmark}  % the 75 % MCS, best (last eliminated) first
\providecommand{\rhmcsAllMseTrangeStableSevenFive}{}\renewcommand{\rhmcsAllMseTrangeStableSevenFive}{yes}  % the 75 % MCS is the same at block lengths (5, 20)
\providecommand{\rhmcsAllMseTrangePMaxOutSevenFive}{}\renewcommand{\rhmcsAllMseTrangePMaxOutSevenFive}{--}  % largest MCS p-value outside the 75 % set
\providecommand{\rhmcsAllMseTrangeNSevenFiveLo}{}\renewcommand{\rhmcsAllMseTrangeNSevenFiveLo}{10}  % ... at block length 5
\providecommand{\rhmcsAllMseTrangeSetSevenFiveLo}{}\renewcommand{\rhmcsAllMseTrangeSetSevenFiveLo}{the random forest on the HAR benchmark, the linear models on the HAR benchmark, XGBoost on the HAR benchmark, XGBoost on the full predictor set, the random forest on the full predictor set, the elastic net on the full predictor set, ridge regression on the full predictor set, LightGBM on the full predictor set, the lasso on the full predictor set, and LightGBM on the HAR benchmark}  % ... at block length 5
\providecommand{\rhmcsAllMseTrangeNSevenFiveHi}{}\renewcommand{\rhmcsAllMseTrangeNSevenFiveHi}{10}  % ... at block length 20
\providecommand{\rhmcsAllMseTrangeSetSevenFiveHi}{}\renewcommand{\rhmcsAllMseTrangeSetSevenFiveHi}{the random forest on the HAR benchmark, the linear models on the HAR benchmark, the random forest on the full predictor set, XGBoost on the full predictor set, the elastic net on the full predictor set, ridge regression on the full predictor set, the lasso on the full predictor set, XGBoost on the HAR benchmark, LightGBM on the full predictor set, and LightGBM on the HAR benchmark}  % ... at block length 20
\providecommand{\rhmcsAllMseTrangeBest}{}\renewcommand{\rhmcsAllMseTrangeBest}{the random forest on the HAR benchmark}  % the last model left (MCS p-value 1)
\providecommand{\rhmcsOosQlikeTmaxPLinHar}{}\renewcommand{\rhmcsOosQlikeTmaxPLinHar}{0.043}  % MCS p-value, 16:00 bar, T_oos, Qlike, Tmax: Linear, HAR
\providecommand{\rhmcsOosQlikeTmaxPRidgeFull}{}\renewcommand{\rhmcsOosQlikeTmaxPRidgeFull}{0.008}  % MCS p-value, 16:00 bar, T_oos, Qlike, Tmax: Ridge, Full
\providecommand{\rhmcsOosQlikeTmaxPLassoFull}{}\renewcommand{\rhmcsOosQlikeTmaxPLassoFull}{0.008}  % MCS p-value, 16:00 bar, T_oos, Qlike, Tmax: Lasso, Full
\providecommand{\rhmcsOosQlikeTmaxPEnetFull}{}\renewcommand{\rhmcsOosQlikeTmaxPEnetFull}{0.024}  % MCS p-value, 16:00 bar, T_oos, Qlike, Tmax: Elastic net, Full
\providecommand{\rhmcsOosQlikeTmaxPLgbmHar}{}\renewcommand{\rhmcsOosQlikeTmaxPLgbmHar}{0.16}  % MCS p-value, 16:00 bar, T_oos, Qlike, Tmax: LightGBM, HAR
\providecommand{\rhmcsOosQlikeTmaxPLgbmFull}{}\renewcommand{\rhmcsOosQlikeTmaxPLgbmFull}{0.18}  % MCS p-value, 16:00 bar, T_oos, Qlike, Tmax: LightGBM, Full
\providecommand{\rhmcsOosQlikeTmaxPXgbHar}{}\renewcommand{\rhmcsOosQlikeTmaxPXgbHar}{0.16}  % MCS p-value, 16:00 bar, T_oos, Qlike, Tmax: XGBoost, HAR
\providecommand{\rhmcsOosQlikeTmaxPXgbFull}{}\renewcommand{\rhmcsOosQlikeTmaxPXgbFull}{0.18}  % MCS p-value, 16:00 bar, T_oos, Qlike, Tmax: XGBoost, Full
\providecommand{\rhmcsOosQlikeTmaxPRfHar}{}\renewcommand{\rhmcsOosQlikeTmaxPRfHar}{0.043}  % MCS p-value, 16:00 bar, T_oos, Qlike, Tmax: Random forest, HAR
\providecommand{\rhmcsOosQlikeTmaxPRfFull}{}\renewcommand{\rhmcsOosQlikeTmaxPRfFull}{0.14}  % MCS p-value, 16:00 bar, T_oos, Qlike, Tmax: Random forest, Full
\providecommand{\rhmcsOosQlikeTmaxPBench}{}\renewcommand{\rhmcsOosQlikeTmaxPBench}{1.00}  % MCS p-value, 16:00 bar, T_oos, Qlike, Tmax: Closing-interval HAR benchmark
\providecommand{\rhmcsOosQlikeTmaxNNinety}{}\renewcommand{\rhmcsOosQlikeTmaxNNinety}{6}  % models in the 90 % MCS (of 11)
\providecommand{\rhmcsOosQlikeTmaxSetNinety}{}\renewcommand{\rhmcsOosQlikeTmaxSetNinety}{the closing-interval HAR benchmark, XGBoost on the full predictor set, LightGBM on the full predictor set, XGBoost on the HAR benchmark, LightGBM on the HAR benchmark, and the random forest on the full predictor set}  % the 90 % MCS, best (last eliminated) first
\providecommand{\rhmcsOosQlikeTmaxStableNinety}{}\renewcommand{\rhmcsOosQlikeTmaxStableNinety}{yes}  % the 90 % MCS is the same at block lengths (5, 20)
\providecommand{\rhmcsOosQlikeTmaxPMaxOutNinety}{}\renewcommand{\rhmcsOosQlikeTmaxPMaxOutNinety}{0.043}  % largest MCS p-value outside the 90 % set
\providecommand{\rhmcsOosQlikeTmaxNNinetyLo}{}\renewcommand{\rhmcsOosQlikeTmaxNNinetyLo}{6}  % ... at block length 5
\providecommand{\rhmcsOosQlikeTmaxSetNinetyLo}{}\renewcommand{\rhmcsOosQlikeTmaxSetNinetyLo}{the closing-interval HAR benchmark, XGBoost on the full predictor set, LightGBM on the full predictor set, XGBoost on the HAR benchmark, LightGBM on the HAR benchmark, and the random forest on the full predictor set}  % ... at block length 5
\providecommand{\rhmcsOosQlikeTmaxNNinetyHi}{}\renewcommand{\rhmcsOosQlikeTmaxNNinetyHi}{6}  % ... at block length 20
\providecommand{\rhmcsOosQlikeTmaxSetNinetyHi}{}\renewcommand{\rhmcsOosQlikeTmaxSetNinetyHi}{the closing-interval HAR benchmark, XGBoost on the full predictor set, LightGBM on the full predictor set, LightGBM on the HAR benchmark, XGBoost on the HAR benchmark, and the random forest on the full predictor set}  % ... at block length 20
\providecommand{\rhmcsOosQlikeTmaxNSevenFive}{}\renewcommand{\rhmcsOosQlikeTmaxNSevenFive}{1}  % models in the 75 % MCS (of 11)
\providecommand{\rhmcsOosQlikeTmaxSetSevenFive}{}\renewcommand{\rhmcsOosQlikeTmaxSetSevenFive}{the closing-interval HAR benchmark}  % the 75 % MCS, best (last eliminated) first
\providecommand{\rhmcsOosQlikeTmaxStableSevenFive}{}\renewcommand{\rhmcsOosQlikeTmaxStableSevenFive}{yes}  % the 75 % MCS is the same at block lengths (5, 20)
\providecommand{\rhmcsOosQlikeTmaxPMaxOutSevenFive}{}\renewcommand{\rhmcsOosQlikeTmaxPMaxOutSevenFive}{0.18}  % largest MCS p-value outside the 75 % set
\providecommand{\rhmcsOosQlikeTmaxNSevenFiveLo}{}\renewcommand{\rhmcsOosQlikeTmaxNSevenFiveLo}{1}  % ... at block length 5
\providecommand{\rhmcsOosQlikeTmaxSetSevenFiveLo}{}\renewcommand{\rhmcsOosQlikeTmaxSetSevenFiveLo}{the closing-interval HAR benchmark}  % ... at block length 5
\providecommand{\rhmcsOosQlikeTmaxNSevenFiveHi}{}\renewcommand{\rhmcsOosQlikeTmaxNSevenFiveHi}{1}  % ... at block length 20
\providecommand{\rhmcsOosQlikeTmaxSetSevenFiveHi}{}\renewcommand{\rhmcsOosQlikeTmaxSetSevenFiveHi}{the closing-interval HAR benchmark}  % ... at block length 20
\providecommand{\rhmcsOosQlikeTmaxBest}{}\renewcommand{\rhmcsOosQlikeTmaxBest}{the closing-interval HAR benchmark}  % the last model left (MCS p-value 1)
\providecommand{\rhmcsOosQlikeTrangePLinHar}{}\renewcommand{\rhmcsOosQlikeTrangePLinHar}{0.034}  % MCS p-value, 16:00 bar, T_oos, Qlike, Trange: Linear, HAR
\providecommand{\rhmcsOosQlikeTrangePRidgeFull}{}\renewcommand{\rhmcsOosQlikeTrangePRidgeFull}{0.007}  % MCS p-value, 16:00 bar, T_oos, Qlike, Trange: Ridge, Full
\providecommand{\rhmcsOosQlikeTrangePLassoFull}{}\renewcommand{\rhmcsOosQlikeTrangePLassoFull}{0.008}  % MCS p-value, 16:00 bar, T_oos, Qlike, Trange: Lasso, Full
\providecommand{\rhmcsOosQlikeTrangePEnetFull}{}\renewcommand{\rhmcsOosQlikeTrangePEnetFull}{0.013}  % MCS p-value, 16:00 bar, T_oos, Qlike, Trange: Elastic net, Full
\providecommand{\rhmcsOosQlikeTrangePLgbmHar}{}\renewcommand{\rhmcsOosQlikeTrangePLgbmHar}{0.034}  % MCS p-value, 16:00 bar, T_oos, Qlike, Trange: LightGBM, HAR
\providecommand{\rhmcsOosQlikeTrangePLgbmFull}{}\renewcommand{\rhmcsOosQlikeTrangePLgbmFull}{0.29}  % MCS p-value, 16:00 bar, T_oos, Qlike, Trange: LightGBM, Full
\providecommand{\rhmcsOosQlikeTrangePXgbHar}{}\renewcommand{\rhmcsOosQlikeTrangePXgbHar}{0.034}  % MCS p-value, 16:00 bar, T_oos, Qlike, Trange: XGBoost, HAR
\providecommand{\rhmcsOosQlikeTrangePXgbFull}{}\renewcommand{\rhmcsOosQlikeTrangePXgbFull}{0.29}  % MCS p-value, 16:00 bar, T_oos, Qlike, Trange: XGBoost, Full
\providecommand{\rhmcsOosQlikeTrangePRfHar}{}\renewcommand{\rhmcsOosQlikeTrangePRfHar}{0.025}  % MCS p-value, 16:00 bar, T_oos, Qlike, Trange: Random forest, HAR
\providecommand{\rhmcsOosQlikeTrangePRfFull}{}\renewcommand{\rhmcsOosQlikeTrangePRfFull}{0.025}  % MCS p-value, 16:00 bar, T_oos, Qlike, Trange: Random forest, Full
\providecommand{\rhmcsOosQlikeTrangePBench}{}\renewcommand{\rhmcsOosQlikeTrangePBench}{1.00}  % MCS p-value, 16:00 bar, T_oos, Qlike, Trange: Closing-interval HAR benchmark
\providecommand{\rhmcsOosQlikeTrangeNNinety}{}\renewcommand{\rhmcsOosQlikeTrangeNNinety}{3}  % models in the 90 % MCS (of 11)
\providecommand{\rhmcsOosQlikeTrangeSetNinety}{}\renewcommand{\rhmcsOosQlikeTrangeSetNinety}{the closing-interval HAR benchmark, XGBoost on the full predictor set, and LightGBM on the full predictor set}  % the 90 % MCS, best (last eliminated) first
\providecommand{\rhmcsOosQlikeTrangeStableNinety}{}\renewcommand{\rhmcsOosQlikeTrangeStableNinety}{yes}  % the 90 % MCS is the same at block lengths (5, 20)
\providecommand{\rhmcsOosQlikeTrangePMaxOutNinety}{}\renewcommand{\rhmcsOosQlikeTrangePMaxOutNinety}{0.034}  % largest MCS p-value outside the 90 % set
\providecommand{\rhmcsOosQlikeTrangeNNinetyLo}{}\renewcommand{\rhmcsOosQlikeTrangeNNinetyLo}{3}  % ... at block length 5
\providecommand{\rhmcsOosQlikeTrangeSetNinetyLo}{}\renewcommand{\rhmcsOosQlikeTrangeSetNinetyLo}{the closing-interval HAR benchmark, XGBoost on the full predictor set, and LightGBM on the full predictor set}  % ... at block length 5
\providecommand{\rhmcsOosQlikeTrangeNNinetyHi}{}\renewcommand{\rhmcsOosQlikeTrangeNNinetyHi}{3}  % ... at block length 20
\providecommand{\rhmcsOosQlikeTrangeSetNinetyHi}{}\renewcommand{\rhmcsOosQlikeTrangeSetNinetyHi}{the closing-interval HAR benchmark, XGBoost on the full predictor set, and LightGBM on the full predictor set}  % ... at block length 20
\providecommand{\rhmcsOosQlikeTrangeNSevenFive}{}\renewcommand{\rhmcsOosQlikeTrangeNSevenFive}{3}  % models in the 75 % MCS (of 11)
\providecommand{\rhmcsOosQlikeTrangeSetSevenFive}{}\renewcommand{\rhmcsOosQlikeTrangeSetSevenFive}{the closing-interval HAR benchmark, XGBoost on the full predictor set, and LightGBM on the full predictor set}  % the 75 % MCS, best (last eliminated) first
\providecommand{\rhmcsOosQlikeTrangeStableSevenFive}{}\renewcommand{\rhmcsOosQlikeTrangeStableSevenFive}{yes}  % the 75 % MCS is the same at block lengths (5, 20)
\providecommand{\rhmcsOosQlikeTrangePMaxOutSevenFive}{}\renewcommand{\rhmcsOosQlikeTrangePMaxOutSevenFive}{0.034}  % largest MCS p-value outside the 75 % set
\providecommand{\rhmcsOosQlikeTrangeNSevenFiveLo}{}\renewcommand{\rhmcsOosQlikeTrangeNSevenFiveLo}{3}  % ... at block length 5
\providecommand{\rhmcsOosQlikeTrangeSetSevenFiveLo}{}\renewcommand{\rhmcsOosQlikeTrangeSetSevenFiveLo}{the closing-interval HAR benchmark, XGBoost on the full predictor set, and LightGBM on the full predictor set}  % ... at block length 5
\providecommand{\rhmcsOosQlikeTrangeNSevenFiveHi}{}\renewcommand{\rhmcsOosQlikeTrangeNSevenFiveHi}{3}  % ... at block length 20
\providecommand{\rhmcsOosQlikeTrangeSetSevenFiveHi}{}\renewcommand{\rhmcsOosQlikeTrangeSetSevenFiveHi}{the closing-interval HAR benchmark, XGBoost on the full predictor set, and LightGBM on the full predictor set}  % ... at block length 20
\providecommand{\rhmcsOosQlikeTrangeBest}{}\renewcommand{\rhmcsOosQlikeTrangeBest}{the closing-interval HAR benchmark}  % the last model left (MCS p-value 1)
\providecommand{\rhmcsOosMseTmaxPLinHar}{}\renewcommand{\rhmcsOosMseTmaxPLinHar}{0.56}  % MCS p-value, 16:00 bar, T_oos, Mse, Tmax: Linear, HAR
\providecommand{\rhmcsOosMseTmaxPRidgeFull}{}\renewcommand{\rhmcsOosMseTmaxPRidgeFull}{0.44}  % MCS p-value, 16:00 bar, T_oos, Mse, Tmax: Ridge, Full
\providecommand{\rhmcsOosMseTmaxPLassoFull}{}\renewcommand{\rhmcsOosMseTmaxPLassoFull}{0.39}  % MCS p-value, 16:00 bar, T_oos, Mse, Tmax: Lasso, Full
\providecommand{\rhmcsOosMseTmaxPEnetFull}{}\renewcommand{\rhmcsOosMseTmaxPEnetFull}{0.49}  % MCS p-value, 16:00 bar, T_oos, Mse, Tmax: Elastic net, Full
\providecommand{\rhmcsOosMseTmaxPLgbmHar}{}\renewcommand{\rhmcsOosMseTmaxPLgbmHar}{0.56}  % MCS p-value, 16:00 bar, T_oos, Mse, Tmax: LightGBM, HAR
\providecommand{\rhmcsOosMseTmaxPLgbmFull}{}\renewcommand{\rhmcsOosMseTmaxPLgbmFull}{0.54}  % MCS p-value, 16:00 bar, T_oos, Mse, Tmax: LightGBM, Full
\providecommand{\rhmcsOosMseTmaxPXgbHar}{}\renewcommand{\rhmcsOosMseTmaxPXgbHar}{0.56}  % MCS p-value, 16:00 bar, T_oos, Mse, Tmax: XGBoost, HAR
\providecommand{\rhmcsOosMseTmaxPXgbFull}{}\renewcommand{\rhmcsOosMseTmaxPXgbFull}{0.56}  % MCS p-value, 16:00 bar, T_oos, Mse, Tmax: XGBoost, Full
\providecommand{\rhmcsOosMseTmaxPRfHar}{}\renewcommand{\rhmcsOosMseTmaxPRfHar}{1.00}  % MCS p-value, 16:00 bar, T_oos, Mse, Tmax: Random forest, HAR
\providecommand{\rhmcsOosMseTmaxPRfFull}{}\renewcommand{\rhmcsOosMseTmaxPRfFull}{0.56}  % MCS p-value, 16:00 bar, T_oos, Mse, Tmax: Random forest, Full
\providecommand{\rhmcsOosMseTmaxPBench}{}\renewcommand{\rhmcsOosMseTmaxPBench}{0.56}  % MCS p-value, 16:00 bar, T_oos, Mse, Tmax: Closing-interval HAR benchmark
\providecommand{\rhmcsOosMseTmaxNNinety}{}\renewcommand{\rhmcsOosMseTmaxNNinety}{11}  % models in the 90 % MCS (of 11)
\providecommand{\rhmcsOosMseTmaxSetNinety}{}\renewcommand{\rhmcsOosMseTmaxSetNinety}{the random forest on the HAR benchmark, XGBoost on the HAR benchmark, LightGBM on the HAR benchmark, the random forest on the full predictor set, the linear models on the HAR benchmark, the closing-interval HAR benchmark, XGBoost on the full predictor set, LightGBM on the full predictor set, the elastic net on the full predictor set, ridge regression on the full predictor set, and the lasso on the full predictor set}  % the 90 % MCS, best (last eliminated) first
\providecommand{\rhmcsOosMseTmaxStableNinety}{}\renewcommand{\rhmcsOosMseTmaxStableNinety}{yes}  % the 90 % MCS is the same at block lengths (5, 20)
\providecommand{\rhmcsOosMseTmaxPMaxOutNinety}{}\renewcommand{\rhmcsOosMseTmaxPMaxOutNinety}{--}  % largest MCS p-value outside the 90 % set
\providecommand{\rhmcsOosMseTmaxNNinetyLo}{}\renewcommand{\rhmcsOosMseTmaxNNinetyLo}{11}  % ... at block length 5
\providecommand{\rhmcsOosMseTmaxSetNinetyLo}{}\renewcommand{\rhmcsOosMseTmaxSetNinetyLo}{the random forest on the HAR benchmark, XGBoost on the HAR benchmark, LightGBM on the HAR benchmark, the linear models on the HAR benchmark, the random forest on the full predictor set, the closing-interval HAR benchmark, XGBoost on the full predictor set, LightGBM on the full predictor set, the elastic net on the full predictor set, ridge regression on the full predictor set, and the lasso on the full predictor set}  % ... at block length 5
\providecommand{\rhmcsOosMseTmaxNNinetyHi}{}\renewcommand{\rhmcsOosMseTmaxNNinetyHi}{11}  % ... at block length 20
\providecommand{\rhmcsOosMseTmaxSetNinetyHi}{}\renewcommand{\rhmcsOosMseTmaxSetNinetyHi}{the random forest on the HAR benchmark, XGBoost on the HAR benchmark, LightGBM on the HAR benchmark, the linear models on the HAR benchmark, the random forest on the full predictor set, the closing-interval HAR benchmark, XGBoost on the full predictor set, LightGBM on the full predictor set, the elastic net on the full predictor set, ridge regression on the full predictor set, and the lasso on the full predictor set}  % ... at block length 20
\providecommand{\rhmcsOosMseTmaxNSevenFive}{}\renewcommand{\rhmcsOosMseTmaxNSevenFive}{11}  % models in the 75 % MCS (of 11)
\providecommand{\rhmcsOosMseTmaxSetSevenFive}{}\renewcommand{\rhmcsOosMseTmaxSetSevenFive}{the random forest on the HAR benchmark, XGBoost on the HAR benchmark, LightGBM on the HAR benchmark, the random forest on the full predictor set, the linear models on the HAR benchmark, the closing-interval HAR benchmark, XGBoost on the full predictor set, LightGBM on the full predictor set, the elastic net on the full predictor set, ridge regression on the full predictor set, and the lasso on the full predictor set}  % the 75 % MCS, best (last eliminated) first
\providecommand{\rhmcsOosMseTmaxStableSevenFive}{}\renewcommand{\rhmcsOosMseTmaxStableSevenFive}{no}  % the 75 % MCS is the same at block lengths (5, 20)
\providecommand{\rhmcsOosMseTmaxPMaxOutSevenFive}{}\renewcommand{\rhmcsOosMseTmaxPMaxOutSevenFive}{--}  % largest MCS p-value outside the 75 % set
\providecommand{\rhmcsOosMseTmaxNSevenFiveLo}{}\renewcommand{\rhmcsOosMseTmaxNSevenFiveLo}{10}  % ... at block length 5
\providecommand{\rhmcsOosMseTmaxSetSevenFiveLo}{}\renewcommand{\rhmcsOosMseTmaxSetSevenFiveLo}{the random forest on the HAR benchmark, XGBoost on the HAR benchmark, LightGBM on the HAR benchmark, the linear models on the HAR benchmark, the random forest on the full predictor set, the closing-interval HAR benchmark, XGBoost on the full predictor set, LightGBM on the full predictor set, the elastic net on the full predictor set, and ridge regression on the full predictor set}  % ... at block length 5
\providecommand{\rhmcsOosMseTmaxNSevenFiveHi}{}\renewcommand{\rhmcsOosMseTmaxNSevenFiveHi}{11}  % ... at block length 20
\providecommand{\rhmcsOosMseTmaxSetSevenFiveHi}{}\renewcommand{\rhmcsOosMseTmaxSetSevenFiveHi}{the random forest on the HAR benchmark, XGBoost on the HAR benchmark, LightGBM on the HAR benchmark, the linear models on the HAR benchmark, the random forest on the full predictor set, the closing-interval HAR benchmark, XGBoost on the full predictor set, LightGBM on the full predictor set, the elastic net on the full predictor set, ridge regression on the full predictor set, and the lasso on the full predictor set}  % ... at block length 20
\providecommand{\rhmcsOosMseTmaxBest}{}\renewcommand{\rhmcsOosMseTmaxBest}{the random forest on the HAR benchmark}  % the last model left (MCS p-value 1)
\providecommand{\rhmcsOosMseTrangePLinHar}{}\renewcommand{\rhmcsOosMseTrangePLinHar}{0.56}  % MCS p-value, 16:00 bar, T_oos, Mse, Trange: Linear, HAR
\providecommand{\rhmcsOosMseTrangePRidgeFull}{}\renewcommand{\rhmcsOosMseTrangePRidgeFull}{0.56}  % MCS p-value, 16:00 bar, T_oos, Mse, Trange: Ridge, Full
\providecommand{\rhmcsOosMseTrangePLassoFull}{}\renewcommand{\rhmcsOosMseTrangePLassoFull}{0.56}  % MCS p-value, 16:00 bar, T_oos, Mse, Trange: Lasso, Full
\providecommand{\rhmcsOosMseTrangePEnetFull}{}\renewcommand{\rhmcsOosMseTrangePEnetFull}{0.56}  % MCS p-value, 16:00 bar, T_oos, Mse, Trange: Elastic net, Full
\providecommand{\rhmcsOosMseTrangePLgbmHar}{}\renewcommand{\rhmcsOosMseTrangePLgbmHar}{0.56}  % MCS p-value, 16:00 bar, T_oos, Mse, Trange: LightGBM, HAR
\providecommand{\rhmcsOosMseTrangePLgbmFull}{}\renewcommand{\rhmcsOosMseTrangePLgbmFull}{0.56}  % MCS p-value, 16:00 bar, T_oos, Mse, Trange: LightGBM, Full
\providecommand{\rhmcsOosMseTrangePXgbHar}{}\renewcommand{\rhmcsOosMseTrangePXgbHar}{0.56}  % MCS p-value, 16:00 bar, T_oos, Mse, Trange: XGBoost, HAR
\providecommand{\rhmcsOosMseTrangePXgbFull}{}\renewcommand{\rhmcsOosMseTrangePXgbFull}{0.56}  % MCS p-value, 16:00 bar, T_oos, Mse, Trange: XGBoost, Full
\providecommand{\rhmcsOosMseTrangePRfHar}{}\renewcommand{\rhmcsOosMseTrangePRfHar}{1.00}  % MCS p-value, 16:00 bar, T_oos, Mse, Trange: Random forest, HAR
\providecommand{\rhmcsOosMseTrangePRfFull}{}\renewcommand{\rhmcsOosMseTrangePRfFull}{0.56}  % MCS p-value, 16:00 bar, T_oos, Mse, Trange: Random forest, Full
\providecommand{\rhmcsOosMseTrangePBench}{}\renewcommand{\rhmcsOosMseTrangePBench}{0.56}  % MCS p-value, 16:00 bar, T_oos, Mse, Trange: Closing-interval HAR benchmark
\providecommand{\rhmcsOosMseTrangeNNinety}{}\renewcommand{\rhmcsOosMseTrangeNNinety}{11}  % models in the 90 % MCS (of 11)
\providecommand{\rhmcsOosMseTrangeSetNinety}{}\renewcommand{\rhmcsOosMseTrangeSetNinety}{the random forest on the HAR benchmark, XGBoost on the HAR benchmark, the linear models on the HAR benchmark, the random forest on the full predictor set, XGBoost on the full predictor set, LightGBM on the full predictor set, the elastic net on the full predictor set, ridge regression on the full predictor set, the lasso on the full predictor set, the closing-interval HAR benchmark, and LightGBM on the HAR benchmark}  % the 90 % MCS, best (last eliminated) first
\providecommand{\rhmcsOosMseTrangeStableNinety}{}\renewcommand{\rhmcsOosMseTrangeStableNinety}{yes}  % the 90 % MCS is the same at block lengths (5, 20)
\providecommand{\rhmcsOosMseTrangePMaxOutNinety}{}\renewcommand{\rhmcsOosMseTrangePMaxOutNinety}{--}  % largest MCS p-value outside the 90 % set
\providecommand{\rhmcsOosMseTrangeNNinetyLo}{}\renewcommand{\rhmcsOosMseTrangeNNinetyLo}{11}  % ... at block length 5
\providecommand{\rhmcsOosMseTrangeSetNinetyLo}{}\renewcommand{\rhmcsOosMseTrangeSetNinetyLo}{the random forest on the HAR benchmark, XGBoost on the HAR benchmark, the linear models on the HAR benchmark, XGBoost on the full predictor set, the random forest on the full predictor set, the elastic net on the full predictor set, LightGBM on the full predictor set, LightGBM on the HAR benchmark, ridge regression on the full predictor set, the lasso on the full predictor set, and the closing-interval HAR benchmark}  % ... at block length 5
\providecommand{\rhmcsOosMseTrangeNNinetyHi}{}\renewcommand{\rhmcsOosMseTrangeNNinetyHi}{11}  % ... at block length 20
\providecommand{\rhmcsOosMseTrangeSetNinetyHi}{}\renewcommand{\rhmcsOosMseTrangeSetNinetyHi}{the random forest on the HAR benchmark, the linear models on the HAR benchmark, the random forest on the full predictor set, XGBoost on the HAR benchmark, XGBoost on the full predictor set, LightGBM on the full predictor set, the elastic net on the full predictor set, ridge regression on the full predictor set, the lasso on the full predictor set, the closing-interval HAR benchmark, and LightGBM on the HAR benchmark}  % ... at block length 20
\providecommand{\rhmcsOosMseTrangeNSevenFive}{}\renewcommand{\rhmcsOosMseTrangeNSevenFive}{11}  % models in the 75 % MCS (of 11)
\providecommand{\rhmcsOosMseTrangeSetSevenFive}{}\renewcommand{\rhmcsOosMseTrangeSetSevenFive}{the random forest on the HAR benchmark, XGBoost on the HAR benchmark, the linear models on the HAR benchmark, the random forest on the full predictor set, XGBoost on the full predictor set, LightGBM on the full predictor set, the elastic net on the full predictor set, ridge regression on the full predictor set, the lasso on the full predictor set, the closing-interval HAR benchmark, and LightGBM on the HAR benchmark}  % the 75 % MCS, best (last eliminated) first
\providecommand{\rhmcsOosMseTrangeStableSevenFive}{}\renewcommand{\rhmcsOosMseTrangeStableSevenFive}{yes}  % the 75 % MCS is the same at block lengths (5, 20)
\providecommand{\rhmcsOosMseTrangePMaxOutSevenFive}{}\renewcommand{\rhmcsOosMseTrangePMaxOutSevenFive}{--}  % largest MCS p-value outside the 75 % set
\providecommand{\rhmcsOosMseTrangeNSevenFiveLo}{}\renewcommand{\rhmcsOosMseTrangeNSevenFiveLo}{11}  % ... at block length 5
\providecommand{\rhmcsOosMseTrangeSetSevenFiveLo}{}\renewcommand{\rhmcsOosMseTrangeSetSevenFiveLo}{the random forest on the HAR benchmark, XGBoost on the HAR benchmark, the linear models on the HAR benchmark, XGBoost on the full predictor set, the random forest on the full predictor set, the elastic net on the full predictor set, LightGBM on the full predictor set, LightGBM on the HAR benchmark, ridge regression on the full predictor set, the lasso on the full predictor set, and the closing-interval HAR benchmark}  % ... at block length 5
\providecommand{\rhmcsOosMseTrangeNSevenFiveHi}{}\renewcommand{\rhmcsOosMseTrangeNSevenFiveHi}{11}  % ... at block length 20
\providecommand{\rhmcsOosMseTrangeSetSevenFiveHi}{}\renewcommand{\rhmcsOosMseTrangeSetSevenFiveHi}{the random forest on the HAR benchmark, the linear models on the HAR benchmark, the random forest on the full predictor set, XGBoost on the HAR benchmark, XGBoost on the full predictor set, LightGBM on the full predictor set, the elastic net on the full predictor set, ridge regression on the full predictor set, the lasso on the full predictor set, the closing-interval HAR benchmark, and LightGBM on the HAR benchmark}  % ... at block length 20
\providecommand{\rhmcsOosMseTrangeBest}{}\renewcommand{\rhmcsOosMseTrangeBest}{the random forest on the HAR benchmark}  % the last model left (MCS p-value 1)
\providecommand{\rhmcsChQlikeTmaxPLinHar}{}\renewcommand{\rhmcsChQlikeTmaxPLinHar}{0.050}  % MCS p-value, 16:00 bar, the benchmark's span, Qlike, Tmax: Linear, HAR
\providecommand{\rhmcsChQlikeTmaxPRidgeFull}{}\renewcommand{\rhmcsChQlikeTmaxPRidgeFull}{0.050}  % MCS p-value, 16:00 bar, the benchmark's span, Qlike, Tmax: Ridge, Full
\providecommand{\rhmcsChQlikeTmaxPLassoFull}{}\renewcommand{\rhmcsChQlikeTmaxPLassoFull}{0.050}  % MCS p-value, 16:00 bar, the benchmark's span, Qlike, Tmax: Lasso, Full
\providecommand{\rhmcsChQlikeTmaxPEnetFull}{}\renewcommand{\rhmcsChQlikeTmaxPEnetFull}{0.102}  % MCS p-value, 16:00 bar, the benchmark's span, Qlike, Tmax: Elastic net, Full
\providecommand{\rhmcsChQlikeTmaxPLgbmHar}{}\renewcommand{\rhmcsChQlikeTmaxPLgbmHar}{0.102}  % MCS p-value, 16:00 bar, the benchmark's span, Qlike, Tmax: LightGBM, HAR
\providecommand{\rhmcsChQlikeTmaxPLgbmFull}{}\renewcommand{\rhmcsChQlikeTmaxPLgbmFull}{0.16}  % MCS p-value, 16:00 bar, the benchmark's span, Qlike, Tmax: LightGBM, Full
\providecommand{\rhmcsChQlikeTmaxPXgbHar}{}\renewcommand{\rhmcsChQlikeTmaxPXgbHar}{0.102}  % MCS p-value, 16:00 bar, the benchmark's span, Qlike, Tmax: XGBoost, HAR
\providecommand{\rhmcsChQlikeTmaxPXgbFull}{}\renewcommand{\rhmcsChQlikeTmaxPXgbFull}{0.22}  % MCS p-value, 16:00 bar, the benchmark's span, Qlike, Tmax: XGBoost, Full
\providecommand{\rhmcsChQlikeTmaxPRfHar}{}\renewcommand{\rhmcsChQlikeTmaxPRfHar}{0.050}  % MCS p-value, 16:00 bar, the benchmark's span, Qlike, Tmax: Random forest, HAR
\providecommand{\rhmcsChQlikeTmaxPRfFull}{}\renewcommand{\rhmcsChQlikeTmaxPRfFull}{0.102}  % MCS p-value, 16:00 bar, the benchmark's span, Qlike, Tmax: Random forest, Full
\providecommand{\rhmcsChQlikeTmaxPBench}{}\renewcommand{\rhmcsChQlikeTmaxPBench}{1.00}  % MCS p-value, 16:00 bar, the benchmark's span, Qlike, Tmax: Closing-interval HAR benchmark
\providecommand{\rhmcsChQlikeTmaxNNinety}{}\renewcommand{\rhmcsChQlikeTmaxNNinety}{7}  % models in the 90 % MCS (of 11)
\providecommand{\rhmcsChQlikeTmaxSetNinety}{}\renewcommand{\rhmcsChQlikeTmaxSetNinety}{the closing-interval HAR benchmark, XGBoost on the full predictor set, LightGBM on the full predictor set, LightGBM on the HAR benchmark, XGBoost on the HAR benchmark, the random forest on the full predictor set, and the elastic net on the full predictor set}  % the 90 % MCS, best (last eliminated) first
\providecommand{\rhmcsChQlikeTmaxStableNinety}{}\renewcommand{\rhmcsChQlikeTmaxStableNinety}{no}  % the 90 % MCS is the same at block lengths (5, 20)
\providecommand{\rhmcsChQlikeTmaxPMaxOutNinety}{}\renewcommand{\rhmcsChQlikeTmaxPMaxOutNinety}{0.050}  % largest MCS p-value outside the 90 % set
\providecommand{\rhmcsChQlikeTmaxNNinetyLo}{}\renewcommand{\rhmcsChQlikeTmaxNNinetyLo}{3}  % ... at block length 5
\providecommand{\rhmcsChQlikeTmaxSetNinetyLo}{}\renewcommand{\rhmcsChQlikeTmaxSetNinetyLo}{the closing-interval HAR benchmark, XGBoost on the full predictor set, and LightGBM on the full predictor set}  % ... at block length 5
\providecommand{\rhmcsChQlikeTmaxNNinetyHi}{}\renewcommand{\rhmcsChQlikeTmaxNNinetyHi}{7}  % ... at block length 20
\providecommand{\rhmcsChQlikeTmaxSetNinetyHi}{}\renewcommand{\rhmcsChQlikeTmaxSetNinetyHi}{the closing-interval HAR benchmark, XGBoost on the full predictor set, LightGBM on the full predictor set, LightGBM on the HAR benchmark, XGBoost on the HAR benchmark, the random forest on the full predictor set, and the elastic net on the full predictor set}  % ... at block length 20
\providecommand{\rhmcsChQlikeTmaxNSevenFive}{}\renewcommand{\rhmcsChQlikeTmaxNSevenFive}{1}  % models in the 75 % MCS (of 11)
\providecommand{\rhmcsChQlikeTmaxSetSevenFive}{}\renewcommand{\rhmcsChQlikeTmaxSetSevenFive}{the closing-interval HAR benchmark}  % the 75 % MCS, best (last eliminated) first
\providecommand{\rhmcsChQlikeTmaxStableSevenFive}{}\renewcommand{\rhmcsChQlikeTmaxStableSevenFive}{yes}  % the 75 % MCS is the same at block lengths (5, 20)
\providecommand{\rhmcsChQlikeTmaxPMaxOutSevenFive}{}\renewcommand{\rhmcsChQlikeTmaxPMaxOutSevenFive}{0.22}  % largest MCS p-value outside the 75 % set
\providecommand{\rhmcsChQlikeTmaxNSevenFiveLo}{}\renewcommand{\rhmcsChQlikeTmaxNSevenFiveLo}{1}  % ... at block length 5
\providecommand{\rhmcsChQlikeTmaxSetSevenFiveLo}{}\renewcommand{\rhmcsChQlikeTmaxSetSevenFiveLo}{the closing-interval HAR benchmark}  % ... at block length 5
\providecommand{\rhmcsChQlikeTmaxNSevenFiveHi}{}\renewcommand{\rhmcsChQlikeTmaxNSevenFiveHi}{1}  % ... at block length 20
\providecommand{\rhmcsChQlikeTmaxSetSevenFiveHi}{}\renewcommand{\rhmcsChQlikeTmaxSetSevenFiveHi}{the closing-interval HAR benchmark}  % ... at block length 20
\providecommand{\rhmcsChQlikeTmaxBest}{}\renewcommand{\rhmcsChQlikeTmaxBest}{the closing-interval HAR benchmark}  % the last model left (MCS p-value 1)
\providecommand{\rhmcsChQlikeTrangePLinHar}{}\renewcommand{\rhmcsChQlikeTrangePLinHar}{\ensuremath{<}0.001}  % MCS p-value, 16:00 bar, the benchmark's span, Qlike, Trange: Linear, HAR
\providecommand{\rhmcsChQlikeTrangePRidgeFull}{}\renewcommand{\rhmcsChQlikeTrangePRidgeFull}{0.008}  % MCS p-value, 16:00 bar, the benchmark's span, Qlike, Trange: Ridge, Full
\providecommand{\rhmcsChQlikeTrangePLassoFull}{}\renewcommand{\rhmcsChQlikeTrangePLassoFull}{0.007}  % MCS p-value, 16:00 bar, the benchmark's span, Qlike, Trange: Lasso, Full
\providecommand{\rhmcsChQlikeTrangePEnetFull}{}\renewcommand{\rhmcsChQlikeTrangePEnetFull}{0.008}  % MCS p-value, 16:00 bar, the benchmark's span, Qlike, Trange: Elastic net, Full
\providecommand{\rhmcsChQlikeTrangePLgbmHar}{}\renewcommand{\rhmcsChQlikeTrangePLgbmHar}{0.008}  % MCS p-value, 16:00 bar, the benchmark's span, Qlike, Trange: LightGBM, HAR
\providecommand{\rhmcsChQlikeTrangePLgbmFull}{}\renewcommand{\rhmcsChQlikeTrangePLgbmFull}{0.31}  % MCS p-value, 16:00 bar, the benchmark's span, Qlike, Trange: LightGBM, Full
\providecommand{\rhmcsChQlikeTrangePXgbHar}{}\renewcommand{\rhmcsChQlikeTrangePXgbHar}{0.007}  % MCS p-value, 16:00 bar, the benchmark's span, Qlike, Trange: XGBoost, HAR
\providecommand{\rhmcsChQlikeTrangePXgbFull}{}\renewcommand{\rhmcsChQlikeTrangePXgbFull}{0.31}  % MCS p-value, 16:00 bar, the benchmark's span, Qlike, Trange: XGBoost, Full
\providecommand{\rhmcsChQlikeTrangePRfHar}{}\renewcommand{\rhmcsChQlikeTrangePRfHar}{\ensuremath{<}0.001}  % MCS p-value, 16:00 bar, the benchmark's span, Qlike, Trange: Random forest, HAR
\providecommand{\rhmcsChQlikeTrangePRfFull}{}\renewcommand{\rhmcsChQlikeTrangePRfFull}{\ensuremath{<}0.001}  % MCS p-value, 16:00 bar, the benchmark's span, Qlike, Trange: Random forest, Full
\providecommand{\rhmcsChQlikeTrangePBench}{}\renewcommand{\rhmcsChQlikeTrangePBench}{1.00}  % MCS p-value, 16:00 bar, the benchmark's span, Qlike, Trange: Closing-interval HAR benchmark
\providecommand{\rhmcsChQlikeTrangeNNinety}{}\renewcommand{\rhmcsChQlikeTrangeNNinety}{3}  % models in the 90 % MCS (of 11)
\providecommand{\rhmcsChQlikeTrangeSetNinety}{}\renewcommand{\rhmcsChQlikeTrangeSetNinety}{the closing-interval HAR benchmark, XGBoost on the full predictor set, and LightGBM on the full predictor set}  % the 90 % MCS, best (last eliminated) first
\providecommand{\rhmcsChQlikeTrangeStableNinety}{}\renewcommand{\rhmcsChQlikeTrangeStableNinety}{yes}  % the 90 % MCS is the same at block lengths (5, 20)
\providecommand{\rhmcsChQlikeTrangePMaxOutNinety}{}\renewcommand{\rhmcsChQlikeTrangePMaxOutNinety}{0.008}  % largest MCS p-value outside the 90 % set
\providecommand{\rhmcsChQlikeTrangeNNinetyLo}{}\renewcommand{\rhmcsChQlikeTrangeNNinetyLo}{3}  % ... at block length 5
\providecommand{\rhmcsChQlikeTrangeSetNinetyLo}{}\renewcommand{\rhmcsChQlikeTrangeSetNinetyLo}{the closing-interval HAR benchmark, XGBoost on the full predictor set, and LightGBM on the full predictor set}  % ... at block length 5
\providecommand{\rhmcsChQlikeTrangeNNinetyHi}{}\renewcommand{\rhmcsChQlikeTrangeNNinetyHi}{3}  % ... at block length 20
\providecommand{\rhmcsChQlikeTrangeSetNinetyHi}{}\renewcommand{\rhmcsChQlikeTrangeSetNinetyHi}{the closing-interval HAR benchmark, XGBoost on the full predictor set, and LightGBM on the full predictor set}  % ... at block length 20
\providecommand{\rhmcsChQlikeTrangeNSevenFive}{}\renewcommand{\rhmcsChQlikeTrangeNSevenFive}{3}  % models in the 75 % MCS (of 11)
\providecommand{\rhmcsChQlikeTrangeSetSevenFive}{}\renewcommand{\rhmcsChQlikeTrangeSetSevenFive}{the closing-interval HAR benchmark, XGBoost on the full predictor set, and LightGBM on the full predictor set}  % the 75 % MCS, best (last eliminated) first
\providecommand{\rhmcsChQlikeTrangeStableSevenFive}{}\renewcommand{\rhmcsChQlikeTrangeStableSevenFive}{yes}  % the 75 % MCS is the same at block lengths (5, 20)
\providecommand{\rhmcsChQlikeTrangePMaxOutSevenFive}{}\renewcommand{\rhmcsChQlikeTrangePMaxOutSevenFive}{0.008}  % largest MCS p-value outside the 75 % set
\providecommand{\rhmcsChQlikeTrangeNSevenFiveLo}{}\renewcommand{\rhmcsChQlikeTrangeNSevenFiveLo}{3}  % ... at block length 5
\providecommand{\rhmcsChQlikeTrangeSetSevenFiveLo}{}\renewcommand{\rhmcsChQlikeTrangeSetSevenFiveLo}{the closing-interval HAR benchmark, XGBoost on the full predictor set, and LightGBM on the full predictor set}  % ... at block length 5
\providecommand{\rhmcsChQlikeTrangeNSevenFiveHi}{}\renewcommand{\rhmcsChQlikeTrangeNSevenFiveHi}{3}  % ... at block length 20
\providecommand{\rhmcsChQlikeTrangeSetSevenFiveHi}{}\renewcommand{\rhmcsChQlikeTrangeSetSevenFiveHi}{the closing-interval HAR benchmark, XGBoost on the full predictor set, and LightGBM on the full predictor set}  % ... at block length 20
\providecommand{\rhmcsChQlikeTrangeBest}{}\renewcommand{\rhmcsChQlikeTrangeBest}{the closing-interval HAR benchmark}  % the last model left (MCS p-value 1)
\providecommand{\rhmcsChMseTmaxPLinHar}{}\renewcommand{\rhmcsChMseTmaxPLinHar}{0.59}  % MCS p-value, 16:00 bar, the benchmark's span, Mse, Tmax: Linear, HAR
\providecommand{\rhmcsChMseTmaxPRidgeFull}{}\renewcommand{\rhmcsChMseTmaxPRidgeFull}{0.47}  % MCS p-value, 16:00 bar, the benchmark's span, Mse, Tmax: Ridge, Full
\providecommand{\rhmcsChMseTmaxPLassoFull}{}\renewcommand{\rhmcsChMseTmaxPLassoFull}{0.42}  % MCS p-value, 16:00 bar, the benchmark's span, Mse, Tmax: Lasso, Full
\providecommand{\rhmcsChMseTmaxPEnetFull}{}\renewcommand{\rhmcsChMseTmaxPEnetFull}{0.53}  % MCS p-value, 16:00 bar, the benchmark's span, Mse, Tmax: Elastic net, Full
\providecommand{\rhmcsChMseTmaxPLgbmHar}{}\renewcommand{\rhmcsChMseTmaxPLgbmHar}{0.59}  % MCS p-value, 16:00 bar, the benchmark's span, Mse, Tmax: LightGBM, HAR
\providecommand{\rhmcsChMseTmaxPLgbmFull}{}\renewcommand{\rhmcsChMseTmaxPLgbmFull}{0.54}  % MCS p-value, 16:00 bar, the benchmark's span, Mse, Tmax: LightGBM, Full
\providecommand{\rhmcsChMseTmaxPXgbHar}{}\renewcommand{\rhmcsChMseTmaxPXgbHar}{0.59}  % MCS p-value, 16:00 bar, the benchmark's span, Mse, Tmax: XGBoost, HAR
\providecommand{\rhmcsChMseTmaxPXgbFull}{}\renewcommand{\rhmcsChMseTmaxPXgbFull}{0.58}  % MCS p-value, 16:00 bar, the benchmark's span, Mse, Tmax: XGBoost, Full
\providecommand{\rhmcsChMseTmaxPRfHar}{}\renewcommand{\rhmcsChMseTmaxPRfHar}{1.00}  % MCS p-value, 16:00 bar, the benchmark's span, Mse, Tmax: Random forest, HAR
\providecommand{\rhmcsChMseTmaxPRfFull}{}\renewcommand{\rhmcsChMseTmaxPRfFull}{0.59}  % MCS p-value, 16:00 bar, the benchmark's span, Mse, Tmax: Random forest, Full
\providecommand{\rhmcsChMseTmaxPBench}{}\renewcommand{\rhmcsChMseTmaxPBench}{0.58}  % MCS p-value, 16:00 bar, the benchmark's span, Mse, Tmax: Closing-interval HAR benchmark
\providecommand{\rhmcsChMseTmaxNNinety}{}\renewcommand{\rhmcsChMseTmaxNNinety}{11}  % models in the 90 % MCS (of 11)
\providecommand{\rhmcsChMseTmaxSetNinety}{}\renewcommand{\rhmcsChMseTmaxSetNinety}{the random forest on the HAR benchmark, XGBoost on the HAR benchmark, LightGBM on the HAR benchmark, the linear models on the HAR benchmark, the random forest on the full predictor set, the closing-interval HAR benchmark, XGBoost on the full predictor set, LightGBM on the full predictor set, the elastic net on the full predictor set, ridge regression on the full predictor set, and the lasso on the full predictor set}  % the 90 % MCS, best (last eliminated) first
\providecommand{\rhmcsChMseTmaxStableNinety}{}\renewcommand{\rhmcsChMseTmaxStableNinety}{yes}  % the 90 % MCS is the same at block lengths (5, 20)
\providecommand{\rhmcsChMseTmaxPMaxOutNinety}{}\renewcommand{\rhmcsChMseTmaxPMaxOutNinety}{--}  % largest MCS p-value outside the 90 % set
\providecommand{\rhmcsChMseTmaxNNinetyLo}{}\renewcommand{\rhmcsChMseTmaxNNinetyLo}{11}  % ... at block length 5
\providecommand{\rhmcsChMseTmaxSetNinetyLo}{}\renewcommand{\rhmcsChMseTmaxSetNinetyLo}{the random forest on the HAR benchmark, XGBoost on the HAR benchmark, LightGBM on the HAR benchmark, the linear models on the HAR benchmark, the random forest on the full predictor set, the closing-interval HAR benchmark, XGBoost on the full predictor set, LightGBM on the full predictor set, the elastic net on the full predictor set, ridge regression on the full predictor set, and the lasso on the full predictor set}  % ... at block length 5
\providecommand{\rhmcsChMseTmaxNNinetyHi}{}\renewcommand{\rhmcsChMseTmaxNNinetyHi}{11}  % ... at block length 20
\providecommand{\rhmcsChMseTmaxSetNinetyHi}{}\renewcommand{\rhmcsChMseTmaxSetNinetyHi}{the random forest on the HAR benchmark, XGBoost on the HAR benchmark, LightGBM on the HAR benchmark, the linear models on the HAR benchmark, the random forest on the full predictor set, the closing-interval HAR benchmark, XGBoost on the full predictor set, LightGBM on the full predictor set, the elastic net on the full predictor set, ridge regression on the full predictor set, and the lasso on the full predictor set}  % ... at block length 20
\providecommand{\rhmcsChMseTmaxNSevenFive}{}\renewcommand{\rhmcsChMseTmaxNSevenFive}{11}  % models in the 75 % MCS (of 11)
\providecommand{\rhmcsChMseTmaxSetSevenFive}{}\renewcommand{\rhmcsChMseTmaxSetSevenFive}{the random forest on the HAR benchmark, XGBoost on the HAR benchmark, LightGBM on the HAR benchmark, the linear models on the HAR benchmark, the random forest on the full predictor set, the closing-interval HAR benchmark, XGBoost on the full predictor set, LightGBM on the full predictor set, the elastic net on the full predictor set, ridge regression on the full predictor set, and the lasso on the full predictor set}  % the 75 % MCS, best (last eliminated) first
\providecommand{\rhmcsChMseTmaxStableSevenFive}{}\renewcommand{\rhmcsChMseTmaxStableSevenFive}{yes}  % the 75 % MCS is the same at block lengths (5, 20)
\providecommand{\rhmcsChMseTmaxPMaxOutSevenFive}{}\renewcommand{\rhmcsChMseTmaxPMaxOutSevenFive}{--}  % largest MCS p-value outside the 75 % set
\providecommand{\rhmcsChMseTmaxNSevenFiveLo}{}\renewcommand{\rhmcsChMseTmaxNSevenFiveLo}{11}  % ... at block length 5
\providecommand{\rhmcsChMseTmaxSetSevenFiveLo}{}\renewcommand{\rhmcsChMseTmaxSetSevenFiveLo}{the random forest on the HAR benchmark, XGBoost on the HAR benchmark, LightGBM on the HAR benchmark, the linear models on the HAR benchmark, the random forest on the full predictor set, the closing-interval HAR benchmark, XGBoost on the full predictor set, LightGBM on the full predictor set, the elastic net on the full predictor set, ridge regression on the full predictor set, and the lasso on the full predictor set}  % ... at block length 5
\providecommand{\rhmcsChMseTmaxNSevenFiveHi}{}\renewcommand{\rhmcsChMseTmaxNSevenFiveHi}{11}  % ... at block length 20
\providecommand{\rhmcsChMseTmaxSetSevenFiveHi}{}\renewcommand{\rhmcsChMseTmaxSetSevenFiveHi}{the random forest on the HAR benchmark, XGBoost on the HAR benchmark, LightGBM on the HAR benchmark, the linear models on the HAR benchmark, the random forest on the full predictor set, the closing-interval HAR benchmark, XGBoost on the full predictor set, LightGBM on the full predictor set, the elastic net on the full predictor set, ridge regression on the full predictor set, and the lasso on the full predictor set}  % ... at block length 20
\providecommand{\rhmcsChMseTmaxBest}{}\renewcommand{\rhmcsChMseTmaxBest}{the random forest on the HAR benchmark}  % the last model left (MCS p-value 1)
\providecommand{\rhmcsChMseTrangePLinHar}{}\renewcommand{\rhmcsChMseTrangePLinHar}{0.60}  % MCS p-value, 16:00 bar, the benchmark's span, Mse, Trange: Linear, HAR
\providecommand{\rhmcsChMseTrangePRidgeFull}{}\renewcommand{\rhmcsChMseTrangePRidgeFull}{0.60}  % MCS p-value, 16:00 bar, the benchmark's span, Mse, Trange: Ridge, Full
\providecommand{\rhmcsChMseTrangePLassoFull}{}\renewcommand{\rhmcsChMseTrangePLassoFull}{0.60}  % MCS p-value, 16:00 bar, the benchmark's span, Mse, Trange: Lasso, Full
\providecommand{\rhmcsChMseTrangePEnetFull}{}\renewcommand{\rhmcsChMseTrangePEnetFull}{0.60}  % MCS p-value, 16:00 bar, the benchmark's span, Mse, Trange: Elastic net, Full
\providecommand{\rhmcsChMseTrangePLgbmHar}{}\renewcommand{\rhmcsChMseTrangePLgbmHar}{0.41}  % MCS p-value, 16:00 bar, the benchmark's span, Mse, Trange: LightGBM, HAR
\providecommand{\rhmcsChMseTrangePLgbmFull}{}\renewcommand{\rhmcsChMseTrangePLgbmFull}{0.60}  % MCS p-value, 16:00 bar, the benchmark's span, Mse, Trange: LightGBM, Full
\providecommand{\rhmcsChMseTrangePXgbHar}{}\renewcommand{\rhmcsChMseTrangePXgbHar}{0.60}  % MCS p-value, 16:00 bar, the benchmark's span, Mse, Trange: XGBoost, HAR
\providecommand{\rhmcsChMseTrangePXgbFull}{}\renewcommand{\rhmcsChMseTrangePXgbFull}{0.60}  % MCS p-value, 16:00 bar, the benchmark's span, Mse, Trange: XGBoost, Full
\providecommand{\rhmcsChMseTrangePRfHar}{}\renewcommand{\rhmcsChMseTrangePRfHar}{1.00}  % MCS p-value, 16:00 bar, the benchmark's span, Mse, Trange: Random forest, HAR
\providecommand{\rhmcsChMseTrangePRfFull}{}\renewcommand{\rhmcsChMseTrangePRfFull}{0.60}  % MCS p-value, 16:00 bar, the benchmark's span, Mse, Trange: Random forest, Full
\providecommand{\rhmcsChMseTrangePBench}{}\renewcommand{\rhmcsChMseTrangePBench}{0.60}  % MCS p-value, 16:00 bar, the benchmark's span, Mse, Trange: Closing-interval HAR benchmark
\providecommand{\rhmcsChMseTrangeNNinety}{}\renewcommand{\rhmcsChMseTrangeNNinety}{11}  % models in the 90 % MCS (of 11)
\providecommand{\rhmcsChMseTrangeSetNinety}{}\renewcommand{\rhmcsChMseTrangeSetNinety}{the random forest on the HAR benchmark, the linear models on the HAR benchmark, XGBoost on the HAR benchmark, the random forest on the full predictor set, XGBoost on the full predictor set, the elastic net on the full predictor set, ridge regression on the full predictor set, the lasso on the full predictor set, LightGBM on the full predictor set, the closing-interval HAR benchmark, and LightGBM on the HAR benchmark}  % the 90 % MCS, best (last eliminated) first
\providecommand{\rhmcsChMseTrangeStableNinety}{}\renewcommand{\rhmcsChMseTrangeStableNinety}{yes}  % the 90 % MCS is the same at block lengths (5, 20)
\providecommand{\rhmcsChMseTrangePMaxOutNinety}{}\renewcommand{\rhmcsChMseTrangePMaxOutNinety}{--}  % largest MCS p-value outside the 90 % set
\providecommand{\rhmcsChMseTrangeNNinetyLo}{}\renewcommand{\rhmcsChMseTrangeNNinetyLo}{11}  % ... at block length 5
\providecommand{\rhmcsChMseTrangeSetNinetyLo}{}\renewcommand{\rhmcsChMseTrangeSetNinetyLo}{the random forest on the HAR benchmark, the linear models on the HAR benchmark, XGBoost on the HAR benchmark, XGBoost on the full predictor set, the random forest on the full predictor set, the elastic net on the full predictor set, ridge regression on the full predictor set, LightGBM on the full predictor set, the lasso on the full predictor set, the closing-interval HAR benchmark, and LightGBM on the HAR benchmark}  % ... at block length 5
\providecommand{\rhmcsChMseTrangeNNinetyHi}{}\renewcommand{\rhmcsChMseTrangeNNinetyHi}{11}  % ... at block length 20
\providecommand{\rhmcsChMseTrangeSetNinetyHi}{}\renewcommand{\rhmcsChMseTrangeSetNinetyHi}{the random forest on the HAR benchmark, the linear models on the HAR benchmark, the random forest on the full predictor set, XGBoost on the full predictor set, the elastic net on the full predictor set, ridge regression on the full predictor set, the lasso on the full predictor set, XGBoost on the HAR benchmark, the closing-interval HAR benchmark, LightGBM on the full predictor set, and LightGBM on the HAR benchmark}  % ... at block length 20
\providecommand{\rhmcsChMseTrangeNSevenFive}{}\renewcommand{\rhmcsChMseTrangeNSevenFive}{11}  % models in the 75 % MCS (of 11)
\providecommand{\rhmcsChMseTrangeSetSevenFive}{}\renewcommand{\rhmcsChMseTrangeSetSevenFive}{the random forest on the HAR benchmark, the linear models on the HAR benchmark, XGBoost on the HAR benchmark, the random forest on the full predictor set, XGBoost on the full predictor set, the elastic net on the full predictor set, ridge regression on the full predictor set, the lasso on the full predictor set, LightGBM on the full predictor set, the closing-interval HAR benchmark, and LightGBM on the HAR benchmark}  % the 75 % MCS, best (last eliminated) first
\providecommand{\rhmcsChMseTrangeStableSevenFive}{}\renewcommand{\rhmcsChMseTrangeStableSevenFive}{yes}  % the 75 % MCS is the same at block lengths (5, 20)
\providecommand{\rhmcsChMseTrangePMaxOutSevenFive}{}\renewcommand{\rhmcsChMseTrangePMaxOutSevenFive}{--}  % largest MCS p-value outside the 75 % set
\providecommand{\rhmcsChMseTrangeNSevenFiveLo}{}\renewcommand{\rhmcsChMseTrangeNSevenFiveLo}{11}  % ... at block length 5
\providecommand{\rhmcsChMseTrangeSetSevenFiveLo}{}\renewcommand{\rhmcsChMseTrangeSetSevenFiveLo}{the random forest on the HAR benchmark, the linear models on the HAR benchmark, XGBoost on the HAR benchmark, XGBoost on the full predictor set, the random forest on the full predictor set, the elastic net on the full predictor set, ridge regression on the full predictor set, LightGBM on the full predictor set, the lasso on the full predictor set, the closing-interval HAR benchmark, and LightGBM on the HAR benchmark}  % ... at block length 5
\providecommand{\rhmcsChMseTrangeNSevenFiveHi}{}\renewcommand{\rhmcsChMseTrangeNSevenFiveHi}{11}  % ... at block length 20
\providecommand{\rhmcsChMseTrangeSetSevenFiveHi}{}\renewcommand{\rhmcsChMseTrangeSetSevenFiveHi}{the random forest on the HAR benchmark, the linear models on the HAR benchmark, the random forest on the full predictor set, XGBoost on the full predictor set, the elastic net on the full predictor set, ridge regression on the full predictor set, the lasso on the full predictor set, XGBoost on the HAR benchmark, the closing-interval HAR benchmark, LightGBM on the full predictor set, and LightGBM on the HAR benchmark}  % ... at block length 20
\providecommand{\rhmcsChMseTrangeBest}{}\renewcommand{\rhmcsChMseTrangeBest}{the random forest on the HAR benchmark}  % the last model left (MCS p-value 1)
\providecommand{\rhmcsBarQAfSpaC}{}\renewcommand{\rhmcsBarQAfSpaC}{\ensuremath{<}0.001}  % SPA p-value (consistent), family BarQAf
\providecommand{\rhmcsBarQAfSpaL}{}\renewcommand{\rhmcsBarQAfSpaL}{\ensuremath{<}0.001}  % ... lower
\providecommand{\rhmcsBarQAfSpaU}{}\renewcommand{\rhmcsBarQAfSpaU}{\ensuremath{<}0.001}  % ... upper
\providecommand{\rhmcsBarQAfNRej}{}\renewcommand{\rhmcsBarQAfNRej}{6}  % StepM rejections at FWER 5 % (of 6)
\providecommand{\rhmcsBarQAfNRejWord}{}\renewcommand{\rhmcsBarQAfNRejWord}{six}  % ... as a word
\providecommand{\rhmcsBarQAfNRaw}{}\renewcommand{\rhmcsBarQAfNRaw}{6}  % single-test bootstrap p < 0.05 (no adjustment)
\providecommand{\rhmcsBarQAfNRawWord}{}\renewcommand{\rhmcsBarQAfNRawWord}{six}  % ... as a word
\providecommand{\rhmcsBarQAfNPos}{}\renewcommand{\rhmcsBarQAfNPos}{6}  % alternatives with a positive estimate (better than the benchmark)
\providecommand{\rhmcsBarQAfPMin}{}\renewcommand{\rhmcsBarQAfPMin}{\ensuremath{<}0.001}  % smallest StepM-adjusted p-value
\providecommand{\rhmcsBarQAfPMax}{}\renewcommand{\rhmcsBarQAfPMax}{\ensuremath{<}0.001}  % largest StepM-adjusted p-value
\providecommand{\rhmcsBarQAfPTwoMin}{}\renewcommand{\rhmcsBarQAfPTwoMin}{\ensuremath{<}0.001}  % smallest two-sided StepM-adjusted p-value
\providecommand{\rhmcsBarQAfPTwoMax}{}\renewcommand{\rhmcsBarQAfPTwoMax}{\ensuremath{<}0.001}  % largest two-sided StepM-adjusted p-value
\providecommand{\rhmcsBarQAfRawMin}{}\renewcommand{\rhmcsBarQAfRawMin}{\ensuremath{<}0.001}  % smallest single-test bootstrap p-value
\providecommand{\rhmcsBarQAfRej}{}\renewcommand{\rhmcsBarQAfRej}{ridge regression on the full predictor set, the lasso on the full predictor set, the elastic net on the full predictor set, LightGBM on the full predictor set, XGBoost on the full predictor set, and the random forest on the full predictor set}  % the alternatives StepM rejects at 5 % (none: none)
\providecommand{\rhmcsBarQAfNRejTwo}{}\renewcommand{\rhmcsBarQAfNRejTwo}{6}  % StepM rejections at FWER 5 %, two-sided
\providecommand{\rhmcsBarQAfNRejTwoWord}{}\renewcommand{\rhmcsBarQAfNRejTwoWord}{six}  % ... as a word
\providecommand{\rhmcsBarQAfNRawTwo}{}\renewcommand{\rhmcsBarQAfNRawTwo}{6}  % single-test bootstrap p < 0.05, two-sided
\providecommand{\rhmcsBarQAfRejTwo}{}\renewcommand{\rhmcsBarQAfRejTwo}{ridge regression on the full predictor set, the lasso on the full predictor set, the elastic net on the full predictor set, LightGBM on the full predictor set, XGBoost on the full predictor set, and the random forest on the full predictor set}  % the alternatives the two-sided StepM rejects at 5 % (none: none)
\providecommand{\rhmcsBarQAfPRidge}{}\renewcommand{\rhmcsBarQAfPRidge}{\ensuremath{<}0.001}  % StepM-adjusted p-value (one-sided), Ridge
\providecommand{\rhmcsBarQAfRawRidge}{}\renewcommand{\rhmcsBarQAfRawRidge}{\ensuremath{<}0.001}  % single-test bootstrap p-value (one-sided), Ridge
\providecommand{\rhmcsBarQAfPTwoRidge}{}\renewcommand{\rhmcsBarQAfPTwoRidge}{\ensuremath{<}0.001}  % StepM-adjusted p-value (two-sided), Ridge
\providecommand{\rhmcsBarQAfRawTwoRidge}{}\renewcommand{\rhmcsBarQAfRawTwoRidge}{\ensuremath{<}0.001}  % single-test bootstrap p-value (two-sided), Ridge
\providecommand{\rhmcsBarQAfPLasso}{}\renewcommand{\rhmcsBarQAfPLasso}{\ensuremath{<}0.001}  % StepM-adjusted p-value (one-sided), Lasso
\providecommand{\rhmcsBarQAfRawLasso}{}\renewcommand{\rhmcsBarQAfRawLasso}{\ensuremath{<}0.001}  % single-test bootstrap p-value (one-sided), Lasso
\providecommand{\rhmcsBarQAfPTwoLasso}{}\renewcommand{\rhmcsBarQAfPTwoLasso}{\ensuremath{<}0.001}  % StepM-adjusted p-value (two-sided), Lasso
\providecommand{\rhmcsBarQAfRawTwoLasso}{}\renewcommand{\rhmcsBarQAfRawTwoLasso}{\ensuremath{<}0.001}  % single-test bootstrap p-value (two-sided), Lasso
\providecommand{\rhmcsBarQAfPEnet}{}\renewcommand{\rhmcsBarQAfPEnet}{\ensuremath{<}0.001}  % StepM-adjusted p-value (one-sided), Elastic net
\providecommand{\rhmcsBarQAfRawEnet}{}\renewcommand{\rhmcsBarQAfRawEnet}{\ensuremath{<}0.001}  % single-test bootstrap p-value (one-sided), Elastic net
\providecommand{\rhmcsBarQAfPTwoEnet}{}\renewcommand{\rhmcsBarQAfPTwoEnet}{\ensuremath{<}0.001}  % StepM-adjusted p-value (two-sided), Elastic net
\providecommand{\rhmcsBarQAfRawTwoEnet}{}\renewcommand{\rhmcsBarQAfRawTwoEnet}{\ensuremath{<}0.001}  % single-test bootstrap p-value (two-sided), Elastic net
\providecommand{\rhmcsBarQAfPLgbm}{}\renewcommand{\rhmcsBarQAfPLgbm}{\ensuremath{<}0.001}  % StepM-adjusted p-value (one-sided), LightGBM
\providecommand{\rhmcsBarQAfRawLgbm}{}\renewcommand{\rhmcsBarQAfRawLgbm}{\ensuremath{<}0.001}  % single-test bootstrap p-value (one-sided), LightGBM
\providecommand{\rhmcsBarQAfPTwoLgbm}{}\renewcommand{\rhmcsBarQAfPTwoLgbm}{\ensuremath{<}0.001}  % StepM-adjusted p-value (two-sided), LightGBM
\providecommand{\rhmcsBarQAfRawTwoLgbm}{}\renewcommand{\rhmcsBarQAfRawTwoLgbm}{\ensuremath{<}0.001}  % single-test bootstrap p-value (two-sided), LightGBM
\providecommand{\rhmcsBarQAfPXgb}{}\renewcommand{\rhmcsBarQAfPXgb}{\ensuremath{<}0.001}  % StepM-adjusted p-value (one-sided), XGBoost
\providecommand{\rhmcsBarQAfRawXgb}{}\renewcommand{\rhmcsBarQAfRawXgb}{\ensuremath{<}0.001}  % single-test bootstrap p-value (one-sided), XGBoost
\providecommand{\rhmcsBarQAfPTwoXgb}{}\renewcommand{\rhmcsBarQAfPTwoXgb}{\ensuremath{<}0.001}  % StepM-adjusted p-value (two-sided), XGBoost
\providecommand{\rhmcsBarQAfRawTwoXgb}{}\renewcommand{\rhmcsBarQAfRawTwoXgb}{\ensuremath{<}0.001}  % single-test bootstrap p-value (two-sided), XGBoost
\providecommand{\rhmcsBarQAfPRf}{}\renewcommand{\rhmcsBarQAfPRf}{\ensuremath{<}0.001}  % StepM-adjusted p-value (one-sided), Random forest
\providecommand{\rhmcsBarQAfRawRf}{}\renewcommand{\rhmcsBarQAfRawRf}{\ensuremath{<}0.001}  % single-test bootstrap p-value (one-sided), Random forest
\providecommand{\rhmcsBarQAfPTwoRf}{}\renewcommand{\rhmcsBarQAfPTwoRf}{\ensuremath{<}0.001}  % StepM-adjusted p-value (two-sided), Random forest
\providecommand{\rhmcsBarQAfRawTwoRf}{}\renewcommand{\rhmcsBarQAfRawTwoRf}{\ensuremath{<}0.001}  % single-test bootstrap p-value (two-sided), Random forest
\providecommand{\rhmcsBarQAfNRejLo}{}\renewcommand{\rhmcsBarQAfNRejLo}{6}  % StepM rejections at 5 %, block length 5
\providecommand{\rhmcsBarQAfSpaCLo}{}\renewcommand{\rhmcsBarQAfSpaCLo}{\ensuremath{<}0.001}  % SPA p-value (consistent), block length 5
\providecommand{\rhmcsBarQAfNRejTwoLo}{}\renewcommand{\rhmcsBarQAfNRejTwoLo}{6}  % ... two-sided, block length 5
\providecommand{\rhmcsBarQAfPMinLo}{}\renewcommand{\rhmcsBarQAfPMinLo}{\ensuremath{<}0.001}  % smallest adjusted p, block length 5
\providecommand{\rhmcsBarQAfPTwoMinLo}{}\renewcommand{\rhmcsBarQAfPTwoMinLo}{\ensuremath{<}0.001}  % ... two-sided, block length 5
\providecommand{\rhmcsBarQAfNRejHi}{}\renewcommand{\rhmcsBarQAfNRejHi}{6}  % StepM rejections at 5 %, block length 20
\providecommand{\rhmcsBarQAfSpaCHi}{}\renewcommand{\rhmcsBarQAfSpaCHi}{\ensuremath{<}0.001}  % SPA p-value (consistent), block length 20
\providecommand{\rhmcsBarQAfNRejTwoHi}{}\renewcommand{\rhmcsBarQAfNRejTwoHi}{6}  % ... two-sided, block length 20
\providecommand{\rhmcsBarQAfPMinHi}{}\renewcommand{\rhmcsBarQAfPMinHi}{\ensuremath{<}0.001}  % smallest adjusted p, block length 20
\providecommand{\rhmcsBarQAfPTwoMinHi}{}\renewcommand{\rhmcsBarQAfPTwoMinHi}{\ensuremath{<}0.001}  % ... two-sided, block length 20
\providecommand{\rhmcsBarMAfSpaC}{}\renewcommand{\rhmcsBarMAfSpaC}{0.64}  % SPA p-value (consistent), family BarMAf
\providecommand{\rhmcsBarMAfSpaL}{}\renewcommand{\rhmcsBarMAfSpaL}{0.22}  % ... lower
\providecommand{\rhmcsBarMAfSpaU}{}\renewcommand{\rhmcsBarMAfSpaU}{0.64}  % ... upper
\providecommand{\rhmcsBarMAfNRej}{}\renewcommand{\rhmcsBarMAfNRej}{0}  % StepM rejections at FWER 5 % (of 6)
\providecommand{\rhmcsBarMAfNRejWord}{}\renewcommand{\rhmcsBarMAfNRejWord}{none}  % ... as a word
\providecommand{\rhmcsBarMAfNRaw}{}\renewcommand{\rhmcsBarMAfNRaw}{0}  % single-test bootstrap p < 0.05 (no adjustment)
\providecommand{\rhmcsBarMAfNRawWord}{}\renewcommand{\rhmcsBarMAfNRawWord}{none}  % ... as a word
\providecommand{\rhmcsBarMAfNPos}{}\renewcommand{\rhmcsBarMAfNPos}{0}  % alternatives with a positive estimate (better than the benchmark)
\providecommand{\rhmcsBarMAfPMin}{}\renewcommand{\rhmcsBarMAfPMin}{0.88}  % smallest StepM-adjusted p-value
\providecommand{\rhmcsBarMAfPMax}{}\renewcommand{\rhmcsBarMAfPMax}{0.89}  % largest StepM-adjusted p-value
\providecommand{\rhmcsBarMAfPTwoMin}{}\renewcommand{\rhmcsBarMAfPTwoMin}{0.28}  % smallest two-sided StepM-adjusted p-value
\providecommand{\rhmcsBarMAfPTwoMax}{}\renewcommand{\rhmcsBarMAfPTwoMax}{0.38}  % largest two-sided StepM-adjusted p-value
\providecommand{\rhmcsBarMAfRawMin}{}\renewcommand{\rhmcsBarMAfRawMin}{0.84}  % smallest single-test bootstrap p-value
\providecommand{\rhmcsBarMAfRej}{}\renewcommand{\rhmcsBarMAfRej}{none}  % the alternatives StepM rejects at 5 % (none: none)
\providecommand{\rhmcsBarMAfNRejTwo}{}\renewcommand{\rhmcsBarMAfNRejTwo}{0}  % StepM rejections at FWER 5 %, two-sided
\providecommand{\rhmcsBarMAfNRejTwoWord}{}\renewcommand{\rhmcsBarMAfNRejTwoWord}{none}  % ... as a word
\providecommand{\rhmcsBarMAfNRawTwo}{}\renewcommand{\rhmcsBarMAfNRawTwo}{0}  % single-test bootstrap p < 0.05, two-sided
\providecommand{\rhmcsBarMAfRejTwo}{}\renewcommand{\rhmcsBarMAfRejTwo}{none}  % the alternatives the two-sided StepM rejects at 5 % (none: none)
\providecommand{\rhmcsBarMAfPRidge}{}\renewcommand{\rhmcsBarMAfPRidge}{0.89}  % StepM-adjusted p-value (one-sided), Ridge
\providecommand{\rhmcsBarMAfRawRidge}{}\renewcommand{\rhmcsBarMAfRawRidge}{0.87}  % single-test bootstrap p-value (one-sided), Ridge
\providecommand{\rhmcsBarMAfPTwoRidge}{}\renewcommand{\rhmcsBarMAfPTwoRidge}{0.32}  % StepM-adjusted p-value (two-sided), Ridge
\providecommand{\rhmcsBarMAfRawTwoRidge}{}\renewcommand{\rhmcsBarMAfRawTwoRidge}{0.23}  % single-test bootstrap p-value (two-sided), Ridge
\providecommand{\rhmcsBarMAfPLasso}{}\renewcommand{\rhmcsBarMAfPLasso}{0.89}  % StepM-adjusted p-value (one-sided), Lasso
\providecommand{\rhmcsBarMAfRawLasso}{}\renewcommand{\rhmcsBarMAfRawLasso}{0.87}  % single-test bootstrap p-value (one-sided), Lasso
\providecommand{\rhmcsBarMAfPTwoLasso}{}\renewcommand{\rhmcsBarMAfPTwoLasso}{0.28}  % StepM-adjusted p-value (two-sided), Lasso
\providecommand{\rhmcsBarMAfRawTwoLasso}{}\renewcommand{\rhmcsBarMAfRawTwoLasso}{0.20}  % single-test bootstrap p-value (two-sided), Lasso
\providecommand{\rhmcsBarMAfPEnet}{}\renewcommand{\rhmcsBarMAfPEnet}{0.89}  % StepM-adjusted p-value (one-sided), Elastic net
\providecommand{\rhmcsBarMAfRawEnet}{}\renewcommand{\rhmcsBarMAfRawEnet}{0.87}  % single-test bootstrap p-value (one-sided), Elastic net
\providecommand{\rhmcsBarMAfPTwoEnet}{}\renewcommand{\rhmcsBarMAfPTwoEnet}{0.32}  % StepM-adjusted p-value (two-sided), Elastic net
\providecommand{\rhmcsBarMAfRawTwoEnet}{}\renewcommand{\rhmcsBarMAfRawTwoEnet}{0.23}  % single-test bootstrap p-value (two-sided), Elastic net
\providecommand{\rhmcsBarMAfPLgbm}{}\renewcommand{\rhmcsBarMAfPLgbm}{0.88}  % StepM-adjusted p-value (one-sided), LightGBM
\providecommand{\rhmcsBarMAfRawLgbm}{}\renewcommand{\rhmcsBarMAfRawLgbm}{0.84}  % single-test bootstrap p-value (one-sided), LightGBM
\providecommand{\rhmcsBarMAfPTwoLgbm}{}\renewcommand{\rhmcsBarMAfPTwoLgbm}{0.38}  % StepM-adjusted p-value (two-sided), LightGBM
\providecommand{\rhmcsBarMAfRawTwoLgbm}{}\renewcommand{\rhmcsBarMAfRawTwoLgbm}{0.36}  % single-test bootstrap p-value (two-sided), LightGBM
\providecommand{\rhmcsBarMAfPXgb}{}\renewcommand{\rhmcsBarMAfPXgb}{0.88}  % StepM-adjusted p-value (one-sided), XGBoost
\providecommand{\rhmcsBarMAfRawXgb}{}\renewcommand{\rhmcsBarMAfRawXgb}{0.84}  % single-test bootstrap p-value (one-sided), XGBoost
\providecommand{\rhmcsBarMAfPTwoXgb}{}\renewcommand{\rhmcsBarMAfPTwoXgb}{0.38}  % StepM-adjusted p-value (two-sided), XGBoost
\providecommand{\rhmcsBarMAfRawTwoXgb}{}\renewcommand{\rhmcsBarMAfRawTwoXgb}{0.35}  % single-test bootstrap p-value (two-sided), XGBoost
\providecommand{\rhmcsBarMAfPRf}{}\renewcommand{\rhmcsBarMAfPRf}{0.89}  % StepM-adjusted p-value (one-sided), Random forest
\providecommand{\rhmcsBarMAfRawRf}{}\renewcommand{\rhmcsBarMAfRawRf}{0.86}  % single-test bootstrap p-value (one-sided), Random forest
\providecommand{\rhmcsBarMAfPTwoRf}{}\renewcommand{\rhmcsBarMAfPTwoRf}{0.35}  % StepM-adjusted p-value (two-sided), Random forest
\providecommand{\rhmcsBarMAfRawTwoRf}{}\renewcommand{\rhmcsBarMAfRawTwoRf}{0.29}  % single-test bootstrap p-value (two-sided), Random forest
\providecommand{\rhmcsBarMAfNRejLo}{}\renewcommand{\rhmcsBarMAfNRejLo}{0}  % StepM rejections at 5 %, block length 5
\providecommand{\rhmcsBarMAfSpaCLo}{}\renewcommand{\rhmcsBarMAfSpaCLo}{0.64}  % SPA p-value (consistent), block length 5
\providecommand{\rhmcsBarMAfNRejTwoLo}{}\renewcommand{\rhmcsBarMAfNRejTwoLo}{0}  % ... two-sided, block length 5
\providecommand{\rhmcsBarMAfPMinLo}{}\renewcommand{\rhmcsBarMAfPMinLo}{0.91}  % smallest adjusted p, block length 5
\providecommand{\rhmcsBarMAfPTwoMinLo}{}\renewcommand{\rhmcsBarMAfPTwoMinLo}{0.21}  % ... two-sided, block length 5
\providecommand{\rhmcsBarMAfNRejHi}{}\renewcommand{\rhmcsBarMAfNRejHi}{0}  % StepM rejections at 5 %, block length 20
\providecommand{\rhmcsBarMAfSpaCHi}{}\renewcommand{\rhmcsBarMAfSpaCHi}{0.63}  % SPA p-value (consistent), block length 20
\providecommand{\rhmcsBarMAfNRejTwoHi}{}\renewcommand{\rhmcsBarMAfNRejTwoHi}{0}  % ... two-sided, block length 20
\providecommand{\rhmcsBarMAfPMinHi}{}\renewcommand{\rhmcsBarMAfPMinHi}{0.84}  % smallest adjusted p, block length 20
\providecommand{\rhmcsBarMAfPTwoMinHi}{}\renewcommand{\rhmcsBarMAfPTwoMinHi}{0.39}  % ... two-sided, block length 20
\providecommand{\rhmcsAllQAfSpaC}{}\renewcommand{\rhmcsAllQAfSpaC}{0.045}  % SPA p-value (consistent), family AllQAf
\providecommand{\rhmcsAllQAfSpaL}{}\renewcommand{\rhmcsAllQAfSpaL}{0.037}  % ... lower
\providecommand{\rhmcsAllQAfSpaU}{}\renewcommand{\rhmcsAllQAfSpaU}{0.045}  % ... upper
\providecommand{\rhmcsAllQAfNRej}{}\renewcommand{\rhmcsAllQAfNRej}{1}  % StepM rejections at FWER 5 % (of 6)
\providecommand{\rhmcsAllQAfNRejWord}{}\renewcommand{\rhmcsAllQAfNRejWord}{one}  % ... as a word
\providecommand{\rhmcsAllQAfNRaw}{}\renewcommand{\rhmcsAllQAfNRaw}{2}  % single-test bootstrap p < 0.05 (no adjustment)
\providecommand{\rhmcsAllQAfNRawWord}{}\renewcommand{\rhmcsAllQAfNRawWord}{two}  % ... as a word
\providecommand{\rhmcsAllQAfNPos}{}\renewcommand{\rhmcsAllQAfNPos}{3}  % alternatives with a positive estimate (better than the benchmark)
\providecommand{\rhmcsAllQAfPMin}{}\renewcommand{\rhmcsAllQAfPMin}{0.045}  % smallest StepM-adjusted p-value
\providecommand{\rhmcsAllQAfPMax}{}\renewcommand{\rhmcsAllQAfPMax}{0.90}  % largest StepM-adjusted p-value
\providecommand{\rhmcsAllQAfPTwoMin}{}\renewcommand{\rhmcsAllQAfPTwoMin}{0.11}  % smallest two-sided StepM-adjusted p-value
\providecommand{\rhmcsAllQAfPTwoMax}{}\renewcommand{\rhmcsAllQAfPTwoMax}{0.81}  % largest two-sided StepM-adjusted p-value
\providecommand{\rhmcsAllQAfRawMin}{}\renewcommand{\rhmcsAllQAfRawMin}{0.020}  % smallest single-test bootstrap p-value
\providecommand{\rhmcsAllQAfRej}{}\renewcommand{\rhmcsAllQAfRej}{XGBoost on the full predictor set}  % the alternatives StepM rejects at 5 % (none: none)
\providecommand{\rhmcsAllQAfNRejTwo}{}\renewcommand{\rhmcsAllQAfNRejTwo}{0}  % StepM rejections at FWER 5 %, two-sided
\providecommand{\rhmcsAllQAfNRejTwoWord}{}\renewcommand{\rhmcsAllQAfNRejTwoWord}{none}  % ... as a word
\providecommand{\rhmcsAllQAfNRawTwo}{}\renewcommand{\rhmcsAllQAfNRawTwo}{0}  % single-test bootstrap p < 0.05, two-sided
\providecommand{\rhmcsAllQAfRejTwo}{}\renewcommand{\rhmcsAllQAfRejTwo}{none}  % the alternatives the two-sided StepM rejects at 5 % (none: none)
\providecommand{\rhmcsAllQAfPRidge}{}\renewcommand{\rhmcsAllQAfPRidge}{0.78}  % StepM-adjusted p-value (one-sided), Ridge
\providecommand{\rhmcsAllQAfRawRidge}{}\renewcommand{\rhmcsAllQAfRawRidge}{0.76}  % single-test bootstrap p-value (one-sided), Ridge
\providecommand{\rhmcsAllQAfPTwoRidge}{}\renewcommand{\rhmcsAllQAfPTwoRidge}{0.55}  % StepM-adjusted p-value (two-sided), Ridge
\providecommand{\rhmcsAllQAfRawTwoRidge}{}\renewcommand{\rhmcsAllQAfRawTwoRidge}{0.50}  % single-test bootstrap p-value (two-sided), Ridge
\providecommand{\rhmcsAllQAfPLasso}{}\renewcommand{\rhmcsAllQAfPLasso}{0.90}  % StepM-adjusted p-value (one-sided), Lasso
\providecommand{\rhmcsAllQAfRawLasso}{}\renewcommand{\rhmcsAllQAfRawLasso}{0.90}  % single-test bootstrap p-value (one-sided), Lasso
\providecommand{\rhmcsAllQAfPTwoLasso}{}\renewcommand{\rhmcsAllQAfPTwoLasso}{0.32}  % StepM-adjusted p-value (two-sided), Lasso
\providecommand{\rhmcsAllQAfRawTwoLasso}{}\renewcommand{\rhmcsAllQAfRawTwoLasso}{0.18}  % single-test bootstrap p-value (two-sided), Lasso
\providecommand{\rhmcsAllQAfPEnet}{}\renewcommand{\rhmcsAllQAfPEnet}{0.67}  % StepM-adjusted p-value (one-sided), Elastic net
\providecommand{\rhmcsAllQAfRawEnet}{}\renewcommand{\rhmcsAllQAfRawEnet}{0.62}  % single-test bootstrap p-value (one-sided), Elastic net
\providecommand{\rhmcsAllQAfPTwoEnet}{}\renewcommand{\rhmcsAllQAfPTwoEnet}{0.81}  % StepM-adjusted p-value (two-sided), Elastic net
\providecommand{\rhmcsAllQAfRawTwoEnet}{}\renewcommand{\rhmcsAllQAfRawTwoEnet}{0.81}  % single-test bootstrap p-value (two-sided), Elastic net
\providecommand{\rhmcsAllQAfPLgbm}{}\renewcommand{\rhmcsAllQAfPLgbm}{0.081}  % StepM-adjusted p-value (one-sided), LightGBM
\providecommand{\rhmcsAllQAfRawLgbm}{}\renewcommand{\rhmcsAllQAfRawLgbm}{0.037}  % single-test bootstrap p-value (one-sided), LightGBM
\providecommand{\rhmcsAllQAfPTwoLgbm}{}\renewcommand{\rhmcsAllQAfPTwoLgbm}{0.18}  % StepM-adjusted p-value (two-sided), LightGBM
\providecommand{\rhmcsAllQAfRawTwoLgbm}{}\renewcommand{\rhmcsAllQAfRawTwoLgbm}{0.081}  % single-test bootstrap p-value (two-sided), LightGBM
\providecommand{\rhmcsAllQAfPXgb}{}\renewcommand{\rhmcsAllQAfPXgb}{0.045}  % StepM-adjusted p-value (one-sided), XGBoost
\providecommand{\rhmcsAllQAfRawXgb}{}\renewcommand{\rhmcsAllQAfRawXgb}{0.020}  % single-test bootstrap p-value (one-sided), XGBoost
\providecommand{\rhmcsAllQAfPTwoXgb}{}\renewcommand{\rhmcsAllQAfPTwoXgb}{0.11}  % StepM-adjusted p-value (two-sided), XGBoost
\providecommand{\rhmcsAllQAfRawTwoXgb}{}\renewcommand{\rhmcsAllQAfRawTwoXgb}{0.050}  % single-test bootstrap p-value (two-sided), XGBoost
\providecommand{\rhmcsAllQAfPRf}{}\renewcommand{\rhmcsAllQAfPRf}{0.26}  % StepM-adjusted p-value (one-sided), Random forest
\providecommand{\rhmcsAllQAfRawRf}{}\renewcommand{\rhmcsAllQAfRawRf}{0.15}  % single-test bootstrap p-value (one-sided), Random forest
\providecommand{\rhmcsAllQAfPTwoRf}{}\renewcommand{\rhmcsAllQAfPTwoRf}{0.48}  % StepM-adjusted p-value (two-sided), Random forest
\providecommand{\rhmcsAllQAfRawTwoRf}{}\renewcommand{\rhmcsAllQAfRawTwoRf}{0.29}  % single-test bootstrap p-value (two-sided), Random forest
\providecommand{\rhmcsAllQAfNRejLo}{}\renewcommand{\rhmcsAllQAfNRejLo}{1}  % StepM rejections at 5 %, block length 5
\providecommand{\rhmcsAllQAfSpaCLo}{}\renewcommand{\rhmcsAllQAfSpaCLo}{0.036}  % SPA p-value (consistent), block length 5
\providecommand{\rhmcsAllQAfNRejTwoLo}{}\renewcommand{\rhmcsAllQAfNRejTwoLo}{0}  % ... two-sided, block length 5
\providecommand{\rhmcsAllQAfPMinLo}{}\renewcommand{\rhmcsAllQAfPMinLo}{0.036}  % smallest adjusted p, block length 5
\providecommand{\rhmcsAllQAfPTwoMinLo}{}\renewcommand{\rhmcsAllQAfPTwoMinLo}{0.091}  % ... two-sided, block length 5
\providecommand{\rhmcsAllQAfNRejHi}{}\renewcommand{\rhmcsAllQAfNRejHi}{0}  % StepM rejections at 5 %, block length 20
\providecommand{\rhmcsAllQAfSpaCHi}{}\renewcommand{\rhmcsAllQAfSpaCHi}{0.051}  % SPA p-value (consistent), block length 20
\providecommand{\rhmcsAllQAfNRejTwoHi}{}\renewcommand{\rhmcsAllQAfNRejTwoHi}{0}  % ... two-sided, block length 20
\providecommand{\rhmcsAllQAfPMinHi}{}\renewcommand{\rhmcsAllQAfPMinHi}{0.051}  % smallest adjusted p, block length 20
\providecommand{\rhmcsAllQAfPTwoMinHi}{}\renewcommand{\rhmcsAllQAfPTwoMinHi}{0.12}  % ... two-sided, block length 20
\providecommand{\rhmcsAllMAfSpaC}{}\renewcommand{\rhmcsAllMAfSpaC}{0.66}  % SPA p-value (consistent), family AllMAf
\providecommand{\rhmcsAllMAfSpaL}{}\renewcommand{\rhmcsAllMAfSpaL}{0.20}  % ... lower
\providecommand{\rhmcsAllMAfSpaU}{}\renewcommand{\rhmcsAllMAfSpaU}{0.66}  % ... upper
\providecommand{\rhmcsAllMAfNRej}{}\renewcommand{\rhmcsAllMAfNRej}{0}  % StepM rejections at FWER 5 % (of 6)
\providecommand{\rhmcsAllMAfNRejWord}{}\renewcommand{\rhmcsAllMAfNRejWord}{none}  % ... as a word
\providecommand{\rhmcsAllMAfNRaw}{}\renewcommand{\rhmcsAllMAfNRaw}{0}  % single-test bootstrap p < 0.05 (no adjustment)
\providecommand{\rhmcsAllMAfNRawWord}{}\renewcommand{\rhmcsAllMAfNRawWord}{none}  % ... as a word
\providecommand{\rhmcsAllMAfNPos}{}\renewcommand{\rhmcsAllMAfNPos}{0}  % alternatives with a positive estimate (better than the benchmark)
\providecommand{\rhmcsAllMAfPMin}{}\renewcommand{\rhmcsAllMAfPMin}{0.91}  % smallest StepM-adjusted p-value
\providecommand{\rhmcsAllMAfPMax}{}\renewcommand{\rhmcsAllMAfPMax}{0.91}  % largest StepM-adjusted p-value
\providecommand{\rhmcsAllMAfPTwoMin}{}\renewcommand{\rhmcsAllMAfPTwoMin}{0.36}  % smallest two-sided StepM-adjusted p-value
\providecommand{\rhmcsAllMAfPTwoMax}{}\renewcommand{\rhmcsAllMAfPTwoMax}{0.37}  % largest two-sided StepM-adjusted p-value
\providecommand{\rhmcsAllMAfRawMin}{}\renewcommand{\rhmcsAllMAfRawMin}{0.86}  % smallest single-test bootstrap p-value
\providecommand{\rhmcsAllMAfRej}{}\renewcommand{\rhmcsAllMAfRej}{none}  % the alternatives StepM rejects at 5 % (none: none)
\providecommand{\rhmcsAllMAfNRejTwo}{}\renewcommand{\rhmcsAllMAfNRejTwo}{0}  % StepM rejections at FWER 5 %, two-sided
\providecommand{\rhmcsAllMAfNRejTwoWord}{}\renewcommand{\rhmcsAllMAfNRejTwoWord}{none}  % ... as a word
\providecommand{\rhmcsAllMAfNRawTwo}{}\renewcommand{\rhmcsAllMAfNRawTwo}{0}  % single-test bootstrap p < 0.05, two-sided
\providecommand{\rhmcsAllMAfRejTwo}{}\renewcommand{\rhmcsAllMAfRejTwo}{none}  % the alternatives the two-sided StepM rejects at 5 % (none: none)
\providecommand{\rhmcsAllMAfPRidge}{}\renewcommand{\rhmcsAllMAfPRidge}{0.91}  % StepM-adjusted p-value (one-sided), Ridge
\providecommand{\rhmcsAllMAfRawRidge}{}\renewcommand{\rhmcsAllMAfRawRidge}{0.87}  % single-test bootstrap p-value (one-sided), Ridge
\providecommand{\rhmcsAllMAfPTwoRidge}{}\renewcommand{\rhmcsAllMAfPTwoRidge}{0.36}  % StepM-adjusted p-value (two-sided), Ridge
\providecommand{\rhmcsAllMAfRawTwoRidge}{}\renewcommand{\rhmcsAllMAfRawTwoRidge}{0.26}  % single-test bootstrap p-value (two-sided), Ridge
\providecommand{\rhmcsAllMAfPLasso}{}\renewcommand{\rhmcsAllMAfPLasso}{0.91}  % StepM-adjusted p-value (one-sided), Lasso
\providecommand{\rhmcsAllMAfRawLasso}{}\renewcommand{\rhmcsAllMAfRawLasso}{0.87}  % single-test bootstrap p-value (one-sided), Lasso
\providecommand{\rhmcsAllMAfPTwoLasso}{}\renewcommand{\rhmcsAllMAfPTwoLasso}{0.36}  % StepM-adjusted p-value (two-sided), Lasso
\providecommand{\rhmcsAllMAfRawTwoLasso}{}\renewcommand{\rhmcsAllMAfRawTwoLasso}{0.26}  % single-test bootstrap p-value (two-sided), Lasso
\providecommand{\rhmcsAllMAfPEnet}{}\renewcommand{\rhmcsAllMAfPEnet}{0.91}  % StepM-adjusted p-value (one-sided), Elastic net
\providecommand{\rhmcsAllMAfRawEnet}{}\renewcommand{\rhmcsAllMAfRawEnet}{0.87}  % single-test bootstrap p-value (one-sided), Elastic net
\providecommand{\rhmcsAllMAfPTwoEnet}{}\renewcommand{\rhmcsAllMAfPTwoEnet}{0.37}  % StepM-adjusted p-value (two-sided), Elastic net
\providecommand{\rhmcsAllMAfRawTwoEnet}{}\renewcommand{\rhmcsAllMAfRawTwoEnet}{0.26}  % single-test bootstrap p-value (two-sided), Elastic net
\providecommand{\rhmcsAllMAfPLgbm}{}\renewcommand{\rhmcsAllMAfPLgbm}{0.91}  % StepM-adjusted p-value (one-sided), LightGBM
\providecommand{\rhmcsAllMAfRawLgbm}{}\renewcommand{\rhmcsAllMAfRawLgbm}{0.86}  % single-test bootstrap p-value (one-sided), LightGBM
\providecommand{\rhmcsAllMAfPTwoLgbm}{}\renewcommand{\rhmcsAllMAfPTwoLgbm}{0.37}  % StepM-adjusted p-value (two-sided), LightGBM
\providecommand{\rhmcsAllMAfRawTwoLgbm}{}\renewcommand{\rhmcsAllMAfRawTwoLgbm}{0.33}  % single-test bootstrap p-value (two-sided), LightGBM
\providecommand{\rhmcsAllMAfPXgb}{}\renewcommand{\rhmcsAllMAfPXgb}{0.91}  % StepM-adjusted p-value (one-sided), XGBoost
\providecommand{\rhmcsAllMAfRawXgb}{}\renewcommand{\rhmcsAllMAfRawXgb}{0.86}  % single-test bootstrap p-value (one-sided), XGBoost
\providecommand{\rhmcsAllMAfPTwoXgb}{}\renewcommand{\rhmcsAllMAfPTwoXgb}{0.37}  % StepM-adjusted p-value (two-sided), XGBoost
\providecommand{\rhmcsAllMAfRawTwoXgb}{}\renewcommand{\rhmcsAllMAfRawTwoXgb}{0.32}  % single-test bootstrap p-value (two-sided), XGBoost
\providecommand{\rhmcsAllMAfPRf}{}\renewcommand{\rhmcsAllMAfPRf}{0.91}  % StepM-adjusted p-value (one-sided), Random forest
\providecommand{\rhmcsAllMAfRawRf}{}\renewcommand{\rhmcsAllMAfRawRf}{0.86}  % single-test bootstrap p-value (one-sided), Random forest
\providecommand{\rhmcsAllMAfPTwoRf}{}\renewcommand{\rhmcsAllMAfPTwoRf}{0.37}  % StepM-adjusted p-value (two-sided), Random forest
\providecommand{\rhmcsAllMAfRawTwoRf}{}\renewcommand{\rhmcsAllMAfRawTwoRf}{0.30}  % single-test bootstrap p-value (two-sided), Random forest
\providecommand{\rhmcsAllMAfNRejLo}{}\renewcommand{\rhmcsAllMAfNRejLo}{0}  % StepM rejections at 5 %, block length 5
\providecommand{\rhmcsAllMAfSpaCLo}{}\renewcommand{\rhmcsAllMAfSpaCLo}{0.73}  % SPA p-value (consistent), block length 5
\providecommand{\rhmcsAllMAfNRejTwoLo}{}\renewcommand{\rhmcsAllMAfNRejTwoLo}{0}  % ... two-sided, block length 5
\providecommand{\rhmcsAllMAfPMinLo}{}\renewcommand{\rhmcsAllMAfPMinLo}{0.97}  % smallest adjusted p, block length 5
\providecommand{\rhmcsAllMAfPTwoMinLo}{}\renewcommand{\rhmcsAllMAfPTwoMinLo}{0.24}  % ... two-sided, block length 5
\providecommand{\rhmcsAllMAfNRejHi}{}\renewcommand{\rhmcsAllMAfNRejHi}{0}  % StepM rejections at 5 %, block length 20
\providecommand{\rhmcsAllMAfSpaCHi}{}\renewcommand{\rhmcsAllMAfSpaCHi}{0.62}  % SPA p-value (consistent), block length 20
\providecommand{\rhmcsAllMAfNRejTwoHi}{}\renewcommand{\rhmcsAllMAfNRejTwoHi}{0}  % ... two-sided, block length 20
\providecommand{\rhmcsAllMAfPMinHi}{}\renewcommand{\rhmcsAllMAfPMinHi}{0.87}  % smallest adjusted p, block length 20
\providecommand{\rhmcsAllMAfPTwoMinHi}{}\renewcommand{\rhmcsAllMAfPTwoMinHi}{0.45}  % ... two-sided, block length 20
\providecommand{\rhmcsOosQAfSpaC}{}\renewcommand{\rhmcsOosQAfSpaC}{0.34}  % SPA p-value (consistent), family OosQAf
\providecommand{\rhmcsOosQAfSpaL}{}\renewcommand{\rhmcsOosQAfSpaL}{0.34}  % ... lower
\providecommand{\rhmcsOosQAfSpaU}{}\renewcommand{\rhmcsOosQAfSpaU}{0.48}  % ... upper
\providecommand{\rhmcsOosQAfNRej}{}\renewcommand{\rhmcsOosQAfNRej}{0}  % StepM rejections at FWER 5 % (of 6)
\providecommand{\rhmcsOosQAfNRejWord}{}\renewcommand{\rhmcsOosQAfNRejWord}{none}  % ... as a word
\providecommand{\rhmcsOosQAfNRaw}{}\renewcommand{\rhmcsOosQAfNRaw}{0}  % single-test bootstrap p < 0.05 (no adjustment)
\providecommand{\rhmcsOosQAfNRawWord}{}\renewcommand{\rhmcsOosQAfNRawWord}{none}  % ... as a word
\providecommand{\rhmcsOosQAfNPos}{}\renewcommand{\rhmcsOosQAfNPos}{3}  % alternatives with a positive estimate (better than the benchmark)
\providecommand{\rhmcsOosQAfPMin}{}\renewcommand{\rhmcsOosQAfPMin}{0.48}  % smallest StepM-adjusted p-value
\providecommand{\rhmcsOosQAfPMax}{}\renewcommand{\rhmcsOosQAfPMax}{0.99}  % largest StepM-adjusted p-value
\providecommand{\rhmcsOosQAfPTwoMin}{}\renewcommand{\rhmcsOosQAfPTwoMin}{0.032}  % smallest two-sided StepM-adjusted p-value
\providecommand{\rhmcsOosQAfPTwoMax}{}\renewcommand{\rhmcsOosQAfPTwoMax}{0.78}  % largest two-sided StepM-adjusted p-value
\providecommand{\rhmcsOosQAfRawMin}{}\renewcommand{\rhmcsOosQAfRawMin}{0.26}  % smallest single-test bootstrap p-value
\providecommand{\rhmcsOosQAfRej}{}\renewcommand{\rhmcsOosQAfRej}{none}  % the alternatives StepM rejects at 5 % (none: none)
\providecommand{\rhmcsOosQAfNRejTwo}{}\renewcommand{\rhmcsOosQAfNRejTwo}{1}  % StepM rejections at FWER 5 %, two-sided
\providecommand{\rhmcsOosQAfNRejTwoWord}{}\renewcommand{\rhmcsOosQAfNRejTwoWord}{one}  % ... as a word
\providecommand{\rhmcsOosQAfNRawTwo}{}\renewcommand{\rhmcsOosQAfNRawTwo}{3}  % single-test bootstrap p < 0.05, two-sided
\providecommand{\rhmcsOosQAfRejTwo}{}\renewcommand{\rhmcsOosQAfRejTwo}{the lasso on the full predictor set}  % the alternatives the two-sided StepM rejects at 5 % (none: none)
\providecommand{\rhmcsOosQAfPRidge}{}\renewcommand{\rhmcsOosQAfPRidge}{0.99}  % StepM-adjusted p-value (one-sided), Ridge
\providecommand{\rhmcsOosQAfRawRidge}{}\renewcommand{\rhmcsOosQAfRawRidge}{0.99}  % single-test bootstrap p-value (one-sided), Ridge
\providecommand{\rhmcsOosQAfPTwoRidge}{}\renewcommand{\rhmcsOosQAfPTwoRidge}{0.052}  % StepM-adjusted p-value (two-sided), Ridge
\providecommand{\rhmcsOosQAfRawTwoRidge}{}\renewcommand{\rhmcsOosQAfRawTwoRidge}{0.019}  % single-test bootstrap p-value (two-sided), Ridge
\providecommand{\rhmcsOosQAfPLasso}{}\renewcommand{\rhmcsOosQAfPLasso}{0.99}  % StepM-adjusted p-value (one-sided), Lasso
\providecommand{\rhmcsOosQAfRawLasso}{}\renewcommand{\rhmcsOosQAfRawLasso}{0.99}  % single-test bootstrap p-value (one-sided), Lasso
\providecommand{\rhmcsOosQAfPTwoLasso}{}\renewcommand{\rhmcsOosQAfPTwoLasso}{0.032}  % StepM-adjusted p-value (two-sided), Lasso
\providecommand{\rhmcsOosQAfRawTwoLasso}{}\renewcommand{\rhmcsOosQAfRawTwoLasso}{0.012}  % single-test bootstrap p-value (two-sided), Lasso
\providecommand{\rhmcsOosQAfPEnet}{}\renewcommand{\rhmcsOosQAfPEnet}{0.98}  % StepM-adjusted p-value (one-sided), Elastic net
\providecommand{\rhmcsOosQAfRawEnet}{}\renewcommand{\rhmcsOosQAfRawEnet}{0.97}  % single-test bootstrap p-value (one-sided), Elastic net
\providecommand{\rhmcsOosQAfPTwoEnet}{}\renewcommand{\rhmcsOosQAfPTwoEnet}{0.11}  % StepM-adjusted p-value (two-sided), Elastic net
\providecommand{\rhmcsOosQAfRawTwoEnet}{}\renewcommand{\rhmcsOosQAfRawTwoEnet}{0.049}  % single-test bootstrap p-value (two-sided), Elastic net
\providecommand{\rhmcsOosQAfPLgbm}{}\renewcommand{\rhmcsOosQAfPLgbm}{0.53}  % StepM-adjusted p-value (one-sided), LightGBM
\providecommand{\rhmcsOosQAfRawLgbm}{}\renewcommand{\rhmcsOosQAfRawLgbm}{0.31}  % single-test bootstrap p-value (one-sided), LightGBM
\providecommand{\rhmcsOosQAfPTwoLgbm}{}\renewcommand{\rhmcsOosQAfPTwoLgbm}{0.76}  % StepM-adjusted p-value (two-sided), LightGBM
\providecommand{\rhmcsOosQAfRawTwoLgbm}{}\renewcommand{\rhmcsOosQAfRawTwoLgbm}{0.62}  % single-test bootstrap p-value (two-sided), LightGBM
\providecommand{\rhmcsOosQAfPXgb}{}\renewcommand{\rhmcsOosQAfPXgb}{0.48}  % StepM-adjusted p-value (one-sided), XGBoost
\providecommand{\rhmcsOosQAfRawXgb}{}\renewcommand{\rhmcsOosQAfRawXgb}{0.26}  % single-test bootstrap p-value (one-sided), XGBoost
\providecommand{\rhmcsOosQAfPTwoXgb}{}\renewcommand{\rhmcsOosQAfPTwoXgb}{0.68}  % StepM-adjusted p-value (two-sided), XGBoost
\providecommand{\rhmcsOosQAfRawTwoXgb}{}\renewcommand{\rhmcsOosQAfRawTwoXgb}{0.53}  % single-test bootstrap p-value (two-sided), XGBoost
\providecommand{\rhmcsOosQAfPRf}{}\renewcommand{\rhmcsOosQAfPRf}{0.59}  % StepM-adjusted p-value (one-sided), Random forest
\providecommand{\rhmcsOosQAfRawRf}{}\renewcommand{\rhmcsOosQAfRawRf}{0.40}  % single-test bootstrap p-value (one-sided), Random forest
\providecommand{\rhmcsOosQAfPTwoRf}{}\renewcommand{\rhmcsOosQAfPTwoRf}{0.78}  % StepM-adjusted p-value (two-sided), Random forest
\providecommand{\rhmcsOosQAfRawTwoRf}{}\renewcommand{\rhmcsOosQAfRawTwoRf}{0.78}  % single-test bootstrap p-value (two-sided), Random forest
\providecommand{\rhmcsOosQAfNRejLo}{}\renewcommand{\rhmcsOosQAfNRejLo}{0}  % StepM rejections at 5 %, block length 5
\providecommand{\rhmcsOosQAfSpaCLo}{}\renewcommand{\rhmcsOosQAfSpaCLo}{0.34}  % SPA p-value (consistent), block length 5
\providecommand{\rhmcsOosQAfNRejTwoLo}{}\renewcommand{\rhmcsOosQAfNRejTwoLo}{2}  % ... two-sided, block length 5
\providecommand{\rhmcsOosQAfPMinLo}{}\renewcommand{\rhmcsOosQAfPMinLo}{0.47}  % smallest adjusted p, block length 5
\providecommand{\rhmcsOosQAfPTwoMinLo}{}\renewcommand{\rhmcsOosQAfPTwoMinLo}{0.017}  % ... two-sided, block length 5
\providecommand{\rhmcsOosQAfNRejHi}{}\renewcommand{\rhmcsOosQAfNRejHi}{0}  % StepM rejections at 5 %, block length 20
\providecommand{\rhmcsOosQAfSpaCHi}{}\renewcommand{\rhmcsOosQAfSpaCHi}{0.43}  % SPA p-value (consistent), block length 20
\providecommand{\rhmcsOosQAfNRejTwoHi}{}\renewcommand{\rhmcsOosQAfNRejTwoHi}{0}  % ... two-sided, block length 20
\providecommand{\rhmcsOosQAfPMinHi}{}\renewcommand{\rhmcsOosQAfPMinHi}{0.46}  % smallest adjusted p, block length 20
\providecommand{\rhmcsOosQAfPTwoMinHi}{}\renewcommand{\rhmcsOosQAfPTwoMinHi}{0.051}  % ... two-sided, block length 20
\providecommand{\rhmcsOosQBmSpaC}{}\renewcommand{\rhmcsOosQBmSpaC}{0.53}  % SPA p-value (consistent), family OosQBm
\providecommand{\rhmcsOosQBmSpaL}{}\renewcommand{\rhmcsOosQBmSpaL}{0.083}  % ... lower
\providecommand{\rhmcsOosQBmSpaU}{}\renewcommand{\rhmcsOosQBmSpaU}{0.69}  % ... upper
\providecommand{\rhmcsOosQBmNRej}{}\renewcommand{\rhmcsOosQBmNRej}{0}  % StepM rejections at FWER 5 % (of 6)
\providecommand{\rhmcsOosQBmNRejWord}{}\renewcommand{\rhmcsOosQBmNRejWord}{none}  % ... as a word
\providecommand{\rhmcsOosQBmNRaw}{}\renewcommand{\rhmcsOosQBmNRaw}{0}  % single-test bootstrap p < 0.05 (no adjustment)
\providecommand{\rhmcsOosQBmNRawWord}{}\renewcommand{\rhmcsOosQBmNRawWord}{none}  % ... as a word
\providecommand{\rhmcsOosQBmNPos}{}\renewcommand{\rhmcsOosQBmNPos}{0}  % alternatives with a positive estimate (better than the benchmark)
\providecommand{\rhmcsOosQBmPMin}{}\renewcommand{\rhmcsOosQBmPMin}{0.98}  % smallest StepM-adjusted p-value
\providecommand{\rhmcsOosQBmPMax}{}\renewcommand{\rhmcsOosQBmPMax}{1.00}  % largest StepM-adjusted p-value
\providecommand{\rhmcsOosQBmPTwoMin}{}\renewcommand{\rhmcsOosQBmPTwoMin}{0.001}  % smallest two-sided StepM-adjusted p-value
\providecommand{\rhmcsOosQBmPTwoMax}{}\renewcommand{\rhmcsOosQBmPTwoMax}{0.17}  % largest two-sided StepM-adjusted p-value
\providecommand{\rhmcsOosQBmRawMin}{}\renewcommand{\rhmcsOosQBmRawMin}{0.92}  % smallest single-test bootstrap p-value
\providecommand{\rhmcsOosQBmRej}{}\renewcommand{\rhmcsOosQBmRej}{none}  % the alternatives StepM rejects at 5 % (none: none)
\providecommand{\rhmcsOosQBmNRejTwo}{}\renewcommand{\rhmcsOosQBmNRejTwo}{4}  % StepM rejections at FWER 5 %, two-sided
\providecommand{\rhmcsOosQBmNRejTwoWord}{}\renewcommand{\rhmcsOosQBmNRejTwoWord}{four}  % ... as a word
\providecommand{\rhmcsOosQBmNRawTwo}{}\renewcommand{\rhmcsOosQBmNRawTwo}{4}  % single-test bootstrap p < 0.05, two-sided
\providecommand{\rhmcsOosQBmRejTwo}{}\renewcommand{\rhmcsOosQBmRejTwo}{ridge regression on the full predictor set, the lasso on the full predictor set, the elastic net on the full predictor set, and the random forest on the full predictor set}  % the alternatives the two-sided StepM rejects at 5 % (none: none)
\providecommand{\rhmcsOosQBmPRidge}{}\renewcommand{\rhmcsOosQBmPRidge}{1.00}  % StepM-adjusted p-value (one-sided), Ridge
\providecommand{\rhmcsOosQBmRawRidge}{}\renewcommand{\rhmcsOosQBmRawRidge}{1.00}  % single-test bootstrap p-value (one-sided), Ridge
\providecommand{\rhmcsOosQBmPTwoRidge}{}\renewcommand{\rhmcsOosQBmPTwoRidge}{0.001}  % StepM-adjusted p-value (two-sided), Ridge
\providecommand{\rhmcsOosQBmRawTwoRidge}{}\renewcommand{\rhmcsOosQBmRawTwoRidge}{\ensuremath{<}0.001}  % single-test bootstrap p-value (two-sided), Ridge
\providecommand{\rhmcsOosQBmPLasso}{}\renewcommand{\rhmcsOosQBmPLasso}{1.00}  % StepM-adjusted p-value (one-sided), Lasso
\providecommand{\rhmcsOosQBmRawLasso}{}\renewcommand{\rhmcsOosQBmRawLasso}{1.00}  % single-test bootstrap p-value (one-sided), Lasso
\providecommand{\rhmcsOosQBmPTwoLasso}{}\renewcommand{\rhmcsOosQBmPTwoLasso}{0.001}  % StepM-adjusted p-value (two-sided), Lasso
\providecommand{\rhmcsOosQBmRawTwoLasso}{}\renewcommand{\rhmcsOosQBmRawTwoLasso}{\ensuremath{<}0.001}  % single-test bootstrap p-value (two-sided), Lasso
\providecommand{\rhmcsOosQBmPEnet}{}\renewcommand{\rhmcsOosQBmPEnet}{1.00}  % StepM-adjusted p-value (one-sided), Elastic net
\providecommand{\rhmcsOosQBmRawEnet}{}\renewcommand{\rhmcsOosQBmRawEnet}{1.00}  % single-test bootstrap p-value (one-sided), Elastic net
\providecommand{\rhmcsOosQBmPTwoEnet}{}\renewcommand{\rhmcsOosQBmPTwoEnet}{0.002}  % StepM-adjusted p-value (two-sided), Elastic net
\providecommand{\rhmcsOosQBmRawTwoEnet}{}\renewcommand{\rhmcsOosQBmRawTwoEnet}{\ensuremath{<}0.001}  % single-test bootstrap p-value (two-sided), Elastic net
\providecommand{\rhmcsOosQBmPLgbm}{}\renewcommand{\rhmcsOosQBmPLgbm}{0.98}  % StepM-adjusted p-value (one-sided), LightGBM
\providecommand{\rhmcsOosQBmRawLgbm}{}\renewcommand{\rhmcsOosQBmRawLgbm}{0.92}  % single-test bootstrap p-value (one-sided), LightGBM
\providecommand{\rhmcsOosQBmPTwoLgbm}{}\renewcommand{\rhmcsOosQBmPTwoLgbm}{0.17}  % StepM-adjusted p-value (two-sided), LightGBM
\providecommand{\rhmcsOosQBmRawTwoLgbm}{}\renewcommand{\rhmcsOosQBmRawTwoLgbm}{0.15}  % single-test bootstrap p-value (two-sided), LightGBM
\providecommand{\rhmcsOosQBmPXgb}{}\renewcommand{\rhmcsOosQBmPXgb}{0.98}  % StepM-adjusted p-value (one-sided), XGBoost
\providecommand{\rhmcsOosQBmRawXgb}{}\renewcommand{\rhmcsOosQBmRawXgb}{0.92}  % single-test bootstrap p-value (one-sided), XGBoost
\providecommand{\rhmcsOosQBmPTwoXgb}{}\renewcommand{\rhmcsOosQBmPTwoXgb}{0.17}  % StepM-adjusted p-value (two-sided), XGBoost
\providecommand{\rhmcsOosQBmRawTwoXgb}{}\renewcommand{\rhmcsOosQBmRawTwoXgb}{0.16}  % single-test bootstrap p-value (two-sided), XGBoost
\providecommand{\rhmcsOosQBmPRf}{}\renewcommand{\rhmcsOosQBmPRf}{1.00}  % StepM-adjusted p-value (one-sided), Random forest
\providecommand{\rhmcsOosQBmRawRf}{}\renewcommand{\rhmcsOosQBmRawRf}{0.98}  % single-test bootstrap p-value (one-sided), Random forest
\providecommand{\rhmcsOosQBmPTwoRf}{}\renewcommand{\rhmcsOosQBmPTwoRf}{0.039}  % StepM-adjusted p-value (two-sided), Random forest
\providecommand{\rhmcsOosQBmRawTwoRf}{}\renewcommand{\rhmcsOosQBmRawTwoRf}{0.030}  % single-test bootstrap p-value (two-sided), Random forest
\providecommand{\rhmcsOosQBmNRejLo}{}\renewcommand{\rhmcsOosQBmNRejLo}{0}  % StepM rejections at 5 %, block length 5
\providecommand{\rhmcsOosQBmSpaCLo}{}\renewcommand{\rhmcsOosQBmSpaCLo}{0.53}  % SPA p-value (consistent), block length 5
\providecommand{\rhmcsOosQBmNRejTwoLo}{}\renewcommand{\rhmcsOosQBmNRejTwoLo}{4}  % ... two-sided, block length 5
\providecommand{\rhmcsOosQBmPMinLo}{}\renewcommand{\rhmcsOosQBmPMinLo}{0.97}  % smallest adjusted p, block length 5
\providecommand{\rhmcsOosQBmPTwoMinLo}{}\renewcommand{\rhmcsOosQBmPTwoMinLo}{\ensuremath{<}0.001}  % ... two-sided, block length 5
\providecommand{\rhmcsOosQBmNRejHi}{}\renewcommand{\rhmcsOosQBmNRejHi}{0}  % StepM rejections at 5 %, block length 20
\providecommand{\rhmcsOosQBmSpaCHi}{}\renewcommand{\rhmcsOosQBmSpaCHi}{0.53}  % SPA p-value (consistent), block length 20
\providecommand{\rhmcsOosQBmNRejTwoHi}{}\renewcommand{\rhmcsOosQBmNRejTwoHi}{4}  % ... two-sided, block length 20
\providecommand{\rhmcsOosQBmPMinHi}{}\renewcommand{\rhmcsOosQBmPMinHi}{0.98}  % smallest adjusted p, block length 20
\providecommand{\rhmcsOosQBmPTwoMinHi}{}\renewcommand{\rhmcsOosQBmPTwoMinHi}{0.003}  % ... two-sided, block length 20
\providecommand{\rhmcsOosQBmAllSpaC}{}\renewcommand{\rhmcsOosQBmAllSpaC}{0.53}  % SPA p-value (consistent), family OosQBmAll
\providecommand{\rhmcsOosQBmAllSpaL}{}\renewcommand{\rhmcsOosQBmAllSpaL}{0.084}  % ... lower
\providecommand{\rhmcsOosQBmAllSpaU}{}\renewcommand{\rhmcsOosQBmAllSpaU}{0.79}  % ... upper
\providecommand{\rhmcsOosQBmAllNRej}{}\renewcommand{\rhmcsOosQBmAllNRej}{0}  % StepM rejections at FWER 5 % (of 10)
\providecommand{\rhmcsOosQBmAllNRejWord}{}\renewcommand{\rhmcsOosQBmAllNRejWord}{none}  % ... as a word
\providecommand{\rhmcsOosQBmAllNRaw}{}\renewcommand{\rhmcsOosQBmAllNRaw}{0}  % single-test bootstrap p < 0.05 (no adjustment)
\providecommand{\rhmcsOosQBmAllNRawWord}{}\renewcommand{\rhmcsOosQBmAllNRawWord}{none}  % ... as a word
\providecommand{\rhmcsOosQBmAllNPos}{}\renewcommand{\rhmcsOosQBmAllNPos}{0}  % alternatives with a positive estimate (better than the benchmark)
\providecommand{\rhmcsOosQBmAllPMin}{}\renewcommand{\rhmcsOosQBmAllPMin}{0.99}  % smallest StepM-adjusted p-value
\providecommand{\rhmcsOosQBmAllPMax}{}\renewcommand{\rhmcsOosQBmAllPMax}{1.00}  % largest StepM-adjusted p-value
\providecommand{\rhmcsOosQBmAllPTwoMin}{}\renewcommand{\rhmcsOosQBmAllPTwoMin}{0.003}  % smallest two-sided StepM-adjusted p-value
\providecommand{\rhmcsOosQBmAllPTwoMax}{}\renewcommand{\rhmcsOosQBmAllPTwoMax}{0.17}  % largest two-sided StepM-adjusted p-value
\providecommand{\rhmcsOosQBmAllRawMin}{}\renewcommand{\rhmcsOosQBmAllRawMin}{0.92}  % smallest single-test bootstrap p-value
\providecommand{\rhmcsOosQBmAllRej}{}\renewcommand{\rhmcsOosQBmAllRej}{none}  % the alternatives StepM rejects at 5 % (none: none)
\providecommand{\rhmcsOosQBmAllNRejTwo}{}\renewcommand{\rhmcsOosQBmAllNRejTwo}{8}  % StepM rejections at FWER 5 %, two-sided
\providecommand{\rhmcsOosQBmAllNRejTwoWord}{}\renewcommand{\rhmcsOosQBmAllNRejTwoWord}{eight}  % ... as a word
\providecommand{\rhmcsOosQBmAllNRawTwo}{}\renewcommand{\rhmcsOosQBmAllNRawTwo}{8}  % single-test bootstrap p < 0.05, two-sided
\providecommand{\rhmcsOosQBmAllRejTwo}{}\renewcommand{\rhmcsOosQBmAllRejTwo}{the linear models on the HAR benchmark, ridge regression on the full predictor set, the lasso on the full predictor set, the elastic net on the full predictor set, LightGBM on the HAR benchmark, XGBoost on the HAR benchmark, the random forest on the HAR benchmark, and the random forest on the full predictor set}  % the alternatives the two-sided StepM rejects at 5 % (none: none)
\providecommand{\rhmcsOosQBmAllPLinHar}{}\renewcommand{\rhmcsOosQBmAllPLinHar}{1.00}  % StepM-adjusted p-value (one-sided), Linear, HAR
\providecommand{\rhmcsOosQBmAllRawLinHar}{}\renewcommand{\rhmcsOosQBmAllRawLinHar}{0.99}  % single-test bootstrap p-value (one-sided), Linear, HAR
\providecommand{\rhmcsOosQBmAllPTwoLinHar}{}\renewcommand{\rhmcsOosQBmAllPTwoLinHar}{0.014}  % StepM-adjusted p-value (two-sided), Linear, HAR
\providecommand{\rhmcsOosQBmAllRawTwoLinHar}{}\renewcommand{\rhmcsOosQBmAllRawTwoLinHar}{0.007}  % single-test bootstrap p-value (two-sided), Linear, HAR
\providecommand{\rhmcsOosQBmAllPRidgeFull}{}\renewcommand{\rhmcsOosQBmAllPRidgeFull}{1.00}  % StepM-adjusted p-value (one-sided), Ridge, Full
\providecommand{\rhmcsOosQBmAllRawRidgeFull}{}\renewcommand{\rhmcsOosQBmAllRawRidgeFull}{1.00}  % single-test bootstrap p-value (one-sided), Ridge, Full
\providecommand{\rhmcsOosQBmAllPTwoRidgeFull}{}\renewcommand{\rhmcsOosQBmAllPTwoRidgeFull}{0.003}  % StepM-adjusted p-value (two-sided), Ridge, Full
\providecommand{\rhmcsOosQBmAllRawTwoRidgeFull}{}\renewcommand{\rhmcsOosQBmAllRawTwoRidgeFull}{\ensuremath{<}0.001}  % single-test bootstrap p-value (two-sided), Ridge, Full
\providecommand{\rhmcsOosQBmAllPLassoFull}{}\renewcommand{\rhmcsOosQBmAllPLassoFull}{1.00}  % StepM-adjusted p-value (one-sided), Lasso, Full
\providecommand{\rhmcsOosQBmAllRawLassoFull}{}\renewcommand{\rhmcsOosQBmAllRawLassoFull}{1.00}  % single-test bootstrap p-value (one-sided), Lasso, Full
\providecommand{\rhmcsOosQBmAllPTwoLassoFull}{}\renewcommand{\rhmcsOosQBmAllPTwoLassoFull}{0.003}  % StepM-adjusted p-value (two-sided), Lasso, Full
\providecommand{\rhmcsOosQBmAllRawTwoLassoFull}{}\renewcommand{\rhmcsOosQBmAllRawTwoLassoFull}{\ensuremath{<}0.001}  % single-test bootstrap p-value (two-sided), Lasso, Full
\providecommand{\rhmcsOosQBmAllPEnetFull}{}\renewcommand{\rhmcsOosQBmAllPEnetFull}{1.00}  % StepM-adjusted p-value (one-sided), Elastic net, Full
\providecommand{\rhmcsOosQBmAllRawEnetFull}{}\renewcommand{\rhmcsOosQBmAllRawEnetFull}{1.00}  % single-test bootstrap p-value (one-sided), Elastic net, Full
\providecommand{\rhmcsOosQBmAllPTwoEnetFull}{}\renewcommand{\rhmcsOosQBmAllPTwoEnetFull}{0.004}  % StepM-adjusted p-value (two-sided), Elastic net, Full
\providecommand{\rhmcsOosQBmAllRawTwoEnetFull}{}\renewcommand{\rhmcsOosQBmAllRawTwoEnetFull}{\ensuremath{<}0.001}  % single-test bootstrap p-value (two-sided), Elastic net, Full
\providecommand{\rhmcsOosQBmAllPLgbmHar}{}\renewcommand{\rhmcsOosQBmAllPLgbmHar}{1.00}  % StepM-adjusted p-value (one-sided), LightGBM, HAR
\providecommand{\rhmcsOosQBmAllRawLgbmHar}{}\renewcommand{\rhmcsOosQBmAllRawLgbmHar}{0.99}  % single-test bootstrap p-value (one-sided), LightGBM, HAR
\providecommand{\rhmcsOosQBmAllPTwoLgbmHar}{}\renewcommand{\rhmcsOosQBmAllPTwoLgbmHar}{0.021}  % StepM-adjusted p-value (two-sided), LightGBM, HAR
\providecommand{\rhmcsOosQBmAllRawTwoLgbmHar}{}\renewcommand{\rhmcsOosQBmAllRawTwoLgbmHar}{0.009}  % single-test bootstrap p-value (two-sided), LightGBM, HAR
\providecommand{\rhmcsOosQBmAllPLgbmFull}{}\renewcommand{\rhmcsOosQBmAllPLgbmFull}{0.99}  % StepM-adjusted p-value (one-sided), LightGBM, Full
\providecommand{\rhmcsOosQBmAllRawLgbmFull}{}\renewcommand{\rhmcsOosQBmAllRawLgbmFull}{0.92}  % single-test bootstrap p-value (one-sided), LightGBM, Full
\providecommand{\rhmcsOosQBmAllPTwoLgbmFull}{}\renewcommand{\rhmcsOosQBmAllPTwoLgbmFull}{0.17}  % StepM-adjusted p-value (two-sided), LightGBM, Full
\providecommand{\rhmcsOosQBmAllRawTwoLgbmFull}{}\renewcommand{\rhmcsOosQBmAllRawTwoLgbmFull}{0.15}  % single-test bootstrap p-value (two-sided), LightGBM, Full
\providecommand{\rhmcsOosQBmAllPXgbHar}{}\renewcommand{\rhmcsOosQBmAllPXgbHar}{1.00}  % StepM-adjusted p-value (one-sided), XGBoost, HAR
\providecommand{\rhmcsOosQBmAllRawXgbHar}{}\renewcommand{\rhmcsOosQBmAllRawXgbHar}{0.99}  % single-test bootstrap p-value (one-sided), XGBoost, HAR
\providecommand{\rhmcsOosQBmAllPTwoXgbHar}{}\renewcommand{\rhmcsOosQBmAllPTwoXgbHar}{0.016}  % StepM-adjusted p-value (two-sided), XGBoost, HAR
\providecommand{\rhmcsOosQBmAllRawTwoXgbHar}{}\renewcommand{\rhmcsOosQBmAllRawTwoXgbHar}{0.006}  % single-test bootstrap p-value (two-sided), XGBoost, HAR
\providecommand{\rhmcsOosQBmAllPXgbFull}{}\renewcommand{\rhmcsOosQBmAllPXgbFull}{0.99}  % StepM-adjusted p-value (one-sided), XGBoost, Full
\providecommand{\rhmcsOosQBmAllRawXgbFull}{}\renewcommand{\rhmcsOosQBmAllRawXgbFull}{0.92}  % single-test bootstrap p-value (one-sided), XGBoost, Full
\providecommand{\rhmcsOosQBmAllPTwoXgbFull}{}\renewcommand{\rhmcsOosQBmAllPTwoXgbFull}{0.17}  % StepM-adjusted p-value (two-sided), XGBoost, Full
\providecommand{\rhmcsOosQBmAllRawTwoXgbFull}{}\renewcommand{\rhmcsOosQBmAllRawTwoXgbFull}{0.16}  % single-test bootstrap p-value (two-sided), XGBoost, Full
\providecommand{\rhmcsOosQBmAllPRfHar}{}\renewcommand{\rhmcsOosQBmAllPRfHar}{1.00}  % StepM-adjusted p-value (one-sided), Random forest, HAR
\providecommand{\rhmcsOosQBmAllRawRfHar}{}\renewcommand{\rhmcsOosQBmAllRawRfHar}{1.00}  % single-test bootstrap p-value (one-sided), Random forest, HAR
\providecommand{\rhmcsOosQBmAllPTwoRfHar}{}\renewcommand{\rhmcsOosQBmAllPTwoRfHar}{0.006}  % StepM-adjusted p-value (two-sided), Random forest, HAR
\providecommand{\rhmcsOosQBmAllRawTwoRfHar}{}\renewcommand{\rhmcsOosQBmAllRawTwoRfHar}{0.002}  % single-test bootstrap p-value (two-sided), Random forest, HAR
\providecommand{\rhmcsOosQBmAllPRfFull}{}\renewcommand{\rhmcsOosQBmAllPRfFull}{1.00}  % StepM-adjusted p-value (one-sided), Random forest, Full
\providecommand{\rhmcsOosQBmAllRawRfFull}{}\renewcommand{\rhmcsOosQBmAllRawRfFull}{0.98}  % single-test bootstrap p-value (one-sided), Random forest, Full
\providecommand{\rhmcsOosQBmAllPTwoRfFull}{}\renewcommand{\rhmcsOosQBmAllPTwoRfFull}{0.039}  % StepM-adjusted p-value (two-sided), Random forest, Full
\providecommand{\rhmcsOosQBmAllRawTwoRfFull}{}\renewcommand{\rhmcsOosQBmAllRawTwoRfFull}{0.030}  % single-test bootstrap p-value (two-sided), Random forest, Full
\providecommand{\rhmcsOosQBmAllNRejLo}{}\renewcommand{\rhmcsOosQBmAllNRejLo}{0}  % StepM rejections at 5 %, block length 5
\providecommand{\rhmcsOosQBmAllSpaCLo}{}\renewcommand{\rhmcsOosQBmAllSpaCLo}{0.53}  % SPA p-value (consistent), block length 5
\providecommand{\rhmcsOosQBmAllNRejTwoLo}{}\renewcommand{\rhmcsOosQBmAllNRejTwoLo}{8}  % ... two-sided, block length 5
\providecommand{\rhmcsOosQBmAllPMinLo}{}\renewcommand{\rhmcsOosQBmAllPMinLo}{0.99}  % smallest adjusted p, block length 5
\providecommand{\rhmcsOosQBmAllPTwoMinLo}{}\renewcommand{\rhmcsOosQBmAllPTwoMinLo}{0.002}  % ... two-sided, block length 5
\providecommand{\rhmcsOosQBmAllNRejHi}{}\renewcommand{\rhmcsOosQBmAllNRejHi}{0}  % StepM rejections at 5 %, block length 20
\providecommand{\rhmcsOosQBmAllSpaCHi}{}\renewcommand{\rhmcsOosQBmAllSpaCHi}{0.53}  % SPA p-value (consistent), block length 20
\providecommand{\rhmcsOosQBmAllNRejTwoHi}{}\renewcommand{\rhmcsOosQBmAllNRejTwoHi}{8}  % ... two-sided, block length 20
\providecommand{\rhmcsOosQBmAllPMinHi}{}\renewcommand{\rhmcsOosQBmAllPMinHi}{0.99}  % smallest adjusted p, block length 20
\providecommand{\rhmcsOosQBmAllPTwoMinHi}{}\renewcommand{\rhmcsOosQBmAllPTwoMinHi}{0.005}  % ... two-sided, block length 20
\providecommand{\rhmcsOosMAfSpaC}{}\renewcommand{\rhmcsOosMAfSpaC}{0.67}  % SPA p-value (consistent), family OosMAf
\providecommand{\rhmcsOosMAfSpaL}{}\renewcommand{\rhmcsOosMAfSpaL}{0.19}  % ... lower
\providecommand{\rhmcsOosMAfSpaU}{}\renewcommand{\rhmcsOosMAfSpaU}{0.67}  % ... upper
\providecommand{\rhmcsOosMAfNRej}{}\renewcommand{\rhmcsOosMAfNRej}{0}  % StepM rejections at FWER 5 % (of 6)
\providecommand{\rhmcsOosMAfNRejWord}{}\renewcommand{\rhmcsOosMAfNRejWord}{none}  % ... as a word
\providecommand{\rhmcsOosMAfNRaw}{}\renewcommand{\rhmcsOosMAfNRaw}{0}  % single-test bootstrap p < 0.05 (no adjustment)
\providecommand{\rhmcsOosMAfNRawWord}{}\renewcommand{\rhmcsOosMAfNRawWord}{none}  % ... as a word
\providecommand{\rhmcsOosMAfNPos}{}\renewcommand{\rhmcsOosMAfNPos}{0}  % alternatives with a positive estimate (better than the benchmark)
\providecommand{\rhmcsOosMAfPMin}{}\renewcommand{\rhmcsOosMAfPMin}{0.92}  % smallest StepM-adjusted p-value
\providecommand{\rhmcsOosMAfPMax}{}\renewcommand{\rhmcsOosMAfPMax}{0.92}  % largest StepM-adjusted p-value
\providecommand{\rhmcsOosMAfPTwoMin}{}\renewcommand{\rhmcsOosMAfPTwoMin}{0.32}  % smallest two-sided StepM-adjusted p-value
\providecommand{\rhmcsOosMAfPTwoMax}{}\renewcommand{\rhmcsOosMAfPTwoMax}{0.34}  % largest two-sided StepM-adjusted p-value
\providecommand{\rhmcsOosMAfRawMin}{}\renewcommand{\rhmcsOosMAfRawMin}{0.86}  % smallest single-test bootstrap p-value
\providecommand{\rhmcsOosMAfRej}{}\renewcommand{\rhmcsOosMAfRej}{none}  % the alternatives StepM rejects at 5 % (none: none)
\providecommand{\rhmcsOosMAfNRejTwo}{}\renewcommand{\rhmcsOosMAfNRejTwo}{0}  % StepM rejections at FWER 5 %, two-sided
\providecommand{\rhmcsOosMAfNRejTwoWord}{}\renewcommand{\rhmcsOosMAfNRejTwoWord}{none}  % ... as a word
\providecommand{\rhmcsOosMAfNRawTwo}{}\renewcommand{\rhmcsOosMAfNRawTwo}{0}  % single-test bootstrap p < 0.05, two-sided
\providecommand{\rhmcsOosMAfRejTwo}{}\renewcommand{\rhmcsOosMAfRejTwo}{none}  % the alternatives the two-sided StepM rejects at 5 % (none: none)
\providecommand{\rhmcsOosMAfPRidge}{}\renewcommand{\rhmcsOosMAfPRidge}{0.92}  % StepM-adjusted p-value (one-sided), Ridge
\providecommand{\rhmcsOosMAfRawRidge}{}\renewcommand{\rhmcsOosMAfRawRidge}{0.88}  % single-test bootstrap p-value (one-sided), Ridge
\providecommand{\rhmcsOosMAfPTwoRidge}{}\renewcommand{\rhmcsOosMAfPTwoRidge}{0.33}  % StepM-adjusted p-value (two-sided), Ridge
\providecommand{\rhmcsOosMAfRawTwoRidge}{}\renewcommand{\rhmcsOosMAfRawTwoRidge}{0.23}  % single-test bootstrap p-value (two-sided), Ridge
\providecommand{\rhmcsOosMAfPLasso}{}\renewcommand{\rhmcsOosMAfPLasso}{0.92}  % StepM-adjusted p-value (one-sided), Lasso
\providecommand{\rhmcsOosMAfRawLasso}{}\renewcommand{\rhmcsOosMAfRawLasso}{0.88}  % single-test bootstrap p-value (one-sided), Lasso
\providecommand{\rhmcsOosMAfPTwoLasso}{}\renewcommand{\rhmcsOosMAfPTwoLasso}{0.32}  % StepM-adjusted p-value (two-sided), Lasso
\providecommand{\rhmcsOosMAfRawTwoLasso}{}\renewcommand{\rhmcsOosMAfRawTwoLasso}{0.22}  % single-test bootstrap p-value (two-sided), Lasso
\providecommand{\rhmcsOosMAfPEnet}{}\renewcommand{\rhmcsOosMAfPEnet}{0.92}  % StepM-adjusted p-value (one-sided), Elastic net
\providecommand{\rhmcsOosMAfRawEnet}{}\renewcommand{\rhmcsOosMAfRawEnet}{0.88}  % single-test bootstrap p-value (one-sided), Elastic net
\providecommand{\rhmcsOosMAfPTwoEnet}{}\renewcommand{\rhmcsOosMAfPTwoEnet}{0.34}  % StepM-adjusted p-value (two-sided), Elastic net
\providecommand{\rhmcsOosMAfRawTwoEnet}{}\renewcommand{\rhmcsOosMAfRawTwoEnet}{0.23}  % single-test bootstrap p-value (two-sided), Elastic net
\providecommand{\rhmcsOosMAfPLgbm}{}\renewcommand{\rhmcsOosMAfPLgbm}{0.92}  % StepM-adjusted p-value (one-sided), LightGBM
\providecommand{\rhmcsOosMAfRawLgbm}{}\renewcommand{\rhmcsOosMAfRawLgbm}{0.86}  % single-test bootstrap p-value (one-sided), LightGBM
\providecommand{\rhmcsOosMAfPTwoLgbm}{}\renewcommand{\rhmcsOosMAfPTwoLgbm}{0.34}  % StepM-adjusted p-value (two-sided), LightGBM
\providecommand{\rhmcsOosMAfRawTwoLgbm}{}\renewcommand{\rhmcsOosMAfRawTwoLgbm}{0.31}  % single-test bootstrap p-value (two-sided), LightGBM
\providecommand{\rhmcsOosMAfPXgb}{}\renewcommand{\rhmcsOosMAfPXgb}{0.92}  % StepM-adjusted p-value (one-sided), XGBoost
\providecommand{\rhmcsOosMAfRawXgb}{}\renewcommand{\rhmcsOosMAfRawXgb}{0.86}  % single-test bootstrap p-value (one-sided), XGBoost
\providecommand{\rhmcsOosMAfPTwoXgb}{}\renewcommand{\rhmcsOosMAfPTwoXgb}{0.34}  % StepM-adjusted p-value (two-sided), XGBoost
\providecommand{\rhmcsOosMAfRawTwoXgb}{}\renewcommand{\rhmcsOosMAfRawTwoXgb}{0.29}  % single-test bootstrap p-value (two-sided), XGBoost
\providecommand{\rhmcsOosMAfPRf}{}\renewcommand{\rhmcsOosMAfPRf}{0.92}  % StepM-adjusted p-value (one-sided), Random forest
\providecommand{\rhmcsOosMAfRawRf}{}\renewcommand{\rhmcsOosMAfRawRf}{0.86}  % single-test bootstrap p-value (one-sided), Random forest
\providecommand{\rhmcsOosMAfPTwoRf}{}\renewcommand{\rhmcsOosMAfPTwoRf}{0.34}  % StepM-adjusted p-value (two-sided), Random forest
\providecommand{\rhmcsOosMAfRawTwoRf}{}\renewcommand{\rhmcsOosMAfRawTwoRf}{0.28}  % single-test bootstrap p-value (two-sided), Random forest
\providecommand{\rhmcsOosMAfNRejLo}{}\renewcommand{\rhmcsOosMAfNRejLo}{0}  % StepM rejections at 5 %, block length 5
\providecommand{\rhmcsOosMAfSpaCLo}{}\renewcommand{\rhmcsOosMAfSpaCLo}{0.73}  % SPA p-value (consistent), block length 5
\providecommand{\rhmcsOosMAfNRejTwoLo}{}\renewcommand{\rhmcsOosMAfNRejTwoLo}{0}  % ... two-sided, block length 5
\providecommand{\rhmcsOosMAfPMinLo}{}\renewcommand{\rhmcsOosMAfPMinLo}{0.97}  % smallest adjusted p, block length 5
\providecommand{\rhmcsOosMAfPTwoMinLo}{}\renewcommand{\rhmcsOosMAfPTwoMinLo}{0.21}  % ... two-sided, block length 5
\providecommand{\rhmcsOosMAfNRejHi}{}\renewcommand{\rhmcsOosMAfNRejHi}{0}  % StepM rejections at 5 %, block length 20
\providecommand{\rhmcsOosMAfSpaCHi}{}\renewcommand{\rhmcsOosMAfSpaCHi}{0.62}  % SPA p-value (consistent), block length 20
\providecommand{\rhmcsOosMAfNRejTwoHi}{}\renewcommand{\rhmcsOosMAfNRejTwoHi}{0}  % ... two-sided, block length 20
\providecommand{\rhmcsOosMAfPMinHi}{}\renewcommand{\rhmcsOosMAfPMinHi}{0.88}  % smallest adjusted p, block length 20
\providecommand{\rhmcsOosMAfPTwoMinHi}{}\renewcommand{\rhmcsOosMAfPTwoMinHi}{0.38}  % ... two-sided, block length 20
\providecommand{\rhmcsOosMBmSpaC}{}\renewcommand{\rhmcsOosMBmSpaC}{0.63}  % SPA p-value (consistent), family OosMBm
\providecommand{\rhmcsOosMBmSpaL}{}\renewcommand{\rhmcsOosMBmSpaL}{0.33}  % ... lower
\providecommand{\rhmcsOosMBmSpaU}{}\renewcommand{\rhmcsOosMBmSpaU}{0.63}  % ... upper
\providecommand{\rhmcsOosMBmNRej}{}\renewcommand{\rhmcsOosMBmNRej}{0}  % StepM rejections at FWER 5 % (of 6)
\providecommand{\rhmcsOosMBmNRejWord}{}\renewcommand{\rhmcsOosMBmNRejWord}{none}  % ... as a word
\providecommand{\rhmcsOosMBmNRaw}{}\renewcommand{\rhmcsOosMBmNRaw}{0}  % single-test bootstrap p < 0.05 (no adjustment)
\providecommand{\rhmcsOosMBmNRawWord}{}\renewcommand{\rhmcsOosMBmNRawWord}{none}  % ... as a word
\providecommand{\rhmcsOosMBmNPos}{}\renewcommand{\rhmcsOosMBmNPos}{1}  % alternatives with a positive estimate (better than the benchmark)
\providecommand{\rhmcsOosMBmPMin}{}\renewcommand{\rhmcsOosMBmPMin}{0.63}  % smallest StepM-adjusted p-value
\providecommand{\rhmcsOosMBmPMax}{}\renewcommand{\rhmcsOosMBmPMax}{0.88}  % largest StepM-adjusted p-value
\providecommand{\rhmcsOosMBmPTwoMin}{}\renewcommand{\rhmcsOosMBmPTwoMin}{0.47}  % smallest two-sided StepM-adjusted p-value
\providecommand{\rhmcsOosMBmPTwoMax}{}\renewcommand{\rhmcsOosMBmPTwoMax}{0.74}  % largest two-sided StepM-adjusted p-value
\providecommand{\rhmcsOosMBmRawMin}{}\renewcommand{\rhmcsOosMBmRawMin}{0.33}  % smallest single-test bootstrap p-value
\providecommand{\rhmcsOosMBmRej}{}\renewcommand{\rhmcsOosMBmRej}{none}  % the alternatives StepM rejects at 5 % (none: none)
\providecommand{\rhmcsOosMBmNRejTwo}{}\renewcommand{\rhmcsOosMBmNRejTwo}{0}  % StepM rejections at FWER 5 %, two-sided
\providecommand{\rhmcsOosMBmNRejTwoWord}{}\renewcommand{\rhmcsOosMBmNRejTwoWord}{none}  % ... as a word
\providecommand{\rhmcsOosMBmNRawTwo}{}\renewcommand{\rhmcsOosMBmNRawTwo}{0}  % single-test bootstrap p < 0.05, two-sided
\providecommand{\rhmcsOosMBmRejTwo}{}\renewcommand{\rhmcsOosMBmRejTwo}{none}  % the alternatives the two-sided StepM rejects at 5 % (none: none)
\providecommand{\rhmcsOosMBmPRidge}{}\renewcommand{\rhmcsOosMBmPRidge}{0.88}  % StepM-adjusted p-value (one-sided), Ridge
\providecommand{\rhmcsOosMBmRawRidge}{}\renewcommand{\rhmcsOosMBmRawRidge}{0.88}  % single-test bootstrap p-value (one-sided), Ridge
\providecommand{\rhmcsOosMBmPTwoRidge}{}\renewcommand{\rhmcsOosMBmPTwoRidge}{0.48}  % StepM-adjusted p-value (two-sided), Ridge
\providecommand{\rhmcsOosMBmRawTwoRidge}{}\renewcommand{\rhmcsOosMBmRawTwoRidge}{0.26}  % single-test bootstrap p-value (two-sided), Ridge
\providecommand{\rhmcsOosMBmPLasso}{}\renewcommand{\rhmcsOosMBmPLasso}{0.88}  % StepM-adjusted p-value (one-sided), Lasso
\providecommand{\rhmcsOosMBmRawLasso}{}\renewcommand{\rhmcsOosMBmRawLasso}{0.88}  % single-test bootstrap p-value (one-sided), Lasso
\providecommand{\rhmcsOosMBmPTwoLasso}{}\renewcommand{\rhmcsOosMBmPTwoLasso}{0.47}  % StepM-adjusted p-value (two-sided), Lasso
\providecommand{\rhmcsOosMBmRawTwoLasso}{}\renewcommand{\rhmcsOosMBmRawTwoLasso}{0.253}  % single-test bootstrap p-value (two-sided), Lasso
\providecommand{\rhmcsOosMBmPEnet}{}\renewcommand{\rhmcsOosMBmPEnet}{0.88}  % StepM-adjusted p-value (one-sided), Elastic net
\providecommand{\rhmcsOosMBmRawEnet}{}\renewcommand{\rhmcsOosMBmRawEnet}{0.88}  % single-test bootstrap p-value (one-sided), Elastic net
\providecommand{\rhmcsOosMBmPTwoEnet}{}\renewcommand{\rhmcsOosMBmPTwoEnet}{0.49}  % StepM-adjusted p-value (two-sided), Elastic net
\providecommand{\rhmcsOosMBmRawTwoEnet}{}\renewcommand{\rhmcsOosMBmRawTwoEnet}{0.26}  % single-test bootstrap p-value (two-sided), Elastic net
\providecommand{\rhmcsOosMBmPLgbm}{}\renewcommand{\rhmcsOosMBmPLgbm}{0.87}  % StepM-adjusted p-value (one-sided), LightGBM
\providecommand{\rhmcsOosMBmRawLgbm}{}\renewcommand{\rhmcsOosMBmRawLgbm}{0.81}  % single-test bootstrap p-value (one-sided), LightGBM
\providecommand{\rhmcsOosMBmPTwoLgbm}{}\renewcommand{\rhmcsOosMBmPTwoLgbm}{0.62}  % StepM-adjusted p-value (two-sided), LightGBM
\providecommand{\rhmcsOosMBmRawTwoLgbm}{}\renewcommand{\rhmcsOosMBmRawTwoLgbm}{0.40}  % single-test bootstrap p-value (two-sided), LightGBM
\providecommand{\rhmcsOosMBmPXgb}{}\renewcommand{\rhmcsOosMBmPXgb}{0.86}  % StepM-adjusted p-value (one-sided), XGBoost
\providecommand{\rhmcsOosMBmRawXgb}{}\renewcommand{\rhmcsOosMBmRawXgb}{0.76}  % single-test bootstrap p-value (one-sided), XGBoost
\providecommand{\rhmcsOosMBmPTwoXgb}{}\renewcommand{\rhmcsOosMBmPTwoXgb}{0.74}  % StepM-adjusted p-value (two-sided), XGBoost
\providecommand{\rhmcsOosMBmRawTwoXgb}{}\renewcommand{\rhmcsOosMBmRawTwoXgb}{0.55}  % single-test bootstrap p-value (two-sided), XGBoost
\providecommand{\rhmcsOosMBmPRf}{}\renewcommand{\rhmcsOosMBmPRf}{0.63}  % StepM-adjusted p-value (one-sided), Random forest
\providecommand{\rhmcsOosMBmRawRf}{}\renewcommand{\rhmcsOosMBmRawRf}{0.33}  % single-test bootstrap p-value (one-sided), Random forest
\providecommand{\rhmcsOosMBmPTwoRf}{}\renewcommand{\rhmcsOosMBmPTwoRf}{0.74}  % StepM-adjusted p-value (two-sided), Random forest
\providecommand{\rhmcsOosMBmRawTwoRf}{}\renewcommand{\rhmcsOosMBmRawTwoRf}{0.72}  % single-test bootstrap p-value (two-sided), Random forest
\providecommand{\rhmcsOosMBmNRejLo}{}\renewcommand{\rhmcsOosMBmNRejLo}{0}  % StepM rejections at 5 %, block length 5
\providecommand{\rhmcsOosMBmSpaCLo}{}\renewcommand{\rhmcsOosMBmSpaCLo}{0.72}  % SPA p-value (consistent), block length 5
\providecommand{\rhmcsOosMBmNRejTwoLo}{}\renewcommand{\rhmcsOosMBmNRejTwoLo}{0}  % ... two-sided, block length 5
\providecommand{\rhmcsOosMBmPMinLo}{}\renewcommand{\rhmcsOosMBmPMinLo}{0.72}  % smallest adjusted p, block length 5
\providecommand{\rhmcsOosMBmPTwoMinLo}{}\renewcommand{\rhmcsOosMBmPTwoMinLo}{0.32}  % ... two-sided, block length 5
\providecommand{\rhmcsOosMBmNRejHi}{}\renewcommand{\rhmcsOosMBmNRejHi}{0}  % StepM rejections at 5 %, block length 20
\providecommand{\rhmcsOosMBmSpaCHi}{}\renewcommand{\rhmcsOosMBmSpaCHi}{0.61}  % SPA p-value (consistent), block length 20
\providecommand{\rhmcsOosMBmNRejTwoHi}{}\renewcommand{\rhmcsOosMBmNRejTwoHi}{0}  % ... two-sided, block length 20
\providecommand{\rhmcsOosMBmPMinHi}{}\renewcommand{\rhmcsOosMBmPMinHi}{0.61}  % smallest adjusted p, block length 20
\providecommand{\rhmcsOosMBmPTwoMinHi}{}\renewcommand{\rhmcsOosMBmPTwoMinHi}{0.58}  % ... two-sided, block length 20
\providecommand{\rhmcsOosMBmAllSpaC}{}\renewcommand{\rhmcsOosMBmAllSpaC}{0.21}  % SPA p-value (consistent), family OosMBmAll
\providecommand{\rhmcsOosMBmAllSpaL}{}\renewcommand{\rhmcsOosMBmAllSpaL}{0.20}  % ... lower
\providecommand{\rhmcsOosMBmAllSpaU}{}\renewcommand{\rhmcsOosMBmAllSpaU}{0.21}  % ... upper
\providecommand{\rhmcsOosMBmAllNRej}{}\renewcommand{\rhmcsOosMBmAllNRej}{0}  % StepM rejections at FWER 5 % (of 10)
\providecommand{\rhmcsOosMBmAllNRejWord}{}\renewcommand{\rhmcsOosMBmAllNRejWord}{none}  % ... as a word
\providecommand{\rhmcsOosMBmAllNRaw}{}\renewcommand{\rhmcsOosMBmAllNRaw}{0}  % single-test bootstrap p < 0.05 (no adjustment)
\providecommand{\rhmcsOosMBmAllNRawWord}{}\renewcommand{\rhmcsOosMBmAllNRawWord}{none}  % ... as a word
\providecommand{\rhmcsOosMBmAllNPos}{}\renewcommand{\rhmcsOosMBmAllNPos}{5}  % alternatives with a positive estimate (better than the benchmark)
\providecommand{\rhmcsOosMBmAllPMin}{}\renewcommand{\rhmcsOosMBmAllPMin}{0.21}  % smallest StepM-adjusted p-value
\providecommand{\rhmcsOosMBmAllPMax}{}\renewcommand{\rhmcsOosMBmAllPMax}{0.88}  % largest StepM-adjusted p-value
\providecommand{\rhmcsOosMBmAllPTwoMin}{}\renewcommand{\rhmcsOosMBmAllPTwoMin}{0.40}  % smallest two-sided StepM-adjusted p-value
\providecommand{\rhmcsOosMBmAllPTwoMax}{}\renewcommand{\rhmcsOosMBmAllPTwoMax}{0.88}  % largest two-sided StepM-adjusted p-value
\providecommand{\rhmcsOosMBmAllRawMin}{}\renewcommand{\rhmcsOosMBmAllRawMin}{0.094}  % smallest single-test bootstrap p-value
\providecommand{\rhmcsOosMBmAllRej}{}\renewcommand{\rhmcsOosMBmAllRej}{none}  % the alternatives StepM rejects at 5 % (none: none)
\providecommand{\rhmcsOosMBmAllNRejTwo}{}\renewcommand{\rhmcsOosMBmAllNRejTwo}{0}  % StepM rejections at FWER 5 %, two-sided
\providecommand{\rhmcsOosMBmAllNRejTwoWord}{}\renewcommand{\rhmcsOosMBmAllNRejTwoWord}{none}  % ... as a word
\providecommand{\rhmcsOosMBmAllNRawTwo}{}\renewcommand{\rhmcsOosMBmAllNRawTwo}{0}  % single-test bootstrap p < 0.05, two-sided
\providecommand{\rhmcsOosMBmAllRejTwo}{}\renewcommand{\rhmcsOosMBmAllRejTwo}{none}  % the alternatives the two-sided StepM rejects at 5 % (none: none)
\providecommand{\rhmcsOosMBmAllPLinHar}{}\renewcommand{\rhmcsOosMBmAllPLinHar}{0.73}  % StepM-adjusted p-value (one-sided), Linear, HAR
\providecommand{\rhmcsOosMBmAllRawLinHar}{}\renewcommand{\rhmcsOosMBmAllRawLinHar}{0.41}  % single-test bootstrap p-value (one-sided), Linear, HAR
\providecommand{\rhmcsOosMBmAllPTwoLinHar}{}\renewcommand{\rhmcsOosMBmAllPTwoLinHar}{0.88}  % StepM-adjusted p-value (two-sided), Linear, HAR
\providecommand{\rhmcsOosMBmAllRawTwoLinHar}{}\renewcommand{\rhmcsOosMBmAllRawTwoLinHar}{0.88}  % single-test bootstrap p-value (two-sided), Linear, HAR
\providecommand{\rhmcsOosMBmAllPRidgeFull}{}\renewcommand{\rhmcsOosMBmAllPRidgeFull}{0.88}  % StepM-adjusted p-value (one-sided), Ridge, Full
\providecommand{\rhmcsOosMBmAllRawRidgeFull}{}\renewcommand{\rhmcsOosMBmAllRawRidgeFull}{0.88}  % single-test bootstrap p-value (one-sided), Ridge, Full
\providecommand{\rhmcsOosMBmAllPTwoRidgeFull}{}\renewcommand{\rhmcsOosMBmAllPTwoRidgeFull}{0.56}  % StepM-adjusted p-value (two-sided), Ridge, Full
\providecommand{\rhmcsOosMBmAllRawTwoRidgeFull}{}\renewcommand{\rhmcsOosMBmAllRawTwoRidgeFull}{0.26}  % single-test bootstrap p-value (two-sided), Ridge, Full
\providecommand{\rhmcsOosMBmAllPLassoFull}{}\renewcommand{\rhmcsOosMBmAllPLassoFull}{0.88}  % StepM-adjusted p-value (one-sided), Lasso, Full
\providecommand{\rhmcsOosMBmAllRawLassoFull}{}\renewcommand{\rhmcsOosMBmAllRawLassoFull}{0.88}  % single-test bootstrap p-value (one-sided), Lasso, Full
\providecommand{\rhmcsOosMBmAllPTwoLassoFull}{}\renewcommand{\rhmcsOosMBmAllPTwoLassoFull}{0.54}  % StepM-adjusted p-value (two-sided), Lasso, Full
\providecommand{\rhmcsOosMBmAllRawTwoLassoFull}{}\renewcommand{\rhmcsOosMBmAllRawTwoLassoFull}{0.253}  % single-test bootstrap p-value (two-sided), Lasso, Full
\providecommand{\rhmcsOosMBmAllPEnetFull}{}\renewcommand{\rhmcsOosMBmAllPEnetFull}{0.88}  % StepM-adjusted p-value (one-sided), Elastic net, Full
\providecommand{\rhmcsOosMBmAllRawEnetFull}{}\renewcommand{\rhmcsOosMBmAllRawEnetFull}{0.88}  % single-test bootstrap p-value (one-sided), Elastic net, Full
\providecommand{\rhmcsOosMBmAllPTwoEnetFull}{}\renewcommand{\rhmcsOosMBmAllPTwoEnetFull}{0.57}  % StepM-adjusted p-value (two-sided), Elastic net, Full
\providecommand{\rhmcsOosMBmAllRawTwoEnetFull}{}\renewcommand{\rhmcsOosMBmAllRawTwoEnetFull}{0.26}  % single-test bootstrap p-value (two-sided), Elastic net, Full
\providecommand{\rhmcsOosMBmAllPLgbmHar}{}\renewcommand{\rhmcsOosMBmAllPLgbmHar}{0.30}  % StepM-adjusted p-value (one-sided), LightGBM, HAR
\providecommand{\rhmcsOosMBmAllRawLgbmHar}{}\renewcommand{\rhmcsOosMBmAllRawLgbmHar}{0.12}  % single-test bootstrap p-value (one-sided), LightGBM, HAR
\providecommand{\rhmcsOosMBmAllPTwoLgbmHar}{}\renewcommand{\rhmcsOosMBmAllPTwoLgbmHar}{0.54}  % StepM-adjusted p-value (two-sided), LightGBM, HAR
\providecommand{\rhmcsOosMBmAllRawTwoLgbmHar}{}\renewcommand{\rhmcsOosMBmAllRawTwoLgbmHar}{0.24}  % single-test bootstrap p-value (two-sided), LightGBM, HAR
\providecommand{\rhmcsOosMBmAllPLgbmFull}{}\renewcommand{\rhmcsOosMBmAllPLgbmFull}{0.87}  % StepM-adjusted p-value (one-sided), LightGBM, Full
\providecommand{\rhmcsOosMBmAllRawLgbmFull}{}\renewcommand{\rhmcsOosMBmAllRawLgbmFull}{0.81}  % single-test bootstrap p-value (one-sided), LightGBM, Full
\providecommand{\rhmcsOosMBmAllPTwoLgbmFull}{}\renewcommand{\rhmcsOosMBmAllPTwoLgbmFull}{0.70}  % StepM-adjusted p-value (two-sided), LightGBM, Full
\providecommand{\rhmcsOosMBmAllRawTwoLgbmFull}{}\renewcommand{\rhmcsOosMBmAllRawTwoLgbmFull}{0.40}  % single-test bootstrap p-value (two-sided), LightGBM, Full
\providecommand{\rhmcsOosMBmAllPXgbHar}{}\renewcommand{\rhmcsOosMBmAllPXgbHar}{0.21}  % StepM-adjusted p-value (one-sided), XGBoost, HAR
\providecommand{\rhmcsOosMBmAllRawXgbHar}{}\renewcommand{\rhmcsOosMBmAllRawXgbHar}{0.094}  % single-test bootstrap p-value (one-sided), XGBoost, HAR
\providecommand{\rhmcsOosMBmAllPTwoXgbHar}{}\renewcommand{\rhmcsOosMBmAllPTwoXgbHar}{0.40}  % StepM-adjusted p-value (two-sided), XGBoost, HAR
\providecommand{\rhmcsOosMBmAllRawTwoXgbHar}{}\renewcommand{\rhmcsOosMBmAllRawTwoXgbHar}{0.16}  % single-test bootstrap p-value (two-sided), XGBoost, HAR
\providecommand{\rhmcsOosMBmAllPXgbFull}{}\renewcommand{\rhmcsOosMBmAllPXgbFull}{0.86}  % StepM-adjusted p-value (one-sided), XGBoost, Full
\providecommand{\rhmcsOosMBmAllRawXgbFull}{}\renewcommand{\rhmcsOosMBmAllRawXgbFull}{0.76}  % single-test bootstrap p-value (one-sided), XGBoost, Full
\providecommand{\rhmcsOosMBmAllPTwoXgbFull}{}\renewcommand{\rhmcsOosMBmAllPTwoXgbFull}{0.82}  % StepM-adjusted p-value (two-sided), XGBoost, Full
\providecommand{\rhmcsOosMBmAllRawTwoXgbFull}{}\renewcommand{\rhmcsOosMBmAllRawTwoXgbFull}{0.55}  % single-test bootstrap p-value (two-sided), XGBoost, Full
\providecommand{\rhmcsOosMBmAllPRfHar}{}\renewcommand{\rhmcsOosMBmAllPRfHar}{0.21}  % StepM-adjusted p-value (one-sided), Random forest, HAR
\providecommand{\rhmcsOosMBmAllRawRfHar}{}\renewcommand{\rhmcsOosMBmAllRawRfHar}{0.100}  % single-test bootstrap p-value (one-sided), Random forest, HAR
\providecommand{\rhmcsOosMBmAllPTwoRfHar}{}\renewcommand{\rhmcsOosMBmAllPTwoRfHar}{0.42}  % StepM-adjusted p-value (two-sided), Random forest, HAR
\providecommand{\rhmcsOosMBmAllRawTwoRfHar}{}\renewcommand{\rhmcsOosMBmAllRawTwoRfHar}{0.18}  % single-test bootstrap p-value (two-sided), Random forest, HAR
\providecommand{\rhmcsOosMBmAllPRfFull}{}\renewcommand{\rhmcsOosMBmAllPRfFull}{0.69}  % StepM-adjusted p-value (one-sided), Random forest, Full
\providecommand{\rhmcsOosMBmAllRawRfFull}{}\renewcommand{\rhmcsOosMBmAllRawRfFull}{0.33}  % single-test bootstrap p-value (one-sided), Random forest, Full
\providecommand{\rhmcsOosMBmAllPTwoRfFull}{}\renewcommand{\rhmcsOosMBmAllPTwoRfFull}{0.85}  % StepM-adjusted p-value (two-sided), Random forest, Full
\providecommand{\rhmcsOosMBmAllRawTwoRfFull}{}\renewcommand{\rhmcsOosMBmAllRawTwoRfFull}{0.72}  % single-test bootstrap p-value (two-sided), Random forest, Full
\providecommand{\rhmcsOosMBmAllNRejLo}{}\renewcommand{\rhmcsOosMBmAllNRejLo}{0}  % StepM rejections at 5 %, block length 5
\providecommand{\rhmcsOosMBmAllSpaCLo}{}\renewcommand{\rhmcsOosMBmAllSpaCLo}{0.14}  % SPA p-value (consistent), block length 5
\providecommand{\rhmcsOosMBmAllNRejTwoLo}{}\renewcommand{\rhmcsOosMBmAllNRejTwoLo}{0}  % ... two-sided, block length 5
\providecommand{\rhmcsOosMBmAllPMinLo}{}\renewcommand{\rhmcsOosMBmAllPMinLo}{0.14}  % smallest adjusted p, block length 5
\providecommand{\rhmcsOosMBmAllPTwoMinLo}{}\renewcommand{\rhmcsOosMBmAllPTwoMinLo}{0.30}  % ... two-sided, block length 5
\providecommand{\rhmcsOosMBmAllNRejHi}{}\renewcommand{\rhmcsOosMBmAllNRejHi}{0}  % StepM rejections at 5 %, block length 20
\providecommand{\rhmcsOosMBmAllSpaCHi}{}\renewcommand{\rhmcsOosMBmAllSpaCHi}{0.31}  % SPA p-value (consistent), block length 20
\providecommand{\rhmcsOosMBmAllNRejTwoHi}{}\renewcommand{\rhmcsOosMBmAllNRejTwoHi}{0}  % ... two-sided, block length 20
\providecommand{\rhmcsOosMBmAllPMinHi}{}\renewcommand{\rhmcsOosMBmAllPMinHi}{0.31}  % smallest adjusted p, block length 20
\providecommand{\rhmcsOosMBmAllPTwoMinHi}{}\renewcommand{\rhmcsOosMBmAllPTwoMinHi}{0.63}  % ... two-sided, block length 20
\providecommand{\rhmcsChQBmSpaC}{}\renewcommand{\rhmcsChQBmSpaC}{0.54}  % SPA p-value (consistent), family ChQBm
\providecommand{\rhmcsChQBmSpaL}{}\renewcommand{\rhmcsChQBmSpaL}{0.11}  % ... lower
\providecommand{\rhmcsChQBmSpaU}{}\renewcommand{\rhmcsChQBmSpaU}{0.67}  % ... upper
\providecommand{\rhmcsChQBmNRej}{}\renewcommand{\rhmcsChQBmNRej}{0}  % StepM rejections at FWER 5 % (of 6)
\providecommand{\rhmcsChQBmNRejWord}{}\renewcommand{\rhmcsChQBmNRejWord}{none}  % ... as a word
\providecommand{\rhmcsChQBmNRaw}{}\renewcommand{\rhmcsChQBmNRaw}{0}  % single-test bootstrap p < 0.05 (no adjustment)
\providecommand{\rhmcsChQBmNRawWord}{}\renewcommand{\rhmcsChQBmNRawWord}{none}  % ... as a word
\providecommand{\rhmcsChQBmNPos}{}\renewcommand{\rhmcsChQBmNPos}{0}  % alternatives with a positive estimate (better than the benchmark)
\providecommand{\rhmcsChQBmPMin}{}\renewcommand{\rhmcsChQBmPMin}{0.95}  % smallest StepM-adjusted p-value
\providecommand{\rhmcsChQBmPMax}{}\renewcommand{\rhmcsChQBmPMax}{1.00}  % largest StepM-adjusted p-value
\providecommand{\rhmcsChQBmPTwoMin}{}\renewcommand{\rhmcsChQBmPTwoMin}{0.004}  % smallest two-sided StepM-adjusted p-value
\providecommand{\rhmcsChQBmPTwoMax}{}\renewcommand{\rhmcsChQBmPTwoMax}{0.22}  % largest two-sided StepM-adjusted p-value
\providecommand{\rhmcsChQBmRawMin}{}\renewcommand{\rhmcsChQBmRawMin}{0.89}  % smallest single-test bootstrap p-value
\providecommand{\rhmcsChQBmRej}{}\renewcommand{\rhmcsChQBmRej}{none}  % the alternatives StepM rejects at 5 % (none: none)
\providecommand{\rhmcsChQBmNRejTwo}{}\renewcommand{\rhmcsChQBmNRejTwo}{4}  % StepM rejections at FWER 5 %, two-sided
\providecommand{\rhmcsChQBmNRejTwoWord}{}\renewcommand{\rhmcsChQBmNRejTwoWord}{four}  % ... as a word
\providecommand{\rhmcsChQBmNRawTwo}{}\renewcommand{\rhmcsChQBmNRawTwo}{4}  % single-test bootstrap p < 0.05, two-sided
\providecommand{\rhmcsChQBmRejTwo}{}\renewcommand{\rhmcsChQBmRejTwo}{ridge regression on the full predictor set, the lasso on the full predictor set, the elastic net on the full predictor set, and the random forest on the full predictor set}  % the alternatives the two-sided StepM rejects at 5 % (none: none)
\providecommand{\rhmcsChQBmPRidge}{}\renewcommand{\rhmcsChQBmPRidge}{1.00}  % StepM-adjusted p-value (one-sided), Ridge
\providecommand{\rhmcsChQBmRawRidge}{}\renewcommand{\rhmcsChQBmRawRidge}{1.00}  % single-test bootstrap p-value (one-sided), Ridge
\providecommand{\rhmcsChQBmPTwoRidge}{}\renewcommand{\rhmcsChQBmPTwoRidge}{0.008}  % StepM-adjusted p-value (two-sided), Ridge
\providecommand{\rhmcsChQBmRawTwoRidge}{}\renewcommand{\rhmcsChQBmRawTwoRidge}{0.003}  % single-test bootstrap p-value (two-sided), Ridge
\providecommand{\rhmcsChQBmPLasso}{}\renewcommand{\rhmcsChQBmPLasso}{1.00}  % StepM-adjusted p-value (one-sided), Lasso
\providecommand{\rhmcsChQBmRawLasso}{}\renewcommand{\rhmcsChQBmRawLasso}{1.00}  % single-test bootstrap p-value (one-sided), Lasso
\providecommand{\rhmcsChQBmPTwoLasso}{}\renewcommand{\rhmcsChQBmPTwoLasso}{0.004}  % StepM-adjusted p-value (two-sided), Lasso
\providecommand{\rhmcsChQBmRawTwoLasso}{}\renewcommand{\rhmcsChQBmRawTwoLasso}{0.001}  % single-test bootstrap p-value (two-sided), Lasso
\providecommand{\rhmcsChQBmPEnet}{}\renewcommand{\rhmcsChQBmPEnet}{1.00}  % StepM-adjusted p-value (one-sided), Elastic net
\providecommand{\rhmcsChQBmRawEnet}{}\renewcommand{\rhmcsChQBmRawEnet}{1.00}  % single-test bootstrap p-value (one-sided), Elastic net
\providecommand{\rhmcsChQBmPTwoEnet}{}\renewcommand{\rhmcsChQBmPTwoEnet}{0.013}  % StepM-adjusted p-value (two-sided), Elastic net
\providecommand{\rhmcsChQBmRawTwoEnet}{}\renewcommand{\rhmcsChQBmRawTwoEnet}{0.005}  % single-test bootstrap p-value (two-sided), Elastic net
\providecommand{\rhmcsChQBmPLgbm}{}\renewcommand{\rhmcsChQBmPLgbm}{0.97}  % StepM-adjusted p-value (one-sided), LightGBM
\providecommand{\rhmcsChQBmRawLgbm}{}\renewcommand{\rhmcsChQBmRawLgbm}{0.91}  % single-test bootstrap p-value (one-sided), LightGBM
\providecommand{\rhmcsChQBmPTwoLgbm}{}\renewcommand{\rhmcsChQBmPTwoLgbm}{0.19}  % StepM-adjusted p-value (two-sided), LightGBM
\providecommand{\rhmcsChQBmRawTwoLgbm}{}\renewcommand{\rhmcsChQBmRawTwoLgbm}{0.17}  % single-test bootstrap p-value (two-sided), LightGBM
\providecommand{\rhmcsChQBmPXgb}{}\renewcommand{\rhmcsChQBmPXgb}{0.95}  % StepM-adjusted p-value (one-sided), XGBoost
\providecommand{\rhmcsChQBmRawXgb}{}\renewcommand{\rhmcsChQBmRawXgb}{0.89}  % single-test bootstrap p-value (one-sided), XGBoost
\providecommand{\rhmcsChQBmPTwoXgb}{}\renewcommand{\rhmcsChQBmPTwoXgb}{0.22}  % StepM-adjusted p-value (two-sided), XGBoost
\providecommand{\rhmcsChQBmRawTwoXgb}{}\renewcommand{\rhmcsChQBmRawTwoXgb}{0.22}  % single-test bootstrap p-value (two-sided), XGBoost
\providecommand{\rhmcsChQBmPRf}{}\renewcommand{\rhmcsChQBmPRf}{1.00}  % StepM-adjusted p-value (one-sided), Random forest
\providecommand{\rhmcsChQBmRawRf}{}\renewcommand{\rhmcsChQBmRawRf}{0.99}  % single-test bootstrap p-value (one-sided), Random forest
\providecommand{\rhmcsChQBmPTwoRf}{}\renewcommand{\rhmcsChQBmPTwoRf}{0.013}  % StepM-adjusted p-value (two-sided), Random forest
\providecommand{\rhmcsChQBmRawTwoRf}{}\renewcommand{\rhmcsChQBmRawTwoRf}{0.008}  % single-test bootstrap p-value (two-sided), Random forest
\providecommand{\rhmcsChQBmNRejLo}{}\renewcommand{\rhmcsChQBmNRejLo}{0}  % StepM rejections at 5 %, block length 5
\providecommand{\rhmcsChQBmSpaCLo}{}\renewcommand{\rhmcsChQBmSpaCLo}{0.54}  % SPA p-value (consistent), block length 5
\providecommand{\rhmcsChQBmNRejTwoLo}{}\renewcommand{\rhmcsChQBmNRejTwoLo}{4}  % ... two-sided, block length 5
\providecommand{\rhmcsChQBmPMinLo}{}\renewcommand{\rhmcsChQBmPMinLo}{0.96}  % smallest adjusted p, block length 5
\providecommand{\rhmcsChQBmPTwoMinLo}{}\renewcommand{\rhmcsChQBmPTwoMinLo}{0.003}  % ... two-sided, block length 5
\providecommand{\rhmcsChQBmNRejHi}{}\renewcommand{\rhmcsChQBmNRejHi}{0}  % StepM rejections at 5 %, block length 20
\providecommand{\rhmcsChQBmSpaCHi}{}\renewcommand{\rhmcsChQBmSpaCHi}{0.54}  % SPA p-value (consistent), block length 20
\providecommand{\rhmcsChQBmNRejTwoHi}{}\renewcommand{\rhmcsChQBmNRejTwoHi}{4}  % ... two-sided, block length 20
\providecommand{\rhmcsChQBmPMinHi}{}\renewcommand{\rhmcsChQBmPMinHi}{0.96}  % smallest adjusted p, block length 20
\providecommand{\rhmcsChQBmPTwoMinHi}{}\renewcommand{\rhmcsChQBmPTwoMinHi}{0.004}  % ... two-sided, block length 20
\providecommand{\rhmcsChQBmAllSpaC}{}\renewcommand{\rhmcsChQBmAllSpaC}{0.54}  % SPA p-value (consistent), family ChQBmAll
\providecommand{\rhmcsChQBmAllSpaL}{}\renewcommand{\rhmcsChQBmAllSpaL}{0.11}  % ... lower
\providecommand{\rhmcsChQBmAllSpaU}{}\renewcommand{\rhmcsChQBmAllSpaU}{0.78}  % ... upper
\providecommand{\rhmcsChQBmAllNRej}{}\renewcommand{\rhmcsChQBmAllNRej}{0}  % StepM rejections at FWER 5 % (of 10)
\providecommand{\rhmcsChQBmAllNRejWord}{}\renewcommand{\rhmcsChQBmAllNRejWord}{none}  % ... as a word
\providecommand{\rhmcsChQBmAllNRaw}{}\renewcommand{\rhmcsChQBmAllNRaw}{0}  % single-test bootstrap p < 0.05 (no adjustment)
\providecommand{\rhmcsChQBmAllNRawWord}{}\renewcommand{\rhmcsChQBmAllNRawWord}{none}  % ... as a word
\providecommand{\rhmcsChQBmAllNPos}{}\renewcommand{\rhmcsChQBmAllNPos}{0}  % alternatives with a positive estimate (better than the benchmark)
\providecommand{\rhmcsChQBmAllPMin}{}\renewcommand{\rhmcsChQBmAllPMin}{0.98}  % smallest StepM-adjusted p-value
\providecommand{\rhmcsChQBmAllPMax}{}\renewcommand{\rhmcsChQBmAllPMax}{1.00}  % largest StepM-adjusted p-value
\providecommand{\rhmcsChQBmAllPTwoMin}{}\renewcommand{\rhmcsChQBmAllPTwoMin}{\ensuremath{<}0.001}  % smallest two-sided StepM-adjusted p-value
\providecommand{\rhmcsChQBmAllPTwoMax}{}\renewcommand{\rhmcsChQBmAllPTwoMax}{0.22}  % largest two-sided StepM-adjusted p-value
\providecommand{\rhmcsChQBmAllRawMin}{}\renewcommand{\rhmcsChQBmAllRawMin}{0.89}  % smallest single-test bootstrap p-value
\providecommand{\rhmcsChQBmAllRej}{}\renewcommand{\rhmcsChQBmAllRej}{none}  % the alternatives StepM rejects at 5 % (none: none)
\providecommand{\rhmcsChQBmAllNRejTwo}{}\renewcommand{\rhmcsChQBmAllNRejTwo}{8}  % StepM rejections at FWER 5 %, two-sided
\providecommand{\rhmcsChQBmAllNRejTwoWord}{}\renewcommand{\rhmcsChQBmAllNRejTwoWord}{eight}  % ... as a word
\providecommand{\rhmcsChQBmAllNRawTwo}{}\renewcommand{\rhmcsChQBmAllNRawTwo}{8}  % single-test bootstrap p < 0.05, two-sided
\providecommand{\rhmcsChQBmAllRejTwo}{}\renewcommand{\rhmcsChQBmAllRejTwo}{the linear models on the HAR benchmark, ridge regression on the full predictor set, the lasso on the full predictor set, the elastic net on the full predictor set, LightGBM on the HAR benchmark, XGBoost on the HAR benchmark, the random forest on the HAR benchmark, and the random forest on the full predictor set}  % the alternatives the two-sided StepM rejects at 5 % (none: none)
\providecommand{\rhmcsChQBmAllPLinHar}{}\renewcommand{\rhmcsChQBmAllPLinHar}{1.00}  % StepM-adjusted p-value (one-sided), Linear, HAR
\providecommand{\rhmcsChQBmAllRawLinHar}{}\renewcommand{\rhmcsChQBmAllRawLinHar}{1.00}  % single-test bootstrap p-value (one-sided), Linear, HAR
\providecommand{\rhmcsChQBmAllPTwoLinHar}{}\renewcommand{\rhmcsChQBmAllPTwoLinHar}{\ensuremath{<}0.001}  % StepM-adjusted p-value (two-sided), Linear, HAR
\providecommand{\rhmcsChQBmAllRawTwoLinHar}{}\renewcommand{\rhmcsChQBmAllRawTwoLinHar}{\ensuremath{<}0.001}  % single-test bootstrap p-value (two-sided), Linear, HAR
\providecommand{\rhmcsChQBmAllPRidgeFull}{}\renewcommand{\rhmcsChQBmAllPRidgeFull}{1.00}  % StepM-adjusted p-value (one-sided), Ridge, Full
\providecommand{\rhmcsChQBmAllRawRidgeFull}{}\renewcommand{\rhmcsChQBmAllRawRidgeFull}{1.00}  % single-test bootstrap p-value (one-sided), Ridge, Full
\providecommand{\rhmcsChQBmAllPTwoRidgeFull}{}\renewcommand{\rhmcsChQBmAllPTwoRidgeFull}{0.008}  % StepM-adjusted p-value (two-sided), Ridge, Full
\providecommand{\rhmcsChQBmAllRawTwoRidgeFull}{}\renewcommand{\rhmcsChQBmAllRawTwoRidgeFull}{0.003}  % single-test bootstrap p-value (two-sided), Ridge, Full
\providecommand{\rhmcsChQBmAllPLassoFull}{}\renewcommand{\rhmcsChQBmAllPLassoFull}{1.00}  % StepM-adjusted p-value (one-sided), Lasso, Full
\providecommand{\rhmcsChQBmAllRawLassoFull}{}\renewcommand{\rhmcsChQBmAllRawLassoFull}{1.00}  % single-test bootstrap p-value (one-sided), Lasso, Full
\providecommand{\rhmcsChQBmAllPTwoLassoFull}{}\renewcommand{\rhmcsChQBmAllPTwoLassoFull}{0.004}  % StepM-adjusted p-value (two-sided), Lasso, Full
\providecommand{\rhmcsChQBmAllRawTwoLassoFull}{}\renewcommand{\rhmcsChQBmAllRawTwoLassoFull}{0.001}  % single-test bootstrap p-value (two-sided), Lasso, Full
\providecommand{\rhmcsChQBmAllPEnetFull}{}\renewcommand{\rhmcsChQBmAllPEnetFull}{1.00}  % StepM-adjusted p-value (one-sided), Elastic net, Full
\providecommand{\rhmcsChQBmAllRawEnetFull}{}\renewcommand{\rhmcsChQBmAllRawEnetFull}{1.00}  % single-test bootstrap p-value (one-sided), Elastic net, Full
\providecommand{\rhmcsChQBmAllPTwoEnetFull}{}\renewcommand{\rhmcsChQBmAllPTwoEnetFull}{0.013}  % StepM-adjusted p-value (two-sided), Elastic net, Full
\providecommand{\rhmcsChQBmAllRawTwoEnetFull}{}\renewcommand{\rhmcsChQBmAllRawTwoEnetFull}{0.005}  % single-test bootstrap p-value (two-sided), Elastic net, Full
\providecommand{\rhmcsChQBmAllPLgbmHar}{}\renewcommand{\rhmcsChQBmAllPLgbmHar}{1.00}  % StepM-adjusted p-value (one-sided), LightGBM, HAR
\providecommand{\rhmcsChQBmAllRawLgbmHar}{}\renewcommand{\rhmcsChQBmAllRawLgbmHar}{1.00}  % single-test bootstrap p-value (one-sided), LightGBM, HAR
\providecommand{\rhmcsChQBmAllPTwoLgbmHar}{}\renewcommand{\rhmcsChQBmAllPTwoLgbmHar}{0.003}  % StepM-adjusted p-value (two-sided), LightGBM, HAR
\providecommand{\rhmcsChQBmAllRawTwoLgbmHar}{}\renewcommand{\rhmcsChQBmAllRawTwoLgbmHar}{\ensuremath{<}0.001}  % single-test bootstrap p-value (two-sided), LightGBM, HAR
\providecommand{\rhmcsChQBmAllPLgbmFull}{}\renewcommand{\rhmcsChQBmAllPLgbmFull}{0.99}  % StepM-adjusted p-value (one-sided), LightGBM, Full
\providecommand{\rhmcsChQBmAllRawLgbmFull}{}\renewcommand{\rhmcsChQBmAllRawLgbmFull}{0.91}  % single-test bootstrap p-value (one-sided), LightGBM, Full
\providecommand{\rhmcsChQBmAllPTwoLgbmFull}{}\renewcommand{\rhmcsChQBmAllPTwoLgbmFull}{0.19}  % StepM-adjusted p-value (two-sided), LightGBM, Full
\providecommand{\rhmcsChQBmAllRawTwoLgbmFull}{}\renewcommand{\rhmcsChQBmAllRawTwoLgbmFull}{0.17}  % single-test bootstrap p-value (two-sided), LightGBM, Full
\providecommand{\rhmcsChQBmAllPXgbHar}{}\renewcommand{\rhmcsChQBmAllPXgbHar}{1.00}  % StepM-adjusted p-value (one-sided), XGBoost, HAR
\providecommand{\rhmcsChQBmAllRawXgbHar}{}\renewcommand{\rhmcsChQBmAllRawXgbHar}{1.00}  % single-test bootstrap p-value (one-sided), XGBoost, HAR
\providecommand{\rhmcsChQBmAllPTwoXgbHar}{}\renewcommand{\rhmcsChQBmAllPTwoXgbHar}{0.002}  % StepM-adjusted p-value (two-sided), XGBoost, HAR
\providecommand{\rhmcsChQBmAllRawTwoXgbHar}{}\renewcommand{\rhmcsChQBmAllRawTwoXgbHar}{\ensuremath{<}0.001}  % single-test bootstrap p-value (two-sided), XGBoost, HAR
\providecommand{\rhmcsChQBmAllPXgbFull}{}\renewcommand{\rhmcsChQBmAllPXgbFull}{0.98}  % StepM-adjusted p-value (one-sided), XGBoost, Full
\providecommand{\rhmcsChQBmAllRawXgbFull}{}\renewcommand{\rhmcsChQBmAllRawXgbFull}{0.89}  % single-test bootstrap p-value (one-sided), XGBoost, Full
\providecommand{\rhmcsChQBmAllPTwoXgbFull}{}\renewcommand{\rhmcsChQBmAllPTwoXgbFull}{0.22}  % StepM-adjusted p-value (two-sided), XGBoost, Full
\providecommand{\rhmcsChQBmAllRawTwoXgbFull}{}\renewcommand{\rhmcsChQBmAllRawTwoXgbFull}{0.22}  % single-test bootstrap p-value (two-sided), XGBoost, Full
\providecommand{\rhmcsChQBmAllPRfHar}{}\renewcommand{\rhmcsChQBmAllPRfHar}{1.00}  % StepM-adjusted p-value (one-sided), Random forest, HAR
\providecommand{\rhmcsChQBmAllRawRfHar}{}\renewcommand{\rhmcsChQBmAllRawRfHar}{1.00}  % single-test bootstrap p-value (one-sided), Random forest, HAR
\providecommand{\rhmcsChQBmAllPTwoRfHar}{}\renewcommand{\rhmcsChQBmAllPTwoRfHar}{\ensuremath{<}0.001}  % StepM-adjusted p-value (two-sided), Random forest, HAR
\providecommand{\rhmcsChQBmAllRawTwoRfHar}{}\renewcommand{\rhmcsChQBmAllRawTwoRfHar}{\ensuremath{<}0.001}  % single-test bootstrap p-value (two-sided), Random forest, HAR
\providecommand{\rhmcsChQBmAllPRfFull}{}\renewcommand{\rhmcsChQBmAllPRfFull}{1.00}  % StepM-adjusted p-value (one-sided), Random forest, Full
\providecommand{\rhmcsChQBmAllRawRfFull}{}\renewcommand{\rhmcsChQBmAllRawRfFull}{0.99}  % single-test bootstrap p-value (one-sided), Random forest, Full
\providecommand{\rhmcsChQBmAllPTwoRfFull}{}\renewcommand{\rhmcsChQBmAllPTwoRfFull}{0.013}  % StepM-adjusted p-value (two-sided), Random forest, Full
\providecommand{\rhmcsChQBmAllRawTwoRfFull}{}\renewcommand{\rhmcsChQBmAllRawTwoRfFull}{0.008}  % single-test bootstrap p-value (two-sided), Random forest, Full
\providecommand{\rhmcsChQBmAllNRejLo}{}\renewcommand{\rhmcsChQBmAllNRejLo}{0}  % StepM rejections at 5 %, block length 5
\providecommand{\rhmcsChQBmAllSpaCLo}{}\renewcommand{\rhmcsChQBmAllSpaCLo}{0.54}  % SPA p-value (consistent), block length 5
\providecommand{\rhmcsChQBmAllNRejTwoLo}{}\renewcommand{\rhmcsChQBmAllNRejTwoLo}{8}  % ... two-sided, block length 5
\providecommand{\rhmcsChQBmAllPMinLo}{}\renewcommand{\rhmcsChQBmAllPMinLo}{0.98}  % smallest adjusted p, block length 5
\providecommand{\rhmcsChQBmAllPTwoMinLo}{}\renewcommand{\rhmcsChQBmAllPTwoMinLo}{\ensuremath{<}0.001}  % ... two-sided, block length 5
\providecommand{\rhmcsChQBmAllNRejHi}{}\renewcommand{\rhmcsChQBmAllNRejHi}{0}  % StepM rejections at 5 %, block length 20
\providecommand{\rhmcsChQBmAllSpaCHi}{}\renewcommand{\rhmcsChQBmAllSpaCHi}{0.54}  % SPA p-value (consistent), block length 20
\providecommand{\rhmcsChQBmAllNRejTwoHi}{}\renewcommand{\rhmcsChQBmAllNRejTwoHi}{8}  % ... two-sided, block length 20
\providecommand{\rhmcsChQBmAllPMinHi}{}\renewcommand{\rhmcsChQBmAllPMinHi}{0.99}  % smallest adjusted p, block length 20
\providecommand{\rhmcsChQBmAllPTwoMinHi}{}\renewcommand{\rhmcsChQBmAllPTwoMinHi}{\ensuremath{<}0.001}  % ... two-sided, block length 20
\providecommand{\rhmcsChMBmSpaC}{}\renewcommand{\rhmcsChMBmSpaC}{0.66}  % SPA p-value (consistent), family ChMBm
\providecommand{\rhmcsChMBmSpaL}{}\renewcommand{\rhmcsChMBmSpaL}{0.36}  % ... lower
\providecommand{\rhmcsChMBmSpaU}{}\renewcommand{\rhmcsChMBmSpaU}{0.66}  % ... upper
\providecommand{\rhmcsChMBmNRej}{}\renewcommand{\rhmcsChMBmNRej}{0}  % StepM rejections at FWER 5 % (of 6)
\providecommand{\rhmcsChMBmNRejWord}{}\renewcommand{\rhmcsChMBmNRejWord}{none}  % ... as a word
\providecommand{\rhmcsChMBmNRaw}{}\renewcommand{\rhmcsChMBmNRaw}{0}  % single-test bootstrap p < 0.05 (no adjustment)
\providecommand{\rhmcsChMBmNRawWord}{}\renewcommand{\rhmcsChMBmNRawWord}{none}  % ... as a word
\providecommand{\rhmcsChMBmNPos}{}\renewcommand{\rhmcsChMBmNPos}{1}  % alternatives with a positive estimate (better than the benchmark)
\providecommand{\rhmcsChMBmPMin}{}\renewcommand{\rhmcsChMBmPMin}{0.66}  % smallest StepM-adjusted p-value
\providecommand{\rhmcsChMBmPMax}{}\renewcommand{\rhmcsChMBmPMax}{0.88}  % largest StepM-adjusted p-value
\providecommand{\rhmcsChMBmPTwoMin}{}\renewcommand{\rhmcsChMBmPTwoMin}{0.51}  % smallest two-sided StepM-adjusted p-value
\providecommand{\rhmcsChMBmPTwoMax}{}\renewcommand{\rhmcsChMBmPTwoMax}{0.77}  % largest two-sided StepM-adjusted p-value
\providecommand{\rhmcsChMBmRawMin}{}\renewcommand{\rhmcsChMBmRawMin}{0.36}  % smallest single-test bootstrap p-value
\providecommand{\rhmcsChMBmRej}{}\renewcommand{\rhmcsChMBmRej}{none}  % the alternatives StepM rejects at 5 % (none: none)
\providecommand{\rhmcsChMBmNRejTwo}{}\renewcommand{\rhmcsChMBmNRejTwo}{0}  % StepM rejections at FWER 5 %, two-sided
\providecommand{\rhmcsChMBmNRejTwoWord}{}\renewcommand{\rhmcsChMBmNRejTwoWord}{none}  % ... as a word
\providecommand{\rhmcsChMBmNRawTwo}{}\renewcommand{\rhmcsChMBmNRawTwo}{0}  % single-test bootstrap p < 0.05, two-sided
\providecommand{\rhmcsChMBmRejTwo}{}\renewcommand{\rhmcsChMBmRejTwo}{none}  % the alternatives the two-sided StepM rejects at 5 % (none: none)
\providecommand{\rhmcsChMBmPRidge}{}\renewcommand{\rhmcsChMBmPRidge}{0.88}  % StepM-adjusted p-value (one-sided), Ridge
\providecommand{\rhmcsChMBmRawRidge}{}\renewcommand{\rhmcsChMBmRawRidge}{0.87}  % single-test bootstrap p-value (one-sided), Ridge
\providecommand{\rhmcsChMBmPTwoRidge}{}\renewcommand{\rhmcsChMBmPTwoRidge}{0.53}  % StepM-adjusted p-value (two-sided), Ridge
\providecommand{\rhmcsChMBmRawTwoRidge}{}\renewcommand{\rhmcsChMBmRawTwoRidge}{0.29}  % single-test bootstrap p-value (two-sided), Ridge
\providecommand{\rhmcsChMBmPLasso}{}\renewcommand{\rhmcsChMBmPLasso}{0.88}  % StepM-adjusted p-value (one-sided), Lasso
\providecommand{\rhmcsChMBmRawLasso}{}\renewcommand{\rhmcsChMBmRawLasso}{0.87}  % single-test bootstrap p-value (one-sided), Lasso
\providecommand{\rhmcsChMBmPTwoLasso}{}\renewcommand{\rhmcsChMBmPTwoLasso}{0.51}  % StepM-adjusted p-value (two-sided), Lasso
\providecommand{\rhmcsChMBmRawTwoLasso}{}\renewcommand{\rhmcsChMBmRawTwoLasso}{0.28}  % single-test bootstrap p-value (two-sided), Lasso
\providecommand{\rhmcsChMBmPEnet}{}\renewcommand{\rhmcsChMBmPEnet}{0.88}  % StepM-adjusted p-value (one-sided), Elastic net
\providecommand{\rhmcsChMBmRawEnet}{}\renewcommand{\rhmcsChMBmRawEnet}{0.87}  % single-test bootstrap p-value (one-sided), Elastic net
\providecommand{\rhmcsChMBmPTwoEnet}{}\renewcommand{\rhmcsChMBmPTwoEnet}{0.53}  % StepM-adjusted p-value (two-sided), Elastic net
\providecommand{\rhmcsChMBmRawTwoEnet}{}\renewcommand{\rhmcsChMBmRawTwoEnet}{0.28}  % single-test bootstrap p-value (two-sided), Elastic net
\providecommand{\rhmcsChMBmPLgbm}{}\renewcommand{\rhmcsChMBmPLgbm}{0.88}  % StepM-adjusted p-value (one-sided), LightGBM
\providecommand{\rhmcsChMBmRawLgbm}{}\renewcommand{\rhmcsChMBmRawLgbm}{0.82}  % single-test bootstrap p-value (one-sided), LightGBM
\providecommand{\rhmcsChMBmPTwoLgbm}{}\renewcommand{\rhmcsChMBmPTwoLgbm}{0.57}  % StepM-adjusted p-value (two-sided), LightGBM
\providecommand{\rhmcsChMBmRawTwoLgbm}{}\renewcommand{\rhmcsChMBmRawTwoLgbm}{0.36}  % single-test bootstrap p-value (two-sided), LightGBM
\providecommand{\rhmcsChMBmPXgb}{}\renewcommand{\rhmcsChMBmPXgb}{0.87}  % StepM-adjusted p-value (one-sided), XGBoost
\providecommand{\rhmcsChMBmRawXgb}{}\renewcommand{\rhmcsChMBmRawXgb}{0.76}  % single-test bootstrap p-value (one-sided), XGBoost
\providecommand{\rhmcsChMBmPTwoXgb}{}\renewcommand{\rhmcsChMBmPTwoXgb}{0.75}  % StepM-adjusted p-value (two-sided), XGBoost
\providecommand{\rhmcsChMBmRawTwoXgb}{}\renewcommand{\rhmcsChMBmRawTwoXgb}{0.55}  % single-test bootstrap p-value (two-sided), XGBoost
\providecommand{\rhmcsChMBmPRf}{}\renewcommand{\rhmcsChMBmPRf}{0.66}  % StepM-adjusted p-value (one-sided), Random forest
\providecommand{\rhmcsChMBmRawRf}{}\renewcommand{\rhmcsChMBmRawRf}{0.36}  % single-test bootstrap p-value (one-sided), Random forest
\providecommand{\rhmcsChMBmPTwoRf}{}\renewcommand{\rhmcsChMBmPTwoRf}{0.77}  % StepM-adjusted p-value (two-sided), Random forest
\providecommand{\rhmcsChMBmRawTwoRf}{}\renewcommand{\rhmcsChMBmRawTwoRf}{0.77}  % single-test bootstrap p-value (two-sided), Random forest
\providecommand{\rhmcsChMBmNRejLo}{}\renewcommand{\rhmcsChMBmNRejLo}{0}  % StepM rejections at 5 %, block length 5
\providecommand{\rhmcsChMBmSpaCLo}{}\renewcommand{\rhmcsChMBmSpaCLo}{0.74}  % SPA p-value (consistent), block length 5
\providecommand{\rhmcsChMBmNRejTwoLo}{}\renewcommand{\rhmcsChMBmNRejTwoLo}{0}  % ... two-sided, block length 5
\providecommand{\rhmcsChMBmPMinLo}{}\renewcommand{\rhmcsChMBmPMinLo}{0.74}  % smallest adjusted p, block length 5
\providecommand{\rhmcsChMBmPTwoMinLo}{}\renewcommand{\rhmcsChMBmPTwoMinLo}{0.34}  % ... two-sided, block length 5
\providecommand{\rhmcsChMBmNRejHi}{}\renewcommand{\rhmcsChMBmNRejHi}{0}  % StepM rejections at 5 %, block length 20
\providecommand{\rhmcsChMBmSpaCHi}{}\renewcommand{\rhmcsChMBmSpaCHi}{0.62}  % SPA p-value (consistent), block length 20
\providecommand{\rhmcsChMBmNRejTwoHi}{}\renewcommand{\rhmcsChMBmNRejTwoHi}{0}  % ... two-sided, block length 20
\providecommand{\rhmcsChMBmPMinHi}{}\renewcommand{\rhmcsChMBmPMinHi}{0.62}  % smallest adjusted p, block length 20
\providecommand{\rhmcsChMBmPTwoMinHi}{}\renewcommand{\rhmcsChMBmPTwoMinHi}{0.65}  % ... two-sided, block length 20
\providecommand{\rhmcsChMBmAllSpaC}{}\renewcommand{\rhmcsChMBmAllSpaC}{0.253}  % SPA p-value (consistent), family ChMBmAll
\providecommand{\rhmcsChMBmAllSpaL}{}\renewcommand{\rhmcsChMBmAllSpaL}{0.23}  % ... lower
\providecommand{\rhmcsChMBmAllSpaU}{}\renewcommand{\rhmcsChMBmAllSpaU}{0.253}  % ... upper
\providecommand{\rhmcsChMBmAllNRej}{}\renewcommand{\rhmcsChMBmAllNRej}{0}  % StepM rejections at FWER 5 % (of 10)
\providecommand{\rhmcsChMBmAllNRejWord}{}\renewcommand{\rhmcsChMBmAllNRejWord}{none}  % ... as a word
\providecommand{\rhmcsChMBmAllNRaw}{}\renewcommand{\rhmcsChMBmAllNRaw}{0}  % single-test bootstrap p < 0.05 (no adjustment)
\providecommand{\rhmcsChMBmAllNRawWord}{}\renewcommand{\rhmcsChMBmAllNRawWord}{none}  % ... as a word
\providecommand{\rhmcsChMBmAllNPos}{}\renewcommand{\rhmcsChMBmAllNPos}{5}  % alternatives with a positive estimate (better than the benchmark)
\providecommand{\rhmcsChMBmAllPMin}{}\renewcommand{\rhmcsChMBmAllPMin}{0.253}  % smallest StepM-adjusted p-value
\providecommand{\rhmcsChMBmAllPMax}{}\renewcommand{\rhmcsChMBmAllPMax}{0.88}  % largest StepM-adjusted p-value
\providecommand{\rhmcsChMBmAllPTwoMin}{}\renewcommand{\rhmcsChMBmAllPTwoMin}{0.48}  % smallest two-sided StepM-adjusted p-value
\providecommand{\rhmcsChMBmAllPTwoMax}{}\renewcommand{\rhmcsChMBmAllPTwoMax}{0.90}  % largest two-sided StepM-adjusted p-value
\providecommand{\rhmcsChMBmAllRawMin}{}\renewcommand{\rhmcsChMBmAllRawMin}{0.11}  % smallest single-test bootstrap p-value
\providecommand{\rhmcsChMBmAllRej}{}\renewcommand{\rhmcsChMBmAllRej}{none}  % the alternatives StepM rejects at 5 % (none: none)
\providecommand{\rhmcsChMBmAllNRejTwo}{}\renewcommand{\rhmcsChMBmAllNRejTwo}{0}  % StepM rejections at FWER 5 %, two-sided
\providecommand{\rhmcsChMBmAllNRejTwoWord}{}\renewcommand{\rhmcsChMBmAllNRejTwoWord}{none}  % ... as a word
\providecommand{\rhmcsChMBmAllNRawTwo}{}\renewcommand{\rhmcsChMBmAllNRawTwo}{0}  % single-test bootstrap p < 0.05, two-sided
\providecommand{\rhmcsChMBmAllRejTwo}{}\renewcommand{\rhmcsChMBmAllRejTwo}{none}  % the alternatives the two-sided StepM rejects at 5 % (none: none)
\providecommand{\rhmcsChMBmAllPLinHar}{}\renewcommand{\rhmcsChMBmAllPLinHar}{0.72}  % StepM-adjusted p-value (one-sided), Linear, HAR
\providecommand{\rhmcsChMBmAllRawLinHar}{}\renewcommand{\rhmcsChMBmAllRawLinHar}{0.41}  % single-test bootstrap p-value (one-sided), Linear, HAR
\providecommand{\rhmcsChMBmAllPTwoLinHar}{}\renewcommand{\rhmcsChMBmAllPTwoLinHar}{0.90}  % StepM-adjusted p-value (two-sided), Linear, HAR
\providecommand{\rhmcsChMBmAllRawTwoLinHar}{}\renewcommand{\rhmcsChMBmAllRawTwoLinHar}{0.86}  % single-test bootstrap p-value (two-sided), Linear, HAR
\providecommand{\rhmcsChMBmAllPRidgeFull}{}\renewcommand{\rhmcsChMBmAllPRidgeFull}{0.88}  % StepM-adjusted p-value (one-sided), Ridge, Full
\providecommand{\rhmcsChMBmAllRawRidgeFull}{}\renewcommand{\rhmcsChMBmAllRawRidgeFull}{0.87}  % single-test bootstrap p-value (one-sided), Ridge, Full
\providecommand{\rhmcsChMBmAllPTwoRidgeFull}{}\renewcommand{\rhmcsChMBmAllPTwoRidgeFull}{0.63}  % StepM-adjusted p-value (two-sided), Ridge, Full
\providecommand{\rhmcsChMBmAllRawTwoRidgeFull}{}\renewcommand{\rhmcsChMBmAllRawTwoRidgeFull}{0.29}  % single-test bootstrap p-value (two-sided), Ridge, Full
\providecommand{\rhmcsChMBmAllPLassoFull}{}\renewcommand{\rhmcsChMBmAllPLassoFull}{0.88}  % StepM-adjusted p-value (one-sided), Lasso, Full
\providecommand{\rhmcsChMBmAllRawLassoFull}{}\renewcommand{\rhmcsChMBmAllRawLassoFull}{0.87}  % single-test bootstrap p-value (one-sided), Lasso, Full
\providecommand{\rhmcsChMBmAllPTwoLassoFull}{}\renewcommand{\rhmcsChMBmAllPTwoLassoFull}{0.61}  % StepM-adjusted p-value (two-sided), Lasso, Full
\providecommand{\rhmcsChMBmAllRawTwoLassoFull}{}\renewcommand{\rhmcsChMBmAllRawTwoLassoFull}{0.28}  % single-test bootstrap p-value (two-sided), Lasso, Full
\providecommand{\rhmcsChMBmAllPEnetFull}{}\renewcommand{\rhmcsChMBmAllPEnetFull}{0.88}  % StepM-adjusted p-value (one-sided), Elastic net, Full
\providecommand{\rhmcsChMBmAllRawEnetFull}{}\renewcommand{\rhmcsChMBmAllRawEnetFull}{0.87}  % single-test bootstrap p-value (one-sided), Elastic net, Full
\providecommand{\rhmcsChMBmAllPTwoEnetFull}{}\renewcommand{\rhmcsChMBmAllPTwoEnetFull}{0.64}  % StepM-adjusted p-value (two-sided), Elastic net, Full
\providecommand{\rhmcsChMBmAllRawTwoEnetFull}{}\renewcommand{\rhmcsChMBmAllRawTwoEnetFull}{0.28}  % single-test bootstrap p-value (two-sided), Elastic net, Full
\providecommand{\rhmcsChMBmAllPLgbmHar}{}\renewcommand{\rhmcsChMBmAllPLgbmHar}{0.48}  % StepM-adjusted p-value (one-sided), LightGBM, HAR
\providecommand{\rhmcsChMBmAllRawLgbmHar}{}\renewcommand{\rhmcsChMBmAllRawLgbmHar}{0.15}  % single-test bootstrap p-value (one-sided), LightGBM, HAR
\providecommand{\rhmcsChMBmAllPTwoLgbmHar}{}\renewcommand{\rhmcsChMBmAllPTwoLgbmHar}{0.68}  % StepM-adjusted p-value (two-sided), LightGBM, HAR
\providecommand{\rhmcsChMBmAllRawTwoLgbmHar}{}\renewcommand{\rhmcsChMBmAllRawTwoLgbmHar}{0.30}  % single-test bootstrap p-value (two-sided), LightGBM, HAR
\providecommand{\rhmcsChMBmAllPLgbmFull}{}\renewcommand{\rhmcsChMBmAllPLgbmFull}{0.88}  % StepM-adjusted p-value (one-sided), LightGBM, Full
\providecommand{\rhmcsChMBmAllRawLgbmFull}{}\renewcommand{\rhmcsChMBmAllRawLgbmFull}{0.82}  % single-test bootstrap p-value (one-sided), LightGBM, Full
\providecommand{\rhmcsChMBmAllPTwoLgbmFull}{}\renewcommand{\rhmcsChMBmAllPTwoLgbmFull}{0.68}  % StepM-adjusted p-value (two-sided), LightGBM, Full
\providecommand{\rhmcsChMBmAllRawTwoLgbmFull}{}\renewcommand{\rhmcsChMBmAllRawTwoLgbmFull}{0.36}  % single-test bootstrap p-value (two-sided), LightGBM, Full
\providecommand{\rhmcsChMBmAllPXgbHar}{}\renewcommand{\rhmcsChMBmAllPXgbHar}{0.255}  % StepM-adjusted p-value (one-sided), XGBoost, HAR
\providecommand{\rhmcsChMBmAllRawXgbHar}{}\renewcommand{\rhmcsChMBmAllRawXgbHar}{0.11}  % single-test bootstrap p-value (one-sided), XGBoost, HAR
\providecommand{\rhmcsChMBmAllPTwoXgbHar}{}\renewcommand{\rhmcsChMBmAllPTwoXgbHar}{0.49}  % StepM-adjusted p-value (two-sided), XGBoost, HAR
\providecommand{\rhmcsChMBmAllRawTwoXgbHar}{}\renewcommand{\rhmcsChMBmAllRawTwoXgbHar}{0.21}  % single-test bootstrap p-value (two-sided), XGBoost, HAR
\providecommand{\rhmcsChMBmAllPXgbFull}{}\renewcommand{\rhmcsChMBmAllPXgbFull}{0.87}  % StepM-adjusted p-value (one-sided), XGBoost, Full
\providecommand{\rhmcsChMBmAllRawXgbFull}{}\renewcommand{\rhmcsChMBmAllRawXgbFull}{0.76}  % single-test bootstrap p-value (one-sided), XGBoost, Full
\providecommand{\rhmcsChMBmAllPTwoXgbFull}{}\renewcommand{\rhmcsChMBmAllPTwoXgbFull}{0.82}  % StepM-adjusted p-value (two-sided), XGBoost, Full
\providecommand{\rhmcsChMBmAllRawTwoXgbFull}{}\renewcommand{\rhmcsChMBmAllRawTwoXgbFull}{0.55}  % single-test bootstrap p-value (two-sided), XGBoost, Full
\providecommand{\rhmcsChMBmAllPRfHar}{}\renewcommand{\rhmcsChMBmAllPRfHar}{0.253}  % StepM-adjusted p-value (one-sided), Random forest, HAR
\providecommand{\rhmcsChMBmAllRawRfHar}{}\renewcommand{\rhmcsChMBmAllRawRfHar}{0.11}  % single-test bootstrap p-value (one-sided), Random forest, HAR
\providecommand{\rhmcsChMBmAllPTwoRfHar}{}\renewcommand{\rhmcsChMBmAllPTwoRfHar}{0.48}  % StepM-adjusted p-value (two-sided), Random forest, HAR
\providecommand{\rhmcsChMBmAllRawTwoRfHar}{}\renewcommand{\rhmcsChMBmAllRawTwoRfHar}{0.21}  % single-test bootstrap p-value (two-sided), Random forest, HAR
\providecommand{\rhmcsChMBmAllPRfFull}{}\renewcommand{\rhmcsChMBmAllPRfFull}{0.71}  % StepM-adjusted p-value (one-sided), Random forest, Full
\providecommand{\rhmcsChMBmAllRawRfFull}{}\renewcommand{\rhmcsChMBmAllRawRfFull}{0.36}  % single-test bootstrap p-value (one-sided), Random forest, Full
\providecommand{\rhmcsChMBmAllPTwoRfFull}{}\renewcommand{\rhmcsChMBmAllPTwoRfFull}{0.90}  % StepM-adjusted p-value (two-sided), Random forest, Full
\providecommand{\rhmcsChMBmAllRawTwoRfFull}{}\renewcommand{\rhmcsChMBmAllRawTwoRfFull}{0.77}  % single-test bootstrap p-value (two-sided), Random forest, Full
\providecommand{\rhmcsChMBmAllNRejLo}{}\renewcommand{\rhmcsChMBmAllNRejLo}{0}  % StepM rejections at 5 %, block length 5
\providecommand{\rhmcsChMBmAllSpaCLo}{}\renewcommand{\rhmcsChMBmAllSpaCLo}{0.21}  % SPA p-value (consistent), block length 5
\providecommand{\rhmcsChMBmAllNRejTwoLo}{}\renewcommand{\rhmcsChMBmAllNRejTwoLo}{0}  % ... two-sided, block length 5
\providecommand{\rhmcsChMBmAllPMinLo}{}\renewcommand{\rhmcsChMBmAllPMinLo}{0.21}  % smallest adjusted p, block length 5
\providecommand{\rhmcsChMBmAllPTwoMinLo}{}\renewcommand{\rhmcsChMBmAllPTwoMinLo}{0.40}  % ... two-sided, block length 5
\providecommand{\rhmcsChMBmAllNRejHi}{}\renewcommand{\rhmcsChMBmAllNRejHi}{0}  % StepM rejections at 5 %, block length 20
\providecommand{\rhmcsChMBmAllSpaCHi}{}\renewcommand{\rhmcsChMBmAllSpaCHi}{0.39}  % SPA p-value (consistent), block length 20
\providecommand{\rhmcsChMBmAllNRejTwoHi}{}\renewcommand{\rhmcsChMBmAllNRejTwoHi}{0}  % ... two-sided, block length 20
\providecommand{\rhmcsChMBmAllPMinHi}{}\renewcommand{\rhmcsChMBmAllPMinHi}{0.39}  % smallest adjusted p, block length 20
\providecommand{\rhmcsChMBmAllPTwoMinHi}{}\renewcommand{\rhmcsChMBmAllPTwoMinHi}{0.69}  % ... two-sided, block length 20
\providecommand{\rhmcsMidShortMeanSpaC}{}\renewcommand{\rhmcsMidShortMeanSpaC}{0.026}  % SPA p-value (consistent), family MidShortMean
\providecommand{\rhmcsMidShortMeanSpaL}{}\renewcommand{\rhmcsMidShortMeanSpaL}{0.026}  % ... lower
\providecommand{\rhmcsMidShortMeanSpaU}{}\renewcommand{\rhmcsMidShortMeanSpaU}{0.026}  % ... upper
\providecommand{\rhmcsMidShortMeanNRej}{}\renewcommand{\rhmcsMidShortMeanNRej}{3}  % StepM rejections at FWER 5 % (of 6)
\providecommand{\rhmcsMidShortMeanNRejWord}{}\renewcommand{\rhmcsMidShortMeanNRejWord}{three}  % ... as a word
\providecommand{\rhmcsMidShortMeanNRaw}{}\renewcommand{\rhmcsMidShortMeanNRaw}{6}  % single-test bootstrap p < 0.05 (no adjustment)
\providecommand{\rhmcsMidShortMeanNRawWord}{}\renewcommand{\rhmcsMidShortMeanNRawWord}{six}  % ... as a word
\providecommand{\rhmcsMidShortMeanNPos}{}\renewcommand{\rhmcsMidShortMeanNPos}{6}  % alternatives with a positive estimate (better than the benchmark)
\providecommand{\rhmcsMidShortMeanPMin}{}\renewcommand{\rhmcsMidShortMeanPMin}{0.026}  % smallest StepM-adjusted p-value
\providecommand{\rhmcsMidShortMeanPMax}{}\renewcommand{\rhmcsMidShortMeanPMax}{0.051}  % largest StepM-adjusted p-value
\providecommand{\rhmcsMidShortMeanPTwoMin}{}\renewcommand{\rhmcsMidShortMeanPTwoMin}{0.034}  % smallest two-sided StepM-adjusted p-value
\providecommand{\rhmcsMidShortMeanPTwoMax}{}\renewcommand{\rhmcsMidShortMeanPTwoMax}{0.083}  % largest two-sided StepM-adjusted p-value
\providecommand{\rhmcsMidShortMeanRawMin}{}\renewcommand{\rhmcsMidShortMeanRawMin}{0.015}  % smallest single-test bootstrap p-value
\providecommand{\rhmcsMidShortMeanRej}{}\renewcommand{\rhmcsMidShortMeanRej}{the lasso on the full predictor set, the elastic net on the full predictor set, and the random forest on the full predictor set}  % the alternatives StepM rejects at 5 % (none: none)
\providecommand{\rhmcsMidShortMeanNRejTwo}{}\renewcommand{\rhmcsMidShortMeanNRejTwo}{2}  % StepM rejections at FWER 5 %, two-sided
\providecommand{\rhmcsMidShortMeanNRejTwoWord}{}\renewcommand{\rhmcsMidShortMeanNRejTwoWord}{two}  % ... as a word
\providecommand{\rhmcsMidShortMeanNRawTwo}{}\renewcommand{\rhmcsMidShortMeanNRawTwo}{3}  % single-test bootstrap p < 0.05, two-sided
\providecommand{\rhmcsMidShortMeanRejTwo}{}\renewcommand{\rhmcsMidShortMeanRejTwo}{the elastic net on the full predictor set and the random forest on the full predictor set}  % the alternatives the two-sided StepM rejects at 5 % (none: none)
\providecommand{\rhmcsMidShortMeanPRidge}{}\renewcommand{\rhmcsMidShortMeanPRidge}{0.051}  % StepM-adjusted p-value (one-sided), Ridge
\providecommand{\rhmcsMidShortMeanRawRidge}{}\renewcommand{\rhmcsMidShortMeanRawRidge}{0.037}  % single-test bootstrap p-value (one-sided), Ridge
\providecommand{\rhmcsMidShortMeanPTwoRidge}{}\renewcommand{\rhmcsMidShortMeanPTwoRidge}{0.083}  % StepM-adjusted p-value (two-sided), Ridge
\providecommand{\rhmcsMidShortMeanRawTwoRidge}{}\renewcommand{\rhmcsMidShortMeanRawTwoRidge}{0.057}  % single-test bootstrap p-value (two-sided), Ridge
\providecommand{\rhmcsMidShortMeanPLasso}{}\renewcommand{\rhmcsMidShortMeanPLasso}{0.038}  % StepM-adjusted p-value (one-sided), Lasso
\providecommand{\rhmcsMidShortMeanRawLasso}{}\renewcommand{\rhmcsMidShortMeanRawLasso}{0.025}  % single-test bootstrap p-value (one-sided), Lasso
\providecommand{\rhmcsMidShortMeanPTwoLasso}{}\renewcommand{\rhmcsMidShortMeanPTwoLasso}{0.059}  % StepM-adjusted p-value (two-sided), Lasso
\providecommand{\rhmcsMidShortMeanRawTwoLasso}{}\renewcommand{\rhmcsMidShortMeanRawTwoLasso}{0.036}  % single-test bootstrap p-value (two-sided), Lasso
\providecommand{\rhmcsMidShortMeanPEnet}{}\renewcommand{\rhmcsMidShortMeanPEnet}{0.031}  % StepM-adjusted p-value (one-sided), Elastic net
\providecommand{\rhmcsMidShortMeanRawEnet}{}\renewcommand{\rhmcsMidShortMeanRawEnet}{0.020}  % single-test bootstrap p-value (one-sided), Elastic net
\providecommand{\rhmcsMidShortMeanPTwoEnet}{}\renewcommand{\rhmcsMidShortMeanPTwoEnet}{0.046}  % StepM-adjusted p-value (two-sided), Elastic net
\providecommand{\rhmcsMidShortMeanRawTwoEnet}{}\renewcommand{\rhmcsMidShortMeanRawTwoEnet}{0.025}  % single-test bootstrap p-value (two-sided), Elastic net
\providecommand{\rhmcsMidShortMeanPLgbm}{}\renewcommand{\rhmcsMidShortMeanPLgbm}{0.051}  % StepM-adjusted p-value (one-sided), LightGBM
\providecommand{\rhmcsMidShortMeanRawLgbm}{}\renewcommand{\rhmcsMidShortMeanRawLgbm}{0.041}  % single-test bootstrap p-value (one-sided), LightGBM
\providecommand{\rhmcsMidShortMeanPTwoLgbm}{}\renewcommand{\rhmcsMidShortMeanPTwoLgbm}{0.083}  % StepM-adjusted p-value (two-sided), LightGBM
\providecommand{\rhmcsMidShortMeanRawTwoLgbm}{}\renewcommand{\rhmcsMidShortMeanRawTwoLgbm}{0.067}  % single-test bootstrap p-value (two-sided), LightGBM
\providecommand{\rhmcsMidShortMeanPXgb}{}\renewcommand{\rhmcsMidShortMeanPXgb}{0.051}  % StepM-adjusted p-value (one-sided), XGBoost
\providecommand{\rhmcsMidShortMeanRawXgb}{}\renewcommand{\rhmcsMidShortMeanRawXgb}{0.039}  % single-test bootstrap p-value (one-sided), XGBoost
\providecommand{\rhmcsMidShortMeanPTwoXgb}{}\renewcommand{\rhmcsMidShortMeanPTwoXgb}{0.083}  % StepM-adjusted p-value (two-sided), XGBoost
\providecommand{\rhmcsMidShortMeanRawTwoXgb}{}\renewcommand{\rhmcsMidShortMeanRawTwoXgb}{0.065}  % single-test bootstrap p-value (two-sided), XGBoost
\providecommand{\rhmcsMidShortMeanPRf}{}\renewcommand{\rhmcsMidShortMeanPRf}{0.026}  % StepM-adjusted p-value (one-sided), Random forest
\providecommand{\rhmcsMidShortMeanRawRf}{}\renewcommand{\rhmcsMidShortMeanRawRf}{0.015}  % single-test bootstrap p-value (one-sided), Random forest
\providecommand{\rhmcsMidShortMeanPTwoRf}{}\renewcommand{\rhmcsMidShortMeanPTwoRf}{0.034}  % StepM-adjusted p-value (two-sided), Random forest
\providecommand{\rhmcsMidShortMeanRawTwoRf}{}\renewcommand{\rhmcsMidShortMeanRawTwoRf}{0.020}  % single-test bootstrap p-value (two-sided), Random forest
\providecommand{\rhmcsMidShortMeanNRejLo}{}\renewcommand{\rhmcsMidShortMeanNRejLo}{3}  % StepM rejections at 5 %, block length 5
\providecommand{\rhmcsMidShortMeanSpaCLo}{}\renewcommand{\rhmcsMidShortMeanSpaCLo}{0.037}  % SPA p-value (consistent), block length 5
\providecommand{\rhmcsMidShortMeanNRejTwoLo}{}\renewcommand{\rhmcsMidShortMeanNRejTwoLo}{0}  % ... two-sided, block length 5
\providecommand{\rhmcsMidShortMeanPMinLo}{}\renewcommand{\rhmcsMidShortMeanPMinLo}{0.037}  % smallest adjusted p, block length 5
\providecommand{\rhmcsMidShortMeanPTwoMinLo}{}\renewcommand{\rhmcsMidShortMeanPTwoMinLo}{0.055}  % ... two-sided, block length 5
\providecommand{\rhmcsMidShortMeanNRejHi}{}\renewcommand{\rhmcsMidShortMeanNRejHi}{3}  % StepM rejections at 5 %, block length 20
\providecommand{\rhmcsMidShortMeanSpaCHi}{}\renewcommand{\rhmcsMidShortMeanSpaCHi}{0.028}  % SPA p-value (consistent), block length 20
\providecommand{\rhmcsMidShortMeanNRejTwoHi}{}\renewcommand{\rhmcsMidShortMeanNRejTwoHi}{1}  % ... two-sided, block length 20
\providecommand{\rhmcsMidShortMeanPMinHi}{}\renewcommand{\rhmcsMidShortMeanPMinHi}{0.028}  % smallest adjusted p, block length 20
\providecommand{\rhmcsMidShortMeanPTwoMinHi}{}\renewcommand{\rhmcsMidShortMeanPTwoMinHi}{0.036}  % ... two-sided, block length 20
\providecommand{\rhmcsMidShortSrSpaC}{}\renewcommand{\rhmcsMidShortSrSpaC}{0.008}  % SPA p-value (consistent), family MidShortSr
\providecommand{\rhmcsMidShortSrSpaL}{}\renewcommand{\rhmcsMidShortSrSpaL}{0.008}  % ... lower
\providecommand{\rhmcsMidShortSrSpaU}{}\renewcommand{\rhmcsMidShortSrSpaU}{0.008}  % ... upper
\providecommand{\rhmcsMidShortSrNRej}{}\renewcommand{\rhmcsMidShortSrNRej}{6}  % StepM rejections at FWER 5 % (of 6)
\providecommand{\rhmcsMidShortSrNRejWord}{}\renewcommand{\rhmcsMidShortSrNRejWord}{six}  % ... as a word
\providecommand{\rhmcsMidShortSrNRaw}{}\renewcommand{\rhmcsMidShortSrNRaw}{6}  % single-test bootstrap p < 0.05 (no adjustment)
\providecommand{\rhmcsMidShortSrNRawWord}{}\renewcommand{\rhmcsMidShortSrNRawWord}{six}  % ... as a word
\providecommand{\rhmcsMidShortSrNPos}{}\renewcommand{\rhmcsMidShortSrNPos}{6}  % alternatives with a positive estimate (better than the benchmark)
\providecommand{\rhmcsMidShortSrPMin}{}\renewcommand{\rhmcsMidShortSrPMin}{0.008}  % smallest StepM-adjusted p-value
\providecommand{\rhmcsMidShortSrPMax}{}\renewcommand{\rhmcsMidShortSrPMax}{0.028}  % largest StepM-adjusted p-value
\providecommand{\rhmcsMidShortSrPTwoMin}{}\renewcommand{\rhmcsMidShortSrPTwoMin}{0.016}  % smallest two-sided StepM-adjusted p-value
\providecommand{\rhmcsMidShortSrPTwoMax}{}\renewcommand{\rhmcsMidShortSrPTwoMax}{0.060}  % largest two-sided StepM-adjusted p-value
\providecommand{\rhmcsMidShortSrRawMin}{}\renewcommand{\rhmcsMidShortSrRawMin}{0.003}  % smallest single-test bootstrap p-value
\providecommand{\rhmcsMidShortSrRej}{}\renewcommand{\rhmcsMidShortSrRej}{ridge regression on the full predictor set, the lasso on the full predictor set, the elastic net on the full predictor set, LightGBM on the full predictor set, XGBoost on the full predictor set, and the random forest on the full predictor set}  % the alternatives StepM rejects at 5 % (none: none)
\providecommand{\rhmcsMidShortSrNRejTwo}{}\renewcommand{\rhmcsMidShortSrNRejTwo}{3}  % StepM rejections at FWER 5 %, two-sided
\providecommand{\rhmcsMidShortSrNRejTwoWord}{}\renewcommand{\rhmcsMidShortSrNRejTwoWord}{three}  % ... as a word
\providecommand{\rhmcsMidShortSrNRawTwo}{}\renewcommand{\rhmcsMidShortSrNRawTwo}{6}  % single-test bootstrap p < 0.05, two-sided
\providecommand{\rhmcsMidShortSrRejTwo}{}\renewcommand{\rhmcsMidShortSrRejTwo}{the lasso on the full predictor set, the elastic net on the full predictor set, and the random forest on the full predictor set}  % the alternatives the two-sided StepM rejects at 5 % (none: none)
\providecommand{\rhmcsMidShortSrPRidge}{}\renewcommand{\rhmcsMidShortSrPRidge}{0.028}  % StepM-adjusted p-value (one-sided), Ridge
\providecommand{\rhmcsMidShortSrRawRidge}{}\renewcommand{\rhmcsMidShortSrRawRidge}{0.019}  % single-test bootstrap p-value (one-sided), Ridge
\providecommand{\rhmcsMidShortSrPTwoRidge}{}\renewcommand{\rhmcsMidShortSrPTwoRidge}{0.060}  % StepM-adjusted p-value (two-sided), Ridge
\providecommand{\rhmcsMidShortSrRawTwoRidge}{}\renewcommand{\rhmcsMidShortSrRawTwoRidge}{0.039}  % single-test bootstrap p-value (two-sided), Ridge
\providecommand{\rhmcsMidShortSrPLasso}{}\renewcommand{\rhmcsMidShortSrPLasso}{0.017}  % StepM-adjusted p-value (one-sided), Lasso
\providecommand{\rhmcsMidShortSrRawLasso}{}\renewcommand{\rhmcsMidShortSrRawLasso}{0.010}  % single-test bootstrap p-value (one-sided), Lasso
\providecommand{\rhmcsMidShortSrPTwoLasso}{}\renewcommand{\rhmcsMidShortSrPTwoLasso}{0.040}  % StepM-adjusted p-value (two-sided), Lasso
\providecommand{\rhmcsMidShortSrRawTwoLasso}{}\renewcommand{\rhmcsMidShortSrRawTwoLasso}{0.021}  % single-test bootstrap p-value (two-sided), Lasso
\providecommand{\rhmcsMidShortSrPEnet}{}\renewcommand{\rhmcsMidShortSrPEnet}{0.012}  % StepM-adjusted p-value (one-sided), Elastic net
\providecommand{\rhmcsMidShortSrRawEnet}{}\renewcommand{\rhmcsMidShortSrRawEnet}{0.007}  % single-test bootstrap p-value (one-sided), Elastic net
\providecommand{\rhmcsMidShortSrPTwoEnet}{}\renewcommand{\rhmcsMidShortSrPTwoEnet}{0.027}  % StepM-adjusted p-value (two-sided), Elastic net
\providecommand{\rhmcsMidShortSrRawTwoEnet}{}\renewcommand{\rhmcsMidShortSrRawTwoEnet}{0.012}  % single-test bootstrap p-value (two-sided), Elastic net
\providecommand{\rhmcsMidShortSrPLgbm}{}\renewcommand{\rhmcsMidShortSrPLgbm}{0.028}  % StepM-adjusted p-value (one-sided), LightGBM
\providecommand{\rhmcsMidShortSrRawLgbm}{}\renewcommand{\rhmcsMidShortSrRawLgbm}{0.021}  % single-test bootstrap p-value (one-sided), LightGBM
\providecommand{\rhmcsMidShortSrPTwoLgbm}{}\renewcommand{\rhmcsMidShortSrPTwoLgbm}{0.060}  % StepM-adjusted p-value (two-sided), LightGBM
\providecommand{\rhmcsMidShortSrRawTwoLgbm}{}\renewcommand{\rhmcsMidShortSrRawTwoLgbm}{0.047}  % single-test bootstrap p-value (two-sided), LightGBM
\providecommand{\rhmcsMidShortSrPXgb}{}\renewcommand{\rhmcsMidShortSrPXgb}{0.028}  % StepM-adjusted p-value (one-sided), XGBoost
\providecommand{\rhmcsMidShortSrRawXgb}{}\renewcommand{\rhmcsMidShortSrRawXgb}{0.020}  % single-test bootstrap p-value (one-sided), XGBoost
\providecommand{\rhmcsMidShortSrPTwoXgb}{}\renewcommand{\rhmcsMidShortSrPTwoXgb}{0.060}  % StepM-adjusted p-value (two-sided), XGBoost
\providecommand{\rhmcsMidShortSrRawTwoXgb}{}\renewcommand{\rhmcsMidShortSrRawTwoXgb}{0.046}  % single-test bootstrap p-value (two-sided), XGBoost
\providecommand{\rhmcsMidShortSrPRf}{}\renewcommand{\rhmcsMidShortSrPRf}{0.008}  % StepM-adjusted p-value (one-sided), Random forest
\providecommand{\rhmcsMidShortSrRawRf}{}\renewcommand{\rhmcsMidShortSrRawRf}{0.003}  % single-test bootstrap p-value (one-sided), Random forest
\providecommand{\rhmcsMidShortSrPTwoRf}{}\renewcommand{\rhmcsMidShortSrPTwoRf}{0.016}  % StepM-adjusted p-value (two-sided), Random forest
\providecommand{\rhmcsMidShortSrRawTwoRf}{}\renewcommand{\rhmcsMidShortSrRawTwoRf}{0.007}  % single-test bootstrap p-value (two-sided), Random forest
\providecommand{\rhmcsMidShortSrNRejLo}{}\renewcommand{\rhmcsMidShortSrNRejLo}{6}  % StepM rejections at 5 %, block length 5
\providecommand{\rhmcsMidShortSrSpaCLo}{}\renewcommand{\rhmcsMidShortSrSpaCLo}{0.010}  % SPA p-value (consistent), block length 5
\providecommand{\rhmcsMidShortSrNRejTwoLo}{}\renewcommand{\rhmcsMidShortSrNRejTwoLo}{3}  % ... two-sided, block length 5
\providecommand{\rhmcsMidShortSrPMinLo}{}\renewcommand{\rhmcsMidShortSrPMinLo}{0.010}  % smallest adjusted p, block length 5
\providecommand{\rhmcsMidShortSrPTwoMinLo}{}\renewcommand{\rhmcsMidShortSrPTwoMinLo}{0.027}  % ... two-sided, block length 5
\providecommand{\rhmcsMidShortSrNRejHi}{}\renewcommand{\rhmcsMidShortSrNRejHi}{6}  % StepM rejections at 5 %, block length 20
\providecommand{\rhmcsMidShortSrSpaCHi}{}\renewcommand{\rhmcsMidShortSrSpaCHi}{0.008}  % SPA p-value (consistent), block length 20
\providecommand{\rhmcsMidShortSrNRejTwoHi}{}\renewcommand{\rhmcsMidShortSrNRejTwoHi}{3}  % ... two-sided, block length 20
\providecommand{\rhmcsMidShortSrPMinHi}{}\renewcommand{\rhmcsMidShortSrPMinHi}{0.008}  % smallest adjusted p, block length 20
\providecommand{\rhmcsMidShortSrPTwoMinHi}{}\renewcommand{\rhmcsMidShortSrPTwoMinHi}{0.015}  % ... two-sided, block length 20
\providecommand{\rhmcsMidTwinMeanSpaC}{}\renewcommand{\rhmcsMidTwinMeanSpaC}{0.027}  % SPA p-value (consistent), family MidTwinMean
\providecommand{\rhmcsMidTwinMeanSpaL}{}\renewcommand{\rhmcsMidTwinMeanSpaL}{0.027}  % ... lower
\providecommand{\rhmcsMidTwinMeanSpaU}{}\renewcommand{\rhmcsMidTwinMeanSpaU}{0.027}  % ... upper
\providecommand{\rhmcsMidTwinMeanNRej}{}\renewcommand{\rhmcsMidTwinMeanNRej}{1}  % StepM rejections at FWER 5 % (of 6)
\providecommand{\rhmcsMidTwinMeanNRejWord}{}\renewcommand{\rhmcsMidTwinMeanNRejWord}{one}  % ... as a word
\providecommand{\rhmcsMidTwinMeanNRaw}{}\renewcommand{\rhmcsMidTwinMeanNRaw}{3}  % single-test bootstrap p < 0.05 (no adjustment)
\providecommand{\rhmcsMidTwinMeanNRawWord}{}\renewcommand{\rhmcsMidTwinMeanNRawWord}{three}  % ... as a word
\providecommand{\rhmcsMidTwinMeanNPos}{}\renewcommand{\rhmcsMidTwinMeanNPos}{6}  % alternatives with a positive estimate (better than the benchmark)
\providecommand{\rhmcsMidTwinMeanPMin}{}\renewcommand{\rhmcsMidTwinMeanPMin}{0.027}  % smallest StepM-adjusted p-value
\providecommand{\rhmcsMidTwinMeanPMax}{}\renewcommand{\rhmcsMidTwinMeanPMax}{0.13}  % largest StepM-adjusted p-value
\providecommand{\rhmcsMidTwinMeanPTwoMin}{}\renewcommand{\rhmcsMidTwinMeanPTwoMin}{0.042}  % smallest two-sided StepM-adjusted p-value
\providecommand{\rhmcsMidTwinMeanPTwoMax}{}\renewcommand{\rhmcsMidTwinMeanPTwoMax}{0.26}  % largest two-sided StepM-adjusted p-value
\providecommand{\rhmcsMidTwinMeanRawMin}{}\renewcommand{\rhmcsMidTwinMeanRawMin}{0.013}  % smallest single-test bootstrap p-value
\providecommand{\rhmcsMidTwinMeanRej}{}\renewcommand{\rhmcsMidTwinMeanRej}{the random forest on the full predictor set}  % the alternatives StepM rejects at 5 % (none: none)
\providecommand{\rhmcsMidTwinMeanNRejTwo}{}\renewcommand{\rhmcsMidTwinMeanNRejTwo}{1}  % StepM rejections at FWER 5 %, two-sided
\providecommand{\rhmcsMidTwinMeanNRejTwoWord}{}\renewcommand{\rhmcsMidTwinMeanNRejTwoWord}{one}  % ... as a word
\providecommand{\rhmcsMidTwinMeanNRawTwo}{}\renewcommand{\rhmcsMidTwinMeanNRawTwo}{2}  % single-test bootstrap p < 0.05, two-sided
\providecommand{\rhmcsMidTwinMeanRejTwo}{}\renewcommand{\rhmcsMidTwinMeanRejTwo}{the random forest on the full predictor set}  % the alternatives the two-sided StepM rejects at 5 % (none: none)
\providecommand{\rhmcsMidTwinMeanPRidge}{}\renewcommand{\rhmcsMidTwinMeanPRidge}{0.13}  % StepM-adjusted p-value (one-sided), Ridge
\providecommand{\rhmcsMidTwinMeanRawRidge}{}\renewcommand{\rhmcsMidTwinMeanRawRidge}{0.082}  % single-test bootstrap p-value (one-sided), Ridge
\providecommand{\rhmcsMidTwinMeanPTwoRidge}{}\renewcommand{\rhmcsMidTwinMeanPTwoRidge}{0.26}  % StepM-adjusted p-value (two-sided), Ridge
\providecommand{\rhmcsMidTwinMeanRawTwoRidge}{}\renewcommand{\rhmcsMidTwinMeanRawTwoRidge}{0.16}  % single-test bootstrap p-value (two-sided), Ridge
\providecommand{\rhmcsMidTwinMeanPLasso}{}\renewcommand{\rhmcsMidTwinMeanPLasso}{0.11}  % StepM-adjusted p-value (one-sided), Lasso
\providecommand{\rhmcsMidTwinMeanRawLasso}{}\renewcommand{\rhmcsMidTwinMeanRawLasso}{0.063}  % single-test bootstrap p-value (one-sided), Lasso
\providecommand{\rhmcsMidTwinMeanPTwoLasso}{}\renewcommand{\rhmcsMidTwinMeanPTwoLasso}{0.21}  % StepM-adjusted p-value (two-sided), Lasso
\providecommand{\rhmcsMidTwinMeanRawTwoLasso}{}\renewcommand{\rhmcsMidTwinMeanRawTwoLasso}{0.12}  % single-test bootstrap p-value (two-sided), Lasso
\providecommand{\rhmcsMidTwinMeanPEnet}{}\renewcommand{\rhmcsMidTwinMeanPEnet}{0.069}  % StepM-adjusted p-value (one-sided), Elastic net
\providecommand{\rhmcsMidTwinMeanRawEnet}{}\renewcommand{\rhmcsMidTwinMeanRawEnet}{0.036}  % single-test bootstrap p-value (one-sided), Elastic net
\providecommand{\rhmcsMidTwinMeanPTwoEnet}{}\renewcommand{\rhmcsMidTwinMeanPTwoEnet}{0.13}  % StepM-adjusted p-value (two-sided), Elastic net
\providecommand{\rhmcsMidTwinMeanRawTwoEnet}{}\renewcommand{\rhmcsMidTwinMeanRawTwoEnet}{0.063}  % single-test bootstrap p-value (two-sided), Elastic net
\providecommand{\rhmcsMidTwinMeanPLgbm}{}\renewcommand{\rhmcsMidTwinMeanPLgbm}{0.060}  % StepM-adjusted p-value (one-sided), LightGBM
\providecommand{\rhmcsMidTwinMeanRawLgbm}{}\renewcommand{\rhmcsMidTwinMeanRawLgbm}{0.025}  % single-test bootstrap p-value (one-sided), LightGBM
\providecommand{\rhmcsMidTwinMeanPTwoLgbm}{}\renewcommand{\rhmcsMidTwinMeanPTwoLgbm}{0.11}  % StepM-adjusted p-value (two-sided), LightGBM
\providecommand{\rhmcsMidTwinMeanRawTwoLgbm}{}\renewcommand{\rhmcsMidTwinMeanRawTwoLgbm}{0.047}  % single-test bootstrap p-value (two-sided), LightGBM
\providecommand{\rhmcsMidTwinMeanPXgb}{}\renewcommand{\rhmcsMidTwinMeanPXgb}{0.13}  % StepM-adjusted p-value (one-sided), XGBoost
\providecommand{\rhmcsMidTwinMeanRawXgb}{}\renewcommand{\rhmcsMidTwinMeanRawXgb}{0.101}  % single-test bootstrap p-value (one-sided), XGBoost
\providecommand{\rhmcsMidTwinMeanPTwoXgb}{}\renewcommand{\rhmcsMidTwinMeanPTwoXgb}{0.26}  % StepM-adjusted p-value (two-sided), XGBoost
\providecommand{\rhmcsMidTwinMeanRawTwoXgb}{}\renewcommand{\rhmcsMidTwinMeanRawTwoXgb}{0.20}  % single-test bootstrap p-value (two-sided), XGBoost
\providecommand{\rhmcsMidTwinMeanPRf}{}\renewcommand{\rhmcsMidTwinMeanPRf}{0.027}  % StepM-adjusted p-value (one-sided), Random forest
\providecommand{\rhmcsMidTwinMeanRawRf}{}\renewcommand{\rhmcsMidTwinMeanRawRf}{0.013}  % single-test bootstrap p-value (one-sided), Random forest
\providecommand{\rhmcsMidTwinMeanPTwoRf}{}\renewcommand{\rhmcsMidTwinMeanPTwoRf}{0.042}  % StepM-adjusted p-value (two-sided), Random forest
\providecommand{\rhmcsMidTwinMeanRawTwoRf}{}\renewcommand{\rhmcsMidTwinMeanRawTwoRf}{0.013}  % single-test bootstrap p-value (two-sided), Random forest
\providecommand{\rhmcsMidTwinMeanNRejLo}{}\renewcommand{\rhmcsMidTwinMeanNRejLo}{1}  % StepM rejections at 5 %, block length 5
\providecommand{\rhmcsMidTwinMeanSpaCLo}{}\renewcommand{\rhmcsMidTwinMeanSpaCLo}{0.036}  % SPA p-value (consistent), block length 5
\providecommand{\rhmcsMidTwinMeanNRejTwoLo}{}\renewcommand{\rhmcsMidTwinMeanNRejTwoLo}{0}  % ... two-sided, block length 5
\providecommand{\rhmcsMidTwinMeanPMinLo}{}\renewcommand{\rhmcsMidTwinMeanPMinLo}{0.036}  % smallest adjusted p, block length 5
\providecommand{\rhmcsMidTwinMeanPTwoMinLo}{}\renewcommand{\rhmcsMidTwinMeanPTwoMinLo}{0.053}  % ... two-sided, block length 5
\providecommand{\rhmcsMidTwinMeanNRejHi}{}\renewcommand{\rhmcsMidTwinMeanNRejHi}{1}  % StepM rejections at 5 %, block length 20
\providecommand{\rhmcsMidTwinMeanSpaCHi}{}\renewcommand{\rhmcsMidTwinMeanSpaCHi}{0.033}  % SPA p-value (consistent), block length 20
\providecommand{\rhmcsMidTwinMeanNRejTwoHi}{}\renewcommand{\rhmcsMidTwinMeanNRejTwoHi}{1}  % ... two-sided, block length 20
\providecommand{\rhmcsMidTwinMeanPMinHi}{}\renewcommand{\rhmcsMidTwinMeanPMinHi}{0.033}  % smallest adjusted p, block length 20
\providecommand{\rhmcsMidTwinMeanPTwoMinHi}{}\renewcommand{\rhmcsMidTwinMeanPTwoMinHi}{0.047}  % ... two-sided, block length 20
\providecommand{\rhmcsMidTwinSrSpaC}{}\renewcommand{\rhmcsMidTwinSrSpaC}{0.007}  % SPA p-value (consistent), family MidTwinSr
\providecommand{\rhmcsMidTwinSrSpaL}{}\renewcommand{\rhmcsMidTwinSrSpaL}{0.007}  % ... lower
\providecommand{\rhmcsMidTwinSrSpaU}{}\renewcommand{\rhmcsMidTwinSrSpaU}{0.007}  % ... upper
\providecommand{\rhmcsMidTwinSrNRej}{}\renewcommand{\rhmcsMidTwinSrNRej}{1}  % StepM rejections at FWER 5 % (of 6)
\providecommand{\rhmcsMidTwinSrNRejWord}{}\renewcommand{\rhmcsMidTwinSrNRejWord}{one}  % ... as a word
\providecommand{\rhmcsMidTwinSrNRaw}{}\renewcommand{\rhmcsMidTwinSrNRaw}{3}  % single-test bootstrap p < 0.05 (no adjustment)
\providecommand{\rhmcsMidTwinSrNRawWord}{}\renewcommand{\rhmcsMidTwinSrNRawWord}{three}  % ... as a word
\providecommand{\rhmcsMidTwinSrNPos}{}\renewcommand{\rhmcsMidTwinSrNPos}{6}  % alternatives with a positive estimate (better than the benchmark)
\providecommand{\rhmcsMidTwinSrPMin}{}\renewcommand{\rhmcsMidTwinSrPMin}{0.007}  % smallest StepM-adjusted p-value
\providecommand{\rhmcsMidTwinSrPMax}{}\renewcommand{\rhmcsMidTwinSrPMax}{0.14}  % largest StepM-adjusted p-value
\providecommand{\rhmcsMidTwinSrPTwoMin}{}\renewcommand{\rhmcsMidTwinSrPTwoMin}{0.012}  % smallest two-sided StepM-adjusted p-value
\providecommand{\rhmcsMidTwinSrPTwoMax}{}\renewcommand{\rhmcsMidTwinSrPTwoMax}{0.27}  % largest two-sided StepM-adjusted p-value
\providecommand{\rhmcsMidTwinSrRawMin}{}\renewcommand{\rhmcsMidTwinSrRawMin}{0.002}  % smallest single-test bootstrap p-value
\providecommand{\rhmcsMidTwinSrRej}{}\renewcommand{\rhmcsMidTwinSrRej}{the random forest on the full predictor set}  % the alternatives StepM rejects at 5 % (none: none)
\providecommand{\rhmcsMidTwinSrNRejTwo}{}\renewcommand{\rhmcsMidTwinSrNRejTwo}{1}  % StepM rejections at FWER 5 %, two-sided
\providecommand{\rhmcsMidTwinSrNRejTwoWord}{}\renewcommand{\rhmcsMidTwinSrNRejTwoWord}{one}  % ... as a word
\providecommand{\rhmcsMidTwinSrNRawTwo}{}\renewcommand{\rhmcsMidTwinSrNRawTwo}{2}  % single-test bootstrap p < 0.05, two-sided
\providecommand{\rhmcsMidTwinSrRejTwo}{}\renewcommand{\rhmcsMidTwinSrRejTwo}{the random forest on the full predictor set}  % the alternatives the two-sided StepM rejects at 5 % (none: none)
\providecommand{\rhmcsMidTwinSrPRidge}{}\renewcommand{\rhmcsMidTwinSrPRidge}{0.14}  % StepM-adjusted p-value (one-sided), Ridge
\providecommand{\rhmcsMidTwinSrRawRidge}{}\renewcommand{\rhmcsMidTwinSrRawRidge}{0.089}  % single-test bootstrap p-value (one-sided), Ridge
\providecommand{\rhmcsMidTwinSrPTwoRidge}{}\renewcommand{\rhmcsMidTwinSrPTwoRidge}{0.27}  % StepM-adjusted p-value (two-sided), Ridge
\providecommand{\rhmcsMidTwinSrRawTwoRidge}{}\renewcommand{\rhmcsMidTwinSrRawTwoRidge}{0.17}  % single-test bootstrap p-value (two-sided), Ridge
\providecommand{\rhmcsMidTwinSrPLasso}{}\renewcommand{\rhmcsMidTwinSrPLasso}{0.12}  % StepM-adjusted p-value (one-sided), Lasso
\providecommand{\rhmcsMidTwinSrRawLasso}{}\renewcommand{\rhmcsMidTwinSrRawLasso}{0.071}  % single-test bootstrap p-value (one-sided), Lasso
\providecommand{\rhmcsMidTwinSrPTwoLasso}{}\renewcommand{\rhmcsMidTwinSrPTwoLasso}{0.22}  % StepM-adjusted p-value (two-sided), Lasso
\providecommand{\rhmcsMidTwinSrRawTwoLasso}{}\renewcommand{\rhmcsMidTwinSrRawTwoLasso}{0.12}  % single-test bootstrap p-value (two-sided), Lasso
\providecommand{\rhmcsMidTwinSrPEnet}{}\renewcommand{\rhmcsMidTwinSrPEnet}{0.082}  % StepM-adjusted p-value (one-sided), Elastic net
\providecommand{\rhmcsMidTwinSrRawEnet}{}\renewcommand{\rhmcsMidTwinSrRawEnet}{0.045}  % single-test bootstrap p-value (one-sided), Elastic net
\providecommand{\rhmcsMidTwinSrPTwoEnet}{}\renewcommand{\rhmcsMidTwinSrPTwoEnet}{0.15}  % StepM-adjusted p-value (two-sided), Elastic net
\providecommand{\rhmcsMidTwinSrRawTwoEnet}{}\renewcommand{\rhmcsMidTwinSrRawTwoEnet}{0.072}  % single-test bootstrap p-value (two-sided), Elastic net
\providecommand{\rhmcsMidTwinSrPLgbm}{}\renewcommand{\rhmcsMidTwinSrPLgbm}{0.065}  % StepM-adjusted p-value (one-sided), LightGBM
\providecommand{\rhmcsMidTwinSrRawLgbm}{}\renewcommand{\rhmcsMidTwinSrRawLgbm}{0.027}  % single-test bootstrap p-value (one-sided), LightGBM
\providecommand{\rhmcsMidTwinSrPTwoLgbm}{}\renewcommand{\rhmcsMidTwinSrPTwoLgbm}{0.12}  % StepM-adjusted p-value (two-sided), LightGBM
\providecommand{\rhmcsMidTwinSrRawTwoLgbm}{}\renewcommand{\rhmcsMidTwinSrRawTwoLgbm}{0.050}  % single-test bootstrap p-value (two-sided), LightGBM
\providecommand{\rhmcsMidTwinSrPXgb}{}\renewcommand{\rhmcsMidTwinSrPXgb}{0.14}  % StepM-adjusted p-value (one-sided), XGBoost
\providecommand{\rhmcsMidTwinSrRawXgb}{}\renewcommand{\rhmcsMidTwinSrRawXgb}{0.11}  % single-test bootstrap p-value (one-sided), XGBoost
\providecommand{\rhmcsMidTwinSrPTwoXgb}{}\renewcommand{\rhmcsMidTwinSrPTwoXgb}{0.27}  % StepM-adjusted p-value (two-sided), XGBoost
\providecommand{\rhmcsMidTwinSrRawTwoXgb}{}\renewcommand{\rhmcsMidTwinSrRawTwoXgb}{0.20}  % single-test bootstrap p-value (two-sided), XGBoost
\providecommand{\rhmcsMidTwinSrPRf}{}\renewcommand{\rhmcsMidTwinSrPRf}{0.007}  % StepM-adjusted p-value (one-sided), Random forest
\providecommand{\rhmcsMidTwinSrRawRf}{}\renewcommand{\rhmcsMidTwinSrRawRf}{0.002}  % single-test bootstrap p-value (one-sided), Random forest
\providecommand{\rhmcsMidTwinSrPTwoRf}{}\renewcommand{\rhmcsMidTwinSrPTwoRf}{0.012}  % StepM-adjusted p-value (two-sided), Random forest
\providecommand{\rhmcsMidTwinSrRawTwoRf}{}\renewcommand{\rhmcsMidTwinSrRawTwoRf}{0.003}  % single-test bootstrap p-value (two-sided), Random forest
\providecommand{\rhmcsMidTwinSrNRejLo}{}\renewcommand{\rhmcsMidTwinSrNRejLo}{1}  % StepM rejections at 5 %, block length 5
\providecommand{\rhmcsMidTwinSrSpaCLo}{}\renewcommand{\rhmcsMidTwinSrSpaCLo}{0.011}  % SPA p-value (consistent), block length 5
\providecommand{\rhmcsMidTwinSrNRejTwoLo}{}\renewcommand{\rhmcsMidTwinSrNRejTwoLo}{1}  % ... two-sided, block length 5
\providecommand{\rhmcsMidTwinSrPMinLo}{}\renewcommand{\rhmcsMidTwinSrPMinLo}{0.011}  % smallest adjusted p, block length 5
\providecommand{\rhmcsMidTwinSrPTwoMinLo}{}\renewcommand{\rhmcsMidTwinSrPTwoMinLo}{0.017}  % ... two-sided, block length 5
\providecommand{\rhmcsMidTwinSrNRejHi}{}\renewcommand{\rhmcsMidTwinSrNRejHi}{1}  % StepM rejections at 5 %, block length 20
\providecommand{\rhmcsMidTwinSrSpaCHi}{}\renewcommand{\rhmcsMidTwinSrSpaCHi}{0.010}  % SPA p-value (consistent), block length 20
\providecommand{\rhmcsMidTwinSrNRejTwoHi}{}\renewcommand{\rhmcsMidTwinSrNRejTwoHi}{1}  % ... two-sided, block length 20
\providecommand{\rhmcsMidTwinSrPMinHi}{}\renewcommand{\rhmcsMidTwinSrPMinHi}{0.010}  % smallest adjusted p, block length 20
\providecommand{\rhmcsMidTwinSrPTwoMinHi}{}\renewcommand{\rhmcsMidTwinSrPTwoMinHi}{0.016}  % ... two-sided, block length 20
\providecommand{\rhmcsBaShortMeanSpaC}{}\renewcommand{\rhmcsBaShortMeanSpaC}{0.025}  % SPA p-value (consistent), family BaShortMean
\providecommand{\rhmcsBaShortMeanSpaL}{}\renewcommand{\rhmcsBaShortMeanSpaL}{0.025}  % ... lower
\providecommand{\rhmcsBaShortMeanSpaU}{}\renewcommand{\rhmcsBaShortMeanSpaU}{0.025}  % ... upper
\providecommand{\rhmcsBaShortMeanNRej}{}\renewcommand{\rhmcsBaShortMeanNRej}{3}  % StepM rejections at FWER 5 % (of 6)
\providecommand{\rhmcsBaShortMeanNRejWord}{}\renewcommand{\rhmcsBaShortMeanNRejWord}{three}  % ... as a word
\providecommand{\rhmcsBaShortMeanNRaw}{}\renewcommand{\rhmcsBaShortMeanNRaw}{6}  % single-test bootstrap p < 0.05 (no adjustment)
\providecommand{\rhmcsBaShortMeanNRawWord}{}\renewcommand{\rhmcsBaShortMeanNRawWord}{six}  % ... as a word
\providecommand{\rhmcsBaShortMeanNPos}{}\renewcommand{\rhmcsBaShortMeanNPos}{6}  % alternatives with a positive estimate (better than the benchmark)
\providecommand{\rhmcsBaShortMeanPMin}{}\renewcommand{\rhmcsBaShortMeanPMin}{0.025}  % smallest StepM-adjusted p-value
\providecommand{\rhmcsBaShortMeanPMax}{}\renewcommand{\rhmcsBaShortMeanPMax}{0.050}  % largest StepM-adjusted p-value
\providecommand{\rhmcsBaShortMeanPTwoMin}{}\renewcommand{\rhmcsBaShortMeanPTwoMin}{0.034}  % smallest two-sided StepM-adjusted p-value
\providecommand{\rhmcsBaShortMeanPTwoMax}{}\renewcommand{\rhmcsBaShortMeanPTwoMax}{0.082}  % largest two-sided StepM-adjusted p-value
\providecommand{\rhmcsBaShortMeanRawMin}{}\renewcommand{\rhmcsBaShortMeanRawMin}{0.015}  % smallest single-test bootstrap p-value
\providecommand{\rhmcsBaShortMeanRej}{}\renewcommand{\rhmcsBaShortMeanRej}{the lasso on the full predictor set, the elastic net on the full predictor set, and the random forest on the full predictor set}  % the alternatives StepM rejects at 5 % (none: none)
\providecommand{\rhmcsBaShortMeanNRejTwo}{}\renewcommand{\rhmcsBaShortMeanNRejTwo}{2}  % StepM rejections at FWER 5 %, two-sided
\providecommand{\rhmcsBaShortMeanNRejTwoWord}{}\renewcommand{\rhmcsBaShortMeanNRejTwoWord}{two}  % ... as a word
\providecommand{\rhmcsBaShortMeanNRawTwo}{}\renewcommand{\rhmcsBaShortMeanNRawTwo}{3}  % single-test bootstrap p < 0.05, two-sided
\providecommand{\rhmcsBaShortMeanRejTwo}{}\renewcommand{\rhmcsBaShortMeanRejTwo}{the elastic net on the full predictor set and the random forest on the full predictor set}  % the alternatives the two-sided StepM rejects at 5 % (none: none)
\providecommand{\rhmcsBaShortMeanPRidge}{}\renewcommand{\rhmcsBaShortMeanPRidge}{0.050}  % StepM-adjusted p-value (one-sided), Ridge
\providecommand{\rhmcsBaShortMeanRawRidge}{}\renewcommand{\rhmcsBaShortMeanRawRidge}{0.036}  % single-test bootstrap p-value (one-sided), Ridge
\providecommand{\rhmcsBaShortMeanPTwoRidge}{}\renewcommand{\rhmcsBaShortMeanPTwoRidge}{0.082}  % StepM-adjusted p-value (two-sided), Ridge
\providecommand{\rhmcsBaShortMeanRawTwoRidge}{}\renewcommand{\rhmcsBaShortMeanRawTwoRidge}{0.056}  % single-test bootstrap p-value (two-sided), Ridge
\providecommand{\rhmcsBaShortMeanPLasso}{}\renewcommand{\rhmcsBaShortMeanPLasso}{0.037}  % StepM-adjusted p-value (one-sided), Lasso
\providecommand{\rhmcsBaShortMeanRawLasso}{}\renewcommand{\rhmcsBaShortMeanRawLasso}{0.025}  % single-test bootstrap p-value (one-sided), Lasso
\providecommand{\rhmcsBaShortMeanPTwoLasso}{}\renewcommand{\rhmcsBaShortMeanPTwoLasso}{0.058}  % StepM-adjusted p-value (two-sided), Lasso
\providecommand{\rhmcsBaShortMeanRawTwoLasso}{}\renewcommand{\rhmcsBaShortMeanRawTwoLasso}{0.035}  % single-test bootstrap p-value (two-sided), Lasso
\providecommand{\rhmcsBaShortMeanPEnet}{}\renewcommand{\rhmcsBaShortMeanPEnet}{0.031}  % StepM-adjusted p-value (one-sided), Elastic net
\providecommand{\rhmcsBaShortMeanRawEnet}{}\renewcommand{\rhmcsBaShortMeanRawEnet}{0.020}  % single-test bootstrap p-value (one-sided), Elastic net
\providecommand{\rhmcsBaShortMeanPTwoEnet}{}\renewcommand{\rhmcsBaShortMeanPTwoEnet}{0.045}  % StepM-adjusted p-value (two-sided), Elastic net
\providecommand{\rhmcsBaShortMeanRawTwoEnet}{}\renewcommand{\rhmcsBaShortMeanRawTwoEnet}{0.025}  % single-test bootstrap p-value (two-sided), Elastic net
\providecommand{\rhmcsBaShortMeanPLgbm}{}\renewcommand{\rhmcsBaShortMeanPLgbm}{0.050}  % StepM-adjusted p-value (one-sided), LightGBM
\providecommand{\rhmcsBaShortMeanRawLgbm}{}\renewcommand{\rhmcsBaShortMeanRawLgbm}{0.041}  % single-test bootstrap p-value (one-sided), LightGBM
\providecommand{\rhmcsBaShortMeanPTwoLgbm}{}\renewcommand{\rhmcsBaShortMeanPTwoLgbm}{0.082}  % StepM-adjusted p-value (two-sided), LightGBM
\providecommand{\rhmcsBaShortMeanRawTwoLgbm}{}\renewcommand{\rhmcsBaShortMeanRawTwoLgbm}{0.067}  % single-test bootstrap p-value (two-sided), LightGBM
\providecommand{\rhmcsBaShortMeanPXgb}{}\renewcommand{\rhmcsBaShortMeanPXgb}{0.050}  % StepM-adjusted p-value (one-sided), XGBoost
\providecommand{\rhmcsBaShortMeanRawXgb}{}\renewcommand{\rhmcsBaShortMeanRawXgb}{0.038}  % single-test bootstrap p-value (one-sided), XGBoost
\providecommand{\rhmcsBaShortMeanPTwoXgb}{}\renewcommand{\rhmcsBaShortMeanPTwoXgb}{0.082}  % StepM-adjusted p-value (two-sided), XGBoost
\providecommand{\rhmcsBaShortMeanRawTwoXgb}{}\renewcommand{\rhmcsBaShortMeanRawTwoXgb}{0.063}  % single-test bootstrap p-value (two-sided), XGBoost
\providecommand{\rhmcsBaShortMeanPRf}{}\renewcommand{\rhmcsBaShortMeanPRf}{0.025}  % StepM-adjusted p-value (one-sided), Random forest
\providecommand{\rhmcsBaShortMeanRawRf}{}\renewcommand{\rhmcsBaShortMeanRawRf}{0.015}  % single-test bootstrap p-value (one-sided), Random forest
\providecommand{\rhmcsBaShortMeanPTwoRf}{}\renewcommand{\rhmcsBaShortMeanPTwoRf}{0.034}  % StepM-adjusted p-value (two-sided), Random forest
\providecommand{\rhmcsBaShortMeanRawTwoRf}{}\renewcommand{\rhmcsBaShortMeanRawTwoRf}{0.019}  % single-test bootstrap p-value (two-sided), Random forest
\providecommand{\rhmcsBaShortMeanNRejLo}{}\renewcommand{\rhmcsBaShortMeanNRejLo}{3}  % StepM rejections at 5 %, block length 5
\providecommand{\rhmcsBaShortMeanSpaCLo}{}\renewcommand{\rhmcsBaShortMeanSpaCLo}{0.037}  % SPA p-value (consistent), block length 5
\providecommand{\rhmcsBaShortMeanNRejTwoLo}{}\renewcommand{\rhmcsBaShortMeanNRejTwoLo}{0}  % ... two-sided, block length 5
\providecommand{\rhmcsBaShortMeanPMinLo}{}\renewcommand{\rhmcsBaShortMeanPMinLo}{0.037}  % smallest adjusted p, block length 5
\providecommand{\rhmcsBaShortMeanPTwoMinLo}{}\renewcommand{\rhmcsBaShortMeanPTwoMinLo}{0.054}  % ... two-sided, block length 5
\providecommand{\rhmcsBaShortMeanNRejHi}{}\renewcommand{\rhmcsBaShortMeanNRejHi}{3}  % StepM rejections at 5 %, block length 20
\providecommand{\rhmcsBaShortMeanSpaCHi}{}\renewcommand{\rhmcsBaShortMeanSpaCHi}{0.027}  % SPA p-value (consistent), block length 20
\providecommand{\rhmcsBaShortMeanNRejTwoHi}{}\renewcommand{\rhmcsBaShortMeanNRejTwoHi}{1}  % ... two-sided, block length 20
\providecommand{\rhmcsBaShortMeanPMinHi}{}\renewcommand{\rhmcsBaShortMeanPMinHi}{0.027}  % smallest adjusted p, block length 20
\providecommand{\rhmcsBaShortMeanPTwoMinHi}{}\renewcommand{\rhmcsBaShortMeanPTwoMinHi}{0.035}  % ... two-sided, block length 20
\providecommand{\rhmcsBaShortSrSpaC}{}\renewcommand{\rhmcsBaShortSrSpaC}{0.007}  % SPA p-value (consistent), family BaShortSr
\providecommand{\rhmcsBaShortSrSpaL}{}\renewcommand{\rhmcsBaShortSrSpaL}{0.007}  % ... lower
\providecommand{\rhmcsBaShortSrSpaU}{}\renewcommand{\rhmcsBaShortSrSpaU}{0.007}  % ... upper
\providecommand{\rhmcsBaShortSrNRej}{}\renewcommand{\rhmcsBaShortSrNRej}{6}  % StepM rejections at FWER 5 % (of 6)
\providecommand{\rhmcsBaShortSrNRejWord}{}\renewcommand{\rhmcsBaShortSrNRejWord}{six}  % ... as a word
\providecommand{\rhmcsBaShortSrNRaw}{}\renewcommand{\rhmcsBaShortSrNRaw}{6}  % single-test bootstrap p < 0.05 (no adjustment)
\providecommand{\rhmcsBaShortSrNRawWord}{}\renewcommand{\rhmcsBaShortSrNRawWord}{six}  % ... as a word
\providecommand{\rhmcsBaShortSrNPos}{}\renewcommand{\rhmcsBaShortSrNPos}{6}  % alternatives with a positive estimate (better than the benchmark)
\providecommand{\rhmcsBaShortSrPMin}{}\renewcommand{\rhmcsBaShortSrPMin}{0.007}  % smallest StepM-adjusted p-value
\providecommand{\rhmcsBaShortSrPMax}{}\renewcommand{\rhmcsBaShortSrPMax}{0.027}  % largest StepM-adjusted p-value
\providecommand{\rhmcsBaShortSrPTwoMin}{}\renewcommand{\rhmcsBaShortSrPTwoMin}{0.015}  % smallest two-sided StepM-adjusted p-value
\providecommand{\rhmcsBaShortSrPTwoMax}{}\renewcommand{\rhmcsBaShortSrPTwoMax}{0.058}  % largest two-sided StepM-adjusted p-value
\providecommand{\rhmcsBaShortSrRawMin}{}\renewcommand{\rhmcsBaShortSrRawMin}{0.002}  % smallest single-test bootstrap p-value
\providecommand{\rhmcsBaShortSrRej}{}\renewcommand{\rhmcsBaShortSrRej}{ridge regression on the full predictor set, the lasso on the full predictor set, the elastic net on the full predictor set, LightGBM on the full predictor set, XGBoost on the full predictor set, and the random forest on the full predictor set}  % the alternatives StepM rejects at 5 % (none: none)
\providecommand{\rhmcsBaShortSrNRejTwo}{}\renewcommand{\rhmcsBaShortSrNRejTwo}{3}  % StepM rejections at FWER 5 %, two-sided
\providecommand{\rhmcsBaShortSrNRejTwoWord}{}\renewcommand{\rhmcsBaShortSrNRejTwoWord}{three}  % ... as a word
\providecommand{\rhmcsBaShortSrNRawTwo}{}\renewcommand{\rhmcsBaShortSrNRawTwo}{6}  % single-test bootstrap p < 0.05, two-sided
\providecommand{\rhmcsBaShortSrRejTwo}{}\renewcommand{\rhmcsBaShortSrRejTwo}{the lasso on the full predictor set, the elastic net on the full predictor set, and the random forest on the full predictor set}  % the alternatives the two-sided StepM rejects at 5 % (none: none)
\providecommand{\rhmcsBaShortSrPRidge}{}\renewcommand{\rhmcsBaShortSrPRidge}{0.027}  % StepM-adjusted p-value (one-sided), Ridge
\providecommand{\rhmcsBaShortSrRawRidge}{}\renewcommand{\rhmcsBaShortSrRawRidge}{0.018}  % single-test bootstrap p-value (one-sided), Ridge
\providecommand{\rhmcsBaShortSrPTwoRidge}{}\renewcommand{\rhmcsBaShortSrPTwoRidge}{0.058}  % StepM-adjusted p-value (two-sided), Ridge
\providecommand{\rhmcsBaShortSrRawTwoRidge}{}\renewcommand{\rhmcsBaShortSrRawTwoRidge}{0.037}  % single-test bootstrap p-value (two-sided), Ridge
\providecommand{\rhmcsBaShortSrPLasso}{}\renewcommand{\rhmcsBaShortSrPLasso}{0.017}  % StepM-adjusted p-value (one-sided), Lasso
\providecommand{\rhmcsBaShortSrRawLasso}{}\renewcommand{\rhmcsBaShortSrRawLasso}{0.010}  % single-test bootstrap p-value (one-sided), Lasso
\providecommand{\rhmcsBaShortSrPTwoLasso}{}\renewcommand{\rhmcsBaShortSrPTwoLasso}{0.039}  % StepM-adjusted p-value (two-sided), Lasso
\providecommand{\rhmcsBaShortSrRawTwoLasso}{}\renewcommand{\rhmcsBaShortSrRawTwoLasso}{0.021}  % single-test bootstrap p-value (two-sided), Lasso
\providecommand{\rhmcsBaShortSrPEnet}{}\renewcommand{\rhmcsBaShortSrPEnet}{0.011}  % StepM-adjusted p-value (one-sided), Elastic net
\providecommand{\rhmcsBaShortSrRawEnet}{}\renewcommand{\rhmcsBaShortSrRawEnet}{0.007}  % single-test bootstrap p-value (one-sided), Elastic net
\providecommand{\rhmcsBaShortSrPTwoEnet}{}\renewcommand{\rhmcsBaShortSrPTwoEnet}{0.026}  % StepM-adjusted p-value (two-sided), Elastic net
\providecommand{\rhmcsBaShortSrRawTwoEnet}{}\renewcommand{\rhmcsBaShortSrRawTwoEnet}{0.012}  % single-test bootstrap p-value (two-sided), Elastic net
\providecommand{\rhmcsBaShortSrPLgbm}{}\renewcommand{\rhmcsBaShortSrPLgbm}{0.027}  % StepM-adjusted p-value (one-sided), LightGBM
\providecommand{\rhmcsBaShortSrRawLgbm}{}\renewcommand{\rhmcsBaShortSrRawLgbm}{0.020}  % single-test bootstrap p-value (one-sided), LightGBM
\providecommand{\rhmcsBaShortSrPTwoLgbm}{}\renewcommand{\rhmcsBaShortSrPTwoLgbm}{0.058}  % StepM-adjusted p-value (two-sided), LightGBM
\providecommand{\rhmcsBaShortSrRawTwoLgbm}{}\renewcommand{\rhmcsBaShortSrRawTwoLgbm}{0.045}  % single-test bootstrap p-value (two-sided), LightGBM
\providecommand{\rhmcsBaShortSrPXgb}{}\renewcommand{\rhmcsBaShortSrPXgb}{0.027}  % StepM-adjusted p-value (one-sided), XGBoost
\providecommand{\rhmcsBaShortSrRawXgb}{}\renewcommand{\rhmcsBaShortSrRawXgb}{0.019}  % single-test bootstrap p-value (one-sided), XGBoost
\providecommand{\rhmcsBaShortSrPTwoXgb}{}\renewcommand{\rhmcsBaShortSrPTwoXgb}{0.058}  % StepM-adjusted p-value (two-sided), XGBoost
\providecommand{\rhmcsBaShortSrRawTwoXgb}{}\renewcommand{\rhmcsBaShortSrRawTwoXgb}{0.043}  % single-test bootstrap p-value (two-sided), XGBoost
\providecommand{\rhmcsBaShortSrPRf}{}\renewcommand{\rhmcsBaShortSrPRf}{0.007}  % StepM-adjusted p-value (one-sided), Random forest
\providecommand{\rhmcsBaShortSrRawRf}{}\renewcommand{\rhmcsBaShortSrRawRf}{0.002}  % single-test bootstrap p-value (one-sided), Random forest
\providecommand{\rhmcsBaShortSrPTwoRf}{}\renewcommand{\rhmcsBaShortSrPTwoRf}{0.015}  % StepM-adjusted p-value (two-sided), Random forest
\providecommand{\rhmcsBaShortSrRawTwoRf}{}\renewcommand{\rhmcsBaShortSrRawTwoRf}{0.006}  % single-test bootstrap p-value (two-sided), Random forest
\providecommand{\rhmcsBaShortSrNRejLo}{}\renewcommand{\rhmcsBaShortSrNRejLo}{6}  % StepM rejections at 5 %, block length 5
\providecommand{\rhmcsBaShortSrSpaCLo}{}\renewcommand{\rhmcsBaShortSrSpaCLo}{0.010}  % SPA p-value (consistent), block length 5
\providecommand{\rhmcsBaShortSrNRejTwoLo}{}\renewcommand{\rhmcsBaShortSrNRejTwoLo}{3}  % ... two-sided, block length 5
\providecommand{\rhmcsBaShortSrPMinLo}{}\renewcommand{\rhmcsBaShortSrPMinLo}{0.010}  % smallest adjusted p, block length 5
\providecommand{\rhmcsBaShortSrPTwoMinLo}{}\renewcommand{\rhmcsBaShortSrPTwoMinLo}{0.027}  % ... two-sided, block length 5
\providecommand{\rhmcsBaShortSrNRejHi}{}\renewcommand{\rhmcsBaShortSrNRejHi}{6}  % StepM rejections at 5 %, block length 20
\providecommand{\rhmcsBaShortSrSpaCHi}{}\renewcommand{\rhmcsBaShortSrSpaCHi}{0.008}  % SPA p-value (consistent), block length 20
\providecommand{\rhmcsBaShortSrNRejTwoHi}{}\renewcommand{\rhmcsBaShortSrNRejTwoHi}{3}  % ... two-sided, block length 20
\providecommand{\rhmcsBaShortSrPMinHi}{}\renewcommand{\rhmcsBaShortSrPMinHi}{0.008}  % smallest adjusted p, block length 20
\providecommand{\rhmcsBaShortSrPTwoMinHi}{}\renewcommand{\rhmcsBaShortSrPTwoMinHi}{0.014}  % ... two-sided, block length 20
\providecommand{\rhmcsBaTwinMeanSpaC}{}\renewcommand{\rhmcsBaTwinMeanSpaC}{0.028}  % SPA p-value (consistent), family BaTwinMean
\providecommand{\rhmcsBaTwinMeanSpaL}{}\renewcommand{\rhmcsBaTwinMeanSpaL}{0.028}  % ... lower
\providecommand{\rhmcsBaTwinMeanSpaU}{}\renewcommand{\rhmcsBaTwinMeanSpaU}{0.028}  % ... upper
\providecommand{\rhmcsBaTwinMeanNRej}{}\renewcommand{\rhmcsBaTwinMeanNRej}{1}  % StepM rejections at FWER 5 % (of 6)
\providecommand{\rhmcsBaTwinMeanNRejWord}{}\renewcommand{\rhmcsBaTwinMeanNRejWord}{one}  % ... as a word
\providecommand{\rhmcsBaTwinMeanNRaw}{}\renewcommand{\rhmcsBaTwinMeanNRaw}{3}  % single-test bootstrap p < 0.05 (no adjustment)
\providecommand{\rhmcsBaTwinMeanNRawWord}{}\renewcommand{\rhmcsBaTwinMeanNRawWord}{three}  % ... as a word
\providecommand{\rhmcsBaTwinMeanNPos}{}\renewcommand{\rhmcsBaTwinMeanNPos}{6}  % alternatives with a positive estimate (better than the benchmark)
\providecommand{\rhmcsBaTwinMeanPMin}{}\renewcommand{\rhmcsBaTwinMeanPMin}{0.028}  % smallest StepM-adjusted p-value
\providecommand{\rhmcsBaTwinMeanPMax}{}\renewcommand{\rhmcsBaTwinMeanPMax}{0.13}  % largest StepM-adjusted p-value
\providecommand{\rhmcsBaTwinMeanPTwoMin}{}\renewcommand{\rhmcsBaTwinMeanPTwoMin}{0.042}  % smallest two-sided StepM-adjusted p-value
\providecommand{\rhmcsBaTwinMeanPTwoMax}{}\renewcommand{\rhmcsBaTwinMeanPTwoMax}{0.26}  % largest two-sided StepM-adjusted p-value
\providecommand{\rhmcsBaTwinMeanRawMin}{}\renewcommand{\rhmcsBaTwinMeanRawMin}{0.013}  % smallest single-test bootstrap p-value
\providecommand{\rhmcsBaTwinMeanRej}{}\renewcommand{\rhmcsBaTwinMeanRej}{the random forest on the full predictor set}  % the alternatives StepM rejects at 5 % (none: none)
\providecommand{\rhmcsBaTwinMeanNRejTwo}{}\renewcommand{\rhmcsBaTwinMeanNRejTwo}{1}  % StepM rejections at FWER 5 %, two-sided
\providecommand{\rhmcsBaTwinMeanNRejTwoWord}{}\renewcommand{\rhmcsBaTwinMeanNRejTwoWord}{one}  % ... as a word
\providecommand{\rhmcsBaTwinMeanNRawTwo}{}\renewcommand{\rhmcsBaTwinMeanNRawTwo}{2}  % single-test bootstrap p < 0.05, two-sided
\providecommand{\rhmcsBaTwinMeanRejTwo}{}\renewcommand{\rhmcsBaTwinMeanRejTwo}{the random forest on the full predictor set}  % the alternatives the two-sided StepM rejects at 5 % (none: none)
\providecommand{\rhmcsBaTwinMeanPRidge}{}\renewcommand{\rhmcsBaTwinMeanPRidge}{0.13}  % StepM-adjusted p-value (one-sided), Ridge
\providecommand{\rhmcsBaTwinMeanRawRidge}{}\renewcommand{\rhmcsBaTwinMeanRawRidge}{0.083}  % single-test bootstrap p-value (one-sided), Ridge
\providecommand{\rhmcsBaTwinMeanPTwoRidge}{}\renewcommand{\rhmcsBaTwinMeanPTwoRidge}{0.26}  % StepM-adjusted p-value (two-sided), Ridge
\providecommand{\rhmcsBaTwinMeanRawTwoRidge}{}\renewcommand{\rhmcsBaTwinMeanRawTwoRidge}{0.16}  % single-test bootstrap p-value (two-sided), Ridge
\providecommand{\rhmcsBaTwinMeanPLasso}{}\renewcommand{\rhmcsBaTwinMeanPLasso}{0.11}  % StepM-adjusted p-value (one-sided), Lasso
\providecommand{\rhmcsBaTwinMeanRawLasso}{}\renewcommand{\rhmcsBaTwinMeanRawLasso}{0.063}  % single-test bootstrap p-value (one-sided), Lasso
\providecommand{\rhmcsBaTwinMeanPTwoLasso}{}\renewcommand{\rhmcsBaTwinMeanPTwoLasso}{0.21}  % StepM-adjusted p-value (two-sided), Lasso
\providecommand{\rhmcsBaTwinMeanRawTwoLasso}{}\renewcommand{\rhmcsBaTwinMeanRawTwoLasso}{0.12}  % single-test bootstrap p-value (two-sided), Lasso
\providecommand{\rhmcsBaTwinMeanPEnet}{}\renewcommand{\rhmcsBaTwinMeanPEnet}{0.069}  % StepM-adjusted p-value (one-sided), Elastic net
\providecommand{\rhmcsBaTwinMeanRawEnet}{}\renewcommand{\rhmcsBaTwinMeanRawEnet}{0.036}  % single-test bootstrap p-value (one-sided), Elastic net
\providecommand{\rhmcsBaTwinMeanPTwoEnet}{}\renewcommand{\rhmcsBaTwinMeanPTwoEnet}{0.13}  % StepM-adjusted p-value (two-sided), Elastic net
\providecommand{\rhmcsBaTwinMeanRawTwoEnet}{}\renewcommand{\rhmcsBaTwinMeanRawTwoEnet}{0.063}  % single-test bootstrap p-value (two-sided), Elastic net
\providecommand{\rhmcsBaTwinMeanPLgbm}{}\renewcommand{\rhmcsBaTwinMeanPLgbm}{0.061}  % StepM-adjusted p-value (one-sided), LightGBM
\providecommand{\rhmcsBaTwinMeanRawLgbm}{}\renewcommand{\rhmcsBaTwinMeanRawLgbm}{0.025}  % single-test bootstrap p-value (one-sided), LightGBM
\providecommand{\rhmcsBaTwinMeanPTwoLgbm}{}\renewcommand{\rhmcsBaTwinMeanPTwoLgbm}{0.12}  % StepM-adjusted p-value (two-sided), LightGBM
\providecommand{\rhmcsBaTwinMeanRawTwoLgbm}{}\renewcommand{\rhmcsBaTwinMeanRawTwoLgbm}{0.048}  % single-test bootstrap p-value (two-sided), LightGBM
\providecommand{\rhmcsBaTwinMeanPXgb}{}\renewcommand{\rhmcsBaTwinMeanPXgb}{0.13}  % StepM-adjusted p-value (one-sided), XGBoost
\providecommand{\rhmcsBaTwinMeanRawXgb}{}\renewcommand{\rhmcsBaTwinMeanRawXgb}{0.103}  % single-test bootstrap p-value (one-sided), XGBoost
\providecommand{\rhmcsBaTwinMeanPTwoXgb}{}\renewcommand{\rhmcsBaTwinMeanPTwoXgb}{0.26}  % StepM-adjusted p-value (two-sided), XGBoost
\providecommand{\rhmcsBaTwinMeanRawTwoXgb}{}\renewcommand{\rhmcsBaTwinMeanRawTwoXgb}{0.20}  % single-test bootstrap p-value (two-sided), XGBoost
\providecommand{\rhmcsBaTwinMeanPRf}{}\renewcommand{\rhmcsBaTwinMeanPRf}{0.028}  % StepM-adjusted p-value (one-sided), Random forest
\providecommand{\rhmcsBaTwinMeanRawRf}{}\renewcommand{\rhmcsBaTwinMeanRawRf}{0.013}  % single-test bootstrap p-value (one-sided), Random forest
\providecommand{\rhmcsBaTwinMeanPTwoRf}{}\renewcommand{\rhmcsBaTwinMeanPTwoRf}{0.042}  % StepM-adjusted p-value (two-sided), Random forest
\providecommand{\rhmcsBaTwinMeanRawTwoRf}{}\renewcommand{\rhmcsBaTwinMeanRawTwoRf}{0.014}  % single-test bootstrap p-value (two-sided), Random forest
\providecommand{\rhmcsBaTwinMeanNRejLo}{}\renewcommand{\rhmcsBaTwinMeanNRejLo}{1}  % StepM rejections at 5 %, block length 5
\providecommand{\rhmcsBaTwinMeanSpaCLo}{}\renewcommand{\rhmcsBaTwinMeanSpaCLo}{0.036}  % SPA p-value (consistent), block length 5
\providecommand{\rhmcsBaTwinMeanNRejTwoLo}{}\renewcommand{\rhmcsBaTwinMeanNRejTwoLo}{0}  % ... two-sided, block length 5
\providecommand{\rhmcsBaTwinMeanPMinLo}{}\renewcommand{\rhmcsBaTwinMeanPMinLo}{0.036}  % smallest adjusted p, block length 5
\providecommand{\rhmcsBaTwinMeanPTwoMinLo}{}\renewcommand{\rhmcsBaTwinMeanPTwoMinLo}{0.054}  % ... two-sided, block length 5
\providecommand{\rhmcsBaTwinMeanNRejHi}{}\renewcommand{\rhmcsBaTwinMeanNRejHi}{1}  % StepM rejections at 5 %, block length 20
\providecommand{\rhmcsBaTwinMeanSpaCHi}{}\renewcommand{\rhmcsBaTwinMeanSpaCHi}{0.034}  % SPA p-value (consistent), block length 20
\providecommand{\rhmcsBaTwinMeanNRejTwoHi}{}\renewcommand{\rhmcsBaTwinMeanNRejTwoHi}{1}  % ... two-sided, block length 20
\providecommand{\rhmcsBaTwinMeanPMinHi}{}\renewcommand{\rhmcsBaTwinMeanPMinHi}{0.034}  % smallest adjusted p, block length 20
\providecommand{\rhmcsBaTwinMeanPTwoMinHi}{}\renewcommand{\rhmcsBaTwinMeanPTwoMinHi}{0.048}  % ... two-sided, block length 20
\providecommand{\rhmcsBaTwinSrSpaC}{}\renewcommand{\rhmcsBaTwinSrSpaC}{0.008}  % SPA p-value (consistent), family BaTwinSr
\providecommand{\rhmcsBaTwinSrSpaL}{}\renewcommand{\rhmcsBaTwinSrSpaL}{0.008}  % ... lower
\providecommand{\rhmcsBaTwinSrSpaU}{}\renewcommand{\rhmcsBaTwinSrSpaU}{0.008}  % ... upper
\providecommand{\rhmcsBaTwinSrNRej}{}\renewcommand{\rhmcsBaTwinSrNRej}{1}  % StepM rejections at FWER 5 % (of 6)
\providecommand{\rhmcsBaTwinSrNRejWord}{}\renewcommand{\rhmcsBaTwinSrNRejWord}{one}  % ... as a word
\providecommand{\rhmcsBaTwinSrNRaw}{}\renewcommand{\rhmcsBaTwinSrNRaw}{3}  % single-test bootstrap p < 0.05 (no adjustment)
\providecommand{\rhmcsBaTwinSrNRawWord}{}\renewcommand{\rhmcsBaTwinSrNRawWord}{three}  % ... as a word
\providecommand{\rhmcsBaTwinSrNPos}{}\renewcommand{\rhmcsBaTwinSrNPos}{6}  % alternatives with a positive estimate (better than the benchmark)
\providecommand{\rhmcsBaTwinSrPMin}{}\renewcommand{\rhmcsBaTwinSrPMin}{0.008}  % smallest StepM-adjusted p-value
\providecommand{\rhmcsBaTwinSrPMax}{}\renewcommand{\rhmcsBaTwinSrPMax}{0.14}  % largest StepM-adjusted p-value
\providecommand{\rhmcsBaTwinSrPTwoMin}{}\renewcommand{\rhmcsBaTwinSrPTwoMin}{0.013}  % smallest two-sided StepM-adjusted p-value
\providecommand{\rhmcsBaTwinSrPTwoMax}{}\renewcommand{\rhmcsBaTwinSrPTwoMax}{0.27}  % largest two-sided StepM-adjusted p-value
\providecommand{\rhmcsBaTwinSrRawMin}{}\renewcommand{\rhmcsBaTwinSrRawMin}{0.002}  % smallest single-test bootstrap p-value
\providecommand{\rhmcsBaTwinSrRej}{}\renewcommand{\rhmcsBaTwinSrRej}{the random forest on the full predictor set}  % the alternatives StepM rejects at 5 % (none: none)
\providecommand{\rhmcsBaTwinSrNRejTwo}{}\renewcommand{\rhmcsBaTwinSrNRejTwo}{1}  % StepM rejections at FWER 5 %, two-sided
\providecommand{\rhmcsBaTwinSrNRejTwoWord}{}\renewcommand{\rhmcsBaTwinSrNRejTwoWord}{one}  % ... as a word
\providecommand{\rhmcsBaTwinSrNRawTwo}{}\renewcommand{\rhmcsBaTwinSrNRawTwo}{2}  % single-test bootstrap p < 0.05, two-sided
\providecommand{\rhmcsBaTwinSrRejTwo}{}\renewcommand{\rhmcsBaTwinSrRejTwo}{the random forest on the full predictor set}  % the alternatives the two-sided StepM rejects at 5 % (none: none)
\providecommand{\rhmcsBaTwinSrPRidge}{}\renewcommand{\rhmcsBaTwinSrPRidge}{0.14}  % StepM-adjusted p-value (one-sided), Ridge
\providecommand{\rhmcsBaTwinSrRawRidge}{}\renewcommand{\rhmcsBaTwinSrRawRidge}{0.089}  % single-test bootstrap p-value (one-sided), Ridge
\providecommand{\rhmcsBaTwinSrPTwoRidge}{}\renewcommand{\rhmcsBaTwinSrPTwoRidge}{0.27}  % StepM-adjusted p-value (two-sided), Ridge
\providecommand{\rhmcsBaTwinSrRawTwoRidge}{}\renewcommand{\rhmcsBaTwinSrRawTwoRidge}{0.17}  % single-test bootstrap p-value (two-sided), Ridge
\providecommand{\rhmcsBaTwinSrPLasso}{}\renewcommand{\rhmcsBaTwinSrPLasso}{0.12}  % StepM-adjusted p-value (one-sided), Lasso
\providecommand{\rhmcsBaTwinSrRawLasso}{}\renewcommand{\rhmcsBaTwinSrRawLasso}{0.071}  % single-test bootstrap p-value (one-sided), Lasso
\providecommand{\rhmcsBaTwinSrPTwoLasso}{}\renewcommand{\rhmcsBaTwinSrPTwoLasso}{0.22}  % StepM-adjusted p-value (two-sided), Lasso
\providecommand{\rhmcsBaTwinSrRawTwoLasso}{}\renewcommand{\rhmcsBaTwinSrRawTwoLasso}{0.12}  % single-test bootstrap p-value (two-sided), Lasso
\providecommand{\rhmcsBaTwinSrPEnet}{}\renewcommand{\rhmcsBaTwinSrPEnet}{0.080}  % StepM-adjusted p-value (one-sided), Elastic net
\providecommand{\rhmcsBaTwinSrRawEnet}{}\renewcommand{\rhmcsBaTwinSrRawEnet}{0.043}  % single-test bootstrap p-value (one-sided), Elastic net
\providecommand{\rhmcsBaTwinSrPTwoEnet}{}\renewcommand{\rhmcsBaTwinSrPTwoEnet}{0.14}  % StepM-adjusted p-value (two-sided), Elastic net
\providecommand{\rhmcsBaTwinSrRawTwoEnet}{}\renewcommand{\rhmcsBaTwinSrRawTwoEnet}{0.070}  % single-test bootstrap p-value (two-sided), Elastic net
\providecommand{\rhmcsBaTwinSrPLgbm}{}\renewcommand{\rhmcsBaTwinSrPLgbm}{0.066}  % StepM-adjusted p-value (one-sided), LightGBM
\providecommand{\rhmcsBaTwinSrRawLgbm}{}\renewcommand{\rhmcsBaTwinSrRawLgbm}{0.027}  % single-test bootstrap p-value (one-sided), LightGBM
\providecommand{\rhmcsBaTwinSrPTwoLgbm}{}\renewcommand{\rhmcsBaTwinSrPTwoLgbm}{0.12}  % StepM-adjusted p-value (two-sided), LightGBM
\providecommand{\rhmcsBaTwinSrRawTwoLgbm}{}\renewcommand{\rhmcsBaTwinSrRawTwoLgbm}{0.050}  % single-test bootstrap p-value (two-sided), LightGBM
\providecommand{\rhmcsBaTwinSrPXgb}{}\renewcommand{\rhmcsBaTwinSrPXgb}{0.14}  % StepM-adjusted p-value (one-sided), XGBoost
\providecommand{\rhmcsBaTwinSrRawXgb}{}\renewcommand{\rhmcsBaTwinSrRawXgb}{0.11}  % single-test bootstrap p-value (one-sided), XGBoost
\providecommand{\rhmcsBaTwinSrPTwoXgb}{}\renewcommand{\rhmcsBaTwinSrPTwoXgb}{0.27}  % StepM-adjusted p-value (two-sided), XGBoost
\providecommand{\rhmcsBaTwinSrRawTwoXgb}{}\renewcommand{\rhmcsBaTwinSrRawTwoXgb}{0.21}  % single-test bootstrap p-value (two-sided), XGBoost
\providecommand{\rhmcsBaTwinSrPRf}{}\renewcommand{\rhmcsBaTwinSrPRf}{0.008}  % StepM-adjusted p-value (one-sided), Random forest
\providecommand{\rhmcsBaTwinSrRawRf}{}\renewcommand{\rhmcsBaTwinSrRawRf}{0.002}  % single-test bootstrap p-value (one-sided), Random forest
\providecommand{\rhmcsBaTwinSrPTwoRf}{}\renewcommand{\rhmcsBaTwinSrPTwoRf}{0.013}  % StepM-adjusted p-value (two-sided), Random forest
\providecommand{\rhmcsBaTwinSrRawTwoRf}{}\renewcommand{\rhmcsBaTwinSrRawTwoRf}{0.003}  % single-test bootstrap p-value (two-sided), Random forest
\providecommand{\rhmcsBaTwinSrNRejLo}{}\renewcommand{\rhmcsBaTwinSrNRejLo}{1}  % StepM rejections at 5 %, block length 5
\providecommand{\rhmcsBaTwinSrSpaCLo}{}\renewcommand{\rhmcsBaTwinSrSpaCLo}{0.011}  % SPA p-value (consistent), block length 5
\providecommand{\rhmcsBaTwinSrNRejTwoLo}{}\renewcommand{\rhmcsBaTwinSrNRejTwoLo}{1}  % ... two-sided, block length 5
\providecommand{\rhmcsBaTwinSrPMinLo}{}\renewcommand{\rhmcsBaTwinSrPMinLo}{0.011}  % smallest adjusted p, block length 5
\providecommand{\rhmcsBaTwinSrPTwoMinLo}{}\renewcommand{\rhmcsBaTwinSrPTwoMinLo}{0.017}  % ... two-sided, block length 5
\providecommand{\rhmcsBaTwinSrNRejHi}{}\renewcommand{\rhmcsBaTwinSrNRejHi}{1}  % StepM rejections at 5 %, block length 20
\providecommand{\rhmcsBaTwinSrSpaCHi}{}\renewcommand{\rhmcsBaTwinSrSpaCHi}{0.012}  % SPA p-value (consistent), block length 20
\providecommand{\rhmcsBaTwinSrNRejTwoHi}{}\renewcommand{\rhmcsBaTwinSrNRejTwoHi}{1}  % ... two-sided, block length 20
\providecommand{\rhmcsBaTwinSrPMinHi}{}\renewcommand{\rhmcsBaTwinSrPMinHi}{0.012}  % smallest adjusted p, block length 20
\providecommand{\rhmcsBaTwinSrPTwoMinHi}{}\renewcommand{\rhmcsBaTwinSrPTwoMinHi}{0.018}  % ... two-sided, block length 20
\providecommand{\rhmcsTwinTwoSidedPRidge}{}\renewcommand{\rhmcsTwinTwoSidedPRidge}{0.17}  % two-sided pooled-SR bootstrap p of the Full - HAR Sharpe difference, Ridge (D19: the check p_bt; the paper's test is Ledoit-Wolf)
\providecommand{\rhmcsTwinTwoSidedPLasso}{}\renewcommand{\rhmcsTwinTwoSidedPLasso}{0.12}  % two-sided pooled-SR bootstrap p of the Full - HAR Sharpe difference, Lasso (D19: the check p_bt; the paper's test is Ledoit-Wolf)
\providecommand{\rhmcsTwinTwoSidedPEnet}{}\renewcommand{\rhmcsTwinTwoSidedPEnet}{0.072}  % two-sided pooled-SR bootstrap p of the Full - HAR Sharpe difference, Elastic net (D19: the check p_bt; the paper's test is Ledoit-Wolf)
\providecommand{\rhmcsTwinTwoSidedPLgbm}{}\renewcommand{\rhmcsTwinTwoSidedPLgbm}{0.050}  % two-sided pooled-SR bootstrap p of the Full - HAR Sharpe difference, LightGBM (D19: the check p_bt; the paper's test is Ledoit-Wolf)
\providecommand{\rhmcsTwinTwoSidedPXgb}{}\renewcommand{\rhmcsTwinTwoSidedPXgb}{0.20}  % two-sided pooled-SR bootstrap p of the Full - HAR Sharpe difference, XGBoost (D19: the check p_bt; the paper's test is Ledoit-Wolf)
\providecommand{\rhmcsTwinTwoSidedPRf}{}\renewcommand{\rhmcsTwinTwoSidedPRf}{0.003}  % two-sided pooled-SR bootstrap p of the Full - HAR Sharpe difference, Random forest (D19: the check p_bt; the paper's test is Ledoit-Wolf)
\fi

\section{Results}
\label{sec:rh_results}
\label{sec:rh_exp_setup}
\begin{itemize}
\item \textbf{Design, benchmark and outline.} Every forecast follows the design, spans and tests of
  Sections~\ref{sec:rh_methods_setup} and~\ref{sec:rh_results_eval}: every bar of $\mathcal T_{\mathrm{oos}}$ is
  scored, and the primary benchmark is the OLS on the HAR benchmark set, the OLS row of Table~\ref{tab:rh_oos}, which
  the HAR-set penalized models reproduce. The window length was fixed after out-of-sample scoring
  (Section~\ref{sec:rh_results_robust}). We report accuracy on every bar (Section~\ref{sec:rh_results_oos}), the
  predictor groups that carry the gain (Section~\ref{sec:rh_results_ablation}) and robustness (Section~\ref{sec:rh_results_robust});
  Section~\ref{sec:rh_fi} asks which series the models use.
\end{itemize}

\subsection{Accuracy on every bar}
\label{sec:rh_results_oos}
\label{sec:rh_results_predictors}
\begin{itemize}
\item \textbf{Overview.} Table~\ref{tab:rh_oos} reports the MSE, the average QLIKE loss and $R^2_{\mathrm{OOS}}$ of
  the six model classes on the HAR benchmark set and on the full predictor set.
\item \textbf{Against the primary benchmark.} The full predictor set lowers the average QLIKE loss of every model
  class below that of the OLS on the HAR benchmark set.
\item \textbf{Within each class.} The paired comparison of $f_{m,\mathrm{Full},\mathcal I_{\mathrm{rh}}}$ with
  $f_{m,\mathrm{HAR},\mathcal I_{\mathrm{rh}}}$ (Section~\ref{sec:rh_results_eval}; the DM column of
  Table~\ref{tab:rh_oos}) isolates what the exogenous predictors add to a given model $m$.
\item \textbf{Squared-error loss.} Under squared error the comparison reverses: on every bar the full set raises the
  MSE in all \rhorAfPairs{} classes, by \rhcaLinMsePctBarSpanMin{} to \rhcaLinMsePctBarSpanMax\% for the penalized
  linear models and by \rhcaTreeMsePctBarSpanMin{} to \rhcaTreeMsePctBarSpanMax\% for the tree ensembles, and lowers
  $R^2_{\mathrm{OOS}}$ (Table~\ref{tab:rh_oos}). A handful of days in the COVID-19 crash account for the difference
  (Appendix~\ref{app:rh_results_calib}).
\end{itemize}

\subsection{Which predictors carry the gain?}
\label{sec:rh_results_ablation}
\ifdefined\rhbyLXgbReading\else% AUTO-GENERATED by writeup/make_rehaul_beyond_ladder.py -- do not edit; re-run the script.
% results/ablations_2014start_w3000/scores_fomcfix; requested by the author after the beyond-base results; not pre-registered
\newcommand{\rhbyLLgbmFullAllPct}{$-$1.6}  % lgbm FULL vs BASE, all
\newcommand{\rhbyLLgbmKOneAllPct}{$-$1.2}  % lgbm k=1 vs BASE, all; requested by the author after the beyond-base results; not pre-registered
\newcommand{\rhbyLLgbmKOneAllDm}{$-$4.76}  % DM vs BASE
\newcommand{\rhbyLLgbmKOneAllP}{\ensuremath{<}0.001}  % p-value
\newcommand{\rhbyLLgbmKOneAllDmFull}{2.59}  % DM vs FULL (positive: FULL lower)
\newcommand{\rhbyLLgbmKOneAllPFull}{0.010}  % p-value
\newcommand{\rhbyLLgbmKOneAllR}{0.74}  % recovery (L_BASE - L_k) / (L_BASE - L_FULL)
\newcommand{\rhbyLLgbmKThreeAllPct}{$-$1.1}  % lgbm k=3 vs BASE, all; requested by the author after the beyond-base results; not pre-registered
\newcommand{\rhbyLLgbmKThreeAllDm}{$-$4.17}  % DM vs BASE
\newcommand{\rhbyLLgbmKThreeAllP}{\ensuremath{<}0.001}  % p-value
\newcommand{\rhbyLLgbmKThreeAllDmFull}{2.97}  % DM vs FULL (positive: FULL lower)
\newcommand{\rhbyLLgbmKThreeAllPFull}{0.003}  % p-value
\newcommand{\rhbyLLgbmKThreeAllR}{0.72}  % recovery (L_BASE - L_k) / (L_BASE - L_FULL)
\newcommand{\rhbyLLgbmKFiveAllPct}{$-$1.2}  % lgbm k=5 vs BASE, all; requested by the author after the beyond-base results; not pre-registered
\newcommand{\rhbyLLgbmKFiveAllDm}{$-$4.52}  % DM vs BASE
\newcommand{\rhbyLLgbmKFiveAllP}{\ensuremath{<}0.001}  % p-value
\newcommand{\rhbyLLgbmKFiveAllDmFull}{2.57}  % DM vs FULL (positive: FULL lower)
\newcommand{\rhbyLLgbmKFiveAllPFull}{0.010}  % p-value
\newcommand{\rhbyLLgbmKFiveAllR}{0.76}  % recovery (L_BASE - L_k) / (L_BASE - L_FULL)
\newcommand{\rhbyLLgbmKTenAllPct}{$-$1.1}  % lgbm k=10 vs BASE, all; requested by the author after the beyond-base results; not pre-registered
\newcommand{\rhbyLLgbmKTenAllDm}{$-$4.20}  % DM vs BASE
\newcommand{\rhbyLLgbmKTenAllP}{\ensuremath{<}0.001}  % p-value
\newcommand{\rhbyLLgbmKTenAllDmFull}{2.78}  % DM vs FULL (positive: FULL lower)
\newcommand{\rhbyLLgbmKTenAllPFull}{0.005}  % p-value
\newcommand{\rhbyLLgbmKTenAllR}{0.72}  % recovery (L_BASE - L_k) / (L_BASE - L_FULL)
\newcommand{\rhbyLLgbmKTwentyAllPct}{$-$1.4}  % lgbm k=20 vs BASE, all; requested by the author after the beyond-base results; not pre-registered
\newcommand{\rhbyLLgbmKTwentyAllDm}{$-$5.25}  % DM vs BASE
\newcommand{\rhbyLLgbmKTwentyAllP}{\ensuremath{<}0.001}  % p-value
\newcommand{\rhbyLLgbmKTwentyAllDmFull}{1.28}  % DM vs FULL (positive: FULL lower)
\newcommand{\rhbyLLgbmKTwentyAllPFull}{0.202}  % p-value
\newcommand{\rhbyLLgbmKTwentyAllR}{0.93}  % recovery (L_BASE - L_k) / (L_BASE - L_FULL)
\newcommand{\rhbyLLgbmFullOrdPct}{$-$1.6}  % lgbm FULL vs BASE, ord
\newcommand{\rhbyLLgbmKOneOrdPct}{$-$1.2}  % lgbm k=1 vs BASE, ord; requested by the author after the beyond-base results; not pre-registered
\newcommand{\rhbyLLgbmKOneOrdDm}{$-$5.13}  % DM vs BASE
\newcommand{\rhbyLLgbmKOneOrdP}{\ensuremath{<}0.001}  % p-value
\newcommand{\rhbyLLgbmKOneOrdDmFull}{2.61}  % DM vs FULL (positive: FULL lower)
\newcommand{\rhbyLLgbmKOneOrdPFull}{0.009}  % p-value
\newcommand{\rhbyLLgbmKOneOrdR}{0.75}  % recovery (L_BASE - L_k) / (L_BASE - L_FULL)
\newcommand{\rhbyLLgbmKThreeOrdPct}{$-$1.1}  % lgbm k=3 vs BASE, ord; requested by the author after the beyond-base results; not pre-registered
\newcommand{\rhbyLLgbmKThreeOrdDm}{$-$4.54}  % DM vs BASE
\newcommand{\rhbyLLgbmKThreeOrdP}{\ensuremath{<}0.001}  % p-value
\newcommand{\rhbyLLgbmKThreeOrdDmFull}{3.13}  % DM vs FULL (positive: FULL lower)
\newcommand{\rhbyLLgbmKThreeOrdPFull}{0.002}  % p-value
\newcommand{\rhbyLLgbmKThreeOrdR}{0.70}  % recovery (L_BASE - L_k) / (L_BASE - L_FULL)
\newcommand{\rhbyLLgbmKFiveOrdPct}{$-$1.2}  % lgbm k=5 vs BASE, ord; requested by the author after the beyond-base results; not pre-registered
\newcommand{\rhbyLLgbmKFiveOrdDm}{$-$4.99}  % DM vs BASE
\newcommand{\rhbyLLgbmKFiveOrdP}{\ensuremath{<}0.001}  % p-value
\newcommand{\rhbyLLgbmKFiveOrdDmFull}{3.03}  % DM vs FULL (positive: FULL lower)
\newcommand{\rhbyLLgbmKFiveOrdPFull}{0.002}  % p-value
\newcommand{\rhbyLLgbmKFiveOrdR}{0.72}  % recovery (L_BASE - L_k) / (L_BASE - L_FULL)
\newcommand{\rhbyLLgbmKTenOrdPct}{$-$1.1}  % lgbm k=10 vs BASE, ord; requested by the author after the beyond-base results; not pre-registered
\newcommand{\rhbyLLgbmKTenOrdDm}{$-$4.68}  % DM vs BASE
\newcommand{\rhbyLLgbmKTenOrdP}{\ensuremath{<}0.001}  % p-value
\newcommand{\rhbyLLgbmKTenOrdDmFull}{2.98}  % DM vs FULL (positive: FULL lower)
\newcommand{\rhbyLLgbmKTenOrdPFull}{0.003}  % p-value
\newcommand{\rhbyLLgbmKTenOrdR}{0.70}  % recovery (L_BASE - L_k) / (L_BASE - L_FULL)
\newcommand{\rhbyLLgbmKTwentyOrdPct}{$-$1.5}  % lgbm k=20 vs BASE, ord; requested by the author after the beyond-base results; not pre-registered
\newcommand{\rhbyLLgbmKTwentyOrdDm}{$-$5.82}  % DM vs BASE
\newcommand{\rhbyLLgbmKTwentyOrdP}{\ensuremath{<}0.001}  % p-value
\newcommand{\rhbyLLgbmKTwentyOrdDmFull}{1.03}  % DM vs FULL (positive: FULL lower)
\newcommand{\rhbyLLgbmKTwentyOrdPFull}{0.303}  % p-value
\newcommand{\rhbyLLgbmKTwentyOrdR}{0.94}  % recovery (L_BASE - L_k) / (L_BASE - L_FULL)
\newcommand{\rhbyLLgbmFullEvPct}{$-$0.6}  % lgbm FULL vs BASE, ev
\newcommand{\rhbyLLgbmKOneEvPct}{$-$0.3}  % lgbm k=1 vs BASE, ev; requested by the author after the beyond-base results; not pre-registered
\newcommand{\rhbyLLgbmKOneEvDm}{$-$0.26}  % DM vs BASE
\newcommand{\rhbyLLgbmKOneEvP}{0.795}  % p-value
\newcommand{\rhbyLLgbmKOneEvDmFull}{0.54}  % DM vs FULL (positive: FULL lower)
\newcommand{\rhbyLLgbmKOneEvPFull}{0.592}  % p-value
\newcommand{\rhbyLLgbmKOneEvR}{0.44}  % recovery (L_BASE - L_k) / (L_BASE - L_FULL)
\newcommand{\rhbyLLgbmKThreeEvPct}{$-$0.8}  % lgbm k=3 vs BASE, ev; requested by the author after the beyond-base results; not pre-registered
\newcommand{\rhbyLLgbmKThreeEvDm}{$-$0.63}  % DM vs BASE
\newcommand{\rhbyLLgbmKThreeEvP}{0.526}  % p-value
\newcommand{\rhbyLLgbmKThreeEvDmFull}{$-$0.38}  % DM vs FULL (positive: FULL lower)
\newcommand{\rhbyLLgbmKThreeEvPFull}{0.701}  % p-value
\newcommand{\rhbyLLgbmKThreeEvR}{1.36}  % recovery (L_BASE - L_k) / (L_BASE - L_FULL)
\newcommand{\rhbyLLgbmKFiveEvPct}{$-$1.6}  % lgbm k=5 vs BASE, ev; requested by the author after the beyond-base results; not pre-registered
\newcommand{\rhbyLLgbmKFiveEvDm}{$-$1.03}  % DM vs BASE
\newcommand{\rhbyLLgbmKFiveEvP}{0.305}  % p-value
\newcommand{\rhbyLLgbmKFiveEvDmFull}{$-$1.72}  % DM vs FULL (positive: FULL lower)
\newcommand{\rhbyLLgbmKFiveEvPFull}{0.086}  % p-value
\newcommand{\rhbyLLgbmKFiveEvR}{2.62}  % recovery (L_BASE - L_k) / (L_BASE - L_FULL)
\newcommand{\rhbyLLgbmKTenEvPct}{$-$1.0}  % lgbm k=10 vs BASE, ev; requested by the author after the beyond-base results; not pre-registered
\newcommand{\rhbyLLgbmKTenEvDm}{$-$0.71}  % DM vs BASE
\newcommand{\rhbyLLgbmKTenEvP}{0.477}  % p-value
\newcommand{\rhbyLLgbmKTenEvDmFull}{$-$1.05}  % DM vs FULL (positive: FULL lower)
\newcommand{\rhbyLLgbmKTenEvPFull}{0.294}  % p-value
\newcommand{\rhbyLLgbmKTenEvR}{1.66}  % recovery (L_BASE - L_k) / (L_BASE - L_FULL)
\newcommand{\rhbyLLgbmKTwentyEvPct}{$-$0.3}  % lgbm k=20 vs BASE, ev; requested by the author after the beyond-base results; not pre-registered
\newcommand{\rhbyLLgbmKTwentyEvDm}{$-$0.18}  % DM vs BASE
\newcommand{\rhbyLLgbmKTwentyEvP}{0.855}  % p-value
\newcommand{\rhbyLLgbmKTwentyEvDmFull}{1.20}  % DM vs FULL (positive: FULL lower)
\newcommand{\rhbyLLgbmKTwentyEvPFull}{0.230}  % p-value
\newcommand{\rhbyLLgbmKTwentyEvR}{0.40}  % recovery (L_BASE - L_k) / (L_BASE - L_FULL)
\newcommand{\rhbyLLgbmReading}{intermediate}  % section 2 reading, ORDINARY days
\newcommand{\rhbyLLgbmTop}{sumabsret\_vwstock, sumabsret\_ewstock, effspread\_vwstock, vix, sumret3\_vwstock}  % the model's top 5 block series
\newcommand{\rhbyLXgbFullAllPct}{$-$1.8}  % xgb FULL vs BASE, all
\newcommand{\rhbyLXgbKOneAllPct}{$-$0.6}  % xgb k=1 vs BASE, all; requested by the author after the beyond-base results; not pre-registered
\newcommand{\rhbyLXgbKOneAllDm}{$-$5.62}  % DM vs BASE
\newcommand{\rhbyLXgbKOneAllP}{\ensuremath{<}0.001}  % p-value
\newcommand{\rhbyLXgbKOneAllDmFull}{4.04}  % DM vs FULL (positive: FULL lower)
\newcommand{\rhbyLXgbKOneAllPFull}{\ensuremath{<}0.001}  % p-value
\newcommand{\rhbyLXgbKOneAllR}{0.33}  % recovery (L_BASE - L_k) / (L_BASE - L_FULL)
\newcommand{\rhbyLXgbKThreeAllPct}{$-$1.4}  % xgb k=3 vs BASE, all; requested by the author after the beyond-base results; not pre-registered
\newcommand{\rhbyLXgbKThreeAllDm}{$-$4.92}  % DM vs BASE
\newcommand{\rhbyLXgbKThreeAllP}{\ensuremath{<}0.001}  % p-value
\newcommand{\rhbyLXgbKThreeAllDmFull}{2.59}  % DM vs FULL (positive: FULL lower)
\newcommand{\rhbyLXgbKThreeAllPFull}{0.010}  % p-value
\newcommand{\rhbyLXgbKThreeAllR}{0.77}  % recovery (L_BASE - L_k) / (L_BASE - L_FULL)
\newcommand{\rhbyLXgbKFiveAllPct}{$-$1.6}  % xgb k=5 vs BASE, all; requested by the author after the beyond-base results; not pre-registered
\newcommand{\rhbyLXgbKFiveAllDm}{$-$5.41}  % DM vs BASE
\newcommand{\rhbyLXgbKFiveAllP}{\ensuremath{<}0.001}  % p-value
\newcommand{\rhbyLXgbKFiveAllDmFull}{1.53}  % DM vs FULL (positive: FULL lower)
\newcommand{\rhbyLXgbKFiveAllPFull}{0.127}  % p-value
\newcommand{\rhbyLXgbKFiveAllR}{0.90}  % recovery (L_BASE - L_k) / (L_BASE - L_FULL)
\newcommand{\rhbyLXgbKTenAllPct}{$-$1.4}  % xgb k=10 vs BASE, all; requested by the author after the beyond-base results; not pre-registered
\newcommand{\rhbyLXgbKTenAllDm}{$-$4.81}  % DM vs BASE
\newcommand{\rhbyLXgbKTenAllP}{\ensuremath{<}0.001}  % p-value
\newcommand{\rhbyLXgbKTenAllDmFull}{3.28}  % DM vs FULL (positive: FULL lower)
\newcommand{\rhbyLXgbKTenAllPFull}{0.001}  % p-value
\newcommand{\rhbyLXgbKTenAllR}{0.77}  % recovery (L_BASE - L_k) / (L_BASE - L_FULL)
\newcommand{\rhbyLXgbKTwentyAllPct}{$-$1.6}  % xgb k=20 vs BASE, all; requested by the author after the beyond-base results; not pre-registered
\newcommand{\rhbyLXgbKTwentyAllDm}{$-$4.87}  % DM vs BASE
\newcommand{\rhbyLXgbKTwentyAllP}{\ensuremath{<}0.001}  % p-value
\newcommand{\rhbyLXgbKTwentyAllDmFull}{1.97}  % DM vs FULL (positive: FULL lower)
\newcommand{\rhbyLXgbKTwentyAllPFull}{0.049}  % p-value
\newcommand{\rhbyLXgbKTwentyAllR}{0.90}  % recovery (L_BASE - L_k) / (L_BASE - L_FULL)
\newcommand{\rhbyLXgbFullOrdPct}{$-$1.8}  % xgb FULL vs BASE, ord
\newcommand{\rhbyLXgbKOneOrdPct}{$-$0.6}  % xgb k=1 vs BASE, ord; requested by the author after the beyond-base results; not pre-registered
\newcommand{\rhbyLXgbKOneOrdDm}{$-$5.89}  % DM vs BASE
\newcommand{\rhbyLXgbKOneOrdP}{\ensuremath{<}0.001}  % p-value
\newcommand{\rhbyLXgbKOneOrdDmFull}{4.42}  % DM vs FULL (positive: FULL lower)
\newcommand{\rhbyLXgbKOneOrdPFull}{\ensuremath{<}0.001}  % p-value
\newcommand{\rhbyLXgbKOneOrdR}{0.33}  % recovery (L_BASE - L_k) / (L_BASE - L_FULL)
\newcommand{\rhbyLXgbKThreeOrdPct}{$-$1.5}  % xgb k=3 vs BASE, ord; requested by the author after the beyond-base results; not pre-registered
\newcommand{\rhbyLXgbKThreeOrdDm}{$-$5.55}  % DM vs BASE
\newcommand{\rhbyLXgbKThreeOrdP}{\ensuremath{<}0.001}  % p-value
\newcommand{\rhbyLXgbKThreeOrdDmFull}{2.32}  % DM vs FULL (positive: FULL lower)
\newcommand{\rhbyLXgbKThreeOrdPFull}{0.020}  % p-value
\newcommand{\rhbyLXgbKThreeOrdR}{0.81}  % recovery (L_BASE - L_k) / (L_BASE - L_FULL)
\newcommand{\rhbyLXgbKFiveOrdPct}{$-$1.6}  % xgb k=5 vs BASE, ord; requested by the author after the beyond-base results; not pre-registered
\newcommand{\rhbyLXgbKFiveOrdDm}{$-$5.95}  % DM vs BASE
\newcommand{\rhbyLXgbKFiveOrdP}{\ensuremath{<}0.001}  % p-value
\newcommand{\rhbyLXgbKFiveOrdDmFull}{1.43}  % DM vs FULL (positive: FULL lower)
\newcommand{\rhbyLXgbKFiveOrdPFull}{0.152}  % p-value
\newcommand{\rhbyLXgbKFiveOrdR}{0.91}  % recovery (L_BASE - L_k) / (L_BASE - L_FULL)
\newcommand{\rhbyLXgbKTenOrdPct}{$-$1.4}  % xgb k=10 vs BASE, ord; requested by the author after the beyond-base results; not pre-registered
\newcommand{\rhbyLXgbKTenOrdDm}{$-$5.42}  % DM vs BASE
\newcommand{\rhbyLXgbKTenOrdP}{\ensuremath{<}0.001}  % p-value
\newcommand{\rhbyLXgbKTenOrdDmFull}{3.18}  % DM vs FULL (positive: FULL lower)
\newcommand{\rhbyLXgbKTenOrdPFull}{0.001}  % p-value
\newcommand{\rhbyLXgbKTenOrdR}{0.77}  % recovery (L_BASE - L_k) / (L_BASE - L_FULL)
\newcommand{\rhbyLXgbKTwentyOrdPct}{$-$1.7}  % xgb k=20 vs BASE, ord; requested by the author after the beyond-base results; not pre-registered
\newcommand{\rhbyLXgbKTwentyOrdDm}{$-$5.48}  % DM vs BASE
\newcommand{\rhbyLXgbKTwentyOrdP}{\ensuremath{<}0.001}  % p-value
\newcommand{\rhbyLXgbKTwentyOrdDmFull}{1.60}  % DM vs FULL (positive: FULL lower)
\newcommand{\rhbyLXgbKTwentyOrdPFull}{0.110}  % p-value
\newcommand{\rhbyLXgbKTwentyOrdR}{0.92}  % recovery (L_BASE - L_k) / (L_BASE - L_FULL)
\newcommand{\rhbyLXgbFullEvPct}{$-$1.3}  % xgb FULL vs BASE, ev
\newcommand{\rhbyLXgbKOneEvPct}{$-$0.5}  % xgb k=1 vs BASE, ev; requested by the author after the beyond-base results; not pre-registered
\newcommand{\rhbyLXgbKOneEvDm}{$-$0.63}  % DM vs BASE
\newcommand{\rhbyLXgbKOneEvP}{0.530}  % p-value
\newcommand{\rhbyLXgbKOneEvDmFull}{0.57}  % DM vs FULL (positive: FULL lower)
\newcommand{\rhbyLXgbKOneEvPFull}{0.568}  % p-value
\newcommand{\rhbyLXgbKOneEvR}{0.35}  % recovery (L_BASE - L_k) / (L_BASE - L_FULL)
\newcommand{\rhbyLXgbKThreeEvPct}{0.0}  % xgb k=3 vs BASE, ev; requested by the author after the beyond-base results; not pre-registered
\newcommand{\rhbyLXgbKThreeEvDm}{0.01}  % DM vs BASE
\newcommand{\rhbyLXgbKThreeEvP}{0.992}  % p-value
\newcommand{\rhbyLXgbKThreeEvDmFull}{1.34}  % DM vs FULL (positive: FULL lower)
\newcommand{\rhbyLXgbKThreeEvPFull}{0.180}  % p-value
\newcommand{\rhbyLXgbKThreeEvR}{$-$0.01}  % recovery (L_BASE - L_k) / (L_BASE - L_FULL)
\newcommand{\rhbyLXgbKFiveEvPct}{$-$0.9}  % xgb k=5 vs BASE, ev; requested by the author after the beyond-base results; not pre-registered
\newcommand{\rhbyLXgbKFiveEvDm}{$-$0.60}  % DM vs BASE
\newcommand{\rhbyLXgbKFiveEvP}{0.549}  % p-value
\newcommand{\rhbyLXgbKFiveEvDmFull}{0.48}  % DM vs FULL (positive: FULL lower)
\newcommand{\rhbyLXgbKFiveEvPFull}{0.632}  % p-value
\newcommand{\rhbyLXgbKFiveEvR}{0.71}  % recovery (L_BASE - L_k) / (L_BASE - L_FULL)
\newcommand{\rhbyLXgbKTenEvPct}{$-$1.0}  % xgb k=10 vs BASE, ev; requested by the author after the beyond-base results; not pre-registered
\newcommand{\rhbyLXgbKTenEvDm}{$-$0.63}  % DM vs BASE
\newcommand{\rhbyLXgbKTenEvP}{0.529}  % p-value
\newcommand{\rhbyLXgbKTenEvDmFull}{0.46}  % DM vs FULL (positive: FULL lower)
\newcommand{\rhbyLXgbKTenEvPFull}{0.648}  % p-value
\newcommand{\rhbyLXgbKTenEvR}{0.79}  % recovery (L_BASE - L_k) / (L_BASE - L_FULL)
\newcommand{\rhbyLXgbKTwentyEvPct}{$-$0.8}  % xgb k=20 vs BASE, ev; requested by the author after the beyond-base results; not pre-registered
\newcommand{\rhbyLXgbKTwentyEvDm}{$-$0.42}  % DM vs BASE
\newcommand{\rhbyLXgbKTwentyEvP}{0.674}  % p-value
\newcommand{\rhbyLXgbKTwentyEvDmFull}{0.79}  % DM vs FULL (positive: FULL lower)
\newcommand{\rhbyLXgbKTwentyEvPFull}{0.427}  % p-value
\newcommand{\rhbyLXgbKTwentyEvR}{0.59}  % recovery (L_BASE - L_k) / (L_BASE - L_FULL)
\newcommand{\rhbyLXgbReading}{concentrated}  % section 2 reading, ORDINARY days
\newcommand{\rhbyLXgbTop}{stocktwits\_sentiment, sumabsret\_vwstock, effspread\_vwstock, sumabsret\_ewstock, vix}  % the model's top 5 block series
\newcommand{\rhbyLAgree}{split}  % the two models' readings
\fi
\ifdefined\rhbyFXgbReading\else% AUTO-GENERATED by writeup/make_rehaul_beyond_factor.py -- do not edit; re-run the script.
% The trees' factor test beyond HAR + RM + FOMC (requested by the author after the beyond-base results; not pre-registered); results/beyond_trees_2026-10-27/factor.
\newcommand{\rhbyFLgbmBaseOrdQlike}{0.1584}  % QLIKE lgbm base ordinary
\newcommand{\rhbyFLgbmBaseOrdDmFull}{5.26}  % DM vs FULL
\newcommand{\rhbyFLgbmBaseOrdPFull}{\ensuremath{<}0.001}  % p vs FULL
\newcommand{\rhbyFLgbmPcaOneOrdQlike}{0.1583}  % QLIKE lgbm pca1 ordinary
\newcommand{\rhbyFLgbmPcaOneOrdPct}{$-$0.04}  % QLIKE change % vs BASE
\newcommand{\rhbyFLgbmPcaOneOrdDm}{$-$0.47}  % DM vs BASE (HLN, daily means, lag 8)
\newcommand{\rhbyFLgbmPcaOneOrdP}{0.641}  % p vs BASE
\newcommand{\rhbyFLgbmPcaOneOrdR}{0.02}  % recovery R
\newcommand{\rhbyFLgbmPcaOneOrdDmFull}{5.26}  % DM vs FULL
\newcommand{\rhbyFLgbmPcaOneOrdPFull}{\ensuremath{<}0.001}  % p vs FULL
\newcommand{\rhbyFLgbmPcaTwoOrdQlike}{0.1583}  % QLIKE lgbm pca2 ordinary
\newcommand{\rhbyFLgbmPcaTwoOrdPct}{$-$0.02}  % QLIKE change % vs BASE
\newcommand{\rhbyFLgbmPcaTwoOrdDm}{$-$0.24}  % DM vs BASE (HLN, daily means, lag 8)
\newcommand{\rhbyFLgbmPcaTwoOrdP}{0.807}  % p vs BASE
\newcommand{\rhbyFLgbmPcaTwoOrdR}{0.02}  % recovery R
\newcommand{\rhbyFLgbmPcaTwoOrdDmFull}{5.53}  % DM vs FULL
\newcommand{\rhbyFLgbmPcaTwoOrdPFull}{\ensuremath{<}0.001}  % p vs FULL
\newcommand{\rhbyFLgbmPcaThreeOrdQlike}{0.1582}  % QLIKE lgbm pca3 ordinary
\newcommand{\rhbyFLgbmPcaThreeOrdPct}{$-$0.13}  % QLIKE change % vs BASE
\newcommand{\rhbyFLgbmPcaThreeOrdDm}{$-$1.13}  % DM vs BASE (HLN, daily means, lag 8)
\newcommand{\rhbyFLgbmPcaThreeOrdP}{0.256}  % p vs BASE
\newcommand{\rhbyFLgbmPcaThreeOrdR}{0.08}  % recovery R
\newcommand{\rhbyFLgbmPcaThreeOrdDmFull}{5.51}  % DM vs FULL
\newcommand{\rhbyFLgbmPcaThreeOrdPFull}{\ensuremath{<}0.001}  % p vs FULL
\newcommand{\rhbyFLgbmPcaFiveOrdQlike}{0.1581}  % QLIKE lgbm pca5 ordinary
\newcommand{\rhbyFLgbmPcaFiveOrdPct}{$-$0.16}  % QLIKE change % vs BASE
\newcommand{\rhbyFLgbmPcaFiveOrdDm}{$-$1.13}  % DM vs BASE (HLN, daily means, lag 8)
\newcommand{\rhbyFLgbmPcaFiveOrdP}{0.257}  % p vs BASE
\newcommand{\rhbyFLgbmPcaFiveOrdR}{0.10}  % recovery R
\newcommand{\rhbyFLgbmPcaFiveOrdDmFull}{5.52}  % DM vs FULL
\newcommand{\rhbyFLgbmPcaFiveOrdPFull}{\ensuremath{<}0.001}  % p vs FULL
\newcommand{\rhbyFLgbmPcaTenOrdQlike}{0.1575}  % QLIKE lgbm pca10 ordinary
\newcommand{\rhbyFLgbmPcaTenOrdPct}{$-$0.57}  % QLIKE change % vs BASE
\newcommand{\rhbyFLgbmPcaTenOrdDm}{$-$2.94}  % DM vs BASE (HLN, daily means, lag 8)
\newcommand{\rhbyFLgbmPcaTenOrdP}{0.003}  % p vs BASE
\newcommand{\rhbyFLgbmPcaTenOrdR}{0.35}  % recovery R
\newcommand{\rhbyFLgbmPcaTenOrdDmFull}{4.66}  % DM vs FULL
\newcommand{\rhbyFLgbmPcaTenOrdPFull}{\ensuremath{<}0.001}  % p vs FULL
\newcommand{\rhbyFLgbmPlsOneOrdQlike}{0.1581}  % QLIKE lgbm pls1 ordinary
\newcommand{\rhbyFLgbmPlsOneOrdPct}{$-$0.19}  % QLIKE change % vs BASE
\newcommand{\rhbyFLgbmPlsOneOrdDm}{$-$1.67}  % DM vs BASE (HLN, daily means, lag 8)
\newcommand{\rhbyFLgbmPlsOneOrdP}{0.096}  % p vs BASE
\newcommand{\rhbyFLgbmPlsOneOrdR}{0.12}  % recovery R
\newcommand{\rhbyFLgbmPlsOneOrdDmFull}{5.31}  % DM vs FULL
\newcommand{\rhbyFLgbmPlsOneOrdPFull}{\ensuremath{<}0.001}  % p vs FULL
\newcommand{\rhbyFLgbmPlsThreeOrdQlike}{0.1575}  % QLIKE lgbm pls3 ordinary
\newcommand{\rhbyFLgbmPlsThreeOrdPct}{$-$0.52}  % QLIKE change % vs BASE
\newcommand{\rhbyFLgbmPlsThreeOrdDm}{$-$2.95}  % DM vs BASE (HLN, daily means, lag 8)
\newcommand{\rhbyFLgbmPlsThreeOrdP}{0.003}  % p vs BASE
\newcommand{\rhbyFLgbmPlsThreeOrdR}{0.32}  % recovery R
\newcommand{\rhbyFLgbmPlsThreeOrdDmFull}{4.30}  % DM vs FULL
\newcommand{\rhbyFLgbmPlsThreeOrdPFull}{\ensuremath{<}0.001}  % p vs FULL
\newcommand{\rhbyFLgbmFullOrdQlike}{0.1558}  % QLIKE lgbm full ordinary
\newcommand{\rhbyFLgbmFullOrdPct}{$-$1.62}  % QLIKE change % vs BASE
\newcommand{\rhbyFLgbmFullOrdDm}{$-$5.26}  % DM vs BASE (HLN, daily means, lag 8)
\newcommand{\rhbyFLgbmFullOrdP}{\ensuremath{<}0.001}  % p vs BASE
\newcommand{\rhbyFLgbmFullOrdR}{1.00}  % recovery R
\newcommand{\rhbyFLgbmBaseAllQlike}{0.1636}  % QLIKE lgbm base all
\newcommand{\rhbyFLgbmBaseAllDmFull}{4.92}  % DM vs FULL
\newcommand{\rhbyFLgbmBaseAllPFull}{\ensuremath{<}0.001}  % p vs FULL
\newcommand{\rhbyFLgbmPcaOneAllQlike}{0.1636}  % QLIKE lgbm pca1 all
\newcommand{\rhbyFLgbmPcaOneAllPct}{0.00}  % QLIKE change % vs BASE
\newcommand{\rhbyFLgbmPcaOneAllDm}{$-$0.01}  % DM vs BASE (HLN, daily means, lag 8)
\newcommand{\rhbyFLgbmPcaOneAllP}{0.992}  % p vs BASE
\newcommand{\rhbyFLgbmPcaOneAllR}{0.00}  % recovery R
\newcommand{\rhbyFLgbmPcaOneAllDmFull}{4.91}  % DM vs FULL
\newcommand{\rhbyFLgbmPcaOneAllPFull}{\ensuremath{<}0.001}  % p vs FULL
\newcommand{\rhbyFLgbmPcaTwoAllQlike}{0.1636}  % QLIKE lgbm pca2 all
\newcommand{\rhbyFLgbmPcaTwoAllPct}{0.00}  % QLIKE change % vs BASE
\newcommand{\rhbyFLgbmPcaTwoAllDm}{$-$0.03}  % DM vs BASE (HLN, daily means, lag 8)
\newcommand{\rhbyFLgbmPcaTwoAllP}{0.974}  % p vs BASE
\newcommand{\rhbyFLgbmPcaTwoAllR}{0.00}  % recovery R
\newcommand{\rhbyFLgbmPcaTwoAllDmFull}{5.17}  % DM vs FULL
\newcommand{\rhbyFLgbmPcaTwoAllPFull}{\ensuremath{<}0.001}  % p vs FULL
\newcommand{\rhbyFLgbmPcaThreeAllQlike}{0.1634}  % QLIKE lgbm pca3 all
\newcommand{\rhbyFLgbmPcaThreeAllPct}{$-$0.10}  % QLIKE change % vs BASE
\newcommand{\rhbyFLgbmPcaThreeAllDm}{$-$0.86}  % DM vs BASE (HLN, daily means, lag 8)
\newcommand{\rhbyFLgbmPcaThreeAllP}{0.391}  % p vs BASE
\newcommand{\rhbyFLgbmPcaThreeAllR}{0.06}  % recovery R
\newcommand{\rhbyFLgbmPcaThreeAllDmFull}{5.26}  % DM vs FULL
\newcommand{\rhbyFLgbmPcaThreeAllPFull}{\ensuremath{<}0.001}  % p vs FULL
\newcommand{\rhbyFLgbmPcaFiveAllQlike}{0.1635}  % QLIKE lgbm pca5 all
\newcommand{\rhbyFLgbmPcaFiveAllPct}{$-$0.09}  % QLIKE change % vs BASE
\newcommand{\rhbyFLgbmPcaFiveAllDm}{$-$0.65}  % DM vs BASE (HLN, daily means, lag 8)
\newcommand{\rhbyFLgbmPcaFiveAllP}{0.515}  % p vs BASE
\newcommand{\rhbyFLgbmPcaFiveAllR}{0.06}  % recovery R
\newcommand{\rhbyFLgbmPcaFiveAllDmFull}{5.14}  % DM vs FULL
\newcommand{\rhbyFLgbmPcaFiveAllPFull}{\ensuremath{<}0.001}  % p vs FULL
\newcommand{\rhbyFLgbmPcaTenAllQlike}{0.1628}  % QLIKE lgbm pca10 all
\newcommand{\rhbyFLgbmPcaTenAllPct}{$-$0.47}  % QLIKE change % vs BASE
\newcommand{\rhbyFLgbmPcaTenAllDm}{$-$2.58}  % DM vs BASE (HLN, daily means, lag 8)
\newcommand{\rhbyFLgbmPcaTenAllP}{0.010}  % p vs BASE
\newcommand{\rhbyFLgbmPcaTenAllR}{0.30}  % recovery R
\newcommand{\rhbyFLgbmPcaTenAllDmFull}{4.46}  % DM vs FULL
\newcommand{\rhbyFLgbmPcaTenAllPFull}{\ensuremath{<}0.001}  % p vs FULL
\newcommand{\rhbyFLgbmPlsOneAllQlike}{0.1634}  % QLIKE lgbm pls1 all
\newcommand{\rhbyFLgbmPlsOneAllPct}{$-$0.15}  % QLIKE change % vs BASE
\newcommand{\rhbyFLgbmPlsOneAllDm}{$-$1.35}  % DM vs BASE (HLN, daily means, lag 8)
\newcommand{\rhbyFLgbmPlsOneAllP}{0.178}  % p vs BASE
\newcommand{\rhbyFLgbmPlsOneAllR}{0.10}  % recovery R
\newcommand{\rhbyFLgbmPlsOneAllDmFull}{4.91}  % DM vs FULL
\newcommand{\rhbyFLgbmPlsOneAllPFull}{\ensuremath{<}0.001}  % p vs FULL
\newcommand{\rhbyFLgbmPlsThreeAllQlike}{0.1628}  % QLIKE lgbm pls3 all
\newcommand{\rhbyFLgbmPlsThreeAllPct}{$-$0.48}  % QLIKE change % vs BASE
\newcommand{\rhbyFLgbmPlsThreeAllDm}{$-$2.86}  % DM vs BASE (HLN, daily means, lag 8)
\newcommand{\rhbyFLgbmPlsThreeAllP}{0.004}  % p vs BASE
\newcommand{\rhbyFLgbmPlsThreeAllR}{0.31}  % recovery R
\newcommand{\rhbyFLgbmPlsThreeAllDmFull}{4.04}  % DM vs FULL
\newcommand{\rhbyFLgbmPlsThreeAllPFull}{\ensuremath{<}0.001}  % p vs FULL
\newcommand{\rhbyFLgbmFullAllQlike}{0.1611}  % QLIKE lgbm full all
\newcommand{\rhbyFLgbmFullAllPct}{$-$1.56}  % QLIKE change % vs BASE
\newcommand{\rhbyFLgbmFullAllDm}{$-$4.92}  % DM vs BASE (HLN, daily means, lag 8)
\newcommand{\rhbyFLgbmFullAllP}{\ensuremath{<}0.001}  % p vs BASE
\newcommand{\rhbyFLgbmFullAllR}{1.00}  % recovery R
\newcommand{\rhbyFLgbmBaseEvQlike}{0.3252}  % QLIKE lgbm base event
\newcommand{\rhbyFLgbmBaseEvDmFull}{0.47}  % DM vs FULL
\newcommand{\rhbyFLgbmBaseEvPFull}{0.642}  % p vs FULL
\newcommand{\rhbyFLgbmPcaOneEvQlike}{0.3271}  % QLIKE lgbm pca1 event
\newcommand{\rhbyFLgbmPcaOneEvPct}{0.58}  % QLIKE change % vs BASE
\newcommand{\rhbyFLgbmPcaOneEvDm}{0.91}  % DM vs BASE (HLN, daily means, lag 8)
\newcommand{\rhbyFLgbmPcaOneEvP}{0.365}  % p vs BASE
\newcommand{\rhbyFLgbmPcaOneEvR}{$-$0.94}  % recovery R
\newcommand{\rhbyFLgbmPcaOneEvDmFull}{0.69}  % DM vs FULL
\newcommand{\rhbyFLgbmPcaOneEvPFull}{0.489}  % p vs FULL
\newcommand{\rhbyFLgbmPcaTwoEvQlike}{0.3263}  % QLIKE lgbm pca2 event
\newcommand{\rhbyFLgbmPcaTwoEvPct}{0.32}  % QLIKE change % vs BASE
\newcommand{\rhbyFLgbmPcaTwoEvDm}{0.77}  % DM vs BASE (HLN, daily means, lag 8)
\newcommand{\rhbyFLgbmPcaTwoEvP}{0.443}  % p vs BASE
\newcommand{\rhbyFLgbmPcaTwoEvR}{$-$0.52}  % recovery R
\newcommand{\rhbyFLgbmPcaTwoEvDmFull}{0.67}  % DM vs FULL
\newcommand{\rhbyFLgbmPcaTwoEvPFull}{0.506}  % p vs FULL
\newcommand{\rhbyFLgbmPcaThreeEvQlike}{0.3267}  % QLIKE lgbm pca3 event
\newcommand{\rhbyFLgbmPcaThreeEvPct}{0.46}  % QLIKE change % vs BASE
\newcommand{\rhbyFLgbmPcaThreeEvDm}{1.00}  % DM vs BASE (HLN, daily means, lag 8)
\newcommand{\rhbyFLgbmPcaThreeEvP}{0.319}  % p vs BASE
\newcommand{\rhbyFLgbmPcaThreeEvR}{$-$0.74}  % recovery R
\newcommand{\rhbyFLgbmPcaThreeEvDmFull}{0.80}  % DM vs FULL
\newcommand{\rhbyFLgbmPcaThreeEvPFull}{0.425}  % p vs FULL
\newcommand{\rhbyFLgbmPcaFiveEvQlike}{0.3286}  % QLIKE lgbm pca5 event
\newcommand{\rhbyFLgbmPcaFiveEvPct}{1.05}  % QLIKE change % vs BASE
\newcommand{\rhbyFLgbmPcaFiveEvDm}{2.06}  % DM vs BASE (HLN, daily means, lag 8)
\newcommand{\rhbyFLgbmPcaFiveEvP}{0.040}  % p vs BASE
\newcommand{\rhbyFLgbmPcaFiveEvR}{$-$1.68}  % recovery R
\newcommand{\rhbyFLgbmPcaFiveEvDmFull}{1.04}  % DM vs FULL
\newcommand{\rhbyFLgbmPcaFiveEvPFull}{0.298}  % p vs FULL
\newcommand{\rhbyFLgbmPcaTenEvQlike}{0.3284}  % QLIKE lgbm pca10 event
\newcommand{\rhbyFLgbmPcaTenEvPct}{0.97}  % QLIKE change % vs BASE
\newcommand{\rhbyFLgbmPcaTenEvDm}{2.23}  % DM vs BASE (HLN, daily means, lag 8)
\newcommand{\rhbyFLgbmPcaTenEvP}{0.026}  % p vs BASE
\newcommand{\rhbyFLgbmPcaTenEvR}{$-$1.56}  % recovery R
\newcommand{\rhbyFLgbmPcaTenEvDmFull}{1.09}  % DM vs FULL
\newcommand{\rhbyFLgbmPcaTenEvPFull}{0.278}  % p vs FULL
\newcommand{\rhbyFLgbmPlsOneEvQlike}{0.3268}  % QLIKE lgbm pls1 event
\newcommand{\rhbyFLgbmPlsOneEvPct}{0.49}  % QLIKE change % vs BASE
\newcommand{\rhbyFLgbmPlsOneEvDm}{0.76}  % DM vs BASE (HLN, daily means, lag 8)
\newcommand{\rhbyFLgbmPlsOneEvP}{0.446}  % p vs BASE
\newcommand{\rhbyFLgbmPlsOneEvR}{$-$0.78}  % recovery R
\newcommand{\rhbyFLgbmPlsOneEvDmFull}{0.67}  % DM vs FULL
\newcommand{\rhbyFLgbmPlsOneEvPFull}{0.505}  % p vs FULL
\newcommand{\rhbyFLgbmPlsThreeEvQlike}{0.3256}  % QLIKE lgbm pls3 event
\newcommand{\rhbyFLgbmPlsThreeEvPct}{0.11}  % QLIKE change % vs BASE
\newcommand{\rhbyFLgbmPlsThreeEvDm}{0.32}  % DM vs BASE (HLN, daily means, lag 8)
\newcommand{\rhbyFLgbmPlsThreeEvP}{0.749}  % p vs BASE
\newcommand{\rhbyFLgbmPlsThreeEvR}{$-$0.18}  % recovery R
\newcommand{\rhbyFLgbmPlsThreeEvDmFull}{0.55}  % DM vs FULL
\newcommand{\rhbyFLgbmPlsThreeEvPFull}{0.583}  % p vs FULL
\newcommand{\rhbyFLgbmFullEvQlike}{0.3232}  % QLIKE lgbm full event
\newcommand{\rhbyFLgbmFullEvPct}{$-$0.62}  % QLIKE change % vs BASE
\newcommand{\rhbyFLgbmFullEvDm}{$-$0.47}  % DM vs BASE (HLN, daily means, lag 8)
\newcommand{\rhbyFLgbmFullEvP}{0.642}  % p vs BASE
\newcommand{\rhbyFLgbmFullEvR}{1.00}  % recovery R
\newcommand{\rhbyFLgbmReading}{dense beyond a few components}  % section 1 reading, ORDINARY days
\newcommand{\rhbyFLgbmRPcaTenOrd}{0.35}  % R of BASE + 10 PCs, ORDINARY
\newcommand{\rhbyFLgbmBestShortOrdR}{0.32}  % best k <= 3 arm (pls3), R
\newcommand{\rhbyFXgbBaseOrdQlike}{0.1590}  % QLIKE xgb base ordinary
\newcommand{\rhbyFXgbBaseOrdDmFull}{5.86}  % DM vs FULL
\newcommand{\rhbyFXgbBaseOrdPFull}{\ensuremath{<}0.001}  % p vs FULL
\newcommand{\rhbyFXgbPcaOneOrdQlike}{0.1587}  % QLIKE xgb pca1 ordinary
\newcommand{\rhbyFXgbPcaOneOrdPct}{$-$0.15}  % QLIKE change % vs BASE
\newcommand{\rhbyFXgbPcaOneOrdDm}{$-$1.96}  % DM vs BASE (HLN, daily means, lag 8)
\newcommand{\rhbyFXgbPcaOneOrdP}{0.049}  % p vs BASE
\newcommand{\rhbyFXgbPcaOneOrdR}{0.08}  % recovery R
\newcommand{\rhbyFXgbPcaOneOrdDmFull}{5.77}  % DM vs FULL
\newcommand{\rhbyFXgbPcaOneOrdPFull}{\ensuremath{<}0.001}  % p vs FULL
\newcommand{\rhbyFXgbPcaTwoOrdQlike}{0.1586}  % QLIKE xgb pca2 ordinary
\newcommand{\rhbyFXgbPcaTwoOrdPct}{$-$0.22}  % QLIKE change % vs BASE
\newcommand{\rhbyFXgbPcaTwoOrdDm}{$-$2.16}  % DM vs BASE (HLN, daily means, lag 8)
\newcommand{\rhbyFXgbPcaTwoOrdP}{0.030}  % p vs BASE
\newcommand{\rhbyFXgbPcaTwoOrdR}{0.12}  % recovery R
\newcommand{\rhbyFXgbPcaTwoOrdDmFull}{5.95}  % DM vs FULL
\newcommand{\rhbyFXgbPcaTwoOrdPFull}{\ensuremath{<}0.001}  % p vs FULL
\newcommand{\rhbyFXgbPcaThreeOrdQlike}{0.1586}  % QLIKE xgb pca3 ordinary
\newcommand{\rhbyFXgbPcaThreeOrdPct}{$-$0.22}  % QLIKE change % vs BASE
\newcommand{\rhbyFXgbPcaThreeOrdDm}{$-$2.17}  % DM vs BASE (HLN, daily means, lag 8)
\newcommand{\rhbyFXgbPcaThreeOrdP}{0.030}  % p vs BASE
\newcommand{\rhbyFXgbPcaThreeOrdR}{0.12}  % recovery R
\newcommand{\rhbyFXgbPcaThreeOrdDmFull}{5.60}  % DM vs FULL
\newcommand{\rhbyFXgbPcaThreeOrdPFull}{\ensuremath{<}0.001}  % p vs FULL
\newcommand{\rhbyFXgbPcaFiveOrdQlike}{0.1585}  % QLIKE xgb pca5 ordinary
\newcommand{\rhbyFXgbPcaFiveOrdPct}{$-$0.33}  % QLIKE change % vs BASE
\newcommand{\rhbyFXgbPcaFiveOrdDm}{$-$2.67}  % DM vs BASE (HLN, daily means, lag 8)
\newcommand{\rhbyFXgbPcaFiveOrdP}{0.007}  % p vs BASE
\newcommand{\rhbyFXgbPcaFiveOrdR}{0.18}  % recovery R
\newcommand{\rhbyFXgbPcaFiveOrdDmFull}{5.81}  % DM vs FULL
\newcommand{\rhbyFXgbPcaFiveOrdPFull}{\ensuremath{<}0.001}  % p vs FULL
\newcommand{\rhbyFXgbPcaTenOrdQlike}{0.1579}  % QLIKE xgb pca10 ordinary
\newcommand{\rhbyFXgbPcaTenOrdPct}{$-$0.67}  % QLIKE change % vs BASE
\newcommand{\rhbyFXgbPcaTenOrdDm}{$-$3.42}  % DM vs BASE (HLN, daily means, lag 8)
\newcommand{\rhbyFXgbPcaTenOrdP}{\ensuremath{<}0.001}  % p vs BASE
\newcommand{\rhbyFXgbPcaTenOrdR}{0.37}  % recovery R
\newcommand{\rhbyFXgbPcaTenOrdDmFull}{4.96}  % DM vs FULL
\newcommand{\rhbyFXgbPcaTenOrdPFull}{\ensuremath{<}0.001}  % p vs FULL
\newcommand{\rhbyFXgbPlsOneOrdQlike}{0.1585}  % QLIKE xgb pls1 ordinary
\newcommand{\rhbyFXgbPlsOneOrdPct}{$-$0.29}  % QLIKE change % vs BASE
\newcommand{\rhbyFXgbPlsOneOrdDm}{$-$2.78}  % DM vs BASE (HLN, daily means, lag 8)
\newcommand{\rhbyFXgbPlsOneOrdP}{0.005}  % p vs BASE
\newcommand{\rhbyFXgbPlsOneOrdR}{0.16}  % recovery R
\newcommand{\rhbyFXgbPlsOneOrdDmFull}{5.71}  % DM vs FULL
\newcommand{\rhbyFXgbPlsOneOrdPFull}{\ensuremath{<}0.001}  % p vs FULL
\newcommand{\rhbyFXgbPlsThreeOrdQlike}{0.1580}  % QLIKE xgb pls3 ordinary
\newcommand{\rhbyFXgbPlsThreeOrdPct}{$-$0.64}  % QLIKE change % vs BASE
\newcommand{\rhbyFXgbPlsThreeOrdDm}{$-$3.56}  % DM vs BASE (HLN, daily means, lag 8)
\newcommand{\rhbyFXgbPlsThreeOrdP}{\ensuremath{<}0.001}  % p vs BASE
\newcommand{\rhbyFXgbPlsThreeOrdR}{0.36}  % recovery R
\newcommand{\rhbyFXgbPlsThreeOrdDmFull}{4.41}  % DM vs FULL
\newcommand{\rhbyFXgbPlsThreeOrdPFull}{\ensuremath{<}0.001}  % p vs FULL
\newcommand{\rhbyFXgbFullOrdQlike}{0.1561}  % QLIKE xgb full ordinary
\newcommand{\rhbyFXgbFullOrdPct}{$-$1.80}  % QLIKE change % vs BASE
\newcommand{\rhbyFXgbFullOrdDm}{$-$5.86}  % DM vs BASE (HLN, daily means, lag 8)
\newcommand{\rhbyFXgbFullOrdP}{\ensuremath{<}0.001}  % p vs BASE
\newcommand{\rhbyFXgbFullOrdR}{1.00}  % recovery R
\newcommand{\rhbyFXgbBaseAllQlike}{0.1641}  % QLIKE xgb base all
\newcommand{\rhbyFXgbBaseAllDmFull}{5.26}  % DM vs FULL
\newcommand{\rhbyFXgbBaseAllPFull}{\ensuremath{<}0.001}  % p vs FULL
\newcommand{\rhbyFXgbPcaOneAllQlike}{0.1638}  % QLIKE xgb pca1 all
\newcommand{\rhbyFXgbPcaOneAllPct}{$-$0.19}  % QLIKE change % vs BASE
\newcommand{\rhbyFXgbPcaOneAllDm}{$-$2.54}  % DM vs BASE (HLN, daily means, lag 8)
\newcommand{\rhbyFXgbPcaOneAllP}{0.011}  % p vs BASE
\newcommand{\rhbyFXgbPcaOneAllR}{0.11}  % recovery R
\newcommand{\rhbyFXgbPcaOneAllDmFull}{4.97}  % DM vs FULL
\newcommand{\rhbyFXgbPcaOneAllPFull}{\ensuremath{<}0.001}  % p vs FULL
\newcommand{\rhbyFXgbPcaTwoAllQlike}{0.1639}  % QLIKE xgb pca2 all
\newcommand{\rhbyFXgbPcaTwoAllPct}{$-$0.15}  % QLIKE change % vs BASE
\newcommand{\rhbyFXgbPcaTwoAllDm}{$-$1.55}  % DM vs BASE (HLN, daily means, lag 8)
\newcommand{\rhbyFXgbPcaTwoAllP}{0.121}  % p vs BASE
\newcommand{\rhbyFXgbPcaTwoAllR}{0.09}  % recovery R
\newcommand{\rhbyFXgbPcaTwoAllDmFull}{5.53}  % DM vs FULL
\newcommand{\rhbyFXgbPcaTwoAllPFull}{\ensuremath{<}0.001}  % p vs FULL
\newcommand{\rhbyFXgbPcaThreeAllQlike}{0.1638}  % QLIKE xgb pca3 all
\newcommand{\rhbyFXgbPcaThreeAllPct}{$-$0.19}  % QLIKE change % vs BASE
\newcommand{\rhbyFXgbPcaThreeAllDm}{$-$1.95}  % DM vs BASE (HLN, daily means, lag 8)
\newcommand{\rhbyFXgbPcaThreeAllP}{0.051}  % p vs BASE
\newcommand{\rhbyFXgbPcaThreeAllR}{0.11}  % recovery R
\newcommand{\rhbyFXgbPcaThreeAllDmFull}{5.19}  % DM vs FULL
\newcommand{\rhbyFXgbPcaThreeAllPFull}{\ensuremath{<}0.001}  % p vs FULL
\newcommand{\rhbyFXgbPcaFiveAllQlike}{0.1636}  % QLIKE xgb pca5 all
\newcommand{\rhbyFXgbPcaFiveAllPct}{$-$0.29}  % QLIKE change % vs BASE
\newcommand{\rhbyFXgbPcaFiveAllDm}{$-$2.40}  % DM vs BASE (HLN, daily means, lag 8)
\newcommand{\rhbyFXgbPcaFiveAllP}{0.016}  % p vs BASE
\newcommand{\rhbyFXgbPcaFiveAllR}{0.17}  % recovery R
\newcommand{\rhbyFXgbPcaFiveAllDmFull}{5.27}  % DM vs FULL
\newcommand{\rhbyFXgbPcaFiveAllPFull}{\ensuremath{<}0.001}  % p vs FULL
\newcommand{\rhbyFXgbPcaTenAllQlike}{0.1631}  % QLIKE xgb pca10 all
\newcommand{\rhbyFXgbPcaTenAllPct}{$-$0.63}  % QLIKE change % vs BASE
\newcommand{\rhbyFXgbPcaTenAllDm}{$-$3.29}  % DM vs BASE (HLN, daily means, lag 8)
\newcommand{\rhbyFXgbPcaTenAllP}{0.001}  % p vs BASE
\newcommand{\rhbyFXgbPcaTenAllR}{0.36}  % recovery R
\newcommand{\rhbyFXgbPcaTenAllDmFull}{4.51}  % DM vs FULL
\newcommand{\rhbyFXgbPcaTenAllPFull}{\ensuremath{<}0.001}  % p vs FULL
\newcommand{\rhbyFXgbPlsOneAllQlike}{0.1637}  % QLIKE xgb pls1 all
\newcommand{\rhbyFXgbPlsOneAllPct}{$-$0.28}  % QLIKE change % vs BASE
\newcommand{\rhbyFXgbPlsOneAllDm}{$-$2.65}  % DM vs BASE (HLN, daily means, lag 8)
\newcommand{\rhbyFXgbPlsOneAllP}{0.008}  % p vs BASE
\newcommand{\rhbyFXgbPlsOneAllR}{0.16}  % recovery R
\newcommand{\rhbyFXgbPlsOneAllDmFull}{5.13}  % DM vs FULL
\newcommand{\rhbyFXgbPlsOneAllPFull}{\ensuremath{<}0.001}  % p vs FULL
\newcommand{\rhbyFXgbPlsThreeAllQlike}{0.1632}  % QLIKE xgb pls3 all
\newcommand{\rhbyFXgbPlsThreeAllPct}{$-$0.59}  % QLIKE change % vs BASE
\newcommand{\rhbyFXgbPlsThreeAllDm}{$-$3.21}  % DM vs BASE (HLN, daily means, lag 8)
\newcommand{\rhbyFXgbPlsThreeAllP}{0.001}  % p vs BASE
\newcommand{\rhbyFXgbPlsThreeAllR}{0.33}  % recovery R
\newcommand{\rhbyFXgbPlsThreeAllDmFull}{4.27}  % DM vs FULL
\newcommand{\rhbyFXgbPlsThreeAllPFull}{\ensuremath{<}0.001}  % p vs FULL
\newcommand{\rhbyFXgbFullAllQlike}{0.1612}  % QLIKE xgb full all
\newcommand{\rhbyFXgbFullAllPct}{$-$1.77}  % QLIKE change % vs BASE
\newcommand{\rhbyFXgbFullAllDm}{$-$5.26}  % DM vs BASE (HLN, daily means, lag 8)
\newcommand{\rhbyFXgbFullAllP}{\ensuremath{<}0.001}  % p vs BASE
\newcommand{\rhbyFXgbFullAllR}{1.00}  % recovery R
\newcommand{\rhbyFXgbBaseEvQlike}{0.3229}  % QLIKE xgb base event
\newcommand{\rhbyFXgbBaseEvDmFull}{0.72}  % DM vs FULL
\newcommand{\rhbyFXgbBaseEvPFull}{0.470}  % p vs FULL
\newcommand{\rhbyFXgbPcaOneEvQlike}{0.3204}  % QLIKE xgb pca1 event
\newcommand{\rhbyFXgbPcaOneEvPct}{$-$0.77}  % QLIKE change % vs BASE
\newcommand{\rhbyFXgbPcaOneEvDm}{$-$2.02}  % DM vs BASE (HLN, daily means, lag 8)
\newcommand{\rhbyFXgbPcaOneEvP}{0.044}  % p vs BASE
\newcommand{\rhbyFXgbPcaOneEvR}{0.59}  % recovery R
\newcommand{\rhbyFXgbPcaOneEvDmFull}{0.26}  % DM vs FULL
\newcommand{\rhbyFXgbPcaOneEvPFull}{0.793}  % p vs FULL
\newcommand{\rhbyFXgbPcaTwoEvQlike}{0.3255}  % QLIKE xgb pca2 event
\newcommand{\rhbyFXgbPcaTwoEvPct}{0.80}  % QLIKE change % vs BASE
\newcommand{\rhbyFXgbPcaTwoEvDm}{1.73}  % DM vs BASE (HLN, daily means, lag 8)
\newcommand{\rhbyFXgbPcaTwoEvP}{0.083}  % p vs BASE
\newcommand{\rhbyFXgbPcaTwoEvR}{$-$0.61}  % recovery R
\newcommand{\rhbyFXgbPcaTwoEvDmFull}{1.36}  % DM vs FULL
\newcommand{\rhbyFXgbPcaTwoEvPFull}{0.173}  % p vs FULL
\newcommand{\rhbyFXgbPcaThreeEvQlike}{0.3234}  % QLIKE xgb pca3 event
\newcommand{\rhbyFXgbPcaThreeEvPct}{0.17}  % QLIKE change % vs BASE
\newcommand{\rhbyFXgbPcaThreeEvDm}{0.23}  % DM vs BASE (HLN, daily means, lag 8)
\newcommand{\rhbyFXgbPcaThreeEvP}{0.817}  % p vs BASE
\newcommand{\rhbyFXgbPcaThreeEvR}{$-$0.13}  % recovery R
\newcommand{\rhbyFXgbPcaThreeEvDmFull}{0.98}  % DM vs FULL
\newcommand{\rhbyFXgbPcaThreeEvPFull}{0.326}  % p vs FULL
\newcommand{\rhbyFXgbPcaFiveEvQlike}{0.3236}  % QLIKE xgb pca5 event
\newcommand{\rhbyFXgbPcaFiveEvPct}{0.24}  % QLIKE change % vs BASE
\newcommand{\rhbyFXgbPcaFiveEvDm}{0.27}  % DM vs BASE (HLN, daily means, lag 8)
\newcommand{\rhbyFXgbPcaFiveEvP}{0.785}  % p vs BASE
\newcommand{\rhbyFXgbPcaFiveEvR}{$-$0.18}  % recovery R
\newcommand{\rhbyFXgbPcaFiveEvDmFull}{0.95}  % DM vs FULL
\newcommand{\rhbyFXgbPcaFiveEvPFull}{0.344}  % p vs FULL
\newcommand{\rhbyFXgbPcaTenEvQlike}{0.3230}  % QLIKE xgb pca10 event
\newcommand{\rhbyFXgbPcaTenEvPct}{0.04}  % QLIKE change % vs BASE
\newcommand{\rhbyFXgbPcaTenEvDm}{0.04}  % DM vs BASE (HLN, daily means, lag 8)
\newcommand{\rhbyFXgbPcaTenEvP}{0.965}  % p vs BASE
\newcommand{\rhbyFXgbPcaTenEvR}{$-$0.03}  % recovery R
\newcommand{\rhbyFXgbPcaTenEvDmFull}{0.98}  % DM vs FULL
\newcommand{\rhbyFXgbPcaTenEvPFull}{0.325}  % p vs FULL
\newcommand{\rhbyFXgbPlsOneEvQlike}{0.3223}  % QLIKE xgb pls1 event
\newcommand{\rhbyFXgbPlsOneEvPct}{$-$0.16}  % QLIKE change % vs BASE
\newcommand{\rhbyFXgbPlsOneEvDm}{$-$0.37}  % DM vs BASE (HLN, daily means, lag 8)
\newcommand{\rhbyFXgbPlsOneEvP}{0.711}  % p vs BASE
\newcommand{\rhbyFXgbPlsOneEvR}{0.12}  % recovery R
\newcommand{\rhbyFXgbPlsOneEvDmFull}{0.65}  % DM vs FULL
\newcommand{\rhbyFXgbPlsOneEvPFull}{0.517}  % p vs FULL
\newcommand{\rhbyFXgbPlsThreeEvQlike}{0.3238}  % QLIKE xgb pls3 event
\newcommand{\rhbyFXgbPlsThreeEvPct}{0.28}  % QLIKE change % vs BASE
\newcommand{\rhbyFXgbPlsThreeEvDm}{0.24}  % DM vs BASE (HLN, daily means, lag 8)
\newcommand{\rhbyFXgbPlsThreeEvP}{0.808}  % p vs BASE
\newcommand{\rhbyFXgbPlsThreeEvR}{$-$0.21}  % recovery R
\newcommand{\rhbyFXgbPlsThreeEvDmFull}{1.06}  % DM vs FULL
\newcommand{\rhbyFXgbPlsThreeEvPFull}{0.289}  % p vs FULL
\newcommand{\rhbyFXgbFullEvQlike}{0.3186}  % QLIKE xgb full event
\newcommand{\rhbyFXgbFullEvPct}{$-$1.31}  % QLIKE change % vs BASE
\newcommand{\rhbyFXgbFullEvDm}{$-$0.72}  % DM vs BASE (HLN, daily means, lag 8)
\newcommand{\rhbyFXgbFullEvP}{0.470}  % p vs BASE
\newcommand{\rhbyFXgbFullEvR}{1.00}  % recovery R
\newcommand{\rhbyFXgbReading}{dense beyond a few components}  % section 1 reading, ORDINARY days
\newcommand{\rhbyFXgbRPcaTenOrd}{0.37}  % R of BASE + 10 PCs, ORDINARY
\newcommand{\rhbyFXgbBestShortOrdR}{0.36}  % best k <= 3 arm (pls3), R
\newcommand{\rhbyFAgreement}{agree}  % the two models' readings
\newcommand{\rhbyFLabel}{requested by the author after the beyond-base results; not pre-registered}
\fi
\begin{itemize}
\item \textbf{Two groups carry the gain.} The FOMC calendar and the futures' realized moments, added to the HAR
  benchmark set, deliver at least \rhabBaseShareMin\% of every class's every-bar gain from the full set. What remains is
  small, only the tree ensembles capture it, and a few series carry it. How much a model uses a series is a poor guide
  to these gains (Section~\ref{sec:rh_fi}).
\item \textbf{Design} (Table~\ref{tab:rh_ablation}; Appendix~\ref{app:rh_results_ablation}). The design of
  Table~\ref{tab:rh_oos}, including the full set's frozen hyperparameters, with only the columns changed:
\item \textbf{Not an augmented HAR.} The futures' realized moments, computed from the same one-minute returns that
  define the target, are the natural augmented HAR. Added to the HAR benchmark set they lower the loss in every class
  (Panel~B; \rhabOlsAddRmAllPct\% for the OLS), by about as much as the three classical extensions estimated together
  (\rhabOlsSharqlAllPct\%; Panel~C). Yet the full set is significantly more accurate than the HAR benchmark set plus
  the realized moments in \rhabNClassesBeatAddRm{} of the \rhabNClassesAddRm{} classes (Panel~D).
\item \textbf{The FOMC calendar and the realized moments} (Panels~A, B and~D).
\item \textbf{What remains} (Panel~D, last row; Appendix~\ref{app:rh_results_beyond}).
\item \textbf{Alternative data on their live windows} (Panel~E). The implied-volatility, StockTwits and
  volatility-demand series start after the data and enter as zero until the rolling window holds \rhorWindowDays{}
  trading days of their history (Section~\ref{sec:rh_predictors}), so over most of the span Panels~A and~B measure
  them through their availability shares. On the bars on which their values are live the gains remain small:
\end{itemize}


\subsection{Robustness}
\label{sec:rh_results_robust}
\begin{itemize}
\item \textbf{Estimation window.} The window length $W$ of \rhorWindowDays{} trading days was chosen after the
  \rhwsNWinWord{} candidate windows from \rhwsWinMin{} to \rhrbWinLong{} trading days had been
  scored out of sample on $\mathcal T_{\mathrm{oos}}$; it is not a pre-sample choice. Appendix~\ref{app:rh_robust_window}
  reports every window.
\item \textbf{Multiple comparisons.} Adjusting for the number of comparisons leaves the every-bar result intact: under
  QLIKE the 90\% model confidence set \citep{HansenLundeNason2011} holds \rhmcsBarQlikeTmaxNNinety{} of the
  \rhmcsNModels{} distinct pooled forecasts (\rhmcsBarQlikeTmaxSetNinety), and StepM \citep{RomanoWolf2005} rejects
  equal accuracy of the full and HAR benchmark sets in \rhmcsBarQAfNRej{} of the \rhmcsNAlt{} classes (SPA
  $p$-value \rhmcsBarQAfSpaC; \citealp{Hansen2005}), against \rhmcsBarMAfNRejWord{} under squared error
  (Tables~\ref{tab:rh_mcs} and~\ref{tab:rh_mcs_spa}).
\end{itemize}









